% ============================================================
%  R. Nielsen, "Grundideernes Logik. II."
%  Kjøbenhavn: Forlagt af den Gyldendalske Boghandel (F. Hegel),
%  Trykt hos J. H. Schultz, 1866.
%  TRANSCRIPTION (Danish) — from the Google Books scan of the British
%  Museum/Library copy.
%
%  Scan: ~/bibliotek/Nielsen, Rasmus/1866-grundideernes-logik-2-bl.pdf (486 PDF pp.)
%        sha256 fa8aaca7d4c5ba9f... (see SCANS.tsv). A second scan, of the
%        Harvard copy (1866-grundideernes-logik-2.pdf, offset +22), crops the
%        right edge of many versos and was abandoned as copy-text.
%  TYPE: Fraktur (Danish body); Greek/Latin technical terms in antiqua.
%  WITNESS for ocrdiff.py: the scan's embedded text layer (USABLE; reads ø/æ).
%
%  OFFSET (text-layer folios, all pages; confirm by eye):
%    * PDF 3 title; PDF 5--7 Indhold; PDF 8 Rettelser.
%    * Roman folios VII--XVI: PDF = roman + 2.
%    * Arabic folios 1--465: PDF = printed + 18 (PDF 19--483), no drift.
%
%  STATUS: transcribed in full 2026-09-22 (front matter III--XVI, pp. 1--465);
%  not yet published. OPEN items below.
%
%  Word-accuracy: every batch diffed against the BL scan's embedded text layer
%  before splicing (ocrdiff.py). 32 batches (front matter + 31 x 15 pp.),
%  799 candidates, 23 transcription errors corrected, 0 unresolved.
%  See TRANSCRIPTION-PLAYBOOK.md §3 step 4. The Indhold and Rettelser were
%  checked word by word at the image instead (the diff reads only folio'd
%  text pages). pp. 1--105 and the front matter were first transcribed from
%  the Harvard scan, whose cropped versos forced conjectured line ends; all of
%  those were then re-read at the BL image (about 420 conjectures confirmed,
%  about 40 corrected, plus words the crop had hidden: p.36 "ikke", "sat",
%  p.40 "da", p.66 "just"), and re-diffed against the BL layer.
%  Emphasis: letterspacing read at the image on every page (spacing.py as a
%  detector). Printer's errors logged at the site as "% PRINTER:" (104).
%  Doubtful readings logged as "% UNSURE:" (24). Quote balance off zero by 11,
%  every instance a logged defect in the print.
%
%  ACCURACY [2026-09-22]: sampled 6 pages at the image (printed 152, 172,
%  309, 329, 386, 390; chosen by REPAIR-PLAYBOOK.md §3's deterministic rule,
%  seed "nielsen/grundideernes-logik-2|v1") by a reader who had not
%  transcribed the book, no OCR consulted: 0 word errors, rate bounded at
%  <=0.5/page (95%). NOT a full collation. The same pass adjudicated 15
%  ocrdiff candidates drawn from elsewhere in the book: 0 of 15 were real
%  transcription errors (all OCR faults of the layer: ellipsis read as "11",
%  the β) heading read as "3", line-duplications, footnote-order artifacts).
%  Observed, not counted, at the sampled pages: the ou/ø class again
%  (p.152 mirakuløs, p.172 afgjørende and fuldførte, p.310 Beviisførelse,
%  p.325 bør, all on weakly inked leaves, with clear ø elsewhere on the same
%  pages); and p.283 "Conseqvens" read from letterform plus the book's habit
%  rather than from a clean impression (p.286 is sharp), which wants a second
%  eye.
%
%  OPEN:
%   * Hebrew (p.55 note): RESOLVED 2026-09-22. Set with cjhebrew, which does
%     pointed Hebrew under plain pdflatex; the build needs that package and no
%     font change. The pointing is still an editorial reading: the BL scan is
%     1-bit and the type has broken up, so the consonants and word order are
%     certain and the individual points are not (see the UNSURE note at the site).
%   * o/ø: early batches logged some faint ø as PRINTER "o" (e.g. pp.318--328
%     dodelig, tankelos, indholdslos, noiagtig); the rule adopted later is that
%     only a positive signal counts. Review these PRINTER notes.
%   * MARK: the \opage{152} marker sits before "kun"; p.152 begins mid-word
%     and the marker belongs at "natur\opage{152}ligviis". Other sampled
%     markers, including the mid-word ones at pp.172/173, are correctly placed.
%   * Heading forms: minor heads (α)/β)/γ) with "Functioner af ...") are set
%     centred on their own line; the αα)/1)/2)/3) heads run in. Only B./a./b.
%     heads have toc entries. A normalisation pass against the Indhold is
%     still owed (TRANSCRIPTION-PLAYBOOK §6: record, don't normalise, any
%     disagreement between Indhold and heads).
%   * Front matter: the unheaded preface (VII--XVI, signed 11 Marts 1866) is
%     set as \chapter*{[Indledning]}; "[Forord]" may be the better editorial
%     title.
%
%  FRONT-MATTER MARKER: the roman front matter is spliced from
%  .parts/pp0-0.texfrag (marker "pp. 0--0"), because splice.py only accepts
%  arabic ranges.
% ============================================================
\documentclass[12pt]{book}
\usepackage[utf8]{inputenc}
\usepackage[T1]{fontenc}
\usepackage[danish]{babel}
\usepackage[protrusion=true,expansion=true]{microtype}
\usepackage[margin=1.4in]{geometry}
\usepackage{setspace}
\usepackage{amsmath}
\usepackage{libertinus}
\usepackage{libertinust1math}
\usepackage{textalpha}   % polytonic Greek typed directly
\usepackage{cjhebrew}    % pointed Hebrew, p. 55 note (pdflatex-safe)
\usepackage{fancyhdr}
\usepackage[colorlinks=true,linkcolor=black,urlcolor=black,
  pdftitle={Grundideernes Logik II},
  pdfauthor={Rasmus Nielsen}]{hyperref}
\setlength{\headheight}{15pt}
\pagestyle{fancy}
\fancyhf{}
\fancyhead[LE,RO]{\thepage}
\fancyhead[RE]{\textit{Grundideernes Logik}}
\fancyhead[LO]{\textit{\leftmark}}
\setlength{\parindent}{1.5em}
\setlength{\parskip}{0pt}
\setstretch{1.3}

\title{\textbf{Grundideernes Logik}\\[4pt]\large II.}
\author{R.\ Nielsen}
\date{Kjøbenhavn: Forlagt af den Gyldendalske Boghandel (F.\ Hegel).\\[2pt]
  \small Trykt hos J.\ H.\ Schultz. 1866.}

% ---- original-page markers (paginate.py) ----
\usepackage{marginnote}
\newcommand{\opage}[1]{\marginnote{\footnotesize\textit{[#1]}}}
\newcommand{\apage}[1]{\marginnote{\footnotesize\textit{[#1]}}}
\begin{document}
\maketitle
\tableofcontents

% FRONT MATTER (roman folios III--XVI), replaces the marker pp. 0--0
% COPY-TEXT: British Library copy (Google scan, 1866-grundideernes-logik-2-bl.pdf;
% roman page = PDF + 2). First transcribed from the Harvard copy; every page
% (PDF 3, 5--18) re-read against the BL copy and repaired.
% ------------------------------------------------------------
% TITLE PAGE (PDF 3), not reproduced as text; \maketitle covers it.
% Wording and spelling exactly as printed (Fraktur, line by line):
%   Grundideernes Logik.
%   II.
%   Af
%   R. Nielsen.            [letterspaced]
%   [owner's armorial stamp of the scanned copy, not part of the print]
%   Kjøbenhavn.            [letterspaced]
%   Forlagt af den Gyldendalske Boghandel (F. Hegel).   ["Hegel" letterspaced]
%   Trykt hos J. H. Schultz.   [letterspaced; final glyph is the tz-ligature,
%                               read at 800 dpi in the Harvard copy and
%                               confirmed at 600 dpi in the BL copy: the
%                               printer is "Schultz", NOT "Schulz" as the
%                               preamble \date has it]
%   1866.
% The title page carries "Grundideernes Logik." with a full stop, and the
% imprint "Kjøbenhavn." (ø).
%
% PDF 4 (title verso): no printed matter, only a library stamp of the
% scanned copy. Nothing to transcribe.
% ------------------------------------------------------------

\newcommand{\tocentry}[3]{\par\noindent\hangindent=\dimexpr#1em+2.2em\relax\hangafter=1\hspace*{#1em}#2\dotfill\ #3\par}

% --- p. III ---
% PDF 5: Indhold, first page. The page bears NO folio; III is inferred.
\chapter*{Indhold.}
\addcontentsline{toc}{chapter}{Indhold}
\opage{[III]}%
\begin{center}---\end{center}
{\setlength{\parskip}{2pt}
\hfill Side.\par
\tocentry{0}{B. \textbf{Logiske Functioner}}{1.}
\tocentry{1.5}{a. Functioner af Dommen: Formalfunctioner}{17.}
\tocentry{3}{α. Dommens Form: Functioner af Videns Ord}{19.}
\tocentry{4.5}{αα) Sprogformende Functioner}{20.}
\tocentry{6.5}{Orddannende Functioner}{30.}
\tocentry{6.5}{Formdannende Functioner}{35.}
% PRINTER: page number "42" printed without its full stop (all others have one); recurs in BL copy
\tocentry{4.5}{ββ) Tankeformende Functioner}{42}
\tocentry{4.5}{γγ) Ideeformende Functioner}{70.}
\tocentry{6.5}{Mythiske Functioner}{73.}
\tocentry{6.5}{Gnostiske Functioner}{80.}
\tocentry{6.5}{Theosophiske Functioner}{86.}
\tocentry{6.5}{1) Ordet og Gud i tvetydig Eenhed}{89.}
\tocentry{6.5}{2) Ordet og Gud i tvetydig Modsætning}{90.}
\tocentry{6.5}{3) Ordet og Gud i udviklet Modsætningseenhed}{92.}
\tocentry{3}{β. Dom og Gjenstand: Functioner af Magtens Ord}{95.}
\tocentry{4.5}{αα) Viden af Dom og Gjenstand: theoretiske Functioner.}{101.}
% PRINTER: printed "Porallelisme", sense "Parallelisme" (confirmed at 600 dpi); recurs in BL copy
\tocentry{6.5}{1) Porallelisme af Dom og Gjenstand: supponerende Functioner}{103.}
\tocentry{6.5}{2) Gjenstanden bestemmer Dommen: abstraherende Functioner}{108.}
\tocentry{6.5}{3) Dommen bestemmer Gjenstanden: determinerende Functioner}{112.}
\tocentry{4.5}{ββ) Magten i Dom og Gjenstand: praktiske Functioner.}{117.}
\tocentry{6.5}{1) I den logiske Dom præformeres Gjenstanden: anticiperende Functioner}{122.}
\tocentry{6.5}{2) Med den antilogiske Dom realiseres Gjenstanden: executive Functioner}{128.}
\tocentry{6.5}{3) I den logisk-antilogiske Dom forklares Gjenstanden: explicative Functioner}{139.}
% --- p. IV ---
\opage{IV}\hfill Side.\par
\tocentry{4.5}{γγ) Ideen i Dom og Gjenstand: ideale Magtfunctioner.}{147.}
\tocentry{6.5}{1) Magtordets Idealitet: imperative Functioner}{148.}
\tocentry{6.5}{2) Magtordets Realitet: indicative Functioner}{157.}
\tocentry{6.5}{3) Det kategoriske Magtsprog: interpreterende Functioner}{166.}
\tocentry{3}{γ. Dommens Gyldighed: Functionernes Modalitet}{176.}
\tocentry{4.5}{αα) Den assertoriske Dom: Functioner af umiddelbar Bestemthed}{180.}
\tocentry{6.5}{1) Det antilogiske Element i den assertoriske Dom: Modalitetens objective Side}{182.}
\tocentry{6.5}{2) Det logiske Moment i den assertoriske Dom: Modalitetens subjective Side}{191.}
\tocentry{6.5}{3) Den assertoriske Doms logisk-antilogiske Væsenhed: dobbeltsidig Modalitet}{201.}
\tocentry{4.5}{ββ) Den problematiske Dom: Functioner af ophævet Bestemthed}{210.}
% lower-case "den" as printed (the other entries begin with a capital)
\tocentry{6.5}{1) den skeptiske Mulighed: Modalitetens subjective Side}{212.}
\tocentry{6.5}{2) Den dogmatiske Mulighed: Modalitetens objective Side}{222.}
\tocentry{6.5}{3) Mulighedernes Ontologie: dobbeltsidig Modalitet.}{234.}
\tocentry{4.5}{γγ) Den apodiktiske Dom: Functioner af reflecteret Bestemthed}{247.}
\tocentry{6.5}{1) Det Logisk-Apodiktiske: væsentlig Nødvendighed.}{253.}
\tocentry{6.5}{2) Det Antilogisk-Apodiktiske: virkelig Nødvendighed.}{264.}
\tocentry{6.5}{3) Det Rationelt-Apodiktiske: absolut Nødvendighed.}{275.}
\tocentry{1.5}{b. Functioner af Slutningen: Modalfunctioner}{282.}
\tocentry{3}{α. Subjectiv Slutning: Modalfunctioner af Viden}{293.}
\tocentry{4.5}{αα) Umiddelbar Slutning: analytiske Modalfunctioner.}{294.}
\tocentry{6.5}{1) Schematisme: de saakaldte umiddelbare Slutningers Schema}{299.}
\tocentry{6.5}{2) De analytiske Modalfunctioner}{302.}
\tocentry{6.5}{3) Den umiddelbare Slutnings Overgang i den middelbare}{309.}
\tocentry{4.5}{ββ) Den middelbare Slutning: synthetiske Modalfunctioner}{312.}
\tocentry{6.5}{1) Syllogistik}{315.}
\tocentry{6.5}{2) Deduction af syllogistiske Former}{319.}
% --- p. V ---
\opage{V}%
\tocentry{6.5}{3) Syllogistikens Opløsning}{324.}
\tocentry{4.5}{γγ) Den immanente Slutning: dialektiske Modalfunctioner}{329.}
\tocentry{6.5}{1) Den immanente Slutning som Form for Subjectiviteten selv}{330.}
\tocentry{6.5}{2) Den immanente Slutning som Form for Subjectivitetens Indhold}{344.}
\tocentry{8.5}{Kategoriske Immanensslutninger: Subsumtionsproblemet}{347.}
\tocentry{8.5}{Disjunctive Immanensslutninger: Disjunctionsproblemet}{351.}
\tocentry{8.5}{Slutningssystemernes Totalisering: Individuationsproblemet}{356.}
\tocentry{6.5}{3) Den immanente Slutning som Form for den frie Nødvendighed}{363.}
\tocentry{3}{β. Objectiv Slutning: Modalfunctioner af Magt}{384.}
\tocentry{4.5}{αα) Den objective Slutnings Formalbegreb}{390.}
\tocentry{4.5}{ββ) Magtens Ord i den objective Slutning}{408.}
\tocentry{6.5}{1) I Magtens Ord er Tyngden en objectiv Slutning.}{409.}
\tocentry{6.5}{2) I Magtens Ord er Affiniteten en objectiv Slutning}{414.}
\tocentry{6.5}{3) I Magtens Ord er den organiske Proces en objectiv Slutning}{420.}
\tocentry{4.5}{γγ) Den objective Slutning og det kategoriske Magtsprog.}{424.}
\tocentry{6.5}{1) Den mechaniske Totalitet i kategorisk Mening}{426.}
\tocentry{6.5}{2) Objectiveringens anden Grundtotalitet bestemt ved det kategoriske Magtsprog}{429.}
\tocentry{6.5}{3) Objectiveringens tredie Grundtotalitet bestemt ved det kategoriske Magtsprog}{432.}
\tocentry{3}{γ. Ideeslutning: Modalfunctioner af Sandhed}{434.}
\tocentry{4.5}{αα) Idealitetsslutning: Ideernes Præexistens}{435.}
\tocentry{4.5}{ββ) Realitetsslutning: Ideernes Realexistens}{450.}
\tocentry{4.5}{γγ) Afslutning: Idee-Existensens Fuldendelse}{456.}
}
% printed double rule closes the Indhold
\begin{center}---\end{center}

% --- p. VI ---
% PDF 8: Rettelser. The page bears NO folio; VI is inferred.
% Transcribed as printed; the errata are NOT applied anywhere in the text.
% The entries are not in page order (282 after 315; 239 after 328) as printed.
\chapter*{Rettelser.}
\addcontentsline{toc}{chapter}{Rettelser}
\opage{[VI]}%
\begin{center}---\end{center}
{\setlength{\parindent}{0pt}
S.\ 79 L.\ 15 f.\ o.\ udvikes, læs: udvikles.\\
--- 125 --- 4 f.\ o.\ en, læs: er.\\
--- 130 --- 11 f.\ n.\ des, læs: das.\\
% PRINTER: printed "læs;" (semicolon), sense "læs:"; recurs in BL copy
--- 139 --- 7 f.\ n.\ I den, læs; 3) I den.\\
--- 194 --- 12 f.\ n.\ høist, læs: høiest.\\
--- 195 --- 6 f.\ n.\ belyste, læs: belyse.\\
--- 305 --- 2 f.\ n.\ De, læs: Det.\\
--- 315 --- 10 f.\ n.\ Syllogistisk, læs: 1) Syllogistik.\\
% Greek in bold antiqua; first form has two acute accents, the correction
% has smooth breathing (coronis) on the second alpha
--- 282 --- 7 f.\ n.\ καλοκαγάθία, læs: καλοκἀγαθία.\\
--- 398 --- 10 f.\ n.\ realiset, læs: realiseret.\\
--- 400 --- 10 f.\ o.\ blot Række, læs: Række blot.\par}
\begin{center}---\end{center}
\begin{center}{\bfseries Rettelser til Grundideernes Logik I.}\end{center}
{\setlength{\parindent}{0pt}
S.\ 327 L.\ 14 f.\ n.\ frembringes, læs: forbruges.\\
--- 328 --- 14 f.\ o.\ (H O) læs: (H\textsubscript{2} O).\\
--- 328 --- 8 f.\ n.\ denne Kjendsgjerning, læs: dette som Kjendsgjerning.\\
--- 239 --- 17 f.\ o.\ dannede, læs: dannende.\par}
\begin{center}---\end{center}

% --- p. VII ---
% PDF 9: the prefatory text. NO heading is printed (the text opens with a
% large initial G), so the heading below is editorial. It is a dated and
% signed preface. The page bears NO folio; VII is inferred.
\chapter*{[Indledning]}
\addcontentsline{toc}{chapter}{[Indledning]}
\opage{[VII]}Grundideernes Logik er et Forsøg i den rene Philosophie, et Forsøg til en udpræget Theisme, en videnskabelig Protest imod al Theologie.

\emph{Den rene Philosophie.} Opgaven er at frigjøre Tænkningen fra uvedkommende Forestillinger, Phantasier og Sandsebilleder, at den kan sættes istand til at udvikle Begreber og under Begrebernes Udvikling erkjende og fastholde Ideens nødvendige Former. Naar det imidlertid efter en eengang vedtagen Sprogbrug hedder, at det i Logiken er Tænkningen alene som Tænkning, der udretter Alt, da gjælder det i dette som i saamange andre Tilfælde: \textit{a potiori fit denominatio.} Den saakaldte rene Tænkning er, naar den tages for sig, en Abstraction som, trods speculative Hexerier med den rene Væren, uden Anskuelsens Hjælp ikke vilde være istand til at røre sig af Stedet. Det er desaarsag en Misforstaaelse, naar „den rene Philosophie“ skal indskrænkes til simpel Afliren af den gamle Melodie paa Væren og Tænken og Tænken og Væren. Der kan i den rene Philosophie indføres mange og forskjelligartede Aandsphænomener, vælges Exempler fra Livets og Videnskabens allerforskjelligste Sphærer; det hvorpaa Alt i denne Henseende væsentlig kommer an, er en Iagttagelse af følgende tre Grundfordringer: 1) Abstractionen maa gjennemføres, saa at Begrebet med dets Conseqvenser paa ethvert Udviklingstrin kan blive tilstrækkelig klart
% --- p. VIII ---
\opage{VIII}og gjennemsigtigt; 2) Kritiken af paagjældende Aandsphænomener og Bearbeidelsen af valgte Exempler maa anlægges saaledes, at kun det Logiske og det, der i særegen Grad tjener til Tydeliggjørelse af ontologiske Problemer og deres Løsning, kommer i Betragtning; 3) Udviklingen maa skride frem i væsentlig Continuitet, saa at det Efterfølgende Trin for Trin kan sees at fremgaae med logisk Nødvendighed af sit Foregaaende: kun i en continuerlig Strøm kan Ideen overalt være med i Bevægelsen.

\emph{Udpræget Theisme.} Ihvorvel en Grundideernes Logik ikke i noget Partie kan falde umiddelbart sammen med Religionsphilosophien, saa er det dog aldeles nødvendigt, at den, medmindre den da skulde kunne hjælpe sig med et atheistisk Princip, maa udvikle saavel Gudsideens som Verdensideens logiske Form, maa logisk bestemme Forholdet mellem disse Ideer og saaledes afgive Grundlaget for en tilsvarende Religionsphilosophie. Umuligheden af at give enten en Ideens eller Ideernes Logik uden med det Samme at angive Grundtrækkene til en philosophisk Verdensanskuelse, maa vel være indlysende for Enhver, der har noget Begreb om Sagen. Men hvorved bestemmes nu egentlig en saadan philosophisk Verdensanskuelse? Psychologisk seet, kunde det næsten tage sig ud, som om den Philosopherende havde en personlig Anskuelse færdig i Forveien, og saa, istedenfor at rette sin philosophiske Verdensanskuelse efter Logiken, tvertimod lempede Logiken efter sin personlige Anskuelse. Paa denne Maade vilde en Atheist da anlægge en atheistisk, en Pantheist en pantheistisk, en Deist en deistisk, og en Theist en theistisk Logik.

At nu virkelig den Philosopherendes personlige Anskuelse maa blive en Factor med i den eiendommelige Maade, hvorpaa han opfatter og bestemmer Ideens logiske Væsenhed, kan ikke
% --- p. IX ---
\opage{IX}negtes; men Philosophien vilde ikke være nogen Videnskab, dersom den ikke var istand til at forskaffe sig en aldeles nøiagtig og uafviselig Maalestok, hvorefter de forskjelligartede, indbyrdes modstridende logiske Systemers større eller mindre videnskabelige Gyldighed lader sig bedømme. Det kommer ved Prøvelsen af et logisk System fornemmelig an paa at undersøge, hvorledes dets Princip, Begyndelse og Methode oprindelig er bestemt. Fortjener det givne System virkelig Navn af System, da vil et Blik paa Anlæget i det Hele i Forening med et grundigt Bekjendtskab til Begyndelsen og de første Partier ordentligviis være tilstrækkeligt til at afgjøre, til hvilken Verdensanskuelse et saadant System med Nødvendighed maa føre. De, der kjende den hegelske Philosophie f.\ Ex., ville allerede af Logikens første Deel, af de første Partier i det første Afsnit om Væren og navnlig af Begyndelsen med den rene Væren kunne see, at den til en saadan Logik svarende Verdensanskuelse nødvendig maa blive en logisk Pantheisme, en Panlogisme, som kun ved Mystificationer af „Begrebet“, det rene Begreb, den immanente Methodes logiske Proteus, vil kunne give sig et theistisk Anstrøg.

At Grundideernes Logik i de to nu foreliggende Forsøg efter Anlæg, Princip og Methode er og maa være Theisme, utvetydig Theisme, kan ligeledes sees af Begyndelsen selv og de første Partier. Er det sandt, at Philosophien ikke kan begynde med Eenhed af Tænken og Væren, dersom den ikke udtrykkelig begynder med Viden; er det sandt, at Vidensprincipet absolut maa falde sammen med Subjectivitetsprincipet; er det sandt, at Subjectiviteten maa være Princip, ikke blot for sig selv, men for al Objectivitet, og kun kan være hiint formedelst den absolute Viden, dette formedelst den absolute Magt og begge Dele ved den Modsætningseenhed af Viden og Magt, der alene giver den absolute Sandhed: da vil det allerede heraf
% --- p. X ---
\opage{X}være klart, at den logiske Bestemmelse af Grundforholdet mellem Gudsideen og Verdensideen maa blive theistisk.

\emph{Videnskabelig Protest imod al Theologie.} At en atheistisk Philosophie protesterer imod al Theologie, er i sin Orden. „Men“, kunde man sige, „hvorledes er det muligt, at en Philosophie, der anmelder sig som Forsøg til en udpræget Theisme, som anerkjender Gudsideen og følgelig fra sit Standpunkt maa lægge an paa at udvikle Gudsbegrebet, kan fremstille sig som en videnskabelig Protest, ikke alene imod Theologien i en vis bestemt Form, men ubetinget imod al Theologie? En videnskabelig Udvikling af Gudsbegrebet maa jo dog altid falde sammen med en objectiv Lære om Gud. Da nu en hvilkensomhelst Lære om Gud --- naar man blot vil erindre den gamle orthodoxe Fiirdeling: \textit{theologia accipitur sensu quadruplici: generalissime pro quavis de Deo doctrina, generaliter pro theologia vera, specialiter pro theologia revelata, specialissime pro doctrina de Deo uno et trino} --- maa kaldes Theologie, saa spørges: hvorledes er det muligt, at en Philosophie skal kunne opstille en Lære om Gud, som dog, for ikke at være Theologie, ikke tør være nogen Lære om Gud?“

Til en Besvarelse af dette Spørgsmaal begynde vi helst med en Reduction af Theologiens fiirdobbelte Form, idet vi sammenfatte de to første Definitioner under Benævnelsen: naturlig Theologie (\textit{theologia naturalis}), de to sidste under Kategorien: Aabenbaringstheologie (\textit{theologia revelata}).

At nu den naturlige Theologie, der oprindelig er en Abstraction af en positiv Kirkelære, maa være meget forskjellig fra sand Philosophie, er let at indsee. Den naturlige Theologie begynder nemlig paa sin Theologisk med en umiddelbar Gudsforestilling. I en saadan umiddelbar Forestilling bliver Gud selv da opfattet som et oververdsligt objectivt Væsen med visse bestemte Egenskaber. Opgaven er: at bevise, at et saadant
% --- p. XI ---
\opage{XI}oververdsligt objectivt Væsen med saadanne Egenskaber virkelig er til, at dette Væsen virkelig har skabt Verden o.\ s.\ fr. Med en saadan umiddelbar Gudsforestilling, en Forestilling, der gjør Gud til en for Betragtningen udvortes Gjenstand, kan Philosophien, saasandt den vil være Videnskab, umulig begynde. Istedenfor at indlade sig med umiddelbare Gudsforestillinger og qvakle med Raisonnementsbeviser for Guds Tilværelse, begynder Philosophien med at forudsætte, hvad der ikke er et Hiinsidigt eller Fraværende, men et for den almene Viden væsentligt Nærværende, nemlig Tilværelsen selv med sit uendeligt rige Indhold. Dette Videns Indhold bestemmes da efter dets indholdsrige Eenhed med sig selv som det Absolute, og Opgaven bliver at udvikle og klare de for Værens absolute Indhold adæqvate Vidensformer.

Denne Opgave er videnskabelig, thi slige Former kunne udvikles, revideres og corrigeres med en Nøiagtighed, der i Dialektiken ikke staaer tilbage for det Exacte i Mathematiken. For Formernes Udvikling er Methoden saaledes anlagt, at Indholdet som ved et Ideens Høidetryk strømmer til af sig selv: er blot Formen sand, saa er enhver Tvivl om Indholdets Realitet med det Samme hævet. Som den fuldkomne Form for det absolute Indhold bliver da Gudsbegrebet udviklet; i det udviklede Begrebs Eenhed med sit Indhold erkjendes Gudsideen; i Gudsideen erkjendes Gud; men denne Erkjendelse er ikke theologisk.

Hvad dernæst Aabenbaringstheologien, den kirkelig positive Theologie, særlig „den speculative Dogmatik“, angaaer, da har Forfatteren af nærværende Skrift jo allerede udtalt sin Anskuelse saa ofte og saa uforbeholdent, at han her kan indskrænke sig til en stereotyp Gjentagelse af, hvad han, for at samle Opmærksomheden paa det første Nødvendige, maa ansee for afgjørende Punkter.

Det kommer da for det Første an paa at angive det i
% --- p. XII ---
\opage{XII}Theologien Theologiske ved en Begrebsbestemmelse, hvis Gyldighed enhver Sagkyndig maa indrømme. Vi fastholde følgende Kriterium: „Det i Theologien Theologiske kjendes derpaa, at den Aarsag, hvoraf alle Virkninger i Naturens som i Aandens Rige deels middelbart, deels umiddelbart afledes og forklares, ikke søges i Tingenes Natur eller Fornuftens Væsen, men i den absolute Villie, i den Almægtige, der skaber af Intet. Men selv den reneste, skarpsindigste og conseqventeste Anvendelse af et saadant Princip (Occasionalisme, Malebranche) gjør Videnskaben uvidenskabelig.“

I Conseqvens af den opstillede Begrebsbestemmelse kommer det saa for det Andet an paa, hvorledes Theologiens og Philosophiens Mellemværende bliver at opgjøre. Vi vælge for Philosophiens Vedkommende følgende Opgjørelse: I Middelalderen var det Theologien, der angreb Philosophien, fordi Philosophien ikke var troende nok; i Nutiden er det omvendt Philosophien, der angriber Theologien, fordi Theologien ikke er videnskabelig nok. Det er da i denne Henseende nødvendigt, at der maa gjøres en skarp Adskillelse imellem videnskabeligt i videre og i snævrere Forstand. „Videnskabeligt i videre Forstand er Alt, hvad der bliver Gjenstand for Videnskabens Behandling: det Videnskabelige ligger i Behandlingen; videnskabeligt i snævrere Forstand er derimod kun det, hvis Gyldighed Videnskaben indseer, hvis Sandhed den kan bedømme, prøve og anerkjende“.

Ved at benytte den Tvetydighed, der ligger i Kategorien: „Videnskabeligt“, lykkes det undertiden Troeslæren at give sig Skin af Videnskab; det er da Videnskabslærens Opgave at paavise de vigtigste herhen hørende Sophismer.

\emph{Første Sophisme.} „Skal Troen ikke beroe paa subjective Indfald og priisgives Individernes Luner, maa det udtrykkelig kunne paavises, hvad man skal troe: Troen maa have
% --- p. XIII ---
\opage{XIII}sin Gjenstand. At undersøge Troesgjenstandene grundig er at undersøge dem videnskabeligt. Troeslæren beroer paa en saadan Undersøgelse; altsaa er Troeslæren en Videnskab“.

Denne Sophisme giver Sagen det Udseende, at Troesgjenstandene sikkres ved Videnskaben, fordi de prøves i Videnskaben. Sagen forholder sig omvendt. Troesgjenstandene opløse sig i Dunst, naar Videnskaben først faaer Lov til at anbringe sin Kritik. Det er en videnskabelig Kritik ligesaa umuligt at anerkjende Evangeliet om den Opstandne, som det er Videnskaben umuligt at anerkjende Mythen om Herkules. At anerkjende et Mirakel i Kraft af videnskabelig Kritik er en Modsigelse, da Muligheden af et Mirakel gjør videnskabelig Kritik umulig: hvor Underet taler, maa Forstanden tie, hvor Miraklet begynder, maa Videnskaben ende.

\emph{Anden Sophisme.} „Enhver eiendommelig Videnskab maa nødvendigviis have sine Forudsætninger. Ligesom Logiken forudsætter og retfærdiggjør Tænkningens Gyldighed, og Æsthetiken det Skjønnes Gyldighed, saaledes maa Troeslæren for sit Vedkommende forudsætte og retfærdiggjøre Gyldigheden af Troen og Troens Indhold. Den, der beviser, at Troeslæren ikke er en Videnskab, fordi den har sine særegne Forudsætninger, beviser med det Samme, at der slet ingen Videnskab gives; men her gjælder det: \textit{qui nimium probat, nihil probat.} Altsaa er Troeslæren en Videnskab“.

Denne Sophisme blender ved den Tvetydighed, der ligger i at medbringe en Forudsætning. Hvad der gjør Troeslæren uvidenskabelig, er ingenlunde den Omstændighed, at den medbringer Forudsætninger, men netop den Omstændighed, at de Forudsætninger, den medbringer, stride mod Videnskaben. Den forudsætter Miraklet, den forudsætter Subjectivitetens og den subjective Forvisnings absolute Gyldighed; med disse Forud%
% --- p. XIV ---
\opage{XIV}sætninger har den saa at sige væltet sig ind paa Videnskaben, og brudt den videnskabelige Maalestok. Skulle Conseqvenserne være videnskabelige, maae Forudsætningerne ogsaa være det; paa uvidenskabelig Grund kan ingen videnskabelig Bygning opføres: Naturvidenskab er mulig, men Mirakelvidenskab er umulig.

\emph{Tredie Sophisme.} „Troeslærens overnaturlige Forudsætninger have, hver for sig, vistnok kun Gyldighed for den Troende; men imellem disse Sætninger indbyrdes gives der en logisk og forsaavidt objectiv Sammenhæng. Da Sammenhængen er objectiv, og Troessætningerne netop finde deres sande Udtydning og dybere Begrundelse i denne Sammenhæng, er Begrundelsen selv objectiv, og Troeslæren altsaa i Kraft af denne Begrundelse en Videnskab“.

Det Blendende i denne Sophisme beroer derpaa, at en logisk Sammenhæng uden videre bliver til en objectiv Sammenhæng, og denne igjen til en objectiv Begrundelse. „Alle Mennesker ere Triangler; Cajus er et Menneske, altsaa er Cajus en Triangel. I enhver Triangel er Vinklernes Sum liig to Rette. Cajus er en Triangel, altsaa er Vinklernes Sum i Cajus liig to Rette“. Imellem disse Sætninger er her unegtelig en logisk Sammenhæng; men er her derfor en objectiv videnskabelig Begrundelse? Fra Livets Standpunkt kan der være en Evighedsforskjel imellem det Hellige og det Profane, det Evangeliske og det Mythologiske; men naar Videnskaben nu eengang ikke kan bedømme dette Hellige, ikke magte dette Evangeliske, er Forskjellen i Videnskabens Øine, som om den slet ikke var. --- Da Troessætningerne hver for sig ere ubetinget subjective, ere de i Videnskaben ligesaa vilkaarlige som de meest eventyrlige Phantasier. Gives der altsaa en videnskabelig Begrundelse af Troeslærdomme, saa gives der ogsaa en videnskabelig Begrundelse af Mythelærdomme, en videnskabelig Begrundelse af Nonsens.

% --- p. XV ---
\opage{XV}\emph{Fjerde Sophisme.} „Aanden i Troeslæren, „den Hellig Aand“, er ikke blot i subjectiv Forstand en Troens Aand, den er i objectiv Forstand Sandhedens absolute Aand, den Aand, der „randsager Guds Dybheder“ og „veileder os til al Sandhed“. Staaer det nu fast, at enhver Videnskab forudsætter en Sandhed; staaer det fast, at de lavere Sandheder altid maae søge deres sidste Forklaring i en høieste Sandhed, saa staaer det ogsaa fast, at den Aand, der veileder os til al Sandhed, maa være den Aand, der veileder os til al Videnskab; Troeslæren er altsaa ikke blot en Videnskab, men den høieste Videnskab, den, hvori al anden Videnskab forklares“.

Det Blendende i denne Sophisme beroer paa en dobbelt Subreption. Det antages, at den Aand, der gaaer igjennem Theologernes objective Troeslære, virkelig er den „Hellig Aand“; denne Antagelse er falsk. Hvert Kirkesamfund har jo, som bekjendt, en egen Troeslære; den ene Troeslære staaer i Strid med den anden; Tilhængerne af den ene modsige og fordømme Tilhængerne af den anden. Var nu Troeslæren væsentlig et Værk af „den Hellig Aand“, da maatte „den Hellig Aand“ jo være splidagtig med sig selv; men en Aand, der er splidagtig med sig selv, kan ikke være Sandhedens Aand. Det antages, at „den Hellig Aand“ i sit Princip maa være en videnskabelig Aand. Denne Antagelse er falsk. Var „den Hellig Aand“ en Videnskabens Aand, saa maatte Christi og Apostlernes Erkjendelse af den evige Sandhed have været den høieste videnskabelige Erkjendelse, saa maatte Christus og hans Apostle, med andre Ord, have været de største Videnskabsmænd paa Jorden. Var Troens Aand i Sandhed en Videnskabens Aand, saa maatte det jo gaae fremad med Theologien, eftersom det i det Hele gaaer fremad med Videnskabeligheden. Men hvordan gaaer det? Videnskabeligheden gaaer frem, men Theologien gaaer tilbage.

% --- p. XVI ---
\opage{XVI}Efter at de Hovedpunkter, hvortil en videnskabelig Protest mod al Theologie i det Hele lader sig henføre, saaledes ere opstillede, kommer det saa for det Tredie an paa at faae disse Punkter \emph{enten videnskabelig anerkjendte eller videnskabelig gjendrevne}.

De, som nu mene, at have „Aandens Vidnesbyrd“ for, at Theologien virkelig er en Videnskab, ville i de opstillede Sætninger finde en beqvem, og, som vi haabe, velkommen Anledning til at forhøie Theologiens Anseelse og lade dens Lys skinne, idet de Punkt for Punkt belyse de Feiltagelser og Misforstaaelser, hvori Forfatteren af nærværende Skrift, være sig ved det opstillede Kriterium eller de paaviste Sophismer, beviislig maatte have gjort sig skyldig. Det kommer, vel at mærke, an paa at eftervise \emph{bestemte} Feiltagelser i de \emph{bestemte} Sætninger; thi paa vulgære Udenomsbemærkninger, improviserede Anskuelser, Ændringer eller Opkog af gamle længst gjendrevne Synsmaader o.\ A.\ dsl.\ vil Undertegnede ikke reflectere.

\bigskip
\noindent{\bfseries Den 11te Marts 1866.}

\begin{flushright}{\bfseries R. Nielsen.}\end{flushright}

% --- p. 1 ---
\begin{center}
{\large\bfseries B.}\\[3pt]
{\Large\bfseries Logiske Functioner.}
\end{center}
\addcontentsline{toc}{chapter}{B. Logiske Functioner}

\opage{1}Det er en gammel Indvending, at Philosophien hænger sig i Ord, at dens Forklaringer ere Ordforklaringer, dens Stridigheder Ordstrid, dens Viisdom Ordkram. Denne Indvending er ikke ganske uden Grund, naar man eensidig fæster Opmærksomheden paa Skyggesiderne, paa de ufrugtbare Subtiliteter, de smagløse Udskeielser i tom Ordkløgt, om hvilke en Johan af Salisbury\footnote{\textit{Nihil est odibilius subtilitate, ubi nihil aliud est quam subtilitas.}}% note *)
, en Erasmus fra Rotterdam, en Baco af Verulam og hos os en Holberg have, hver i sin Tidsalder, holdt Rettergang med den speculative Videns Clienter. Grundløs og intetsigende bliver Indvendingen derimod, naar man i blind Iver for det Reelle ringeagter det Ideelle og bilder sig ind, at Forstanden allerbedst seer Tingen, naar den overseer Ordet. At Philosophien umulig kan gaae ind paa en saadan Betragtningsmaade, forstaaes af sig selv. Philosophien er Ordets Videnskab i særegen Forstand; den maa nødvendig holde paa Ordet, saasandt den vil holde paa Tanken og fastholde Begrebet.

Eller hvorledes skulde det vel være muligt at skille Begreber fra tilsvarende Ord, Begrebsudviklinger fra Ordforklaringer? Begrebet Væren og Ordet Væren, Begrebet Intet og Ordet Intet, Begrebet Noget og Ordet Noget ere jo saa uadskillelige, at man ikke kan forandre Ordet uden at ændre Begrebet, ikke opgive Ordet uden at opgive Begrebet. Ligesom
% --- p. 2 ---
\opage{2}en Sjæl ikke kan være uden et Legeme, saaledes kan Begrebet som det i sig bestemte Begreb heller ikke være uden et bestemmende Ord. Det er forsaavidt ret betegnende, at Logiken har sit Navn af Logos, og at Logos er baade Ord og Tanke; thi dersom Tanken ikke var, saa var der intet Ord, men dersom Ordet ikke var, saa var der heller ingen bestemt Tanke.

Hvad der i denne Henseende gjælder om Philosophien, gjælder paa tilsvarende Maade om Fagvidenskaben. Det er utænkeligt, at nogen Realvidenskab virkelig skulde see sig istand til at erkjende Tingen uden formedelst Ordet, opstille Begreber, der ikke beroede paa Ord, give os Sagforklaringer, der ikke tillige vare Ordforklaringer. I enhver Videnskab maa der nødvendig gjøres Forskjel imellem en Lære om Tingene og Tingene selv; det er Læren om Tingene, ikke Tingene selv, der udgjør Realvidenskabens Indhold. Kunde Videnskaben gjøre Læren om Tingene til Eet med Tingene, da maatte man jo f.\ Ex.\ gjenfinde den virkelige Stjernehimmel i Astronomien, Planteverdenen i Botaniken, Dyreverdenen i Zoologien o.\ s.\ v.; men da det ikke er Gjenstanden, men Læren om Gjenstanden, som Videnskaben byder os, saa følger det af sig selv, at Læren maa holde sig til Ordet. Selv i Mathematiken, hvor de exacte Betegnelser dog paa saa mange Maader træde i Ordets Sted, er Ordet nødvendigt for at hæve Betegnelserne til Begrebsbestemmelser.

Er det Ordforklaringer, naar Logikeren udvikler Forskjellen imellem Noget og Andet, fordi Begrebet ikke kan indeholde Andet, end hvad der ligger i Ordet, da er det ogsaa Ordforklaring, naar Mathematikeren gjør Forskjel imellem Vinkel og Bue, Astronomen mellem Parallaxe og Aberration, Botanikeren mellem Blad og Green, Zoologen mellem Pote og Hovfod; thi Bue, Parallaxe, Green, Hovfod ere Ord ligesaavel som Noget og Andet. Hvad der nøiagtig skal tænkes ved saadanne Ord som Pote og Hovfod maa, selv om vi havde Poter og Hovfødder i tusinde ja titusinde Exemplarer
% --- p. 3 ---
\opage{3}for os, ligesaa nødvendig forklares ved andre Ord, som det ved andre Ord maa forklares, hvad der logisk ligger i Ordene Noget og Andet.

At Ordstrid heller ikke er til at undgaae i Realvidenskaben, kan da oplyses ved utallige Exempler; vi vælge et Exempel af Botaniken. \textit{Taxus baccata} (Taxtræet) har aabenbart sit Tilnavn deraf, at man har anseet dens Frugt for et Bær (\textit{bacca}). Er det nu virkelig et Bær, eller skulle vi bruge en anden Benævnelse? Et saadant Spørgsmaal kan ikke afvises med en Bemærkning som den, at det jo ikke kommer an paa et Navn, et Ord, men paa Sagen; thi en Sagforklaring er her umulig uden en tilsvarende Ordforklaring. --- At den foregivne Frugt har Lighed med et Bær, er kjødet, tilrundet, farvet og er Resultat af en Blomst, skal ikke negtes; men den foregivne Frugt er desuagtet ingen Frugt. En Frugt er den efter Befrugtningen udviklede Frugtknude; men her er ingen Frugtknude. Er her ingen Frugtknude, saa er her heller ingen Frugt; er her ingen Frugt, saa er her heller intet Bær, thi et Bær er en Frugt. Maaskee er det saakaldte Bær da et aabent Frugtblad, et saadant, som i Almindelighed findes hos de Nøgenfrøede? Heller ikke denne Antagelse er holdbar. Et Frugtblad maa nødvendigviis være et eget Blad (\textit{folium sui generis}). Men det, vi her have med at gjøre, er intet Blad; det er en Frøstol, en Frøstreng, altsaa en Axe med en Udvikling, ganske analog med den hos Figenen, Rosen o.\ s.\ v. Man seer heraf, hvorledes Striden om det foregivne Bær bliver en Strid om Ord; thi Bær er et Ord, Frugt er et Ord, Frugtknude, Frugtblad, Frøstol o.\ s.\ v.\ ere lutter Ord. Striden om Tingen bliver altsaa en Ordstrid; uden en saadan Ordstrid vilde en videnskabelig Begrebsudvikling være umulig.

„Men“, siger man, „der er dog, hvad Ordstrid angaaer, en væsentlig Forskjel imellem Fagvidenskaberne og Philosophien. Vistnok maa enhver Fagvidenskab udtrykke sine Begreber i
% --- p. 4 ---
\opage{4}Ord, men de i Ord udtrykte Begreber blive, hvad ovenstaaende Exempel da ogsaa tilstrækkelig viser, i den Grad prøvede og sikkrede ved ubestridelige Kjendsgjerninger, at det som Følge heraf ikke bliver Begrebet, der maa rette sig efter Ordet, men tvertimod Ordet, der maa rette sig efter Begrebet. I Philosophien er Opmærksomheden derimod saa eensidig henvendt paa Tanken og ved Tanken paa Ordet, at Tingen bliver afhængig af Tanken og Tanken af Ordet“. Modsætningen mellem Philosophie og Realvidenskab vilde da fra det herved angivne Synspunkt være at opfatte saaledes: i Philosophien er det Ordet, der bestemmer Tingen, i Realvidenskaben er det omvendt Tingen, der bestemmer Ordet. Men hvor ugrundet denne Synsmaade er, skal i det Følgende vises.

At Ting formedelst Ord bestemmes ved Ting, giver en Sagforklaring; at Ord formedelst Ting bestemmes ved Ord, en Ordforklaring. Det skal nu undersøges, hvorvidt man paa endelig Maade er istand til at holde Sagforklaringen udenfor Ordforklaringen og dog tilfredsstille Begrebets Fordring. Hvad Ordet angaaer, er Philologien den af Fagvidenskaberne, der staaer Philosophien nærmest; vi vælge da et Exempel fra Philologien. „For Betragtningen af den Række af Betegnelser af \emph{Forholdene i Sætningen}, dem, som angaae de enkelte substantiviske Forestillingers Forhold i Anskuelsen“, siger Madvig\footnote{Om de grammatikalske Betegnelsers Tilblivelse og Væsen, 1ste Stykke. Indbydelsesskrift til Kjøbenhavns Universitets Fest i Anledning af Hans Majestæt Kongens Fødselsdag den 6 October 1856. S.~20--22.}% note *)
, „er det af den meest afgjørende Vigtighed at tilegne sig og fastholde den Erkjendelse, at de grammatikalske Betegnelser slet Intet have at gjøre med Tingenes egne reale Forhold, men alene med de Forhold i Anskuelsen, hvori de enten i indholdsløs Almindelighed henføres til hinanden eller vise sig som Led i et bestemt enkelt Billede af Handling eller Tilstand, hvorfra Forholdet faaer Indhold. Herrens Forhold
% --- p. 5 ---
\opage{5}til Slaven opfattes efter sit Indhold ikke mindre som det modsatte af Slavens til Herren, fordi jeg betegner den Sammenhæng, hvori Forestillingerne \emph{Herre} og \emph{Slave} træde for Anskuelsen, naar den ene bestemmes ved sit Forhold til den anden, ved een og samme Casus fra begge Sider: \emph{Herrens Slave} og \emph{Slavens Herre}. I Grammatiken er Slaven, om man vil, ligesaavel Herrens Eier som Herren Slavens, det vil sige, Grammatiken har slet Intet at gjøre med Eiendomsforholdet, der er et realt Forhold, og \emph{Eieform} er et aldeles ligesaa forkeert Navn som de gamle latinske Casusnavne, der hentes fra reale Forhold; det bringer, ligesom hine, een af de utallige Skikkelser frem for Sindet, der i Anskuelsen subsumeres under Sammenhængsforholdet, og derved fremkalder det indirect og ubevidst denne sande Forestilling om Formens Betydning: Sammenhøren.“ --- Med Hensyn til nogle tydske Philologers speculativt forvirrede Opfattelser af Casusbetydningen, navnlig af Genitiv som den Casus, der „betegner Noget (den Forestilling, hvis Navn staaer i Genitiv,) som „Ursprung oder Bedingung des Werdens eines Andern“,
% PRINTER: the quotation opened at „betegner is never closed; left as printed.
tilføies i en Anmærkning: „Hvo seer nu ikke, naar han paamindes, at vi her ere midt inde i reale Forhold og i den allerplumpeste Forvirring? I „des Vaters Sohn“ er vistnok Faderen Ursprung und Bedingung des Werdens des Sohnes, men er i „des Sohnes Vater“ Sønnen ogsaa Ursprung des Vaters? Maaskee vilde Forfatteren være kløgtig nok til at svare, at Sønnen er Bedingung des Werdens des Vaters, fordi denne kun formedelst ham er Fader. Han vilde da kun forvexle Tilblivelsen af et Væsen og af et Navn. Men hvorledes vilde han slippe f.\ Ex.\ fra „des Vaters Hand.“ I saadan Uklarhed tumler man sig og af Sligt lader mange sig imponere, fordi der engang er Forvirring paa dette Omraade og fordi man fattes Abstractionens Kraft til at gribe Sagen som den er.“

Det er et stærkt Lys, som med disse faa, men opklarende
% --- p. 6 ---
\opage{6}Ord falder paa Modsætningen mellem Grammatik og Logik. At de grammatikalske Betegnelser slet Intet have at gjøre med Tingenes reale Forhold, vil da med andre Ord sige, at de ikke have Noget at gjøre med Tingenes reale Begreber. Ordforklaring og Sagforklaring træde paa denne Maade ret øiensynlig ud fra hinanden. Forsaavidt her i Talen om Herren og Slaven sees paa det Grammatikalske, er Slaven unegtelig ligesaavel Herrens Eier som Herren Slavens; og hvorfor skulde han ikke ogsaa være det? de høre jo begge i samme Ordklasse, ere begge af lige Rang, thi Slaven er jo et Substantiv ligesaavel som Herren. Anderledes, naar der i Talen udtrykkelig sigtes paa Sagen. I Realiteten er nu Herren A et virkeligt Individ, Slaven B ligeledes; men imellem de to Individer er det reale Forhold af Herre og Slave aldeles tilfældigt. A som Herre kunde nemlig ligesaavel tænkes at have en anden Slave, og B i Egenskab af Slave at have en anden Herre; ja, det kunde jo, hvad den umiddelbare Kjendsgjerning angaaer, ligesaa godt have været B, der var Herre, og A Slave, som omvendt. Der er da, som heraf sees, hverken i den rene Ordforklaring eller i den umiddelbare Sagforklaring, Tale om egentlig Begriben. Fra Ordforklaringens grammatikalske Synspunkt mangler Begrebet, fordi her mangler ethvert Hensyn til Tingenes reale Forhold; Forbindelser, som „Slavens Herre,“ „Kjolens Tjener,“ „Maanens Buxer,“ give i saa Henseende lige god Mening. Fra Sagforklaringens umiddelbart reale Synspunkt mangler Begrebet, fordi det reale Forhold i dets rene Facticitet mangler det til Begrebet Nødvendige, det Almeengyldige: nu er „A Herre og B Slave,“ nei! „nu er det B, der er Herre, og A, der er Slave“. Som det blotte Individ er Slaven intet Begreb, som det blotte Substantiv heller ikke; skal her vise sig Spor til Begreb, maae Tingen og Ordet, Individet og Substantivet i Proces med hinanden.
% PRINTER: printed "Individet og Substantivet i Proces", verb lacking; sense "Substantivet være i Proces" (or "træde i").

Begrebsudviklingen fordrer en Vexelbestemmelse af Ord\opage{7}%
% --- p. 7 ---
forklaring og Sagforklaring, en objectiv Reflexion, hvori Tingen bestemmer Ordet, og Ordet igjen Tingen: Tingen er antilogisk, kun Ordet er logisk. Forsaavidt Tingen bestemmer Ordet, er Slaven afhængig af Herren, ikke omvendt; men forsaavidt Ordet bestemmer Tingen, er Herren som Herre selv væsentlig afhængig af Slaven. Forsaavidt Tingen bestemmer Ordet, er Faderen „Ursprung und Bedingung des Werdens des Sohnes,“ ikke omvendt; forsaavidt Ordet derimod bestemmer Tingen, er Sønnen tillige væsentlig „Ursprung des Vaters.“ Hvad Virkeligheden udstykker, beholder Væsenet i reflecteret Eenhed; for at erkjende de modsatte Bestemmelsers væsentlige Eenhed, maa Virkelighedsforholdet idelig vendes om. At Forholdet vendes om, betegner i logisk Forstand ingenlunde, at Tilblivelsen af et Væsen blot ombyttes med Tilblivelsen af et Navn; det betyder tvertimod, at Anskuelsen af et Virkeligt ombyttes med Indsigten i et Væsentligt. Det er ved Virkningen, at Aarsagen ikke blot faaer Navn af, men væsentlig er Aarsag; det er ved Slaven, at Herren ikke blot faaer Navn af, men væsentlig er Herre; det er ved Sønnen, at Faderen ikke blot faaer Navn af, men væsentlig er Fader.

I Reflexionen af de omvendte Forhold sees Dialektiken. Dialektiken er væsentlig Ordforklaring, et Charakteermærke, der er træffende betegnet ved det schleiermacherske Udtryk, at Dialektiken er en „Sprogets Kunst.“ Men den udvortes, endelige Adskillelse af Ordforklaring og Sagforklaring afstedkommer Forvirring i Videnskaben. I Følelsen af, at Dialektiken er Ordforklaring og det ved Forholdenes Omvending endog en drillende Ordforklaring, slaaer man den uden videre i Hartkorn med Ordspil, Ordkløveri o.\ s.\ v. Man taler om det Væsentlige og mener det Virkelige; naar Dialektiken saa, for at gjøre det Væsentlige kjendeligt, vender Virkelighedsforholdet om, kan den Ordskarpsindige naturligviis i Ordet ikke see Andet end det bare Ord, protesterer altsaa,
% --- p. 8 ---
\opage{8}uden at vide deraf, mod det Væsentlige, og saa skal Dialektikeren strax være en Ordkløver.

En formildet Form af den samme Misforstaaelse giver sig tilkjende derved, at man vel indrømmer en Dialektik af de formelle og aprioriske, men ikke af de reelle, de empiriske Begreber. „Aprioriske Begreber, saasom: Noget og Andet, Væsen og Skin, Grund og Følge o.\ dsl., kunne ved Modsigelsens Anvendelse let frembringes af hinanden. Den, der sætter Grunden, underforstaaer jo allerede Følgen, Grunden er altsaa i Grunden baade sig selv og Følgen; men naar Grunden først paa Grund af Følgen bliver en egentlig Grund, saa er jo Grunden Følge, og Følgen Grund. Det Samme er ikke blot i een Henseende Grund, i en anden Følge, men de to Henseender falde tillige sammen i een Henseende: at dialektisere Begreberne ud af hinanden, er, om jeg maa bruge dette Udtryk, det Samme som at dialektisere \textit{diversi respectus} ind i hinanden. Hvad de aposterioriske Begreber, Erfaringsbegreberne, derimod angaaer, er det umuligt at dialektisere Existensen ud af Begrebet; eller hvem vilde vel forsøge at dialektisere den virkelige Hest udaf Hestens Begreb, eller den virkelige Snegl udaf Sneglens Begreb? Men er det umuligt at dialektisere det Existerende ud af dets Begreb, saa er det vel af samme Grund ogsaa umuligt at dialektisere det ene Existensbegreb ud af det andet, Hestebegrebet f.\ Ex.\ af Fluebegrebet, dette igjen af Sneglebegrebet. I Erfaringens Verden falde \textit{diversi respectus} udenfor hinanden, og hvor \textit{diversi respectus} falde udenfor hinanden, spiller Dialektiken en meget underordnet Rolle.“\footnote{Forelæsning over „Phil.\ Propæd.“ Universitetsaar 1860--61 S.~25.}% note *)

Hele Misforstaaelsen hæves, naar det indsees, at et Realbegreb kun kan dannes og udvikles paa den Betingelse, at Bestemmelser af det antilogisk Virkelige og det logisk Væsentlige articulere hinanden. Den existerende Syre f.\ Ex.\ er
% --- p. 9 ---
\opage{9}vistnok i Virkeligheden selvstændig og forsaavidt ligegyldig imod den existerende Base. At ville fra Virkelighedens Synspunkt paastaae, at det er Syren, der gjør Basen til Base, og omvendt, ligesom det i Causalitetsforholdet er Aarsagen, der gjør Virkningen til Virkning, og omvendt; at Syren altsaa selv er Base, og Basen Syre, ligesom Aarsagen selv er Virkning, og Virkningen Aarsag, er urimeligt og meningsløst: i det Antilogisk-Virkelige falde \textit{diversi respectus} udenfor hinanden. Men at det i væsentlig Forstand, altsaa naar Begrebet skal med som Begreb, absolut maa være Syren, der gjør Basen til Base, og omvendt, er ikke blot i Reflexionen uimodsigeligt, men bliver endog i Virkelighedsphænomenet antydet derved, at det samme Stof kan i en Forbindelse spille Syrens i en anden Basens Rolle: i den kiselsure Leerjord er Leerjorden Base, men i den leersure Magnesia er den Syre; Syren som Kiselsyre gjør altsaa Leerjorden til Base, medens Basen Magnesia gjør Leerjorden til Syre. Istedenfor at bebreide Philosophien dens dialektiske Ordforklaringer, burde man til Videnskabens Gavn omsider lære at indsee, at de dialektiske Ordforklaringer ere nødvendige til Realbegrebernes alsidige Udvikling.

Dersom der i Sandhed kunde gjøres en saa endelig Adskillelse mellem Ordet og Tingen, at Tingen lod sig opfatte som det reent Objective, Ordet som det reent Subjective, vilde Dialektiken unegtelig tabe sin høiere Gyldighed; men hvad er vel et Ords væsentlige Betydning, dersom det ikke er dets Objectivitet, og hvad er det for en Ting, hvis Objectivitet ikke lader sig opfatte og bestemme ved et tilsvarende Ord? Vel kan der være noget Subjectivt i Sprogbrugen, og det beroer da forsaavidt paa os selv, hvilken Betydning vi i en bestemt Sammenhæng ville tillægge et bestemt Ord. Men i al Sprogbrug er der ogsaa noget Objectivt; har et Ord engang erholdt en bestemt Betydning, da maa Betydningen respecteres. Nye Ord kunne dannes, men ikke vilkaarlig:
% --- p. 10 ---
\opage{10}Vexelvirkningen mellem Folkesproget og det videnskabelige Sprog er da ogsaa forlængst anerkjendt.\footnote{\textit{Ce sont les opinions du peuple et le point de vûe sous lequel il envisage les objets, qui donnent la forme au Langage. A mesure que le savoir et la politesse se répandent dans une nation, et à proportion de la durée de leur regne ils portent leurs influences dans la Langue. C'est qu' alors le peuple s'approprie plusieurs expressions que les savans ont inventées; comme ceux-ci adoptent souvent à leur tour, des expressions populaires.}

\textit{De l'influence des opinions sur le langage. Par Mr.\ J.\ D. Michaëlis. Traduit de l'Allemand. à Breme. MDCCLXII. S.~9.}}% note *)
Vel er Tingen uafhængig af Ordet, forsaavidt Ordet er en vilkaarlig Benævnelse; men en Ting saa objectiv, at der til dens Objectivitet ikke svarede noget Ord, er utænkeligt. Det Objective maatte da paa denne Maade forvandle sig til et Unævneligt, et i sig Ubestemmeligt, et Noget, der dog ikke var Noget, eftersom Noget jo er en Benævnelse, og det Unævnelige falder udenfor enhver Benævnelse.

For at Ordet skal kunne bestemmes ved Tingen, maa Tingen igjen bestemmes ved Ordet; for at Ordets Væsen kan være objectivt i Tingen, maa Tingens Væsen være subjectivt i Ordet. Tingen og Ordet ere altsaa i Væsen Eet og dog virkelig forskjellige: dette er Dialektikens Mening.

Ligesom Ordet, dialektisk seet, er Eet med og dog forskjelligt fra Tingen, saaledes er det ogsaa Eet med og forskjelligt fra Tanken. Sættes Ordet mod Tingen, saa bliver Tingen det Objective, Ordet det Subjective; sættes Ordet imod Tanken, saa bliver Ordet det Objective, Tanken det Subjective. Umiddelbart kunne Tanke og Ting aldrig enes; først med Ordet kommer Forklaringen og med Forklaringen en Forstaaelse. Som en almindelig Logoslære, en Lære om den rene Logos, er Logiken kun en Lære om Tanken, forsaavidt den tillige er en Lære om Ordet i dets væsentlige
% --- p. 11 ---
\opage{11}Eenhed med Tanken. Dog her, som overalt, er den væsentlige Eenhed naturligviis begrundet paa Forskjel og Modsætning; Læren om Ordets Eenhed med Tanken er just en Lære om Modsætningen af Ord og Tanke. Er det vildledende at ville fastholde Tanken i endelig Adskillelse fra Ordet, da er det ikke mindre vildledende at ville ligefrem identificere Ordet med Tanken. Den sidste Vildfarelse charakteriserer Speculationen, som overseer Ordets Umiddelbarhed og lægger eensidig Vagt paa, at Logos væsentlig er Tanken.

At Logiken, den rationelle Logoslære, fra Først til Sidst maa bruge Ordet, er da en Selvfølge; men Et er det at bruge Ordet, et Andet at reflectere paa Ordets Brug. Denne Reflexion indstiller sig just i nærværende Sammenhæng, hvor Opgaven er at udvikle de logiske Functioner. Abstrahere vi fra det Eiendommelige i Forholdet imellem Tanke, Ord og Ting, da bliver det umuligt at underkaste de logiske Functioner nogen særegen Behandling. Man maa da enten forudsætte Ordets Eenhed med Tanken og reflectere paa Subjectiviteten, eller omvendt Ordets Eenhed med Tingen og reflectere paa Objectiviteten. I første Tilfælde gaae Functionerne saa ganske op i det Almene, at en særlig Udvikling af deres egen indre Eiendommelighed og gjensidige Forhold bliver umulig. I sidste Tilfælde blive de objective Forskjelligheder saa paatrængende, at det ikke kan være de logiske Functioner selv, de subjective Handlinger, men Functionernes Logik, et System af objective Forholdsbestemmelser, der kommer til Udvikling. Vel have vi i Læren om den \emph{reflecterende Subjectivitet} afhandlet en Fleerhed af særegne Reflexionsformer: Forstandsreflexion, Væsensreflexion, Virkelighedsreflexion med tilsvarende Underafdelinger; men disse Reflexioner, ihvorvel de fra en Side maae opfattes som logiske Functioner, ere dog i en anden Henseende saa aldeles sammensmeltede med Indholdets saavel objective som subjective Formninger, at de umulig kunne træde frem i deres Forskjel fra
% --- p. 12 ---
\opage{12}Indholdet selv. Denne Forskjel melder sig derimod øieblikkelig, naar der reflecteres paa Ordet og Ordets Forhold til Tanken og Tingen.

I Forskjelseenheden af Ordet og Tanken sees en Forskjelseenhed af Reflexionen og Dommen: den Reflecterende dømmer, den Dømmende reflecterer; men Reflecteren og Dømmen ere derfor ikke det Samme. At Reflexionen, der altid svæver mellem modsatte Bestemmelser, og har sin Hvilen i denne Svæven, ogsaa maa hvile i en Svæven mellem det Subjective og Objective, og forsaavidt være ligesaa oplagt til at give Formbestemmelsen et objectivt som subjectivt Præg, er en Selvfølge. I Leddenes Speiling, i Momenternes Gjennemsigtighed er Reflexionen objectiv; men saaledes objectiv kan Dommen ikke være. I Dommen som Dom sætter den dømmende Subjectivitet en udtrykkelig Forskjelseenhed af Objectivt og Subjectivt, idet Dommens Gjenstand vel er objectiv, men Domsacten subjectiv; Forskjelseenheden bestemmes ved Ordet: i Dommens Subject sees Eenheden af Ordet og Tingen, i Prædicatet Eenheden af Ordet og Tanken, i Copula Tingens og Tankens Eenhed ved Ordet. Kun ved at reflectere paa Dommen som Dom, kan Tanken reflectere paa det i dens egen mangesidige Virksomhed Eiendommelige; men Udviklingen af det i Tankens formelle Virksomhed Eiendommelige falder jo netop sammen med en Udvikling af de logiske Functioner. Grundfunctionen er Dommen som Dom: al bestemt og bestemmende Tænken er væsentlig en Dømmen; denne Dømmen udtaler sig ikke blot i enkelte Domme, men i Forbindelser af Domme, ikke blot i Bestemmelser af Dommens formelle, men ogsaa af dens reelle Gyldighed. Dommen gaaer med indre Nødvendighed over i Slutning, Slutning over i Beviis. Det i den almene Grundfunction Almindelige articulerer sig da i tre særskilte Grundfunctioner: at \emph{dømme}, at \emph{slutte}, at \emph{bevise}.

I enhver af disse særskilte Grundfunctioner er Modsæt\opage{13}%
% --- p. 13 ---
ningseenheden af det Subjective og det Objective let at paavise; thi Modsætningsenheden bestemmes ved Ordet.

Den „formelle Logik“ tager Sagen subjectivt; den har forsaavidt Ret, som Functionerne: dømme, slutte og bevise nødvendig maae være subjective, absolut subjective Virksomheder. Men den formelle Logik har Uret, forsaavidt den opfatter det Subjective i en saa eensidig og udialektisk Adskillelse fra det Objective, at de tomme Tankeformer stilles op, som om de vare ligegyldige mod Indholdet. At en saadan Abstraction kun kan gjennemføres ved Hjælp af Ordet, er begribeligt; men idet den formelle Logik behandler Tænkningen som en psychologisk betinget Aandsvirksomhed, opfatter den tillige Ordet som en blot psychologisk Frembringelse, en i sig afmægtig Betegnelse, der falder udenfor Tingens Væsen. Ved en saadan Adskillelse af Tingen og Ordet og, hvad hermed falder sammen, af Tingen og Tanken, bliver Logiken ideeløs. Det er just det Betænkelige ved den formelle Logik, at den under Foregivende af at udvikle Tænkningens „Natur og Væsen“, i Virkeligheden kun driver det til en Udvikling af de triviale Abstractioners, de fade Bemærkningers og endeløse Raisonnementers „Natur og Væsen“.

Den „speculative Logik“ tager Sagen objectivt, ja saa objectivt, at den objective Dom maa dømme sig selv, den objective Slutning slutte sig selv, og Methodens objective Beviis omsider ende med at bevise sig selv. Den speculative Logik har den udødelige Fortjeneste, at have opløst de endelige Tanker i Ideens indre Uendelighed og sat Fordringen af en Dømmen, Slutten og Bevisen, der overskrider de psychologiske Betingelsers snevre Grændseskjel, i Kraft. Men naar Objectiviteten i den Grad absorberer det Subjective, at Tanken bliver Ord og Ordet Ting, saa reagerer Subjectiviteten, og opløser igjen det Objective i den Grad, at Tingen bliver Ord og Ordet Tanke. Kun en Tydeliggjørelse af Tanke, Ord og Ting i deres oprindelige Eenhed og Modsætning vil
% --- p. 14 ---
\opage{14}formaae at hæve den Mystification, hvori den speculative Logik har villet drukne Dualismen af det i sig selv Subjective og det i sig Objective; men en saadan Tydeliggjørelse maa da beroe paa en særlig Udvikling af de logiske Functioner.

Forskjellen mellem Ting og Tanke forvandler sig i den logiske Dom til en Forskjel mellem Subject og Prædicat. Subjectet er det Enkelte, det Enkelte i sin Umiddelbarhed, Dommens d.\ v.\ s.\ Domsactens Object, altsaa Tingen; Prædicatet, det Almene, nemlig det i Reflexionen Almene, er derimod Tanken, altsaa det Subjective. I den umiddelbare Dom bliver nu Tanken saaledes forenet med Tingen, at det seer ud, som om Tingen var Tanke, og Tanken Ting; det kommer af, at Tanken og Tingen enes i Ordet. I Domme som: dette er en Rose, denne Rose er rød, --- er Ordet Rose som Prædicat jo et Udtryk for Begrebet, altsaa for Tanken, som Subject derimod en Betegnelse af Objectet, altsaa af Tingen. Men efter Meningen maa Prædicatet Rose være ligesaa objectivt som Subjectet dette, og Prædicatet rød ligesaa objectivt som Subjectet Rose; thi Copula \emph{er} holder paa Objectiviteten.

Det er altsaa ikke en Deel af Dommen, men hele Dommen, der angiver det Objective; men just af samme Grund er det heller ikke en Deel af Dommen, men hele Dommen, der beroer paa den dømmende Subjectivitet. Den dømmende Subjectivitet bestemmer Objectet som Subject, for at dette Subject ved sit Prædicat kan blive objectivt bestemt d.\ v.\ s.\ bestemt i Qvalitet af Object. Domsacten er, fra denne Side seet, en Objectiveringsact. Den hele Dom er objectiv, thi den dømmende Subjectivitet bestemmes ved Objectet, og udtrykker Bestemtheden kategorisk, idet den gjør Dommens Object til det inhærerende Prædicats Subject. Den hele Dom er subjectiv, thi den dømmende Subjectivitet maa selv bestemme sit Object. Tingen som Ting er i sin objective Umiddelbarhed ligesaalidt et Subject, som den i Tingen tilværende Egenskab
% --- p. 15 ---
\opage{15}i sin Umiddelbarhed er Prædicat: kun i Objectiveringen bliver den existerende Ting til Subject formedelst Ordet; kun ved det for Objectiveringen afgjørende Ord kan den tilværende Egenskab forvandles til Prædicat, til det Almene, hvorunder Subjectet kan subsumeres.

Forsaavidt Dommen, den hele Dom, ifølge sin Natur er subjectiv, er her altsaa i endelig Forstand en Mulighed til, at den dømmende, sit Object oversvævende, Subjectivitet paa vilkaarlig Maade kan forbinde Subject med Prædicat: her kan da opkomme Tvivl om Dommens Gyldighed. Skal Tvivlen hæves, maa Gyldigheden begrundes; denne Begrundelse har igjen en subjectiv og objectiv Side.

Tages Begrundelsen fra den subjective Side, maa Forbindelsen mellem Subject og Prædicat beroe paa en Tankeudvikling og Tankeudviklingen paa de omspurgte Factorers formelle Sammenhæng. Til fuld Belysning af denne Sammenhæng fordres, at Subject og Prædicat træde ud som modstaaende Extremer; imellem Extremerne indskydes et Mellemled, der er saa dobbeltsidigt, at det i en Dom kan vise sig forenet med Subjectet, i en anden med Prædicatet: som en af de saaledes dannede Domme afledet Dom er den omtvistede Doms Gyldighed formelt begrundet; Begrundelsen er en Slutningsact, og den derved begrundede Dom en Slutning. At Slutningsacten er en logisk Function, en Tænkningsact, der ligesaa oprindelig beroer paa Ordet som den umiddelbare Dom, er da en Sag, der forstaaes af sig selv. Skulde der, hvad Hegelianerne paastaae, kunne gives objective Slutninger, Slutninger, som, uden al Medvirken af en tænkende, dømmende, sluttende Subjectivitet, foregik i den objective Fornufts objective Væsen: da maatte Tankerne i den objective Fornuft jo blive til Domme, der kunde dømme sig selv; de sig tænkende Tanker, de sig selv dømmende Domme, maatte da forvandle sig til tænkende og dømmende Subjectiviteter:
% --- p. 16 ---
\opage{16}hvilken Vrimmel af Subjectiviteter i den objective Fornufts
objective Væsen!

Tage vi Sagen fra den objective Side, maa Dommens
Gyldighed godtgjøres middelbart ved umiddelbar Paaviisning.
Under den umiddelbare Paaviisning bliver den dømmende Subjectivitet
og derved Domsacten, den dømmende Function,
bestemt ved Dommens Object: Dommen er objectiv. Med
denne Objectivitet træde nu de Domme, hvori ethvert af
Slutningens Extremer forbindes med Mellemledet, særlig frem,
den ene efter den anden, snart sideordnede, snart i Afhængighedsforhold;
dog altid saaledes, at den deraf resulterende
Slutning bliver bestemt som en objectiv gyldig Dom. Det
Umiddelbare bliver ved denne Slutningsmaade udstykket i
særskilte Led; Paaviisningen er umiddelbar: den umiddelbare,
ved sine Articulationer dog middelbare, Paaviisning er Beviisfunctionens
oprindelige Særkjende. Da Beviisfunctionen
imidlertid er subjectiv som en Function af Dom og Slutning,
saa kan det Middelbare ikke forblive som Udstykning og Sammenstykning
af det Umiddelbare, men maa opløse sig i Reflexionen.
Det Middelbare er her det Objective formedelst
det Umiddelbare; Objectiviteten er netop kjendelig paa den Vexel
af Middelbart og Umiddelbart, hvortil Paaviisningen svarer.
Reflexionen er det Subjective; Reflexionen forvandler Middelbarhedens
umiddelbare Led til Reflexer i Tankebevægelsen, til
en Totalitet af gjennemsigtige, for Indsigten bestemmende
Momenter; ved denne Forvandling bliver Paaviisningen Begrundelse;
i Begrundelsens og Paaviisningens reflecterede Eenhed
haves Beviset.

Sammenstille vi de tre logiske Functioner under eet
Synspunkt, er det klart, at de betegne en fremskridende Objectivering.
I Dommen vexler det Subjective umiddelbart
med det Objective; i Slutningen bliver det Objective Moment
for en formel Udvikling af det Subjective; i Beviset bliver
det Subjective Moment for en Udvikling af det i Viden Ob%
% --- p. 17 ---
\opage{17}jective; denne Fordobling er for hver Function betinget af
Ordet. At virkelige Objectivitetsbestemmelser kunne reflecteres
i subjective Former, skyldes Ordet som et \emph{Videns} Ord; at
subjective Formbestemmelser kunne overgaae i Indhold af virkelig
Existens, skyldes Ordet som et \emph{Magtens} Ord; kun
igjennem Modsætningen af Ordet i Viden og Ordet i Magt
og denne Modsætnings indholdsrige Eenhed i Sandhedens
Ord lade de logiske Functioner sig alsidig udvikle og bestemme.

Da Dommen som Dom udtrykkelig viser, hvorledes Tanken
former sig selv ved at udforme sit ideelle Stof i Ordet, ville
vi i det Følgende afhandle Domsfunctionerne under Benævnelsen:
\emph{Formalfunctioner}; da endvidere Slutningen
ifølge sit Begreb ikke føier nye eller særegne Bestemmelser
til Dommen i dens Enkelthed, men netop angiver den Maade,
hvorpaa den ene Dom deels forudsætter, deels medfører og
stadfæster den anden, opfatte vi Slutningsfunctionerne i det
Hele som \emph{Modalfunctioner}. Det eiendommelige Element,
Beviset som Beviis indfører i Slutningsformen, nemlig Realmomentet,
Ordets objective Reflex, leder saa endelig til
Indsigt i Sagens og Fremstillingens articulerede Vexelbestemmelse;
den Maade, hvorpaa det Eensidige under Beviisførelsen
gaaer op i det Alsidige, det Ufuldstændige suppleres,
det i Begrundelsen Mangelfulde integreres, tjener altsaa til
at vise \emph{Modalfunctionernes Realitet}.

\medskip
\noindent\hspace*{1em}a) \emph{Functioner af Dommen: Formalfunctioner.}
\medskip

Ligesom Tanken er en Function af Tænkningsacten, og omvendt,
saaledes er Dommen en Function af Domsacten: den
Tænkningsact, hvoraf den bestemte Tanke resulterer, er oprindelig
bestemt ved den resulterende Tanke, og ligesaa er den
ved Domsakten bestemte Dom bestemmende for Domsacten
selv. For den endelige Tænken og den endelige Dømmen maa
der naturligviis gives endelige Udgangspunkter, endelig bestemte
Forudsætninger; men de endelige Tænknings- og Domsacter
% --- p. 18 ---
\opage{18}ere kun betingede Yttringer af den væsentlige, med Selvobjectiveringen
identiske Dømmen, den almene med Subjectivitetens
Væren identiske Tænken. For den væsentlige Dømmen, den
uendelige Tænken, gives der ligesaalidt et første fast Udgangspunkt,
som for den absolute Viden eller den ontologiske Subjectivitet.

Adskiller man de enkelte Tankehandlinger fra Tænkningens
Princip, de særlige Domsacter fra den i alle Domme sig
yttrende „Dømmekraft“, da maa man enten leve hen i al
Umiddelbarhed, og forsaavidt i relativ Tankeløshed, eller tilsidst
støde paa et stort Problem, det Problem, der saa charakteristisk
plager Licentiaten i en Students Eventyr: „Seer De,
min Ven! En Bevægelse forudsætter en Retning. Forstanden
kan ei skride frem uden at bevæge sig i en vis Linie, men
inden den følger denne Linie, maa den jo i Forveien have
tænkt sig den. Altsaa har man jo bestandig tænkt enhver
Tanke, førend man tænker den. Enhver Tanke, der synes at
være Minuttets Værk, forudsætter jo saaledes en heel Evighed.
Det kunde gjøre mig næsten forrykt. Hvorledes kan da nogen
Tanke fødes, da den altid maa have været til, før den bliver
til“.\footnote{Efterladte Skrifter af Paul M. Møller. Første Bind. Tredie Udgave.
S. 120.}% note *)

Den uendelige Tanke endeliggjøres ved det Endelighedens
Stof, hvorover den tænker. Det er imidlertid ikke nok, at
Tanken i det Endelige har Noget at tænke over, den maa med
det Samme have Noget at tænke ved, Noget at tænke i: det,
hvori og hvorved der tænkes, er Sproget. Først naar Sproget
forvandler Sandsningens umiddelbare Billeder til Forestillinger,
bliver det Reflexionen muligt at forvandle Forestillingerne til
Tanker; kun ved en saadan Forvandling kan Tænkningen
blive opmærksom paa sig selv og sine Functioner. Ordet,
som ved at forvandle Billede til Forestilling tjener til at be%
% --- p. 19 ---
\opage{19}stemme Tanken, tjener da ved at omsætte Forestillingen i
Billede tillige til at bestemme Tankens Gjenstand. I Sætningen:
dette er et Træ, vækker Ordet ikke nærmest Forestilling
om sig selv som Ord, men Forestilling om det tilværende
Træ, altsaa om Gjenstanden.

Forsaavidt nu Dommen som Dom opfattes fra den subjective
Side, staaer Ordet i Tankens Tjeneste og bliver fra
dette Synspunkt at opfatte som et Videns Ord; Dommens
subjective Væsenhed er Form: alle under Dommens Form
hørende Functioner ere \emph{Functioner af Videns Ord}.
Forsaavidt Dommen betragtes fra den objective Side, altsaa
i Forhold til Gjenstanden, staaer Ordet i Gjenstandens Tjeneste.
Men at Ordet staaer i Gjenstandens Tjeneste, forudsætter
just, at Gjenstanden tillige staaer i Ordets Tjeneste; Ordet
er i denne Phase et Magtens Ord: alle til Forholdet af
Dom og Gjenstand objectivt svarende Functioner ere \emph{Functioner
af Magtens Ord}. En fuldstændig Udvikling af
Videns og Magtens Indhold i Sandhedens Ord kan imidlertid
ikke gives i Dommens Form; som særlig Grundfunction
har Dommen endnu sin Slutning udenfor sig. Dommens
Udvikling maa da ende med en Bestemmelse af dens subjective
Gyldighed, en foreløbig Slutning, som just indleder
Slutningen selv, idet den leder os til Indsigt i \emph{Functionernes
Modalitet}.

\medskip
\noindent\hspace*{1em}\textit{α)} \emph{Dommens Form: Functioner af Videns Ord.}
\medskip

Saavel det, hvorom der dømmes, som det, hvori og
hvorved der dømmes, er Ord; der dømmes om Ord ved Ord;
her gaaer Tanken i Eet med Sproget. Det Spørgsmaal,
hvorvidt Dyrene have Dømmekraft og „naturlig Logik“, hænger
paa det Nøieste sammen med det Spørgsmaal, hvorvidt de
have et egentligt Sprog. Det gjælder en Bestemmelse af Forholdet
mellem den dyriske Sandsebevidsthed og den menneskelige
Selvbevidsthed. Er Forskjellen mellem Dyrets og
% --- p. 20 ---
\opage{20}Menneskets Bevidsthed kun en Gradforskjel, saa er Forskjellen
mellem Dyrets og Menneskets „naturlige Logik“, mellem
Dyrets og Menneskets Sprog ogsaa kun en Gradforskjel. De,
der i Tillid til Dyrenes Perfectibilitet lægge Vind paa at
udvikle Dyrets naturlige Tænkeevne, maae da med det Samme
lægge Vind paa at udvikle og fuldkommengjøre Dyrets naturlige
Sprog. Vil det Sidste kun lykkes i meget indskrænket
Forstand, da vil ogsaa den første Bestræbelse kun give ringe Udbytte.
Dyrets Sprog har Lyd nok, men ingen Ord; fra den
meest betegnende Lyd, hvori Dyret udtaler sin øieblikkelige
Tilstand, til det logiske, i egentlig Forstand menneskelige, Ord
er der samme Afstand som fra dyrisk Sandsebevidsthed til
menneskelig Selvbevidsthed. Til nærmere Forstaaelse af Forholdet
mellem Tanke og Sprog tjener en særlig Udvikling
af Formalfunctionerne som \emph{sprogformende}, \emph{tankeformende}
og \emph{ideeformende Functioner}.

\textit{αα)} \emph{Sprogformende Functioner}. Hvad er Sproget?
Tankens ideelle Legemlighed, Mediet for Subjectivitetens Selvaabenbarelse,
det indre væsentlige Andet, hvori Aandens Lys
kan bryde sine Straaler. Som Tankens første umiddelbare
Aabenbarelse er Sproget vel et System af Tegn, men ikke af
vilkaarlige Tegn; her kan umulig begyndes med en særlig bevidst,
ved Overeenskomst bestemt, Adskillelse af Tegnet selv
og det Betegnede. Skulde Sprogdannelsen begynde med
vilkaarlige Betegnelser, da maatte Tanken forud for enhver
saadan Betegnelse være til som den bestemte Tanke, thi
kun den i sig bestemte Tanke kan vilkaarlig vælge sit Tegn;
men hvorledes skulde Tanken bestemme sig forud for Sproget,
naar Sproget er Ord, og Ordet selv Tankens første umiddelbare
Bestemthed. For at forme Sproget, maa Tanken virke,
for at virke, maa den være; forsaavidt gaaer vistnok Tankevirksomheden
forud for Sproget: Sproget er en Function af
Tanken. Men Tanken kan kun bestemme, forsaavidt den selv
% --- p. 21 ---
\opage{21}er bestemt; for at virke bestemmende paa Sproget, maa den
virksomme Tanke bestemmes ved Sproget: Tanken er i denne
Forstand en Function af Sproget. Som Function af Sproget,
idet Sproget er en Function af Tanken, er Tanken en Function
af en Function af sig selv: ved dette Kredsløb er Uendeligheden
betegnet.

Dog, vi tør ikke glemme, at den saaledes fremhævede Identitet
af Tanken og Ordet er, fra Philologiens Standpunkt, ikke
uden Antydning af selvbevidst philosophisk Sikkerhed, temmelig
skarp afviist. „Det fremsættes undertiden,“ siger Madvig,\footnote{Indbydelsesskrift
til Kjøbenhavns Universitets Fest i Anledning af
Hans Majestæt Kongens Fødselsdag, den 18de September 1842.
Om Sprogets Væsen, Udvikling og Liv. S. 5--6.}% note *)
„som en dybere Opfatning af Sproget, at det ikke er til for
Meddelelsens, men for Tankens egen Skyld og at Mennesket
ikke tænker uden i og med Sproget. Det Sande heri er, at
intet Menneske udvikler sig isoleret, men i et gjennem det tilværende
Sprog beredet Forestillingsfællesskab med Andre, hvoraf
følger, at Ordet, Fællesskabsmærket, er uadskilleligt knyttet til
de fælles Forestillinger, ledsager dem som allerede givne og
bevarer dem just i \emph{denne} Form og Articulation; hvorved dog
vel maa mærkes, at selv for den levende Forestilling kan det
bekjendte Ord stundom svigte. De nye Forestillinger, som den
Enkelte uddanner og afsondrer, vil han da stræbe at give
samme Gyldighed og Plads i Rækken, og dog fremkommer
denne Stræben først under Forsøget paa (mundtlig eller
skriftlig) Meddelelse som Søgen efter et nyt Ord \dots\ Selv
Nedskrivningen af et nyt Ord til Erindring om Forestillingen
for sig selv \emph{paa en anden Tid} har aldeles Meddelelsens
Form. Sproget taber Intet af dets væsentlige Charakteer
som Meddelelsesmiddel, fordi det tillige bliver et Middel til
sikkrere og lettere Beherskelse af de igjennem Meddelelsen
fællesgjorte og objectiverede Forestillinger. Den almindelige
% --- p. 22 ---
\opage{22}Tænkningen ledsagende Følelse af, at Forestillingerne kunne
udtales i Sprog og Ens egne Forestillinger just i Ens eget Sprog,
er ingenlunde en Tænken i og ved Ordet; thi Ordet er lydende.
(Al Tale med sig selv i virkelige Ord er en Følge af Meddelelsesdriften
eller en Art Prøve paa Ens egen Sikkerhed i
Forestillingsrækken eller paa dennes Evne til at udtales i de
Ord, man besidder)“.

To Indrømmelser skulle her strax gjøres: 1) „at intet
Menneske udvikler sig isoleret“, og 2) at Sproget er et „Meddelelsesmiddel“;
men den Paastand, at Sproget ikke skulde
være „for Tankens egen Skyld“, at „den almindelige Tænkningen
ledsagende Følelse af, at Forestillingerne kunne udtales
i Sprog“, ingenlunde er „en Tænken i og ved Ordet“, er i den
Grad vilkaarlig og stridende mod grundig Psychologie, at vi
ikke indsee, hvorledes en saa udmærket Mand som Forfatteren af
anførte Program har kunnet opstille en saadan Paastand. Naar
Spørgsmaalet er, hvorvidt der tænkes eller ikke tænkes i og
ved Ordet, maa her udtrykkelig sees paa den articulerede, i
tydelige Tanker sig udviklende Tænkning. Hvorledes bærer
man sig nu ad med at opfatte en Tanke tydelig og bestemt
uden at opfatte den i og ved Ordet? Man besinde sig blot
paa den første nødvendige Skilsmisse mellem Billede, Forestilling
og Tanke! Et Billede (Billedet af et Dyr, et Menneske,
en Sky o. s. v.) kan Bevidstheden fastholde uden Ordets
Hjælp, just fordi Billedet er en umiddelbar Efterligning af
Gjenstanden; men Et er Billede, et Andet Forestilling.
Billederne i Bevidsthedskredsen kunne adskilles og kjendes paa
samme umiddelbare Maade som Gjenstandene i Omverdenen,
men det kunne Forestillingerne ikke. Just idet det Billedlige
opløses, maae Forestillingerne, for ikke at flyde sammen, gjøres
kjendelige ved Tegn, ubilledlige Tegn, det vil sige ved Ord.
„Saaledes er Ordet \emph{Tegn} for en Forestilling, der \emph{kunde} have
faaet et andet Tegn, ligesom denne Lyd ogsaa \emph{kunde} være
% --- p. 23 ---
\opage{23}blevet Tegn for en anden Forestilling“.\footnote{Anførte Program 1842. S. 11.}% note *)
At Ordet som
Tegn selv er ubilledligt, at Lydsymboliken netop faaer sin
Betydning derved, at Ordet er og skal være „Tegn uden Efterligning“,
derom har Programmets Forfatter da ogsaa tydelig
nok udtalt sig i Overeensstemmelse med Hegel.

Hvad følger nu heraf? Ja, heraf følger uimodsigeligt,
at Forestillingernes, de ubilledlige, i Tanker overgaaende Forestillingers
indbyrdes Adskillelse og Forbindelse er saa oprindelig
betinget af Ordet, at den Forestillende og Tænkende
ligesaa vist trænger til Ordet, for at forstaae sig selv, som
for at gjøre sig forstaaelig for Andre. Det er altsaa Forestillingernes
og Tankernes Tydeliggjørelse i Bevidsthedskredsen,
ikke deres Meddelelse, hvorpaa Eftertrykket falder. Den Antydning,
hvormed Programmets Forfatter ledsager sin Paastand,
at Tænkningen staaer i et tilfældigt Forhold til Sproget,
efterdi „Følelsen af, at Forestillingerne kunne udtales i Sprog
og Ens egne Forestillinger just i Ens eget Sprog“, ingenlunde
er „en Tænken i og ved Ordet“, er høist besynderlig, den
lyder nemlig paa, at „Ordet er lydende“. Hvor høit skal
da nu Ordet lyde, for at være Ord? Dersom det er det
Sandselige i Lyden, der gjør Ordet til Ord, saa maa det
Ord, der udtales med Stentorstemme, altsaa være Ord i
høiere Forstand end det Ord, der hviskes. Med en saadan
Opfattelse stemmer det jo rigtignok, at det høitlydende Ord er
et kraftigere „Meddelelsesmiddel“ end det svagtlydende; men
hvor er nu Grændsen? Dersom det Ord, der hviskes, endnu
er et Ord, om end et svagt Ord, hvor svag tør da en Hvisken
være, uden at Ordet forsvinder? „Al Tale med sig selv“
maa, ifølge den opstillede Synsmaade, enten være hørlig,
hvad man hos Dannede gjerne anseer for et Tegn paa, at den
Talende gaaer lidt i Barndom, eller ogsaa en meget sagte
Hvisken. At tale med sig selv er altsaa at hviske til sig selv:
% --- p. 24 ---
\opage{24}hvor der ingen Hvisken er, er der ingen Ord, hvor der ingen
Ord er, er der intet Sprog, hvor der intet Sprog er, er der
ingen Talen.

Men hvad skulle vi nu sige om den, der sidder ganske
taus og læser i en Bog? Sidder maaskee en saadan Læsende
og hvisker, eller, dersom han ikke hvisker, ikke bruger Stemmen,
men kun Øinene, hvad er det da egentlig han læser? Ord
kan det naturligviis ikke være; dersom det, vel at mærke, er
bogstavelig at forstaae, at „Ordet er lydende“, at det er den
sandselige Lyd, der gjør Ordet til Ord.

Den i Programmet fremsatte Anskuelse er saameget besynderligere,
som Programmatarius selv i Forordet (S. 4)
har lagt fortrinlig Vægt paa „Hegels i al Korthed betydelige
Yttringer om Sproget (Encyclopædie, 2te Ausgabe, § 459),“
og ydet disse en udtrykkelig Anerkjendelse. Med denne Anerkjendelse
maa det imidlertid have en egen Sammenhæng; enten
kan Anerkjendelsen ikke gjælde hele § 459, eller ogsaa maa
Philologen ikke ganske have forstaaet Philosophen. Det er
nemlig saa langt fra, at Hegel i den anførte Paragraph er
gaaet ind paa den Synsmaade, at Ordet, for at være Ord,
maa være lydende, at han tvertimod vil have Lyden anseet
som noget for Ordet Uvæsentligt. Ja, han gaaer i denne
Henseende saa vidt, at han i den anpriste Paragraph endog
sætter Skriftsproget høiere end Talesproget. „Die Buchstabenschrift“,
siger Hegel, „ist an und für sich die intelligentere; in
ihr ist das \emph{Wort}, die der Intelligenz eigenthümliche würdigste
Art der Aeußerung ihrer Vorstellungen, zum Bewußtseyn gebracht,
zum Gegenstande der Reflexion gemacht“ \dots\ Hvor
Hegel vil hen med denne Udtalelse, er ikke vanskeligt at opdage:
han vil nemlig gjøre den udadgaaende Brug af Sproget
som „Meddelelsesmiddel“ til det Uvæsentlige, den indadgaaende
Anvendelse, hvorved Tanken articuleres, til det Væsentlige.
„Die erlangte Gewohnheit“, vedbliver han, „tilgt auch später
die Eigenthümlichkeit der Buchstabenschrift, im Interesse des
% --- p. 25 ---
\opage{25}Sehens als ein Umweg durch die Hörbarkeit zu den Vorstellungen
zu erscheinen, und macht sie für uns zur Hieroglyphenschrift,
so daß wir beim Gebrauche derselben die Vermittlung
der Töne nicht im Bewußtseyn vor uns zu haben
bedürfen; Leute dagegen, die eine geringe Gewohnheit des
Lesens haben, sprechen das Gelesene laut vor, um es in seinem
Tönen zu verstehen“.

Dette hører nu, som sagt, Altsammen til „Hegels i
al Korthed betydelige Yttringer om Sproget (Encyclopædie
§ 459)“; men at disse Yttringer indeholde den meest afgjørende
Protest imod det af Programmatarius angivne Synspunkt for
Sprogets Betydning som blot „Meddelelsesmiddel“, er vist.
Naar det da strax i Begyndelsen af Programmet hedder: „Det
fremsættes undertiden som en dybere Opfatning af Sproget,
at det ikke er til for Meddelelsens, men for Tankens egen
Skyld og at Mennesket ikke tænker uden i og med Sproget,“
saa spørger den opmærksomme Læser: hvilke videnskabelige
Fremtoninger mon det vel være, hvortil her sigtes? Efter det
umiddelbare Indtryk at dømme, skulde man mene, at det kun var
nogle dunkle Hoveder, gjærende Mystikere af ringe Rang,
til hvis Anskuelse der i det Hele ikke kunde tages videre Hensyn.
Til en saadan Opfattelse leder deels det Flygtige i Udtrykket:
„det fremsættes undertiden“, deels den Omstændighed, at den
hele Synsmaade afvises med et Par Bemærkninger uden
videre Undersøgelse eller Begrundelse. Imidlertid seer man
dog (S. 6), at det ikke blot er Phantaster som (S. 8) en
Schmidthenner eller K. F. Becker, men selv en Wilhelm
Humboldt, der holder paa Ordets og Tankens oprindelige
Eenhed. Maa man nu end indrømme, at Humboldt trods sine
store Fortjenester af Sprogforskningen kan i sin „Omhu for
Sprogets Værdighed“ her og der være saa uklar, at han af den
Grund ikke ret frembyder faste Holdningspunkter for motiverende
Kritik, da kan dette, hvad Grundspørgsmaalet angaaer,
dog paa ingen Maade gjælde om Hegel.

% --- p. 26 ---
\opage{26}Blandt de Tænkere, der udtrykkelig opfatte Sproget saaledes,
„at det ikke er til for Meddelelsens, men for Tankens egen
Skyld,“ staaer Hegel, hvis „betydelige Yttringer om Sproget“
Programmatarius selv i Forordet har anerkjendt, i første
Række. Ordene blive, ifølge Hegel,\footnote{Encyclopædie. Berlin 1845, § 462, S. 349.}% note *)
„zu einem vom Gedanken
belebten Daseyn. Dies Daseyn ist unseren Gedanken
absolut nothwendig. Wir wissen von unseren Gedanken nur
dann, --- haben nur dann bestimmte, wirkliche Gedanken, wenn
wir ihnen die Form der \emph{Gegenständlichkeit}, des \emph{Unterschiedenseyns}
von unserer \emph{Innerlichkeit}, --- also die
Gestalt der \emph{Aeußerlichkeit} geben, --- und zwar einer \emph{solchen}
Aeußerlichkeit, die zugleich das Gepräge der höchsten \emph{Innerlichkeit}
trägt. Ein so innerliches Aeußerliches ist allein der
\emph{articulirte Ton}, das \emph{Wort}. Ohne Worte denken zu wollen,
--- wie \emph{Mesmer} einmal versucht hat, --- erscheint daher als
eine Unvernunft, die jenen Mann, seiner Versicherung nach,
beinahe zum Wahnsinn geführt hätte. Es ist aber auch lächerlich,
das Gebundenseyn des Gedankens an das Wort für
einen Mangel des Ersteren und für ein Unglück anzusehen;
denn, obgleich man gewöhnlich meint, das \emph{Unaussprechliche}
sey gerade das Vortrefflichste, so hat diese von der Eitelkeit
gehegte Meinung doch gar keinen Grund, da das Unaussprechliche
in Wahrheit nur etwas Trübes, Gährendes ist, das erst,
wenn es zu Worte zu kommen vermag, Klarheit gewinnt“.\footnote{I
sit Skrift \textit{De l'origine du langage. Paris 1858. Préface.} S. 56
har E. Renan forsvaret Ordets oprindelige Betydning for Bestemmelsen
af Tanken mod H. Ritter. \textit{M. Ritter n'attribue au langage qu'un
seul rôle, celui de communiquer la pensée; il méconnaît une
autre fonction non moins importante de la parole, qui est de
servir de formule et de limite à la pensée.}}% note **)

At forklare Sprogets Oprindelse er med andre Ord at
forklare Selvbevidsthedens, den menneskelige Subjectivitets
% --- p. 27 ---
\opage{27}Oprindelse. Grundfunctionen er en ideel Objectivering. Kun
igjennem Objectivering af Andet kan Bevidstheden objectivere
sig for sig selv, kan fra Sandsebevidsthed hæve sig til Selvbevidsthed.
At Dyrets Bevidsthed kun er en Sandsebevidsthed
forudsætter, at den dyriske Subjectivitet ikke formaaer at hæve
sig over en blot sandselig Objectivering, en Objectivering,
hvori det Ideelle selv har Præg af sandselig Umiddelbarhed.
Dette sees deraf, at Dyret vel kan yttre sig ved en til Sandseindtrykket
svarende Lyd, men at et heelt System af dyriske
Lydmodificationer dog ikke vilde udgjøre et Sprog, da den
dyrisk sandselige Lyd ikke kan blive til Ord. For at Lyden
kan blive til Ord, maa den „\emph{løsrives fra hiin}“ --- Sandsebevidsthedens
--- „\emph{Umiddelbarhed og frigjøres af det
ikke dyrisk betagne, men forestillende Menneske}“;\footnote{Madvig.
Indbydelsesskrift til Kjøbenhavns Universitetsfest 1842.
S. 10. Tegnets Frigjørelse fra den første saa at sige umiddelbare
Umiddelbarhed er i dette Skrift paa det Eftertrykkeligste fremhævet
f. Ex. S. 11, hvor det hedder: „Mange krympe sig ved i Ordet
at erkjende et blot \emph{Tegn} uden reelt Aftryk af Forestillingens Væsen
og Charakteer; de ville ikke, som de sige, lade Sproget nedsynke til denne
nøgne og abstracte Skikkelse fra sandselig, malerisk Fylde og Farve“.
Hvortil Forf. bemærker, „at de spicifike, Ordet constituerende Lydarticulationer
% PRINTER: printed "spicifike", sense "specifike"
ikke have noget nødvendigt Forhold til Forestillingens
særlige Beskaffenhed“.}% note *)
altsaa: for at Lyden kan blive til Ord, maa Forestillingen om
Gjenstanden adskilles fra det blotte Sandsebillede, Forestillingen
om Fornemmelsen fra den blotte Fornemmelse, den indre
Objectivering fra den ydre. At nu den indre, ideelle Objectivering
er betinget af en Vexelvirkning mellem de Talende,
og at Sproget forsaavidt danner sig som et Meddelelsesmiddel,
viser kun, at den psychologiske Subjectivitet fra Begyndelsen
af er sandselig udadvendt, og kun ved udenfra, navnlig
fra andre Subjectiviteter kommende Paavirkninger bøies indad
til Selvreflexion og ideel Objectivering. Men heraf følger
% --- p. 28 ---
\opage{28}ingenlunde, at det sig under Meddelelsen dannende Sprog
blot skulde have Betydning som Meddelelsesmiddel; nei, heraf
følger just, at Meddelelsen, væsentlig seet, er Betingelse og
Forudsætning for Udviklingen af den indadvendte Bevidsthed,
en Udvikling, hvorunder Meddelelsesmidlet er at opfatte som
Objectiveringsmiddel.

Dersom Meddelelsen, og ikke den indre, ideelle Objectivering,
var Hovedsagen ved Sprogdannelsen, saa maatte den Talende
jo antages at kunne klare sine Forestillinger og Tanker paa
en anden Maade og ved andre Midler for sig selv end for
Andre, nemlig for sig selv uden Sprog, for Andre ved Hjælp
af Sproget; men dette er beviislig ikke Tilfældet. Den, der
i et uklart og forhutlet Sprog meddeler en Anden sine Tanker,
om ham gjælder Eet af To: enten have Tankerne i hans Bevidsthed
en uklar og forhutlet Form, eller, hvis de ere klare
og bestemte, tænker han dem i et Sprog og udtaler dem i et
andet. Det Sprog, hvori den Talende tænker Tankerne, er
ham da det væsentlige og overordnede; det, hvori han meddeler
dem, det uvæsentlige og underordnede. Den Talende
kan ikke objectivere sine Forestillinger for Andre, uden at
han med det Samme maa objectivere dem for sig selv; men
derimod kan han nok objectivere dem for sig selv, uden at
objectivere dem for Andre: Sproget kan altsaa være Objectiveringsmiddel
uden at være Meddelelsesmiddel, men ikke
omvendt.

Den Maade, hvorpaa Talen bestemmes, falder i Principet
sammen med den Maade, hvorpaa Tanken bestemmes.
Der tales i Sætninger, fordi der tænkes i Sætninger; det
lydende Ord i Talen og det lydløse Ord i Tanken er væsentligt
det samme Ord. Kan Ordet „Lyn“, naar det lyder, gjælde for
en Talesætning „det lyner“, da maa det samme Ord ogsaa,
uden at lyde, kunne gjælde for en Tankesætning. Den Subjectivitet,
der om det paagjældende Phænomen taler til en
Anden i Ordet „Lyn“, vil, naar Ingen er tilstede, tale til sig
% --- p. 29 ---
\opage{29}selv, objectivere sig det samme Phænomen ved det samme Ord;
kun med den Forskjel, at medens Tanken i det lydende Ord
er en Talen, er Talen i det lydløse Ord en Tænken. Indskrænket
til det lydløse Ord \emph{Lyn} alene er Tankesætningen
vistnok uden Articulation, Subjectet gaaer i Eet med Prædicatet,
Substantivet med Verbet; ikke destomindre er Forskjellen
mellem Subject og Prædicat, Substantiv og Verbum
paa ubevidst Maade underforstaaet. Det Samme gjælder,
naar Ordet lyder som Talesætning; uden en uvilkaarlig Adskillelse
af Noget, hvorom der tales, og det, der udsiges derom,
en Adskillelse, som Ordet vel ikke betegner, men foranlediger,
vilde her ikke være nogen Sætning,\footnote{„At Sproget begynder med Sætningen, er aldeles sandt i den Betydning,
at det er den stillede Opgave, hvortil Alt sigter og stræber:
thi Mennesket vil tale, ikke benævne isolerede Forestillinger; det er
fremdeles sandt, at enkelte Ord i Begyndelsen involverede Sætninger,
idet det forudsattes, at den Talende vilde udsige Noget,
ikke hensigtsløst henkaste et Navn. Men naar i den nyeste Tid
dette er blevet fremstillet, som om Sproget virkelig havde begyndt
sin Dannelse med den \emph{hele} Sætning, da ligger heri den besynderligste
Forvirring \dots\ Sætningen er kun heel, naar den indeholder
sine Bestanddele satte i deres Enkelthed og dernæst sammenfattede.
Vil man nu lade Sproget begynde med en saadan Sætning, da
antager man, at der ved det første Forsøg fremtraadte, forstodes
og fastholdtes to adskilte Ord og deraf det ene som Tegn for en
concret Subjectsforestilling (Gjenstand), og i det næste Forsøg, den
næste Sætning gjentager dette sig. Herfra vender man sig da til
Forestillingen om Sætninger, hvori Bestanddelene laae indviklede,
men siden udskiltes, ligesom af en Rodknold, det er, til Ikke-Sætninger,
en Uting, der hverken kunde udsiges eller forstaaes eller
kunne tænkes. Thi al Tale og al Forstaaelse gaaer ud fra, at den
Lyd, som man ikke kan opløse i særskilt betydende Bestanddele, er Tegn
for en enkelt Forestilling, og af denne udeelte Masse kan ingen Proces
udbringe adskilte, aldeles forskjellige Ord, eller af mange slige Sætninger
% note breaks here: continues at the foot of p. 30
altid det samme Ord for samme Forestilling; ere der derimod
særskilt betydende Lyd deri, da ere de der som Ord forud for Sætningen
eller vi falde tilbage i den ovenfor antydede Urimelighed.“ Madvig.
Anførte Program. S. 27--28. Renan synes imidlertid ikke endnu
at være bleven opmærksom paa dette vistnok meget vigtige Punkt,
men lader det blive ved temmelig almindelige Reflexioner, saasom:
\textit{Dans nos langues modernes, le sujet, le verbe, ainsi que plusieurs
des relations de temps, de modes et de voix, sont exprimés par
des mots isolés et indépendants. Dans les langues anciennes,
au contraire, ces idées sont le plus souvent accumulées dans un
mot unique et exprimées par une flexion. De l'origine du langage.}
S. 154.}% note *)
og naar her ikke var
nogen Sætning, da heller ingen Tale.

% --- p. 30 ---
\opage{30}Som Forskjelseenhed af Talen og Tænken er Sætningen
et Produkt af ideel Objectivering, og Objectiveringens Idealitet
kjendelig derpaa, at det Objectiverede undergaaer en Forvandling.
I Sætninger som: Hesten løber, Træet grønnes, er det ikke
de sandselige Gjenstande, Hesten og Træet, for sig, heller ikke
de almene Kategorier, løbe og grønnes, for sig, der blive objectiverede;
i den ved Sætningen bestemte Objectivering er
Hesten det Løbende, og det Løbende en Hest. Den fuldstændige
Sætning har sine særskilte Bestanddele, men bestaaer selv
i deres Sammenfatten til Eenhed. Objectiveringens Grundfunction
betinges altsaa af to særskilte Functioner: Bestemmelse
af Sætningens Led og Angivelse af den Maade, hvorpaa
Leddene ere forbundne. Ved Bestemmelse af Sætningens
Led kommer det væsentlig an paa Orddannelse; ved Angivelse
af Forbindelsesmaaden paa Formdannelse. De sprogformende
Functioner opfatte vi \emph{da} med Hensyn hertil fra to Synspunkter,
nemlig som \emph{orddannende} og \emph{formdannende},
dog saaledes, at den ene af disse Virksomheder uophørlig forudsætter
den anden.

\emph{Orddannende Functioner}. For at forklare Oprindelsen
til den Mangfoldighed af Lydforskjelligheder, et Sprog
behøver til Betegnelse af de mangfoldige Tilværelsesphæno%
% SEAM: word continues on p. 31 as "mener," -- written whole here as "Tilværelsesphænomener,""
% --- p. 31 ---
% SEAM: p. 30 ends "Tilværelsesphæno-"; this batch begins with the rest of
% the word, "mener,". The previous batch writes the whole word; drop the
% duplicated "mener" when joining.
\opage{31}mener, har man især lagt Vægt paa Naturmenneskets skarpe
Sands for sandselige Forskjelligheder og store Gaver til at
efterligne. Den orddannende Function kom da, ifølge denne
Synsmaade, oprindelig til at beroe paa en Efterligning,
især af Lydphænomener, saa at de onomatopoietiske Ord
maatte blive de egentlige Rodord (\textit{τὰ γὰρ ονόματα μιμητικα ἐστι.}
% PRINTER: printed "ονόματα μιμητικα" (no breathing on omicron, no accent on μιμητικα), sense "ὀνόματα μιμητικά"
Aristoteles). Men hvor stor en Rolle den uvilkaarlige
Efterligning end kan have spillet ved den begyndende
Orddannelse, saa maae vi dog give de Philologer
Ret, der nedsætte saavel de onomatopoietiske Ord som
Interjectionerne til underordnede Momenter.\footnote{% note *)
„I ethvert Tilfælde bliver ogsaa her Ordet først Ord ved at aflægge
denne Umiddelbarhed, just at betegne det, der lyder saaledes.
Herfra maa kun gjøres den bestemte og begrændsede Undtagelse,
som beroer paa selve Sprogmidlets Forhold til en særegen Side
af den sandselige Natur, det Lydende i den. Opfattelsen af
Lydforestillingen selv i dens individuelle Skikkelser (ikke som almindelig
Kategorie) udtalte sig naturligen igjennem de Lyd, der vare
Forestillingens Indhold og Gjenstand, og saaledes opstode de
\emph{onomatopoietiske} Ord. De ere saa langt fra at høre til Sprogets
Kjerne og Grundlag, at de, fordi de blive staaende i hiint
umiddelbare Forhold, paa en Maade ere uægte Ord, ofte selv i det
ellers uddannede og befæstede Sprog uderkastede legende Forandringer
% PRINTER: printed "uderkastede", sense "underkastede"
og Dannelser“. Madvig. Anførte Program S. 16. Men
Interjectionerne og de onomatopoietiske Efterligninger vilde jo være
fortræffelige Ord, dersom Sproget var et blot Meddelelsesmiddel. Blandt
de nyere Phil. har Renan især taget Ordet for Lydefterligningen.
\textit{La langue des premiers hommes ne fut donc, en quelque sorte,
que l'écho de la nature dans la conscience humaine \ldots\ La rupture,
par exemple, pouvait-elle s'exprimer d'une manière plus
pittoresque que par la racine ῥαγ (ῥήγνυμι, ῥήσσω, ῥώξ); sanscrit:
\emph{rug}; celto-breton: \emph{rogan}; ou par sa forme latine \emph{frac};
allemand: \emph{brechen}? \emph{Frem}, \emph{strep}, \emph{strid}, ne sont-ils pas
également la peinture naturelle du bruit dans ses diverses nuances?}
Anf. Skr. S. 136--37.} Det er ikke Efter%
% --- p. 32 ---
\opage{32}ligningen, det er tvertimod den Betegnelse, hvorved
Sandsebilledet bliver til Forestilling, en Betegnelse, hvorved enhver
Efterligning falder bort, der gjør Lyden til Ord; istedenfor
at være blot \emph{efterlignende} er den orddannende Function
actuelt \emph{bestemmende}.

At den orddannende Function er actuelt bestemmende,
maa ikke forstaaes, som om den selvvirksomme Subjectivitet i
den Forstand lagde Noget til af sit Eget, at den umiddelbart
frembragte ideelle Betegnelser for det Ideelle, særegne Ord
for Forhold, Benævnelser paa rene Kategorier o. dsl. En
saadan Antagelse er ligefrem stridende mod Sagens Natur.
Vel er det ikke det Sandselige i Lyden, der gjør Ordet til
Ord; men Ordet kan til sin første umiddelbare Bestemthed ikke
undvære et Element af sandselig Punktualitet. „At al
Sprogdannelse udgaaer fra Betegnelsen af sandselige Forestillinger,
er ikke blot det Resultat, hvortil alle Forskninger over Ordenes
Betydningers Sammenhæng selv i de meest udviklede Sprog
føre, heller ikke en blot Formodning, grundet paa en
Forudsætning om de første sprogdannende Menneskers Udvikling og
deres Forestillingskredses Grændser, men en Sætning, der
følger ligefrem af Sprogets Natur og dets nødvendige
Tilblivelsesmaade. Bekræftelsen af Forestillingernes
Fællesskab ved Hjelp af deres sandselig nærværende Gjenstand
var Udgangspunktet for den hele Grundovereenskomst, hvori
Sproget har sin Rod“.\footnote{% note *)
Madvig. Anførte Program. S. 18. Hvor conseqvent Forfatteren
fastholder og gjennemfører denne Synsmaade, sees blandt Andet af
følgende Udtalelse (Program 1856 S. 11). „Naar Humboldt
etsteds, kæmpende med sig selv, yttrer, at han dog ikke seer, hvorfor
der ikke kan gives oprindelige Ord for Forhold ligesaavel som for
Forestillinger, glemmer han, at Forudsætningen for Ordet er
Forestillingen (og oprindelig Forestillingen om det sandselig
Paaviselige) \ldots\ oprindelig og ifølge den første Forudsætning, hvorved Lyd
% note continues on p. 33
bliver gjort til Ord, bliver Alt, hvad der fremkommer i Sproget,
ved sin Fremkomst til Forestillingsnorm eller pronominalsk
Forestillingspaapegen.“} Den orddannende Function
er altsaa ved Betegnelser for det Almindelige ikke en
% --- p. 33 ---
\opage{33}selvstændig skabende, men en forklarende, idealiserende
Virksomhed.

Den Kjendsgjerning, at ethvert Ord i Sproget oprindelig
har havt en sandselig Betydning, er af Sensualisterne
naturligviis bleven misforstaaet i den Grad, at de have anseet
Subjectivitetens aprioriske Selvvirksomhed for en tom
Indbildning. Men Subjectivitetens Selvvirksomhed under
Sprogdannelsen er netop kjendelig af den Maade, hvorpaa det
Udvortes omsættes i sit tilsvarende Indvortes, det Sandselige
i sit Usandselige, det Individuelle i sit Almene.
Subjectiviteten kommer ikke abstract ved sig selv til sig selv; den maa
begynde med sit Andet, men ved Opfattelse af sit Andet,
det Objective, viser den just sit oprindelig subjective Væsen.
Saaledes kan Tænkningen heller ikke begynde med en abstract
Bestemmen af sig selv; den maa begynde med sit Andet, sin
Realnegation, det Sandselige. Saalænge Ordene endnu kun
have sandselig Betydning, er Tænkningen fortabt i
Sandsningen; men allerede den Kjendsgjerning, at Lyden er præget
til Ord og derved til Kategorie, beviser, at Tænkningen selv
i sit Andet har Tanken med. Kun efterhaanden som
Opfattelsen af Forholdene i det Udvortes articuleres ved Sproget,
foregaaer der en tilsvarende Articulation af det Indvortes.
Hvorfor den Talende ikke oprindelig vælger egne Ord til
Betegnelse af det Indvortes og Ideelle, er, forsaavidt Sproget
opfattes som Meddelelsesmiddel, klart af vort Foregaaende; men
hvorfor vælger den Tænkende nu ikke særegne, af alt
Sandseligt uafhængige Tegn, abstract ideelle Privattegn, til at
betegne Tankerne for sig selv? Svar: fordi det er umuligt.
Hvad skulde slige abstract ideelle Privattegn vel være? Egne
Forestillinger! men saa maatte disse Forestillinger, som ovenfor
% --- p. 34 ---
\opage{34}er viist, jo selv have deres Tegn. Særegne Anskuelser! men
hvad er en særegen Anskuelse uden et særeget Billede, uden
Skin af noget Sandseligt. Det er altsaa ikke blot for
Meddelelsens Skyld, men for Tænkningens og Tankernes egen
Skyld, at Ord for det Sandselige faae en overført Betydning
og anvendes til Betegnelse af det Usandselige.\footnote{% note *)
\textit{Ceux qui ont tiré le langage exclusivement de la sensation se
sont trompés, aussi bien que ceux qui ont assigné aux idées une
origine purement matérielle \ldots\ Le \emph{transport} ou la métaphore
a été de la sorte le grand procédé de la formation du langage.
Une analogie en a entraîné une autre, et ainsi le sens des
mots a voyagé de la manière en apparence la plus capricieuse;
souvent même la signification primitive a disparu, et n'a laissé
subsister que les acceptions dérivées, .. Le verbe \emph{être}, dont M.
Cousin disait hardiment en 1829: „Je ne connais aucune langue
où le mot français \emph{être} soit exprimé par un correspondant qui
représente une idée sensible“; le verbe \emph{être}, dis je, dans
presque toutes les langues, se tire d'une idée sensible. L'opinion des
philologues qui assignent pour sens premier au verbe hébreu \emph{haia}
ou \emph{hawa} (\emph{être}), celui de \emph{respirer}, et cherchent dans ce mot des
traces d'onomatopée, n'est pas dénuée de vraisemblance. En arabe
et en éthiopien, le verbe \emph{kâna}, qui joue le même rôle, signifie
primitivement \emph{se tenir debout} (\emph{exstare}). \emph{Koum} (\emph{stare})
en hébreu passe aussi dans ses derivés au sens \emph{d'être}
% PRINTER: printed "derivés", sense "dérivés"
(\emph{substantia}). Quant aux langues indo-européennes, elles ont
composé leur verbe substantif avec trois verbes differents: 1) \emph{as}
% PRINTER: printed "differents", sense "différents"
(sanscr. \emph{asmi}, ἐμμί, εἰμί, sum); 2) bhû (φύω, \emph{fui}, allem. \emph{bin},
persan \emph{bouden}); 3) \emph{sthâ} (\emph{stare}, persan \emph{hestem}), devenu
partie du verbe \emph{être}, au moins comme auxiliaire, dans les langues
modernes de l'Inde et dans les langues romanes (stato, été). De
ces trois verbes, le troisième est notoirement un verbe physique
et signifie se \emph{tenir debout}. Le deuxième a eu
très-vraisemblablement le sens primitif de \emph{souffler}. Quant au troisième, il
% PRINTER: printed "Quant au troisième" (a second "troisième"), sense "Quant au premier" (i.e. "as"; Renan's argument requires it)
parâit se rattacher au pronom de la troisième personne; mais ce
% PRINTER: printed "parâit", sense "paraît"
pronom lui-même, quelque abstrait qu'il soit, semble se rapporter à
un sens primitivement concret.} Renan. Anf. Skr. S. 122--130.}

% --- p. 35 ---
\opage{35}\emph{Formdannende Functioner.} Hvor inderlig
Bestemmelsen af Stof og af Form hører sammen i Sproget, sees
af Orddannelsen. Allerede Lydens første Articulation,
hvorved den bliver Betegnelse for en bestemt tilsvarende
Forestilling, forudsætter Formvirksomhed. Da det nu ovenikjøbet
kan være et Spørgsmaal, hvorvidt de særskilte Lydelementer,
hvoraf en Sætning dannes, i deres blot elementære
Forskjellighed ere at opfatte som oprindelige „Ord“ eller snarere som
„Rødder,“ hvoraf Ord i egentlig Forstand kunne dannes, saa
indsees heraf, at de i de indo-europæiske Sprog til bestemte
Klasser henhørende Ord, maae have gjennemgaaet en dobbelt
Formning.\footnote{% note *)
Die Sätze der sanskritischen Sprachen bestehen aus Wörtern; in
diesen sind die Wurzeln aufgehoben oder enthalten, etwa in der
Weise, wie ein chemisches Element (z. B. Sauerstoff) in einem
chemisch zusammengesetzten Körper (z. B. Wasser), und die Wörter
sind die Glieder des Satzes. Nirgends, auch im Englischen nicht,
treten Wurzeln, als solche, in der Rede auf. Denn es liegt im
Wesen und im Begriff der Wurzel, etwas Vereinzeltes zu sein; in
der Rede aber hat alles Zusammenhang. Die Wurzel ist also
allemal nur ein abstractes Product der Analyse. Sie kann zwar
unter Umständen lautlich unverändert und ohne Zusatz Glied der
Rede werden, und im Chinesischen tritt allerdings die Wurzel in
ihrer vollen Nacktheit in die Rede, dann genügt es aber schon, daß
sie mit andern Gliedern des Satzes, wenn diese auch selbst wiederum
bloß nackte Wurzeln sind, zusammengesprochen werden, um ihr
abstractes Wesen als Wurzel abzulegen und lebendiges Element
der Rede zu werden. Durch das Zusammenfassen zweier oder
mehrerer Wurzeln in einem bestimmten Verhältnisse hört die Wurzel
auf eine solche zu sein. Es ist also ungenau, zu sagen, der
chinesische Satz bestehe aus Wurzeln; denn in den Satz aufgenommen
verschwindet das Wesen der Wurzel. In Bezug auf letztere
unterscheiden sich also das Chinesische und die sanskritischen Sprachen
dadurch von einander, daß in diesen die Wurzel im Worte, in
jenem aber erst im Satze aufgeht. H. Steinthal. Charakteristik
der hauptsächlichsten Typen des Sprachbaues. Berlin 1860. S. 112.}

% --- p. 36 ---
\opage{36}Ligefrem og direct har den formdannende Function ikke
begyndt med at adskille Substantiver, Adjectiver, Verber, o. s. v.;
men allerede ved Dannelsen af Sætningen, i hvilken enhver af
Rødderne fik sin Function, var en saadan Adskillelse oprindelig
forudsat. „Naar man oftere har stridt, om Substantivet
eller Verbet var ældst i Sproget, viste der sig sædvanlig i
hele Striden nogen Uklarhed i Forestillingerne. I og for sig
var Verbalroden ligesaavel et Substantiv som Navnet for
nogen anden Forestilling, og ligesaa de oprindelige Adjectivrødder;
men da man i det begyndende Sprog ikke fastholdt
Virksomheden eller Egenskaben som abstracte Forestillinger og
talte om dem, kom Rødderne ikke til Anvendelse som Substantiver,
men bleve strax Verber eller Adjectiver, idet de
brugtes udsigende; thi dette er Verbet, Virksomhedsnavnet sat
som Prædicat \ldots\ Ved Betragtningen af Rodens Hensætten
som Verbum til et Subjects-Substantiv (naar dette var dannet)
maa man først komme paa det Rene med den saakaldte \textit{copula}
mellem Subject og Prædicat \ldots\ Denne \textit{copula} ligger som
den Act, hvorved Forestillingerne henføres til hinanden, bagved
Sproget og er Talens Forudsætning; i Sproget selv hverken
er eller kan den være direct udtrykt anderledes end ved
Talens Sammenhæng. Ethvert Tegn vilde selv trænge til et
nyt Tegn for at det henførtes; allermindst er det almindeligste
Prædicat, \emph{er}, et Tegn for \textit{copula}“.\footnote{% note *)
Madvig. Anf. Program 1842. S. 29--30.}

Ligefrem kunne Tanken og Ordet aldrig blive commensurable;
deres Eenhed er en Forskjelseenhed. Den bestemte,
sit bestemte Udtryk iboende Tanke er altid reflecteret i en oversvævende
Tænkning og kan kun i denne Reflexion først egentlig
fattes. Betydningen af et Ord være ved Sprogbrugen nok saa
bestemt, den vil dog altid ændres efter den Sammenhæng,
hvori Ordet forekommer, og det Samme gjælder om Sætningen.
Hvad der da ikke kan opnaaes direct, maa opnaaes
% --- p. 37 ---
\opage{37}indirect; men for at Opnaaelsen kan finde Sted, maa der
foregaae Articulationer af det Indirecte ved det Directe, og
omvendt; det er paa disse Articulationer, at de formgivende
Virksomheder særlig skulle kjendes. Overalt, hvor
Tankeformerne ikke chrystalliserede i et System af ufuldendte Betegnelser,
maatte Tilfredsstillelsen af de første elementære Fordringer,
paa Grund af Misforholdet, den altfor store Kloft mellem
det Directe og det Indirecte, uvilkaarlig fremkalde nye
Fordringer, disses Tilfredsstillelse igjen nye o. s. fr. „Sprogets
Opgave, til Opfattelse og Efterdannelse at fremstille de
Forbindelser af Enkeltforestillinger til Anskuelser, i hvilke det
tænkende Væsen bevæger sig, fandt ikke sin fulde Løsning i det,
der udgjør den sprogdannende Virksomheds første Resultat, de
Ord, der betegne Forestillingerne efter deres Indhold eller
paavise dem efter deres Relation (de benævnende og
pronominale Ord), idet den anskuende Bevidsthed indeholder Mere,
end der udtrykkes ved disse Ord og deres Opfatning efter
Talens ingen anden Antydning behøvende Grundforudsætning
om den prædicative Forbindelse til Udsagn, og forbinder flere
Forestillinger i een Anskuelse, end der efter denne Forudsætning
alene kunne opfattes i deres Sammenhæng. Til at bevirke
Opfatningen af det, der saaledes er Mere i den anskuende
Bevidsthed, tjene de grammatikalske Betegnelser (Ordstilling,
hjelpende Ord, Bøining af de benævnende Ord)“.\footnote{% note *)
Madvig. Om de grammatikalske Betegnelsers Tilbliven og Væsen,
1ste Stykke. Universitetsprogram. 6te October 1856. S. 7.}

Men just ved at reflectere paa Forholdet imellem det
Directe og det Indirecte i „de grammatikalske Betegnelser“
og mærke os „Motiverne“, saaledes som den skarpsindige
Philolog selv nøiagtig angiver dem,\footnote{% note **)
Anf. Program S. 7--8.} faae vi jo Punkt for Punkt
Beviset for, at Sproget er uadskilleligt fra Tænkningen, og
at Sprogudviklingen just af den Grund --- ikke for
Meddelelsens Skyld --- maa falde sammen med den tænkende
% --- p. 38 ---
\opage{38}Subjectivitets Selvudvikling. „De \emph{Motiver}, der enten
kræve eller kunne foranledige grammatikalske Betegnelser, lade
sig efter Anskuelsesformens Natur henføre til bestemte
Hovedklasser, idet de enten (A) angaae \emph{Enkeltforestillingens
Omfang og den Form}, hvorunder den i bestemt Function
(som Substantiv, Verbum, Adjectiv) optages i Anskuelsen i
Forhold til Ordets umiddelbare Betydning (saasom Betegnelser
for Substantivernes Tal, for deres Opfatning som individuelt
bestemte eller ubestemte, for Adjectivforestillingens Grader, for
Verbets Brug i passiv eller paa lignende Maade modificeret
Betydning), eller (B) \emph{Forestillingernes Forhold i Forbindelse}
til Heelhed i Anskuelsen (\emph{Forhold i Sætningen}), hvilket
Forhold deels (a) er prædicativ eller attributiv Forbindelse
(Betegnelser ved Prædicatet eller Attributet efter dets Henhøren
til Subjectet, \emph{prædicative eller attributive
Betegnelser}), deels (b) et Bestemmelsesforhold af de
substantiviske Forestillinger til og i Udsagnet enten som omprædicerede
eller trædende til i den rumlige Anskuelses Form (\emph{specielle
Forholdsbetegnelser}. Casus, Præpositioner), eller endelig
(C) \emph{hele Anskuelsens Stilling for Bevidstheden}
med Hensyn til Virkelighed, Præsents og Sammenhæng med
andre Anskuelser til sammensat Anskuelse (\emph{Sætningens
Forhold}; Betegnelser for Maade og Tid, Conjunctioner!“\footnote{% note *)
Madvig. Anf. Program 1856. S. 8.}
% PRINTER: printed "Conjunctioner!“", sense "Conjunctioner)“" (the parenthesis is never closed)

Det er igjen Grundforholdet, det psychologisk-ontologiske
Forhold mellem Billede, Forestilling og Tanke, hvorpaa det
væsentlig kommer an, naar Spørgsmaalet er, hvad der kan
udtrykkes direct, og hvad indirect. Ligesom Forestillingen er
det i logisk Forstand idealiserede Billede, saaledes er Tanken
igjen den idealiserede Forestilling. Jo mindre Idealitet en
Forestilling har, desto mere direct kan Fremstillingen være;
omvendt, jo mere abstract ideel Forestillingen er, desmere
indirect er den fra Begyndelsen af indeholdt i den directe
Fremstilling af det Concrete. „Den, der taler om \emph{Mænd}, har
% --- p. 39 ---
\opage{39}ikke en særlig abstract Forestilling om Fleerhed, men han har
og forlanger opfattet Forestillingen \emph{Mand} under Fleerhedens
Form; den, der vil tale om en Hests Løb som forbigangent,
har ikke for Bevidstheden en abstract Forestilling om
Forbigangenhed, heller ikke vilde han, hvis han havde den, tale om
den, som om Hesten og Løbet; ved at \emph{nævne}
Forbigangenheden, vilde han forvirre hele Meningen; den, der siger:
„Manden ligger i Huset,“ har ikke en Forestilling om \emph{Væren}
i fremhævet i Bevidstheden eller vil stille den ved Siden af
Forestillingen om Mand og Liggen som en tredie; men begge
have Fornemmelsen af en Skuen af de benævnte Forestillinger
under en særegen Form, den, der, naar den fastholdes ved
Abstraction og gjøres til Gjenstand for Omtale, vil kræve et
Navn som \emph{Forbigangenhed} eller \emph{Væren} i, og tillige et
Forlangende om, at Tilhøreren skal efterdanne og skue de
benævnte Forestillinger (i første Tilfælde: „Hesten løb“, hele
Udsagnet) under samme Form“.\footnote{% note *)
Madvig. Anf. Program 1856 S. 11.} --- Den, der taler om \emph{Mænd},
gjør altsaa, være sig bevidst eller ubevidst, en Forskjel imellem
det, der skal forestilles, og det, der skal tænkes. Skikkelsen
\emph{Mand} skal forestilles, \emph{Fleerheden} skal tænkes: idet
Forestillingen \emph{Mand} opfattes „under \emph{Fleerhedens} Form“, bliver
det Directe bestemt ved det Indirecte, d. v. s. Forestillingen
bestemmer Tanken og Tanken Forestillingen; begges gjensidige
Bestemmelse gives i og ved Sproget, og kan ikke gives ved
noget andet Middel. Vil Nogen forsøge at angive et andet
Middel, da er Resultatet Eet af To: enten er Midlet
uadæqvat, og da forfeilet, eller adæqvat, og da et Sprog. Hvad
der i Talen udtrykkes indirect, er ligesaa vel indirect for den
Talende selv som for Tilhøreren.

Forsøger man at sætte Adskillelsen mellem Tanke og
Ord deri, at den ene Talende kan have en mere udviklet
Bevidsthed end den anden, uagtet de begge sige det Samme,
% --- p. 40 ---
\opage{40}og altsaa, hvad Meddelelsen angaaer, staae paa lige Trin: da
vil ogsaa dette Forsøg strande paa Philologiens Klipper. Lad
To meddele sig i den samme Sætning: „Hesten løb“; lad den Ene
have, „en abstract Forestilling om Forbigangenhed,“ den Anden
ikke: er det nu herved godtgjort, at Sproget som Meddelelsesmiddel
staaer i et tilfældigt Forhold til Tanken? Paa ingen
Maade. For det Første have vi jo Philologens udtrykkelige
Erklæring, at den, der, idet han udtalte Sætningen: „Hesten
løb,“ havde en abstract Forestilling om Forbigangenheden, dog
ikke vilde indblande den i Talen, for ikke at forvirre Meningen.
Men forvirrer han Meningen ved at indblande den uvedkommende
Forestilling, da forvirrer han den jo ikke blot for
Tilhøreren, men ligesaa meget for sig selv. Dernæst, og dette
maa eftertrykkelig fremhæves, kan Ingen have en abstract
Forestilling om Forbigangenhed og fastholde den i en Abstraction,
uden at han, netop for Abstractionens Skyld, maa have
et tilsvarende Ord. Forbigangenheden er en ideel Forestilling;
dens Væsen er negativt; den kan ikke anskues, den kan kun
tænkes, og Tanken bestemmes ved det ubilledlige Tegn, ved
Ordet.

Det, vi i denne Sammenhæng skulle indsee, er, at Sproget
oprindelig „er til for Tankens egen Skyld,“ og „at Mennesket
ikke tænker uden i og med Sproget“. Hvilket af de for
grammatikalske Betegnelser anførte Motiver vi end betragte, saa
komme vi dog til det Resultat, at Tanken bestemmes ved
det Sprog, den selv bestemmer, at Tanken formedelst Sproget
er en Function af en Function af sig selv.

See vi hen til Betegnelser for „Enkeltforestillingens
Omfang og dens Form,“ altsaa til Motivet (A), hvorfor er her
da ikke i Sproget, ved Siden af Navnet for Enkeltforestillingen
selv, et Navn for dens Omfang og et eget Navn for dens
Form? Bestemmelserne: Omfang og Form, skulle, hedder det,
ikke forestilles direct, men indirect, det vil sige, de skulle ikke
forestilles, men tænkes. Desuagtet kunne slige Bestemmelser
% --- p. 41 ---
\opage{41}ikke i den sprogdannende Bevidsthed begynde som klare Tanker;
de maae tvertimod begynde som dunkle Antydninger.
Hvorledes hæves de dunkle Antydninger da op til Bevidsthedens
Lys? Derved, at de som flygtige Forestillinger (i dette
Tilfælde Biforestillinger om et bestemt Omfang, en bestemt
Form) hæfte sig ved en Gjenstandsforestilling; denne
Tilhæftning er kun mulig ved Ordet. Det er ikke Mennesket
som Individ, men Mennesket som Ord, hvortil der kan føies
en Artikel; det er ikke Individet Mand, men Ordet Mand, der
ved en Omlyd gaaer over i Pluralis; men er Individet ombyttet
med Ordet, saa er den dunkle Antydning af \emph{bestemt} eller
\emph{ubestemt}, af \emph{enkelt} eller \emph{flere} uvilkaarlig med i
Biforestillinger. Idet nu disse Hovedforestillingens Ord ledsagende
Biforestillinger mærkes som væsentlige, maae de som indirecte
Forestillinger bestemmes enten paa direct eller indirect Maade.
Først naar Betegnelsen ved Sproget har sikkret de indirecte
Forestillinger, blive de Gjenstand for Abstractionen og ved
Abstractionen for den bestemmende Tænken. Uden Sprog
vilde der ikke være Bevidsthed om Forskjel imellem et bestemt
og et ubestemt Individ, mellem Enkelttal og Fleertal, mellem
Activ og Passiv; uden Sprog vilde der for Tanken hverken
være Bestemthed eller Ubestemthed, efterdi Tanken Bestemthed
ligesaa vist forudsætter Ordet, som Ordet Bestemthed
forudsætter Tanken.

Kaste vi et Blik paa Motivet (B), skal vist Ingen negte,
at Forskjellen mellem Subjectet og dets Prædicat, mellem
Substantivet og dets Attribut, er en Begrebsforskjel, altsaa en
Tankeforskjel. Uden en saadan Forskjel vilde den tilsvarende
Tanke ikke være, men uden et tilsvarende Sprog vilde
Forskjellen selv ikke være. Tingen, antilogisk som Ting, er hverken
Subject eller Prædicat, hverken Substantiv eller Attribut; kun
i Sproget er der Subject og Prædicat, kun i Sproget er der
Substantiv og Attribut, kun i Sproget kan der tænkes.

Hvad endelig Motivet (C) angaaer, hvor er det da muligt
% --- p. 42 ---
\opage{42}at tænke sig „en Anskuelsens Stilling for Bevidstheden,“ en
af Anskuelser sammensat Anskuelse, altsaa en Anskuelse af
Anskuelser, uden Ord, uden Sprog? En Anskuelsens Stilling for
Bevidstheden er jo noget ganske Andet end en Gjenstandens
Stilling for Bevidstheden. En Anskuelse af en Anskuelse, en
Forestilling om en Forestilling er ikke som et Billede af et
Billede, et Syn paa et Syn, i sandselig Umiddelbarhed; her
er en Fordobling, der beroer paa Biforestillingers Forvandling
til negative Formbestemmelser, Bestemmelser af Billedligt ved
Ubilledligt, af Ubilledligt ved Billedligt, Reflexionsbestemmelser
af Ord ved Sætning og af Sætning ved Ord, af
Hovedsætning ved Bisætning o. s. v. „I den Magt, hvormed de
enkelte Led af Tanken samtidig trænge sig frem, ligger
Opfordring til at udvide Sætningsforbindelsen ved at knytte
flere Bisætninger til samme Hovedsætning og Bisætninger
til Bisætninger, en Opfordring, der bliver stærkere, jo rigere
Tankeindholdet, jo livligere og mere omfattende Meddelelsen,
jo større Trangen til skarp Begrændsning af den enkelte Tanke
ved Bibestemmelser bliver“.\footnote{% note *)
Madvig. Program 1856. S. 39.} --- Den, der erkjender dette og
endda antager, at der kan tænkes bestemte Tanker uden Ord,
uden Sætninger, følgelig uden alt Hensyn til Forskjellen mellem
Hovedsætninger og Bisætninger, maa virkelig, dersom hans
Antagelse skal gjælde for Mere end en tom Paastand, paatage
sig philologisk at vise os, hvorledes en Subjectivitet bærer sig
ad med at tænke enten en Tanke uden Sætning, eller en
Sætning uden Ord, eller en Sætning i Ord uden Sprog.

\textit{ββ)} \emph{Tankeformende Functioner.} Idet Tanken
bearbeider sit Objectiveringsmiddel, bearbeider den væsentlig
sig selv; de tankeformende Functioner kunne da heller ikke
holdes i en Rubrik for sig ligeoverfor de sprogformende.
Men Tanken er forskjellig fra Ordet som det, der betegnes,
er forskjelligt fra Betegnelsen; den Opgave, at uddanne et
% --- p. 43 ---
\opage{43}System af Betegnelser, er altsaa forskjellig fra den Opgave at
udforme og gjennemklare det i Betegnelserne reflecterede
Indhold: med Hensyn paa den første Opgave ere Functionerne
sprogformende, med Hensyn paa den sidste tankeformende.
Det er for Psychologie og Logik af stor Betydning, at
Philologien i vor Tid har naaet det Vendepunkt, at den med fuld
Bevidsthed er istand til at behandle de sprogformende Functioner
for sig, uden Indblanding af uvedkommende Speculationer.

„Spørgsmaalet om, \emph{hvad} og \emph{hvormeget} af Tanke og
Forestilling der ligger i Sprogets grammatikalske Betegnelser“,
siger Madvig,\footnote{% note *)
Program 1856. S. 9.} „og hvorledes det ligger deri, ad hvilken Vei
det er knyttet dertil, paatrænger sig som saa vigtigt for den
videnskabelige Fremstilling af grammatikalske Begreber og dets
Afgjørelse er en saa væsentlig Forudsætning for Gyldigheden af
Betragtninger og Domme over Sprogenes grammatikalske
Fortrin og Mangler, deres grammatikalske Bygnings Forhold til
Aandslivet, at man skulde vente, at der for begge disse Hensyns
Skyld var arbeidet ivrigt paa at klare det, at man i det Ringeste
havde søgt at stille det og besvare det ret bestemt“. At Opgaven
virkelig lader sig stille og besvare bestemt, har Forfatteren af
Programmet fra 1856 viist ved at angive og bestemme
Motiverne. Men da Sproget „er til for Tankens egen Skyld“,
da Mennesket „ikke tænker uden i og med Sproget“, maae de
sprogformende Functioner ligesaa vist reflectere sig i de
tankeformende, som disse i hine. Her kan ikke være Tvivl om,
\emph{at} det skeer; men Spørgsmaalet er, \emph{hvorledes} det skeer.

Det er, eller har idetmindste tidligere været, en temmelig
udbredt Mening, at de gamle classiske Sprog, Græsk og
Latin, især Latinen, skulde udmærke sig ved streng Logik, og
at den strenge Logik skulde lægge sig for Dagen i deres
Grammatik. Men hvormegen Uklarhed der kan indblande sig i
slige Bedømmelser, kan Philologien selv allerbedst godtgjøre.\footnote{% note **)
„Stundom henvises man til de fine Nuancer og Vendinger, hvoraf
hvert Sprog har sine, og som, saaledes som det ovenfor er
% note continues on p. 44
udviklet, mere vidne om naiv Modtagelighed end om conseqvent Logik;
men almindeligere paaberaabes den strenge Conseqvents, hvormed
Sætningens Dele sammenføies i nøie Overeensstemmelse med og
Afhængighed af hverandre, Prædicatet rettende sig efter Subjectet i
Person, Tal, Kjøn, Substantiverne indordnede i deres bestemte
Casus o. s. v., med andre Ord, det synes at opfattes som særlig
logisk, at Sproget iagttager sine egne Regler og anvender sine
Betegnelsesmidler hvert efter dets Betydning, som om de simplere
Sprog ikke indenfor deres Systems Omraade gjorde ganske det
Samme“ \ldots\ „Til denne Forestilling om det mere Logiske i de
gamle Sprog sluttede sig dernæst en uklar Overførelse af andre
Egenskaber, som man, hver efter sit Skjøn, fandt hos de gamle
Folk og ved deres Litteraturer, paa Sprogene, tildeels fordi man
af Sprogenes egen grammatikalske Beskaffenhed vilde bevise, hvad
der kun af de store historiske Culturforhold lader sig bevise, at
Kundskab i de gamle Sprog er af den største og meest udbredte
Vigtighed for Dannelsen, eller fordi man overhovedet søgte
Anbefalinger for det Sprogstudium, man drev. Medens Latin
beholdt Anseelsen for særlig logisk Bestemthed, søgtes andre
Prædicater for det Græske. Det er muligt, at der ved en
Sammenligning mellem de to klassiske Sprog indbyrdes i reen grammatikalsk
Henseende, som forresten, naar Momenterne skulle rigtig og ublandet
afveies, er yderst vanskelig, vil vise sig nogen Overvægt i
Modtagelighed for Bimotivers (Perturbationers) Indflydelse paa
Hovedanalogien i Græsk, (see Fortalen til den tydske Udgave af min
græske Syntax. S. XI); meget stor vil den neppe være, og en
Slutning deraf til hele Folkets aandelige Disposition vil være
saameget usikkrere, som ikke alle særlige i Udviklingsvilkaarene
liggende Indflydelser kunne medtages“. Madvig. Om de
grammatikalske Betegnelsers Tilblivelse og Væsen. Andet og sidste Stykke.
Universitetsprogram 1857. S. 83--86.}
% --- p. 44 ---
\opage{44}„Sproget“, siger Madvig, „har ikke blot en grammatikalsk
\emph{Form}, men det har ogsaa indenfor denne Forms Grændser
et Standpunkt, en \emph{Tilstand}, en raa ubearbeidet Tilværelse
og en høiere Udvikling, en tilveiebragt stilistisk Evne, og heri
ligger baade et \emph{Vidnesbyrd} om, at der er arbeidet meer
% --- p. 45 ---
\opage{45}eller mindre i og med Sproget, at der har været aandelig
Bevægelse hos Folket, og en \emph{Beqvemhed} for fremtidigt
Arbeide. Her, ikke i Grundformen, ikke engang i de
Retninger og den Stilling, der under Udviklingen er givet enkelte
Led af Grundformen, men i dens fyldige og bøielige, og ikke
blot i enkelt Retning bøielige, Udvikling ligger Sprogets
grammatikalske Fortrin \ldots\ Men medens det af den græske
Grammatik, navnlig Syntaxen, kan skjønnes, at Grækerne
have havt og pleiet en udførlig Meddelelse i Tale og Skrift,
at Folkets Medlemmer, i det Ringeste mange af dem, have
været istand til at bevæge sig i store og sammensatte
Forestillingsgrupper, og medens det af Metriken, som Deel af
Grammatiken, kan sees, at de have havt en i sandselig skjønne
og mangfoldige rhytmiske Former indklædt Poesie, kan der
af Sprogets grammatikalske Bygning slet ikke drages
nogensomhelst videre Slutning enten til Indholdsfylden af den i
Sproget pleiede Meddelelse eller til den Tænkningens Retning
(f. E. mod det Ideale eller det Materielle o. s. v.), som deri
har gjort sig gjeldende. Her indtræder som det langt rigere,
langt bestemtere sproglige Vidnesbyrd om hvad der er
foregaaet hos Folket, om man vil, om dets Verdensanskuelse (det
vil sige, om de Begreber og Forestillinger, der \emph{overhovedet}
have været i Bevægelse og ere fremkomne hos Folket)
Ordforraadet og Betydningsforraadet i Ordene. Det er her, at
det mærkes, at ikke blot Forhandlinger, som de, Demosthenes
førte, men ogsaa som de, for hvilke Platon og Aristoteles stode
i Spidsen, have fundet Sted hos Grækerne, at alle
Aandslivets Egne ere gjennemvandrede af dem. Det Lexikalske, det
Stofagtige, det Materielle i Sprogene, som man stundom
kalder det, indeholder det eneste, indtil en vis Grad bestemte
Afbillede af Folkets \emph{Aandsretning}“.\footnote{% note *)
Madvig. Program 1857. S. 88--89.}

% --- p. 46 ---
% p. 46 opens a new paragraph (p. 45 ends with a full stop and footnote).
\opage{46}Men saa klar, bestemt og slaaende denne Udtalelse end ved første Øiekast synes at være, saa er her dog et Punkt, og det et Hovedpunkt, der vækker Betænkeligheder. Det skal villig indrømmes, at der „af Sprogets grammatikalske Bygning ikke kan drages nogensomhelst videre Slutning enten til Indholdsfylden af den i Sproget pleiede Meddelelse eller til den Tænkningens Retning, som deri har gjort sig gjeldende“;
men --- maae vi spørge --- gaaer det nu ogsaa an at gjøre en saadan Adskillelse af det Grammatikalske og det Lexikalske, naar Talen er om Sprogudviklingens Forhold til Aandslivet? Det er dog vist ligesaa umuligt at læse sig til det bestemte Afbillede af et Folks Aandsretning i Sprogets Lexikon som i dets Grammatik. I begge, det Lexikalske saavelsom det Grammatikalske, naar hver af disse Hovedsider betragtes for sig, seer man jo kun en Mængde Betingelser, men ikke deres Anvendelse. Først i Anvendelsen, i den bestemte Fremstilling af et Aandsindhold, viser det sig, hvad et Sprog i aandelig Henseende formaaer. I Fremstillingen virker det Grammatikalske Punkt for Punkt sammen med det Lexikalske; det er i denne articulerende Samvirken, at Vexelvirkningen af sprog- og tankeformende Functioner egentlig skal kjendes.

At det ikke blot er det Lexikalske, men ogsaa det Grammatikalske, der i ethvert Sprog indvirker paa Tankebevægelsen, er især paafaldende i de Tilfælde, hvor man tænker i et Sprog og meddeler sig i et andet; men det kan desuden bogstavelig vises ved Forsøg med Oversættelser især af logiske Skrifter. Man forsøge engang at oversætte Hegels Logik paa Fransk, og man vil overbevise sig om, at Forandringer i det Lexikalske ogsaa medføre Forandringer i det Grammatikalske, at de herved motiverede Omskrivninger foraarsage Ændringer i Tanken, idet de foraarsage Forandringer i Fremstillingsmaaden, og at de til den forandrede Fremstilling svarende Tankeforandringer indeholde logiske Misviisninger, som for Hovedtankens Skyld idelig maae corrigeres. Man vil end%
% --- p. 47 ---
\opage{47}videre overbevise sig om, at den Forskjel, der med Hensyn paa Tankegangen er en Følge af Sprogenes Forskjellighed, er af en ganske anden Natur, end den, der er en Følge af det Eiendommelige hos de Forfattere, der tænke i det samme Sprog. Man vil endelig overbevise sig om, at forskjellige grammatikalske Vendinger kunne give samme Mening, uden at de derfor i logisk Forstand give nøiagtig den samme Tanke: det tydske \emph{ist gewesen} og det danske \emph{har været} give samme Mening, men ikke samme Tanke. „I det Hele“, siger Heiberg,\footnote{\emph{Perseus.} Journal for den speculative Idee. Nr.~2. 1838. S.~44--45.}% note *)
{} „vil man indsee, at det Sprog, hvori man philosopherer, ikke kan være uden Indflydelse paa selve Philosopheringen. Thi Sproget bestaaer af de, saa at sige, legemliggjorte Kategorier, saa at Grammatiken maa betragtes som en Samling af oplysende Figurer til Logiken.
Derfor har Hegel ogsaa Kategorier, som han ved Tænkningen i et fremmed Sprog ikke vilde være kommen til, og hvoraf nogle heller ikke paa Dansk have noget Tilsvarende“.

Der er i Sprogenes Dannelse og Udvikling en Forskjel, der ingenlunde kan udledes af bestemte Culturtrin, saa lidt som af det Tilfældige i et første heldigere eller uheldigere Greb paa Tingen, men viser hen til det Eiendommelige i Folkenaturen. Man tænke sig det hebraiske Sprog cultiveret, saameget man vil, det faaer dog aldrig det logiske Præg som det græske. Den hebraiske Tænkning er væsentlig psychologisk. Hebræerens Tænken er i særlig Forstand en Talen med sig selv; han tænker med \emph{Hjertet}.\footnote{\textit{Prenons pour exemple l'hébreu, qui nous représente un état fort ancien du langage. On sent que le phénomène qui a servi d'occasion à la création des radicaux de cette langue, et en général des langues sémitiques, a été presque toujours physique. „Je conviens, dit Herder, que le penseur abstrait ne doit pas % note breaks here: continues at the foot of p. 48
trouver la langue hébraique très parfaite; mais sa forme agissante en fait l'instrument le plus favorable au poëte. Tout en elle nous crie: Je vis, je me meus, j'agis! je n'ai pas été créée par le penseur abstrait, par le philosophe profond, mais par les sens, par les passions!} Renan. Anførte Skr. S.~124.}% note **) -- PRINTER: anchor printed "**)", but the note itself is marked "*)"
% PRINTER: the „ opened before "Je conviens" is never closed
{} Ikke blot Jødefolkets %
% --- p. 48 ---
\opage{48}Culturhistorie, ogsaa Sprogets oprindelige Characteer bærer Præg af, at dette Folks høieste Tænkning var og skulde være, ikke en philosophisk, men en religiøs Tænkning. Anderledes med det græske Sprog; thi Grækeren veed, at han tænker med \emph{Hovedet}.\footnote{Grækeren veed da ogsaa, hvad det vil sige, at Sproget er til „for Tankens egen Skyld, og at Mennesket ikke tænker uden i og med Sproget“: Οὐκοῦν διάνοια μὲν καὶ λόγος ταὐτόν· πλὴν ὁ μὲν ἐντὸς τῆς ψυχῆς πρὸς αὑτὴν διάλογος ἄνευ φωνῆς γιγνόμενος τοῦτ᾽ αὐτὸ ἡμῖν ἐπωνομάσθη, διάνοια; . . Τὸ δέ γ᾽ ἀπ᾽ ἐκείνης ῥεῦμα διὰ τοῦ στόματος ἰὸν μετὰ φθόγγου κέκληται λόγος. Plato. Sophist. 263~E. Jvfr. Theæt. 189.~E.}% note *)
{} Længe, længe før her kunde være Tale om en philosopherende Tænkning, der skulde bearbeide Sproget, var Sproget selv eiendommelig forberedt paa philosophisk Udvikling.

Et Sprogs logiske Anlæg kjendes da ikke uden videre derpaa, at det i al Umiddelbarhed afgiver Midler for en Tankeudvikling, thi det gjør jo til en vis Grad ethvert Sprog; men udtrykkelig derpaa, at det under sin Udvikling betinger og fremmer en Tænkning over Tænkningen, en Dømmen om Dommen, at det, med andre Ord, egner sig til at udtrykke rene Tankeformer ved bestemte Afskygninger i bestemte Overgange. Her er da en væsentlig Forskjel imellem den psychologisk naive og den sig bevidste logiske Tænkning. I begge er der en Skjelnen og Skjønnen, i begge en Adskillen og Forbinden af Forestillinger; men den psychologiske Tænkning, der skjønner og skjelner umiddelbart, blander uvilkaarlig det Tilfældige med det Nødvendige, det Specielle med det Almene, Billedet med Begrebet; den logiske Tænkning, der er opmærksom paa sig selv, er derimod sikker paa Abstractionen og udøver en stadig Selvkritik.

% --- p. 49 ---
\opage{49}Den reneste Forstaaelse af sin egen formelle Virksomhed, som Tanken kan opnaae, erholder den ved en Betragtning af de Sprogformer, hvori den bliver sig selv gjennemsigtig. Naar den formelle Logik gjør Forskjel imellem „dømmende“ og „ikke dømmende“ Sætninger, kan denne Adskillelse paa en Maade anerkjendes, forsaavidt Eftertrykket udelukkende falder paa Formen. Men nu viser Uklarheden sig just deri, at Logiken, inden man ret veed deraf, blander det logisk Almene med det logisk Individuelle, det abstract Objective med det abstract Subjective, Dommens Form i Sproget med Domsacten i den dømmende Subjectivitet: dette giver en stor Confusion.

Sætninger som: „det regner“, „han nyser“, „der er Nogen, der ringer“, ere, hedder det, „ikke egentlig dømmende;“ „de ere blot fortællende, berettende, angivende, uden at der er Tale om en Antagen af Et om et Andet.“ Seer man paa det reent Formelle, da maa dette indrømmes, thi Formen er umiddelbar, og giver ingen klar Afspeiling af den logiske Function; men naar der saa tilføies, at slige Sætninger „ville kunne forvandle sig til egentlig dømmende, saasnart der opstaaer Spørgsmaal om Udsagnets eller Angivelsens Rigtighed“, saa er Forvirringen i fuld Gang. Forsaavidt Sætningen: der er Nogen, der ringer, giver Anledning til Tvivl, Undersøgelse og Discussion, sees her alene paa Indholdet, ikke paa Formen. Man forlanger af Sætningen en Garantie for, at den Talende nu ogsaa virkelig tænker ved, hvad han siger; men en saadan Garantie kan Sætningen som logisk Dom, saalænge den dømmende Subjectivitet sees i psychologisk Vilkaarlighed, aldeles ikke give. En Sætning som: Gud er Kjærlighed, kan af en Papegøie udsiges ligesaa flygtigt og tankeløst som Sætningen: der er Nogen, der ringer. Vil man nu, hvad den sidste Sætning angaaer, fordre en nærmere Bestemmelse: \emph{det er saa, det er virkelig saa}, at der er Nogen, der ringer, da maa man, hvad den første Sætning angaaer, ligeledes tilføie Forsikkringen: \emph{det er saa, det er virkelig saa}, at Gud er %
% --- p. 50 ---
\opage{50}Kjærlighed. Men slige Forsikkringer kunne jo igjen anbringes paa en tankeløs Maade; man maa altsaa forvisse sig om Tanken i de givne Forsikkringer ved nye Forsikkringer: det er i Sandhed saa, at det virkelig er saa, at det er saa, o.~s.~v.

Den Forvirring, som den formelle Logik afstedkommer derved, at den, ikke nøiet med Dommens objective logiske Form, streifer ind paa det abstract Subjective og i Udsagnet selv vil have en aparte Bekræftelse paa, at den Talende nu ogsaa virkelig tænker Noget ved, hvad han siger, den samme Forvirring har da ogsaa, om end paa omvendt Maade og af modsat Grund, trængt sig ind i den speculative Logik. Dommene hos Hegel ere nemlig, som vi have seet, saa objective, at de af bare Objectivitet ere nærved at dømme sig selv og gjøre sig selv subjective; det kan da ikke være Andet, end at de i „Subjectiviteten“ objective Domme maae udvikle sig og skride frem med det i sin Subjectivitet objective Begreb, hvis Indhold er Eet med Formen. Da nu Hegel, hvad Dommene og deres Inddeling angaaer, har beholdt den formelle Logiks Former og sat sig til Opgave at mætte de tomme Former med immanent Indhold: saa møde vi ogsaa hos ham den Misforstaaelse, at Formen, istedenfor at have sit logiske Indhold i Reflexion af rene Formbestemmelser, maa fylde sig med et fra en fyldigere, allerede tilbagelagt Kategorie-Udvikling erobret Indhold;\footnote{--- „Wenn Beispiele von Prädikaten der Reflexions-Urtheile gegeben werden sollen, so müssen sie von anderer Art seyn, als für Urtheile des Daseyns. Im Reflexions-Urtheil ist eigentlich erst ein \emph{bestimmter Inhalt}, d.~h. ein Inhalt überhaupt vorhanden; denn er ist die in die Identität reflektirte Formbestimmung, als von der Form, insofern sie unterschiedene Bestimmtheit ist, --- wie sie es noch als Urtheil ist, unterschieden. Im Urtheil des Daseyns ist der Inhalt nur ein unmittelbarer, oder abstracter, unbestimmter. --- Als Beispiele von Reflexions-Urtheilen können daher dienen: Der Mensch ist \emph{sterblich}, die Dinge sind \emph{vergänglich}, % note breaks here: continues at the foot of p. 51
dieß Ding ist \emph{nützlich}, \emph{schädlich}; \emph{Härte}, \emph{Elasticität} der Körper, die \emph{Glückseligkeit} u.~s.~f. sind solche eigenthümliche Prädikate.“ Hegel. Wissenschaft der Logik. Zweiter Theil. S.~92. --- Fremdeles S.~101--102, hvor Talen er om „Das Urtheil der Nothwendigkeit“, hedder det: \emph{Das kategorische Urtheil} hat nun eine solche Allgemeinheit zum Prädikate, an dem das Subjekt seine \emph{immanente} Natur hat. Es ist aber selbst das erste oder \emph{unmittelbare} Urtheil der Nothwendigkeit; daher die Bestimmtheit des Subjekts, wodurch es gegen die Gattung oder Art ein Besonderes oder Einzelnes ist, insofern der Unmittelbarkeit äußerlicher Existenz angehört. \dots\ Was daran \emph{nothwendig} ist, ist die \emph{substantielle Identität} des Subjekts und Prädikats, gegen welche das Eigene, wodurch sich jenes von diesem unterscheidet, nur als ein unwesentliches Gesetzseyn, --- oder auch nur ein Namen ist; das Subjekt ist in seinem Prädikate in sein An- und Fürsichseyn reflektirt. --- Ein solches Prädikat sollte mit den Prädikaten der bisherigen Urtheile nicht zusammengestellt werden, wenn z.~B. die Urtheile:
\begin{center}\begin{tabular}{r@{\ }l}
 & die Rose ist roth,\\
 & die Rose ist eine Pflanze,\\
oder: & dieser Ring ist gelb,\\
 & er ist Gold,
\end{tabular}\end{center}
in Eine Klasse zusammengeworfen, und eine so äußerliche Eigenschaft, wie die Farbe einer Blume als ein gleiches Prädikat mit ihrer vegetabilischen Natur genommen wird, so wird ein Unterschied übersehen, der dem gemeinsten Auffassen auffallen muß.“ --- Men at Dommens logiske Form, naar det gjælder den rene immanente Formudvikling, ikke har Noget med Erkjendelsens iøvrigt forskjellige Indhold at gjøre, er let at indsee. En Dom som: Mennesket er dødeligt, er i formel Henseende ligesaavel en assertorisk, kategorisk og positiv Dom som den er en „Reflexions-Dom“, ja hvilken Dom er ikke ifølge sin Form en Reflexions-Dom? Den positive Dom er jo kun positiv i Reflexion af den negative, den singulære i Reflexion af den particulære, den kategoriske i Reflexion af den hypothetiske o.~s.~v.}% note *)
% PRINTER: the Hegel quotation from S. 101--102 ("Das kategorische Urtheil ... auffallen muß.“) has a closing quote but no opening one
{} som om det ikke fulgte af sig selv, at alle Former %
% --- p. 51 ---
\opage{51}af logiske Domme strax maatte kunne komme til Anvendelse under Udvikling af Systemets allerførste og fattigste Kategorier.

Den Forvirring, der saaledes har indsneget sig i Logikens %
% --- p. 52 ---
\opage{52}Bestemmelse af de særlige Dommes Eiendommelighed med Hensyn paa Formen, medfører saa igjen en anden og tilsvarende Forvirring, naar Spørgsmaalet er om Dommenes Inddelinger.

I sin „Udsigt over Dommenes og Dømmelsernes Inddelingsmaader“ gaaer den formelle Logik omtrent tilværks, som om det gjaldt en Opstilling af Meubler, en Fordeling af Pindebrænde, en Classification af Sko og Støvler, o.~dsl. Den formelle Logik springer fra en Inddelingsgrund til en anden, opregner snart flere, snart færre, uden at den ene paa nogen Maade udvikles af den anden.

Saaledes kunne da Dommene inddeles med Hensyn til deres Materie i mathematiske, physiske, astronomiske, astrologiske, nekromantiske, chiromantiske o.~s.~v. o.~s.~v. Domme.

Dernæst kunne Dommene inddeles med Hensyn paa den reale Nexus imellem Subject og Prædicat i „ikke egentlig dømmende“ og „egentlig dømmende Sætninger“. De „ikke egentlig dømmende Sætninger“ kunne igjen inddeles i de „blot opstillende“ og de „tautologiske“. De „blot opstillende Sætninger“ kunne igjen inddeles i de „blot fortællende,“ de „blot beskrivende“ og de „blot begrebsopstillende.“ Men de „blot opstillende Sætninger“ kunne igjen inddeles i dem, hvori der ikke ligger noget „Indbegreb af Dømmelser,“ og dem, hvori der ligger et „Indbegreb af Dømmelser uden at træde frem som Dømmen“. Hvad de „tautologiske Sætninger“ angaaer, kunne disse efter Omstændighederne inddeles i de „aldeles intetsigende“ og de „ikke aldeles intetsigende“. Disse sidste kunne igjen inddeles i dem, der tjene til „at hæve Sagen“, og i dem, der ere „blot begrebsopstillende“ o.~s.~v. Fremdeles kunne de „tautologiske Sætninger“ inddeles i de „direct og de indirect tautologiske“, og disse igjen inddeles o.~s.~v. Kunne nu de „ikke egentlig dømmende Sætninger“ taale saamange Inddelinger, Underinddelinger og Underinddelingers Underinddelinger, hvad maae de „egentlig dømmende“ da ikke kunne taale?

% --- p. 53 ---
\opage{53}Endvidere kunne Dommene inddeles med Hensyn paa den logiske Nexus imellem Subject og Prædicat, og under dette Hensyn igjen med Hensyn paa Qvalitet, Qvantitet o.~s.~v. Endelig kunne disse Inddelinger „overhovedet inddeles“ i dem, der i Grunden ikke vedkomme Logiken, nemlig de materielle, og dem, der vedkomme Logiken. De Inddelinger, der vedkomme Logiken, kunne igjen inddeles i „de logiske“ og saa i „visse Inddelinger, som ere en Følge af, at to eller flere Domme kunne staae i forskjelligt indbyrdes Forhold,“ Inddelinger, som til Syvende og Sidst da ogsaa høre til de logiske. o.~s.~v. o.~s.~v. --- „Saaledes“, siger Heiberg, „forbauses man ved at see 20 til 30 Arter af Omdømmer opstillede, hver udmærket med sit barokke Navn. Da endvidere Inddelingen af disse Arter skeer efter forskjellige Inddelingsprinciper, saa kan enhver Art af en vis Inddeling forbinde sig successiv med alle de andre Arter af de andre Inddelinger, hvorved der fremkommer en Hærskare af Distinctioner, som ere aldeles intetsigende, eftersom der i Omdømmet ikke ligger noget Princip for denne Mangfoldighed.“\footnote{Den speculative Logik §~144 Anm.~3.}% note *)

At den formelle Logik med sit Ruskumsnusk af Afdelinger, Rubrikker og Inddelinger har virket frastødende paa et saa klart og logisk Hoved som Heiberg, er forstaaeligt; mere Opmærksomhed vækker det, naar vi see ham vende sig endog imod Hegels eget Forsøg med en immanent Deduction. „Hegel, som gav den scholastiske Logik sit Banesaar, har“ --- siger Heiberg „ved at ville deducere de forskjellige Arter af Omdømmer som forskjellige Udviklingstrin, gjennem hvilke Omdømmet bliver til Slutning, ikke desmindre skjænket dem en Opmærksomhed og tildeelt dem en Vigtighed, som de ere langt fra at fortjene.“

Men her hugger Heiberg nu Knuden over, istedenfor at løse den. For at blive Formalismen qvit, fordrer han, at %
% --- p. 54 ---
\opage{54}Logiken aldeles ikke skal indlade sig paa nogen nærmere Bestemmelse af Dommenes indbyrdes Forskjellighed, men gaae fra Fastsættelsen af Dommens Begreb lige over til Slutningen, saa at den videre Udvikling af Dommen falder sammen med en Udvikling af Slutningen. Herved overseer han, at den logiske Udvikling af de forskjellige Former i Dommens Form er en væsentlig Formudvikling; skulde man i Logiken overspringe væsentlige Formudviklinger for at undgaae Formalisme, da maatte det vel ende med, at man tilsidst oversprang den hele Logik.

„Ethvert Omdømme“ --- hedder det\footnote{Den speculative Logik §~144. Anm.~3.}% note *)
{} --- „er positivt, thi enten det Almindelige, hvorunder det Enkelte subsumeres, er bestemt positivt eller negativt, saa er selve Subsumtionen altid en positiv Act.“ Dette maa indrømmes; men i Indrømmelsen sees da ogsaa strax Fordringen paa en væsentlig Formudvikling. Reflecteres her nemlig paa Domsacten, altsaa paa det i Functionen Subjective, da viser det sig, at Subsumtionens positive Act væsentlig er tre Positioner: Positionen af Subjectet, Positionen af Prædicatet, Positionen af Subjectets og Prædicatets Eenhed. Vi kunne nu gjerne begynde med den Forudsætning, at her væsentlig kun er een eneste logisk Dom; men da Tankeformen altid ændres med Sprogformen, og ingen enkelt Form er absolut adæqvat, saa indsees, at den ene Dom nødvendig maa udvikle sig til en Fleerhed af Domme, idet den ene Form udvikler sig til en Fleerhed af Former. Reflectere vi paa Positionen af Subjectets og Prædicatets Eenhed, altsaa paa \emph{Copula}, da kan man vel indrømme, at Forskjellen mellem Subject og Prædicat er „en Reflexionsforskjel, hvori det Ene kun er, hvad det Andet er“; men Reflexionsforskjellen er derfor ligefuldt en væsentlig Forskjel: det Ene er alligevel ikke, hvad det Andet er. Vistnok siger Dommen, at Subjectet er, hvad Prædicatet er, at Subjectet %
% --- p. 55 ---
\opage{55}er sit Prædicat; men Dommen siger jo ogsaa, at Subjectet ikke er Prædicat, Prædicatet ikke Subject. Sættes nu Subjectets og Prædicatets Væren umiddelbart som een og samme Væren, er Dommen positiv, og den positive Dom bestemt ved Formen: S er P. Sættes omvendt Subjectets og Prædicatets Væren som en med sig ueens, i sig selv forskjellig Væren, saa er Subjectet ikke, hvad Prædicatet er; Dommen er altsaa negativ, og den negative Dom bestemt ved Formen: S er ikke P. Umiddelbart og efter Erkjendelsens Indhold kan den positive Dom jo vistnok stilles selvstændig imod den negative: her er forsaavidt to forskjellige Domme. Det er da ogsaa paa denne Forskjellighed, at Logiken reflecterer, naar den inddeler Dommene efter deres \emph{Qvalitet}: i \emph{positive}, \emph{negative} og \emph{uendelige}. Men, væsentlig seet, er her kun een Dom. Den positive Dom faaer kun logisk Betydning i Reflexion af den negative, ligesom omvendt den negative Dom maa blive uendelig og i sin Uendelighed meningsløs, dersom den ikke kan reflectere sig i sin tilsvarende positive. Hvorledes de sprogformende Functioner forberede de tankeformende, sees da i denne Sammenhæng deraf, at der gives et sprogligt Udtryk for et Subjects Væren eller Ikke-Væren af sit Prædicat.\footnote{At Copula „er“ ikke findes i det hebraiske og aramaiske Sprog, er i Religionsstridighedernes Historie særlig fremhævet med Hensyn paa de sakramentale Ord: \cjRL{hU' dAmiy}, \cjRL{hU' gUpiy}.
% *** HEBREW: SET AS HEBREW (was transliterated until 2026-09-22). ***
% The page sets the pointed phrase הוּא דָמִי, הוּא גוּפִי ("this is my blood, this is
% my body"). It is now reproduced with cjhebrew, which typesets pointed Hebrew
% under plain pdfLaTeX and needs no font change to the Libertinus body text.
% cjhebrew input is a transliteration of its own: ' = alef, U = vav with dagesh
% (shuruq), A = qamats, i = hiriq, y = yod; final forms are automatic. So
% hU' dAmiy = הוּא דָמִי and hU' gUpiy = הוּא גוּפִי. Each \cjRL group is internally
% right-to-left while the surrounding Danish stays left-to-right, which puts the
% two phrases, the comma and the full stop exactly where the 1866 page has them.
% Other books in this repo still transliterate this sort of phrase
% (speculative-methode pp. 33, 39, 159; treschow/menneskelige-natur p. 115);
% they can be converted the same way.
 Et Sprog kan have den abstracte Copula, uden at det just derfor har et til Kategorien \emph{Væren} svarende Abstractum (saaledes er f.~Ex. det franske \textit{être} ikke Væren, men Substans), men et Sprog, der ikke har den abstracte Copula, vil vanskelig frembyde adæqvate Udtryk for Kategorier som Væren og Ikke-Væren.}% note *)
% BUILD: set with cjhebrew (TeX Live package, pdflatex-safe). The Unicode Hebrew
% in the comment above is for reading the source only and is never typeset.
% UNSURE: Hebrew pointing. BL copy (1-bit, 600 ppi, broken type) confirms the letters (הוּא twice with shuruq; ד-מ-י; ג-ו-פ-י) and word order, but the qamets under dalet and the hiriqs under mem and pe cannot be made out with certainty
% NOTE: "S", "P" here and below are set in antiqua capitals in the print; left roman, not \textit.

Reflectere vi særlig paa \emph{Positionen af Subjectet}, da kan man vel indrømme, at det Positive i Domsacten væsentlig bestemmes ved det Negative, og omvendt; men Eenheden af den dobbeltsidige Bestemthed faaer dog først sin punk%
% --- p. 56 ---
\opage{56}tuelle Betydning, naar den udtrykkelig sættes i Dommens Form. Dette S er P; denne Dom er, umiddelbart seet, positiv; men det Negative forvandler sig til Grændse og skinner igjennem, naar der lægges Vægt paa, at det indtil videre kun er \emph{dette, ikke flere} S, hvorom P udsiges. Ligesom nu den positive Dom: dette S (idetmindste) er P, hviler paa Muligheden af, at der er S, som ikke ere P, saaledes hviler omvendt den negative Dom, at der er S, som ikke ere P, paa Muligheden af, at nogle, ja flere, og stedse flere S kunne være P. Forsaavidt det endnu er uvist, hvormange S, der ere P, er det jo ligesaa muligt, at de alle ere det, som at intet er det. Det er da heraf indlysende, at den formelle Logiks Inddeling af Dommene efter \emph{Qvantiteten}: i \emph{singulære}: dette S er P, \emph{particulære}: nogle S ere P, og \emph{generelle}: alle S ere P, har væsentlig Gyldighed med Hensyn til Formen, men bliver vildledende, naar der lægges Vægt paa de enkelte Dommes umiddelbare Selvstændighed. Det er den samme, væsentlig den selvsamme logiske Dom, som i Qvantiteten \emph{fordeler} det Negative og Positive, og som i Qvaliteten lader det Positive udelukke det Negative, og omvendt.

Reflectere vi særlig paa \emph{Positionen af Prædicatet}, da kan man gjerne indrømme, at Eenheden af den dobbeltsidige Bestemthed, af Negativt ved Positivt, er afgjørende for Dommen i enhver Form; men det kommer da tillige væsentlig an paa, hvorvidt Formen selv angiver den Tænkningsact, hvorved Eenheden er sat. At Eenheden af Subject og Prædicat er en Forskjelseenhed, vil da med andre Ord sige, at Prædicatet som Prædicat \emph{tilkommer} Subjectet, forsaavidt det baade er i Subjectet og \emph{kommer til} Subjectet. At Prædicatet tilkommer Subjectet, vil da sige, at Prædicatsbestemmelsen er Subjectets egen oprindelige Bestemmelse; her sættes altsaa Intet i Prædicatet, som ikke allerede er i Subjectet; Prædicatsbestemmelsen er en i Formen gjennemsigtig Subjectsbestemmelse: \emph{Dommen er analytisk}. Men at Prædicatet væsentlig tilkommer Sub%
% --- p. 57 ---
\opage{57}jectet forudsætter da ogsaa, at det ved Prædiceringen virkelig bringes sammen med Subjectet; der føres altsaa ved Prædiceringsacten en ny Bestemmelse til Subjectet selv: \emph{Dommen er}, med andre Ord, \emph{synthetisk}. Fra Endelighedens Synspunkt tager det sig jo vistnok ud, som om den analytiske Dom var fuldstændig for sig ligeoverfor den synthetiske, og omvendt. Men dette er i logisk Forstand en reen Illusion, og den formelle Logik maa bære Skylden for, at denne for Videnskaben saa fordærvelige Illusion bestandig faaer ny Næring. Hvorledes skulde det, naar vi udtrykkelig see paa Begrebet, være muligt at finde en Dom saa analytisk, at deri ikke fandtes et Moment af Synthese. Man lade Prædicatet nok saa meget være i Subjectet selv; Subjectet kan umulig være Subject uden i formel Modsætning til sit Prædicat: ved Dommens Form er Prædicatets formelle Selvstændighed uafviselig hævdet. Omvendt kan Subjectet efter sin umiddelbare Realitet have hvilken Selvstændighed, det vil, --- i Modsætning til sit Prædicat har det, i Kraft af den ved Prædiceringsacten bestemte Form, kun formel Selvstændighed. Er det nu paa den ene Side ikke muligt at tænke sig en Dom saa analytisk, at den ikke skulde indeholde noget Synthetisk, saa er det paa den anden Side heller ikke muligt at tænke sig en Dom saa synthetisk, at den ikke tillige var analytisk. En Dom saa synthetisk, at det Analytiske ganske var udelukket, maatte da være en Dom, hvori Subject og Prædicat optraadte med saa umiddelbar Selvstændighed hvert for sig, at deres Forening var at ansee for aldeles tilfældig; men en saadan Dom er i sit Væsen meningsløs. Subjectet i Dommen er kun Subject ved at reflectere sit Prædicat, ligesom Prædicatet kun er Prædicat ved at reflectere sit Subject. Hvad Prædicatet udsiger, maa Subjectet, var det end kun i et eneste Øieblik, være; men hvad Subjectet som Subject selv er, kan jo dog umulig være dets Væren uvedkommende: den synthetiske Dom er væsentlig %
% --- p. 58 ---
\opage{58}analytisk, efterdi den analytiske Dom væsentlig er synthethisk.\footnote{Det var ved sandt Genieblik, at Kant blev opmærksom paa, at ogsaa aprioriske Domme kunne være synthetiske; men han blev staaende paa Halvveien. At synthetiske Domme ere \textit{a priori} mulige, beroer just paa, at det i Dommen Synthetiske er \textit{a priori} Eet med det Analytiske.}% note *)
% PRINTER: printed "synthethisk" (divided "synthe-|thisk"), sense "synthetisk"

Den positive Dom er altsaa i Begrebet analytisk-synthetisk ligesaa vel som den negative, den singulære analytisk-synthetisk ligesaavel som den generelle. Det er i logisk Forstand saa langt fra, at det Analytiske og Synthetiske her skulde være Noget for sig ligeoverfor det Positive, det Negative, det Singulære, det Particulære og det Generelle, at det tvertimod begrunder Forskjelseenheden af det Positive og det Negative, det Singulære og det Particulære o.~s.~v. Heri ligger da ogsaa Grunden til, at Forskjellen imellem den analytiske og synthetiske Dom ikke, saaledes som Forskjellen imellem den positive og den negative, den particulære og den generelle Dom, har noget direct Udtryk i Sproget.

Da det nu ikke destomindre giver en væsentlig Forskjel, om det er Dommens synthetiske eller analytiske Side, der i Tankens Øieblik foresvæver den Dømmende, saa vækkes allerede herved en Forventning om, at Sproget vist maa have, om ikke et direct, saa dog et indirect Udtryk for Adskillelsen mellem det blot Synthetiske, det Analytiske og det Analytisk-Synthetiske.

Hvorledes dette Udtryk maa være beskaffent, kan, ihvorvel den formelle Logik ikke lader os ane det Mindste, dog neppe være vanskelig at udfinde. Det er da for det Første aabenbart, at det umiddelbart Synthetiske ligefrem bestemmes ved det umiddelbart Givne, S er virkelig P (Muren er virkelig rød; hans Øine ere virkelig blaae): \emph{Dommen er assertorisk}. Men hvorledes skal nu Analysen og Bevidstheden om det abstract %
% --- p. 59 ---
\opage{59}Analytiske udtrykkes? Naturligviis paa negativ og indirect Maade. Naar Prædicatet, umiddelbart seet, kun tilkommer Subjectet, fordi det i Virkeligheden findes forenet dermed, saa ligger jo heri, at det Modsatte kunde være Tilfældet, at nemlig Foreningen ikke fandt Sted; Subjectet svæver da imellem at være og ikke at være det bestemte Prædicat: S er maaskee P, maaskee ikke: \emph{Dommen er problematisk}. Men, naar Dommen er problematisk, saa paatrænger sig uvilkaarlig det Spørgsmaal: hvilket Prædicat maa det omspurgte Subject da, ifølge sit Begreb, uimodsigeligt have? For at afgjøre, hvad heri er nødvendigt og hvad tilfældigt, maa Analysen foretages. Og hvormed ender saa Analysen? Ikke med simpelthen at udskyde det Tilfældige og beholde det abstract Nødvendige (thi et saadant Forsøg røber kun ringe Dømmekraft), men med at optage det Tilfældige som Moment i det concret Nødvendige: S er nødvendigviis P; her er Eenheden af det Synthetiske og det Analytiske: \emph{Dommen er apodiktisk}.

Den formelle Logik inddeler Dommene med Hensyn paa \emph{den reale Nexus} imellem Subject og Prædicat i \emph{analytiske} og \emph{synthetiske}, og med Hensyn paa \emph{Modaliteten} i \emph{assertoriske}, \emph{problematiske} og \emph{apodiktiske}; men den lader enhver af Inddelingerne staae som en Rubrik for sig, uden Antydning af, at den ene skulde i nogen Maade vedkomme den anden. Det er denne Udstykning, denne halv tankeløse, men dog „tænksomme“ Rubricering, der gjør saa ubodelig Skade. Ved at oversee Sammenhængen mellem de to Inddelingsmaader, overseer man naturligviis Sammenhængen imellem Formalfunctionen og Dommens Form. Den samme Dom, der for den Ene stiller sig som assertorisk, stiller sig da for den Anden som problematisk, for den Tredie som apodiktisk o.~s.~v. Indseer man derimod Sammenhængen mellem Formalfunctionen og Formen, saa indseer man ogsaa Sammenhængen mellem de forskjellige Inddelingsmaader. Den assertoriske Dom er i sit Begreb assertorisk, just fordi den %
% --- p. 60 ---
\opage{60}som umiddelbart positiv er umiddelbart synthetisk. I den analytiske Dom gaaer Prædicatet op i sit Subject, i den problematiske løsner det sig fra sit Subject; men det ene af disse Extremer forudsætter det andet. Skal Analysen begynde, maa det være forbi med den umiddelbare Vished: den problematiske Dom svarer til en dømmende Bevidsthed, som ved det umiddelbart Synthetiske drives mod det Analytiske. Endelig er den apodiktiske Dom da heller ikke apodiktisk, fordi det falder en Tankeløs ind, at bruge Formen: S er nødvendigviis P, hvor han helst skulde bruge Formen: S er maaskee P; Formen: S er nødvendigviis P, forudsætter udtrykkelig, at Dommen hviler paa en udtømmende, til Synthesen svarende Analyse. I den apodiktiske Dom er det Assertoriske i Synthesen blevet problematisk, og dette Problematiske paa Grund af Analysen igjen assertorisk, og derved meer end assertorisk.

Reflectere vi endelig paa \emph{den ponerede Eenhed af alle tre Positioner}, altsaa paa Ligevægten imellem Subjectets og Prædicatets logiske Selvstændighed, saa komme vi til den concreteste Formfuldendelse, som Tanken ved Sprogets Hjælp er istand til at give den logiske Dom.

Den positive Dom er i sin fulde Bestemthed \emph{kategorisk}; den kategorisk positive Dom kan ikke umiddelbart modsættes den negative, thi det Positive har optaget det Negative som Moment i sig. Den kategoriske Dom udtrykker Eenheden af det Singulære og det Generelle, hvilket blandt Andet fremlyser deraf, at den for den generelle Dom charakteristiske Pluralis (Alle Mennesker ere dødelige) forvandles til en væsentlig, for den universelle Dom betegnende Singularis (Mennesket er, som saadant, dødeligt). At den kategoriske Dom ogsaa er bestemt som analytisk-synthetisk, sees af Subjectets og Prædicatets udviklede Begrebsforhold; thi Subjectet sættes her i Eenhed med sit Prædicat, som det Enkelte (Existerende) i Eenhed med det Almene (Begrebet). Den kategoriske Dom kan da som assertorisk heller ikke have den problematiske uden%
% SEAM: word continues on p. 61 as "for" ("uden-|for"); written whole here, so the p. 61 batch should begin after "for".
% --- p. 61 ---
\opage{61}for sig; det Tilværende kan jo ikke falde udaf sin Kategorie:
den kategoriske Dom er væsentlig apodiktisk.

Men just fordi den kategoriske Dom har alle Dommens
Grundbestemmelser i sig, kan den ikke have nogen særlig Sprogform:
den er i sit Udtryk ligefrem Dommen som Dom.
Dens Eiendommelighed udvikler sig først, og bliver kjendelig i
Sproget, naar der spørges om dens egne indvortes Modsætninger.
S er (kategorisk) P; det Enkelte (S) subsumeres under
sit Almene (P); det ene og samme Begreb (P) omfatter en
Mangfoldighed af særskilte S, nemlig S\textsubscript{1}, S\textsubscript{2}, S\textsubscript{3}, \dots.\ Altsaa:
S\textsubscript{1} er P, S\textsubscript{2} er P, S\textsubscript{3} er P. o. s. v. Men kategorisk
gjælder ogsaa det Omvendte, at nemlig mange Almeenbestemmelser
findes forenede i det ene og samme existerende Enkelte,
at, med andre Ord, Prædicatet inhærerer sit Subject. Dette
giver os da følgende Række: S er P\textsubscript{1}, S er P\textsubscript{2}, S er P\textsubscript{3} \dots\
Heraf sees, at Prædicatet i den kategoriske Dom kan være
deels eensidigt som en Egenskab ved Tingen, deels alsidigt som
det Totalitetsbegreb, der existerer i Subjectet, den Art, der er
realiseret i Individet.

Er Grundforholdet mellem Subject og Prædicat, det
Existerende og Kategorien, bragt til Bevidsthed, saa er det med
det Samme bragt til Bevidsthed, at den kategoriske Dom
umulig kan sees i sin Betydning, uden at der danner sig en
Række af kategorisk sammenhængende Domme. Dersom nu
f. Ex. S (kategorisk) er P, hvad saa? Ja, saa maa s (kategorisk)
være p; eller endnu abstractere og mere formelt:
dersom S er Subject, saa maa P være Prædicat; saaledes
have vi da den \emph{hypothetiske} Dom. I den hypothetiske
Dom er den kategoriske fordoblet; det seer ved første
Øiekast ud, som var her to selvstændige Domme, een i Forsætningen:
dersom S er P, og een i Eftersætningen: saa er
s p; men Eenheden bestaaer kategorisk i det mellem de tvende
Sætninger, Betingelsen og det Betingede, givne Afhængigheds%
% --- p. 62 ---
\opage{62}forhold.\footnote{% note *)
Til Adskillelse af den kategoriske og den hypothetiske Dom fremhæver
Kant, at hiin bestemmes ved Kategorierne: Substans og
Inhærens, denne ved Kategorierne: Aarsag og Virkning. Hertil bemærker
Trendelenburg (Logische Untersuchungen. Zweiter Band. Berlin
1840. S. 178.), at disse Begreber ikke danne nogen Modsætning.
„Die Substanz ist in der Eigenschaft causal. Die Eigenschaft ist
die an das Ding gebundene Thätigkeit. Die strengste Form der
Inhärenz ist das Verhältniß der im Ganzen inwohnenden Theile.
Da aber das Ganze die Theile trägt und zu dem macht, was sie
sind, so schlägt auch hier die Inhärenz in die Causalität über.“ Men
hvad beviser nu dette? Ikke Andet end, at man med Hensyn paa
den samme Gjenstand kan anvende snart Kategorierne Substans og
Inhærens, snart Kategorierne Aarsag og Virkning. Det, her skulde
bevises, er, at man, naar Kategorierne Aarsag og Virkning danne
Grundlaget, ligesaa træffende anvender den kategoriske som den
hypothetiske Form; men dette skal man vel lade være at bevise. Til
Forsvar for Kants Opfattelse erklærer Apelt (Die Theorie der Induction.
Leipzig 1854. S. 9.) det for en plat Umulighed, at forvandle
den ene af de to Slags Domme til den anden. „Was man
von der \emph{Verwandelung} kategorischer und hypothetischer Urtheile
sowie kategorischer und hypothetischer Schlüsse in einander gesagt
hat, ist eine logische Fabel. Denn jede solche angebliche Verwandlung
besteht nur in einer Veränderung des \emph{grammatischen
Ausdrucks}, aber nie in einer wirklichen Veränderung \emph{der logischen
Form}.“ Dette er nu ligeledes en stor Misforstaaelse;
thi med Forandringen i det Grammatikalske indtræder just en Ændring
i det Logiske. Trendelenburg har da ogsaa (Anfr. Skr. S.
% PRINTER: printed "Anfr. Skr.", sense "Anf. Skr." (cf. p. 65)
179.) ved Exempler oplyst, hvorledes en kategorisk Dom lader sig
forvandle til en hypothetisk: „Wenn ein hypothetisches Urtheil verschiedene
Gegenstände zusammenbringt, so sind sie immer unter
einem höhern Begriff eins, da sie als Bedingung und Bedingtes,
als Grund und Folge eins gesetzt werden. Wermöge dieser in
% PRINTER: printed "Wermöge", sense "Vermöge"
einer Thätigkeit oder einer Eigenschaft ruhenden Einheit kann es
geschehen, daß auch in diesem Falle das hypothetische Urtheil ein fast
gleichbedeutendes kategorisches neben sich hat. Man vergleiche z. B.
% note breaks here, p. 62/63
die Urtheile: Wenn Bernstein gerieben wird, so entwickelt sich Electricität.
Der geriebene Bernstein entwickelt Electricität. Ferner:
Wenn ein Kegel durch eine Ebene geschnitten wird, so entstehen
regelmäßige Curven. Der Schnitt durch einen Kegel erzeugt regelmäßige
Curven“.

Altsaa: fra den ene Side søger man at udjævne Forskjellen
mellem den kategoriske og den hypothetiske Dom, fordi man indseer,
at forskjellige Kategorier kunne anvendes paa samme Gjenstand, og
mærker ikke, at Dommens Form forandrer sig med de forandrede
Kategorier; fra den anden Side vil man reise en Skillemuur imellem
den kategoriske og den hypothetiske Dom, fordi der nu vitterlig er
en Forskjel imellem det Logiske og det Grammatikalske, og mærker
ikke, at det Logiske, ifølge Vexelbestemmelsen af de sprogformende og
tankeformende Functioner, ændres ved det Grammatikalske. Et
tredie Standpunkt repræsenteres af Ulrici (Compendium der Logik.
Leipzig 1860. S. 177--178.): „Hypothetische Urtheile giebt es in
Wahrheit gar nicht. Denn daß sie auf dem Verhältniß der Ursache und
Wirkung beruhen, ist (z. B. von dem Urtheil: wenn dieß Mineral
ein Metall ist, so ist es schmelzbar) nicht nachweisbar und würde
keinen Unterschied begrüuden. Und ihre anderweitige Verschiedenheit
% PRINTER: printed "begrüuden", sense "begründen"
von den kategorischen ist entweder nur eine rein äußerliche des sprachlichen
Ausdrucks, oder beruht bloß auf einer rein subjectiven Beziehung
des Urtheilenden zu dem Gegenstande (Inhalt) des Urtheils,
--- ist also ohne alle logische Bedeutung.“ Den dømmende Bevidstheds
subjective Forhold til Dommen og Dommens Gjenstand
er altsaa --- „ohne alle logische Bedeutung?“

Til at rede sig ud af disse Qvaklerier gives der kun een Udvei:
at fastholde Dommens Form i dens Eenhed med det abstract
logiske Indhold og mærke sig Forskjelseenheden af Ord og Tanke.}
At den hypothetiske Dom kan betegne noget Uvist,
noget endnu Uafgjort, i Lighed med den problematiske, maa
% --- p. 63 ---
\opage{63}indrømmes, men Eftertrykket falder dog ingenlunde paa Uvisheden,
tvertimod! det falder kategorisk paa den afgjørende Bestemthed,
hvormed Forsætningen medfører Eftersætningen: dersom
S er P, saa er (kategorisk) s p, eller dersom s ikke er
p, saa er (kategorisk) S ikke P.

Idet hver existerende Ting udtrykker sin Kategorie, existere
Tingene i kategorisk Sammenhæng. Den kategoriske Sammen%
% --- p. 64 ---
\opage{64}hæng er tilfældig-nødvendig, tilfældig med Hensyn paa de tilværende
Leds Umiddelbarhed, nødvendig med Hensyn paa
deres indbyrdes Afhængighed. Tingenes tilfældig-nødvendige
Sammenhæng bringes udtrykkelig til Bevidsthed i den hypothetiske
Dom; men den hypothetiske Dom er i denne Henseende
kun en eensidig Udvikling af den kategoriske. Den hypothetiske
Dom har sit Støttepunkt i den Forudsætning, at det Enkelte
subsumeres uden det Almene,\footnote{% note *)
See vi bort fra Deductionen og indblande Forestillinger om et
eller andet reelt Indhold, kunne Totalitetsprædicater naturligviis
ogsaa optræde som Subjecter: dersom Planten er en Rose,
saa er Frugten o. s. v.}
% NOTE: "subsumeres uden det Almene" as printed (cf. "under sit Almene", p. 61)
at det altsaa er de enkelte
Ting i deres Relationer og gjensidige Betingethed, hvorpaa
det under Totalitetsudviklingen væsentlig kommer an: Totalitetseenheden
er endnu kun Sammenhængseenhed. Men hermed
er den kategoriske Dom ingenlunde udtømt.

Ifølge den kategoriske Dom kan Tingen først fuldstændig
svare til sin Kategorie og Kategorien til Tingen, naar begge
hver i sin Form udtrykke den samme Totalitet af kategoriske
Bestemmelser. Men naar Prædicatet har samme Totalitet
som Subjectet, kun med den Forskjel, at Prædicatsbestemmelserne
ere væsentlige og logiske, Subjectsbestemmelserne existentielle og
antilogiske, saa kan Forholdet jo vendes om: Subjectet, det
Existerende, bliver da Prædicat, og Prædicatet, det Existerendes
Begreb, Subject. Begrebet, det i sig totaliserede Væsen,
virkeliggjøres i en Mangfoldighed af Existentser; disse Existenser
udgjøre nu Prædicatets Led, medens Begrebseenheden
afgiver Subjectet. Saaledes fremkommer da, naar vi see paa
Subjectets Væsen, den conjunctive Dom: S er (i sit Væsen)
baade P\textsubscript{1} og P\textsubscript{2} og P\textsubscript{3}, og, naar vi see paa Prædicatets existerende
Led, den \emph{kategorisk-disjunctive} Dom\footnote{% note **)
At den disjunctive Dom iøvrigt ogsaa kan have det Existerende
som Subject og Prædicat i eensidige Kategorier, (Tingen er enten
% note breaks here, p. 64/65
sort eller hvid eller rød) er i denne Sammenhæng, hvor det gjælder
om Deductionen, uden Betydning.}: S er
% --- p. 65 ---
\opage{65}(i sin Existens) enten P\textsubscript{1} eller P\textsubscript{2} eller P\textsubscript{3} \dots.\ Da nu hvert
P i sin Totalitet er et virkeliggjort S, saa er ikke alene hvert
existerende Led en relativ Totalitet for sig, men tillige Moment
i den omfattende, alle Leddene gjennemtrængende Totalitet.
Den gjennemtrængende Totalitetseenhed er en Grundeenhed;
ved Vexelbestemmelsen af den hypothetiske, den conjunctive og
den disjunctive Dom bringes Formen til Bevidsthed, den
kategoriske Form, der tjener til Udtryk for, at de hinanden udelukkende,
men dog betingende Leds Sammenhæng væsentlig beroer
paa en Grundeenhed.\footnote{% note *)
Alle disjunctiven Urtheile \dots.\ (die auf dem Verhältniß der Wechselwirkung
oder des Ganzen und Theils beruhen sollen, was wiederum
nicht einmal richtig ist) sind ebenfalls \emph{nur} sprachlich von den
s. g. positiven (kategorischen) unterschieden, und lassen sich daher
unbeschadet ihres Sinnes in kategorische umwandeln: Das Urtheil
z. B.: Cajus ist entweder gesund oder krank, ist nur sprachliche Abkürzung
für: Es ist unmöglich, daß Cajus zugleich gesund und
krank sey, oder: Es ist nothwendig, daß er entweder krank oder gesund
ist. Und das Urtheil: Die Kegelschnitte sind entweder Kreise
oder Ellipsen oder Parabeln oder Hyperbeln, läßt sich auch so ausdrücken:
die Kegelschnitte sind theils Kreise, theils Ellipsen oder:
Die Kegelschnitte können nur die Form von Kreisen, Ellipsen etc.
haben, nur diese Formen sind möglich. Ulrici. Anf. Skr. S. 178.
Det Sande i denne Reduction er, at Dommens særskilte Former
ere Modificationer af een og samme Grundform; men det Usande
er, at de særskilte Former af den Grund ikke skulde have nogen
logisk Betydning. Hvor vilkaarlig Reductionen er, falder i Øinene,
naar Dommen: Cajus er enten sund eller syg, skal være „sprachliche
Abkürzung“ for: det er umuligt, at Cajus paa een Gang kan være
baade sund og syg. Det vilde virkelig være interessant at see en
Smule Beviis for, at Dommen i den sidste Form er mere logisk
end i den første. Hvorfor kan Cajus ikke paa een Gang være
baade sund og syg? Fordi her forudsættes et afgjørende Enten-Eller:
Enten sund eller syg; hvis der reises Tvivl om denne Forudsætning,
% note breaks here, p. 65/66 (Forudsætning divided For-|udsætning)
reiser der sig med det Samme Tvivl om Dommen:
det er umuligt, at Cajus paa een Gang kan være baade sund og
syg. Det Logiske i den sidste Dom kan altsaa ligesaa vel siges at
beroe paa Logiken i den første, som omvendt.}

% --- p. 66 ---
\opage{66}Naar den formelle Logik indskærper Nødvendigheden af
at gjennemgaae det Formelle ved Tænkningen paa en udførlig
og indtrængende Maade, maae vi give den Ret. Læren om
Dommens Inddelingsmaader burde i saa Henseende ligesaa
lidt afskrække os som Læren om Dommenes „Subalternation“,
„Æqvipollens“, „Opposition“, „Conversion“ „Contraposition“
o. s. v., fordi den ikke er „aandrig og tiltrækkende“,
naar den blot var videnskabelig frugtbar. Men her er just
den store Misforstaaelse, hvori Forsvarerne af den formelle
Logik gjøre sig skyldige. Det er ikke Udførligheden og Skarpsindigheden,
men det Golde, Ufrugtbare i Behandlingen af de
logiske Former, hvorimod her opponeres. Istedenfor at paavise
et Princip, der selv sætter Tanken i Bevægelse, og under
Tankebevægelsen lader den ene Formbestemmelse i en given Sammenhæng
udvikle sig af den anden, gaaer Skolelogiken tvertimod
ud paa at afsætte og opstille værende Tankeresultater
med Tilsidesættelse af ethvert Hensyn til Princip og Bevægelse.
De frosne Bølger i en Bæk bære dog altid Spor af Bevægelse,
nemlig af en standset Bevægelse; men den formelle
Logiks Distinctioner og Rubrikker bære ikke Spor af Bevægelse;
dens Sætninger ere som færdige Figurer, Streger og
Kryds, hvis Tal og Orden maae indpræntes i Hukommelsen.

Heri er en stor Misviisning. Vel synes det, som om
de mange Abstractioner, Begrebsopstillinger, nøiagtige Definitioner
o. s. v. netop maatte tjene til at udvikle Forstanden,
øve Skarpsindigheden og, som Herbart udtrykker det, „vænne
os til vedholdende Opmærksomhed paa Tænkningens naturlige
Gang“; men disse Fordele ere af meget tvetydig Art. Hvad
kan det hjælpe, at Forstanden udvikles, naar den udvikles i en
% --- p. 67 ---
\opage{67}skjæv Retning, idet den udvikles til at adskille det Tænkte fra
Tænkningen, de logiske Former fra deres tilsvarende Functioner.

Det er ikke et een Gang for alle færdigt Schema over
Dommenes Inddelinger, men en Udvikling og Motivering af
den Formforskjellighed, hvorunder Dommen vexelviis fremtræder,
der aabner os Indblik i de logiske Domsfunctioner.
Allermindst kan den scholastiske Formalisme „vænne
os til vedholdende Opmærksomhed paa Tænkningens naturlige
Gang“; det er tvertimod Tænkningens unaturlige Gang, hvormed
Formalismen fixerer Opmærksomheden. Thi hvad kan være
unaturligere end at hexe de levende Former ind i en krystalliseret
Væren? hvad kan vel være unaturligere og mere pedantisk
end den Maade, hvorpaa den formelle Logik ved sine Distinctioner
sætter Abstractum mod Abstractum, uden Anelse om, at
den levende Tanke i hvert Nu ligesaa nødvendig fordrer Abstractionens
Ophævelse som Abstractionen selv? Dersom det
virkelig var Tankens naturlige, ikke dens unaturlige Gang,
hvorpaa den scholastiske Logik med sine barbariske Memorialvers
henvendte Opmærksomheden, saa maatte man i videnskabelige
Fremstillinger, navnlig i enhver philosophisk Afhandling,
jo see Frugterne af Logikens Lærdomme; men den Forfatter,
der for Alvor vilde lægge an paa at lade Logikens
Lærdomme skinne igjennem en logisk Behandling af et concret
Emne, hos ham vilde Tanken nok omtrent komme til
at gaae en ligesaa naturlig Gang som hos en Magister Stygotius
eller Else David Skolemesters\footnote{% note *)
Man kan ligesaa vel blive pedantisk derved, at man med smaalig
ufrugtbar Skarpsindighed efterjager det abstract Logiske, og vilkaarlig
fastsætter et Schema for det „udi egen Indbildning“ philosophisk
Stringente, som derved, at man synlig overanstrenger sig med at
udtrykke Alt ret ceremonielt, figurligt, sindrigt og ziirligt. Er det
komisk, naar en Stygotius som Frier retter sit Udtryk: „trækker
% note breaks here, p. 67/68
mit Hjertes Jern“ til „støder mit Hjertes Jern“, da er det ogsaa
komisk, naar en Stygotius som Logiker bliver sky for en kategorisk
Dom, og bilder sig ind, at der lurer en Tautologie i den; Expectorationer
som: „Guld er, at jeg skal tillade mig en implicite tautologisk
Sætning --- thi naar vi see paa den Maade, hvorpaa Begrebet
Metal har dannet sig, og paa det, der i Almindelighed forstaaes ved
Metal, ville vi ikke kunne Andet end kalde den Sætning „Guld er et
Metal“, en implicite tautologisk Sætning, da den i \emph{Grunden}
eller, hvad vilde jeg sige, \emph{overhovedet} hverken siger Andet eller
Mere, end at Guldet er en saadan Ting som Guld, Sølv, Kobber
o s. v. --- et Metal“, vilde være Stiil for en Stygotius.
% PRINTER: printed "o s. v." (no point after "o"), sense "o. s. v."

Redacteurer af den formelle Logik (thi en Langhalm, der kun
tærskes igjen, fordi den er tærsket saa tidt, kan da vel neppe benyttes
til productivt Forfatterskab) gjøre sig ikke sjeldent skyldige i den komiske
Modsigelse, at medens de i det saakaldte Identitetsprincip, hvor Tautologien
er iøinefaldende, ikke kunne opdage det Tautologiske, men
alene for „Vishedens“ og „Sandhedens“, „Determinationens“,
„Statueringens“, „Hvadhedens“, „Eenstemmighedens“ og „Samstemmighedens“
Skyld fryde sig over „det dybere Logiske“, der ligger i
Logikens øverste Sætning, at A er hvad A er, faae de derimod
Skrupler, naar de skulle erkjende det for en Dom, at Guld er et
Metal, thi det kunde jo muligviis være en Tautologie. Det falder
naturligviis ikke den pussilanime Skarpsindighed ind at tænke
igjennem og lade Sagens Afgjørelse komme an paa en Dialektik;
thi da vilde det snart klare sig. Ere følgende Domme: Guld er
et Metal, Sølv er et Metal, Kobber er et Metal, implicit tautologiske,
hvad følger saa heraf? Saa følger heraf, at Guld er
Sølv, og at Sølv er Kobber, en Consequens, hvis Urimelighed er
den naturlige Forstand ligesaa indlysende som dens Nødvendighed.}.

% --- p. 68 ---
\opage{68}Til Tankens naturlige Gang hører, at de sprogformende og
de tankeformende Functioner indgaae i levende Vexelvirkning;
kun saaledes kan Forestillingen gaae over i Begreb, og det
Billedlige belyses ved det Ubilledlige. Logikens Betydning er,
at Tanken faaer Øie paa sig selv, at Abstractionen gjennemføres.
Først dømmes der om Gjenstanden, og Dommens
Form afsættes i Sproget; den tankeformende Function er
% --- p. 69 ---
\opage{69}latent i den sprogformende\footnote{% note *)
De sprogformende Functioner kjendes nemlig, hvad Stilen angaaer,
ikke umiddelbart derpaa, at her dannes enten nye Ord eller nye grammatikalske
Vendinger, men derpaa, at de givne Sprogmidler benyttes
paa eiendommelig Maade. Til enhver eiendommelig
Sprogvending svarer en eiendommelig Tankevending. Den, der
siger: „alt Metal er smelteligt“, fører et andet Sprog og tænker
Tanken anderledes, end den, der siger: „dersom dette er et Metal,
saa er det smelteligt“; denne fører igjen et andet Sprog og tænker
Tanken anderledes end den, der siger: „dersom denne Masse ikke
lader sig smelte, kan det umulig være Metal“.}. I Dommens sproglige Form
seer Tanken saa sit eget Spor, et Aftryk af sit usynlige Væsen,
og begynder nu med en Dom over Dommen: den sprogformende
Function indleder den tankeformende. Er Tanken først ved
egne Sprogvendinger bleven opmærksom paa det Abstracte og
har vundet Sikkerhed i Abstractionen, da er det saa langt fra,
at fyldige, mere figurlige Udtryk skulde svække Idealitetens Lysvirkning,
at Reflexionen, den gjennemførte Reflexion, tvertimod
har en Tilfredsstillelse af at slaae ind paa Anskuelsens Vei
og forvandle det Billedlige til en klar, i Farven gjennemsigtig
Form for det Ubilledlige.\footnote{% note **)
Det er vidunderligt, hvorledes Anskuelsen og Tanken, det Billedlige
og det Ubilledlige, kunne ved Ordets Trylleri indslynges i
hinanden, saa at Tankens Klarhed og Billedets Dristighed staae i
ligefremt Forhold. Vexelvirkningen af sprog- og tankeformende Functioner
er da æsthetisk. --- „Min Sorg det er min Ridderborg, der som
en Ørnerede ligger høit oppe paa Bjergenes Spidse mellem Skyerne;
Ingen kan storme den. Fra den flyver jeg ned i Virkeligheden og
griber mit Bytte; men jeg bliver ikke dernede, mit Bytte bringer
jeg hjem, og dette Bytte er et Billede, jeg væver ind i Tapeterne
paa mit Slot. Da lever jeg som en Afdød. Alt, hvad der er
oplevet, dukker jeg ned i Glemsels Daab til Erindringens Evighed.
Alt Endeligt og Tilfældigt er glemt og udslettet. Da sidder jeg
som en gammel, graahærdet Mand tankefuld og forklarer Billederne
med sagte Stemme, næsten hviskende, og ved min Side sidder
et Barn og lytter til, skjøndt det husker Alt, førend jeg fortæller
% note breaks here, p. 69/70 (fortæller divided for-|tæller)
det.“ Victor Eremita. Enten-Eller. Kjøbenhavn 1849.
S. 21. --- Hvad skulle vi nu fra Logikens Standpunkt dømme
om en saadan Fremstilling? Er her kun Billeder uden Tanker,
saa er vel det Hele en fad Opdigtelse, og vi maae med Mester
Mannegrimm bede om „en Løgn endnu“. Er her derimod ligesaa
mange, ja endnu flere Tanker end Billeder, ere de farverige
Billedsætninger væsentlig klare Tankesætninger, skulde Tankerne
da maaskee vinde ved at afføres den rige Dragt og fremstilles
i deres logiske Nøgenhed? Men dersom dette Sidste hverken
er ønskeligt eller i strengere Forstand muligt, saa have vi jo her
et Exempel paa æsthetisk Vexelvirkning af sprog- og tankeformende
Functioner.}

% --- p. 70 ---
\opage{70}Ogsaa det logiske Udtryk kan have sin væsentlige Skjønhed.
En Fremstilling, der i Eet og Alt skulde tilfredsstille
Logikens æsthetiske Fordring, er vistnok en stor Sjeldenhed,
da den forudsætter en egen, lykkelig Forening af Reflexion og
Phantasie; men at Logiken, den strenge Logik, gjør æsthetiske
Fordringer, er vist: det Æsthetiske er her begrundet i den
levende Tænknings Tilegnelse af Alt, hvad der rører sig i
Bevidsthedslivet; det er dette Æsthetiske, der beaander Tanken,
giver Afskygningerne Farveskjær og gjør Productiviteten kjendelig
paa den tiltrækkende Harmonie af sprog- og tankeformende
Functioner.

\textit{γγ)} \emph{Ideeformende Functioner.} Mellem Subjectiviteten
og Objecternes virkelige Verden er Kløften saa dyb,
at den aabner sig paany, hver Gang Reflexionen drømmer
om at have faaet den udfyldt. \emph{Functionernes Logik}
endte med en Concretion af Ideal- og Realsystemets teleologiske
Eenhed, nemlig som Slutningseenhed af Natur og
Aand, en Udvikling, hvori Dualismen af det Objective og det
Subjective foreløbig fandt sin Forklaring. Fra flere Sider
er det da ogsaa forud blevet viist, hvorledes det Antilogiske
ved Sandsningens Mægling lod sig omsætte i det Logiske;
hvorledes det samme Indhold kunde antage forskjellig Form,
den samme Form forskjelligt Indhold o. s. v. Men under
% --- p. 71 ---
\opage{71}alle disse Vendinger, Omsætninger og Forklaringer er der
underforstaaet et Noget, som, just fordi det var i umiddelbar
Anvendelse, ikke strax kunde blive Gjenstand for særlig Reflexion,
nemlig Ordet.

Hvad er da Ordet? At Ordet ikke uden videre er det
Samme som Sproget, er klart af vort Foregaaende; thi der
gives jo Sprog uden fritstaaende Ord, og det Sprog, der har
særlig formede Ord, har jo mange og forskjellige Ord. Kategorien:
\emph{Ord} staaer da med Hensyn hertil deels lavere,
deels høiere end Kategorien: \emph{Sprog}. Lavere! thi ligesom Bogstaverne
ere Elementer i Stavelser, Stavelserne Elementer i
Ord, saaledes er det enkelte Ord igjen Element i Sætningen,
og Sætningen Element i Sproget. Men som Sætningselement
er Ordet endnu ikke væsentlig Ord; et egentligt Sprog
kan ikke undvære Sætninger, men en Sætning kan forvandle
Rødder til Ord: at der kan være Sprog uden Ord, er da
sandt, forsaavidt der kan dannes Sætninger uden Ord. Høiere!
thi Ordet er Tankens Udtryk. Forsaavidt Sætningen giver
en Mening, og Meningen en Tanke, er Sætningen væsentlig
et Ord. „Der er et Ord, et gammelt Ord, som siger“ o. s. v.
At Sætningen efter sin Mening og formedelst Tanken er Ord,
vil da med andre Ord sige, at Ordet som Ord i dybere Forstand
er høiere end Sætningen. Fra dette Synspunkt kan
der altsaa ikke gives Sprog uden Ord; det maatte da være
et Sprog uden Sætninger, et reent Dyresprog, eller et Sprog
af Sætninger uden Mening, et meningsløst Sprog.

Staaer nu Ordet i dybere Forstand høiere end Sætningen,
saa staaer det ogsaa høiere end Sproget. Tage vi
Ordene som Sætningselementer, altsaa endeligt og i Fleertallet,
saa kan man jo vistnok sige, at det er Sproget, som
ved en særegen Udvikling gjør sine Elementer til Ord; men
tage vi Ordet som Ord i absolut Betydning, altsaa i Enkelttallet,
saa maae vi sige, at det er Ordet, der er det Oprindelige,
Sproget det Afledede; at det er Ordet, der gjør
% --- p. 72 ---
\opage{72}Sproget til Sprog. Hvor urimeligt det er at ville gjøre
Sproget til det absolut Oprindelige, --- hvad de jo dog gjøre,
der med F. Schlegel henføre Sprogenes Oprindelse til en
overnaturlig Act, et Mirakel, --- kan allerede føles, naar man
istedenfor det Johanneiske: „I Begyndelsen var Ordet, og
Ordet var hos Gud, og Ordet var Gud“, forsøger at sætte:
I Begyndelsen var Sproget, og Sproget var hos Gud, og
Sproget var Gud.

Er det nu ikke Sproget som endeligt Sprog, men Ordet
i absolut Forstand, der er det Oprindelige, saa indsees heraf,
at hvad der i endelig og relativ Forstand udrettes ved Sproget,
maa i absolut og oprindelig Forstand tilskrives Ordet.

Hvad skylder den talende Subjectivitet da egentlig Sproget?
Svaret lyder fra det Foregaaende: en Objectiveringsfordobling,
hvorved Modsætningen af den indvortes og den udvortes
Verden væsentlig articuleres! Vistnok vilde Skikkelserne fra
Omverdenen ogsaa uden Sprogets Mægling kunne speile sig
i Bevidstheden --- vi see det jo paa Dyrene; --- men denne
Speiling maatte da foregaae: deels saaledes, at hvert Billede,
om end kun for et Øieblik, blev staaende i antilogisk Omrids,
forholdende sig til Gjenstanden som Copien til sin Original;
deels saaledes, at Billederne opløste sig i Forestillinger ligesom
Skyer i Taager, idet Forestillingernes Afløb paa den ene Side
foregik med aandelig bestemt Naturnødvendighed, paa den
anden Side tabte sig og svandt i tomme, svævende Reflexioner.
Kun ved Sprogets Hjælp bliver det Bevidstheden muligt at
articulere Tilværelsens objective Indhold i subjectiv Form og,
omvendt, at give Subjectivitetens Indhold objectiv Form; kun
ved Sprogets Hjælp kan det samme Indhold indgaae i dobbelt
Form, og den samme Form i dobbelt Indhold; kun ved
Sproget bliver det muligt at forme Ideen, det absolute
Indhold.

Hvad der i relativ Forstand udrettes ved Sproget, maa,
som sagt, i absolut Forstand henføres til Ordet: Forholdet
% --- p. 73 ---
\opage{73}mellem Indholdet og Sproget forudsætter et Grundforhold
mellem Ideen og Ordet. Da Ideens absolute Indhold paa
een Gang er baade subjectivt og objectivt, blive de ideeformende
Functioner særlig at bestemme ved Forskjelseenheden af det
Subjective og det Objective. Forsaavidt denne Eenhed viser
sig som en \emph{magisk} Eenhed af Idee og Ord, ere de ideeformende
Functioner \emph{mythiske}; forsaavidt Ordets Eenhed med
Ideen har et \emph{phantastisk} Præg, ere de \emph{gnostiske}; sættes
Eenheden i \emph{Guddommens Væsen}, blive de \emph{theosophiske}.

\emph{Mythiske Functioner.} Det er Forskjelseenheden af
det Objective og Subjective ved Ordet, hvorpaa Eftertrykket
falder. Væsentlig kommer det ved denne Forskjelseenhed an
paa Indholdet. Er Indholdet endeligt, og Adskillelsen mellem
Subjectivt og Objectivt uklar, saa have vi en i Fornemmelser
og dunkle Forestillinger gjærende, famlende Bevidsthed for os,
et Bevidsthedstrin, der svarer til Barnets i den tidlige Alder:
dette Standpunkt kan ikke her komme i Betragtning. Er
Indholdet endeligt, og Adskillelsen mellem Subjectivt og Objectivt
klar, saa have vi det kritisk-forstandige Bevidsthedstrin.
Paa dette Trin er Ordet Et og Tingen et Andet, Forestillingen
Et og Virkeligheden et Andet, og Sproget mægler saa paa
endelig Maade imellem Factorerne i denne Udstykning: dette
Standpunkt kan heller ikke komme i Betragtning. Det for
de ideeformende Functioner Afgjørende er, at Bevidsthedsindholdet
maa være absolut.

Da nu et absolut Indhold umulig kan bevæge sig gjennem
Subjectiviteten, uden at reflectere sig i relative Modsætninger,
saa indsees endvidere, at den Bevidsthed, der skal objectivere
sig Ideen i og ved Ordet, og deri erkjende det Absolute,
allerede forud maa være naaet saavidt, at den i endelig
Forstand kan skjelne Tingene i Omverdenen, og i det Hele
gjøre Forskjel paa Virkelighed og Sandsebedrag, paa hvad der
er udvortes, og hvad der kun er indvortes. Den psychologiske
Subjectivitet, som, efter at have naaet en inden snevre
% --- p. 74 ---
\opage{74}Grændser klar Forstaaelse af Tilværelsen i det Endelige,
gribes og bevæges af Ideen paa saa umiddelbar Maade, at
Endelighedens Modsætninger af Drøm og Virkelighed, af
Indvortes og Udvortes forsvinde, for at gjentage sig i høiere,
til det Absolute formedelst Ordet svarende Former, er \emph{mythisk}.

Kjendemærket paa det Mythiske er, i Principet, ingenlunde
Uklarhed over det udvortes Endelige; det er ikke Sandselighedens
Taagebilleder, men Ideens Lynslag, ikke Forstandens
Nat, men Fornuftens Daggry, der laaner Bevidsthedsindholdet
de vidunderlige mythiske Farver. Hvor eiendommelig Farveblandingen
er, falder i Øinene, naar man ret vil charakterisere
Mythologien og det mythiske Standpunkt. „For Mythologien“,
siger Heiberg, „gives der intet andet, Alt omfattende Synspunkt,
end det poetiske. Mythologien er Poesie, og reent oprindelig
Folkepoesie“. Senere hen bliver saa denne Folkepoesie
tillige opfattet som en Art speculativ Philosophie. „Enhver
Mythologie“, hedder det, „begynder med Kosmogonie;
thi den philosopherende Menneskeheds første Bestræbelse er at
forklare Tilværelsens Under. Det første Spørgsmaal er:
hvorledes er Verden blevet til? Hvad var der, førend Alt
dette blev til? Eller hvad var der, da der endnu Intet var?
Forstanden kan saalidet besvare disse Spørgsmaal, at den endogsaa
maa erklære dem for taabelige og sig selv modsigende.
Men Phantasien svæver dristigt hen over Tingenes Begyndelse,
uden at frygte for at tabe sig i Intet. Den befolker Tomheden,
den lader Lyset opstaae af Mørket, Tilværelsen af Ikketilværelsen“\footnote{% note *)
J. L. Heibergs nordiske Mythologie. Paa Dansk udgivet af Chr.
Winther. Kbhvn. 1864. S. 4 og S. 24.} o. s. v.

At Mythologien paa een Gang er baade Poesie og Philosophie,
vil da med andre Ord sige, at den i Grunden hverken
er Poesie eller Philosophie; ikke Poesie, fordi den er Philosophie;
ikke Philosophie, fordi den er Poesie. Det Nærmere sees
% --- p. 75 ---
\opage{75}af Functionerne. Det Poetiske bliver her ikke opfattet som
Poesie. Poesien forudsætter en paa det absolute Indhold
lydende Adskillelse af Illusion og Virkelighed, en Adskillelse, der
udtrykkelig sætter Illusionen over Virkeligheden, fordi den
sætter Ideen over Endeligheden; men ved en saa reflecteret Adskillelse
kan Mythen som Mythe umulig bestaae: de mythiske
Skikkelser holde sig ved Ideens Magie paa en høiere Virkeligheds
Trin; her er paa dette Trin end ikke Anelse om Illusion.
Ligesaalidt bliver det Philosophiske i Mythedannelsen
opfattet som Philosophie. Hvad skulde det vel være for en
Art Philosophie, som begyndte med en udtrykkelig Splid mellem
Forstand og Phantasie, og som paa Grund af denne Splid
opgav Forstanden og lod Phantasien raade? Istedenfor at
begynde med en Dualisme af Phantasie og Forstand, begynder
Mythedannelsen tvertimod med en af Forstand og Fornuft
imprægneret Phantasie; man kan om den mythiske Ideebevægelse
med fuld Ret sige, at Tanken phantaserer, fordi
\emph{Phantasien tænker}.

Naar da den af Forstand og Fornuft imprægnerede
Phantasie „svæver dristigt hen over Tingenes Begyndelse,
uden at frygte for at tabe sig i Intet“, hvad styrer da dens
Flugt, hvad er det, der bevarer den fra at styrte i Tomhedens
Afgrund? Væsentlig er det \emph{Ideen}; umiddelbart seet, er det
\emph{Ordet}. Begyndelsens Afgrund er lige dyb, enten vi begynde
med \emph{Chaos} eller med \emph{Ginungagap}, den er i begge uendelig
dyb; men i det uendelige Dyb formæles Ideen med
Ordet.

Begynde vi med Chaos, da er det naturligviis ikke den
tomme Benævnelse, hvorpaa her lægges Vægt, thi Benævnelserne
kunne vexle. Men de eenstydige Benævnelser lyde
alle paa det samme Grundord. At Chaos nu som Grundord
er en Ideens Aabenbarelse, sat og bevæget af Ideens Magt,
kan let paavises. Chaos er i mythisk Forstand hverken et
abstract logisk Begreb eller et Udtryk for en Tilstand af physisk
% --- p. 76 ---
\opage{76}Forvirring, men en i Ordet fastholdt Totalitetseenhed, et
subjectivt-objectivt Første, der i billedlig og dog ubilledlig
Form anskueliggjør, hvad ikke kan anskues, nemlig Tingenes
Oprindelse, Ideen af al virkelig Begyndelse eller den virkelige
Begyndelses Idee.

At Chaos som Begreb ikke er i Begrebets, men i et
ubilledligt Billedes Form, følger udtrykkelig af Anskuelsens
Eenhed med Ordet. Hvad er der i Chaos? Intet, og dog
Alt! Denne Tankens Selvmodsigelse med de deraf udspringende
Conseqvenser opfattes ikke i Form af en Tanke, men
sees i et Billede, et mægtigt Tankebillede, dannet af Ordet.
Idet Tanken ved Ordets Hjælp phantaserer over Intet og Alt,
tænker Phantasien væsentlig paa Alt og anskuer dog i Chaos
Intet, det Tomme; den tænker væsentlig paa Intet, og anskuer
dog i Chaos Alt, nemlig Spirerne til Alt (\textit{non bene junctarum
discordia semina rerum}) Indholdet, det i Chaos % PRINTER: no full stop printed after "rerum)", sense "rerum). Indholdet"
indholdsløse Indhold, er væsentlig subjectivt formedelst Ordet:
et Product af phantaserende Tænkning, af tænkende Phantasie;
men det anskues tillige objectivt, slaaer den anskuende Bevidsthed
med Objectivitetens Magt i Ordet, og antager Skikkelse
af physisk umiddelbar Realitet.

Men at Chaos i objectiv Mening desuagtet ikke har physisk
umiddelbar Realitet, er netop en Følge af det i Billedet Ubilledlige,
af Anskuelsens mythiske Eenhed med Ordet. I objectiv
physisk Betydning skulde Chaos jo være den virkelige Opløsning,
hvori Altilværelsen maatte befinde sig, naar alle Former
vare forstyrrede, og Atomernes Vrimmel i regelløs Forvirring.
For en forstandig kritisk Bevidsthed om et saadant Chaos har
Ordet tabt sin magiske Kraft. Ordet Chaos er i dette Tilfælde
kun et Navn; Chaos derimod som udvortes virkelig Tilstand er
Tingen, og Tingen, ikke Ordet, er det, hvorpaa det kommer
an. Adskilt fra den paa Ordet hvilende Phantasie, vil Forstanden
da paa egen Haand gjennemforske det virkelige Chaos,
Tingenes udvortes virkelige Begyndelse. Men dette For%
% --- p. 77 ---
\opage{77}standens Forehavende vil totalt mislykkes. Som virkelig
Tilstand er Chaos aldeles uvirkelig; der er virkelig ingen Tilværelse,
der kan begynde med et virkeligt Chaos; hvad Forstanden
skal faae ud af et saadant Chaos, maa den i Virkeligheden
selv indsmugle deri: det er en ørkesløs Leg, Forstanden
leger med sig selv og det Virkelige.

I det mythiske Chaos er Objectiviteten subjectiv; det objective
Indhold er ved Ordet omsat i subjectiv Form. Hvor
forskjelligt det mythiske Chaos er fra det udvortes physiske, er
da ogsaa kjendeligt paa de Skikkelser, der fremgaae af Chaos.
„Først er Chaos, dernæst den udstrakte Jord, det mørke Tartaros
--- og Eros, den skjønneste blandt de udødelige Guder“.
Den, der seer de gamle Guder udtræde af Chaos, kan ikke
være i Tvivl om, hvad Chaos er for et Væsen. Idet her
begyndes med Chaos, tages her paa en Maade Tilløb til en
objectiv Fremstilling af Verdens Tilblivelse, af Tingenes, de
endelige Realiteters Dannelse; men Synet vender sig. Kosmogonien
slaaer om og bliver Theogonie; Chaos bliver en
Guddom, den Guddom, af hvis Skjød de særlige Guder
fødes. Denne Tilbagegang af det Kosmogoniske i det Theogoniske
er for Mythedannelsen i høi Grad charakteristisk. For
at Verden kan blive til, maae Guderne først blive til; for at
de endelige Realiteter i Virkeligheden kunne blive til, maae Realideerne
først blive til: Realideernes endelige Vorden reflecterer
i sig Grundideernes evige Væren.

Chaos er altsaa i sin Idealitet evig, i sin umiddelbare
Realitet timelig, Guddom i Væren, Naturgrund i Vorden,
Totalitet i oprindelig Eenhed, Totaliteternes Grund i oprindelig
Fleerhed; i Chaos er det Relative absolut, det Absolute
relativt; i Chaos-Ordet er Bevægelsen Hvile og Hvilen
igjen Bevægelse, Fylden Tomhed, og Tomheden Fylde, det
Indvortes et Udvortes, og dette Udvortes igjen et Indvortes:
i Chaos er Ideen Ord, og Ordet Idee.

% --- p. 78 ---
\opage{78}Det er da ikke en Tilfældighed, men en væsentlig Nødvendighed,
at de ældste Guder udgaae af Chaos. I logisk
Form er Nødvendigheden kjendelig paa, at Fleerheden fremgaaer
af Eenheden, efterdi Eenheden selv oprindelig er Fleerhed,
at det Særskilte i Totaliteten fremgaaer af det Almene,
efterdi det Almene oprindelig er en Totalitet af Særskilte.
Men i den logiske Form mangler Indholdet umiddelbar Realitet;
Realiteten, den mythiske Realitet, kommer med Ordet,
thi det ved Ideen bestemte Ord bestemmer Ideen.

Saadanne Skikkelser som Gæa, Tartaros, Erebos, Nyx
og Eros ere i deres Oprindelighed mere end tomme Personificationer.
„Skikkelser af det Skikkelsesløse“, „elementære Former
af det første Ene, i sig selv Formløse“, slige fastende Abstractioner
falde ligesaa langt under Ideen som under Virkligheden. % PRINTER: printed "Virkligheden", sense "Virkeligheden"
Nei, ved Ideens Magie i det mythiske Ord ere de ældste
Guder, trods det i Formen Uformelige, typiske Skikkelser,
fulde af Realitet, hver Guddom en Totalitet i sig, et eget
Væsen, et bestemt Physiognomie. Det er Ideen, der i særskilte
Totalitetseenheder gjennemtrænger, bevæger og bestemmer
Ordet; men det er i umiddelbar Forstand Ordet, det særskilte
Grundord, der i sin Eenhed optager, bestemmer og samler
Mangfoldigheden af de ideelle og reelle, ubilledlige og billedlige
Momenter, der udgjøre Ideens absolute Indhold.

Gudinden Gæa, der faaer Navn efter Jorden, optager
ved dette Navn al Jordens Realitet i sig, bliver ved dette
Navn en Guddom med umaadelige Evner, vældige Kræfter
og store Muligheder; men at Gudinden har faaet Navn efter
den virkelige Jord, siger med det Samme, at det er den
virkelige Jord, der har faaet Navn efter Gudinden. Jordgudinden
og den virkelige Jord have ikke blot eet Navn, men
de ere eet Væsen, betegnede ved eet Ord: de ere i Ordet Eet,
efterdi de ere i Ideen Eet. At den virkelige Jord er Gudinde,
viser os ret paa en slaaende Maade, hvad det vil sige,
at Ideens objective, reelle Indhold er --- formedelst Ordet ---
% --- p. 79 ---
\opage{79}omsat i subjectiv, ideel Form; at Gudinden med alle Træk
af elementær Subjectivitet er Jorden selv, den virkelige Jord,
anskueliggjør paa den anden Side, hvad Ordet formaaer, naar
det gjælder om at oversætte et subjectivt Indhold, --- thi Gudindens
Væsen er i sin Idealitet et subjectivt Væsen, --- i
umiddelbar, objectiv Form.

Ikke blot en positiv Størrelse som Jorden, selv saadanne
negative Sandseformer som Mørket og Natten, Væsenheder
som Tartaros, Erebos og Nyx indsuge Ideens Realitet
gjennem Ordet, individualiseres i Ordet og blive levende
Skikkelser formedelst Ordet. Gjennemaandet og bevæget
af Ideen har Ordet ikke blot den eiendommelige Magt,
at det som Sprogets Høieste og Dybeste formaaer at samle
de Relativitetsbestemmelser, der i Talens Sætninger deelviis
fremstilles, beskrives og udvikes, i prægnante Eenhedsudtryk; % PRINTER: printed "udvikes", sense "udvikles"
men som Led i Sætningen, som Element i Sproget har Ordet
tillige den Smidighed, Bøielighed, Bevægelighed og Relativitet,
at det i de forskjelligste Forbindelser kan være snart for
Sandsen, snart for Tanken, snart figurligt, snart bogstaveligt,
kan forene, hvad det selv har adskilt, adskille, hvad det selv
har forenet, kan, Led for Led og Punkt for Punkt, opklare det
Reelle ved det Ideelle, det Billedlige ved det Ubilledlige, det
Umiddelbare ved det Reflecterede. Er Ordet først bleven en
Gud, og Guden et Ord, da staaer det i Subjectivitetens
Magt at lade denne Ordets Gud optræde i de eventyrligste
Situationer, undergaae de mærkværdigste Forvandlinger, vise
sig fra de forskjelligste Sider, tænkende, attraaende, handlende,
lidende.

Det mythiske Ord, der forvandler de særskilte Ideer til
Gudeskikkelser, forvandler med det Samme de ideale Bevægelser
til virkelige Begivenheder; denne Forvandling er med
Hensyn paa de ideeformende Functioners mythiske Væsen afgjørende
charakteristisk. Idet Ideernes Strid med den materielle
Virkelighed og gjennem Virkelighedens elementære
% --- p. 80 ---
\opage{80}Modsætninger med hinanden indbyrdes forvandles til en Historie
om Gudernes Kampe, er den magiske Eenhed af Idee
og Virkelighed paa Ordets Veie ført hen til det Vendepunkt,
da den mythiske Identitet af det Ideelle og Reelle slaaer over
i Dualisme, og Illusionen viser sig.

Den colossale Dristighed, hvormed den tænkende, af
Ideen imprægnerede Phantasie ved Syn af virkende Guder
trodser Sandsynligheden, er et talende Vidnesbyrd om Ordets
Magt i de objectiverede Begivenheder. Gudealderen, den
fjerne Fortid, er Ideernes, de ideale Aarsagers gyldne Tid;
Nærværelsen med sine Kjendsgjerninger er Virkelighedens,
Betingelsernes og de materielle Aarsagers Tid: kunne de to
Tider vel blive til een Tid? I Ordets, d.\ v.\ s.\ i Fortællingens,
Erindringens, Speil reflecterer det Forbigangne sig
jo vistnok i det Nærværende som de store Grundideer i Udviklingens
virkelige Resultater; men Virkeligheden med sine
Tingsaarsager reagerer med Nærværelsens Magt mod hine
ideale Forudsætninger; Conseqvensen er en Splid mellem
Tingen og Ordet, en Splid mellem det Objective og det
Subjective, en Splid mellem Indhold og Form, Phantasie
og Tanke, en Splid, der vækker Kritiken, ophæver Standpunktet
og stiller de ideeformende Functioner nye og høiere
Opgaver.

\emph{Gnostiske Functioner}. I Principet er Mythologie
forskjellig fra Poesie; thi Mythen, der har magisk Objectivitet,
er uden Illusion, medens Poesien som Poesie væsentlig forudsætter
Bevidsthed om Illusion. Da Bevidstheden imidlertid
kan skride continuerlig frem fra dunkel Anelse til klar Erkjendelse,
kan som i en vaagen Drøm paa een Gang være
i Illusionens Magt og have en Viden om Illusionen: saa er
det forklarligt nok, at Mythologie ligesaa umærkeligt har kunnet
gaae over i Folkepoesie, som Folkepoesien igjen i Kunstpoesie.
Men derfor bliver Adskillelsen af Illusion og Virkelighed dog
% --- p. 81 ---
\opage{81}ligefuldt afgjørende for det Bevidsthedstrin, hvormed Poesien
begynder.

Alt efterhaanden som det mythiske Indhold forklares i
Poesiens Former, vil Forstanden, der seer Illusionen, løsrive
sig fra Phantasien, gjøre Front mod det Virkelige og
stræbe at orientere sig i Tingenes Verden. Men Tingenes
Verden er sandselig, endelig, materiel og forsaavidt ideeløs.
Skilt fra Phantasien er Forstanden tillige skilt fra Ideen.
Men den ideeforladte Forstand har væsentlig Trang til Idealitet;
Forstanden er i Principet fornuftig: her begynder Philosophien.
Poesien med sine ideeopfyldte Illusioner og Philosophien
med sine formale Begreber danne nu Extremer med
den sandselig concrete Verden imellem sig. For at hævde
Ideen, arbeide de ideale Magter, Poesien og Philosophien, hver
fra sin Side, paa at opløse, forvandle og forklare det udvortes
Virkelige; men Virkeligheden er stærk, og Kampen haard.
Til Forstærkning i den haarde Kamp føler da den i Almeenbegreber
afmattede Reflexion en Trang til at fylde de tomme
Former med anskueligt Indhold ved at fornye den gamle Pagt
mellem Anskuelse og Tænkning, Reflexion og Phantasie; paa
disse Vilkaar blive de ideeformende Functioner \emph{gnostiske}.

Det er Gnosticismens Væsen, og ikke en bestemt historisk
Skikkelse af Gnosticismen, vi her have for Øie; kun forsaavidt
kan der tages Hensyn til det historiske Phænomen, som dette
Phænomen er i særegen Forstand egnet til at gjøre Væsenet
aabenbart. Den retablerede Forening af Forstand og Phantasie
er kjendelig paa Ordet, Ordet kjendeligt paa de Kategorier,
der bære Anskuelsen. Her begyndes ikke væsentlig med Chaos.
Chaos er platonisk reduceret til den rene Tomhed, et Udtryk
for det Sandseliges væsentlige Intethed; saaledes er Chaos
Intetheden, den negative Væren, \emph{Kenoma}. At det Materielle
og Sandselige i Altilværelsen er reduceret til Kenoma,
er i Gnosticismen en Antydning af, at Begyndelsen, Udviklingens
Realbegyndelse, ikke kan søges i det physisk Objective, i
% --- p. 82 ---
\opage{82}sandselig Realitet, men maa søges i det væsentlig Subjective,
i overgribende Idealitet. I kategorisk Modsætning til \emph{Kenoma},
det objectivt sandselige Tomme, staaer da \emph{Pleroma},
Ideeriget, Magtens og Fyldens Region.

Forsaavidt den gnostiske Tænkning hengiver sig til Phantasieanskuelser,
er Pleroma et objectivt Tankebillede i Lighed
med det mythiske Chaos; men en nærmere Betragtning af
dette Tankebillede viser os snart, at vi i Pleroma have med
et Subjectivitetens, et Aandelighedens og Inderlighedens Rige
at gjøre. Pleromas oprindelige Eenhed er i \emph{Dybet} og
\emph{Tausheden} (βυθὸς og σιγὴ). I og for sig kan nu vistnok
Dybet ligesaa vel betegne noget Objectivt som noget Subjectivt;
men allerede \emph{Tausheden}, ikke Stilheden, lader os
ane, at Dybet er i Subjectiviteten. Anelsen bliver til Vished,
idet Modsætningens første abstracte Form nærmere bestemmes
ved en følgende, concretere Modsætning, Modsætningen
mellem \emph{Urfaderen} og \emph{Tanken} (προπάτωρ og
ἔννοια). Kaste vi derhos et Blik paa de stedse tiltagende
Concretioner, hvori den oprindelige Modsætning trinviis gjentager
sig, Modsætningen af \emph{Fornuften} og \emph{Sandheden}
(νοῦς og ἀλήθεια), \emph{Ordet} og \emph{Livet} o.\ s.\ v., saa erkjende
vi i Pleroma et Rige af Subjectiviteter og ontologisk subjectivt
Indhold.

Hvor eiendommelig Sammensmeltningen af Billede og
Begreb er i det gnostiske Pleroma, falder i Øinene, naar vi
mærke os den Umiddelbarhed, hvori Ordbetydningen af de
opstillede Kategorier, ja endog Ordenes Kjøn er opfattet og
benyttet. At Kategorierne: Dybet og Tausheden, Urfaderen
og Tanken, Fornuften og Sandheden, o.\ s.\ v. ere hypostaserede
Væsenheder, logiske Personificationer, kan sees, ogsaa af en
Oversættelse; men de sproglige Tilfældigheder, som den phantaserende
Tænkning har benyttet ved disse Personificationer,
opdages kun ved et Blik paa Grundtexten. Benævnelserne
paa de pleromatiske Hypostaser ere nemlig valgte saaledes, at
% --- p. 83 ---
\opage{83}den ene altid er et Hankjønsord, den anden tilsvarende et
Hunkjønsord: βυθὸς, προπάτωρ, νοῦς, λόγος \dots\ ere
Hankjønsord, hvorimod σιγὴ, ἔννοια, ἀλήθεια, ζωὴ \dots\
ere Hunkjønsord.

Ved denne Deling, en Grunddeling, der i sin Art svarer
til Hegels objective Dom (das Urtheil, afledet af urtheilen)
er der tilveiebragt en vis Ligevægt imellem Reflexion og Phantasie.
Reflexionen fordrer, at Pleromas Grunddeling skal i
alle Gjentagelser afgive Modsætninger af Totalitetsbegreber.
Dette er paa en vis Maade naaet ved de angivne Kategorier;
λόγος, νοῦς, ἔννοια, ἀλήθεια ere sande Totalitetskategorier.
Ligeledes er her ved de kategoriske Modsætninger betegnet en
Fremadskriden fra det Abstracte mod det Concrete, fra det
Almene gjennem det Særskilte mod det Individuelle; thi det
maa indrømmes, at βυθὸς og σιγὴ ere mere abstracte end
νοῦς og ἀλήθεια, disse igjen mere abstracte end ἄνθρωπος
og ἐκκλησία. Ja, saa sindrig er Rækken anlagt, at Reflexionen
endog, hvis det forlanges, kan gjøre logisk Rede for Fremskridtet.
De abstractere Fordoblinger danne nemlig Grundlag
for de concretere; disse ere, for at tale i Hegels Sprog, Mediationer
af hine. Saaledes ere λόγος og ζωὴ aabenbart
Mediationer af νοῦς og ἀλήθεια, idet νοῦς formedelst ἀλήθεια
bliver til λόγος, maa ἀλήθεια formedelst νοῦς naturligviis
blive til ζωὴ, o.\ s.\ v.

Men er Reflexionen paa aandrig Maade tilfredsstillet, da
er Phantasien det ikke mindre: Sindrighed og Aandrighed ere
overhovedet Væsensmærker for al Gnosticisme, den nyere saavel
som den ældre. At Phantasien har havt sin Deel i det
aandrige Valg af Modsætningskategorier, sees af den Maade,
hvorpaa de valgte Hankjøns- og Hunkjønsord understøtte
Personificationen. Ifølge Personificationen er hvert kategorisk
Par et Ægtepar, Mand og Hustru: forgudede Former, og
dog ikke mythiske Guder!

Forskjellen mellem Mythisk og Gnostisk er nu iøinefaldende;
% --- p. 84 ---
\opage{84}hvilken Afstand imellem en pleromatisk Hypostase med sin
abstract-metaphysiske Habitus, i sin af Reflexionen sindrigt tilmaalte
Idealitet, og en af Mythologiens Guder eller Gudinder,
der individualiserer sig i det Endelige, leger med det Tilfældige
og bugner af Realitet! I Gnosticismen har Phantasien
ikke længere Tanken i sin Magt, og ligesaalidt kan Tanken
siges at have Phantasien i sin Magt. Det Forbund, hvori
begge, Phantasie og Tanke, ret skulde styrke hinanden, er betinget
af en Accommodation, hvorved de gjensidig svække hinanden:
Phantasien bliver tænksom, og Tanken phantastisk.

At den gnostiske Reflexion med en ganske anden Bevidsthed
sigter paa Idee og Idealitet end den mythiske, er en Selvfølge.
Det Sindrige og Aandrige sees da ogsaa i den Maade,
hvorpaa de særskilte Ideer repræsenteres af Hypostaserne i
Pleroma. Ogsaa i denne Henseende gjør Ordet sin Umiddelbarhed
gjældende ved Navnet. Enhver af de pleromatiske Hypostaser
fører, som bekjendt, Navn af Æon (αἰὼν). Hvad
betyder Æon? En Verdensalder og en Evighed! altsaa en
af Tiden omsluttet, i Tidens Ring indfattet Evighed, et individualiseret
Absolutum. Som individualiseret Absolutum er
da enhver af Pleromas Æoner en særlig Totalitet, et relativt
Pleroma i sig, med andre Ord: en Idee. At den ene Æon
har den anden til Modsætning, vil da sige, at den ene Idee
har den anden til Modsætning: de særskilte Æoner have hver
sit Sted i Pleroma, de særskilte Ideer have hver sin Plads i
Ideernes System.

Det er et Særkjende for de gnostiske Functioner, at de
ikke kunne forlige det Ubegrændsede med Grændsen eller magte
Concretionen. Hvad det Ubegrændsede og Grændsen angaaer,
svæver Gnosticismen uophørlig imellem den orientalske og den
græske Synsmaade. Den vil i græsk Aand individualisere
det Absolute, og man skulde da forsaavidt mene, at de meest
concrete Æoner maatte indtage den høieste Plads i Pleroma;
men tvertimod! Gnosticismen med sin orientalske Respect for
% --- p. 85 ---
\opage{85}det Ubegrændsede, det Uendelige, har en afgjort Forkjærlighed
for det Abstracte; det er netop de abstracte Æoner (βυθὸς
og σιγὴ), der indtage den høieste eller, om man saa vil, den
dybeste Plads, Centralpladsen, i „Fyldens“ Rige. Striden
imellem det Begrændsede og det Ubegrændsede, en Strid, der
oprindelig ulmer i hele Pleroma, kommer da ogsaa omsider
til aabenbart Gjennembrud i Katastrophen med Sophia.

Det er et sindrigt, dybsindigt Træk i den valentinianske Gnosis,
at Sophia, Æonernes dybt grundende Æon, denne Ideerne begribende
Idee, gribes af en uimodstaaelig Længsel efter at skue
Propator. Istedenfor en Fordybelse i det Concrete vælger
Viisdommens qvindelige, receptive Æon altsaa en Tilbagevenden
til det Abstracte, som til det oprindelige Productive;
istedenfor at erkjende det Ubestemte som negativ Bestemthed i
den concrete og derved bestemte Idee, vil Sophia, den Begrændsede,
umiddelbart favne den Grændseløse: Grændsevogteren
(ὅρος) er ikke istand til at holde hende tilbage. Sophias
Fald og Opreisning, Materien, der i Kenomas Mørke dannes
af hendes Angst, hendes Taarer og Latter; Demiurgens
(Verdensæonens) Fødsel af den faldne Sophia; Verdens Skabelse
ved Lysets Fængsling i Mørket; Forløsningen ved Sjælelysets
Tilbagevenden til Pleroma o.\ s.\ v., dette tilsammen
frembyder en Mangfoldighed af sindrige Træk, som alle vise,
hvorledes den phantaserende Tænkning har forvandlet Bevægelserne
i Ideeriget til illusoriske Begivenheder.

Med de gamle gnostiske Systemer er den videnskabelige
Kritik forlængst bleven færdig; men Gnosticismens Kjendemærker:
en tvetydig Eenhed af tænksom Phantaseren og phantastisk
Tænken, sindrige Vendinger og aandrige Udtryk istedenfor
klare Begreber og kategorisk afgjørende Bestemmelser;
Anviisning paa dybere Skuen og dramatiserende Udmalen af
overnaturlige Begivenheder istedenfor indtrængende Undersøgelse
af evige Grundforhold: disse og lignende Egenheder vedblive
% --- p. 86 ---
\opage{86}endnu\footnote{Som f.\ Ex.\ i den speculative Dogmatik.} % note *)
at hendrage Logikerens Opmærksomhed paa de gnostiske
Functioner.

\emph{Theosophiske Functioner}. I det mythiske Ord er
det Phantasien, der ubetinget har Tanken i sin Magt; i det gnostiske
Ord er der en i Tilnærmelse reflecteret Eenhed af Phantasie
og Tanke; i det theosophiske Ord er det Tanken, der
umiddelbart magter Phantasien, dog saaledes, at den, hvor det
gjælder om logiske Bestemmelser, helst gjør Brug af Phantasien.
At den theosophiske Tanke ved en Ideens Magie har
Phantasien i sin Magt, fremlyser strax deraf, at, medens det
mythisk Absolute uvilkaarlig udstykkes i relative Skikkelser,
fortæres den endeliggjørende Relativitet af det theosophisk Absolute.
Men naar Tænkningen affører det Absolute alle relative
Bestemmelser, forsvinder jo Formen: Formen for det
Absolute, den absolute Form, er formløs.

Det Theosophiske er i denne Henseende nærmere beslægtet
med det Mystiske end med det Mythiske. Vel har det Mythiske
ogsaa en Strid mellem Formen og det Formløse; men
det er Phantasien, ikke Tanken, hvem Striden nærmest angaaer.
For Theosophien er det udtrykkelig ikke Phantasien, men Reflexionen,
som væsentlig er interesseret i at forme det Formløse.
Det Formløse har i Tankens Absolute ikke umiddelbar,
elementær Tilværen, saaledes som, hvad Mythylogien angaaer, % PRINTER: printed "Mythylogien", sense "Mythologien"
i Chaos og de ældste Guder. Det Formløse er for Theosophen
tvertimod reflecteret i den Negation, der har fjernet
det Endelige. Og dog kan den theosophiske Bevidsthed ikke
mættes med tom Uendelighed; Theosophiens Blik er fæstet
paa det i sig Absolute, paa Gud som Gud, altsaa paa en
affirmativ Bestemthed.

Den theosophiske Reflexion er naiv; Naiveteten kjendes
paa det Negatives og Affirmatives urolige Eenhed: denne
% --- p. 87 ---
\opage{87}Formens Reflexion i Formløshed minder om Mystiken\footnote{Mystiken gjør Ordet selv negativt. „Gott ist ein lauter Nichts,
Ihn rührt kein Nun noch Hier, Je mehr Du nach ihn greifst, Je
mehr entwird er Dir.“ Angelus Silesius.}. % note *)
Men Theosophien er dog deri forskjellig fra den rene Mystik,
at den lægger Vægt paa Bestemtheden og vil en absolut Form
for det Absolute. Tankens Absolute er Ideen; Sprogets Absolute
er Ordet; Opgaven er at finde et umiddelbart Udtryk
for den absolute Idealitet, eller --- i naiv Reflexion --- at
finde det Navn, hvormed Gud nævner sig selv, det Ord, hvori
Gud aabenbarer sig selv, det Sprog, som Guddommen taler.

Men er det nu ogsaa vist, at Gud selv taler et Sprog?
Er Guddommen ikke det uendelige Dyb (βυθὸς)? Maa ikke
Dybet grunde i evig Taushed (σιγὴ)? Hertil maa Theosophen
svare: „En Guddom uden Sprog, en Gud uden Ord er en
drømmende Guddom, en tankeløs Gud; thi Sproget er Tankens
Aabenbarelse, al Tænken er væsentlig en Talen“. --- Men
er det et Menneske muligt at danne sig noget Begreb om
Guddommens Sprog, at tænke sig den guddommelige Talen?
See, det menneskelige Sprog har i Lyden, i Luftens Zittren,
et sandseligt Element, et umiddelbart Stof, som Tanken
naturligt kan forme til Ord og Sætninger. Men nu er der
dog vel Ingen, der kan forestille sig, at Guddommens Mund
skulde sætte enten Luften eller Vandet eller Ætheren i Svingninger,
for i slige Svingninger at have et sandseligt Element,
et umiddelbart Stof, hvori den guddommelige Tanke kunde
articulere sin første Aabenbarelse? Hertil maa Theosophen
svare: „Det guddommelige Sprogelement har sin Umiddelbarhed
i den ideelle Magt, men denne Umiddelbarhed er ikke
udvortes sandselig. Hvorvidt det i Guddommens Talen oprindelige
Sprogelement nærmest bør sammenlignes med noget
Sandseligt, med Lysphænomener eller med Lydphænomener eller
med begge Dele, er ikke afgjørende for Sprogets Ontologie;
% --- p. 88 ---
\opage{88}men dette er afgjørende, at de sprogformende Functioner maae
betinge de tankeformende --- ligesaa oprindeligt i den ontologiske
Subjectivitet som i den psychologiske; dette er afgjørende, at
Sprogelementet i dets guddommelige Umiddelbarhed maa
bøie sig for Tanken selv til Ord og Sætninger og Domme“.
Men hvorledes er det muligt, at det guddommelige Sprog kan,
i Lighed med det menneskelige, bestaae i et System af Sætninger
og Domme, naar Guddommen deri skal udtale sit
Væsens Indre, en oprindelig Eenhed af Viden og Magt?
Er Ordet hos Gud og selv Gud fra Begyndelsen af, saa
maa dette Ord i guddommelig Eenhed af Viden og Magt jo
være absolut substantielt; men Menneske-Ordet, det afmægtige
Ord, er uden Substantialitet. Hertil maa Theosophen svare:
„Guddomsordet og det guddommelige Sprog ere i Ideen det
Samme og dog forskjellige. I Sproget udtales det Mangfoldige,
det Relative, de mange Syner, Forestillinger, Reflexioner
og Domme; fra denne Relativitetens og Mangfoldighedens
Side har Menneskesproget i formel Henseende unegtelig
Lighed med det guddommelige. Men hvad der i Sproget er
et Relativt, er i Ordet et Absolut. Guddomsordet er det
Absolute selv. Her sees den dybe Kløft mellem Menneske-Ordet
og den guddommelige Logos. I Menneske-Ordet, de
menneskelige Tankers ved Sproget udviklede Indholdsudtryk,
kan Ideen jo vistnok speile sig; men Speilingen er kun formel:
det i Menneskesproget udviklede Ord kan ikke blive een
Substans med Ideen. Anderledes med det almægtige Guddomsord:
det udvikler sit Indhold i Viisdommens Sprog,
det udfolder sin Rigdom i Systemer af Systemer af guddommelige
Domme; men medens det under Udviklingen bestemmer
alle Domme, bevæger alle Sætninger, gjennemtrænger,
opfylder og fuldender den uendelige Totalitet af Begreber og
Formbestemmelser, er det selv alene Totalitetens absolute og
dog umiddelbare Formeenhed. Med denne absolute, men
umiddelbare Formeenhed er Ordet Ideens første Aabenbaring,
% --- p. 89 ---
\opage{89}er i absolut Substantialitet Eet med Ideen, Ideen selv: Ordet
er Gud.“

I ovenstaaende Replikker have vi ladet Theosophen udtale
sin ledende Tanke efter Meningen, for Forstaaelsens Skyld,
men ikke i sit eget Sprog, altsaa ikke i theosophisk Form.
Til Belysning af den theosophiske Form kaste vi et Blik paa
Philo.

1) \emph{Ordet og Gud i tvetydig Eenhed}. Den philoniske
Hovedkategorie er, som bekjendt, Kategorien Logos; men
om Philo iøvrigt opfatter Logos som identisk med Gud eller
som et fra Gud forskjelligt Væsen, om dette Væsen skal antages
at være personligt eller upersonligt, derom have Fortolkerne
hidtil ikke kunnet blive enige\footnote{Jvfr.\ G.\ Steenberg. Om Synspunktet for Opfattelsen af Philos
Gudserkjendelse. Kjøbenhavn 1849. S.\ 84--125.}. % note *)
Den theosophiske
Reflexion vælger i al Naivetet forskjellige Udgangspunkter,
uden at ændse de Modsigelser, som fremkomme, naar
ueensartede Synsmaader skulle gaae op i hinanden. Paa
den ene Side fremstilles da Logos som et med Gud identisk
Væsen, som Guds Tænkning (λογισμὸς τοῦ θεοῦ), som Guds
Ord (ῥῆμα), Guds Eed (ὅρκος), Guds Lov (νόμος), Guds
Pagt (διαθήκη). Men ligesaa umiddelbart opstilles Logos
paa den anden Side som et Væsen for sig, som det Baand
(δεσμὸς), der sammenbinder alle Ting (περιέχων τὰ πάντα),
som, „for at hindre Alt fra at skeie ud, tænkes udspændt fra
den ene Ende af Verden til den anden“; Logos er „den guddommelige
Viisdom, der gjennemstrømmer Verden; Verden % PRINTER: quotation opened at „den guddommelige is not closed
er denne Viisdoms Klædebaand; Viisdommen selv er ὁ τῆς
φύσεως ὀρθὸς λόγος. Uagtet Logos i begge disse Opfattelser
er upersonlig, saa indstille sig dog igjen andre Synspunkter,
fra hvilke Logos udtrykkelig sees som personligt Væsen,
som en Gud selv underordnet Gud (δεύτερος θεὸς, πρωτόγονος
υἱὸς), o.\ s.\ v.

% --- p. 90 ---
\opage{90}Den theosophiske Reflexion kjendes overalt derpaa, at
Tanken, det Grundbevægende, skyder sig ind under abstracte
Anskuelser, og lader Totalitetsmomenterne blive staaende i
umiddelbar Selvstændighed. Den theosophiske Reflexion er
i sin Naivetet ukritisk; svinder nu Naiveteten, medens den
ukritiske Brug af intellectuelle Anskuelser forceres op til en
høiere, paa Troen hvilende Aabenbaringsvidenskab, saa have
vi den dogmatiske Speculation. I den dogmatiske Speculation
er da ogsaa Logos som „Gud Faders Hjerte tillige det evige
Verdenshjerte, hvorigjennem det guddommelige Liv strømmer
ind i Skabningen. Han er Grund og Kilde til al Fornuft
i Skabningen“, o.\ s.\ v.\ o.\ s.\ v.

2) \emph{Ordet og Gud i tvetydig Modsætning}. Den
theosophiske Reflexion, der sætter Eenheden af Gud og Logos
umiddelbart, sætter tillige Modsætningen umiddelbart, og ophæver
igjen den satte Modsætning ligesaa umiddelbart; herfra
den Vrimmel af Modsigelser, som Kritiken saa let opdager.
Skulle disse Modsigelser løses, da maa den ubestemte Eenhed
udvikles til Modsætning, for at Modsætningen, Totalitetsdelingen,
den ideeopfyldte Dom, igjen kan gaae op i udviklet Eenhed.
Det er, vel at mærke, Begrebsudviklinger og ikke nye Aabenbaringer,
hvorpaa det i denne Henseende væsentlig kommer an.

Den speculative Dogmatik forsikkrer, at den paa Grund
af sin nye Aabenbaring finder sig høit ophøiet over den gamle
„alexandrinske Religionsphilosophie“. „Man har ofte“, hedder
det, „sammenstillet Aabenbaringens Logoslære med den alexandrinske
Religionsphilosophie, ja der er ofte udtalt den Paastand,
at Johannes skulde have laant sin Logoslære fra Jøden
Philo. Men her finder den bestemteste Modsætning Sted.
For Philo forvandler Historien sig til en Skinverden, de
historiske Kjendsgjerninger til en Allegorie for almindelige
Tanker og Ideer. For Philo er der et svælgende Dyb befæstet
mellem „Ord“ og „Kjød“ (λόγος og σὰρξ); thi Sandseligheden
og Naturen ere ham kun et Aandens Fængsel, og
% NOTE: the quotation opened at „sammenstillet (p. 90) closes on p. 91.
% --- p. 91 ---
\opage{91}kun ved en mystisk Bortrykkelse fra Virkeligheden, kun ved
Sandselighedens Dødelse kan Tankeren forenes med den guddommelige Fornuft.“

Dersom det nu i Videnskaben blot kom an paa at bringe
„Kjød“ ind i Ordet og „Ord“ ind i Kjødet; dersom Knuden
var løst, naar Speculationen med Pathos blot vedblev at
gjentage sit λόγος σὰρξ ἐγένετο: saa vilde Dogmatiken
unegtelig feire en stor Triumph over den philoniske „Religionsphilosophie“.
Men dersom her virkelig kræves Begrebsudvikling,
vil det være let at vise, at Dogmatiken, trods de
mange triumpherende Forsikkringer, ikke er kommet et Haarsbred
videre end Philo.

For at „Sandseligheden“ og „Naturen“ skal kunne opgaae
som integrerende Momenter i Logos, maa Logosbegrebet
oprindelig være bestemt saaledes, at det Reelle er med i det
Ideelle, at, med andre Ord, Umiddelbarheden fra Begyndelsen
af slaaer igjennem. Dette er nu ikke Tilfældet hos Philo.
Misligheden hos Philo er just den, at Logos eensidig opfattes
som Tanke (λόγος νοητὸς), ikke tillige som Ordet, det Ord,
der i Umiddelbarhedens Identitet kun er Eet med Tanken,
forsaavidt det er Tankens Andet, Tanken modsat. Men den
samme Mislighed gjentager sig i den speculative Dogmatik;
her er da heller ingen Antydning af Sprogelementets ideelle Realitet,
eller Spor til Ordets primitive Umiddelbarhed; Umiddelbarheden
kommer først, naar vi komme til „Kjødet“. Misligheden
hos Philo er endvidere den, at Idealverdenen, det oprindelige
Logosindhold (κόσμος νοητὸς), mangler væsentlig
Articulation, fordi her mangler al Bevidsthed om den for en
saadan Articulation gjældende Methode\footnote{Jvfr. Grundideernes Logik I. S. 178--220.}. % note *)
Men den samme,
den selvsamme Mangel gjentager sig i den speculative Dogmatik.
Ogsaa her speide vi midt imellem „legende Muligheder“,
„plastiske Forbilleder“ og „magiske Virkeligheder“
% --- p. 92 ---
\opage{92}forgjæves efter et articulerende Princip; thi en Henviisning
til „Aandens frie, kunstneriske Virksomhed“ er naturligviis
ørkesløs, saalænge man ikke ved Noget, der ligner et Begreb,
seer sig istand til at angive, hvad her da egentlig menes med
Aandens Plastik og magiske Kunster. Hos Philo bliver Viisdommen
(σοφία), Logos og Magterne (δυνάμεις) adskilte og
blandede med al mulig Vilkaarlighed: snart er Sophia den
Kilde, hvoraf Logos udspringer, snart er Logos igjen Kilden
til Sophia. Det samme, det selvsamme Forestillingernes Histop
og Herned gjenfinde vi i den speculative Dogmatik, naar „de
magiske Syner, der lege for Guds Ansigt i den indre Selvaabenbaring,
maae tænkes at bestemme sig for Gud som
Raadslutninger til Skabelsen“ o. s. v. Philo mangler Bevidsthed
om, at det Logiske i Idealsystemet (κόσμος νοητὸς)
kun lader sig udvikle i Reflexionen af Realsystemets antilogiske
Momenter\footnote{Grundideernes Logik I. S. 224--229.}. % note *)
Det Svage hos Philo er desaarsag den Uklarhed,
hvormed han hypostaserer Logos snart som guddommelig
Person (δεύτερος θεὸς), snart igjen som Verdensvæsen, det
kosmiske Princip, (ὁ τῆς φύσεως ὀρθὸς λόγος), Verdens
iboende Fornuft, der udsaaer Livets Spirer (λόγος σπερματικὸς).
Men den samme, den selvsamme Uklarhed hersker i
den speculative Dogmatik, og af samme, selvsamme Grund.
Uden theosophisk Naivetet veed den dogmatiske Speculation paa % BL copy prints "Speculation" (the Harvard copy was read as "Specnlation")
en egen naiv Maade at operere baade med „Faderens Logos“
og med „Verdenslogos“, og med „Verdensspirerne“ (λόγος
σπερματικὸς) og med „den anden Gud“ (δεύτερος θεὸς)
og med „de oprindende Ideer“ og med Andet deslige.

3) \emph{Ordet og Gud i udviklet Modsætningseenhed.}
Logik er ikke Theosophie; men i de theosophiske
Functioner arbeider der noget Rationelt, som Logiken maa
stræbe at klare. Naar Dommens Form skal rumme Ideen,
det absolute Indhold, viser sig en Grunddeling, hvorved den
% --- p. 93 ---
\opage{93}ene og samme Totaliternes Totalitet adskiller sig i to oprindelig % PRINTER: printed "Totaliternes", sense "Totaliteternes"
modsatte Former, Selvhedens eller Subjectivitetens
og Andethedens eller Objectivitetens Form. Med Selvhedens
Form er det absolute Indhold reflecteret i Eenhed,
Væsenhed, Inderlighed og Grund; i denne Inderlighed er
Dybet den tause Tanke, der bryder Tausheden. Som den
uendelige Subjectivitet, der i den dybe Taushed hæver Tausheden
og udtaler sit Indhold, er Guddommen ikke Ordet selv,
men \emph{Ordets Ophav}. Med Andethedens Form er det absolute
Indhold articuleret i Mangfoldighed, i Forskjellighed,
altsaa en i sine Led udpræget Totalitet; Talens, det vil sige
Tankeaabenbaringens, logiske Element, en primitiv Umiddelbarhed,
hvori Dommen adskiller det Forenede og forener det
Adskilte, er som totaliseret Eenhed \emph{Ordet}. Ordet i sin Umiddelbarhed
er ikke Guddommen, men dens Aabenbaring; ikke
Subjectiviteten selv, men dens objectiverede Indhold. Den
absolute Totalitet er ved de to modsatte Former blevet til to
modsatte Totaliteter: Guddommen og Ordet.

Men naar Dommens Form skal rumme Ideen, det absolute
Indhold, maa den Grunddeling, hvorved den subjective,
i Selvheden hvilende Totalitet sættes imod den objective i
Andethedens umiddelbare-middelbare Udfoldning, nødvendig
have et Bindeled, en Midte, en Modsætningseenhed, hvori
Extremerne mødes. I Kraft af den totaliserede Modsætningseenhed
er Ordet i sin Modsætning til Gud Eet med Gud:
Ordet er Gud. I sin Eenhed med Gud er Ordet den objectiverede
Gud (ὁ δεύτερος θεὸς).

Som det sig objectiverende Selv er Gud Ordets Ophav;
som det objectiverede Selv er Gud identisk med Ordet. Men
som den i sig selv objectiverede Gud er Ordet ikke blot Eet
med Gud, men formedelst Objectiveringen tillige modsat Gud:
i sin objectiverede Modsætning til Gud er Ordet ikke Gud,
men en Verden i Gud, den ideale Verden (κόσμος νοητὸς).

% --- p. 94 ---
\opage{94}Men hvordan gaaer det nu egentlig til, at det samme
absolute Indhold, denne samme i Guddommen objectiverede
Totalitet kan blive baade til Gud selv (πρωτόγονος υἱὸς, δεύτερος
θεὸς) og til en i Gud fra Gud forskjellig Verden?
Det ligger i Dommens Idee! Her er i den objectiverede Totalitet
en afgjørende Spænding mellem Selvhed og Andethed,
Reflexion og Umiddelbarhed, Indholdets Eenhed og Eenhedens
Indhold. Falder Eftertrykket paa Selvheden, og den i Selvheden
reflecterede Indholdseenhed, saa have vi i den objectiverede
Totalitet den objectiverede Guddom, Gud selv i egen
Objectivering. Falder Eftertrykket derimod paa Andetheden
og den Udvikling, hvormed Eenhedens Indhold i Delenes
umiddelbare-middelbare Forholdsbestemmelser er totaliseret, saa
have vi den objectiverede Totalitet som en i Mangfoldighed
og Forskjellighed articuleret Sammenhængseenhed, altsaa en
Verden i ideel Tilværen, en Verden i Gud, forskjellig fra Gud,
en Idealverden.

Men Idealverdenen, den ideelle Verdenstotalitet, vilde
umulig kunne danne en væsentlig Modsætning til den deri
objectiverede Guddom, dersom denne Idealverden ikke oprindelig
var i totaliseret Modsætning til sin Realverden. Den
ideale og den reale Verden have i Realiteten samme Indhold,
men i Idealiteten ved den til Dommens Idee, til
Grunddelingen, svarende Formforskjel have de hver sit Indhold:
i sin Idealverden er Ordet et Videns Ord, men i sin
Realverden er det et Magtens Ord. Kun ved Modsætningseenhed
med Magtens Ord faaer Vidensordet, og dermed den
ideelle Verden, ideel Selvstændighed.

Det er i Læren om Subjectiviteten\footnote{Grundideernes Logik I. S. 210 flgd.} % note *)
fra flere Sider
udviklet, at „Tankehandlingerne“ i Idealsystemet umulig kunne
falde sammen med „Naturhandlingerne“ i Realsystemet. Men
under Problemets første Behandling kunde det endnu ikke sees,
% --- p. 95 ---
\opage{95}paa hvad Maade Ideebevægelsen i den ontologiske Subjectivitet
kunde frigjøre sig fra Realbevægelsen i de virkelige Objecters
Yderverden; nu derimod er Problemet løst. Er Ordet i Articulation
af Sætninger og Domme et guddommeligt Sprog;
have Guddomstankerne indre, ideel Tilværelse i det ideale
Sprog, hvori de have deres første Aabenbarelse; er det Ord,
hvori den uendelige Rigdom af udtalte Tanker har totaliseret
Eenhed, den guddommelige Videns Ord: saa følger jo heraf,
at den ideale Verden maa have indre Selvstændighed i en guddommelig
Tænken og Talen.

Den ideale og den reale Verden ere imidlertid ikke blot
i Modsætning, men --- hvad da ogsaa ligger i Dommens Idee ---
tillige i Modsætningseenhed. Fordrer den til Ideen svarende
Dom en Modsætning af Ordet i Viden og Ordet i Magt,
da maa Dommens Idee nødvendigt ogsaa fordre en totaliseret
Eenhed af Videns og Magtens Ord.

\begin{center}{\bfseries β)\quad Dom og Gjenstand: Functioner af Magtens Ord.}\end{center}

Det er i „Subjectivitetens Logik“\footnote{Grundideernes Logik I. S. 145--166; 201--204; 255--263.} % note *)
fra flere Sider
blevet paaviist, hvorledes det samme Indhold kan i Videns og
Magtens, i Begrebets logiske og Existensens antilogiske Form
forvandle sig og blive et dobbelt Indhold. Denne Fordobling
blev da gjennemreflecteret ved Ordet, uden at der endnu kunde
reflecteres paa Ordet. I nærværende Sammenhæng, hvor
der udtrykkelig reflecteres paa Ordet som Ord, vender Opgaven
tilbage med forhøiet Eftertryk; Opgaven er at vise,
hvorledes Begrebsindholdet i Videns Ord og Existensindholdet
i Magtens Ord forudsætte og bestemme hinanden i det Uendelige.
Idet Begrebet faaer ideel Selvstændighed i Vidensordet,
medens Existensen optræder som en ved Magtens Ord bestemt
Gjenstand for sig, falde de ved første Øiekast saaledes ud fra
% --- p. 96 ---
\opage{96}hinanden, at det synes, som vare de aldeles ligegyldige mod
hinanden. Det seer jo paa den ene Side ud, som om Begrebet,
efter at det i Ordet har faaet ideel Tilværen, kunde
dannes, udvikles og fastholdes enten for sig eller i Sammenhæng
med andre Begreber, uden Hensyn til, om der i Existensen
fandtes nogen til dette Begreb svarende Gjenstand, eller
ikke. Paa den anden Side tager det sig ud, som vare de
existerende, i Magtordet bestaaende Gjenstande aldeles ligegyldige,
vel ikke imod at blive objectiverede, men dog mod
den særegne Begrebsudvikling, hvortil deres Objectivering kan
give Subjectiviteten Anledning. Hvorledes bestemme vi da
Forholdet mellem Begreb og Ting, Dom og Gjenstand?

Den formelle Logik behandler Opgaven paa sin Maade,
idet den til Orientering paa Tænkningens Omraade gjør Adskillelse
mellem vort Begreb om Gjenstanden og Gjenstandens
Begreb, men derimod ikke imellem vor Dom om Gjenstanden
og Gjenstandens Dom; den skjelner imellem Gjenstandens
Begreb og Gjenstanden selv, men ikke mellem Gjenstandens
Dom og Gjenstanden; altsaa mellem Begrebet i subjectiv og
Begrebet i objectiv Betydning, men ikke imellem Dommen i
subjectiv og Dommen i objectiv Betydning. Vi skulle nu see,
hvad Værd der kan tillægges disse Distinctioner.

Forskjellen imellem Begrebet i subjectiv og Begrebet i
objectiv Betydning kan, naar her sees paa Begrebet selv,
umulig være nogen logisk Forskjel. Skulde her være nogen
logisk Forskjel, da maatte den jo være betinget af, at vort
subjective Begreb om Tingen indeholdt enten flere eller færre
eller dog andre Bestemmelser end Tingens objective Begreb;
men saaledes blev det \emph{subjective} Begreb da det Samme
som et \emph{falsk} Begreb, hvilket ikke kan være Meningen. Men
dersom det subjective Begreb har samme logiske Indhold som
det objective, saa falder jo Forskjellen bort: det subjective Begreb
er i sin Sandhed objectivt; vort sande Begreb om Tingen
er Tingens Begreb.

% --- p. 97 ---
\opage{97}Ligesaa uholdbar er den opstillede Distinction mellem
Gjenstandens Begreb og Gjenstanden selv. Gjenstandens Begreb,
hedder det, er det „Intellectuelle“ og „til Grund Liggende“
i Gjenstanden, det, der udgjør Gjenstandens „Natur og
Væsen.“ Men hvorledes er det muligt at skille en Gjenstands
Natur og Væsen fra Gjenstanden selv? Hvad bliver der tilbage
af en Gjenstand, naar dens Natur og Væsen tænkes
borte? Dersom det i Gjenstanden, der ikke udgjør dens Natur
og Væsen, ikke hører til dens Begreb, saa maa det jo være
noget saa Tilfældigt, at det egentlig slet ikke vedkommer Gjenstanden;
maa her derimod foruden det, der udgjør Gjenstandens
Natur og Væsen, nødvendigt være endnu noget Andet, Noget,
som dog, uden at vedkomme dens „Natur og Væsen,“ væsentlig
og nødvendig hører med til dens Existens, saa ligger det
jo i Gjenstandens Natur og Væsen at have Noget, der ikke
vedkommer dens „Natur og Væsen:“ hvad mon det da være
for Noget?

Det hedder, at Gjenstanden selv er „det reelle Udvortes“,
det sandselige Ydre, hvorimod Begrebet er „det Indre,
Ideelle og til Grund Liggende“. Men naar „det reelle Udvortes“
udgjør Gjenstanden selv, den virkelige Gjenstand, saa
maa „det Indre, Ideelle og til Grund Liggende“ jo i
Virkeligheden falde udenfor Gjenstanden\footnote{% note *) -- runs over from p. 97 to p. 98
„A. Her er i een og samme Existens altsaa to Existerende, nemlig
den existerende Gjenstand og det i Gjenstanden existerende Begreb?
C. Nei, saadan kan det ikke tages; det er ikke Gjenstanden selv,
men dens Materiale, der er Begrebet modsat. A. Vi faae vel
altsaa tre Existerende: det existerende Materiale, det existerende Begreb
og den Gjenstandens Existens, hvori begge existere? C. Nei,
saadan kan man heller ikke tage det. A. Hvordan skal man da
tage det? C. Materialet for sig er ikke et særligt Existerende, saalidt
som Begrebet, men Materialet og Begrebet ere to Sider, nemlig
to Existenssider af det samme existerende Ene, Gjenstanden. A.
Det existerende Ene er altsaa i Modsætning til og ligeoverfor de
% note breaks here: p. 97 | p. 98
to Existenssider det i og for sig Existerende? C. Nei, saadan kan
man heller ikke tage det, thi dersom vi lade Abstractionen borttage
de to Existenssider, vil den existerende Eenhed jo selv blive et Ab%
stractum, en Existensside, der ikke existerer. A. Men hvordan skal
man da tage det? C. Saaledes, at den existerende Eenhed sees at
existere i de to Existenssider, som da just ogsaa derved komme til
at existere i uadskillelig Eenhed \dots. A. Men det i Gjenstanden
Materielle er dog vel et Sandseligt, der som sandseligt ikke kan
tænkes, og det i Gjenstanden Intellectuelle et Usandseligt, et Tænke%
ligt, der ikke kan sandses? C. Fuldkommen rigtigt. A. I den
existerende Eenhed er det Utænkelige altsaa dog et Tænkeligt, det
Usandselige et Sandseligt. C. Det forstaaer jeg ikke!“ A. Men
derimod forstaaer Du jo nok, at det Ideelle, det Indre og det
til Grund Liggende, der udgjør Gjenstandens „Natur og Væsen“,
ikke er det Reelle, det Ydre, det Existerende, der netop udgjør
Gjenstandens, den virkelige Gjenstands Natur og Væsen! Jvfr.
Forelæsninger over phil. Propæd. Universitetsaar 1860--61.
S. 26--27.}. % note: the quotation opened at „A. closes after "ikke!“" though A. speaks on; as printed
Heller ikke kan
% --- p. 98 ---
\opage{98}det hjælpe, at man vil gjøre det Indvortes til „det Consti%
tutive,“ det Udvortes til „det Consecutive“; thi i Gjenstanden
som Gjenstand, og her er kun Tale om den virkelige Gjen%
stand, er det Indvortes Gjenstand ligesaa vel som det Ud%
vortes: her er ikke en Gjenstand, der gaaer uden om et Be%
greb, ikke et Begreb, der sidder indeni en Gjenstand; her er
i Gjenstanden forinden og foruden kun Gjenstand.

Hvor uklar Adskillelsen mellem Gjenstanden og det i
Gjenstanden existerende Begreb er, sees allerbedst, naar For%
holdet mellem Begreb og Dom nærmere skal bestemmes.
Man kan --- forsikkres der i den formelle Logik --- ikke bruge
Ordet Dom i objectiv Betydning, saaledes som Ordet Begreb:
„man kan ikke sige, at der paa samme Maade, som der i en
Gjenstand kan træde os et Begreb imøde, ogsaa kan træde
os en Dom imøde.“ Og hvorfor ikke? Fordi man nu een
Gang har sat sig i Hovedet, at der indeni den objective Gjen%
stand sidder et objectivt værende Begreb. Men saasnart man
% --- p. 99 ---
\opage{99}ved at gjennemføre Reflexionen har faaet Bugt med denne Fixidee,
og indseet, at det Begreb, der „i en Gjenstand kan træde
os imøde“, netop er vort eget til Gjenstanden svarende Begreb,
vil man ogsaa strax indsee, at Dommen, den sande, til Gjenstanden
svarende Dom, har samme Objectivitet som det sande
Begreb. At der i Gjenstanden ikke kan træde os en Dom
imøde, er jo vist nok, forsaavidt der ikke kan gives en i Gjenstanden
objectivt værende Dom om Gjenstanden; men det
Samme gjælder jo om Begrebet. Kan Gjenstanden ikke
møde med en Dom om sig selv, med en sig objectivt dømmende
Dom, da kan den heller ikke møde med et Begreb om
sig selv, med et sig objectivt begribende Begreb. Falder
Dommen om Gjenstanden i den dømmende Subjectivitet, da
falder Begrebet om Gjenstanden, Gjenstandens objective Begreb,
tilvisse ogsaa i den begribende Subjectivitet; taber Begrebet
ikke sin objective Gyldighed, fordi det absolut maa begribes
af Subjectiviteten, saa taber Dommen heller ikke sin Objectivitet,
fordi det er en Subjectivitet, der dømmer.

Forholdet mellem Begreb og Existens, mellem Dom og
Gjenstand, er i Principet et Grundforhold mellem Ord og
Ord, mellem Videns og Magtens Ord. Indholdet er det
samme; men Formen er forskjellig: i Vidensordet er Indholdet
Begreb, og Begrebet udtales i Dommen; i Magtens
Ord er Indholdet Existens, og det Existerende Dommens Gjenstand.
Seet fra Begrebets Side er det ikke destomindre Formen,
der bliver den samme, medens Forskjellen falder i Indholdet;
og det Samme gjælder, naar Sagen sees fra Gjenstandens
Synspunkt. Vel maa der i Begrebet om Gjenstanden,
om man vil, i Gjenstandens Begreb, gjøres Forskjel
imellem det, der i Begribningen tilhører Begrebet, og det, der
paa Grund af Existensen tilhører Gjenstanden; men saalænge
Adskillelsen af Begrebs- og Gjenstandsbestemmelser falder ind
under Begrebet, maa Begrebet om Gjenstandsbestemmelserne
paany opløse sig i Begrebsbestemmelser. Formen for begge
% --- p. 100 ---
\opage{100}Arter af Bestemmelser bliver altsaa den samme; Forskjellen
falder paa Indholdet: i Begrebsbestemmelsernes Totalitet er
Indholdet logisk; i Begrebet af Gjenstandsbestemmelsernes To%
talitet er Indholdet antilogisk.

„Men det antilogiske Indhold maa jo dog have en til%
svarende antilogisk Form; Gjenstandsbestemmelserne ere an%
tilogiske saavel i Form som i Indhold;“ saaledes lyder Dommen
om Gjenstanden fra Gjenstandens Synspunkt. Da nu Dommen
om Gjenstanden ikke kan adskilles fra Dommen om Dommen,
uden at Adskillelsen bliver ørkesløs, saa lyder Dommen til%
lige paa det Samme fra Subjectivitetens Synspunkt. De i
Dommen udtrykte Gjenstandsbestemmelser ville da for den
dømmende Bevidsthed stille sig skarpt mod det Almene, ville i
Gjenstandens Navn individualisere sig med antilogisk Punk%
tualitet, og Formen altsaa være antilogisk. Den mod det
Logiske vendte Form staaer da nu for Bevidstheden i en logisk
Speiling af et antilogisk og dog logisk Indhold. Indholds%
fordoblingen viser sig igjen: i sig selv, som Gjenstandsind%
hold, er Indholdet antilogisk ligesom Formen, men i den \emph{logiske} Speiling af den \emph{antilogiske} Form er Indholdet logisk.

At bestemme Forholdet mellem Dom og Gjenstand er
ingenlunde det Samme som at paavise Identiteten af Dom og
Indhold. Den sidste Opgave er i den Forstand logisk, at
den kan indskrænke sig til en Udvikling af rene Begrebsbestemmelser;
den første Opgave er derimod logisk i den Forstand,
at det Antilogiske selv fordrer en logisk Forklaring.
Det er de antilogiske Formers Logik, et af de vanskeligste
Partier i den hele Philosophie, hvorpaa det i denne Sammenhæng
kommer an. Dersom vi kun havde med logiske Former
og tilsvarende Indholdsbestemmelser at gjøre, vilde Functionerne
alle uden Undtagelse være Functioner af Videns
Ord; det Høieste vilde være naaet, naar Dommens Form
var saa vidt udviklet, at Grunddelingen svarede til Ideen.
Den antilogiske Braad, der, saa at sige, føleligt afficerer Ind%
% --- p. 101 ---
\opage{101}holdet og bevirker, at Dommens Gjenstand stiller sig skarpt
imod Dommen, uden at Identiteten af Dom og Indhold
tør opgives, vækker derimod overalt en Viden af Modsætning
mellem Magt og Viden, og gjør Magtens Ord til Functionernes
Grundlag. Functionerne ere logiske; Magtens Ord
er antilogisk; den uendelige Vexelbestemmelse af det Logiske og
det Antilogiske, som herved fremkommer, kræver en Belysning
af Dom og Gjenstand fra Videns, fra Magtens og, ved
Modsætningseenhed af Viden og Magt, fra Ideens Synspunkt.
Vi betegne disse Synspunkter i det Følgende ved:
\emph{Viden af Dom og Gjenstand}, \emph{Magten i Dom og
Gjenstand}, \emph{Ideen i Dom og Gjenstand}.

\textit{αα}) \emph{Viden af Dom og Gjenstand: theoretiske
Functioner.} I psychologisk Dualisme staaer Omverdenen
med sine Gjenstande udvortes imod den betragtende, tænkende,
dømmende Bevidsthed. At Gjenstanden i sin umiddelbare
Selvstændighed falder udenfor Dommen; at Dommen kun
angaaer Bevidstheden om Gjenstanden, medens Gjenstanden
selv er ligegyldig mod Bevidstheden: denne eensidige Betragtningsmaade
er i den Grad rodfæstet i den saakaldte sunde
Menneskeforstand, at Philosophien netop af den Grund fra
alle Sider maa finde sig opfordret til at modarbeide dens
Eensidighed. Vistnok falder i umiddelbar Forstand den virkelige
Gjenstand udenfor den dømmende Bevidsthed; men hvad
der væsentlig gjør Skilsmisse mellem Gjenstanden og Bevidstheden,
er jo netop Dommen om Bevidsthed og Gjenstand.
At det med Hensyn paa Gjenstandens Selvstændighed ligesaa
væsentlig kommer an paa Dommen som paa Gjenstanden, er
Noget, der undgaaer den umiddelbare Sands. For den
umiddelbare Sands tager Sagen sig ud, som om Gjenstanden
ved sig selv udenfra kom Bevidstheden imøde, medens det dog
ligesaa oprindelig er Bevidstheden, der indenfra kommer Gjenstanden
imøde; „den sunde Forstand“ overseer, at Gjenstanden
kun er udenfor Bevidstheden, fordi den er dømt og ved Dommen
% --- p. 102 ---
\opage{102}sat til at staae imod Bevidstheden. Tag Dommen bort, den
Instinctets Dom, hvorpaa al Skilsmisse mellem det Subjec%
tive og det Objective, det Udvortes og det Indvortes væsent%
lig beroer, og Gjenstanden vil øieblikkelig falde saaledes sammen
med Bevidstheden, at her hverken er Gjenstand eller Be%
vidsthed.

Sætter den dømmende Subjectivitet Skilsmisse mellem
sig selv og Gjenstanden i Dommen, da sees ogsaa heraf, at
Adskillelsen af Dom og Gjenstand ligeledes maa beroe paa
Dommen. Men at Adskillelsen af Dom og Gjenstand beroer
paa Dommen, forudsætter igjen, at Identiteten af Dom
og Gjenstand beroer paa Dommen. At her i Dommen er
en Viden af Gjenstanden i dens udtrykkelige Modsætning til
Dommen, grunder sig da paa, at her i Dommen er en
Viden af Gjenstandens Identitet med Dommen: den dobbelte
Viden er i Dommen selv een og samme Viden. I Reflexion
af denne dobbeltsidige, Dommen og Gjenstanden oversvævende,
Viden ere Functionerne \emph{theoretiske}.

Det er et Særkjende for al Theorie, at Videns Gjenstand
bestemmer Viden, og ikke omvendt; men her er nu en
dobbelt Viden, en Viden af Dommen i dens Forhold til
Gjenstanden, og en Viden af Gjenstanden i dens Forhold
til Dommen; Viden har, med andre Ord, en dobbelt Gjenstand:
en \emph{ideel}, nemlig Dommen som Dom, og en \emph{reel,}
Gjenstanden som Gjenstand. Theorien vil ikke vilkaarlig forandre
Noget enten ved Dommen eller ved Gjenstanden, men
tage dem i deres Gjensidighed, som de i sig selv ere. Problemet
sees af følgende Spørgsmaal: Er det Dommen, der
væsentlig bestemmer Gjenstanden, eller Gjenstanden, der bestemmer
Dommen, eller er her en Vexelbestemmelse? Beroer
Erkjendelsens Sandhed paa den blotte Lighed imellem Bestemmelserne
i Dommen og Gjenstandsbestemmelserne, eller maa
her foregaae en Forvandling, for at Gjenstandsbestemmelserne
% --- p. 103 ---
\opage{103}kunne blive identiske med de i Dommen indeholdte Vidensbestemmelser?

Gaaer Theorien ud fra den Forudsætning, at Bestemmelserne
i Dom og Gjenstand ere ueensartede, og at der
imellem ueensartede Bestemmelser kun kan gives en Parallelisme
paa Grund af Ligheden, ikke en ved Forvandling
betinget Overgang, saa træder Overeensstemmelsen som blot
og bar Parallelisme i Identitetens Sted, og den psychologiske
Dualisme maa da søge sit Erkjendelsesprincip i en over Bevidstheden
staaende Forudsætning. Altsaa:

1) \emph{Parallelisme af Dom og Gjenstand: supponerende
Functioner.} Den existerende Gjenstand er i sin
udvortes Umiddelbarhed noget Sandseligt, Dommens Indhold
noget Tænkeligt; det Sandselige er ueensartet med det Tænkelige:
det ligger da nær at opfatte Sandsephænomenerne og
Tankebestemmelserne som to Rækker, der vel kunne afgive en
Parallelisme, men ingen Identitet.

Det er saaledes, ifølge Plato, ikke det Sandselige i
Tingene, men de Ideer, hvori Sandsetingene ere blevne deelagtige,
der lade sig tænke. At Dom og Gjenstand kunne i
den sande Tænkning stemme overeens, er betinget af, at de
samme Ideer, der aabenbare sig subjectivt i Tænkningen,
aabenbare sig objectivt i Tingene. Hvoraf veed nu den dømmende
Subjectivitet, at Ideerne aabenbare sig objectivt i
Tingene, eller, hvad der er det Samme, at Tingene ere deelagtige
i de objective Ideer? Antagelsen beroer paa en af
Tanken selv umiddelbart underskudt Forudsætning. Hvor umiddelbar
og objectiv Underskydningen er, sees af Præexistensmythen
med den tilsvarende Lære om Erindringen (ἀνάμνησις).

Det for Suppositionen Charakteristiske er den Naivetet,
hvormed den subjective Reflexion kommer objectivt til sine
egne Antagelser. Først antages, at Tanken alene kan tænke
Ideerne, efterdi Ideerne alene ere Tankevæsener. Ved at fastholde
denne Antagelse vilde Tænkningen aabenbart ikke
% --- p. 104 ---
\opage{104}komme ud over sin egen Subjectivitet; Dommen vilde med
andre Ord kun have et logisk Indhold, ikke en antilogisk Gjenstand.
Da Gjenstanden i sin ydre Selvstændighed imidlertid
paatrænger sig uvilkaarligt, saa bliver Tanken, uden ret at
vide deraf, bevæget til at lægge sit subjective Indhold objectivt
ind i Tingene og saaledes gjøre disse deelagtige i Ideerne.
Ideerne ere nu i Tingene selv objective, ikke subjective, efterdi
de i Tænkningen ere subjective, ikke objective. Den oprindelige
Overeensstemmelse af Subjectivt og Objectivt, hvorpaa
Tanken bygger som paa et objectivt Givet, har den subjectivt
givet sig selv.

De supponerende Functioners underskudte Grundlag træder
endnu stærkere frem hos Cartesius og efter ham hos Malebranche
og Leibnitz. Forudsætningen er en Adskillelse af det
Intellectuelle og det Materielle, bestemt ved en for de relative
Substanser uoverstigelig Kløft mellem Tænkning og Udstrækning.
En ligefrem Vexelvirkning af Forestillingerne i Bevidsthedskredsen
med Tingene i Omverdenen er fra dette Synspunct
aldeles umulig. Saalidt som en Forestilling eller et
Begreb kan blive virkende Aarsag til en Forandring i Tingenes
Verden, ligesaa lidt kan et physisk Indtryk bevirke Dannelsen
enten af en tilsvarende Forestilling eller af et Begreb. At
Forestillingerne i Bevidstheden og Tingene i Omverdenen,
uden gjennemgaaende Vexelvirkning, blot danne Paralleler, er
for Theorien lige afgjørende, enten vi med Malebranche
appellere til en gjennem Anledningsaarsager uophørlig virkende
Almagt, eller med Leibnitz til en \textit{harmonia præstabilita.}\footnote{% note *) -- runs over pp. 104, 105 and 106
„Materielle Gjenstande“, siger Malebranche, „kunne umulig opfattes
ved sig selv, de opfattes kun ved de i Sjælen tilsvarende Ideer.
Peripatetikernes Mening, at de udvortes Gjenstande skulde tilsende
os deres Afbilleder (\textit{des espèces, qui leur resemblent}), og at disse
saa igjennem de ydre Sandser skulde bringes til den fælles Sands
(\textit{au sens commun}) er urimeligt; disse Indtryksbilleder (\textit{espèces
% note breaks here: p. 104 | p. 105
impresses}) maatte da selv være smaa Legemer, og disse Legemer,
hvoraf Rummet vilde være opfyldt, maatte da ifølge Materiens
Uigjennemtrængelighed knuse og støde hverandre, idet nogle gik til
en, andre til en anden Side, og kunne altsaa ikke gjøre Gjenstandene
synlige.“\par
Deres Mening, som antage, at Gjenstandenes Ideer ere skabte
med os, saa at Gud ikke behøver i hvert Øieblik at vække dem i
Sjælen, naar vi have dem behov, gjendriver Malebranche med
følgende Betragtning: „Vi have, selv hvad simple Figurer angaaer,
en uendelig Mængde Ideer: Ellipser, Triangler, Polygoner o. dsl.
kunne i Henseende til Størrelsesforhold imagineres paa utallige
Maader. Det er da høist usandsynligt, at Gud skulde have skabt
en saadan Vrimmel af Ideer i og med den menneskelige Aand.
Men sæt, at vi ogsaa havde et saadant Oplag af medskabte Ideer
(\textit{un magazin de toutes les idées}), saa bliver det dog aldeles uforklarligt,
hvorledes vi komme til at anvende den rette Idee, naar
den tilsvarende Gjenstand viser sig. Vi see f. Ex. Solen. Det
Billede, Solen præger i Hjernen, har ingen Lighed med den Idee,
vi danne os om Solen; heller ikke mærker Sjælen det allermindste
til det Indtryk, Solen gjør enten paa Øiets Grund eller i Hjernen:
det er altsaa umuligt at forklare, hvorledes vi blandt de utallige
Ideer just komme til at vælge Solens Idee, naar vi see Solen“ \dots\
Ved disse og lignende Raisonnements kommer Malebranche saa til
det Resultat, at Aanden kan see, hvad det er, der er i Gud, som
fremstiller „de skabte Ting“ o. s. v. Jvfr. Om Theologiens Naturbegreb
med særligt Hensyn til Malebranche: \textit{De la recherche de
la vérité.} Af R. Nielsen. I Anledning af Reformationsfesten
1855. S. 13--15.\par
Hvad Malebranche henfører til Gud umiddelbart, henfører
Leibnitz til Gud middelbart. Forudsætningen er, at Monadernes
Monade har anlagt hver enkelt Monade saaledes, at den som et
levende Speil (\textit{miroir vivant}) kan afspeile hele Universet. Men
ved denne Afspeiling maa der dog ikke tænkes paa nogen umiddelbar
Communication imellem Monaden og Omverdenen. Det physiske
Speil modtager fra Gjenstanden selv det Lys, ved hvis Tilbagekastning % note breaks inside "Tilbage-|kastning": p. 105 | p. 106
Billedet dannes; men det metaphysiske Speil danner sine
Billeder af sit indre Anlæg uden umiddelbar Paavirkning udenfra:
imellem disse Billeder og Omverdenens Gjenstande kan der altsaa
kun være en ved Parallelisme sikkret Overeensstemmelse, ikke Cau%
salforhold.}

% --- p. 105 ---
\opage{105}Endskjøndt Cartesius maa tilstaae sig selv, at en Triangels
Væsen f. Ex. er saa fuldkommen bestemt, at de mange
% SEAM: p. 105 ends at a word boundary ("at de mange"); the body continues on p. 106 with "Egenskaber".
% SEAM: the footnote anchored on p. 104 ends on p. 106 ("... ikke Causalforhold.") and is set here in full;
% the next batch must NOT transcribe the footnote lines at the foot of p. 106.
% --- p. 106 ---
\opage{106}Egenskaber, der kunne bevises om den, saaledes afhænge af
Tanken, at den Tænkende maa erkjende dem, enten han vil
eller ikke, saa tør han dog, naar Spørgsmaalet er om Sandheden
af de mathematiske Sætninger, ikke ligefrem stole paa
Tænkningens Vidnesbyrd, men bygger paa den Forvisning, at
\emph{Gud er og ikke bedrager.} I sin Stolen paa Gud, altsaa
i at være vis, ikke paa sig selv, men paa et Andet end
sig selv, er Tanken umiddelbart vis paa sig selv: dette er
Modsigelsen. At Tanken selv har underskudt det Absolute,
hvortil den i sidste Instans appellerer, er imidlertid en Kjendsgjerning,
hvorom de supponerende Functioner ikke selv kunne
give nogen Oplysning.

Skal Tænkningen søge Sandhedens Garantie udenfor sig
selv, naar den beskjæftiger sig med sit eget logiske Indhold,
hvormeget mere da, naar Talen er om de i Legemverdenen
existerende Ting. At de legemlige Ting --- saaledes raisonnerer
Cartesius --- kunne existere, fremgaaer deraf, „at
Gud kan skabe Alt, hvad jeg saaledes er istand til at opfatte.
Men nu er der i mig en vis passiv Evne til at sandse eller til at
modtage og erkjende Sandsetingenes Ideer. For denne Evne
vilde jeg ingen Brug have, dersom der ikke tillige existerede en vis
activ Magt, enten hos mig selv eller hos Andre, til at frembringe
disse Ideer. Denne Magt kan imidlertid ikke være hos mig
selv, \dots\ da hine Ideer frembringes uden min Medvirken, ofte imod % UNSURE: a broken sort (like "¦") stands between "selv," and the dots; set as "selv, \dots"
min Villie. Den maa altsaa være i en anden Substans,
hvori al den Realitet indeholdes objectivt (\textit{formaliter}) eller i
høiere Grad (\textit{eminenter}), som subjectivt er i de af den frembragte
Ideer. Denne Substans maa enten være Gud, eller
% foot of p.106: end of p.104 note, set in pp91-105
% --- p. 107 ---
\opage{107}en Skabning, ædlere end det Legemlige, i hvilken disse Ideer
indeholdes i høiere Grad, eller selve det Legemlige, i hvilket
de netop indeholdes objectivt. Da der nu aldeles ikke er givet
mig nogen Evne til at erkjende, at Ideerne skulde bringes
mig af Gud selv, enten umiddelbart eller gjennem en anden
Skabning, men derimod givet mig en stor Tilbøielighed til at
troe, at de udsendes af de legemlige Ting, kan jeg ikke indsee
Andet, end at Gud maatte være en Bedrager, hvis de ikke
virkeligt udsendtes af de legemlige Ting: de legemlige Ting
existere altsaa.“\footnote{Jvfr. \textit{Meditatio Sexta.} S. 35; 39--40. (\textit{Opera philos.})} % note *)

Paa de supponerende Functioners Grundlag staae Ting
og Tanke, Dom og Gjenstand jo vistnok udvortes imod hinanden;
det Mechaniske og Physiske er ueensartet med det Intellectuelle
og Metaphysiske, det Materielle med det Aandelige:
men, hvad der dog er Pointen i Problemet, Modsætningen
mellem Begrebet og dets Existens, mellem det samme Indhold
i logisk og i antilogisk Form, mellem den samme Form med
logisk og antilogisk Indhold, kan paa dette Standpunkt alligevel
ikke komme til klar Bevidsthed. Antithesen af Begrebet og
Gjenstanden sløves, naar Kategorien Idee ligesaa umiddelbart
anvendes om Gjenstanden som om dens Begreb. Det Spørgsmaal,
hvorledes Ideen er istand til at optræde i to saa uligeartede
Skikkelser, som en legemlig Ting og en Forestilling, forflygtiges,
naar begge Skikkelser dog antages at ligne hinanden
deri, at de begge ere Ideer. Hertil kommer, at man jo i
Consequens af det underskudte Grundlag ved enhver Vanskelighed
har Almagten, altsaa et Magtens Ord, et guddommeligt
Magtsprog i al Vilkaarlighed at henvise til.

I de supponerende Functioner er Tanken oprindelig gaaet
fra sig selv, og dog mener den at være hos sig selv; denne
Modsigelse maa blive aabenbar.

2) \emph{Gjenstanden bestemmer Dommen: abstraherende
% --- p. 108 ---
\opage{108}Functioner.} Gaaer Tanken bogstavelig fra sig
selv, for at begynde med et Andet, da er det Første og Nærmeste,
den støder paa, ikke det Absolute, men det Relative, ikke
en Gud, men en sandselig Ting. Forsaavidt her da begyndes
med de fra Gjenstanden udgaaende Sandseindtryk, begyndes
her med det existentielle Extrem, det Antilogiske. De aprioriske
Begreber, „de medfødte Ideer“, ere her forsvundne; kun
de umiddelbare Existensformer, de sandselige Gjenstandsbestemmelser
staae tilbage. Vi have da igjen den sensualistiske
Erfaringslære. At Locke lader Sjælen begynde som en \textit{tabula
rasa}, faaer i nærværende Sammenhæng den Betydning, at
han gjør Forsøg med at lade alle Begrebsbestemmelser fremkomme
ved Omdannelse af individuelle Gjenstandsbestemmelser.
Det er da ikke Dommen, som ved Forudsætninger \textit{a priori}
skal bestemme Gjenstanden; det er tvertimod Gjenstanden, som
ved Existensens Magtord skal bestemme Dommen. Betragte
vi Theorien, saaledes som den anlægges af Locke og udvikles
af Hume, noget nærmere, da er det fornemmelig to Factorer,
der spille Hovedrollen, nemlig: \emph{Impressionen} og \emph{Abstractionen}.

Ved Impressionen skulde den existerende Gjenstand da
afpræge sin individuelle Form i Sjælen paa lignende Maade,
som Seglet afpræger sine Charakterer i Voxet. Dersom
det nu blev ved blotte Impressioner, da vilde der jo findes
ligesaa mange Individualformer i Sjælen, som der er Gjenstande,
Sjælen har opfattet. Formen for Gjenstandens
Billede i Sjælen og Gjenstanden selv i Tilværelsen vilde paa
denne Maade være den samme; men Indholdet vilde være
forskjelligt, forsaavidt Gjenstandsbilledet i Sjælen kun var og
blev et Billede, en Copie af det Virkelige, medens Virkeligheden
selv, Copiens Original, maatte søges i den existerende
Gjenstand. Men hvor passiv man end vil gjøre den Impressionerne
underkastede Sjæl, er Sjælen dog væsentlig
% --- p. 109 ---
\opage{109}Subjectivitet, og Subjectiviteten Selvvirksomhed.\footnote{Jvfr. Grundideernes Logik. I. S. 25.} % note *)
Under den
sjælelige Passivitet kommer Selvvirksomheden tilsyne i den
Forvandling, Gjenstandsformen undergaaer, idet den optages
i Sjælen: i Gjenstanden selv har Existensformen Enkelthedens
individuelle Præg; opfattet af Sjælen antager den
Almindelighedens Præg; det er en gjennemgaaende Forvandling,
denne gjennemgaaende Forvandling skyldes Abstractionen.

Da Abstractionen ikke lægger Noget til, men tvertimod
drager Noget fra, saa synes det ved første Øiekast, som var
det Existensformen selv, der blev bevaret reen og uforfalsket i
Bevidsthedsformen. Hvad er det vel, som ved Abstractionen
falder bort? Det Tilfældige, det blot Individuelle, det Foranderlige,
det Uvæsentlige! Og hvad er det, som ved Abstractionen
udtrykkelig fastholdes? Det Fælles, det Almene,
det Bestandige, det Væsentlige! Men Gjenstandsformen er jo
i sin objective Bestemthed antilogisk, Bevidsthedsformen i
subjectiv Bevægelighed væsentlig logisk; ved Abstractionen maa
der da være skeet det Vidunder, at en antilogisk Form er
bleven logisk; altsaa en \emph{anden} Form, og dog \emph{den samme}
Form, hvorledes er dette at forstaae?

Abstractionen, denne receptive, negative Function, der
alene skulde gaae ud paa at fastholde Gjenstandsformen som
reen, uforvandsket, fra det Tilfældige befriet Existensform,
har aabenbart gjort Noget ved Formen, har, saa at sige, paa
egen Haand omarbeidet den, har forvandlet den fra antilogisk
Individualform til logisk Almeenbestemmelse, fra Existensform
til Begrebsform. Abstractionen er da ikke saa abstract receptiv,
at den jo ogsaa kan være productiv. Abstractionen er
en Function af apriorisk Selvhed; men den i abstraherende
Functioner fortabte Tænkning opdager det ikke.

Theoriens Resultat er, at baade Begrebs- og Existensformen
opløses, Dommen suspenderes og Gjenstanden bliver
% --- p. 110 ---
\opage{110}borte. Begrebet maa nødvendig forsvinde, naar det Almene,
det i sig Nødvendige, forsvinder. Nødvendighed og Almeengyldighed
ere nemlig Særkjender for alle sande Begrebsbestemmelser;
men det Nødvendige og Almeengyldige lader sig
ikke abstrahere af Kjendsgjerningernes umiddelbare Enkeltheder.
Vel paastaaer Hume, at Kategorier som Lighed, Grad,
Størrelse o. s. v. give umiddelbar Vished ved Abstractionen
og den tilsvarende Reproduction; „thi naar et Object er liigt
et andet, vil Ligheden strax afficere Øiet eller snarere Sindet,
og sjeldent fordrer den anden Prøvelse; ved Bestemmelsen af
Størrelser eller Tal kunne vi gaae frem omtrent paa samme
Maade“ o. s. v. Med Kategorierne Aarsag og Virkning skal
det derimod forholde sig paa ganske anden Maade. Dette
er nu aabenbart en Vildfarelse; thi dersom Causaliteten ikke
lader sig abstrahere af Kjendsgjerninger, saa gjælder det
Samme om Ligheden; og dersom Ligheden lader sig abstrahere,
saa gjælder det Samme om Causaliteten.

Hvorfor lader Aarsagsforholdet sig ikke abstrahere af
Kjendsgjerninger? Fordi dette Forhold som Forhold ikke kan
blive Gjenstand for umiddelbar Erfaring, fordi med andre
Ord Forholdet selv ikke er en sandselig Gjenstand. Men nu
Ligheden! er Ligheden maaskee en sandselig Gjenstand? Hume
antager, at Ligheden i det Lige umiddelbart afficerer Øiet;
det er en barnagtig Forestillingsmaade. Kategorierne Lighed
og Ulighed opstaae kun ved Sammenligning; men Tingene
sammenligne sig ikke selv med hinanden: Sammenligningen er
en Function af den dømmende, abstraherende Subjectivitet.
Sammenligning er en Reflexion, der kun finder Lighed ved
at forudsætte Lighed; for at opdage Ligheden i Tingen, maa
den reflecterende Subjectivitet selv medbringe Ligheden, gaae
ud fra Lighedens Kategorie. Det Samme, det Selvsamme
gjælder om Causaliteten. Forskjellen er kun den, at Causalitetsreflexionen
er indviklet, medens Lighedsreflexionen falder
som af sig selv.

% --- p. 111 ---
\opage{111}Hvad vil det nu sige, at Ligheden abstraheres af de lige
Gjenstande, Uligheden af de ulige? Dog nei! det er i denne
Sammenhæng vigtigere at bestemme, hvad det ikke vil sige.
Thi det vil paa ingen Maade sige, at den sandselige Ting,
foruden dens andre Egenskaber, ogsaa skulde have Egenskaben
Lighed, og at Abstractionen saa ved at fjerne de andre Egenskaber
kunde beholde Ligheden. To Ting være hinanden lige
i Farve, de være f. Ex. begge røde; der er dog Ingen, der
antager, at der i den lige Farve er to sandselige Egenskaber,
Farven og Ligheden; Ligheden i Farven er ikke sandselig, saaledes
som Farven. Farvefornemmelsen beroer paa et Sandseindtryk,
men Ligheden beroer paa en Dom. Abstractionen
bestemmes ved Dommen; i en Dom som: „dette er en Farve“
er Farven i Subjectet sandselig, i Prædicatet usandselig, i
Subjectet antilogisk, i Prædicatet logisk; kun saaledes kan
Kategorien i Prædicatet være en Abstraction af Gjenstanden
i Subjectet.

Dette have nu Hume og Sensualisterne ganske overseet.
Af Farven paa A og Farven paa B ville de uden videre abstrahere
Ligheden, ret som om Abstractionen bestod i at skrabe
Ligheden af Farven, ligesom man i visse Tilfælde kan skrabe
Farven af Gjenstanden. Kun naar Abstractionen reflecteres
i Dommen, forsvinder Illusionen: A og B ere i Farve hinanden
lige; her viser Bestemmelsen „Farve“ hen til umiddelbar
Sandsning, men Prædicatet „lige“ tilhører Reflexionen.

Det er, logisk seet, Sensualismens Axiom, at Gjenstanden
skal bestemme Dommen. Heri ligger den stiltiende Forudsætning,
at den existerende Gjenstand i og for sig har en
sandselig objectiv, af enhver Dom uafhængig Selvstændighed.
Men ogsaa denne Forudsætning opløser sig for de abstraherende
Functioner. Allerede Locke var i Forlegenhed med
Substansen; hans begyndende Tvivl blev udviklet hos Hume:
for Hume er der ingen Substans. „Er“, siger han, „Begrebet
Substans tilført os gjennem Sandserne, saa spørger
% --- p. 112 ---
\opage{112}jeg, hvilken Sands giver os dette Begreb, og hvorledes giver
den os det? Skulde det være modtaget ved Øinene, da
maatte det være en Farve; ved Ørene, da en Lyd; ved Ganen,
da en Smag, --- og saaledes ved alle øvrige Sandser.“

Hume kommer consequent til det Resultat, at der ved
Begrebet Substans ikke kan forstaaes Andet end en Samling
af enkelte Egenskaber (\textit{a collection of particular qualities}),
en Sammenfatten af eenfolde Begreber, der ere forenede ved
Imaginationen og betegnede med et eget Navn. „Saaledes
opdage vi først hos Guldet den gule Farve, Tyngden, Strækbarheden,
Smelteligheden; siden opdage vi, at det lader sig
opløse i Kongevand; denne Egenskab forbinde vi saa med de
øvrige, og henregne den ligesaa vel til Substansen, som hvis
Begrebet derom fra Begyndelsen af havde udgjort en Deel
af det sammensatte Hele.“ Hvad er da Guldet? Som Gjenstandseenhed
er Guldet et Navn, et Aggregat af forskjelligartede
Egenskaber, som ikke Virkeligheden, ikke Tilværelsen selv,
men „Imaginationen“ har forenet. Ved at isolere Egenskaberne
og opfatte hver enkelt Egenskab for sig, ledes Abstractionen til
at berøve Gjenstanden Substantialitet; men berøvet sin Substantialitet
er Gjenstanden tillige berøvet sin objective Selvstændighed.

Substansen er en imaginær Størrelse, der, ifølge Hume,
„fremkommer derved, at man tilskriver de særlige Egenskaber
et ubekjendt \emph{Noget}, i hvilket de antages at inhærere.“ Det
Magtens Ord, der skulde sikkre Gjenstanden en factisk bestemt,
substantiel Objectivitet, er altsaa ved Abstractionen forvandlet
til et Afmagtens Ord, et Navn, thi Substansen, dette Abstractionens
\textit{caput mortuum}, der bliver tilbage, naar alle Egenskaber
falde bort, er kun et Navn: skal Gjenstanden paa denne Maade
bestemme Dommen, saa er her hverken Dom eller Gjenstand.

3) \emph{Dommen bestemmer Gjenstanden: determinerende
Functioner.} Ved at gaae fra sig selv og søge
sine Indholdsbestemmelser i den tilværende Gjenstand har
% --- p. 113 ---
\opage{113}Tanken fra Grunden af forstyrret saavel sit eget som Gjenstandens
Væsen. For at raade Bod paa denne Grundforstyrrelse,
vender Tanken da fra Gjenstandens Opløsning tilbage
til sit aprioriske Udgangspunkt. Begyndelsen skeer fra
det subjective Extrem, idet Dommen bestemmer Gjenstanden.

At Dommen bestemmer Gjenstanden maa naturligviis
ikke forstaaes, som kunde nu Subjectiviteten vilkaarlig dømme
efter eget Tykke, og saa skulde Gjenstanden uden videre rette
sig efter Dommen. Sandheden er tvertimod den, at de for
Gjenstandens Erkjendelse afgjørende Væsensbestemmelser, for
at være Gjenstandsbestemmelser, absolut maae være Begrebsbestemmelser;
det er i Conseqvens af denne Theorie, at Functionerne
af Dommen vise sig som \emph{determinerende}.

Hvor umuligt det er at antage Gjenstandsformer, der ikke
oprindelig ere Subjectivitetsformer, har Kant tilstrækkelig godtgjort.
Det Eensidige hos Kant ligger deri, at han ved Subjectivitet
kun kan tænke sig en psychologisk Subjectivitet, medens
Opgaven omvendt er, ved den psychologiske Subjectivitet kun
at tænke paa det subjectivt Almeengyldige, paa det, der er
væsentligt for al Subjectivitet. Naar Kant erklærer Rums-
og Tidsanskuelserne for subjective Former, saa give vi ham
Ret; men naar han saa med det Samme antager, at disse
Former ikke kunne være objective, efterdi de ere subjective; at
det kun er Former for \emph{vor} Subjectivitet, uden at deraf følger,
at disse Former ogsaa have Gyldighed for en Subjectivitet
af høiere Orden, hvis en saadan gives: da kunne vi ikke give
ham Ret.\footnote{Grundideernes Logik. I. S. 77--80.} % note *)

Den Skarpsindighed, hvormed Kant i sin „transcendentale
Analytik“ udleder Forstandskategorier af Dommene, har vakt
Beundring og fortjener Beundring. Det er herved een Gang
for alle godtgjort, at Ingen kan dømme om Tingens objective
Væren, uden at han ved at reflectere over Dommen tilsidst
% --- p. 114 ---
\opage{114}maa blive sig bevidst, at den objective Væren sættes ved en i
Forstanden selv liggende Kategorie. Den værende Gjenstand
\emph{er}, forsaavidt Subjectiviteten tildømmer den Væren. Vil
Nogen mene, at Gjenstandens objective Væren i Realiteten er
uafhængig af den subjective Væren i Dommen, lad ham gaae
tilbage til Hume og tage, hvad der følger paa. Men naar
Kant saa i den for Dommens Qvalitet afgjørende Adskillelse
af Væren og Ikke-Væren --- „Realitet og Negation“ --- kun
seer en Fremstilling af psychologisk-subjective Former, Forstandsformer,
der ere saa subjective, at de umulig kunne passe
paa Tingen i sig, da er han i Vildfarelse.

Naar der dømmes om Tingen som Phænomen: denne
Ting er haard, rød o. s. v., da veed Kant meget vel, at den
subjective Væren i Dommens Copula ikke gjør den virkelige
Tilværelse Afbræk, men tvertimod determinerer Phænomenets
objective Væren. Thi det forholder sig jo ikke saaledes, at
Phænomenet, Erkjendelsens Gjenstand, har en af Dommen
uafhængig og derfor mod Dommen ligegyldig Tilværen, og
at denne Gjenstandsværen saa, naar det træffer sig, bliver opfattet
og betegnet ved Dommens Copula. Langtfra! Forholdet
er just det omvendte, at Phænomenet dømmes subjectivt
til at være objectivt: Phænomenet har sin Objectivitet
fra den dømmende Subjectivitet.

Anderledes, naar der tales om Tingen, ikke som en Erkjendelsens
Gjenstand, men som „Tingen i sig.“ Her vil nu
Kant absolut have en Væren saa objectiv, at den paa ingen
Maade skal kunne lade sig opfatte ved nogen subjectiv Forstandskategorie,
altsaa heller ikke ved den til Dommens Copula
svarende Væren. Ved denne Antagelse kommer Fornuftkritikens
Forfatter i uendelig Modsigelse med sig selv. Naar
Kant siger om Tingen i sig, at den \emph{er}, saa fælder han en
Dom. Han maa nu anbringe saamange Distinctioner imellem
Realitet og Væren, mellem objectiv og subjectiv, ubegribelig
og begribelig Væren, som han vil: saa staaer det dog fast, at
% --- p. 115 ---
\opage{115}de Distinctioner, hvorved Tingens Væren skal lægges udenfor
Forstandskategorien, beroe paa lutter Forstandskategorier; endvidere
staaer det fast, at han ikke kan erklære det for en Umulighed
at dømme om, hvad Tingen i sig er, uden netop derved,
at han dømmer om, hvad Tingen i sig er. Siger han:
Tingen i sig er et X, saa har han ved denne Dom ligesaa
vist ladet en Forstandskategorie være determinerende for Tingen
i sig, som han ved at sige: Tingen er haard, har ladet den
være determinerende for Tingen som Phænomen.

Hvad her er sagt om de Kategorier, der, ifølge Kant,
ligge til Grund for Dommene efter deres Qvalitet, gjælder
ogsaa om dem, der udledes af Dommenes Inddeling efter
Qvantiteten, nemlig: \emph{Alhed}, \emph{Fleerhed}, \emph{Eenhed}; af Inddelingen
efter Relationen: \emph{Substans} og \emph{Inhærens},
\emph{Aarsag}, \emph{Virkning} og \emph{Vexelvirkning}, og endelig efter
Modaliteten: \emph{Mulighed}, \emph{Virkelighed} og \emph{Nødvendighed}.

Er det først beviist, at Forstandskategorien \emph{Væren} maa
være determinerende for Tingen i sig, da gjælder det Samme
om alle de øvrige Kategorier. Hører Tingen i sig til en
Sphære, hvori Kategorien \emph{Væren} har Gyldighed, da hører
den netop til en Sphære, hvori Kategorierne \emph{Ikke-Væren},
Fleerhed, \emph{Eenhed}, \emph{Aarsag}, \emph{Virkning} o. s. v., have Gyldighed.
Omvendt: have Kategorierne \emph{Ikke-Væren}, \emph{Fleerhed}, % printed with "Fleer-" letterspaced, "hed" (next line) not
\emph{Eenhed}, \emph{Aarsag}, \emph{Virkning} o. s. v., ingen Gyldighed,
naar Talen er om Tingen i sig, da har Kategorien Væren heller
ingen Gyldighed. Man kan da ligesaa lidt sige om Tingen i
sig, at den er, som, at den ikke er; idet Tingen i sig sætter sig
op imod Dommen, ophører den at være Gjenstand.

Den kantiske „Ting i sig“ kan betragtes som en eiendommelig
Repræsentant for det Antilogiske, forsaavidt det
Antilogiske staaer Tanken imod og af den Grund ikke lader sig
tænke. Men hos Kant er Opfattelsen af det Antilogiske dog
i høi Grad forfeilet. For det Første er „Tingen i sig“ an%
% --- p. 116 ---
\opage{116}tilogisk i den Forstand, at den ikke lader sig opfatte i nogensomhelst
Tankebestemmelse, og derfor holder sig som et for
den menneskelige Viden utilgængeligt X. Dette er nu for
Bestemmelsen af det Antilogiske baade Formeget og Forlidet.
Formeget; thi vel staaer det Antilogiske Tanken imod, men
saaledes, at det just ved Modsætningen baade bestemmer Tanken
og lader sig bestemme af Tanken. Forlidet; thi „Tingen
i sig“ hjemfalder til Tanken med dens determinerende
Functioner paa to aldeles forskjellige Maader. Her forudsættes
nemlig, at dette os Ubekjendte, for os Utænkelige, muligviis
kan tænkes af en guddommelig Tænken: Tingen i sig
er altsaa ikke i sig selv absolut utænkelig. Her forudsættes
endvidere, at Tingen i sig kan paa en vis Maade gaae i
Eet med vore Erkjendelsesformer og subjective Forstandskategorier.
Thi hvad er det vel, der udfylder Erkjendelsesformerne
og giver Erfaringen sit bestemte Indhold? Aabenbart „Tingen
i sig“. Men saa maa Phænomenet jo, hvad i en tidligere
Sammenhæng er viist, netop være en Identitet af Tingen i
sig og Erkjendelsesformerne; eller i Phænomenet maa Begrebet
og Tingen være identiske. --- Dernæst, og dette er
vor anden Hovedindvending, har Kant forfeilet det Antilogiske
derved, at han, med alt det Væsen, han iøvrigt gjør af Adskillelsen
mellem Ting og Begreb, dog i al Ubestemthed
lader Tingen i Erfaringen og Begrebet i Tænkningen have
samme Form, kun at Begrebets Form endnu mangler Erfaringens
Indhold. Derved har han overseet det egentlige
Problem. Indeholder det Virkelige i Begrebet ikke Mere end
det blot Mulige; er her i Begrebets Form Intet, slet Intet,
der tyder paa Forskjel mellem Mulighed og Virkelighed, saa
spørge vi: ved hvilken Form og paa hvilken Maade skal da
Bestemmelsen af det bestemte Virkelige lade sig adskille fra
Bestemmelsen af det alsidig bestemte Mulige? Her maa nødvendig
være et Skillemærke i Formen selv; thi at addere den
empiriske Anskuelse til det rene Begreb er dog nok omtrent
% --- p. 117 ---
\opage{117}ligesaa vanskeligt, som at substrahere Himlen fra den blaa
Himmel.

Den eiendommelige Maade, hvorpaa Kant i sin Fornuftkritik
behandler Forholdet imellem Erfaringens formelle og
dens materielle Betingelser, egner sig ret til at hidse den metaphysiske
Reflexion i eensidig Retning. Overalt seer det
ud, som var der i Erfaringsgjenstanden, foruden Erkjendelsesformen,
et ubestemmeligt Noget, der satte Erkjendelsen Skranker,
et fremmed Element, som Tanken ikke kunde faae Bugt med;
og saa er det dog ikke Noget --- Modstanden viger allevegne for
den i Analysen sig fremarbeidende Tanke. Saaledes blev nu
Opmærksomheden hos Fichte udelukkende henvendt paa Tankens
overgribende Magt: Tanken kom til sig selv og saae i Alt
kun sig selv. Havde Reflexionen først opløst Tingen i sig og
trods Virkelighedens Protest tilbagetrængt enhver Forestilling
om Existensens antilogiske Characteer, saa kunde man heller
ikke afvise den subjective Idealisme ved at beraabe sig paa den
Umiddelbarhedens Magt, hvormed det udvortes Virkelige optræder;
det Virkeliges umiddelbare Magt blev da i sit Væsen
erkjendt, ikke for en antilogisk Factor, men for en logisk Magt,
en reen Tankemagt: det Objective vendte tilbage, ikke som
Dommens Gjenstand, men som Dommens Indhold, ja, hos
Hegel, som Dommen selv.

\textit{ββ)} \emph{Magten i Dom og Gjenstand: praktiske
Functioner.} Dersom den hegelske Udvikling af Begrebets
Magt var alsidig tilfredsstillende, vilde Modsætningen af Dom
og Gjenstand opløse sig i en til det psychologiske Standpunkt
svarende Illusion. Istedenfor antilogiske Magtbestemmelser fik
vi da en existentiel Umiddelbarhed; denne Umiddelbarhed,
skjøndt den i Phænomenet synes at være Begrebet modsat,
var da væsentlig kun et Moment i Begrebet selv, efterdi Begrebets
Realitet, dets logisk overgribende Magt just bestaaer i
at sætte og forudsætte de Elementer, gjennem hvis Mediation
det giver sit Indhold Uendelighedens Form og aabenbarer den
% --- p. 118 ---
\opage{118}absolute Idee. Det i denne Theorie Eensidige og Uholdbare
er paa tidligere Stadier af Grundideernes Logik allerede paaviist
fra flere Sider.\footnote{Grundideernes Logik. I. S. 148--155 og S. 162--168 o. s. v.} % note *)
Ikke destomindre vender Theorien
paa ethvert nyt Stadium tilbage med fornyet Eftertryk, og
ganske naturligt, saasandt den hegelske Philosophie er en i sin
Art conseqvent Afslutning af alle forudgaaende Systemer.

Er der ingen af Fortidens Tænkere, der med Rette kan
siges at have reist Problemet om det Antilogiske; er Kant
blandt Alle den, der er kommen Tingen nærmest, uden dog at
have grebet Tingen; er Alt, hvad tidligere Systemer under
forskjellige Benævnelser have sat som Realitetens Kjendemærke
imod Tanken, forfeilet, deels saaledes, at det i sin absolute
Utænkelighed maa blive selvmodsigende og meningsløst, deels
saaledes, at det som en tvetydig Mulighed kun kan tilhøre en
underordnet Sphære og altsaa være bestemt til at forsvinde:
hvad Under da, at en logisk Gjennemførelse af Reflexionen
ender med at forvandle Tilværelsen til et Begreb.

Kun i de antilogiske Formers Logik kan Panlogismen
finde sit væsentlige Correctiv. Tidligere er Vidensbegrebet
under forskjellig Belysning stillet imod Magtbegrebet; Modsætningen
forhøies nu derved, at hvert af Begreberne faaer
sin tilsvarende Udvikling i sit tilsvarende Ord, Videns Begreb
i Videns Ord, Magtbegrebet i Magtens Ord. Kunde Magtbegrebets
Indhold udvikles i Videns Ord, da maatte Modsætningen
af Existens og Begreb opløse sig i et Skin; kun
i Forestillingens, altsaa i Usandhedens, Lys vilde Existensbestemmelserne
tage sig ud som antilogiske Gjenstandsbestemmelser;
seete i Sandhedens Lys vilde de derimod vise sig at
være rene Begrebsbestemmelser; Panlogismen vilde da i Principet
have fuldkommen Ret.

Skal Panlogismen overvindes, maa Magtbegrebet udvikles
i Magtens Ord. Men hvor komme vi nu hen? Som
% --- p. 119 ---
\opage{119}Magten er, maa Ordet være; Magten er antilogisk, Ordet
ligesaa. Som Ordet er, maa Sproget være; Ordet er antilogisk,
Sproget ligesaa. Som Sproget er, maa dets Sætninger,
Domme og Begreber være; Sproget er antilogisk,
Sætninger, Domme og Begreber ligesaa. Antilogiske Begreber,
antilogiske Sætninger, antilogiske Domme! og alt dette skal
behandles i en Logik, det Antilogiske altsaa behandles logisk:
er her ikke et Væv af uopløselige Modsigelser? --- Før vi gaae % UNSURE: Før or For (stroke of ø not visible; text layer "For")
ind paa Sagens Realitet, ville vi betragte Vanskelighederne
fra den formelle Side.

Den formelle Logik gjør, som bekjendt, en endelig Adskillelse
imellem Tanken og det, der tænkes, imellem Begrebet
og det, der begribes, en Adskillelse, der finder sin Stadfæstelse
i Talemaader som: at tænke \emph{paa} Noget, at tænke \emph{over} Noget,
danne sig et Begreb \emph{om} Noget o. dsl. I Uvillie mod det
saaledes endeliggjorte Forhold mellem Tanke og Ting har
den objective Logik saa igjen fra sin Side fastholdt Identiteten
af Tænken og Væren som saa immanent, at det Værende ved
at tænkes blev til Tanke, det Existerende ved at begribes til
Begreb. Den objective Logik tænker ikke paa Noget; den
tænker Noget, og dette Noget er Tanken selv (ὥστε ταὐτόν % UNSURE: final accent printed acute (ταὐτόν), not grave; breathing on αυ faint
νοῦς καὶ νοητόν. Aristoteles.). Den objective Logik danner
sig ikke et Begreb om Noget, --- det vilde være langt under
dens Værdighed, --- men den begriber Noget, og det Noget,
den begriber, er Begrebet selv.

Begge Synsmaader ere, naar de fastholdes hver for sig,
lige eensidige og ved deres Eensidighed lige skadelige for
sand Begrebsudvikling. Den formelle Logiks Adskillelse af
Tanken og dens Gjenstand er eensidig, forsaavidt den overseer
Identiteten; den objective Logiks Identitet af Tanken og det
Tænkte er eensidig, forsaavidt Modsætningen oversees. For
at forsvare det Paradox, at Tankens Gjenstand kun kan være
Tanken selv, beraaber den objective Logik sig paa den immanente
Negation. Ved sin immanente Negation er Tanken istand til
% --- p. 120 ---
\opage{120}at hæve sig over sig selv og gjøre Tanken til Gjenstand; ved
sin immanente Negation er Tanken omvendt istand til at
forvandle den tænkte Gjenstand til Tanke. Gjenstanden er,
med andre Ord, Tankens Andet; men dette Andet er ved den
immanente Negation Tanken selv: her er Tanke af Tanke af
Tanke i Tænken af Tænken af Tænken, og videre naaer man
i Grunden ikke.

For at komme ud over den grændseløse Formalisme,
maae vi nødvendigt skjelne imellem den blotte Formalnegation
og Realnegationen. Saalænge Tankens Andet kun er Tanken
selv, og alle Formbestemmelser forsaavidt logiske, er Negationen
formel; først naar Tankens Andet ikke længere kan forvandles
til Tanke, altsaa naar det Antilogiske indtræder, er Negationen
reel. Vil man nu i Kraft af Realnegationen erklære det
Tankens Andet, der ikke selv er Tanke, for noget i sig Utænkeligt,
da vil man ogsaa let kunne forsvare det Paradox, at
Tanken i Virkeligheden tænker det Utænkelige. Sagen er, at
dette saakaldte Utænkelige i Realiteten er tænkeligt: tænkeligt
derved, at det opfattes af den Tænkende, utænkeligt derved, at
det ikke er Tanke og ikke kan tænke sig selv. Paa Realnegationen,
paa en Tænken af det, der ikke er Tanke, altsaa paa
en logisk Bestemmelse af det Antilogiske, beroer al Tænknings
Realitet.

At Tænkningens Realitet maa falde sammen med Sagens
Realitet, er en Selvfølge. I Sagens Realitet maa her altsaa
--- det ligger i Realnegationen --- være en virkelig Modsætning
af Begrebet i Vidensordet og Existensen i Magtordet;
Begrebet i Vidensordet er Dommen, Existensen i Magtordet
er Gjenstanden: den virkelige Modsætning af Begreb og
Existens er en virkelig Modsætning af Dom og Gjenstand.
Skal den virkelige Modsætning af Dom og Gjenstand nu
ikke synke ned i Endelighedens Udstykning og fixeres i psychologisk
Dualisme, da maa der til den virkelige Modsætning af
Videns og Magtens Ord nødvendigt svare en virkelig Eenhed;
% --- p. 121 ---
\opage{121}en saadan virkelig Eenhed kan ikke være blot theoretisk, den
virkelige Eenhed er her kun, forsaavidt den virkelig tilveiebringes;
den tilveiebringes kun ved en virkelig Overgang, en
Forvandling: Tilveiebringelsen er væsentlig praktisk. For at
Logiken nu skal kunne give en ved Vidensordet bestemt Afspeiling
af den til Magtordet svarende Praxis, er det nødvendigt
at begynde med Forholdet mellem Tanken og Magten.

Theorien, den rene Theorie, maa nøies med Tankebevægelser;
dens Functioner ere kun Tankehandlinger. Ved rene
Tankehandlinger er det, ifølge Realnegationen, umuligt at
forene Viden og Magt. Den virkelige Eenhed, der i det
Uendelige er tilveiebragt og skal tilveiebringes, forudsætter en
væsentlig Eenhed, der hverken kan eller skal tilveiebringes,
fordi den evig er; den er i Subjectiviteten selv, det absolute,
i sig reflecterede Princip. Ifølge den endelige Modsætning
mellem Viden og Magt er Tankehandlingen afmægtig og
Magthandlingen i sin Umiddelbarhed tankeløs; men i Kraft
af Magtens og Videns oprindelige Eenhed er Tankehandlingen
en forvandlet Magthandling, og derved selv mægtig; Magthandlingen
omvendt en forvandlet Tankehandling, og selv
tankefast. Idet Magthandlingen omsættes i Tankehandling,
og Vidensordet assimilerer Magtbestemmelserne, gaaer Bevægelsen
fra det Antilogiske til det Logiske, Gjenstandsformerne
blive Begrebsformer, og Dommen logisk. Idet Tankehandlingen
omvendt gaaer over i Magthandlingen, og Magtordet
indsuger Begrebsbestemmelsernes Idealitet, skeer Bevægelsen
fra det Logiske til det Antilogiske, og Dommen selv bliver
antilogisk. De praktiske Functioner ere i begge Bevægelser
samtidige, forsaavidt Bevægelserne danne et uendeligt Kredsløb;
de til dette Kredsløb svarende Vexelbestemmelser af logiske
og antilogiske Momenter forklares i en begribende Reflexion,
som gjør Dommen paa een Gang logisk og dog
antilogisk.

1) \emph{I den logiske Dom præformeres Gjenstan%
% --- p. 122 ---
\opage{122}den: anticiperende Functioner.} Dersom Dommen
virkelig skal kunne bestemme Gjenstanden, maa Viden blive
til Eet med Magten; denne Eenhed tilveiebringes ved et
System af anticiperende Functioner. Der tales i Teleologien
ofte nok om Anticipationer og Præformationer, uden at det
just derfor bliver klart, hvorledes det egentlig gaaer til med
den præformerende Virksomhed, og hvorvidt Anticipationen
oprindelig er en theoretisk eller praktisk Function. Hvilket
Forhold er her da imellem Præformation og Anticipation?
hvad maa Anticipationen nødvendigt forudsætte? hvorledes
foregaaer den?

Da Præformationen, som Ordet lyder, gaaer ud paa at
give Gjenstanden en indre, forud for den ydre gaaende, Tilværelse
og da med det Samme en indre, forud for den ydre
gaaende, Skikkelse, saa følger det af sig selv, at al Præformation
maa beroe paa Anticipation; hvad er Anticipation vel
Andet end en første Formen af en sig selv forudsættende Form,
en vordende Form uden virkelig Tilværelse. I Præformationen
er det umuligt at paavise en umiddelbart første Form:
thi Præformationen siger i sit Begreb, at enhver tilværende
Form maa resultere af en forudgaaende Form. Kunde der i
Virkeligheden gives en første tilværende Form, da kunde den
ene tilværende Form gaae forud for den anden, forberede den
anden og afløses af den anden uden nogen egentlig Præformeren.
I Præformationen er den forberedende Form derimod
saa væsentlig bestemt ved den resulterende, at den paa een
Gang baade er dens Afbillede og dens Forbillede, og kan kun
være det Sidste ved at være det Første. Den virkelige Skikkelse
er altid, om dette Ord kan bruges, en Denneskikkelse,
ligesom det Tilværende altid er et Detteværende. Er det nu
ikke en Skikkelse i ubestemt Almindelighed, men denne, just
denne Skikkelse, der skal præformeres, da maa den nødvendigviis
forudsættes, for at blive sat: den maa anticiperes. At
en Skikkelse, som endnu ikke er til, og følgelig ikke har det
% --- p. 123 ---
\opage{123}Tilværendes Bestemthed, skal indvirke bestemmende paa den
Begyndelsesform, hvorved den bestemmes, denne Præformationens
Forudsætten af det, der skal sættes, er netop Anticipationen.

Af den i Anticipationen reflecterede Uendelighed fremlyser,
at dens Væren selv er en Reflecteren, og at denne
Reflecteren nødvendigt maa forudsætte en reflecterende Subjectivitet.
Alle anticiperende Functioner ere Subjectivitetsfunctioner.
Men ligesom Selvheden med Nødvendighed forudsætter
Andetheden, saaledes maa Subjectivitetens anticiperende
Virksomhed nødvendigt forudsætte Objectiviteten. At
Anticipationen bestemmes ved Hensynet til en tilkommende
Objectivitet, er øiensynligt; tænkes det Efterfolgende bort, saa % PRINTER: printed "Efterfolgende", sense "Efterfølgende"
ophører det Foregaaende at være et Forudgaaende. Men
ikke blot i formel, ogsaa i reel Forstand maa Anticipationen
forudsætte en given Tilværelse. Hvad der i Anticipationen
sættes subjectivt, sættes jo netop med den udtrykkelige Bestemmelse,
at det engang skal blive objectivt. Men objectivt kan
Intet blive, uden det, som i sin subjective Begyndelse allerede
har Betingelser for Objectivitet. Betingelser for Objectivitet
kan kun det have, der oprindelig er bestemt ved objective
Hensyn; men objective Hensyn, Hensyn til det Objective, forudsætte
en Verden af allerede tilværende Objecter. Deelt
mellem det Subjective og det Objective, eller rettere i sin
udeelbare Eenhed deelviis bestemt ved begge, hører Anticipationen
altsaa under de relative Begreber: en absolut Anticipation
er ligesaa umulig, som en absolut Begyndelse.

Af Anticipationsbegrebets Relativitet følger, at Virkeligheden
i uendeligt Omfang umulig kan anticiperes paa een
Gang. Vel maa enhver relativ Totalitet, ethvert virkeligt
Led fremgaae af sin Mulighed; men derfor kan dog hele
Virkeligheden ikke ligefrem have sin Begyndelse i en forud for
alle virkelige Ting gaaende Mulighed: Muligheden har ikke
blot Virkeligheden foran sig, men ogsaa bagved sig; den virke%
% --- p. 124 ---
\opage{124}lige Mulighed beroer paa Virkelighed. Ligesaa med Anticipationen:
den forbereder et Tilværende ved et allerede Tilværende.
Af Forholdet mellem det Virkelige og Mulige, det Objective
og Subjective i Anticipationen, sees tillige Forholdet mellem
det Theoretiske og Praktiske.

Forsaavidt Anticipationen gaaer ud fra det Objective, ere
de anticiperende Functioner væsentlig theoretiske. Det Theoretiske
er, at Dommen bestemmes ved Gjenstanden, Tanken
ved det Givne, at Virkeligheden, det Objective, raader. Med
Blikket paa det objectivt Virkelige erkjender den anticiperende
Subjectivitet sin virkelige Grændse, sin Skranke. Ligeoverfor
det virkelige Objective er det abstract Subjective afmægtigt.
Eller hvad skulde den rene Subjectivitet vel være istand tll % PRINTER: printed "tll", sense "til"
at forandre i den objective Virkelighed? De virkelige Gjenstande
maatte jo ikke være Objecter, men pure Indbildninger,
rene Drømmebilleder, dersom de uden videre kunde opløses
og forvandles ved en reent subjectiv Bestemmen. Hvad der
gjælder om de objective Led, gjælder naturligviis ogsaa om de
Relationer, hvori Leddene væsentlig lægge deres Objectivitet
for Dagen. Enhver Forandring med det Objective maa skee
paa objective Betingelser, i Conseqvens af objective Forudsætninger.
Vel maa det Objective objectiveres; men det bestemmer
selv sin Objectiverings Indhold, og beviser derved sin
Realitet.

Men dersom Subjectiviteten virkelig i enhver Henseende
er afmægtig ligeoverfor det Objective, saa er Anticipationen
jo en Umulighed: der fordres en Eenhed af Viden og Magt,
en Eenhed, som Anticipationen selv baade forudsætter og tilveiebringer.
Den anticiperende Virksomhed er præformerende;
Præformationen forudsætter just hos Subjectiviteten en overgribende
Magt over det Objective. Uden en saadan overgribende
Magt vilde den subjective Formen ikke være bestemmende
for den objective Form. Som den subjective Virksomhed,
der forbereder og danner en tilkommende Virkelighed,
% --- p. 125 ---
\opage{125}maa Anticipationen altsaa være en ideel Praxis; men fra
hvilken Side kan da det Objective hjemfalde til at bestemmes
af det Subjective? Fra Mulighedens, den reale Muligheds!
Thi den reale Mulighed en paa een Gang baade objectiv og % PRINTER: printed "en", sense "er"
subjectiv.

Forsaavidt Muligheden er objectiv, kunne Objecterne paa
utallige Maader omdannes og forandres ved hinanden indbyrdes.
Som existerende Gjenstand, som ydre Virkelighed,
har Objectet i sig Mulighed til, paa forskjellige Betingelser,
at spille en høist forskjellig Rolle. I sin umiddelbare Virkelighed
er Objectet af bestemt Eiendommelighed, stiller sig mod
Subjectiviteten og gjør sin Selvstændighed gjældende; men
efter Muligheden er det Objective afmægtigt, efter sin Mulighed
er Objectet stofagtigt og ubestemt; at kjende Muligheden,
er at kjende den Side, hvorfra Objectiviteten subjectivt lader
sig angribe: ved sin Viden af objective Muligheder faaer den
anticiperende Subjectivitet en ideel Magt over Objecterne.

Hvor stofagtig Objectiviteten er, sees af de objective
Muligheder. Der dannes aldrig saamange individuelle Skikkelser,
der kunde jo, paa Grund af Mulighederne, dannes
flere; Virkeligheden formes aldrig i saa særlige Typer, den
kunde jo, paa Grund af Mulighederne, formes i andre Typer:
her bliver i de objective Dannelser overalt et uendeligt Overskud
af Muligheder. Det er dette Overskud af Muligheder,
der altid staaer til Subjectivitetens Raadighed; Subjectiviteten
dømmer, og den dømmende Subjectivitet veed at behandle
Muligheden, thi Subjectiviteten \emph{frigjør} Muligheden.
I sin objective Form er Muligheden bunden, og den bundne
Mulighed priisgiven Tilfældigheden; først i den subjective
Form bliver Muligheden fri: den frigjorte Mulighed har
sit Udtryk i Dommen.

Det Virkelige kan som Virkelighed udtrykkes baade i
Gjenstanden og i Dommen om Gjenstanden; men den i
Gjenstanden latente Mulighed kan kun udtrykkes i Mulig%
% --- p. 126 ---
\opage{126}hedens Form, forsaavidt den udtrykkes i Dommen. Vistnok
hviler den objective Mulighed altid i en tilsvarende Virkelighed;
men i Virkeligheden er Muligheden ikke til; var Muligheden
virkelig til, saa var den jo netop en Virkelighed, ingen Mulighed.
I den ydre Virkelighed, i Gjenstanden selv, maa Muligheden
nøies med en Væren uden Tilværen; i den dømmende
Subjectivitet faaer Muligheden derimod ideel Tilværelse, dens
ideelle Tilværelse er Dommen selv, den Dom, hvori den udtales. % UNSURE: period after "udtales" barely inked (gap and faint trace only)
Her sees nu Anticipationens oprindelige Eenhed med
Dommen.

Der kan i den vidende Subjectivitet danne sig mange
Billeder, ogsaa Billeder af Ting, som endnu ikke ere; men
intet Billede kan i sin Umiddelbarhed som Billede gjælde for
en Anticipation. Anticipationen fordrer udtrykkelig en Viden
om, at det givne Billede betegner en endnu ikke tilværende
Gjenstand, en Gjenstand, hvis Tilværelse netop indledes og
forberedes ved Billedet; men en saadan Viden er nadskillelig % PRINTER: printed "nadskillelig", sense "uadskillelig"
fra en Dom; thi kun i en Dom kan Forholdet mellem
Billedet og Gjenstanden udtrykkelig bestemmes. At Anticipationens
Billede ikke bliver staaende i billedlig Umiddelbarhed,
men opløser sig i en Forestilling og gaaer over i Dommen,
sees deraf, at den ved Billedet antydede Fordring, at virkeliggjøres
som Gjenstand, i Dommen udtrykkelig sættes som en
Opgave. Det er ikke Billedet, men Opgaven, den ved Dommen
om Billedet bestemte Opgave, hvorom de anticiperende Functioner
væsentlig bevæge sig. Forvandlet ved Dommen til
Opgave bliver Anticipationsbilledet nærmere bestemt ved Kategoriens
prægnante Eenhedsudtryk, altsaa ved Ordet, det praktiske
Vidensord, der bærer den hele Opgave.

Med Opgaven i det praktiske Vidensord arbeide de anticiperende
Functioner nu paa en nærmere Bestemmelse af
den paagjældende Gjenstands Mulighed. Ved Bestemmelsen
af denne Mulighed articuleres det Subjective Punkt for Punkt
af objective Hensyn. Articulationen udtrykkes i Dommen,
% --- p. 127 ---
\opage{127}idet Dommen igjen udtrykker et eiendommeligt Vexelforhold
af Frihed og Nødvendighed, Vilkaarlighed og Vilkaar. Er
Opgaven, kategorisk bestemt, at construere f.\ Ex.\ et Uhr, da
er Valget, hvorved det just blev denne og ikke en anden Opgave,
jo vistnok vilkaarligt; men Vilkaaret mælder sig strax,
naar der sees paa Opgavens Løsning. Grundfordringen er
at tilveiebringe en Række af isochrone Begivenheder og angive
disse i Tal. Af den saaledes stillede Opgave sees, at dens
Løsning beroer paa Virkelighedens Vilkaar, paa de Betingelser,
Objectiviteten tilbyder. Uhrets Billede staaer da heller ikke
fast; det er endnu kun en Mulighed, og Muligheden vexler
i Alternationen. „Et Uhr er (kan være) baade et Sanduhr
og et Vanduhr og et mechanisk Uhr, o.\ s.\ v. % PRINTER: opening „ before "Et Uhr" has no closing mark
Den conjunctive
Dom: \textit{A} er (kan være) baade \textit{B} og \textit{C} og \textit{D}, stiller
flere Muligheder paa Valg; den disjunctive Dom: \textit{A} er (maa
være) enten \textit{B} eller \textit{C} eller \textit{D}, fordrer udtrykkelig et Valg
imellem flere Muligheder.

Idet den anticiperende Subjectivitet i Dommen overskuer
de af Objectiviteten tilbudte Muligheder og besinder sig paa
et Valg, maa den ligeledes i Dommen see, hvorledes den
bedst afpasser sit Valg efter Virkelighedens Betingelser.
Dommen antager nu den hypothetiske Form: dersom Uhret
skal være et Sanduhr, saa \dots\ dersom \textit{A} er \textit{B}, hvad saa?
Saa er, svarer den formelle Logik, \textit{A} hverken \textit{C} eller \textit{D}.
Dette Svar er i nærværende Sammenhæng ørkesløst; det
Ørkesløse i den theoretiske Formalisme er allerføleligst, naar
Spørgsmaalet gjælder Løsningen af en bestemt Opgave; den
bestemte Opgave er væsentlig praktisk. Altsaa: dersom \textit{A} er
(skal være) \textit{B}, hvad saa? saa gjælder et eget System af Betingelser,
Systemet \textit{b}; dersom \textit{A} er (skal være) \textit{C}, saa gjælder
derimod et andet System af Betingelser, Systemet \textit{c},
o.\ s.\ fr.

Til Dommen, den afgjørende Dom, der gaaer igjennem
et System af Domme, svarer den anticiperede Mulighed, der
% --- p. 128 ---
\opage{128}med sin Bestemthed bestaaer i et System af Muligheder.
For at bestemme Muligheden arbeider den anticiperende
Subjectivitet med lutter imaginære Elementer, men i ethvert
af disse imaginære Elementer sees en Reflex af det Reelle,
af Realiteten i de objective Betingelser. De objective Betingelser
i deres Realitet ere lutter Magtbestemmelser; den i
Anticipationens Overveielser fordybede Subjectivitet har altsaa
gjennem disse Overveielser tilegnet sig et System af reflecterede
Magtbestemmelser; med dette System af reflecterede
Magtbestemmelser articuleres Muligheden til en indre,
imaginær Virkelighed; med den fuldstændige Udvikling af den
indre, imaginære Virkelighed er Præformationen, og med den
Anticipationen, gjennemført.

2) \emph{Med den antilogiske Dom realiseres Gjenstanden:
executive Functioner.} I den ved Anticipationen
bestemte Opgave er Virkeligheden forud sat som
Begreb; ved Opgavens Løsning skal Begrebet sættes som
Virkelighed. Hvorledes et Begreb omsættes i Virkelighed,
sees overalt, hvor et Subjectivitetens Værk objectivt fuldbyrdes,
hvor et Formaal realiseres. Det Væsentlige i Begrebets
Virkeliggjørelse bliver imidlertid misforstaaet fra to
lige eensidige Synspunkter, af Hverdagsbevidstheden og af den
immanente Methodes Philosophie.

For Hverdagsbevidstheden er det de menneskelige Fabrikater,
ikke Naturprodukterne, der maae tjene til Exempler paa
Begrebers Virkeliggjørelse i materielle Gjenstande, Huse,
Huusgeraad, Maskiner o.\ dsl. Det Orienterende ved disse
Exempler er, at de udtrykkelig henpege paa en anticiperende
Subjectivitet, en virkelig Omverden, hvori de til Formaalet
svarende Midler maae søges, og en virksom Individualitet, som
ved vilkaarlig Brug af sine Lemmer skaffer sig Værktøi og
ved vilkaarlig Brug af Værktøiet bearbeider sit Materiale
paa en til Opgaven svarende Maade, o.\ s.\ v.

% --- p. 129 ---
\opage{129}Ved slige Exemplificationer bliver nu Opgaven udstykket i
den Grad, at man, saa at sige, anatomisk kan adskille de
enkelte Momenter. Først har Værkmesteren hele Begrebet af
det paagjældende Værk anticiperet i Formaalet; dernæst bliver,
formedelst Tankens Eenhed med Villien, Udførelsen af den i
Formaalet liggende Plan til virkeligt Forsæt; i Villiens
umiddelbare Indvirkning paa Organerne, i den vilkaarlige
Brug af Arm og Haand, sees endelig Identiteten af det Ideelle
med det Reelle, det Subjective med det Objective, Begrebet
med Virkeligheden. Idet Villien formedelst Tanken bestemmes
af Begrebet, bestemmes Bevægelsen i Haand og Fingre igjen
af Villien og formedelst Villien af Begrebet; hver Fingerbevægelse
svarer nu til en Tankebevægelse; Begrebets ideelle
Magt har forvandlet sig til reel Magt: Overgangen fra Begrebsform
til Gjenstandsform, fra det Logiske til det Antilogiske,
er begyndt.

Men saa oplysende slige Exemplificationer end i propædeutisk
Forstand kunne være, saa ere de dog for Ideernes
Logik utilstrækkelige. Saalænge Opmærksomheden eensidig
dvæler ved den individuelle Værkmester og hans Værk, vil her
altid være en Mængde Biforestillinger, der kunne virke forstyrrende.
Det seer bestandig ud, som om Planen og Subjectiviteten
stode i et blot tilfældigt Forhold; vistnok maa
Planen være i Værkmesterens Bevidsthed, men Værkmesterens
Bevidsthed kan jo meget vel existere uden en saadan Plan.
Ligesaa vilkaarligt og just derfor tilfældigt er Forholdet mellem
Anticipation og Virkeliggjørelse, mellem Plan og Udførelse;
vistnok kan Værkmesteren ikke skride til Udførelse uden Plan,
men han kan jo gaae omkring med Planen uden at skride til
Udførelse. Endelig, og dette er i Hovedsagen det meest Vildledende,
bliver ved slige Exempler Coincidenspunktet af det
Subjective med det Objective i Villiens Brug af Lemmerne
opfattet paa en saa indskrænket og particulær Maade, at
Grundforholdet mellem Begreb og Virkelighed, mellem For%
% --- p. 130 ---
\opage{130}maal, Midler og Formaalets Opnaaelse, det Grundforhold,
der gaaer gjennem hele Tilværelsen, ikke bliver oplyst: her
er og bliver en viid Kløft mellem \emph{Fabrikater} og \emph{Naturprodukter}.

I skarp Modsætning til de populære Exemplificationer,
har Hegel i sin Lære om Begrebets Virkeliggjørelse ved sig
selv, navnlig i sin Lære om Formaalet, i den Grad fjernet
alle Endelighedshensyn, at Modsætningen mellem det subjective
og objective Formaal, mellem Begrebet i sig og Objecternes
virkelige Verden, falder ind under Begrebets almindelige
Selvbevægelse. Mystificationen som i „\emph{Subjectivitetens
Logik}“\footnote{Grundideernes Logik I. S. 226.} % note *)
allerede er paaviist, kan i denne Sammenhæng
belyses fra en ny Side. --- „Die Bewegung des Zwecks“ ---
siger Hegel\footnote{Wissenschaft der Logik. Zweiter Theil. S. 220.} % note **)
--- „kann \dots\ so ausgedrückt werden, daß sie
darauf gehe, seine \emph{Voraussetzung} aufzuheben, das ist die
Unmittelbarkeit des Objekts, und es zu \emph{setzen} als durch den
Begriff bestimmt. Dieses negative Verhalten gegen das Objekt
ist ebenso sehr ein negatives gegen sich selbst, ein Aufheben
der Subjektivität des Zwecks. Positiv ist es die Realisation
des Zwecks, nämlich die Vereinigung des objektiven
Seyns mit demselben, so daß dasselbe, welches als Moment
des Zwecks unmittelbar die mit ihm identische Bestimmtheit
ist, als \emph{äußerliche} sey, und umgekehrt des Objective als % PRINTER: printed "des Objective", sense (Hegel) "das Objective"
\emph{Voraussetzung} vielmehr als durch Begriff bestimmt, \emph{gesetzt} werde.“

Ifølge den hegelske Opfatning forsvinder da al Forskjel
mellem den anticiperende Subjectivitet og Planen, det subjective,
ved Anticipationen bestemte Formaal. Forsaavidt Formaalet
forudsætter Anticipation, og Anticipationen Subjectivitet,
er Formaalsbegrebet hos Hegel denne Subjectivitet. Begrebet
har som „det bestemte Begreb“ alle de Requisiter, det
% --- p. 131 ---
\opage{131}væsentlig trænger til. Her er forsaavidt en tilsnegen Anticipation,
men intet Anticipationsproblem. „Maalets Bevægelse
gaaer ud paa at ophæve sin Forudsætning, nemlig Objectets
Umiddelbarhed, at Objectet udtrykkelig kan blive bestemt ved
Begrebet;“ hvad betyder nu dette? Objectet er som Object
allerede Begreb; idet Maalet som Begreb har Objectets
Umiddelbarhed til Forudsætning, vil dette med andre Ord
sige, at det subjective, i sig reflecterede Begreb har det objective,
umiddelbare Begreb til Forudsætning. Lægges nu Eftertrykket
derpaa, at det reflecterede Begrebs Forudsætning selv
er Begreb, saa er den hele Bevægelse, hvorved den af Begrebet
resulterende Umiddelbarhed igjen ophæves i sit reflecterede
Begreb, kun en Gjentagelse af den forslidte Historie
med Reflexionens Omsætning i Umiddelbarhed, og Umiddelbarhedens
i Reflexion; her kan da ikke være Tale om, at
Maalet skulde realisere sig, thi her er slet ingen Realitet, og
selvfølgelig heller ingen Realisering. Skal Eftertrykket derimod
falde paa Objectets objective Umiddelbarhed, altsaa paa
Realiteten, da spørges, hvorfra Maalet som Begreb faaer den
Selvstændighed, at det i fri Idealitet kan træde op imod en
given Realitet, gribe ind i Tingenes Sammenhæng og magte
det udvortes Objective. Begrebet som Maal er en subjectiv
Tanke, Abstractionen har løsrevet fra den tænkende Subjectivitet,
og stillet op som et Væsen for sig. Det saaledes
opstillede Væsen, „det subjective Maal,“ er et Phantom, en
magtesløs Skygge, der vel skal lade være at realisere sig selv.

Det hegelske Begreb mangler i al dets Guddommelighed
aabenbart den væsentlige Modsætning af ideel Tilværelse i
Dommen og reel Tilværelse i Gjenstanden. Begrebets Selvrealisering
er en fingeret Bevægelse, der skyldes Abstractionen
og Abstractionens trinvise Ophævelse. At „det subjective Formaal
vender sig negativt mod Objectet,“ at det med det Samme
„vender sig negativt mod sig selv,“ at det, med andre Ord,
„ophæver sin Subjectivitet og sætter sig negativt som realiseret
% --- p. 132 ---
\opage{132}Formaal,“ vil kun sige, at den tænkende Tilskuer skal ved
Hjælp af en explicativ Dialektik fæste Blikket snart paa det
ene af disse Momenter, snart paa det andet, snart paa dem
begge tilsammen.

Er Hegel med hele sit System ikke istand til at angive
Forskjellen mellem en Gjenstands Begreb og Gjenstanden selv,
mellem Begrebet i Viden og Begrebet i Magt, hvorledes
skulde han da være istand til at paavise en Overgang fra
Begreb til Gjenstand? Formelt er her Overgange nok; men
her er intet Begreb af reel Overgang. En reel Overgang
forudsætter nemlig en reel Afstand; hvor der ingen reel Afstand
er, er heller ingen, i logisk Forstand, reel Overgang.
Overgangsillusionen, en Illusion ikke ulig den, at den Seilende
seer Kysten bevæge sig, vedligeholdes fornemmelig derved,
at Begrebet i det \emph{Almindelige} snart træder i Modsætning
til, snart igjen smelter sammen med Begrebet i det \emph{Særskilte};
men Begreb er dog altid Begreb: den dialektiserende Reflexion
sætter og hæver Adskillelsen med lige Lethed. Er her
først i Conseqvens af Begrebet i det Særskilte sat Adskillelse
mellem det Objective og det Subjective, som nu Tilfældet er
med Objecterne og „det subjective Formaal“, saa bliver en
saadan Adskillelse igjen hævet ved Begrebet i al Almindelighed.
Om Begrebet i al Almindelighed kan det jo vistnok gjælde,
at det vender sig negativt mod sig selv, idet det vender sig
negativt mod Objectet; thi Begrebet af det Objective hører
unegtelig ligesaa væsentlig under Begrebet i al Almindelighed
som Begrebet af det Subjective; men hvorledes skal Sligt
kunne løse et Realiseringsproblem, hvorledes skal det kunne
forklare os Overgangen fra „det subjective“ til „det realiserede
Formaal“?

Hvad vi i denne Sammenhæng trænge til at faae forklaret,
er jo ikke en Tautologie som den, at det subjective Formaal
i sit Begreb er subjectivt og derved modsat Objectet og
det Objective; nei, det, der skulde forklares os, er det Problem,
% --- p. 133 ---
\opage{133}hvorledes det i Virkeligheden gaaer til, at Begrebet, ikke det
subjectiv-objective Begreb i logisk Tvetydighed, men det bestemte
Begreb, Begrebet af dette særlige, bestemte Formaal kan
bevirke en Overgang fra ideel Tilværelse til virkelig Existens,
idet her da ved Overgangen ikke tænkes paa en logisk Omskrivning
af det Subjective ved det Objective, men paa en
Begivenhed, en Frembringelse, hvori Begrebet sees at bestemme
Virkeligheden og selv blive virkeligt. Til Besvarelse af herhen
hørende Spørgsmaal indeholder den hegelske Logik ingen
Anviisning, end ikke et Vink.

Hos Hegel kan Modsætningen mellem Begreb og Gjenstand,
naar den drives til sit Yderste, ikke naae videre end til
en Modsætning af Væsen og Phænomen, kun med den udtrykkelige
Tilføining, at Totaliteten fra begge Sider er fuldstændig
udviklet, ja under Udviklingen endog bleven til to
Totaliteter. Da nu Totalitetseenheden af de to Totaliteter
sættes som Dommen, den logisk objective Dom (Grunddelingen,
„das Urtheil“), saa skulde man vente, at Hegel i Dommen
fastholdt Modsætningen mellem Begrebet i Væsenets og Begrebet
i Phænomenets Form, og saaledes kom til en Erkjendelse
af Gjenstanden; men nei! i Dommen reflecteres der
paa Udfyldning af Copula, altsaa paa Totalisering af den
Identitet, hvori de særskilte Totaliteter ophæves. „So ist die
Form des Urtheils untergegangen, erstens, weil Subjekt und
Prädikat \emph{an sich} derselbe Inhalt sind; aber zweitens, weil das
Subjekt durch seine Bestimmtheit über sich hinausweist, und
sich auf das Prädikat bezieht, aber ebenso drittens ist \emph{dieß
Beziehen} in das Prädikat übergegangen, macht nur dessen Inhalt
aus, und ist so die \emph{gesetzte} Beziehung oder das Urtheil
selbst. --- So ist die konkrete Identität des Begriffs \dots\ im
\emph{Ganzen} hergestellt“ \dots\footnote{Wissenschaft der Logik. Zweiter Theil. S. 117.} % note *)

Med den udfyldte Copula gaaer Dommen nu over i
% --- p. 134 ---
\opage{134}Slutningen.\footnote{Jvfr. Grundideernes Logik. I. S. 214.} % note *)
Endte Dommen i sin immanente Udvikling
med, at de to Totaliteter, Subjectet og Prædicatet, gik sammen
i Begrebets totaliserede Eenhed, saa maatte man jo nu vente,
at den indholdsrige Totalitetseenhed, som under Slutningens
Selvudvikling skilles og enes i tre Totaliteter, skulde lede til
Indsigt i Modsætningen mellem Begrebet som Begreb og Begrebet
som Gjenstand; men nei! Slutningen ender saaledes:
„Die ganze Formbestimmung des Begriffs ist in ihrem bestimmten
Unterschied und zugleich in der einfachen Identität
des Begriffes gesetzt.“ Og nu tilføies: „Dadurch hat sich nun
\emph{der Formalismus des Schließens}, hiermit die Subjektivität
des Schlusses und des Begriffes überhaupt aufgehoben“.\footnote{Wissenschaft der Logik. Zweiter Theil. S. 169 \& 71.} % note **)

Vil Nogen mene, at her dog omsider maa vise sig en
gjennemgaaende Totalitetsmodsætning, siden det udtrykkelig
hedder, at Begrebets hele Formbestemmelse er sat, ikke blot i
eenfold Identitet, men ogsaa i bestemt Forskjel, da lægge han
Mærke til, hvorledes Kapitlet om Slutningen ender: „Das
Resultat ist \dots\ eine \emph{Unmittelbarkeit}, die durch \emph{Aufheben
der Vermittelung} hervorgegangen, ein Seyn, das ebenso
sehr identisch mit der Vermittelung und der Begriff ist, der
aus und in seinem Andersseyn sich selbst hergestellt hat. Dieß
Seyn ist daher eine \emph{Sache, die an und für sich ist} --- die
\emph{Objektivität}“. Og hvad er saa herved vundet? Et Begreb,
der selv er Gjenstand, fordi det er Object, men et Object,
der dog ikke er Gjenstand, fordi det er Begreb; altsaa
et Begreb, der hverken er Begreb eller Gjenstand; en Gjenstand,
der hverken er Gjenstand eller Begreb; en dialektisk
Historie, der kan blive ved, saalænge det skal være.

Skal Gjenstanden, seet fra Logikens Side, hævde sin
Selvstændighed imod Begrebet; skal Begrebet, seet fra Existensens
Side, hævde sin Selvstændighed imod Gjenstanden:
% --- p. 135 ---
\opage{135}da maae de begge have hver sin Tilværelse, Begrebet en ideel
Tilværelse i en begribende Subjectivitet, Gjenstanden en reel
Tilværelse i virkelig Objectivitet. Begrebets Tilværelse i
Subjectiviteten er bestemt ved et Vidensord; Gjenstandens
objective Tilværelse er bestemt ved et Magtord. Begrebets
Realisation, den virkelige Overgang, hvorved det bestemte Begreb
kommer til Existens som bestemt Gjenstand, skeer ved en
intellectuel-reel Praxis, hvis eiendommelige Charakteer maa
beroe paa de executive Functioner.

De executive Functioner bestemmes selv ved Forskjelseenhed
af antilogiske Domme og antilogiske Sætninger. Den
antilogiske Dom, der har Begrebet i sit Subject og Existensen
i sit Prædicat\footnote{Grundideernes Logik. I. S. 160--62.}, % note *)
faaer først sin fulde Betydning ved Tankens
Overgang i Handling, Theoriens Overgang i Praxis.

Fra Theoriens Synspunkt vil den antilogiske Dom kun
tage sig ud som en bagvendt Fremstilling af den logiske.
Thi naar Nogen, istedenfor at sige: dette er en Axe, dette er
et Hjul, dette en Valse, siger: en Axe er dette, et Hjul er
dette o.\ s.\ v., lyder det jo bagvendt nok. Anderledes, naar
Opgaven er at construere et Hjul, en Valse o.\ dsl. Constructionen
forudsætter Begrebet: Begrebet er det Oprindelige,
Gjenstanden det Afledede. Begrebet haves i Dommens Subject;
den Dømmende veed, hvad et Hjul er i sit Begreb, han
vil altsaa ikke have Begrebet igjen bestemt som Begreb i Prædicatet,
saa at Prædicatet skulde give ham en logisk exact
Definition. Slige theoretiske Formaliteter ere her overflødige.
Vistnok vil den construerende Værkmester definere sit Begreb,
men Definitionen skal være praktisk, ikke theoretisk: den praktiske
Definition er Eet med en tilsvarende Frembringelse, en
Construction.

Det hedder i Logiken, at det ikke er det Enkelte, men
det Særskilte, ikke Gjenstanden selv, men dens Begreb, der
% --- p. 136 ---
\opage{136}lader sig definere (\textit{species sola definitur, non individuum});
dette er nu i theoretisk Forstand ganske rigtigt; den logiske
Dom kan ikke definere Gjenstanden. Men i Praxis forholder
Sagen sig omvendt; her gaaer Bestræbelsen netop ud paa at
definere Begrebets Realitet ved at frembringe den til Begrebet
svarende Gjenstand: \emph{den frembragte Gjenstand} er
\emph{Begrebets praktiske Definition}. Paa saadanne Spørgsmaal
som: hvad er en Trekant? et Qvadrat? en Cirkel? kan
En svare praktisk ved at construere en Trekant, et Qvadrat,
en Cirkel. Hvorvidt slige praktiske Definitioner ville tilfredsstille
Theorien, skal her ikke undersøges; kun dette skal fremhæves,
at den praktiske Definition ledes og bestemmes ved en
antilogisk Dom.

Den, der besvarer Spørgsmaalet: hvad er en Cirkel?
ved at construere en Cirkel, oversætter ligefrem Cirklens Begreb
i Existens. „En Cirkel er en Figur som \emph{denne}, er
\emph{denne Figur}“; med en saadan Dom begynder Constructionen,
men en saadan Dom er antilogisk. Den logiske Dom
falder umiddelbart sammen med sin tilsvarende Sætning; i
Vidensordet er Sætningen Dom, forsaavidt den reflecteres i
Tanken, og Dommen Sætning, forsaavidt den reflecteres i
Sproget: begge Reflexioner ere theoretisk den samme Reflexion.
Anderledes, naar Spørgsmaalet er om Forholdet
mellem Dom og Sætning i det Antilogiske; her er Identiteten
praktisk, og derfor maa her ogsaa være en praktisk
Forskjel.

Den praktiske Forskjel er en Forskjel mellem Tanken,
der leder Handlingen, og Handlingen selv. Den construerende,
altsaa handlende, Subjectivitet maa \emph{tænke}; men Tanken er
praktisk, dens Functioner executive. De executive Functioner
gaae fra Begrebets Synspunkt allesammen ud paa at
vise de praktiske Definitioners Mulighed. Her adskiller den
antilogiske Dom sig i en Mangfoldighed af særskilte Domme.
Den construerende, altsaa tænkende Subjectivitet maa \emph{handle}.
% --- p. 137 ---
\opage{137}Den virkelige, ved Tanken bestemte, Handlen er ikke selv en
Tænken, men tvertimod Tanken saaledes modsat, at den netop
ved Modsætningen kan svare til Tanken. Den virkelige, til
Tanken svarende, Handling kan, da den ikke selv er Tanke,
umulig falde sammen med Dommen som Dom; den maa
nødvendigt have sit Udtryk i en ved Dommen bestemt antilogisk
Sætning. Idet den antilogiske Dom forgrener sig i
Mangfoldigheden af særskilte Domme, forgrener den antilogiske
Sætning sig under Værkets Udførelse i en Mangfoldighed af
særskilte Sætninger; i den articulerede Vexelbestemmelse af
antilogiske Domme med praktisk-antilogiske Sætninger sees den
fulde Betydning af de executive Functioners Logik.

At den i og ved Dommen, den antilogiske Dom i sit
System af tilsvarende Sætninger, frembragte Gjenstand nu
ogsaa virkelig er Dommen modsat, kan sees paa de executive
Functioner. I enhver executiv Function er den for Sætningen
afgjørende Dom et Udtryk for Identiteten af Magt
og Viden i Videns Form: Magten i Dommen er ideel.
Men i enhver executiv Function er den ved Dommen bestemte
Sætning, den af Tanken styrede Handling, tillige et
Udtryk for Identiteten af Viden og Magt i Magtens Form;
Magten i Handlingen er reel: saamange Handlinger, saamange
Magtbestemmelser. Da nu Sætningens Eenhed ved Dommen
Punkt for Punkt giver en Magtbestemmelsens Eenhed med sin
tilsvarende Begrebsbestemmelse, saa maa i Udførelsen selv
Totaliteten af Magtbestemmelserne i den frembragte Gjenstand
nødvendig svare til Totaliteten af Vidensbestemmelserne i den
frembringende Subjectivitet. Gjenstanden i sit System af
antilogiske Sætninger staaer som et Magtens Ord imod det
System af Domme, der i det praktiske Vidensord har givet
Begrebet Tilværelse.

Er Modsætningseenheden af Tænken og Handlen, af
Dom og Sætning, en Modsætningseenhed af Viden og Magt,
saa sees heraf, at den frembringende Subjectivitet maa være
% --- p. 138 ---
\opage{138}paa een Gang baade en formaaende og en vidende Subjectivitet.
Er Magten indskrænket, men Viden udvidet, vil det
sees paa Frembringelsen; er omvendt Magten udvidet, men
Viden indskrænket, vil det ligeledes sees paa Frembringelsen.
Her er et Vendepunkt, hvor Forholdet mellem \emph{Fabrikater}
og \emph{Naturprodukter} nærmere kan belyses.

Alle reale Frembringelser, selv de tilsyneladende meningsløse,
bestemmes ved Begreber, og virkeliggjøre Begreber. At
antage, at der i Virkeligheden skulde kunne gives Frembringelser
saa forvirrede, at de faldt udenfor enhver Tanke, enhver
Begrebsbestemmelse, er ligesaa urimeligt, som at antage, at
der i Sand, Vand eller Luft skulde gives Hvirvler, der faldt
udenfor alle Love. Er det nu sandt, at enhver Frembringelse
beroer paa Modsætningseenhed af Begreb og Virkelighed, og
er Virkeligheden Frembringelsens objective, Begrebet dens subjective
Side, saa indsees jo heraf, at Frembringelsen er en
tvesidig Kategorie, og at det gjør en stor Forskjel, om denne
Kategorie betragtes fra et subjectivt eller fra et objectivt
Standpunkt.

„Fabrikaterne, de menneskelige Frembringelser, ere“, siger
man, „Frembringelser, der skee med Bevidsthed, men Naturprodukterne
ere ubevidste Frembringelser“. En saadan Adskillelse
af Fabrikater og Naturprodukter er vildledende, forsaavidt
den giver det Skin af, at her skulde være to i Principet
forskjellige Frembringelser, subjective og objective, medens
Sandheden er, at enhver Frembringelse er i sidste Instans
baade subjectiv og objectiv.

At de menneskelige, bevidste Frembringelser ere subjective,
falder i Øinene; men de ere dog ligefuldt, med alt det Vilkaarlige,
der udmærker dem, objectivt betingede. Den, der
vil abstrahere fra Værkmesterens Subjectivitet og under Arbeidet
fæste sin Opmærksomhed paa Haanden, Værktøiet og
Materialet, vil overbevise sig om, at det i Frembringelsen
Subjective Led for Led bestemmes objectivt. At Natur%
% --- p. 139 ---
\opage{139}frembringelserne ere objective, falder i Øinene; men den, der vil
lægge Mærke til, hvad de Betragtende sige om „Naturformer“,
„Naturbegreber“, „Naturtanker“, „Naturformaal“ o. s. v.,
vil snart indsee, at ingen Tænkende kan betragte Naturfrembringelsernes
objective Side, uden at see sig om efter den
frembringende Subjectivitet.

„Men“, siger man, „Naturfrembringelserne foregaae jo
med objectiv Nødvendighed; hvorledes kan da en Frembringelse,
der skeer saa objectivt, beroe paa Subjectivitet?“ Vi
svare: Ikke paa en endelig psychologisk Subjectivitet, som med
en begrændset Viden og en endnu mere begrændset Magt befinder
sig i endelig Modsætning til Objecternes Verden; men
vel paa en Subjectivitet, som med oprindelig Eenhed af
Viden og Magt er i et uendelig articuleret Grundforhold til
Objectiviteten i alt Objectivt. Skulde Naturfrembringelserne
tabe deres Realitet, fordi de i Begrebet frembringes ideelt og
paa subjectiv Maade, saa maatte Naturgjenstandene jo ogsaa
tabe deres Realitet, fordi de kun kunne være Objecter, idet
de objectiveres af en uendelig Subjectivitet. Er det nu saa
langt fra, at Objecterne blive mindre objective ved at objectiveres,
at de tvertimod maae have deres Objectivitet i Objectiveringen
selv: da vil det ogsaa kunne forstaaes, at Naturfrembringelserne
ikke tabe deres objective Gyldighed ved at
frembringes af en dem iboende og dog oversvævende Subjectivitet;
det er just det Subjective i Frembringelsen, der begrunder
det Objective.\footnote{Grundideernes Logik. I. S. 444--456.} % note *)

% Indhold (p. 0-0 fragment) numbers this head "3)"; no numeral is printed here.
\emph{I den logisk-antilogiske Dom forklares Gjenstanden:
explicative Functioner.} Det er en egen Sag
med Naturbetragtninger; Alt, hvad den aandrige Betragter
lægger ind i Naturen, faaer han ud igjen, hverken Mere eller
Mindre. Den bekjendte Strid mellem den hallerske og goetheske
Anskuelse angaaende „Naturens Indre“ har Videnskaben da
% --- p. 140 ---
\opage{140}hidtil heller ikke kunnet udjævne. Vel have Naturprodukterne
i deres umiddelbare Objectivitet ingen frigjort Inderlighed,
som kunde her være et Ydre, der liig en Skal omgav og
skjulte det Indre; men Umiddelbarheden, hvormed al Naturvirksomhed
foregaaer, giver dog Frembringelsen et saa objectivt
Præg, at det Subjective paa een Gang baade savnes og
forudsættes. Naar den ene Masse tiltrækker den anden, den
ene Steen trykker den anden, det ene Stof forbinder sig med
det andet: da er her Punkt for Punkt en Bestemmen af
Andet ved Andet, af Objectivt ved Objectivt paa objective
Muligheder, og det Samme gjælder, naar Frøet spirer, og
Organismen trives; --- hvad er saa Naturens Indre? I sin
umiddelbare Legemlighed er Organismen en Gjenstand ligesaavel
som Stenen; at her i Gjenstanden er en Modsætning af
Indre og Ydre maa indrømmes; men Gjenstandens Indre er,
hvad vi alt ovenfor have seet, ligesaa vel Gjenstand, som det
Ydre. At der i Naturens Ydre overalt skulde træde os et
Indre imøde, er en Talefigur, der kan være ligesaa tvetydig,
som naar det i kluntet Umiddelbarhed forsikkres, at der „i
Gjenstanden træder os et Begreb imøde“.

Det gaaer med Naturens Indre som med Naturtankerne;
naar vi see paa det Ydre, saa er der et Indre; men ville vi
see det Indre selv, saa er det forsvundet. Det Indre har
væsentlig en reflecteret, ikke en umiddelbar Væren; men
Naturen reflecterer ikke, og kan i sin umiddelbare Naturlighed
ikke have et væsentligt Indre. Og dog har Naturen et Indre;
naar Goethe siger: „die Natur hat weder Kern noch Schale,
Alles ist sie mit einem Male“, da lyder det slaaende, men er
dog forblommet; Digterens klare Ord trænge høilig til en
Forklaring.

At Naturen er uden et Indre, uden fri Inderlighed,
betyder, at den er uden Subjectivitet; at Naturen har sin
Inderlighed, sit væsentlige Indre, betyder, at den dog har
Subjectivitet; at Naturens Indre er dunkelt, dens Inderlig%
% --- p. 141 ---
\opage{141}hed en Gaade, betyder, at den Subjectivitet, Naturen har,
er en Subjectivitet, der har Naturen, altsaa en i sig selv fordoblet
Subjectivitet; en Subjectivitet, der naturligviis er i
Naturen selv, og en Naturen oversvævende Subjectivitet, en
Subjectivitet, hvori de frembringende Handlinger reflectere sig
som bestemmende Tanker.

Til Naturfrembringelsernes fordoblede Subjectivitet svarer
Forskjelseenheden af antilogiske Domme og praktisk-antilogiske
Sætninger. Uden et System af antilogiske Domme vilde de
for Frembringelserne afgjørende Begreber umulig kunne være.
Begreberne ere kun i en Begriben, Begriben kun i en begribende
Subjectivitet. De antilogiske Domme bestemmes i
Anticipationen ved de logiske; Reflexionen af de logiske og de
antilogiske Domme er i den vidende, Frembringelserne oversvævende,
sig selv i Alt og Alt i sig reflecterende Subjectivitet,
i den guddommelige Selvhed, hvis Viden er Magt, og hvis
Magt er Viden. Men uden et System af antilogiske Sætninger,
virkelige, ved den praktiske Videns anticiperede Begreber
udtrykkelig bestemte Magthandlinger, vilde Naturfrembringelserne
umulig kunne have Andethedens, Virkelighedens,
den reelle Tilværelses reelle Charakteer. Hvor finde vi nu
de antilogiske Sætningers oprindelige Eenhed med de tilsvarende
Domme? I den oversvævende Subjectivitets Væsens-Eenhed
% hyphen at line end kept: capital E marks a hyphenated compound
med den mægtige, Frembringelserne iboende Subjectivitet,
endnu nærmere bestemt, i denne Subjectivitets reflecterede
Eenhed med den umiddelbart frembringende Aarsag,
det i objectiv Virken værende Naturprincip.

Det actuelle Forhold mellem Naturprincipet og den
uendelige Subjectivitet maa i Logiken bestemmes ved \emph{explicative
Functioner}.

Enhver Naturgjenstand er et Frembragt. At erkjende
det Frembragte er at erkjende Frembringelsen; at erkjende
Frembringelsen er at gjentage den i Erkjendelsen;\footnote{Grundideernes Logik I. S. 169--70 flg.} % note *)
at gjen%
% --- p. 142 ---
\opage{142}tage den i Erkjendelsen er at gjennemreflectere det System af
antilogiske Domme, der have ledet de frembringende Handlinger,
og at efterligne disse Handlinger i særlige, til Begrebets
Fordringer svarende Imaginationer. Til Handlingernes,
de antilogiske Sætningers, System svarer Naturprincipet.
Principet --- dette gjælder baade om Principet i Almindelighed
og om de særskilte Principer --- virker, men tænker ikke, gjør
Begrebets Gjerninger, men begriber det ikke selv; Principet
er som Naturprincip Selvhed i objectiv Formen, Subjectivitetsanlæg,
en selvløst virkende Selvhed, de antilogiske Sætningers
productive Eenhed. Til Dommenes, de logisk-antilogiske
Dommes, System svarer Selvhed i Selvets Form,
den reflecterende, vidende Subjectivitet. Den vidende Subjectivitet
--- dette gjælder baade om den mægtige, i sin ontologiske
Viden uendelige, og om den afmægtige, i sin psychologiske
Viden indskrænkede Subjectivitet --- kan i sin rene
Viden, det ligger i Videns Natur, ikke handle, thi Viden er
ikke Handlen; den Vidende kan i sin Viden kun tænke og
dømme.

Til den gjennemgaaende Forskjel mellem Tænken og
Handlen, Viden og Magt, Theorie og Praxis, svarer altsaa
en gjennemgaaende Modsætning af den begribende Subjectivitet,
der tænker Handlingen, men ikke handler, og det virkende
Princip, der gjør Tankens Gjerninger, men ikke tænker.
En Dualisme af Subjectivitet og Naturprincip kan ikke afskrække
den logiske Forsker; thi en saadan Dualisme tjener
just til at begrunde Identiteten. Identiteten hviler i Magten.
Forsaavidt Magten efter sin Idealitet er oprindelig Eet med
Viden, tilhører den Subjectiviteten: den ontologiske Subjectivitet
er i uendelig Viden uendelig mægtig; forsaavidt Magten
i sin umiddelbare Realitet er Viden oprindelig modsat, grundlægger
den Andetheden, sætter Objectiviteten og er identisk
med Principet.

I Subjectiviteten er den objective Magt subjectiv paa
% --- p. 143 ---
\opage{143}Grund af dens Idealitet, dens Eenhed med Viden; i Principet
er den subjective Magt objectiv paa Grund af dens
Andethed, dens Realitet, dens Eenhed med Virken. Her er
da ikke to Magter, en ideel og subjectiv ved Siden af en
reel og objectiv; det er den ene og samme Magt, der i mægtig
Selvmodsætning er ideel i Tænken og Dømmen, reel i Virken
% UNSURE: Dømmen or Dommen (no ø-stroke visible in scan; layer and sense Dømmen)
og Frembringen.

I sit System af antilogiske Domme magter Subjectiviteten
altsaa Principet med alle derfra udgaaende Virksomheder
ideelt. I Kraft af den ideelle Magt bliver hver Naturhandling
Led for Led bestemt af den tilsvarende Dom. Principet
gjør i alle Magtyttringer Tankens Gjerning. I sit System
af antilogiske Sætninger har Principet de i Dommen reflecterede
Begrebsbestemmelser som reelle Magtmotiver, saa at
Frembringelsen netop ved Gjenstandens Modsætning til Begrebet
gjør Begrebet Fyldest.

Enhver Naturfrembringelse, den organiske ligesaavel som
den chemiske og mechaniske, skeer altsaa objectivt ifølge det i
Begrebet sig bevægende Princip, men bestemmes subjectivt af det
i Subjectiviteten reflecterede Begreb. Det Objective i Frembringelsen
er altsaa ikke uden videre subjectivt, men faaer sin
Betydning ved det Subjective; det Subjective i Frembringelsen
er ikke selv objectivt, men faaer sin Betydning ved det Objective.
At det Objective punktlig forklares af det Subjective,
og dette af hiint, at det Objective kun er objectivt ved at betyde
det Subjective, og omvendt, er et theoretisk Moment i
det praktiske Grundforhold, der giver de explicative Functioner
deres eiendommelige Præg.

De explicative, d. e. forklarende, bedømmende, udtydende
Functioner forudsætte netop en Dualisme af Subjectivt og
Objectivt, af Indre og Ydre, Begreb og Gjenstand. Gjenstanden
er som Gjenstand ikke Begreb, men betyder Begrebet,
og skal derfor udtydes og forstaaes ved Begrebet. Begrebet
er Frembringelsen iboende (immanent) og dog oversvævende
% --- p. 144 ---
\opage{144}(transcendent), iboende som et Magtens Ord, oversvævende
som et Videns Ord; paa det sande Forhold mellem det
iboende og oversvævende Begreb beroer Udtydningens Sandhed.

I al Naturvirksomhed er Begrebet immanent; men derfor
kan det dog som Begreb umulig falde sammen med Virksomheden
selv. Gik Begrebet saaledes op i Virksomheden, at
vi i det Virksomme kun fik Virksomheden selv, da blev Virksomheden
begrebløs, altsaa meningsløs. For at det Naturvirksomheden
iboende Begreb kan være Begreb, maa det nødvendig
være oversvævende; for at bevæge sig i Realsystemet
som et Magtens Ord, maa det reflecteres i Idealsystemet som
et Videns Ord.

Der er en væsentlig Forskjel imellem at \emph{være} Noget
og at \emph{betyde} Noget; en Gjenstands Betydning kan, forsaavidt
den vilkaarlig tillægges Gjenstanden, være dens Tilværelse
aldeles ligegyldig; men en Gjenstands Betydning, er,
% the comma after "Betydning" is printed (slightly raised); text layer omits it
naar den netop svarer til Grund og Væsen, afgjørende for
Gjenstandens Væren.

Den Betydning, som ved de explicative Functioner tilkjendes
Naturgjenstanden, er Eet med Gjenstandens eget
Væsen. I saadanne Domme, som: dette er en Steen, en
Plante, et Dyr, udtaler den Dommende, at han ved enhver
% UNSURE: Dommende or Dømmende (ø-stroke not visible in scan; layer Dommende)
af de modende Gjenstande kjender Begrebet: Dommene ere
% UNSURE: modende or mødende (ø-stroke not visible in scan; layer modende)
logiske. Ifølge den logiske Dom er Gjenstanden sit Begreb,
dens Betydning angives umiddelbart i Dommen som Dom;
her kan forsaavidt ikke være Tale om nogen Forskjel imellem
at betyde og at være: hvad Gjenstanden er, betyder den, og
hvad den betyder, det er den. Anderledes, naar Spørgsmaalet
er om Naturgjenstandens Oprindelse, naar det Frembragte skal
forstaaes af Frembringelsen. Her vender Forholdet sig om:
istedenfor at begynde med Gjenstanden og fatte den theoretisk
i Begrebet, gaaer den dømmende Subjectivitet her ud fra
Begrebet, og vil praktisk forstaae, hvorledes et givet Begreb
er blevet realiseret i den bestemte Gjenstand. „Stenen er et
% --- p. 145 ---
\opage{145}Naturprodukt som dette“; „Planten er et Naturprodukt som
dette“; slige antilogiske Domme forudsætte, at den Dommende
% UNSURE: Dommende or Dømmende (ø-stroke not visible in scan; layer Dommende)
oprindelig er vis paa Begrebets objective Gyldighed, og
Dommen afgjør nu i hvert enkelt Tilfælde, hvorledes det
paagjældende Begreb har realiseret sig i den givne Gjenstand.
Fra dette Synspunkt kan her være Tale om, at Gjenstanden
repræsenterer sit Begreb, betyder Begrebet, faaer sin Betydning
og dermed sin Udtydning fra Begrebet. „Er denne Dannelse
et Blad, eller er det maaskee en Green?“ her er i Botaniken
en Sag, der skal nærmere undersøges. Hvad gjør nu Botanikeren?
Han stiller Begrebet Green og den Maade, hvorpaa
dette Begreb morphologisk realiserer sig, skarpt imod Begrebet
Blad og Bladbegrebets morphologiske Fremtræden, og
Dommen falder saa, at dette f. Ex. „maa være en Green!“
Her sees det udtrykkelig, at Gjenstanden betyder sit Begreb,
og faaer sin Udtydning og Betydning fra Begrebet.

Vexelbestemmelsen af Dom og Gjenstand har viist sig at
være baade logisk og antilogisk. I den logiske Dom er Bestemmelsen
anticiperende, i den antilogiske er den realiserende,
i den logisk-antilogiske er den explicerende. „Men“ --- vil
man spørge --- „hvorledes kan en Dom paa een Gang være
baade logisk og antilogisk? Den logiske Dom har sin Form,
den antilogiske har sin; men har Logiken nogen Form for en
Dom, der baade er logisk og antilogisk?“ Vi svare: en særlig
Form for den logisk-antilogiske Dom har Logiken jo vistnok
ikke, forsaavidt man ved særlig Form udelukkende vil
tænke paa en særlig Sætning. Men paa den særlige Sætning
lægges her ingen Vægt; den hele Vægt maa falde paa
den særlige Tanke. Den særlige Tanke udtaler sig i alle Domme,
der udtrykkelig betegne, at den Dømmende i den frembragte Gjenstand
gjenkjender Frembringelsens Begreb. Slige Domme give
sig tilkjende under meget forskjellige Vendinger: „dette maa
være en Green“; „saaledes maa en Bygning være“; „dette er et
% --- p. 146 ---
\opage{146}maadeligt Exemplar, et fortrinligt Exemplar af“ „en Abnormitet
af“ o. s. v.

Reflecteres nu særligt paa det i Frembringelsen immanente
Begreb, saa bliver Andet Punkt for Punkt bestemt ved
Andet, Objectivt ved Objectivt paa objective Muligheder;
igjennem de objective Muligheder og deres betingede Virkeliggjørelse
seer Betragteren Gjenskin af de Tanker, der lede
Principets Virksomhed; det Indre gaaer op i det Ydre, og vi
maae i objectiv Forstand sande Digterens Ord: „die Natur
hat weder Kern noch Schale“. Reflecteres derimod paa Forskjellen
mellem Begrebet og Virksomheden, saa møder den
menneskelige Viden en dobbelt Skranke. For det Første maa
den erkjende det for umuligt at forfølge Reflexionen af det
immanente Naturbegreb i Transcendensens Overgangsformer
gjennem alle Mellembestemmelser ind i Idealsystemets uendelige
Dyb. For det Andet maa den erkjende det for umuligt
at eftergjøre Frembringelser, der forudsætte skabende Eenhed
af Viden og Magt, af antilogiske Domme og praktisk-antilogiske
Sætninger.

Hvorledes Begrebet om en Trekant leder Constructionen af
en Trekant, kan blive os klart, thi vi kunne baade tænke Begrebet
og construere Trekanten; men hvorledes Organismens Begreb kan
bevæge det organiske Princip, og det af Begrebet bevægede Princip
de materielle Dele, vil paa den menneskelige Erkjendelses nuværende
Vilkaar aldrig blive gjennemsigtig: vi kunne indsee \emph{at}, men
ikke \emph{hvorledes}. Enten vi altsaa betragte Begrebets Immanens
eller dets Transcendens, saa har Naturen i begge Henseender
sit Mysterium, sit Indre, et Indre, der taber sig i
det Uudgrandskelige; med denne Forklaring lyder det, ikke
som en Philisterbemærkning, men som en Røst fra Dybet:
„Ins Inn're der Natur dringt kein erschaffner Geist“.

Men ihvorvel Udtydningen af Begrebets Forhold til den
derved frembragte Gjenstand maa blive ufuldkommen, forsaavidt
Eenheden af Tænken og reel Frembringen er ufuld%
% --- p. 147 ---
\opage{147}kommen, saa staaer det dog ligefuldt fast, at den frembragte
Gjenstand kun kan forklares af en Subjectivitet, der i et
logisk System af antilogiske Domme formaaer at eftergjøre
Frembringelsen. De explicative Functioner have ikke blot
Betydning for den menneskelige Viden, de have absolut Betydning
for den guddommelige: den vidende Magt, der frembringer
Gjenstanden, er evig Eet med den mægtige Viden,
der forklarer Frembringelsen og i det Frembragte seer Begrebet
fyldestgjort.\footnote{Det hedder i Religionens Sprog: „Og Gud saae Alt, hvad han
havde gjort, og see, det var Altsammen saare godt“. At Explicationen
i logisk Forstand maa optage Bestemmelser af saadanne Kategorier
som: Bedømmelse, Anerkjendelse, Gjenkjendelse o. s. v., er
en Selvfølge.} % note *)

\textit{γγ}) \emph{Ideeen i Dom og Gjenstand: ideale Magtfunctioner.}
Den reale Verden er en Verden af Gjenstande,
den ideale en Verden af Tanker; Tanken har som
Tanke sit Udtryk i Dommen, derfor er Forholdet mellem
Dom og Gjenstand fra Totalitetens Synspunkt afgjørende
ved Bestemmelsen af Forholdet mellem den ideale og den
reale Verden. Den ideale Verden har indre Tilværelse i
Videns Ord, den reale ydre Tilværelse i Magtens Ord: ved
Ordets Magt er Dualismen gjennemført; ved Ordets Magt
maa Dualismen atter hæves. I Dualismens Ophævelse
sees igjen det Absolute, Ideen; men Ideen, det Absolute, kan
med sin dybe Selvmodsætning umulig komme til Aabenbarelse
i nogen blot endelig og relativ Form, det Absolute fordrer
en uendelig og absolut Form; men en uendelig og dog
absolut Form kan kun gives i Ordet.

Gjennem Ordet, Almagtens Ord, bevæger det Absolute
sig i et dobbelt Kredsløb, et subjectivt og et objectivt. I det
subjective Kredsløb er Magten ideel-reel; den reale Verden
udgaaer fra den ideale og vender i det Uendelige tilbage til
% --- p. 148 ---
\opage{148}den: det Absolute er sat i subjectiv Form, og dette Absolute,
det Absolute i subjectiv Form, svarer til Magtordets Idealitet.
I det objective Kredsløb er Magten reel-ideel; den ideale
Verden afspeiler den reale, saaledes som denne gjennem Subjectiviteten
indkredser sig selv: det Absolute er sat i objectiv
Form, og dette Absolute, det Absolute i objectiv Form, svarer
til Magtordets Realitet. Saaledes er Selvmodsætningen af
det Absolute. Det Absolutes Selvmodsætning viser sig saa
efter sin virkelige Eenhed objectiveret i et til Fordoblingsprincipet
svarende Magtsprog. Altsaa:

1) \emph{Magtordets Idealitet: imperative Functioner.}
Under Udviklingen af Forholdet mellem Dom og
Gjenstand har Magtordet allerede frembudt saavel sin ideelle
som reelle Side: Udviklingen er foregaaet saaledes, at den fra
anticiperende gjennem executive og explicative Functioner løber
tilbage i sig selv: det begynder og ender med Anticipationen.
Flygtig seet, skulde det synes, som om de explicative Functioner
stillede sig i ligefrem Modsætning, saavel til de executive
som til de anticiperende, og paa denne Maade dannede
en endelig Afslutning; men saaledes forholder det sig ikke,
naar Sagen sees i Principet: de explicative Functioner danne
tvertimod en væsentlig Forudsætning for de anticiperende.

At Anticipationen bestemmes ved objective Hensyn, Blik
paa det Virkelige, er ovenfor viist; det er de objective Muligheder,
hvorom det gjælder. Men de objective Muligheder
udfindes kun ved Udtydning af det Virkelige, altsaa kun ved
explicative Functioner. At Explicationen ikke ved endelige
Resultater kan tilfredsstille den uendelige Subjectivitet, er da
ogsaa let at indsee. Det i endelig Forstand tilfredsstillende
Resultat skulde jo være en bestemt Gjenstand, et fuldendt
Værk, hvori den dømmende, Værkets Betydning explicerende,
Subjectivitet kunde see Begrebet fuldkommen realiseret; men
den Gjenstand, det Object, som her repræsenterer det fyldestgjorte
Begreb, er i Qvalitet af virkelig Gjenstand Forandring
% --- p. 149 ---
\opage{149}underkastet; ved Forandringen bliver den for Begrebet adæqvate
Gjenstand jo netop uadæqvat: den ved Begrebet bestemte
Virkelighed er, saa at sige, svanger med Muligheder.
Skal Virkeliggjørelsen af de nye Muligheder nu ogsaa tilfredsstille
Begrebet, maa den begribende Subjectivitet ved
anticiperende Functioner igjen bestemme disse Muligheder og
præformere den Virkelighed, hvori Begrebet skal realiseres.
Den nye Virkelighed er da atter en for Begrebet adæqvat
Gjenstand, som ved at forandres bliver uadæqvat; den uadæqvate
Gjenstand betinger da en ny Anticipation, og saaledes
i det Uendelige.

Anskues nu Anticipationernes evige Kredsløb i Belysning
af det Absolute, da bliver den hele Virkelighed med alle
til dens uendelige Totalitet hørende Elementer gjennemskinnet
af Idealitet, et Medium for Ideens Selvaabenbarelse; hvert
Virkelighedselement er ved Anticipationen et Ideens Moment
fra Begyndelsen af: hvert Led i den hele Virkelighed, hver
virkelig Skikkelse har efter sin objectiv-subjective Mulighed
som Ideemoment oprindelig været anticiperet og objectivt
præformeret i den uendelige Subjectivitet. I Anticipationernes
evige Kredsløb sees den reale Verden uophørlig at fremgaae
af den ideale; Realsystemet præformeres i Idealsystemet: det
Absolute er i sit subjective Kredsløb.

De virkelige Frembringelser befinde sig altsaa i et dobbelt
System af Relationer; fra objectiv Side ere de Endelighedens
Betingelser underkastede og Relativiteten hjemfaldne;
fra subjectiv Side ere de Tilkjendegivelser af det Ubetingede,
Udtalelser af det Absolute. Forsaavidt nu en Frembringelse
direct udtaler det Absolute i Magtens Ord, viser den sig som
en primitiv Naturfrembringelse; forsaavidt den i Kredsen af
sine Betingelser tilhører Relativiteten og den endelige Causalitet,
hører den i Klasse med ordinære Naturprodukter. Den
primitive Naturfrembringelse kaldes da ogsaa, forsaavidt her i
særlig Forstand lægges Eftertryk paa Magtordets Idealitet,
% --- p. 150 ---
\opage{150}en \emph{Skabelse}. Da nu Religionen i sit Sprog har paatrykt
Kategorien Skabelse et ganske eiendommeligt Præg, saa
opstaaer i denne Sammenhæng det Spørgsmaal, hvorledes
Kategoriens religiose Betydning forholder sig til den videnskabelige,
% UNSURE: religiose or religiøse (no ø-stroke visible in scan; layer religiose; also below)
og hvad Dogmatiken, navnlig den speculative, vel
har gjort for at opklare dette Forhold.

Hvorledes „den aabenbarede Religion“ udtrykker sig i
Fortællingen om Skabelsen, er bekjendt. Enten vi oversætte
første Vers af Genesis ved: I Begyndelsen \emph{skabte} Gud,
eller ved: I Begyndelsen \emph{havde} Gud \emph{skabt} Himlen og
Jorden, saa bliver Skabelsen i begge Tilfælde opfattet som en
Begivenhed, hvorfor den ogsaa træffende fremstilles i Form
af en Fortælling: til Kategorien Begivenhed svarer Kategorien
Fortælling. Skal nu den religiose Fortælling omdannes til
Led i en Lærebygning, da maa den i Fortællingen udtalte
Forestilling, om ikke just forvandles til Begreb, saa dog
fastsættes ved en Begrebsbestemmelse, for at kunne opstilles
som Dogme. Saaledes gaaer Læren om Skabelsen fra Religionen
gjennem Theologien over i Dogmatiken.

Vælge vi til Exempel en af de ældre orthodoxe Definitioner
paa Skabelsen: \textit{Actio Dei triuni externa qua is
res omnes visibiles ex nihilo, sex dierum spatio, solo
liberrimæ voluntatis suæ imperio omnipotenter et sapienter
produxit, in nominis sui laudem et hominum
utilitatem} (Quenstedt), saa sees heraf, at Dogmatiken som
kirkeligt Lærebegreb opfatter Skabelsen som en Begivenhed i
Tidens Begyndelse; hertil svarer saa Skabelsen af Intet, en
Skabelse, der udelukkende henføres til Guds frie, almægtige
Villie.

Imellem Kirkelærens fastsatte Skabelsesdogme og den
rationelle Skabelseskategorie, hvoraf Ideelæren kan betjene sig,
naar den i primitive Naturfrembringelser seer en bestemt Udtalelse
af det Absolute, en Præformation af den uendelige
Subjectivitet, er der en gabende Kløft, et svælgende Dyb.
% --- p. 151 ---
\opage{151}Dette Dyb, denne Kløft har den speculative Dogmatik ikke
destomindre foregivet at kunne --- forstaaer sig, til en vis
Grad --- fylde og udjævne; vi skulle nu see, hvad „Fylde“
den bruger til Fyldning, og hvorledes den bærer sig ad med
at udjævne Kløften.

Ifølge Kirkelæren er Skabelsen en mirakuløs Begivenhed,
ikke et Grundsorhold; % PRINTER: printed "Grundsorhold" (long s for f), sense "Grundforhold"
enhver Biforestilling om et evigt Skabelsesforhold,
en Skabelse fra Evighed af, maa afvises som
et ufordrageligt Kjætteri. Ifølge Ideelæren er Skabelsen et
Grundforhold, ikke en mirakuløs Begivenhed; enhver Biforestilling
om Mirakel og mirakuløs Begivenhed maa afvises som
en Phantasie.

Det er nu i charakteristisk Overeensstemmelse med den
speculative Dogmatiks tvetydige Standpunkt, --- thi Standpunktet
er fra Først til Sidst en Tvetydighed ---, at den i
alle Hovedpunkter bevæges af ueensartede Hensyn, Motiver,
som modstræbe hinanden og ende med at afsløve og neutralisere
hinanden indbyrdes. Dette sees især i Behandlingen af
Skabelsesdogmet. Den dogmatiske Speculation vil ikke bryde
med Orthodoxien, vil for enhver Priis undgaae Skin af
Kjætteri; her er det første objective Hensyn, det kirkelige
Motiv. I Kraft af det kirkelige Motiv gjør Speculationen
Mine til, at det skal være Alvor med Skabelsen som mirakuløs
Begivenhed, med Skabelsen af Intet, med Skabelsen af
ubetinget Frihed. „Den mosaiske Skabelseshistorie udtrykker
Skabelsens Grundtanke derved, at Verden er bleven til ved
Guds almægtige Ord. Han talede: „Det blive Lys“! og
det blev Lys. Enhver af de sex Skabelsesdage oprinder kun
i Kraft af det almægtige Skaberord“. Til en vis Grad
betones da ogsaa i Overeensstemmelse hermed Skabelsen af
Intet: „Idet Gud skaber, kalder han det Ikke-Værende til
Tilværelse. Dette er Betydningen af den gamle Lære, at
Gud skabte Verden af Intet“.

Men, det forstaaer sig, disse Betoninger gjælde naturligviis
% --- p. 152 ---
\opage{152}kun til en vis Grad. Her er, foruden det objective
Hensyn til Kirken, endnu et andet objectivt Hensyn, det objective
Hensyn til Videnskaben. Bevæget af det videnskabelige
Motiv dreier Speculationen lemfældigt, for Uindviede næsten
umærkeligt, af fra det kirkelige. Skabelsen er nu ikke blot
en mirakuløs Begivenhed, men et Grundforhold, det Forhold,
hvori Gud oprindelig maa staae til Verden. I Kraft af det
videnskabelige Motiv bliver da „enhver af de sex Skabelsesdage“
til „enhver ny Epoche i Verdenssystemet“. Det Intet,
hvoraf Verden er skabt, er ikke = 0, mod hvilken Forestilling
det gjælder: \textit{ex nihilo nihil fit}“; % PRINTER: closing quote without opening quote
det Intet, hvoraf Verden
er skabt, er da heller ikke det orthodoxe Intet, hiint \textit{nihil
negativum, negatio omnis entitatis}, ifølge hvilket en Hollazius
kan imødegaae Aristoteles med en Sætning som denne: \textit{ex
nihilo nihil per vires naturæ generatur; at per vim Dei
infinitam omnia creata sunt ex nihilo}. Heller ikke kan
det videnskabelige Motiv godt enes med det orthodoxe, naar
Talen er om den ubetingede Frihed (\textit{solum liberrimæ voluntatis
imperium}), hvormed Gud i kirkelig Forstand har skabt
Verden. Vistnok har Gud skabt Verden paa den ene Side
med fuld Frihed; men Verden er dog paa den anden Side
bleven til med uafviselig Nødvendighed. Vistnok er det en
Kjærlighedsnødvendighed; men Nødvendighed er dog altid
Nødvendighed. „Just fordi Gud er Kjærlighed, \emph{kan} han
ikke afslutte sig selv som „Ideernes“ Gud, men \emph{maa} bestemme
sig o. s. v. Ja, han maa bestemme sig, thi det er
„Frihedens overvættes Rigdom, der ikke kan Andet, end at
ville“ \dots\ See, nu er Resultatet ligesaa tvetydigt som Udgangspunktet.

Naar Kirken nu spørger efter Friheden, saa gaaer Nødvendigheden
op i „Frihedens overvættes Rigdom“; og naar
Videnskaben fordrer Nødvendigheden, ja, saa gaaer „Frihedens
overvættes Rigdom, der ikke kan Andet end at ville“ \dots,
under i Nødvendigheden. „Fra dette Synspunkt indlyser det,
% --- p. 153 ---
\opage{153}i hvilken Betydning vi benegte, og i hvilken Betydning vi
bekræfte den Sætning: Uden Verden er Gud ikke Gud“.

I hvilken Betydning benegte vi da Sætningen? I kirkelig
orthodox Betydning, naturligviis! thi Nødvendigheden
gaaer jo under i Friheden. I hvilken Betydning bekræfte vi
Sætningen? I videnskabelig, det er at sige: i speculativ
videnskabelig Betydning, naturligviis! thi Friheden gaaer jo
under i Nødvendigheden. Benegte vi da virkelig Sætningen?
Ja, thi det kirkelige Eftertryk falder paa Friheden. Men er det
nu ogsaa vist, at vi virkelig benegte Sætningen? Nei, thi det
videnskabelige Eftertryk falder jo dog paa Nødvendigheden.
Saaledes blive da Ja og Nei tilsidst en god Theologie.

„Men“, kunde man indvende, „Ja og Nei ere jo i denne
Sammenhæng dialektiske Bestemmelser for det Samme; Friheden
er en dialektisk Bestemmelse af Nødvendigheden, Nødvendigheden
en dialektisk Bestemmelse i Friheden: hvorfor
skulde Dogmatiken ikke have samme videnskabelige Ret til at
benytte sig af Dialektiken, som Philosophien har“? Hertil
maa svares, at Philosophien, hvad Dialektiken angaaer, er
forpligtet til at gjøre Forskjel paa det Tvesidige og det Tvetydige.
Denne Forskjel kan i nærværende Sammenhæng,
hvor Alt dreier sig om Bestemmelser af Idealiteten i Magtens
Ord, tydeliggjøres ved en Belysning af Magtordets \emph{imperative}
Functioner.

I de executive Functioner viste sig en praktisk Forskjel
mellem den bestemmende Tanke og den virkelige, af Tanken
bestemte Handling; den til Forskjellen svarende Eenhed er at
søge i den praktiske Selvbestemmen, i Villien. Men Villiens
ideelle Indvirken paa det legemlig Reelle maa umiddelbart
beroe paa en med Virken identisk Byden. Man kan omskrive
den rene Villen paa mange Maader, men den fra Villien
udgaaende Bestemmen er væsentlig Befalen og Byden. Forsaavidt
den virkelige Frembringelse opfattes som et realiseret
Begreb, staaer den tillige som objectivt Udtryk for et
% --- p. 154 ---
Sub\opage{154}jectivitetens Bliv, et Tegn paa Idealitetens Magt. Kjendsgjerningen
er et Factum; hvert Factum forudsætter et \textit{fieri},
hvert \textit{fieri} et \textit{fiat}: i det Absolutes subjective Kredsløb gaae de
executive Functioner op i de imperative.

I Consequens af de imperative Functioner føres altsaa
ethvert Factum tilbage til et \textit{fiat}, det Objective i Frembringelsen
tilbage til det Subjective, Magtordet i Gjenstanden tilbage til
Vidensordet i Dommen; men at henføre Factum til \textit{fiat}, det
Objective med dets Betingelser til den i Begrebet bestemmende
Subjectivitet er, med andre Ord, at underordne Nødvendigheden
under Friheden, at føre Naturvirksomheden tilbage til
Skabelsen. De imperative Functioners Gyldighed maa i
denne Henseende ligesaavel anerkjendes af Ideelæren som af
Troeslæren; men naar Anerkjendelsen nærmere bestemmes,
aabner Kløften sig paany, den dybe Kløft, der sætter Skilsmisse
mellem Tro og Viden.

I Troeslæren tages det guddommelige Bliv absolut;
Skabelsesacten er ubetinget fri, den guddommelige Magt og
Villie veed ikke, vil ikke vide af nogensomhelst tilskyndende
Nødvendighed at sige. Gud er Gud fra Evighed af; men
Verden er ikke fra Evighed af; Gud er Gud fra Evighed af
--- uden Verden. Motivet til Skabelsen, Mulighed til Skabelsen,
Plan og Formaal for Skabelsen, Alt er i eet eneste
Øieblik samlet og med Uendelighedens Energie concentreret i
eet Ord, eet eneste Ord, det almægtige Ord: \emph{Bliv}! Med
dette Ord begynder Tiden, forud for dette Ord er Evigheden;
med dette Ord begynder Verden, forud for dette Ord er Gud
alene, Gud uden Verden.

Denne Troens absolute Opfattelse af det absolute Bliv
formaaer Ideelæren ikke at tilegne sig. At Gud er Gud fra
Evighed af, betyder, naar Ideelæren skal oversætte det i Begrebet,
at det guddommelige Væsen, den i sin indre Uendelighed
frie Subjectivitet, er med Nødvendighed. Intet almægtigt
Bliv kan byde Guddommen selv at blive til Noget;
% --- p. 155 ---
\opage{155}intet almægtigt Bliv kan byde Guddommen selv at blive til
Intet. Men i Begrebet som Begreb kunde Selvheden umulig
være, dersom Andetheden, Selvets Modsatte, ikke oprindelig
var; seet i Begrebet som Begreb kunde Subjectiviteten, den
i sin indre Uendelighed frie Subjectivitet, ikke være, dersom
Objectiviteten, den reelle Objectivitet, Gjenstandens udvortes
virkelige Verden, ikke virkelig var. Det almægtige Bliv, der
skulde kalde den reelle Objectivitet fra Ikke-Væren til Væren,
dette samme almægtige Bliv maate da ogsaa kalde den uendelige % PRINTER: printed "maate", sense "maatte"
Subjectivitet fra Ikke-Væren til Væren; det almægtige
Bliv, der skulde kalde den udvortes virkelige Verden fra Intet
til Noget, maatte med det Samme kalde det indre, guddommelige
Væsen selv fra Intet til Noget. Dette almægtige, i sin
Absoluthed ubetingede Bliv er Videnskaben et Ubegribeligt.

Men naar Videnskaben ikke kan begribe en Tilblivelse af
Intet, hvorledes kan den da vedkjende sig Kategorien Skabelse?
Her maa igjen lægges Vægt paa de imperative Functioner.
Det Rationelle i Skabelsen beroer, ligesom det Rationelle i
Magtordets Idealitet, det almægtige Bliv, paa Grundforholdet
af det Subjective og det Objective: i dette Grundforhold
er Selvheden, den uendelige Subjectivitet, der igjennem
Andet forholder sig til sig selv, evigt det Overmægtige, det i
Alt Overgribende. Ved den i det Uendelige overgribende
Magt faaer den subjective Bestemmen, det guddommelige Bliv
oprindelig Betydning for alle Naturfrembringelser; af denne
Betydning fremlyser det Rationelle i Skabelsens Kategorie.
Ideelæren kan fra dette Synspunkt erkjende, at Verden er
skabt, at den i oprindelig Modsætning til Gud er bleven til
af Intet; men den Tilbliven af Intet, hvorom Ideelæren
taler, er i kategorisk Strid med den Skabelse af Intet, hvorpaa
Troeslæren bygger.

Det er i vort Foregaaende paa flere Steder omhandlet,
hvorledes Philosophien har stridt med sig selv, for at faae
Materiens Realitet til at stemme med Fornuftens Fordring;
% --- p. 156 ---
\opage{156}det er viist, at Gaadens Løsning maa søges i Magtens Begreb:
Masserne med deres Atomer ere Kjendsgjerninger,
punktuelle Udslag af Magtens Umiddelbarhed. Da nu % PRINTER: printed "Da- nu" (stray raised mark), sense "Da nu"
Magtens reale Umiddelbarhed er Moment i Magtens egen
Realitet, og den reelle Magt i Subjectiviteten selv identisk med
den ideelle, saa forstaaes heraf, at den alt Virkeligt anticiperende,
alle tilværende Skikkelser præformerende Subjectivitet
ikke kan have enten evigt selvstændige Atomer eller Stof eller
Chaos udenfor sin, i alle Ting overgribende, Magt. Alle
Stofbestemmelser, alle materielle Elementer ere Realnegationer
i Guddommens ideelle Magt, antilogiske Momenter i den altformaaende
Idealitet; Alt dannes og frembringes ved Magten
selv, den guddommelige Magt: udenfor den guddommelige
Magt er der Intet, slet Intet, hvoraf Noget kan dannes.

Men med et saadant Intet kan Troeslæren ikke være
tilfreds. For Troen betyder Skabelsen af Intet, at den hele
Verden paa Almagtens Bud gaaer frem af Intet; herved forudsættes,
at Verden, før Skaberen udtalede sit Bliv, var
Intet, slet Intet, at der, som ovenfor viist, gaaer en Evighed
forud for den endelige Skabelse, og at der i hiin Evighed
ingen virkelig Verden var. Denne religiøse Grundtanke vil
ingen Ideelære formaae at hæve til videnskabeligt Begreb.
Skulde Gud tænkes at være en Evighed forud for Verden,
da maatte den ideelle Magt jo være en Evighed forud for den
reelle. Med den reelle Magt, med den guddommelige Magts
reelle Udslag er Verden øieblikkelig til; --- thi hvad er Verden,
seet fra Skabelsens Synspunkt, vel Andet end det reelle Udslag
af guddommelig Magt? --- Forløber der altsaa en heel
Evighed, inden det skabende Bliv lader Verden fremstaae, saa
forløber der ogsaa en heel Evighed, inden det skabende Bliv
giver den ideelle Magt Realitet. Men den ideelle Magt for
sig, uden den reelle Magt, er kun Halvmagt, ikke Almagt; kun
Identiteten af ideel og reel Magt er Almagt.

Seet fra dette Synspunkt bliver Skabelsesmiraklet nu
% --- p. 157 ---
\opage{157}til to Mirakler, det ene, om muligt, vidunderligere end det
andet. Først er det et Mirakel, at Skaberens absolute Bliv
sætter en Verden i Intets Sted, en Verden, hvor forhen
Intet var; dette Mirakel forudsætter Almagten. Dernæst er
Almagtens Tilblivelse selv et Mirakel. For at gjøre Miraklet
med Verden, maa Skaberen gjøre et Mirakel med sig selv,
idet han maa kalde den reelle Magt fra Ikke-Væren til
Væren; for at gjøre Almagtens Mirakel, maa Skaberen først
ved et Mirakel gjøre sig selv almægtig.

Og hvad beviser saa dette? det beviser ikke, at Religionen
har Uret imod Videnskaben, thi Religionen indlader sig ikke
paa at begribeliggjøre Skabelsens Mirakel; men det beviser,
at ethvert videnskabeligt Forsøg med at forklare det almægtige
Bliv, ved at tillægge de imperative Functioner abstract ubetinget
Gyldighed maa ende i det Meningsløse.

2) \emph{Magtordets Realitet}: \emph{indicative Functioner.}
I Systemet af de imperative Functioner har Subjectiviteten
Raadighed over det Frembragte paa alle Frembringelsens
Stadier. Frembringelsen begynder med Explicationens Overgang
eller rettere Tilbagevenden til Anticipation; fra begge
Sider er Begyndelsen bestemt ved den bestemmende Subjectivitet.
Den begyndende Frembringelse underlægges de executive
Functioner; de virkelige Handlinger, de practisk-antilogiske
Sætninger gaae frem af antilogiske Domme, begge i
deres Eenhed af den tænkende, virkende Subjectivitet. Resultatet
af den fuldendte Frembringelse, det Frembragte, som
ifølge Bedømmelsen faaer Betydning derved, at det baade tilfredsstiller
og dog ikke tilfredsstiller den sin Frembringelse explicerende
Subjectivitet, bliver da Moment for nye Anticipationer,
nye Frembringelser o. s. v., altsaa bestandig Subjectiviteten
underlagt. Saaledes har det Absolute optaget alle
Frembringelsens Elementer og Betingelser som Momenter i
sit eget subjective Kredsløb: Subjectivitetens Overlegenhed
culminerer i et altformaaende Bliv.

% --- p. 158 ---
\opage{158}Men at Frembringelsens objective Elementer ere Subjectiviteten
immanente, forudsætter tillige, at Begrebsbestemmelserne,
Frembringelsens subjective Momenter, ere immanente
i Objectiviteten. Hele Kredsbevægelsen maa altsaa, for at
sees i sin Realitet, betragtes fra den objective Side. Saavel
ved de anticiperende som ved de explicerende Functioner er
det ligesaameget Gjenstanden, der bestemmer Dommen, som
omvendt Dommen, der bestemmer Gjenstanden; men at Gjenstanden
bestemmer Dommen, vil jo med andre Ord sige, at
det Objective bestemmer det Subjective. Den uendelige Vexelbestemmelse
af Objectivt ved Subjectivt, og omvendt, maa da
i objectiv Forstand ende med, at det Objective gjennem det
Subjective, fra Frembringelsens Begyndelse til dens Fuldendelse,
Punkt for Punkt bestemmer sig selv. Hermed er Objectiviteten
udtrykkelig sat som det Absolutes egen Form, som
Ideens immanente Form: det Absolute er i sit objective Kredsløb,
og dette er Magtordets Realitet.

Ved Magtordets Idealitet blive Functionerne imperative,
ved Magtordets Realitet blive de indicative; Opgaven er at
vise, hvorledes det Indicative forholder sig til det Imperative.
Overalt, hvor det Objective bestemmer det Subjective, gaaer
Dommen ud paa en ligefrem Angivelse af, hvad det Objective
er; i denne Angivelse sees det Indicative. De i den
formelle Logik afhandlede Domme ere alle indicative. At Bestemmelsen
af det Indicative desuagtet ikke kan træde frem i
den sædvanlige Logik, har sin Grund i, at det mangler sin
Modsætning, nemlig det Imperative: vi skulle nu see, hvorledes
det Imperative fortoner sig og kommer til Ro i det Indicative.

Det Objective maa objectiveres, de objective Frembringelser
præformeres; men hvormegen Deel Subjectiviteten end
har i Præformation og Udførelse, saa kan dens Medvirken
dog aldrig gaae videre end til at bestemme det Objective ret
objectivt. At bestemme det Objective objectivt er ikke at
% --- p. 159 ---
be\opage{159}stemme det umiddelbart paa noget Punkt, men lade det i alle
Punkter bestemme sig selv. Det Objective i Frembringelsen
træder nu fra alle Sider saaledes frem, at den frembringende
Subjectivitet synes at blive aldeles borte. Det Eneste, Subjectiviteten
endnu har at gjøre, er at objectivere sig den objective
Frembringelse; men denne Objectiveren er saa theoretisk,
saa contemplativ, at Intet forandres i hvad der objectivt
foregaaer.

Forskjellen mellem den endelige og den uendelige Subjectivitet
er ogsaa i dette Punkt afgjørende. Naar den endelige,
den psychologiske Subjectivitet skal indvirke bestemmende
paa objective Frembringelser, da mærker man strax Adskillelsen
mellem det subjectivt Vilkaarlige og det objectivt Nødvendige;
just fordi det Objective sætter den afmægtige Subjectivitet saa
umiddelbare Skranker, maa denne igjen ligesaa umiddelbart i
enkelte Punkter gjennembryde disse Skranker; her opstaaer en
endelig Conflict imellem det, hvad Subjectet vil, og det, hvad
der umiddelbart ligger i det Objective; derfor er i alle Fabrikater
og Kunstprodukter Subjectivitetens Indflydelse saa
øiensynlig.

Anderledes med den uendelige Subjectivitet; dens Indvirken
paa det Objective foregaaer ikke paa endelig Maade og
er af den Grund ikke umiddelbar kjendelig. Uagtet Alt er
anticiperet, præformeret og subjectivt bestemt ved det mægtige
Bliv, uagtet der ikke er et Element i det Objective, uden at
dets Væren er forberedt subjectivt, saa seer det dog i Virkeligheden
ud, som om det Objective gik ved sig selv. Ja, det
Objective gaaer virkelig ved sig selv objectivt; her er i Tilsyneladelsen
ingen Skuffelse, intet Bedrag: i alle Frembringelser
er den virkelige Bestemmelse en Bestemmelse af Objectivt
ved Objectivt paa objective Muligheder; Subjectiviteten
har kun at see til, dens Functioner ere indicative.

Det Indicative i Functionerne har alligevel ikke fortrængt,
men forklaret det Imperative: det Indicative viser, at
% --- p. 160 ---
Ud\opage{160}viklingen af det Objective er objectiv; men det Imperative
viser, at det Objective kun er objectivt formedelst Subjectiviteten.
Det Indicative viser, at Naturhandlingerne ledes af
Principer og immanente Begreber; men det Imperative viser,
at Principerne ere Selvhedsfactorer og de immanente Begreber
logisk-antilogiske Udtalelser af en begribende Subjectivitet.
Det Indicative viser, at enhver Skabelse virkelig er en
Naturfrembringelse; men det Imperative viser, at enhver
Naturfrembringelse oprindelig er en Skabelse. Her er Identitet
af Natur og Skabelse.

Men den speculative Dogmatik vil jo ogsaa Identitet af
Natur og Skabelse; lader os da igjen høre lidt paa den speculative
Dogmatik. „Som den johanneiske Anskuelse indeholder,
at Verdens Tilværelse har sin Grund i en skaberisk
\emph{Frembringen}, saaledes indeholder den, at Verden er til
ifølge en \emph{Overgang} fra Ikke-Væren til Væren, ved en \emph{Vorden},
en Tilblivelse, en Fødsel, et \textit{fieri}, et \textit{γίγνεσθαι}. Verden
har altsaa en dobbelt Begyndelse, en kosmogonisk og en skaberisk,
en naturlig og en overnaturlig. Den kosmogoniske \dots\
er den relative, den endelige, som derfor ogsaa er splittet i \emph{sporadisk}
Mangfoldighed \dots\ Men de utalligt mange naturlige
Begyndelser ere alle begrundede i den ene skaberiske, overnaturlige
Begyndelse, i den guddommelige Logosvillie, der i
sig indeslutter Livets og Lysets Væld, og af sit aandelige Dyb
løslader den hele Mangfoldighed af Livskræfter til fri Selvbevægelse.
Og kun fordi denne overnaturlige Begyndelse, kun fordi
den skabende Villie vedbliver at røre sig i de mange endelige Begyndelser
og i fri Allestedsnærværelse gjennemlyser og gjennemvirker
den naturlige Udvikling, kan denne grundigt komme ud
over den chaotiske Gjæring, kunne de sporadiske Modsætninger
forenes og harmoniseres til et organisk og systematisk Hele.
Verden maa derfor i ethvert Moment af sin Tilværelse betragtes
som \textit{natura} eller som organisk Selvudvikling, og som % natura, creatura: antiqua, letterspaced
\textit{creatura} eller som fremskridende Aabenbaring af den
% --- p. 161 ---
guddom\opage{161}melige \emph{Villie}, og den er kun det Ene, \emph{fordi} den er det
Andet“.

Imod denne Udtalelse gjælde to Indvendinger: først at
den er afpasset efter det rationelle Skabelsesbegreb, men
staaer i Strid med det kirkelige Dogme; dernæst, at den er
holdt i en blot beskrivende Form, og er desaarsag begrebløs.

Hvad det første Punkt angaaer, da er Skabelsen af Intet
(\textit{creatio ex nihilo}) i kirkelig Forstand et Udtryk for den
guddommelige Villies absolute, ubetingede Magt over Alt, hvad
vi i Virkelighedens Forstand kalde Andethed, Objectivitet,
Tingsaarsager o. s. v. Det er saa langt fra, at det guddommelige
Selv væsentlig skulde behøve en saadan Andethed, at det
tvertimod kan afskaffe den, naar det vil; det er saa langt fra,
at den guddommelige Subjectivitet skulde til sin egen Virkelighed
behøve et virkelig Objectivt, at den, hvad Øieblik den
vil, kan afskaffe Objectiviteten. At Tingene og Tingsaarsagerne
ere skabte af Intet, betyder, at de igjen, hvad Øieblik det behager
Almagten, kunne blive til Intet; til \textit{creatio rerum}
svarer i den gamle Dogmatik --- og det med sand Consequens
--- \textit{annihilatio rerum}.

„Der er i Theologiens Historie ingen Dogmatiker, som
med større Strenghed har fastholdt Begrebet om Guds Almagt,
end Malebranche, Ingen, som med større videnskabelig
Consequens har udtalt, hvad der ligger i Skabelsesdogmet,
end netop Malebranche.\,\dots\ Vel har Malebranche ikke paa
speculativ Maade kunnet see ind i Faderens Urgrund, opdage
Verdensspirer i Sønnens Pleroma og mediere Skabelsens
Proces i Aandens plastiske Kraft; men Dogmet har han, det
rene Dogme om Skabelsen af Intet. Han har om dette
Dogme ikke blot erklæret, at det maa tænkes (thi fordi det engang
imellem siges om et Dogme, at det maa tænkes, derfor
er det jo ikke sagt, at det virkelig bliver tænkt); nei,
han har i Alvor forsøgt at tænke det, holde Tanken fast og
anvende den til en videnskabelig Forstaaelse af det Skabte.
% --- p. 162 ---
\opage{162}At anvende Skabelsestanken til en Forstaaelse af det Skabte
vil sige: at see den hele Verden under et dobbelt Synspunkt:
\emph{Intet og dog virkelig Noget}. Det Begreb, hvori denne
Modsigelse løses, er Begrebet Almagt; det Skabte er i sig
Intet, men paa Almagtens Bud er det blevet til Noget.
Den, der nu seer paa det Tilblevne saa eensidigt, at han kun
seer dets Intethed, for ham er Verden Skin og Skygge. Om
Theologien vil kalde ham Pantheist eller Akosmist, er ligegyldigt;
Sagen er, at han fornegter Skabelsen, idet han fornegter
Almagten. Den, der seer paa det Tilblevne saa eensidigt,
at han kun seer dets Realitet, glemmer dets Intethed
og tillægger det Væsenhed, for ham er Tilblivelsen en Naturproces
og Verden et Naturens Rige. Om Dogmatiken forøvrigt
maa henregne ham blandt Deister, Naturalister eller
Atheister, er i denne Sammenhæng ligegyldigt; Sagen er, at
han fornegter Skabelsen, idet han fornegter Almagten.\,\dots\ Man
har nu rigtignok meent at kunne gaae et Skridt videre og
forbedre Skabelseslæren ved at tillægge det Skabte dog nogen
Selvstændighed, noget Anlæg til indre Selvudvikiing, nogen % PRINTER: printed "Selvudvikiing", sense "Selvudvikling"
% UNSURE: Selvudvikiing or Selvudvikling with damaged l
Evne til Virksomhed \dots\ Det er imidlertid en mislig Forbedring
\dots\ Thi hvad er en relativ Selvstændighed? En Grad
af Tilværen udenfor Almagten, en Selvvirken som, om den
end for Resten er nok saa ringe, dog ikke er Almagtens
Virken, altsaa: en virkelig Grændse for Almagten. Er det
nu Alvor med denne virkelige Grændse, saa er jo Almagten
en virkelig begrændset Magt og ikke længere Almagt. Har den
Almægtige derimod endnu, som for, Grændsen i sin Magt,
saa er den virkelige Grændse jo endnu, som før, et Intet i % PRINTER: printed "som for, Grændsen" (first instance), sense "som før"
sig, og det mirakuløse Forhold mellem Skaber og Skabning
maa vedblive. Et Forhold, der er begyndt ved et Mirakel,
maa ogsaa fortsættes ved et Mirakel: ikke blot nogle Øieblikke,
ikke blot, om jeg saa maa sige, en 8 à 14 Dage efter
Skabelsen, men i Aarhundreder, i Aartusinder, i Evigheders
Evighed maae de skabte Ting vedblive at være i sig, hvad de
% --- p. 163 ---
\opage{163}oprindelig ere og have været, et Intet, et reent Intet for den
Almægtige. Det er denne Consequens, Malebranche fastholder,
og derfor kan han ikke tænke sig Selvvirksomhed i nogen
Skabning, være sig Engel eller Menneske; thi væsentlig Selvvirksomhed
er en guddommelig Sag, og Gud selv kan ikke
skabe Guder: \textit{il n'en peut faire de véritables causes, il
n'en peut faire des Dieux.}“\footnote{Om Theologiens Naturbegreb af R. Nielsen. I Anledning af Reformationsfesten 1855. S. 22 flg.} % note *)

Ere Naturaarsagerne, de virkelige Tingsaarsager, skabte,
nemlig i dogmatisk-kirkelig Mening, eller ere de uskabte? Dette
er Spørgsmaalet, det uafviselige Grundspørgsmaal. Besvares
Spørgsmaalet med Ja: Naturaarsagerne ere skabte! saa har
Verden ganske rigtig „en overnaturlig,“ en „skaberisk“ Begyndelse;
men saa er det Nonsens at tale om, at den \emph{tillige}
skulde have en „kosmogonisk“, en „naturlig“ Begyndelse. Besvares
Spørgsmaalet med Nei: Naturaarsagerne ere uskabte,
saa har Verden ganske rigtig en „naturlig“ og „kosmogonisk“
Begyndelse; men saa er det Nonsens at tale om, at Verden
\emph{tillige} skulde have en „overnaturlig“ og „skaberisk“ Begyndelse.

Kun paa den Betingelse, at Skabelseskategorien opfattes
rationelt, kan Identiteten af Natur og Skabelse begribes. Den
Modsigelse, at det Skabte skulde være baade Noget og dog i
sig Intet, løses da derved, at Magten kun er Selvmagt
ved at have Magt over Andet, og kun kan have Magt over
Andet, naar dette Andet selv staaer ved Magt. Adskilt fra
den uendelige Selvmagt er det bestaaende Andet et Skin, et
Intet; men i Eenhed med Selvmagten er det existerende
Andet substantielt; dets Substantialitet er ved Modsætning
til Selvmagten dennes reale Moment.\footnote{Grundideernes Logik. I. S. 112--113.} % note **)
Den Modsigelse,
at her skulde være en virkelig Grændse for Almagten, løses
% --- p. 164 ---
\opage{164}derved, at den virkelige Grændse, Magtens egen Realnegation,
netop udtrykker Magtens Realitet: uden Realnegation vilde
Magten kun være ideel, ikke reel; men den abstract ideelle
Magt er kun Halvmagt, ikke Almagt. Istedenfor at gjøre
Almagten til en virkelig begrændset, altsaa indskrænket Magt,
er den virkelige Grændse tvertimod et kategorisk Udtryk for
den oprindelige Magtfordobling, for Magtordets evige Realitet.

Naar den speculative Dogmatik her og der tager Tilløb
til at forklare det Skabtes „givne Selvstændighed“ deraf, at
Gud giver Afkald paa sin Magt, eller at han løslader „de
magiske Muligheder,“ da er dette et umiskjendeligt Vidnesbyrd
om, at „det dogmatiske Talent“ har givet Afkald paa
Tanken og løsladt allehaande phantastiske Indfald. Man
forsøge blot at tænke Noget ved et Indfald som dette, at den
Almægtige ved Skabelsen virkelig, --- forstaaer sig, til en vis
Grad --- har givet Afkald paa sin Magt. Hvad Begreb skal
man nu gjøre sig om en saadan Given Afkald; den maa jo
være skeet enten ved Magt eller ved Afmagt? Er den skeet
ved Afmagt, altsaa af Mangel paa Magt, er Gud under
Skabelsen paa Grund af den det Skabte „givne Selvstændighed“
falden i en vis frivillig Afmagt, saa bliver Skabelsesmiraklet
jo ikke et Almagtens, men et Afmagtens Mirakel.
Er den derimod skeet ved Almagt; har Gud under Skabelsen
givet Afkald paa sin Magt saaledes, at denne hans Given
Afkald netop er bleven et Beviis paa hans uendelige, i alle
Henseender uindskrænkede Magt, saa har han da ogsaa, ved at
give Afkald paa sin Magt, beholdt sin Magt, og følgelig slet
ikke givet Afkald paa sin Magt.

At speculere over Skabelsen uden Indsigt i de herhen
hørende logiske Functioner er en mislig Sag; kun ved at reducere
Forholdet mellem Skaber og Skabning til et Grundforhold
mellem Subjectiviteten og den objective Virkelighed
og derved til et logisk Forhold imellem imperative og
% --- p. 165 ---
indi\opage{165}cative Functioner, vil det være muligt at komme til Klarhed
over saa indholdsrige Kategorier som \emph{Natur} og \emph{Skabelse}.

I Consequens af de imperative Functioner ere Naturobjecterne,
de virkelige Gjenstande, Momenter for den frembringende
Subjectivitet, den uendelig overgribende Selvmagt;
i Conseqvens af de indicative Functioner have Naturobjecterne,
de existerende Gjenstande, væsentlig Selvstændighed
ligeoverfor Subjectiviteten; i Conseqvens af de imperative Functioner
er den tænkende, dømmende Subjectivitet \emph{handlende},
i Conseqvens af de indicative Functioner er den handlende
Subjectivitet tænkende, \emph{dømmende}. Grundforholdet mellem
de indicative Functioner og Naturobjecternes virkelige Verden
lader sig da reducere til Ideen i Dom og Gjenstand.

Dersom den virkelige Gjenstand virkelig kunde bestemme
Dommen uden selv at bestemmes ved Dommen, da maatte
Subjectiviteten i sit System af indicative Functioner forholde
sig til Naturobjecterne som et Speil til de speilede Gjenstande.
Men selv i den reneste Contemplation er den subjective
Speilen dog meer end en Speilen; den subjective Speilen
er i sin Idealitet en uophørlig Tænken og Dømmen. Udenfor
Ideen, altsaa paa Endelighedens Trin, staaer Gjenstanden,
hvorover der dømmes, i et udvortes og ligegyldigt Forhold
til Dommen og den dømmende Subjectivitet. Gjenstanden
bestemmer Dommen, men Dommens indicative Function gjør
det ikke kjendeligt, at Dommen med det Samme bestemmer
Gjenstanden. Hvor dette Moment psychologisk fremhæves,
bliver, som vi ovenfor have seet, Afstanden mellem Subjectiviteten
og „Tingen i sig“ kun desto større.

Anderledes, naar Forholdet mellem Dom og Gjenstand
sees i Ideens Lys. Her bliver det klart, at Gjenstanden kun
kan bestemme Dommen, forsaavidt Dommen bestemmer Gjenstanden.
Det meest afgjørende Udtryk for, at Dommen væsentlig
bestemmer Gjenstanden, haves i de imperative Functioner.
Paa de imperative Functioners Grundlag kan den
% --- p. 166 ---
\opage{166}guddommelige Subjectivitet betragte Naturgjenstandene i deres
objective Umiddelbarhed som Frembringelser af Magtens Ord.
Det Objective i disse Frembringelser, der ere betegnende for
Magtordets Realitet, er bestemmende for Dommens indicative
Functioner.

Den dømmende Subjectivitet, der contemplativt lader
Gjenstanden bestemme Dommen, hviler med uendelig Overlegenhed
i Contemplationen af det Objectives objective Udvikling,
fordi den i sin Selvfordybelse evig veed, at Gjenstanden,
der bestemmer Dommen, oprindelig selv er bestemt ved Dommen;
den contemplative Subjectivitet, som i Systemet af indicative
Functioner skjuler sig under Naturfrembringelserne, veed sig
da identisk med den praktiske Subjectivitet, som gjennemaander
Principerne og i enhver Naturfrembringelse gjenkjender sin
egen Skabning.

3) \emph{Det kategoriske Magtsprog: interpreterende
Functioner.} Hvad det betyder, at der er Tanker i Tingene,
Begreber i Gjenstandene, maa nu være klart af det Foregaaende.
Men dersom der er Tanker i Tingene, saa maa
der vel ogsaa være \emph{Mening} i Tingene; er det da vist, at der
i objectiv Forstand er Mening i Tingene? I Forhold til det
rene Begreb staaer Menen dybt under Viden; men med Hensyn
paa Forholdet mellem Begreb og Existens kan en Indsigt
i Tingenes Mening staae over en objectiv Viden af
Tingenes Væsen. Til Erkjendelse af Tingenes objective
Væsen er det tilstrækkeligt at lægge den endelige Causalitet til
Grund og udforske Tingenes Love; til Opfattelse af Tingenes
Mening udfordres et Blik paa den Opgave, Existensen har
løst, en Indsigt i det Formaal, som ved Tingene skulde naaes,
et Formaal, der først kan siges at give dem, ikke en umiddelbar
Væren, men en væsentlig Betydning; at der er \emph{Tanker}
i Tingene, lærer Physiken; men hvad det egentlig vil sige, at
der er \emph{Mening} i Tingene, maa opklares i Logiken.

Man har i det logisk Explicative gjort Adskillelse mellem
% --- p. 167 ---
\opage{167}en Begriben af Gjenstanden som Hele og en Begriben af
Gjenstandens Dele. Her skal i denne Henseende være „en
\emph{dobbelt Art af Begriben}.“ --- „Hvad nemlig den givne
Gjenstand i det \emph{Hele} angaaer, da indsees,“ hedder det,
„\emph{hvorledes} Sammenhængen deri er beskaffen, eller \emph{hvorledes}
det Mangfoldige deri er forbundet til Eenhed. Hvad
derimod \emph{Enkelthederne} og de enkelte Forhold angaaer, da
indsees der, a) \emph{hvorfor de ere} og b) \emph{hvorfor de ere}
saaledes beskafne, \emph{som de ere}, idet der indsees, hvorledes
den hele Gjenstands Sammenhængsmaade \emph{fører eller kan
føre dem saaledes med sig}; saa at de, med den Beskaffenhed,
de have, sees som \emph{medfornødne} (\emph{integrerende})
og tillige \emph{medindgribende} eller dog \emph{medhenhørende},
eller naturligen \emph{medførte}, Led i det Hele. Naar altsaa,
hvad Gjenstanden i det Hele angaaer, blot den \emph{indvortes
Beskaffenhed} eller dens \emph{Tilværelsesmaade} begribes,
saa indsees derimod, hvad Enkelthederne og Forholdene angaaer,
tillige deres \emph{Tilværelsesgrunde} (\textit{rationes essendi}),
eller i al Fald deres Mulighedsgrunde, og tilmed deres \emph{Tilværelsesbetydning}
eller deres \emph{Realbetydning} i det
givne Hele. Det Hele bliver her kun \emph{fremstillet} eller \emph{beskrevet}
efter dets indvortes Beskaffenhed; Enkelthederne blive
derimod tillige \emph{udtydede} og \emph{forklarede}. Den første Art
Begriben ville vi“ --- hedder det endvidere --- „kalde en \emph{Sammenhængsbegriben},
maaskee og en \emph{formal Begriben}, og
Begrebet et \emph{Sammenhængsbegreb}, maaskee og et \emph{Formalbegreb};
den sidste derimod, med Hensyn til Brugen af
Ordet \emph{Realitet}, en \emph{real Begriben}, og Begreberne \emph{Realbegreber}.
En Begriben af hiin Art ligger her til Grund
for Begribningen af sidste Art; thi det er Sammenhængsbegrebet
om den hele Gjenstands Hvorledes, dens indvortes
Hvorledes nemlig (der ikke maa forvexles med dens udvortes
Beskaffenhed), hvoraf Indsigten i alt det \emph{Enkeltes} Hvorfor,
% --- p. 168 ---
\opage{168}hvoraf altsaa den reale Begriben i Henseende til alt dette
Enkelte, gaaer frem.“

At den saaledes opstillede Adskillelse mellem Sammenhængsbegreb
og Realbegreb er overfladisk, ubehjælpsom og
eensidig, maa ved en Smule Eftertanke let falde i Øinene.
Det Hele kan umulig bestemme Delene, uden at Delene omvendt
maae bestemme det Hele; skal Delenes Hvorfor udledes
af det Heles Hvorledes, da maa det Heles Hvorfor ogsaa
udledes af Delenes Hvorledes. Overalt, hvor der kan være
Tale om en Begriben, maa der gives en Begrunden, altsaa
et Hvorfor. Uden Hvorfor ingen Grund, uden Grund ingen
Begrunden, uden Begrunden ingen Begriben, uden Begriben
intet Begreb. Et Sammenhængsbegreb om hele Gjenstandens
indre Hvorledes uden Hvorfor kan ikke lægges til Grund for
nogen Indsigt i det Enkeltes Hvorfor, og det af den simple
Grund, at et saadant Sammenhængsbegreb ikke kan gives.
At der kan gives f.\ Ex.\ to Slags Skræddere: Dameskræddere,
der ikke sye for Herrer, og Herreskræddere, der ikke give sig
af med at sye for Damer, er ikke stridende mod sund Menneskeforstand
og skal ikke paatales i Logiken; men at der nu ogsaa
skulde gives saadanne to Slags Begreber som: rene Hvorledes-Begreber,
der ikke befatte sig med noget Hvorfor, og eiendommelige
Hvorfor-Begreber, der flyde ovenpaa „det indre
Hvorledes“, er stridende mod sund Menneskeforstand, og maa
paatales i Logiken.

Sammenhængsbegrebet om den hele Gjenstands indre
Hvorledes er ikke Andet og kan jo ikke være Andet end Begrebet
om de \emph{enkelte Deles} indbyrdes Sammenhæng; men
Begrebet om Delenes indbyrdes Sammenhæng er ikke Andet
og kan jo, væsentlig seet, ikke være Andet end Begrebet
om deres i Sammenhængen bestemte Beskaffenhed: som
Sammenhængen er, saaledes er Beskaffenheden, som Beskaffenheden
er, saaledes er Sammenhængen. Er det nu
umuligt at danne sig et Sammenhængsbegreb om det Hele
% --- p. 169 ---
\opage{169}uden at begribe Delenes Sammenhæng, da er det ogsaa umuligt
at begribe Delenes Sammenhæng uden at begribe, \emph{hvorfor}
hver enkelt Deel just har denne og ingen anden Beskaffenhed.
„Indsigten i alt det Enkeltes Hvorfor“ kan kun
lægges til Grund for „en Sammenhængsbegriben“ af Gjenstandens
indre Hvorledes, forsaavidt denne Sammenhængsbegriben
selv lægges til Grund for en Indsigt i alt det Enkeltes
Hvorfor. --- Men betragte vi den opstillede Adskillelse
mellem „Sammenhængsbegriben“ og „Realbegriben“ som en
famlende Antydning af den for Existenserkjendelsen afgjørende
Forskjel imellem det objective Begreb og den i Begrebet udtalte
Mening, lader den sig dog paa en Maade anvende.

„Den begrundende Reflexion“ har allerede lært os, at den
tilstrækkelige Grund kun er objectiv som Existensgrund, forsaavidt
den tillige er subjectiv som Bevæggrund. Til den
dobbeltsidige Grund svarer her et dobbeltsidigt Hvorfor. Det
objective Hvorfor reflecterer sig i det objective Begreb, det
Sammenhængsbegreb, der udtrykker Gjenstandens indre Hvorledes;
af dette Hvorledes fremlyser da ligesaavel det Heles
som Delenes Hvorfor: denne Reflexion af Hvorfor og Hvorledes
giver os Begrebet, men ikke dets Mening. Man tænke
paa Naturprodukter som Feldtspath, Granit o.\ dsl. Det vil
ikke give nogen Mening at lægge Begrebet om Feldtspathens
indre Hvorledes til Grund for en Indsigt i dens enkelte Deles
Hvorfor, naar dette Hvorfor ikke skal gjælde Existensgrunden,
men Motivet. Som det gaaer med Feldtspathen, saaledes
med den hele Jord, saalænge Begribningen holder sig til det
reent Objective. Spørge vi derimod om Meningen med
Feldtspathen og med Jordklodens Existens, da mærkes det
strax, at her ikke sigtes paa den objective Grund, Existensgrunden,
men paa det Existerendes subjective Grund, Motivet,
den for Feldtspathens og Jordklodens Frembringelse gjældende
Bevæggrund.

Er Gjenstanden ikke et Naturprodukt, men et Fabrikat,
% --- p. 170 ---
\opage{170}f.\ Ex.\ et Uhr, da ligger det jo vistnok nær at lægge Sammenhængsbegrebet
om det Heles indre Hvorledes til Grund for
Erkjendelsen af de enkelte Deles Hvorfor, thi det Heles
Hvorledes antager uvilkaarlig Præg af en vilkaarlig Opgave.
I Uhret har Værkmesteren løst en bestemt Opgave paa en
bestemt Maade; derfor er Uhret blevet et saadant, just et
saadant Hele, med saadanne Dele i en saadan Sammensætning.
Det Hele er fra Værkmesterens Side Formaal og Delene
Midler: det er altsaa et subjectivt, ikke et abstract objectivt
Hvorfor, hvorpaa her reflecteres. Dog selv fra det teleologiske
Synspunkt gjælder det, at det ikke blot er Midlernes Hvorfor,
der maa søges i Formaalets Hvorledes, men ogsaa omvendt
Formaalets Hvorfor, der maa søges i Midlernes Hvorledes:
ved Anticipationen er det Formaalet, der motiverer Valget af
Midlerne, forsaavidt de bestemmende Hensyn ere subjective;
men det er omvendt de sig tilbydende Midler, der motivere
Valget af Formaalet, forsaavidt de bestemmende Hensyn ere
objective.

Overalt, hvor Forholdet mellem et Hele og dets Dele
lader sig opfatte teleologisk, vil det paagjældende Værk tale
til Betragteren i et Sprog, der tyder ham Værkmesterens
Mening. Saasandt der er Mening i Tingen, maa det ligesaavel
være den productive Subjectivitet muligt at udtale sig
i Kjendsgjerninger, som at udtale sig i Ord; men Talen i
Kjendsgjerninger er forskjellig fra Talen i Ord: Kjendsgjerningernes
Sprog er et \emph{kategorisk Magtsprog}.

Et svingende Pendul forestiller et objectivt Begreb, men
dette Begreb er endnu uden Mening; er det svingende Pendul
derimod ved Hæmværket sat i Vexelvirkning med et Uhrs
Hjulværk, da er Foreningen af Pendul og Hæmværk i den
virkelige Sammensætning ligesaa talende som Foreningen af
Subject og Prædikat i den logiske Sætning: Pendulets
Svingning giver da en bestemt Mening. At en Valse, omviklet
med en Snor, hvis frie Ende bærer et Lod, sættes i
% --- p. 171 ---
\opage{171}Bevægelse ved Tyngden, fremstiller for den Reflecterende et
objectivt Begreb; men Begrebet for sig er endnu uden Mening.
Er Valsen derimod forbundet med Uhrværkets Valsehjul og
det saaledes, at det ikkun følger Valsen, naar denne dreies
formedelst Loddets Vægt, men ikke, naar den under Optrækningen
dreies i modsat Retning: da have vi i den mechaniske
Forbindelse af slige Led en ligesaa tydelig Mening som i en
grammatikalsk Forbindelse af Sprogets Ord. Ligesom det
enkelte mechaniske Led for sig ikke kan give nogen Mening,
saaledes kan det enkelte Ord for sig jo heller ikke give nogen
Mening. Det enkelte Ord for sig kan ved den blotte Benævnelse
betegne et Begreb; det enkelte Led for sig kan ved
sin individuelle Skikkelse, sit blotte Udseende, repræsentere et
Begreb. Sprogets Sætninger kunne efter Meningen være
deels Hovedsætninger, deels Bisætninger, deels Forsætninger,
deels Eftersætninger o.\ s.\ v.; det Samme gjælder paa tilsvarende
Maade om Sammensætningerne, de praktisk-antilogiske
Sætninger i den sin Opgave løsende Maskine. Hver Sammensætning,
hver praktisk-antilogisk Sætning er bestemt ved en
logisk-antilogisk Dom. Hvad Værkmesteren gjennem Systemet
af slige Domme har sagt til sig selv under Frembringelsen af
sit Værk, det siger han nu til den tænkende Betragter gjennem
det System af reale Sætninger, hvori det fuldførte Værk
giver Mening.

Til Kjendsgjerningernes Sprog, det kategoriske Magtsprog,
svarer et System af \emph{interpreterende} Functioner.
De interpreterende Functioner have en ved første Øiekast paafaldende
Lighed med de explicative. Ved begge gjælder det
om en Gjenkjendelse af det i Frembringelsen udtalte Begreb,
en Bedømmelse af det Frembragte, en Udvikling af dets Betydning;
og dog kan Forskjellen mellem explicative og interpreterende
Functioner udhæves med kategorisk Eftertryk.

Ved den explicative Function bliver den frembragte
Gjenstand nærmest maalt med Frembringelsens iboende
% --- p. 172 ---
Be\opage{172}greb: „dette Hjul er slet gjort“; „dette Hjul er godt gjort“;
„saaledes skal et Hjul være“. Explicationen motiverer alle
Domme ved Begrebet som Begreb; den holder sig til Frembringelsesbegrebet,
ikke til den frembringende Subjectivitet.
Ved den interpreterende Function gaaer Sagen omvendt.
Frembringelsesbegrebet kunde, som ovenfor viist, ikke være
Frembringelsen iboende, dersom det ikke tillige var oversvævende.
Det oversvævende Begreb, Subjectivitetens egentlige
Meningsbegreb, er afgjørende, og det i dobbelt Forstand, for
Interpretationen.

For det Første maa den frembringende Subjectivitet
nødvendigt anticipere sin Opgaves Løsning saavel efter objective
som subjective Hensyn; det ligger da i Sagens Natur,
at Betingelserne i det Objective kunne paa mange Maader
ændre og indskrænke Udførelsen af den oprindelige Plan.
Her viser sig en Differens mellem det Opnaaede og det Tilsigtede,
mellem Frembringelsens iboende og oversvævende Begreb,
mellem den i Værket udtalte Tanke og Værkmesterens
egentlige Mening. Fra dette Synspunkt kunne Naturfrembringelserne
i deres umiddelbare Endelighed vel endog siges
at være i et vist Misforhold til den guddommelige Mening,
et Misforhold, som den altformaaende Subjectivitet dog igjen
ophæver ved guddommelig Interpretation.

Er der nu fra Værkmesterens Side en endelig Differens
mellem det fuldførte Værk og det oprindelig tilsigtede, da
maa --- og dette er det næste Hovedpunkt --- denne Differens
jo faae forhøiet Betydning for den udenfor Staaende, der
lægger an paa at forstaae Frembringelsens factiske Sprog og
læse Værkmesterens Mening ud af Værket: fra dette Standpunkt
er Forskjellen mellem de blot explicative og de interpreterende
Functioner umiskjendeligt. Gjennem de explicative
Functioner har den betragtende Subjectivitet ligefrem med
Forholdet mellem Værket og dets Begreb at gjøre, Hensynet
til Værkmesteren er et Bihensyn; ved de interpreterende
% --- p. 173 ---
Func\opage{173}tioner vil Betragteren derimod i Værket selv tyde Værkmesterens
Tanker, fatte hans oprindelige Mening; saaledes
faaer han da gjennem Værket nærmest med Værkmesterens
Subjectivitet at gjøre. --- Her er et Vendepunkt, hvor de
interpreterende Functioners Betydning for Teleologien ret maa
falde i Øinene. Hvorfor har vel den videnskabelige Bevidsthed
saa vanskeligt ved at blive enig med sig selv angaaende
Naturens Formaal, den i Naturen sig aabenbarende Viisdom?
Aabenbart fordi to stærke Hensyn her komme i Strid med
hinanden: Hensynet til det Objective, thi Naturen er fra
Først til Sidst objectiv, og Hensynet til det Subjective, thi
Teleologien er væsentlig subjectiv: Midler og Formaal ere
kun til for en dømmende Subjectivitet. De, som udelukkende
fæste Blikket paa det Objective, maae da consequent forkaste
Teleologiens Betydning for Naturerkjendelsen; de, som eensidig
fæste Opmærksomheden paa det Subjective, udsætte sig
for at underskyde Naturfrembringelserne deres egen individuelle
Mening, en Mening, som saa igjen kommer i Strid med det
objective, Frembringelserne iboende, Begreb. Nærmere seet er
Striden begrundet i Teleologiens eget Væsen; thi Teleologien
vil paa ingen Maade give Slip paa det Objective; den vil
kun lægge det ind under den overgribende Subjectivitet. Saalænge
Adskillelsen mellem det Objective og det Subjective fastsættes
kritisk: enten et Subjectivt eller et Objectivt, vil Teleologiens
Problem aldrig kunne løses. Dette har Kant følt,
da han, istedenfor at gjøre Teleologiens Princip til et constitutivt
Princip, gjorde det til et regulativt: de interpreterende
Functioner skinne igjennem, men Interpretationen er endnu
kun psychologisk subjectiv.

Videnskaber som Botanik og Zoologie kunne ikke undlade
at anvende Kategorier af typisk-teleologisk Oprindelse; men
saalænge Subjectivitetens Ontologie ikke er gjennemklaret, vil
det være umuligt at afgjøre, hvorvidt det teleologiske Princip
er „constitutivt“ eller blot „regulativt“. Man forlanger en
% --- p. 174 ---
\opage{174}objectiv Viden af det Subjective: dette er Modsigelsen; man
anvender indicative Functioner, hvor man absolut maa bruge
de interpreterende: dette er Uklarheden.

„Fastholde vi“ --- for at anføre et Exempel af Naturlæren
--- „den Forudsætning, at der ligesaavel med Hensyn
paa de elementære som paa de organiske Heelheder maa skjelnes
mellem det Fornuftige og det Meningsløse: saa falde Conseqvenserne
som af sig selv. Ligesom et Dyrs Organisation
nødvendigviis maatte blive meningsløs, dersom dets Organer
indbyrdes stode i Strid med hinanden, dersom det f.\ Ex.
havde Tænder som et Rovdyr, men Lemmer som et Hovdyr,
saaledes maatte en Totalitet som Solsystemet ogsaa blive
meningsløs, dersom Planeterne, istedenfor at bevæge sig i en
lukket Linie omkring Solen, kom til at beskrive Curver, hvori
de fjernede sig i det Uendelige fra den \dots\ Ja, dersom man
--- forudsat, hvad da heller ikke uden videre kan beroe paa
tilfældige Omstændigheder, at Banen absolut maatte være et
Keglesnit --- kun havde Muligheder for Ellipse og Parabel,
saa kunde man endda lade Tilfældet raade; thi, det forstaaer
sig, medens Sandsynligheden for en Parabel er uendelig lille,
efterdi der blandt de uendelig mange mulige Tilfælde kun kan
gives et eneste for Parablen gunstigt, er Sandsynligheden for
Ellipsen, der har en uendelig Mængde gunstige Tilfælde, saa
afgjort overveiende, at selv det Tilfældige i denne Sammenhæng
maatte tage sig ud som det i sig Fornuftige. Men
nu gives her som bekjendt ikke blot et Valg mellem Ellipse og
Parabel, men tillige, hvad særlig kommer i Betragtning, et
Valg mellem Ellipse og Hyperbel; men en Hyperbel gjør
Systemet ligesaa meningsløst som en Parabel. Gaae vi altsaa
tilbage til Begyndelsen, vise sig Alternativerne: enten en
Ellipse eller en Hyperbel, det vil sige: enten en fornuftig eller
en meningsløs Totalitet; her er en væsentlig Forskjel mellem
Tilfældets og Fornuftens Herredømme“.\footnote{Forelæsn. over „Phil. Propæd.“ Universitetsaaret 1861--62.
S. 96--101.} % note *)

% --- p. 175 ---
\opage{175}Her er i Udtydningen af Phænomenet udtrykkelig lagt
an paa at give Beviset for det Teleologiske en saa objectiv
Charakteer som muligt, og dog er det Objective i Beviisførelsen
netop det Utilstrækkelige. Beviset er væsentlig subjectivt,
bygget nemlig paa en \emph{Fortolkning}. Skulde Beviset være
fuldstændig objectivt, da maatte det netop have været Parablen,
ikke Ellipsen, der svarede til det i sig Fornuftige. Beviset
kom da til at lyde saaledes: „Parablen er, forsaavidt Banen
afhænger af tilfældige Omstændigheder, den allerusandsynligste
Curve; da nu Parablen ikke desto mindre er det eneste Fornuftige,
saa sees heraf, at det Fornuftige ikke kan forklares
af blinde Naturaarsager, men maa henføres til en guddommelig,
i Tingenes Gang indgribende Viisdom“.

Hvad fordrer man altsaa, for at Beviset kan være objectivt?
Det Meningsløse! Først skal der sættes Splid
mellem Naturaarsagerne og det Fornuftige; dernæst skal den
guddommelige Subjectivitet benytte Spliden, for paa endelig
Maade at gribe ind og lægge sin Viisdom for Dagen; man
fordrer, med andre Ord, en Splid mellem de indicative og
de imperative Functioner i Guddommen selv: til Viisdommens
Aabenbarelse fordrer man Viisdommens Selvmodsigelse.

Det videnskabelige Grundlag for al Naturteleologie er
den Erkjendelse, at de objective Frembringelser ere Frembringelser
af en guddommelig Subjectivitet; at den guddommelige
Subjectivitet i disse Frembringelser udtaler evige Tanker
med Præg af guddommelig Mening. Hver Gjenstand er et
Magtens Ord, hvert Led en Stavelse, et Element, en Tøddel % UNSURE: Tøddel or Toddel (ø not distinguishable in scan; layer "Toddel")
i det System af praktisk-antilogiske Sætninger, der udgjøre
det kategoriske Magtsprog. Hvorvidt en menneskelig Fortolkning
er istand til at træffe den oprindelige Mening i dette
Virkelighedens kategoriske Magtsprog, bliver et Spørgsmaal
for sig --- her kan i al Fald kun være Tale om en Tilnærmelse;
men dette er det Absolute, at der i Tingene er
en Mening, en guddommelig Mening, at den guddommelige
% --- p. 176 ---
\opage{176}Mening punktlig svarer til en guddommelig Fortolkning, at
de interpreterende Functioner for Ideen selv have ontologisk
Gyldighed.

\begin{center}{\bfseries γ)\quad Dommens Gyldighed: Functionernes Modalitet.}\end{center}

Under hele den foregaaende Behandling af Dommen, saavel
med Hensyn paa dens Form som paa dens Gjenstand, er
det naturligviis bestandig forudsat, at Dommen som Dom
maatte have Gyldighed. Kun den gyldige Dom kan Logiken
anerkjende; den ugyldige maa den forkaste. Naar her desuagtet
i det Følgende særlig bliver Spørgsmaal om Dommens
Gyldighed, da staaer Gyldighed ikke umiddelbart imod Ugyldighed;
men Gyldighed af en Art staaer imod Gyldighed af
en anden. Forsaavidt der dømmes med fuld Bevidsthed, vil
der ved en Dom om Gjenstanden altid foresvæve den Dømmende
en Dom om Dommens Gyldighed. Denne Dom vil,
i Forhold saavel til Gjenstandens som til Erkjendelsens Beskaffenhed,
være forskjellig i væsentlig forskjellige Tilfælde; at
undersøge denne Forskjellighed er, med andre Ord, at undersøge
\emph{Functionernes Modalitet}.

At der til alsidig Belysning af Tankens Virksomhed
gjøres Adskillelse mellem forskjelligartede Functioner, maa ikke
forstaaes, som om de forskjelligartede Functioner havde særlig
Selvstændighed og en egen Tilværelse i Bevidstheden ved
Siden af hverandre; tvertimod! de ere forskjellige Former,
hvorunder Tankevirksomheden fremtræder, Maader, hvorpaa
den lægger sig for Dagen. At Tænkningen ved de forskjellige
Virkemaader maa blive forskjellig, er en Selvfølge; det
er jo, som man siger, Maaden hvorpaa det i alle Ting kommer
an. Men at den ved sine forskjellige Virkemaader forskjelligartede
Tænkning desuagtet er een og samme Tænkning, viser
tillige, at Maaden ikke er Sagen selv, at det væsentlig kommer
an paa Tænkningen, ikke paa Maaden. Den Modsigelse, at
% --- p. 177 ---
\opage{177}Maaden paa een Gang baade er væsentlig og dog uvæsentlig,
er et slaaende Beviis paa Tænkningens Substantialitet; thi
kun i Tænkningens Substantialitet finder Modsigelsen sin
Løsning.

I de imperative, indicative og interpreterende Functioner,
f.\ Ex., er Tænkningen den samme og dog ikke den samme; det
vil sige: Tænkningen bliver i de forskjellige Functioner væsentlig
forskjellig fra sig selv. Forskjellen er saa gjennemgaaende,
at det seer ud, som om hvert Functionssystem maatte
have sit eget Princip; og dog er her kun eet Princip: Tænkningens
Væsen er i Principet eet og samme Væsen. Gjennem
en væsentlig Selvadskillelse formaaer Tanken at udvikle og
fastholde sin væsentlige Eenhed; Tænkningens ideelle Magt
over sine Modsætninger giver Eenheden Substantialitet: jo
dybere Modsætningerne ere, desto stærkere er den substantielle
Eenhed. Da nu Tænkningen med sin logiske Substantialitet
ikke holder sig bagved, men gaaer op i sine Functioner, maae
Functionerne trods deres indbyrdes væsentlige Forskjellighed
med det Samme gaae op i Tænkningens Eenhed.

I Forhold til hinanden ere, for at blive ved Exemplet,
de imperative, de indicative og de interpreterende Functioner
i Grund og Virksomhed forskjellige. Hvad der udrettes ved det
ene System, kan ikke udrettes ved det andet; det Forhold,
hvori den altformaaende Subjectivitet sætter sig til Objecterne i
Kraft af de imperative Functioner, er i Virkeligheden et ganske
andet, end det, der bestemmes ved Systemet af de indicative.

Men med Hensyn paa Tænkningens substantielle Eenhed
er det dog saa langt fra, at Functionerne af de forskjellige
Systemer skulde udelukke hinanden eller staae i Strid
med hinanden, at de tvertimod, indeholdes i hinanden, gaae
paa immanent Maade over i hinanden og harmonere. Saaledes
gik jo de anticiperende Functioner gjennem de executive
over i de explicative, disse igjen over i hine; saaledes kom de
imperative Functioner til Ro i de indicative, idet den af
% --- p. 178 ---
\opage{178}begge Systemers Forskjelseenhed betingede Mening fandt sin
Opklaring i de interpreterende Functioner.

Af Functionernes gjensidige Forhold sees, at deres relative
Selvstændighed er Eet med deres Modalitet: til Tænkningens
Substantialitet svarer Functionernes Modalitet; til
Functionernes Modalitet svarer igjen Dommens Gyldighed.

I den formelle Logik bliver Dommens Gyldighed afhandlet
under „Dommenes Inddeling efter deres Modalitet;“
men Modaliteten selv bliver ikke undersøgt. Hos Kant har
Modaliteten en blot subjectiv Betydning; hos Hegel behandles
den deels som Maal, deels som Modus af det Absolute\footnote{Jvfr. Die objektive Logik. Erstes Buch. Dritter Abschnitt. S. 395:
\emph{Das Maaß}. Zweites Buch. Dritter Abschnitt. S. 191: \emph{Der
Modus des Absoluten}. I 2den Bogs 3die Afsnit. S. 192
udtaler Hegel sig blandt Andet saaledes: „Der Modus, die
\emph{Aeußerlichkeit} des Absoluten, ist \dots\ die als Aeußerlichkeit \emph{gesetzte}
Aeußerlichkeit, eine bloße \emph{Art} und \emph{Weise}; somit der Schein
als Schein, oder die \emph{Reflexion der Form in sich}; somit die
\emph{Identität mit sich, welche das Absolute ist}. In der
That ist also erst im Modus das Absolute als absolute Identität
gesetzt; es ist nur, was es \emph{ist}, nämlich Identität mit sich, als sich
auf sich beziehende Negativität, als \emph{Scheinen}, das \emph{als Scheinen}
gesetßt ist“.} % note *) % PRINTER: printed "gesetßt", sense "gesetzt"
udelukkende fra den objective Side; det skal nu vise
sig, at Modalitetens Betydning for Dommen og dens Functioner
baade er objectiv og subjectiv.

Modalitetens subjective Betydning med Hensyn paa
Dommens Gyldighed fremlyser umiddelbart af Adskillelsen
mellem assertoriske, problematiske og apodiktiske Domme, forsaavidt
Adskillelsen udelukkende bestemmes ved den dømmende
Subjectivitets tilsvarende Forvisning.

Saalænge man nu eensidigt holder sig til det Subjective,
kan det Assertoriske, det Problematiske og det Apodiktiske paa
mangfoldige Maader løbe imellem hinanden: den Ene kan for sin
% --- p. 179 ---
\opage{179}Part have Grund til at yttre sig problematisk over det, hvorom
den Anden ligefrem dømmer assertorisk og omvendt; ja, saa
dunkel kan Subjectiviteten være, at den Dømmende indbilder
sig at vide med apodiktisk Vished, hvad han dog kun besidder
i Form af Paastande og umiddelbare Forsikkringer.

Anderledes, naar Sagen sees fra den objective Side.
Herved faaer Dommens Materie, snart i Form af Gjenstand,
snart i Form af Indhold, særegen Betydning. Det, hvorom
der dømmes, kan nemlig deels være af den Beskaffenhed, at
det ene og alene egner sig for en assertorisk Dom, medens et
Andet, ifølge sin Natur, udtrykkelig svarer til en problematisk,
et Andet igjen til en apodiktisk. Er det en Feil at bruge den
assertoriske Dom i en Sag, der fordrer en apodiktisk, da er
det ogsaa en Feil at fordre en apodiktisk Dom, hvor der kun
kan og skal gives en assertorisk.

Naar det „dogmatiske Talent“ forsikkrer os, at Troessandheden
er ham „det Visseste af Alt“, da udtaler han sin
Overbeviisning i en assertorisk Dom, og herimod kan i og
for sig Intet indvendes; men naar saa samme „Talent“, af
inderlig Trang til Speculationen, erklærer, at Troeslæren alvorlig
trænger til videnskabelig Begrundelse, at dens Indhold,
med andre Ord, egner sig for den apodiktiske Dom: saa løber
det rundt for „Talentet“. Thi af hvad Art er nu Forvisningen?
Giver Troen som Tro den fulde Forvisning, saa
er Vidnesbyrdet det Høieste, og den Troende udtaler sig da
i assertoriske Domme. „Sandelig, sandelig siger jeg Eder“,
en saadan Indgang bebuder just det Assertoriske. Kan det,
hvad Forvisningen angaaer, ikke være Dogmatikeren nok med
Troen som Tro, uagtet Troessandheden er ham „det Visseste
af Alt“, saa maa han ved videnskabelig Begrundelse altsaa
ville søge en endnu høiere Forvisning, en Forvisning, hvorved
den i Troen erkjendte Sandhed kan blive ham endnu
vissere end det Visseste af Alt, en Søgen, der uimodsigelig
er meningsløs.

% --- p. 180 ---
\opage{180}Det Meningsløse i den foregivne Søgen bliver endnu
mere iøinefaldende, naar vi betænke, at Veien fra Troens
assertoriske til Videnskabens apodiktiske Dom gaaer igjennem
det Problematiske.

Naar Anselm med sit \textit{credo, ut intelligam} bygger paa
det Assertoriske, saa peger Cartesius med sit \textit{de omnibus
dubitandum est} udtrykkelig hen paa det Problematiske. Er
den ene af disse Anviisninger sand, saa maa den anden, hvad
Troeserkjendelsen angaaer, aabenbart være falsk. Har Anselmus
Ret, naar han siger: \textit{fides præcedit intellectum}, saa maa
Cartesius nødvendigt have Uret, naar han siger: \textit{dubitatio
præcedit intellectum}; medmindre den speculative Dogmatik
vil opfatte Forholdet mellem Troen og Tvivlen paa samme
dybsindige Maade, som den opfatter Forholdet imellem Skabelse
og Kosmogonie, saa at Troen væsentlig er Tro, fordi den er
Tvivl, ligesom Skabelsen væsentlig er Skabelse, fordi den er
Kosmogonie, og Tvivlen omvendt Tvivl, fordi den er
Tro, ligesom Kosmogonien er Kosmogonie, fordi den er
Skabelse.

Det være nu med Speculationen, som det vil; vi skulle
i det Følgende vise Modalitetens objective Betydning, idet vi
skulle vise, at den væsentlige Forskjel mellem assertorisk, problematisk
og apodiktisk Viden oprindelig afhænger af Forskjellen
imellem det umiddelbart Bestemte, det Ubestemte og det i Reflexionen
Bestemte. Til det Assertoriske svare Functioner af
umiddelbar Bestemthed, til det Problematiske Functioner af
ophævet Bestemthed, til det Apodiktiske Functioner af reflecteret
Bestemthed.

\textit{αα)} \emph{Den assertoriske Dom: Functioner af
umiddelbar Bestemthed.} At der gives en umiddelbar
Bestemthed, der ikke lader sig opløse i Reflexion, er i vort
Foregaaende blevet fremhævet hver Gang vi fra en eller anden
Side have havt Anledning til at omhandle det Antilogiske.
Det Antilogiske har sin Existensgrund i sig selv, slaaer igjen%
% SEAM: word continues on p. 181 as "nem"; written here in full as "igjennem"
% --- p. 181 ---
% SEAM: p. 181 begins with the tail of "igjen-nem" (p. 180 ends "igjen-"); this batch starts with "nem"
\opage{181}nem ved Magt, er, fordi det er, og lader sig forsaavidt ikke
begrunde. De sandselige Qvaliteter i deres existentielle Umiddelbarhed
lade sig ikke begrunde: Ingen Fornuftig kan spørge,
hvorfor det Røde er rødt, det Haarde haardt, det Tunge
tungt, saalidt som Nogen kan spørge, hvorfor det Lige er liigt,
Høiheden høi og Dybden dyb. Der er i al Reflexion Reflexer
af umiddelbar Bestemthed, og i al umiddelbar Bestemthed
Noget, der maa opfattes ligefrem saaledes, som det viser
sig: til enhver saadan Opfattelse svarer en assertorisk Dom.
Det er charakteristisk, at den assertoriske Dom i sin Form
har en Anviisning paa Virkeligheden, at denne Anviisning
umiddelbart falder sammen med en Forsikkring, og Forsikkringen
med den umiddelbare Vished: S er virkelig (det forsikkrer jeg)
P; den virkelige Bestemthed er nemlig ikke blot umiddelbar,
men i sin Umiddelbarhed antilogisk, og ingen antilogisk Bestemthed
kan give anden Forvisning end den, der udtrykkes i
den assertoriske Dom.

Den assertoriske Dom er logisk-antilogisk i særegen
Forstand. Dens Form er logisk, dens Indhold baade logisk
og antilogisk; den giver os i Virkeligheden altid et antilogisk
Indhold i logisk Form, men den logiske Form giver igjen
Indholdet et logisk Præg. „Stenen er virkelig graa“; denne
Dom synes ved første Øiekast at være logisk saavel i Indhold
som Form, men see vi nærmere til, opdage vi let, at
det logiske Indhold kun er en i Formen reflecteret Bestemthed:
det virkelige Indhold er i sin Virkelighed en existerende Gjenstand
med en existerende, d. e. med en Existensen inhærerende
Egenskab. Prædikatet „graa“ betegner her aldeles ikke det
Almindelige, men tvertimod det Enkelte. En almindelig Dom
som: „Stenen er graa“, kan være et Udtryk for, at Tingen,
det Enkelte, er subsumeret under en Egenskab (det Almindelige);
men den assertoriske Dom: „Stenen er virkelig graa“, passer
ikke for slige Betegnelser. Den, der siger: „det Enkelte er
virkelig det Almene“, er tankeløs; thi i Virkeligheden er det
% --- p. 182 ---
\opage{182}Virkelige aldrig det Almene, og det af den gode Grund, at
det Virkelige som Virkeligt altid er et umiddelbart Bestemt,
og det umiddelbart Bestemte ikke det Almene. Den assertoriske
Dom: „Stenen er virkelig graa“, sigter ikke paa Graat i Almindelighed,
men paa det Graa i denne umiddelbare Bestemthed,
det virkelige Graa, den graa Steen.

Da nu den assertoriske Doms Indhold altid maa være af
en dobbeltsidig Charakteer, \emph{antilogisk} formedelst Gjenstanden,
\emph{logisk} formedelst Formen, saa følger det af sig selv, at Udviklingen
af Dommens Gyldighed anlægges med særligt Hensyn
til hver af Hovedsiderne og deres indbyrdes Forhold.
Saaledes komme vi da til særskilt Behandling af det Antilogiske
eller af Modalitetens objective Side, af det Logiske
eller af Modalitetens subjective Side, og tilsidst af Indholdets
logisk-antilogiske Væsenhed, hvorved Modalitetens Gyldighed
alsidig bestemmes.

1) \emph{Det antilogiske Element i den assertoriske
Dom: Modalitetens objective Side}
% PRINTER: no full stop after "Side"
Den assertoriske
Dom omfatter en Sphære for sig, en Sphære af uendeligt
Omfang, nemlig den hele Virkelighed fra Magtbestemmelsernes,
de antilogiske Elementers Synspunkt. Vistnok har Virkelighedens
Verden sine Væsenheder, sine Grunde, sine Muligheder,
Love, Arter, Slægter o. s. v., Kategorier, der i sidste
Instans unddrage sig den assertoriske Videns Umiddelbarhed;
men her er i denne Sammenhæng udtrykkelig Tale, ikke om
Væsenheder i det Virkelige, men om Virkelighedens virkelige
Kjendsgjerninger i deres givne Forskjellighed og uendelige
Mangfoldighed. Vel maa det indrømmes, at de forskjellige
Virkelighedsbestemmelser med deres antilogiske Umiddelbarhed
uophørlig vexle, ændres og forvandles; det Nærværende forsvinder,
det Tilværende forandres, det Existerende gaaer til
Grunde, o. s. v.; men i alle Omvexlinger bliver det ene
antilogiske Element Punkt for Punkt afløst af det andet. Det
Nærværende, der forsvinder, giver kun Plads for et andet
% --- p. 183 ---
\opage{183}Nærværende; ethvert Nærværende er nu virkelig \emph{saaledes},
ikke anderledes; det bestemmer den assertoriske Dom og bestemmes
igjen ved denne. Det Existerende forgaaer, det
gaaer til Grunde; men Grunden danner derfor ingen Kløft
mellem det Existerende, der forgaaer, og det Existerende, der
opstaaer; paa hvert Punkt i den virkelige Overgang følger
Existerende paa Existerende, overalt \emph{dette}! og \emph{dette}! og
\emph{dette}! Her er en Strøm af antilogiske Elementer, der ikke
kjender nogen Afbrydelse: der er i Virkelighedens Væg ingen
Sprækker, ingen Revner.

Dersom Sandhed var umiddelbart Eet med Virkelighed,
og Virkeligheden umiddelbart Eet med Totaliteten af de antilogiske
Bestemmelser, da vilde den assertoriske Dom være
Sandhedens logisk adæqvate Udtryk; nu derimod, da Virkelighedsbestemmelsen
kun er et Moment, om end et Realmoment,
i det virkelige Sande, kan den assertoriske Dom kun angive
en Side af Sandheden, ikke Sandheden selv, den hele Sandhed.
Det er en Vildfarelse, naar den umiddelbare Virkelighedsbestemmelse
uden videre gjøres til Eet med Sandheden;
men det er ogsaa en Vildfarelse, naar Virkelighedsbestemmelsen
i sin Umiddelbarhed forflygtiges og forvandles til Skin, fordi
den ikke er Sandheden selv. Forholdet mellem Sandsning
og Tænkning vil i denne Forbindelse kunne belyses fra et nyt
Synspunkt.

Sophisten Protagoras gik, som bekjendt, ud fra den
Hovedsætning, at „Mennesket er alle Tings Maal, baade de
værendes, at de ere, og de ikkeværendes, at de ikke ere“. Det
er et eget Forhold mellem Sandsning og Tænkning, mellem
Dom og Gjenstand, som herved kommer tilsyne. At Protagoras
gaaer ud fra det Sandselige og oprindelig lægger
Eftertryk paa Sandsefornemmelsen, er vist; ligesom det ogsaa
er umiskjendeligt, at han sophistisk nedsætter Tænkningen og
undergraver al Sandhedserkjendelse. Men det er ogsaa vist,
at Plato, hvormeget han end gjør Tænkningens Ret gjældende,
% --- p. 184 ---
\opage{184}dog ikke ganske har tænkt sig ind i det, som her ligger til
Grund for den sophistiske Synsmaade.

„Skulde Mennesket“, hedder det i Theætet, „blot som
sandsende Individ være Tingenes Maal, da kunde et Dyr
ligesaa fuldt som det viseste Menneske, ja som en Gud, erkjende
det Sande (ὅτι πάντων χρημάτων μέτρον ἐστὶν ὗς ἢ
κυνοκέφαλος ἢ τι ἄλλο ἀτοπώτερον τῶν εχόντων
ἄισθησιν).
% PRINTER: printed "εχόντων", sense "ἐχόντων"; printed "ἄισθησιν", sense "αἴσθησιν"; the quotation opened before "blot som" is never closed

Ved denne Indsigelse ere to meget væsentlige Punkter
overseete. For det Første er Forholdet mellem det Sandselige
og det Virkelige og derhos mellem det Sandselige og det Sande
overseet. Idet Plato miskjender det Sandselige, maa han,
som i første Kapitel\footnote{Grundideernes Logik. I. S. 182--183.} % note *)
er viist, nødvendig miskjende Virkeligheden,
og, idet han miskjender Virkeligheden, miskjende Sandheden.
For det Andet er Forholdet mellem Dom og Gjenstand
og derved Forholdet mellem en menneskelig og en dyrisk
Sandsning overseet. Idet Protagoras gjør Mennesket, ikke
Svinet, ikke Varulven, ikke Frøen, til Tingenes Maal, opfatter
han dog, være sig med klar eller uklar Bevidsthed,
Mennesket udtrykkelig som dømmende Subjectivitet. Naar den
Ene siger: „Vinden er virkelig kold“, den Anden: „Vinden er
virkelig varm“, da er det to assertoriske Domme, der frembringe
Modsigelsen, og derved fremkalde Spørgsmaalet om
det Sande. Den assertoriske Dom er paa denne Maade
formedelst Sandseindtrykket sat, ikke i et tilfældigt, men i et
væsentligt Forhold til Spørgsmaalet. Ved blot at stille
Svinet imod Frøen, Fuglen imod Fisken, kan Sandhedsspørgsmaalet
slet ikke fremkomme: Spørgsmaalet fremkommer
udtrykkelig ved Modsigelsen, Modsigelsen ved den assertoriske
Dom; men Svinet, Frøen, Fuglen, Fisken ere kun sandsende,
ikke dømmende Subjectiviteter.

Først naar Forholdet mellem Virkeligheden og den
% --- p. 185 ---
asser\opage{185}toriske Dom er opfattet klart, kan Protagoras's Sophistik
tilfulde belyses. Det er Vexelbestemmelsen af Gjenstand og
Dom, hvorpaa Eftertrykket falder. Idet her begyndes med
det Sandselige, begyndes her umiddelbart med det Virkelige:
„Vinden er kold, den er virkelig kold“; ved et saadant Udsagn
tænker den dømmende Subjectivitet aldeles ikke paa at
foreskrive Vinden, om den skal være kold eller varm, om den
skal \emph{være} til eller ikke. Den Dømmende vil ikke indvirke paa
Vinden, men erkjender sig paavirket af Vinden; det er, med
andre Ord, ikke Dommen, der bestemmer Gjenstanden, men
omvendt Gjenstanden, der umiddelbart bestemmer Dommen:
Dommen er correct assertorisk. Hvori ligger saa det Sophistiske?
I en svigagtig Brug af Dommen, thi Dommen
skjuler en Modsigelse. Den Dømmende paavirkes af den
virkelige Gjenstand; men han paavirkes kun paa den Maade,
der nøiagtig svarer til hans Individualitet: idet han assertorisk
vil dømme om det Paavirkende, kommer han uvilkaarlig
til at dømme om den Paavirkning, han modtager, derved om
sin egen subjective Tilstand, altsaa bogstavelig om sig selv.
Det er just det Charakteristiske ved den assertoriske Dom, at
den fra at være umiddelbart objectiv, pludselig bliver umiddelbart
subjectiv, at den fra at være en Dom om Gjenstanden
forvandler sig til en Dom om den Dømmende, den Dømmendes
Dom om sig selv.

Den protagoriske Sophistik har fortræffelig blottet Modsigelsen,
men ikke løst den. Istedenfor at erkjende det Virkelige
som et Moment i det Sande og søge det Sande i
Loven, bliver Sophisten staaende ved Virkelighedens antilogiske
Punktualitet, og benytter nu den Modsigelse, hvori Individerne
gjennem assertoriske Domme lægge deres tilfældige Subjectivitet
for Dagen, til et Spil med Sandheden.

Fra en anden Side og paa en anden Maade er Forholdet
mellem Virkeligheden og den assertoriske Dom bleven
miskjendt af Epikuræerne. Ogsaa i Epikurs Kanonik er
% --- p. 186 ---
\opage{186}Sandhedens Kriterium Fornemmelsen eller Sandsningen
(ἄισθησις).
% PRINTER: printed "ἄισθησις", sense "αἴσθησις"
„Der frembringes en Bevægelse i os, som er
Sandsningen, og hvor der er en Bevægelse, maa være et
Bevægende. Derfor ere alle Sandsninger sande; thi Tingene
bevæge os virkelig saaledes, som vi sandse dem. Ogsaa den
Vanvittiges og den Drømmendes Syner ere sande; thi der
tilføres dem virkelig saadanne Billeder. Orestes saae virkelig
Eumeniderne, men han feilede deri, at han antog dem for
palpable Gjenstande“.\footnote{Efterladte Skrifter af Poul Møller. 4de Bd. 3die Udg. Kjøbenhavn 1856. S. 256.} % note *)

Men trods den tilsyneladende Overeensstemmelse mellem
Epikur og Protagoras, hvad Udgangspunktet angaaer, er
Epikur dog kjendelig gaaet et Skridt videre. Fremskridtet er
betegnet derved, at her lægges Vægt paa Forskjellen mellem
Virkning og Aarsag, mellem det Objective og det Subjective
i Tilsyneladelsen. Havde de Eumenider, Orest saae, virkelig
været palpable, saa havde de ogsaa havt Virkelighed for Andre
end Orest; de havde, med andre Ord, havt objectiv Virkelighed.
Allerede heraf skjønnes, hvilke Betingelser Virkeligheden
selv tilbyder til nærmere Bestemmelse af det objectivt Gyldige.
Istedenfor at overile sig med Slutninger fra hvad der sandses
til hvad der ikke sandses, opfordres Betragteren til at lade den
ene Art Sandsning understøtte den anden, indtil han med de
forskjellige Sandsers forenede Vidnesbyrd har saamange Kjendemærker
paa det Sandsedes Objectivitet, at han med fuld
Forvisning tør sige: det Sandsede er virkelig en tilværende
Gjenstand.

Den Maade, --- thi ved Modaliteten kommer jo Alt an
paa Maaden --- hvorpaa forskjelligartede Sandsninger correspondere
ved Opfattelsen af een og samme objective Eenhed,
taler i Virkeligheden for Eenhedens substantielle Væren, dens
Uafhængighed af den sandsende Subjectivitet; dette
% --- p. 187 ---
Virkelig\opage{187}hedens Vidnesbyrd bliver endnu nærmere bekræftet derved, at
flere Sandsende iagttage det samme Sandselige under samme
Betingelser paa samme Maade. Gjøres nu tillige den tilværende
Gjenstands Relationer til andre tilværende Gjenstande,
altsaa Tingenes indbyrdes Vexelvirkning, til Gjenstand for
umiddelbar Iagttagelse, saa har Virkeligheden ydet den betragtende
Subjectivitet alle de Bidrag til Fastsættelse af det
Objective og Sande, som den igjennem Totaliteten af umiddelbare
Indtryk er istand til at yde. Men hvorledes disse
Bidrag end for Resten opfattes og combineres, saa er det,
naar vi lægge tilbørlig Vægt paa de antilogiske Elementer,
dog indlysende, at det logiske Grundlag, hvorpaa Forvisningen
hviler, er en Totalitet af assertoriske Domme. Det Assertoriske
kjendes just derpaa, at den ene Dom har samme Oprindelighed
som den anden; her er i denne Henseende ingen
Overgang fra Assertorisk til Apodiktisk.

Den slaaende Umiddelbarhed, hvormed virkelige Indtryk
af virkelige Gjenstande bestemme den assertoriske Dom, forklarer
den Sikkerhed, hvormed den dømmende Subjectivitet,
trods alle humeske Raisonnementer, orienterer sig i den virkelige
Verden. Lade vi her Eftertrykket falde paa Tingenes
Objectivitet, deres objective Forskjelligheder og objective Forhold,
saa have vi Empirismens Standpunkt. Den empiriske Dom er
væsentlig en Indtryksdom, Indtryksdommen væsentlig assertorisk.
Men de til det Antilogiske bundne, ved Kjendsgjerningerne
sikkrede Indtryksdomme oversvæves tillige af en Reflexion,
der frigjør sig fra den antilogiske Punktualitet og i
logisk Forstand forlanger det Almene: den Dualisme, som
herved paatrænger sig Bevidstheden, lader os see Hume's
Misforstaaelse fra en ny Side.

Alle sande Begreber maae, dette er Hume's Grundtanke,
være Virkelighedsbegreber, altsaa oprindelige Impressionsbegreber.
Men hvorfra komme da de usande Begreber? Fra
den Forvanskning, Impressionsbegreberne undergaae i den
% --- p. 188 ---
paa\opage{188}følgende Reflexion! Først er altsaa Impressionen, saa følger
Reflexionen; først Sandseindtrykket, saa Dommen; først den
assertoriske Dom, saa den almeenlogiske. Her er en øiensynlig
Confusion af det Principielle med det Phænomenale. Hvad det
Phænomenale angaaer, kan Hume psychologisk have Ret; men i
Principet har han Uret. Efter Phænomenet maa det vel
synes, som om Sandseindtrykket allerede stod klart for Bevidstheden,
førend nogen Dom dannede sig; men i Principet er
det, som vi ovenfor have seet hos Protagoras, ligesaa oprindelig
Dommen, der bestemmer Indtrykket, som Indtrykket, der bestemmer
Dommen. Efter Phænomenet skulde det vel synes,
som om den assertoriske Dom, Dommen om denne enkelte,
bestemte Gjenstand i dette enkelte bestemte Tilfælde, gik forud
for den almeenlogiske Dom: Dommen om, hvad der gjælder
for alle mulige Tilfælde; men i Principet er det ligesaa oprindelig
den almeenlogiske Dom, Dommen som Dom, der
maa bestemme den assertoriske, som omvendt. Dersom der
altsaa feiles i Dommen om denne Substans med disse Accidenser,
da er Feilen ikke, hvad Hume paastaaer, den, at
Kategorierne Substans og Accidens lægges til Grund ved
Bedømmelsen af virkelige Gjenstande; denne Feil vedkommer i
al Fald ikke den assertoriske Dom. Hvad her vedkommer den
assertoriske Dom, er ene og alene Anvendelsen af de almindelige
Kategorier paa det enkelte Tilfælde. Er det usandt,
at A i et givet Tilfælde er Aarsag, B Virkning, da er
Feilen ikke, hvad Hume paastaaer, den, at Kategorierne: Aarsag
og Virkning lægges til Grund ved Bedømmelsen af noget
Virkeligt, men Feilen er, at den assertoriske Dom har i det
givne Tilfælde gjort en falsk Anvendelse af Aarsagskategorierne;
hvad enten den iøvrigt har anvendt dem paa en Relation,
hvori de ikke passe, eller forvexlet den virkelige Aarsag med
noget Andet, f. Ex. med Betingelser, Biomstændigheder o. s. v.

Ifølge Hume skulde det være den almeenlogiske Dom,
der havde den assertoriske til Grundlag; men det er omvendt
% --- p. 189 ---
\opage{189}den assertoriske Dom, der har den almeenlogiske til Grundlag.
Ifølge Hume skulde Virkelighedskategorierne godtgjøre deres
Realitet ved at udelukke Aarsagskategorien; men Tingen er
lige det Modsatte: det maa, endog ifølge Hume's egen Methode,
netop være Aarsagskategorien, der ligger til Grund for Virkelighedskategorierne.
Hvad er vel det ubedragelige Mærke paa
en sand Virkelighedskategorie? At den middelbart eller umiddelbart
har sin Oprindelse fra en tilsvarende Impression! Hermed
stemmer det saa, at den assertoriske Dom lægges til
Grund for den almeenlogiske, thi den assertoriske Dom med
sit antilogiske Element stempler just ethvert Virkelighedsbegreb
som Impressionsbegreb. Men herved forudsættes jo netop et
Forhold mellem Impression og Begreb, der for Begrebets
Vedkommende maa erkjendes for et ligesaa nødvendigt som
virkeligt Afhængighedsforhold: Impressionen forholder sig til
Begrebet som Aarsag til Virkning. Tænke vi Aarsagsforholdet
bort, staaer Begrebet i et tilfældigt Forhold til Impressionen;
kan et bestemt Virkelighedsbegreb ligesaa vel tænkes
at følge af den ene Impression, som af den anden; kan den
Impression, der i et Øieblik vækker Forestillingen Guult, i
et andet Øieblik vække Forestillingen Blaat, og den Impression,
der den ene Gang giver en Farve, den anden Gang give en
Tone: hvor er saa den Maalestok, der skal udvise Forskjellen
mellem sande og falske Begreber?

Det Sande i Hume's Anskuelse er det for Empirismen
og den assertoriske Dom charakteristiske Moment, at nemlig
enhver Existens i ethvert Tilfælde individualiserer og just derved
modificerer Begrebet, saa at hvert enkelt Phænomen i hvert
Øieblik nødvendig maa kræve sin umiddelbare Opfattelse, sin
særskilte Bedømmelse, dersom Erkjendelsen af det Virkelige
skal være ganske adæqvat. Spørgsmaalet om en adæqvat
Virkelighedserkjendelse staaer da i allernøieste Forbindelse med
den logiske Behandling af Alvidenhedsproblemet. I Forhold
til Virkelighedens antilogiske Punktualitet maa Alvidenheden
% --- p. 190 ---
\opage{190}være en collectiv Viden, en særlig umiddelbar Viden om alle
tilværende Ting i alle Tilværelsens Enkeltheder: „Ja, endog
alle Eders Hovedhaar ere talte“; disse Ord maa Logiken i
denne Sammenhæng vedkjende sig bogstaveligt.

„Men“, kunde Nogen indvende, „en Alvidenhed, der
tæller Hovedhaar, er os jo ligesaa ubegribelig, som en Almagt,
der ved sit blotte Ord skaber Verden af Intet; kan Videnskaben
vedkjende sig det religiøse Udtryk for Alvidenheden,
hvorfor skulde den da ikke ligesaavel vedkjende sig det religiøse
Udtryk for Almagten? Skal Aabenbaringens Ord om Almagten
afvises, fordi det er ubegribeligt, da maa det Samme jo
gjælde, naar Talen er om Alvidenheden“. Vi svare: Grundideernes
Logik gjør overalt Forskjel mellem det for os Ufattelige
og det i Ideen Ubegribelige. En Almagt, der ubetinget
skaber Verden af Intet, er ikke blot for os ufattelig, men i
Ideen ubegribelig; en Alvidenhed derimod, der umiddelbart
gjælder det Enkelte, er i Ideen begribelig, thi dens Begreb
er en nødvendig Conseqvens af Videns og Virkelighedens articulerede
Vexelbestemmelse.

Allerede paa et tidligere Stadium\footnote{Grundideernes Logik. I. S. 257--263.} % note *)
er Forholdet mellem
Alvidenhedskategorien og Vidensidealet særlig belyst med Hensyn
paa Sandsningens Ontologie\footnote{Jvfr. Grundideernes Logik I. S. 423--427.}. % note **)
I nærværende Sammenhæng
forhøies Belysningen derved, at den uendelige Mangfoldighed
af Sandsninger og Sandsebestemmelser erholde deres
ideelle Form i en tilsvarende Mangfoldighed af assertoriske
Domme. De assertoriske Functioner ere da ogsaa, hvad Modalitetens
objective Side angaaer, oprindelig beslægtede med
de explicative og indicative, og fremtræde som Modificationer
af begge.

Den assertoriske Function modificerer den explicative, idet
den med Eftertryk paa det Givne lader den stramme
% --- p. 191 ---
Mod\opage{191}sætning af Begreb og Existens, hvortil Explicationen støtter
sig, forsvinde i umiddelbar Eenhed. Lade vi Dommen: „dette
er virkelig en Green“, erholde sin Betydning fra de explicative
Functioner, da begynder det, som vi have seet, med, at Begrebet
afgiver Maalestok for Existensen, og ender med, at
den givne Existens anerkjendes for det virkeliggjorte Begrebs
Repræsentant: i dette Existerende er Begrebet Green virkeliggjort.
En saadan Middelbarhed,. en saadan Modsætten og
% PRINTER: printed ",.", sense ","
Mæglen mærkes ikke paa det assertoriske Standpunkt; vistnok
gjenkjendes Begrebet i det Existerende, men Begrebet veies
ikke som Begreb, Reflexionen er som et Blik: Gjenkjendelsen
er ganske umiddelbar.

Den assertoriske Function modificerer den indicative, idet
den med Eftertryk paa det Givne gjør den logiske Identitet
af Begreb og Existens kjendelig ved at antyde Modsætningen.
Fra det indicative Synspunkt svare de virkelige Existenser,
seete i Reflexionens Lys, til de evige Kategorier: Tilfredsstillelsen
er, at det Logiske klarer det Antilogiske og beriges
ved antilogiske Reflexer. Den fuldstændige Ligevægt mellem
det Logiske og det Antilogiske, der er de indicative Functioners
Resultat og Forudsætning, vilde bringe Ueensartetheden af
Begreb og Existens i Forglemmelse, dersom ikke den assertoriske
Function ved sit antilogiske Element gav det umiddelbart
Virkelige en Overvægt.

2) \emph{Det logiske Moment i den assertoriske Dom:
Modalitetens subjective Side.} Vibrerende mellem det
Logiske og det Antilogiske vil den assertoriske Dom bestandig
give det antilogiske Element Overvægt, saalænge den paagjældende
Gjenstand er udvortes objectiv. Er Dommens Gjenstand
derimod indvortes, og i sin Indvorteshed altsaa af subjectiv
Natur, vil Forholdet vende sig om: det logiske Moment
vil blive det overveiende; men denne Overvægt vil saa bevirke,
at det Assertoriske faaer en tvetydig Charakteer. Det Indre
er, som Indre, med al sin Umiddelbarhed dog altid et
% --- p. 192 ---
Re\opage{192}flecteret. Smerten er umiddelbar og dog i Reflexion; Nydelsen,
Frygten, Haabet, Glæden o. s. v. ligesaa. Her kan
ved disse Inderlighedsbestemmelser gives meget forskjellige Forbindelser
af Umiddelbarhed og Reflexion, men i dem alle er
Umiddelbarheden reflecteret, Reflexionen umiddelbar. Den,
der siger: „jeg er virkelig glad“, udtaler sig i en assertorisk
Dom; hans Sindstilstand er given, han har en umiddelbar
Forvisning om den, og meer end en umiddelbar Forvisning
om sin øieblikkelige Sindstilstand kan Ingen have.

Det Indre har sit Ydre og reflecterer sig deri. Som
Reflexion er Reflexionen altid logisk; men det Ydre staaer
med antilogisk Bestemthed Reflexionen imod; det kommer under
Brydningen an paa, hvilken Part der i Virkeligheden faaer
Overhaand, enten Umiddelbarheden over Reflexionen, eller omvendt
denne over hiin.

Faaer Umiddelbarheden Magt med Reflexionen, saa vil
det Indre ligefrem staae malet i det Ydre, være kjendeligt
paa Yttringen. I en Dom som: „Mennesket er virkelig bedrøvet“,
svarer den Dømmendes umiddelbare, skjøndt reflecterede,
Forvisning til Bedrøvelsens umiddelbare, skjøndt reflecterede,
Fremtræden i den Bedrøvedes Ydre. Vil Reflexionen svække
Forvisningen, saa reagerer Umiddelbarheden strax med forhøiet
Eftertryk: „det hjælper ikke, hvad I sige, han er
virkelig bedrøvet“.

Under Spændingen mellem Reflexionen og Umiddelbarheden,
en Spænding, som her er betinget af Virkelighedens
egen Fordobling til Indre og Ydre, sees en Sammensmeltning
af assertoriske og interpreterende Functioner. Naar Dommens
Gjenstand ikke er den Dømmendes eget, men en Andens
Indre, saa beroer Dommen jo paa en Fortolkning af de
Yttringer, hvorigjennem de to Subjectiviteter correspondere.
„Manden er virkelig vred, det forsikkrer jeg, jeg har jo hørt
det af hans egen Mund; ja, han er vred tilgavns, det kan
jeg see af hans Mine“. At slige Domme udtrykkelig
% --- p. 193 ---
forud\opage{193}sætte
en bestemt Udtydning af den Talendes Ord og Mine,
mærkes uvilkaarlig, naar En møder med den modsatte Forsikkring:
„Manden er virkelig ikke vred, han mener ikke Noget
med det“. Forsaavidt Reflexionen sættes i Bevægelse, synes
det vel, som om de Dømmende maatte gribe til Beviser og
Grunde; men paa Grund af den herskende Umiddelbarhed er
Begrundelsen et Skin; Begrundelsen bestaaer i, at den ene
assertoriske Dom ombyttes med den anden. „Manden“, hvis
Vrede omtvistes, „har virkelig meent, hvad han sagde“; „nei,
han har virkelig ikke meent, hvad han sagde.“ „Hans Miner
forraade virkelig Vrede“; „nei, hans Miner forraade virkelig
ingen Vrede“; „det er virkelig Alvor!“ „det er virkelig Spøg det
Hele!“ At den assertoriske Function imidlertid ikke ligefrem
falder sammen med den interpreterende, men bliver en Ændring
af denne, en subjectiv Modalitet, en af de Maader, den interpreterende
Function kan yttre sig paa, følger ligefrem deraf,
at den Begrundelse, Interpretationen væsentlig fordrer, for
den assertoriske Functions Vedkommende, væsentlig udebliver.
Paa den umiddelbare Virkeligheds Trin kan her ikke gives
og skal ikke gives nogen væsentlig Begrundelse: den assertoriske
Interpretation er en ligefrem Oversættelse af Virkelighedens
eget Sprog, og Virkelighedens Sprog er et kategorisk
Magtsprog.

Hvor slaaende, utvetydigt dette kategoriske Magtsprog er,
kan sees overalt, hvor den sunde Sands faaer Ret imod forskruede
Skoletheorier. Ifølge den cartesianske Lære om Forholdet
mellem Sjæl og Legeme kunde Dyret, som bekjendt,
ikke have nogen Sjæl: ethvert Dyr var at betragte som en
Automat. Naar man slog en Hund, og den peb, saa var
dens Piben aldeles ikke nogen Yttring af Smerte; det var
mechanisk en Lyd, som naar man blæser paa en Fløite. Det
er altsaa, efter denne Synsmaade, en Latterlighed at udstede
Forbud mod Dyreplageri, man kunde ligesaa godt udstede
Forbud imod Locomotivplageri. Og dog ansees Dyreplageri
% --- p. 194 ---
\opage{194}for noget Uædelt og Strafværdigt: har man da virkelig faaet
Cartesius gjendrevet? Ja, men ikke theoretisk (de theoretiske
Gjendrivelser komme nemlig ikke her i Betragtning), men
praktisk, det vil sige: ved et Magtsprog; ikke ved et vilkaarligt,
nei ved et kategorisk, Virkelighedens kategoriske Magtsprog.
Idet Kjendsgjerningernes Sprog er kategorisk, maa den umiddelbare
Forstaaelse være assertorisk. „Hunden har virkelig en
Sjæl, det lyser jo ud af dens Øine; det stakkels Dyr kan
virkelig føle Smerte“. Det er kun assertoriske Domme, men
Domme, der, saa at sige, med Dommedagsslag slaae Speculationen
til Jorden.

Faaer Reflexionen omvendt Overhaand over Umiddelbarheden,
da bliver Fortolkningens subjective Side fremherskende;
den assertoriske Interpretation giver paa dette Standpunkt,
hvad man i Ordets sædvanlige Forstand forstaaer ved
en \emph{Mening}.

At Kategorien Mening paa en egen Maade svæver
mellem Subjectivt og Objectivt, mellem Reflexion og Umiddelbarhed,
kan let paavises. Den umiddelbare Opfattelse af
det umiddelbart Objective er endnu ikke en Mening; thi der
er dog vel Ingen, der vil kalde den Antagelse, at Vandet er
flydende, eller at Solen skinner, en Mening. Vilde Nogen
pege paa Solen, naar den stod høist paa Himlen og sige: „see,
det er nu min Mening, at Solen skinner“, saa vilde man
antage det enten for en Yttring af Lune eller et Tegn paa
Forskruethed. Paa den anden Side fører en gjennemgaaende
Reflexion af væsentlige Forhold heller ikke til Mening, men
til det, som er mere end Mening: den fører til Viden. At
Skin maa forudsætte Væsen, at Grunden maa have sin Følge,
Aarsagen sin Virkning, er hos den dømmende Subjectivitet
ikke en Mening, men en Viden. Hvor Reflexionen derimod
sættes i Bevægelse saaledes, at den hverken negter Umiddelbarheden
eller ubetinget bøier sig under den, der er den rette
Plads for Meningen, der har den sit egentlige Spillerum.
% --- p. 195 ---
\opage{195}Stille vi de to assertoriske Domme: „Manden er virkelig
vred“; „nei, Manden er virkelig ikke vred“, imod hinanden,
da er det, egentlig talt, ikke en Modsigelse mellem to Begreber,
men en Strid mellem to Meninger, vi have for os.

Det er Umiddelbarheden af en i sig reflecteret Virkelighed,
der udgjør Meningernes objective Grundlag. Saadanne
i sig reflecterede Virkeligheder have vi i Stat, Kunst og Religion.
Paa Pietetens, Naivetetens og Eenfoldighedens Trin
er Dommen om disse Virkeligheder endnu umiddelbar assertorisk.
„Det er virkelig Ret, hvad Øvrigheden byder“; „det
er virkelig skjønt, hvad Mesteren der har frembragt“; „det er
virkelig sandt, hvad Præsten siger“. I disse Domme er
det Umiddelbares Magt saa overveiende, at den subjective
Reflexion holdes bunden: det er dominerende Fordomme,
sikkre Autoritetssætninger, faste Dogmer, vi her forefinde; men
det er endnu ikke egentlige Meninger. Den Subjectiviteternes
Frigjørelse, den Bevægelighed i Reflexion, den Lethed i oversvævende
Betragtninger, der begunstiger Meningernes Fremkomst
i Samfundet, ligesom Varmen i Juni begunstiger Fluernes
Fremkomst, indtræder først paa et senere Stadium.

Det er Culturstadiet, de subjective Reflexioners Stadium,
der egentlig er at opfatte som Meningernes Stadium. „Nu er det
virkelig ikke altid Ret, hvad Øvrigheden byder“; „det er virkelig
heller ikke Altsammen skjønt, hvad en Kunstner producerer“; „man
kan da virkelig heller ikke troe Alt, hvad Præsten siger“. Hvad
er da Ret? hvad er da skjønt? hvad skal man da troe? Ja,
derom er det just, at der er saa forskjellige Meninger. Medens
de interpreterende Functioner i Kraft af Begrundelsen belyste
den objective Tingenes Mening, vende de assertoriske derimod
Opmærksomheden --- ikke mod det Objective, mod Meningen
i Tingene --- men mod det Subjective, mod Individernes Meninger
om Tingene: vi faae en Hærskare af politiske, æsthetiske
og religiøse Meninger; hver Mening er er assertorisk Dom,
% PRINTER: printed "er er", sense "er en"
bestemt ved subjectiv Reflexion.

% --- p. 196 ---
\opage{196}I Meningen ere de objective og subjective Elementer paa
eiendommelig Maade frigjorte og dog forbundne. Subjectiviteten
er emanciperet; thi hvad er den Enkelte lettere end at danne sig
en Mening, og hvad er den Enkeltes aandelige Eiendom,
dersom det ikke er hans „uforgribelige Mening“: hvert Individ,
selv en Skoledreng, har jo paa dette Trin en uforgribelig
Mening. Og dog peger Meningen objectivt lige paa Tingen;
det er Meningen om Tingen, der afgjør, hvad Tingen er.
Tingen er det Virkelige, og den Menende lægger Vægt paa
det Virkelige, paa Sagen; det er ham naturligviis ikke om sin
individuelle Mening, men om Sagen, den vigtige Sag, at gjøre.
Kunde Nogen overbevise ham om, at hans Mening var urigtig,
„skulde han sandelig for Sagens Vigtigheds Skyld gjerne forandre
den“; men det er umuligt, thi, da det nu een Gang er
hans uforgribelige Mening, at Sagen virkelig forholder sig saaledes,
som han mener det, saa vilde det jo være Forræderi
imod den gode Sag, om han forandrede sin Mening.

Meningen er en assertorisk Dom; naar det Assertoriske
betones stærkt, stiger Meningen til Paastand; i Paastanden
vil da Subjectiviteten selv udtrykkelig indestaae for det Objective:
den Menende bliver paastaaelig. Idet den Menende
bliver paastaaelig, seer det unegtelig ud, som om han vilde
paatrænge Andre sin tilfældige Subjectivitet; men dette er nu
ikke Meningen. Meningen er, at den Paastaaelige ivrer for
Sagen, er den Eneste, der seer Sagen, som den virkelig er.
Den Paastaaelige mener kun, at hans Mening ubetinget bør være
\emph{Alles} Mening, og det af den simple Grund, at det er den
eneste rigtige Mening. De Menende ere liberale, thi Enhver
har jo Ret til at beholde sin Mening; men naar de mange
Meninger gaae imod hverandre, og Enhver holder paa sin,
fordi det er den ene rigtige, saa maae de mange ulige Meninger
dog tilsidst geraade i en betænkelig Strid: Opgaven
er da at finde den for Alle gjældende Mening, den alminde%
% --- p. 197 ---
\opage{197}lige Mening, nærmere bestemt den offentlige Mening, \emph{Opinionen}.

Det er af logisk Interesse at charakterisere Meningernes
Kampe, Kampene for Opinionen, fordi disse Kampe, seete
paa Afstand, tage sig ud, som om det gjaldt Sandheden, og saa
ere de dog i Principet indifferente mod Sandheden. Det er
ikke Sandheden, men Virkeligheden, ikke Virkeligheden som
den udtaler sig i Kjendsgjerningernes kategoriske Magtsprog,
men Virkeligheden saaledes som den, underlagt Interesserne,
omsløret af Reflexionernes Spindelvav og farvet af Meningerne,
tiltaler Subjectivitéterne, hvorom Striden er. % PRINTER: printed "Subjectivitéterne" (accent), sense "Subjectiviteterne"

Den Maade, hvorpaa Striden føres, er i mange Henseender
lærerig. Her undersøges, her overveies, her anføres
Grunde og Modgrunde; Logikens hele Rustkammer med
Domme, Slutninger, Beviser o. s. v. gjennemsøges for at
skaffe de Stridende Vaaben: man raisonnerer, discuterer,
argumenterer; man gaaer hinanden paa Klingen. Men, see
vi nærmere til, saa er Striden flygtig; Reflexionen spiller
paa Virkelighedens Overflade, ofte glimrende, som Solstraaler
paa en kruset Vandflade, men Grunden er flad, Fordybelsen
mangler. Det Charakteristiske er, at hvormeget her end tales
om Principer og Principspørgsmaal, saa kommer det dog
aldrig, og kan, i Meningernes, Paastandenes, Forsikkringernes,
de assertoriske Dommes skarpt afgrændsede Sphære, aldrig
komme til nogen egentlig Principernes Undersøgelse. Reflexionen
har en Behændighed i at bevæge sig, som om den
gik til Bunds i Tingen og naaede Grunden; men denne
Bund er ingen væsentlig Grund, det er en Skorpe af krystalliserede
Meninger. Reflexionen har Færdighed i at stable de
assertoriske Domme saaledes, at det for den Uindviede seer
ud som Batterier af slaaende Argumenter, hele Fæstningsværker
af uigjendrivelige Beviser. Det er Skin, men et Reflexionens
Skin, der svarer til det Umiddelbare i Virkelighedens Væsen.

Ved enhver Strid, den bevæge sig iøvrigt paa Politikens,
% --- p. 198 ---
\opage{198}Kunstens eller Religionens Enemærker, gjør det en væsentlig
Forskjel, om det, der strides for, egentlig er en Sandhed eller
en Opinion. Dette kan nu aldrig ligefrem høres paa Feldtraabet,
thi Feldtraabet lyder naturligviis altid paa Sandheden;
men det kan sees paa Methoden, det kan kjendes af den
Maade, hvorpaa Striden føres. Thi en Methode, der maa
forkastes, naar det gjælder Sandheden, kan netop være den,
der anbefaler sig som den eneste rette, naar det gjælder Opinionen.
Vi kunne oplyse dette ved et Par Exempler.

Er Striden politisk, og en af de Bannerførende optræder
med den Erklæring, at han, hvad Modpartens Standpunkt
angaaer, vilde gjøre sig en Samvittighed af at trætte
Tilhørernes Taalmodighed med at føre Beviis for, hvad
Alle, der ere begavede med sund Sands, forlængst ere enige
om, at det hører til de forlorne Standpunkter, og derfor
skal indskrænke sig til en Paaviisning af Kjendsgjerninger,
der ere saa iøinefaldende, saa talende o. s. v. --- da høre vi
strax, at det ikke er Sandheden, men Opinionen, det er ham
om at gjøre. Var det Ordføreren om Sandheden at gjøre,
saa maatte det jo være ham af den største Vigtighed, at
Standpunkternes principielle Forskjel ret kunde blive Tilhørerne
klar, thi de paaberaabte Kjendsgjerninger ville naturligviis
tage sig forskjelligt ud, eftersom Lyset falder ind fra den ene
eller fra den anden Side. I en Kamp for Opinionen er
Talerens Fremgangsmaade derimod praktisk---taktisk fuldkommen
rigtig. Det er en Paastand, en assertorisk Dom, at Modpartens
Standpunkt hører blandt de forlorne Standpunkter; en Paastand,
en assertorisk Dom, som hviler paa den Paastand, at Alle,
der have sund Sands, uden videre kunne see det. Denne Paastand,
denne assertoriske Dom, hviler igjen, hvis Nogen iblandt
Tilhørerne skulde faae Skrupler og mene, at man dog ogsaa
kunde være af en anden Mening, paa den Paastand, den
assertoriske Dom, at den, der ikke kan see, at Modpartens
Standpunkt er et forlorent, virkelig ikke kan siges at have
% --- p. 199 ---
\opage{199}sund Sands. Denne Paastand, denne assertoriske Dom, hviler
saa til Syvende og Sidst paa den Paastand, at et Individ,
der i Virkeligheden ikke har sund Sands, ikke er berettiget
til at tale med om politiske Standpunkter. Med et saadant
Underlag af Paastande, paa et saadant Bjælkeværk af assertoriske
Domme er Taleren sikker paa sig selv, sikker paa sit
Standpunkt, sikker paa sin Strid: han strider for Opinionen.

Er Striden æsthetisk, maa Methoden ligeledes variere,
eftersom Stridens Maal er en Sandhed eller en Opinion.
Er Maalet en Sandhed, saa kan Veien stundom være lang
og besværlig nok. Thi til en Bedømmelse af Kunstens Værker
er det naturligviis ikke nok at have naturlig Smag og besøgt
adskillige Museer; det kommer ved Bedømmelsen, naar den skal
have videnskabeligt Værd og tjene Sandheden, fornemmelig
an paa, at Dommen motiveres. For at kunne motivere sin
Dom, maa den Dommende fremfor Alt være sikker paa Forholdet
mellem Begreb og Phænomen, mellem det Almene og
det Individuelle; men dertil hører et ligesaa omhyggeligt og
dybtgaaende Studium af Begrebet i Videnskaben som af Phænomenet
i Kunsten. Anderledes, naar Striden gjælder en Opinion.
Her kan en Bannerfører og Basunblæser spare sig Uleiligheden
med Videnskabens Analyser. „Hvad kommer der vel ud af
disse tørre Abstractioner, disse farveløse philosophiske Begreber?
Nei, Phænomen og Følelse og Følelse for Phænomenet, det er
og bliver dog Hovedsagen“. Da nu den finere Følelse først udvikles
ved Berøring med de mangfoldige Kunstphænomener,
medens Forstandens Øie oplades ved at læse Beskrivelser og
affatte Beskrivelser over de mange Kunstsager, Kunstarter,
Kunstnere og Kunstdommere fra de forskjellige Tidsaldere,
saa maa den bannerførende Kunstkritiker, for Opinionens Skyld,
indrette sin Optræden saaledes, at Publikum baade kan have
Tillid til hans rene Smag, hans fine Sands, og tillige være
imponeret af hans megen Lærdom og udførlige Beskrivelser.
I et saadant Apparat, et Apparat, som et hurtigt Hoved
% --- p. 200 ---
\opage{200}vil kunne skaffe sig, især paa en Udenlandsreise, i Løbet af
nogle Maaneder, har Opinionens Bannerfører Alt, hvad han
æsthetisk behøver til Meninger, Paastande, assertoriske Domme.

Dog er her for Methodens og Opinionens Skyld endnu
Eet at iagttage. Et Apparat af Meninger, Paastande og
assertoriske Domme kan være saa stort og imponerende det vil,
saalænge man holder sig til Phænomenet: det viser sig snart
utilstrækkeligt, naar man for Alvor indlader sig paa noget
Videnskabeligt. Methoden kræver her en vis Varsomhed; thi
Methoden er, at man, uden ligefrem at indlade sig med
Videnskaben, bevarer et vist videnskabeligt Anstrøg, det Anstrøg,
der netop passer for den assertoriske Dom. Naar den assertoriske
Dom, istedenfor i al Plumphed at opstilles som Paastand,
ved Tinktur af „Følelse“ og Lærdom faaer et videnskabeligt
Anstrøg, forvandler den sig og avancerer fra simpel Mening
til fiin Bemærkning: man virke især ved fine Bemærkninger!

At en Bemærkning maa kunne have videnskabeligt Anstrøg,
ja næsten blive fuldvægtig videnskabelig, kan da sees af
den Anseelse, Opinionen indrømmer, naturligviis ikke de philosophiske
Systemer, men de philosophiske Bemærkninger. Ja,
Opinionen har virkelig Smag for philosophiske Bemærkninger.
Intet Under, at der til Opinionens Behov leveres Philosophier
alenviis, Philosophier, der ere sammenvævede af lutter
Bemærkninger. Ved slige Bemærkningsphilosophier kan her
da i mange Ting være Adskilligt at blive vaer og Meget at
bemærke, thi Vaerblivelsen gaaer paa Umiddelbarheden, og
Bemærkningen løber jo netop paa Grændsen, den assertoriske
Grændse, mellem Menen og Viden. Nyder nu Bemærkningsphilosophien
vitterlig stor Anseelse i den offentlige Mening, hvorfor
skulde en Bemærkningsæsthetik, en Kunstkritik af fine Bemærkninger,
især naar det er Udstillingsbemærkninger, ikke
nyde samme Anseelse?

Er Striden theologisk, da kan man ikke noksom prise,
% --- p. 201 ---
\opage{201}anbefale og beundre Opinionsmethoden. Her er nu f. Ex.
et Skrift, der med Rette betragtes som et af de mærkeligste
Aandsprodukter i vort Aarhundrede. Det er et christeligt philosophisk
Skrift, yndigt skrevet til en Forsoning af Tro og Viden.
Som talende Beviis paa Skriftets dybe philosophiske Aand tjener,
at uagtet det er forkastet af dem, der just forstaae sig paa Philosophie,
har det dog, trods den strenge Kritik, vundet uhyre
Udbredelse og videnskabelig Anerkjendelse hos alle dem, der
ikke forstaae sig paa Philosophie; thi dets Philosophie er theologisk.
„Vi tør“, siger Anmelderen, „frit paastaae, at dette
Skrift er i Hænderne paa Alle, som med høiere Dannelse forene
dybere Sands for sand Christendom. Vi kunne bevise vor
Paastand; thi vi paastaae, at dersom der virkelig gives Dannede,
i hvis Hænder Skriftet ikke er, da have disse Dannede
ikke Sands for sand Christendom. Vi kunne bevise vor Paastand,
thi vi paastaae, at dersom der virkelig skulde gives
sande Christne, i hvis Hænder Skriftet ikke er, da ere disse
Christne uden høiere Dannelse“. Slige Paastande ere forkastelige,
ja foragtelige, naar der strides for en Sandhed;
men som assertoriske Domme med forstærket Eftertryk gjøre
de ypperlig Virkning, naar der strides for en Opinion.

3) \emph{Den assertoriske Doms logisk-antilogiske
Væsenhed: dobbeltsidig Modalitet.} At articulere en
Forskjelseenhed af det Antilogiske med det Logiske er ingenlunde
noget udelukkende Særkjende for de assertoriske Functioner;
thi det Samme kan jo i en vis Forstand siges at
gjælde om alle Functioner. Nei, det for de assertoriske Functioner
Eiendommelige er netop det Forhold imellem Virkelighed
og Reflexion, at det Virkelige i sin antilogiske Umiddelbarhed
enten maa beherske Tanken ved et kategorisk Magtsprog
eller, forsaavidt Tanken antager Præg af subjectiv Selvstændighed,
da ved æggende Modsætning bevirke en saadan
Reflexionens Stigen, at det pointerende Vexelspil af logiske og
antilogiske Formbestemmelser, der betinger den reale Viden,
% --- p. 202 ---
\opage{202}bliver synligt fra alle Sider; kun saaledes kan Modalitetens
saavel objective som subjective Gyldighed tilfulde belyses.

Paa tre Maader vil Virkeligheden kunne indgaae i ideel
Proces med Reflexionen, eftersom den enten lader sig forvandle
til et for sit individualiserede Indhold umiddelbart
gjennemsigtigt Phænomen og derved til Ideens Afspeiling; eller
den, hvad Charakteren angaaer, vel beholder sin endelige Realitet,
men undergaaer en Forvandling derved, at den ikke
opfattes som et Nærværende, men som et Forbigangent, og
saaledes afdæmper de antilogiske Kanter; eller den med Præg
af en guddommelig Realitet, der er hævet over den endelige Forskjel
mellem Nutid og Fortid, paatrænger sig Subjectiviteten
som en Virkelighed af høiere Orden, en overvirkelig Virkelighed,
en Virkelighed, der gjør Brud paa det Virkelige.
I første Tilfælde er den assertoriske Function \emph{æsthetisk}, i
andet Tilfælde \emph{historisk}, i tredie Tilfælde \emph{religiøs}.

Den \emph{æsthetiske} Eenhed af Væsen og Phænomen er assertorisk;
heri ligger, at den Forskjel, hvortil Eenheden svarer,
ligeledes er assertorisk. Betragte vi Phænomenet, saa er
Formen antilogisk bestemt som det Virkelige selv: denne Jason
i Marmor har f. Ex. en ligesaa umiddelbar Individualitet,
som dette Ladegaardslem i Kjød og Blod; men „Jason er
virkelig skjøn“, „den stakkels fattige Mand er virkelig ikke
skjøn': den ene assertoriske Dom har her samme Gyldighed % PRINTER: printed closing mark "'" after "skjøn", sense "“"
som den anden.

Men hvorpaa beroer nu Gyldigheden? At begge Skikkelser
virkelig ere til, den ene som Figur, den anden som existerende
Individ, stiller dem, hvad Virkelighedens antilogiske Umiddelbarhed
angaaer, i en vis Forstand begge paa lige Trin: den
assertoriske Dom, der opfatter det umiddelbart Virkelige, viser
os Modalitetens objective Side. Forskjellen mellem det
Skjønne og det Uskjønne er altsaa en Virkelighedsforskjel; det
er Virkelighedsformen selv, der gjør Udslaget. Virkelighedsformen
er i sin objective Umiddelbarhed uafhængig af den
% --- p. 203 ---
\opage{203}betragtende Subjectivitet; det virkelig Skjønne er altsaa
virkelig objectivt: Skjønheden er i det Skjønne Eet med det
Virkelige.

Men Uskjønheden er i det Uskjønne jo ogsaa Eet med
det Virkelige. Altsaa: det Virkelige er Eet med det Skjønne,
og det Virkelige er Eet med det Uskjønne; altsaa: det Virkelige
er forskjelligt fra det Uskjønne, thi det er jo Eet med det
Skjønne; det Virkelige er forskjelligt fra det Skjønne, thi det
er jo Eet med det Uskjønne; altsaa: det Virkelige er som
Virkeligt forskjelligt baade fra det Skjønne og fra det Uskjønne,
kan ligesaa vel være uskjønt, som skjønt, er i et udvortes og
ligegyldigt Forhold til Skjønheden. Hvad bliver der da af
Skjønheden? Den forvandler sig til Idee, opgaaer i Phantasien
og reflecterer sig i Subjectiviteten. Det er altsaa Betragteren,
ikke den betragtede Gjenstand, der selv afgjør, hvad
der er skjønt og uskjønt. Fra Modalitetens subjective Side
er Mennesket atter „Tingenes Maal“, baade de skjønnes, at
de ere skjønne, og de uskjønnes, at de ere uskjønne.

Paa det æsthetiske Omraade er den assertoriske Functions
oprindelige Slægtskab med de explicative og interpreterende
ret iøinefaldende. Fra Reflexionens, altsaa fra Subjectivitetens
logiske Synspunkt er Opgaven at gjenkjende det i
Gjenstanden virkeliggjorte Begreb; Functionen er forsaavidt
explicativ. Men i æsthetisk Forstand maa Gjenkjendelsen være
umiddelbar. Ved den umiddelbare Gjenkjendelse gaaer Formen
i Eet med sit Inhold; Begrebet forsvinder som Begreb og % PRINTER: printed "Inhold", sense "Indhold"
forvandler sig for Anskuelsen til Forbillede: at gjenkjende Begrebet
i Gjenstanden er at maale den med Forbilledet. Siger
nu det umiddelbare Skjøn, at Gjenstanden svarer til Forbilledet,
saa falder Dommen, at den virkelig er skjøn; afviger
Gjenstanden mærkelig fra Forbilledet, saa falder Dommen,
at den virkelig er uskjøn; staaer Gjenstanden i paafaldende
Strid med Forbilledet, lyder Dommen paa, at den er hæslig.
Paa Grund af Formens antilogiske Bestemthed og Gjenkjen%
% --- p. 204 ---
\opage{204}delsens Umiddelbarhed ere de æsthetiske Domme oprindelig
assertoriske.

De æsthetiske Dommes assertoriske Charakteer gjør det
forklarligt, at der om Skjønhedsphænomenerne og det Skjønne
kan herske en ligesaa paafaldende Enighed som paafaldende
Uenighed.

Den assertoriske Dom er umiddelbart subjectiv; de æsthetiske
Domme ere Smagsdomme: Enhver af de Dømmende
har sin Smag; hvad hjælper det at ville disputere om Forskjellen
mellem en god og en slet Smag, naar Dommen om
Smagen selv er en Smagsdom? Og dog er her virkelig Enighed;
Folk af den gode Smag slutte sig sammen, forstaae
hinanden og støtte hinanden, idet de bekæmpe den slette: det
er en paafaldende Enighed. Den paafaldende Enighed maa
tilskrives et kategorisk Magtsprog; men denne Gang er det
ikke Virkeligheden i sin Tilfældighed, men den af Ideen imprægnerede
Virkelighed, Ideen i Virkelighedens antilogiske
Form, der taler til Subjectiviteten.

Den assertoriske Dom er umiddelbart objectiv. De æsthetiske
Domme ere Anskuelsesdomme; det er ikke Dommen, der
bestemmer Gjenstanden, men den anskuede Gjenstand, der bestemmer
Dommen. Anskuelsens Gjenstand er fælles for de
mange Anskuende: „see, denne Skikkelse, hvor er den dog skjøn!
disse Træk, disse Linier, hvor ere de yndige!“ --- „Og nu
denne Figur, hvor klodset, uproportioneret!“ o. s. v. Men
uagtet Gjenstandene ere de samme, uagtet de æsthetiske Udbrud
hos Sindige og Sagkyndige ombyttes med Udtydninger, hvorved
de assertoriske Functioner ret bevise deres Slægtskab, ikke
blot med de explicative, forsaavidt Udtydningen skeer efter et
antaget Forbillede, men ogsaa med de interpreterende, forsaavidt
den ene Subjectivitet gjennem Kunstværket vil tiltale den
anden, saa er her dog Uenighed, just som Interpretationerne
komme i Gang, en paafaldende Uenighed; selv de Sagkyndige
kunne nu med deres bedste Villie ikke blive enige.

% --- p. 205 ---
\opage{205}Har da Skjønheden ikke objectiv Gyldighed? Jo, ganske
vist! det er just det Eiendommelige i dens objective Gyldighed,
der gjør Dommen om den saa umiddelbart subjectiv. Men
saa maa Skjønhedens Gyldighed vel beroe paa Subjectiviteten?
Ja, ganske vist! det er just det Eiendommelige ved dens subjective
Gyldighed, at Dommen om den er saa umiddelbart
objectiv.

Forskjelseenheden af Objectivt og Subjectivt har sin
almene Gyldighed i Ideen; det er hverken det Objective for
sig eller det Subjective for sig, men Ideen selv, der er Skjønhedens
Maal; derfor taler Plato om Selvskjønheden, den
usynlige Skjønhed, hvorved alle synlige skjønne Ting ere
skjønne. Og dog seer Plato feil; thi Ideen i sin usynlige
almene Væsenhed er ikke Skjønheden; først som en Magtens
Idee i anskuelig Form har Skjønhedsideen Tilværelse.
Som Skjønhedsidee er Ideen, det Existerendes Forbillede,
selv existerende Skikkelse: den existerende Skikkelse bestemmer
umiddelbart den assertoriske Dom. Saaledes vedbliver det
pointerende Vexelspil af Subjectivt med Objectivt i logisk-antilogiske
Momenter uophørlig; men det forklares af Ideen,
thi den æsthetiske Virkelighed er i Ideen gjennemsigtig.

Den \emph{historiske} Eenhed af Væsen og Phænomen er assertorisk;
men paa en anden Maade end den æsthetiske. For
Ideen maa den æsthetiske Virkelighed være fuldkommen gjennemsigtig;
dette er nu ikke saaledes Tilfældet med den historiske.
Det historisk Virkelige bærer Endelighedens og derved Tilfældighedens
Præg. Kjendsgjerningen er ved sin Tilfældighed
vistnok i Ideens Magt, men ingenlunde dens adæquate Phænomen;
forsaavidt kan den af antilogisk Umiddelbarhed frastødte
Reflexion da heller ikke ligefrem gaae til Ideen for at
finde Forklaringen. Derimod har Virkeligheden, idet den
bliver historisk, det Eiendommelige at omsættes i Forbigangenhed;
ved denne Omsætning beskjæftiger den Reflexionen uden
at opgive Umiddelbarheden.

% --- p. 206 ---
\opage{206}At det Forbigangne, uagtet det ikke har Nærværelsens
Umiddelbarhed, ikke desmindre har en Umiddelbarhed, skjønnes
uvilkaarlig af saadanne Domme som: „Cæsar blev virkelig
myrdet“, „Knud den Store erobrede virkelig England“.
Hvor mange Beviser og Grunde man end kan opdrive til
Stadfæstelse af slige Domme, saa vil dog Umiddelbarhedselementet
altid være afgjørende. Vel har det Forsvundne tabt
sin Umiddelbarhed, forsaavidt det er forsvundet; men det har
bevaret den, forsaavidt det ikke er forsvundet. I den nærværende
Fortælling er det Forsvundne jo ikke forsvundet; den forsvundne
Virkelighed har i Fortællingen en reflecteret Nærværelse
og virker i denne Nærværelse med reflecteret Umiddelbarhed.

Men det Opdigtede kan i en poetisk anskuelig Fortælling
jo ligesaa vel antage Præg af Virkelighed som det forsvundne
Virkelige: hvorpaa beroer nu Adskillelsen mellem det Opdigtede
og det forbigangne Virkelige? Aabenbart paa en Maugfoldighed % PRINTER: printed "Maugfoldighed", sense "Mangfoldighed"
af umiddelbare og middelbare Virkelighedsindtryk! Er
Læseren ikke selv Forsker, saa tager han Fortællingen paa den
Fortællendes Autoritet; men et Autoritetsindtryk er med alle
deri indbefattede Umiddelbarhedsbestemmelser altid et Virkelighedsindtryk.
Er Læseren Forsker, ville mangfoldige Hensyn,
--- til Kilden eller Kilderne, til Begivenhedens Charakteer,
dens Forhold til andre Begivenheder o. s. v. --- komme i
Betragtning; men Hovedsagen er, at de mange særskilte Indtryk
tilsidst samle sig i et Totalindtryk, og at dette for Begivenheden
talende Totalindtryk er et afgjørende Virkelighedsindtryk.

Gjennem Totalindtrykket virker det Forsvundne paa
middelbar og dog umiddelbar Maade; det virker umiddelbart,
thi Mindesmærket staaer der, Beretningen fængsler Opmærksomheden,
og den opmærksomme Tilhører seer Begivenheden
for sine Øine; det virker middelbart, thi Mindesmærket
fjerner, Fortællingen minder i alle Enkeltheder om Afstanden
% --- p. 207 ---
\opage{207}mellem det Nærværende og det Forbigangne: Reflexionen er
i Bevægelse.

Det pointerende Vexelspil af Reflexion og Umiddelbarhed
er her af eiendommelig Art. Paa det æsthetiske Grundlag kan
Reflexionen smelte Umiddelbarheden; thi Phænomenet er væsentlig
Skin; at den assertoriske Function desuagtet spiller
Hovedrollen, ligger deri, at den ophævede Umiddelbarhed øieblikkelig
restitueres til antilogisk Bestemthed, efterdi Phænomenet
er en tilværende Skikkelse. Paa det historiske Grundlag
kan Reflexionen derimod ikke smelte det Umiddelbare, thi
Kjendsgjerning er dog Kjendsgjerning, om den end tilhører det
Forbigangne; Umiddelbarhedens ideelle Proces med Reflexionen
maa da sees i det Middelbare. Imellem Betragteren og Begivenheden
er Forholdet uroligt: bestandig en Fjernelse, der
uafbrudt vexler med en tilsvarende Nærmelse. Middelbart
nærmer Betragteren sig til Begivenheden, og Begivenheden fra
det Forbigangne kommer ham imøde; under Tilnærmelsen
voxer Umiddelbarheden: den virkelige Begivenhed gjør et virkeligt
Indtryk. Men det saaledes Nærværende er dog et
Forbigangent, det fjerner sig igjen; det maa tilbage til sit
Sted, naar Reflexionen begynder, ligesom Gjengangeren maa
tilbage, naar Morgenen dæmrer. Det er dette Urolige, der
kommer til Hvile i Troen, som historisk Tro (\textit{fides historica}).

Den historiske Tro er en assertorisk Function, eiendommelig
bestemt ved Reflexionens Omsætning af Middelbart i
Umiddelbart. Den historiske Tro bestemmes ved Virkelighedsindtryk,
dens Grundlag er Umiddelbarhed, dens Dom er
assertorisk. Ingen Reflexion kan give umiddelbar Forvisning;
kun det Virkelige kan tale for det Virkelige; dette er Modalitetens
objective Side: den assertoriske Dom er umiddelbart
objectiv. Men det Virkelighedsindtryk, der grundlægger den
historiske Tro, er paa Grund af Betragterens og Begivenhedens
indbyrdes Afstand dog kun et middelbart Indtryk, et
Indtryk af en virkelig fraværende, kun i Reflexionen nær%
% --- p. 208 ---
\opage{208}værende Virkelighed: dette er Modalitetens subjective Side;
den assertoriske Dom er umiddelbart subjectiv.

Den \emph{religiøse} Eenhed af Væsen og Phænomen er i en
ganske egen Forstand assertorisk. Den historiske og den religiøse
Tro ligne hinanden deri, at de begge beroe paa umiddelbare
Indtryk af et umiddelbart Fraværende: Troen er en
assertorisk Forvisning om en ikke nærværende Virkelighed.
Først naar Spørgsmaalet opkommer, hvad det da egentlig er,
der hindrer det Virkeliges umiddelbare Nærværelse, kommer
Forskjellen mellem den historiske og den religiøse Tro tilsyne.
Hvad der forhindrer det historisk Virkelige fra at virke med
det Nærværendes hele Umiddelbarhed, er ene og alene Forbigangenheden;
kunde Forbigangenheden hæves, da vilde det
Forbigangne naturligviis virke med samme Umiddelbarhed
som det Nærværende, thi i Væsen ere de eensartede. Anderledes
med det religiøst Virkelige. Det religiøst Virkelige kan,
om det end er et Nuværende, dog ikke fremtræde som et
umiddelbart Nærværende. Et Mirakel kan ikke anskues som
umiddelbart nærværende; dets Nærværelse forkynder sig tvertimod
som et Brud med alt virkeligt Nærværende; et saadant
Brud forfærder, og Forfærdelsen gjør det Nærværende til et
Fraværende. Det saaledes Fraværende er i sin forfærdelige
Nærværelse det Skjulte, det Hemmelighedsfulde, Uforklarlige.

„Hvad er det da?“ Samtidig med Miraklet er dette
Spørgsmaal kun et Udraab; besindig Overveielse kan ikke
være samtidig med Miraklet, thi Miraklets Umiddelbarhed
overvælder Reflexionen. Naar Besindelsen indtræder, og
Overveielsen begynder, er Miraklet forbi; Overveielsen træffer
derfor, selv hos de Samtidige, aldrig et Nærværende, men
altid et Forbigangent.

„Hvad var det da?“ Her viser det sig, at Troen, ogsaa
den religiøse, er en assertorisk Forvisning om et ikke nærværende
Virkeligt. Det er ikke en Begrundelse af Troen,
men Troens logiske Charakteer som assertorisk Forvisning,
% --- p. 209 ---
\opage{209}hvorom her handles.\footnote{Hvad Begrebets religiøse Gyldighed angaaer, henvise vi til S.
Kierkegaards Skrifter, navnlig til „Uvidenskabelig Efterskrift“ og
„Indøvelse i Christendom“. --- „Skal der være Mening i, at
Miraklet beviser, hvo Christus er, saa maae vi jo begynde med, at
vi ikke veed, hvo han er, altsaa i Samtidighedens Situation med
et enkelt Menneske, der er som andre Mennesker, paa hvem der
ikke \emph{ligefrem} er Noget at see, et enkelt Menneske, der saa gjør
Miraklet og siger sig selv at gjøre Miraklet. Hvad vil dette sige?
Det vil sige, at dette enkelte Menneske gjør sig til mere end
Menneske, gjør sig noget nær til Gud: er dette ikke forargeligt?
Du seer noget Uforklarligt, Mirakuløst (mere heller ikke), han selv
siger, at det er et Mirakel --- og Du seer for Dine Øine et enkelt
Menneske. Miraklet kan Intet bevise; thi troer Du ikke, at han er
Den, han siger sig at være, saa benegter Du Miraklet. Miraklet kan
gjøre opmærksom --- nu er Du i Spændingen, og det beroer paa
hvad Du vælger, Forargelsen eller Troen, det er Dit Hjerte, som
skal blive aabenbart“. Indøvelse i Christendom. S. 106.} % note *)
Hvad Kjendsgjerningen er objectivt,
maa umiddelbart afgjøres subjectivt. Mellem Kjendsgjerningen
og den dømmende Subjectivitet er Afstanden uendelig, formedelst
Miraklet. At ophæve den uendelige Afstand er her
udtrykkelig sat som Subjectivitetens uendelige Opgave; og dog
er Løsningen af denne Opgave umiddelbart bestemt ved Indtryk
af Miraklet, altsaa af et Virkeligt.

Seet fra dette Synspunkt har Troen ikke Tvivlen, men
Forargelsen til afgjørende Modsætning. Forargelsen er, logisk
seet, en assertorisk Dom, ligesaavel som Troen „Dette er % PRINTER: no punctuation printed after "Troen" (wide gap), sense ":" or "---"
virkelig et Mirakel!“ „Dette er virkelig en Øienforblindelse!“
Den Troende og den Forargede beraabe sig begge paa Kjendsgjerningen,
det Virkelige i dets Umiddelbarhed, thi den assertoriske
Dom er umiddelbart objectiv. „Dette er virkelig (det
kan jeg jo føle paa mig selv) en Guds Gjerning!“ --- „Dette
er virkelig (det lærer den sunde Sands) en naturlig Begivenhed!“
Den Troende og den Forargede beraabe sig begge paa
indre Sands for Forskjellen mellem Skin og Virkelighed,
% --- p. 210 ---
\opage{210}Løgn og Sandhed; thi den assertoriske Dom er umiddelbart
subjectiv.

I Striden mellem Tro og Forargelse er det ikke blot
Kjendsgjerning, der staaer imod Kjendsgjerning, men i uendelig
Forstand staaer Subjectivitet imod Subjectivitet, Forklaring
mod Forklaring, Forstaaelse mod Forstaaelse: det er de assertoriske
Functioner, som her afgive Grundlag for de interpreterende.
Saalænge Grundlaget er assertorisk, har Troen
aandelig Overlegenhed over Forargelsen, forsaavidt den troende
Subjectivitet med sine positive Kræfter lever et høiere Liv end % UNSURE: stray mark after "Liv" (period?) or none
den Forargede, der kun kan anvende Kræfter og Lidenskaber
negativt. Men skal Grundlaget forandres, skal \textit{credo, ut
intelligam} slaaes sammen med \textit{de omnibus dubitandum est},
skal Modsætningen ikke være Tro og Forargelse, men Tro
og Tvivl, saa gaaer Veien fra det Assertoriske til det Problematiske,
og saa har Forargelsen seiret.

At Forargelsen seirer med Tvivlen, forstaaer sig; thi
Tvivlen, der er Troens Qval, er Forargelsens Lindring.
Er det assertoriske Grundlag opgivet, da maa den reflecterende
Subjectivitet enten blive i det Problematiske eller arbeide sig
igjennem til det Apodiktiske. Det Apodiktiske hviler i rationel
Nødvendighed, men rationel Nødvendighed er kun i Videnskaben.
Kan Troesvidenskaben da naae det Apodiktiske? Den
speculative Dogmatik svarer: til en vis Grad! men det Apodiktiske
til en vis Grad er Nonsens, et assertorisk Nonsens,
der staaer dybt under det Problematiske.

\textit{ββ)} \emph{Den problematiske Dom: Functioner af
ophævet Bestemthed.} Dersom vi, efter Kants Exempel,
blot skulde anvende Modalitetsdommen til at forvisse os om
de tilsvarende Forstandskategoriers aprioriske Gyldighed, vilde
en udførlig Behandling af disse Domme paa Videnskabens
nuværende Trin vistnok være aldeles overflødig. Men det
er ikke Forvisningen om Kategoriernes Gyldighed, det er en
alsidig Udvikling af de i sig gyldige Kategorier, hvorom her
% --- p. 211 ---
\opage{211}handles. Hegel har nu rigtignok meent, at Kategorierne: Virkelighed, Mulighed og Nødvendighed maatte erholde deres fuldstændige Udvikling i „den objective Logik“, en Udvikling, der i hans System gaaer forud for Læren om Dommen. I denne Anledning maae vi bemærke, at Virkelighedsreflexionen\footnote{Grundideernes Logik. I. S. 105--133.}% note *)
{} med alle derunder hørende Kategorier meget vel lader sig afhandle uafhængig af Læren om Dommen, kun at den ikke ved en vildledende Abstraction behandles som uafhængig af al tænkende og dømmende Subjectivitet; men at den Udvikling af Kategorierne: Virkelighed, Mulighed og Nødvendighed, der svarer til Virkelighedsreflexionen, ingenlunde gjør en til Læren om Dommene svarende Udvikling overflødig, vil allerede skjønnes af, hvad der er anført under den assertoriske Dom, og paatrænge sig endnu stærkere, naar Talen er om den problematiske.

Fra hvilket Synspunkt vi end opfatte Mulighedsproblemet, være sig som almindeligt Væsensproblem, som Objectiveringsproblem, eller særlig som Anticipationsproblem: fra alle Sider viser det sig, at Mulighed og Subjectivitet staae i saa nøie Forbindelse, at Muligheden umulig kan udvikles uden i Reflexion af Subjectiviteten. Det er den i sig frie Subjectivitet, der frigjør Muligheden og giver den Tilværelse; men kun i Ordet ved Dommen er Subjectiviteten fri, i logisk Forstand; kan nu Læren om Dommen først udvikle Subjectivitetens Frihed, og beroer Mulighedens Frigjørelse til ideel Tilværen paa den frie Subjectivitet, saa er det jo indlysende, at Læren om Muligheden maa staae i inderlig Sammenhæng med Læren om Dommen.

Ifølge den assertoriske Dom staaer Virkeligheden antilogisk imod Tænkningen: den antilogiske Umiddelbarhed kan paa mange Maader frastøde, tiltrække, bevæge og bestemme Reflexionen, men den bevægede Reflexion er ikke istand til at
% --- p. 212 ---
\opage{212}opløse Virkeligheden. Ifølge den problematiske Dom faaer Reflexionen fuldstændig Magt med Virkeligheden, intet Virkeligt modstaaer Opløsningen: her abstraheres fra det Antilogiske. Charakteermærket for det Virkelige er Bestemtheden, den umiddelbare, antilogiske Bestemthed; Charakteermærket for det Mulige er Opløsningen af det Virkeliges Bestemthed, altsaa det Ubestemte.

Seet fra den subjective Side er det Ubestemte det Samme som det Uvisse. Uvishedens Standpunkt repræsenteres ved Skepticismen; den til Uvisheden svarende Mulighed ville vi derfor kalde den skeptiske. Fra den objective Side er Ubestemtheden et kategorisk Udtryk for den ikke værende, altsaa enten forsvundne eller endnu ikke tilblevne Virkelighed, Virkeligheden skjult i objectiv Væsenhed; denne som objectiv Væsenhed opfattede Mulighed kalde vi den dogmatiske. Ved at bestemme Ubestemtheden, saavel fra den subjective som fra den objective Side, indsee vi, at Ubestemtheden, og med den Muligheden, oprindelig maa være bestemt til at ophæve sig selv; paa denne Indsigt beroer Mulighedens Ontologie.

1) \emph{Den skeptiske Mulighed: Modalitetens subjective Side.} Paa det assertoriske Grundlag er Umiddelbarhedens Indvirkning paa Reflexionen deels tiltrækkende, bestemmende og bindende, deels frastødende, løsnende og frigjørende; det Dobbeltsidige sees allerede deraf, at den assertoriske Dom er paa een Gang baade umiddelbart objectiv og umiddelbart subjectiv: gjennem det umiddelbart Subjective gaaer det Assertoriske over i det Problematiske. Saalænge det Umiddelbare imidlertid fra en eller anden Side er istand til at magte, bestemme og binde Reflexionen, har det Assertoriske afgjørende Overvægt: det Positive, det Virkelige, sikkres derved, at Tænkningen gaaer i Kjendsgjerningens Ledebaand.

Skal Tænkningen frigjøres, og den reflecterende Subjectivitet besinde sig paa sit uendelige Forhold til sig selv gjennem Andet, da kan den virkelige Begyndelse kun skee derved,
% --- p. 213 ---
\opage{213}at Virkeligheden selv viser sig frastødende mod Reflexionen. Det Problematiske har ikke sit Udgangspunkt i sig selv; det springer ved umiddelbar Conseqvens ud fra det Assertoriske. Man kan egentlig ikke sige, at Skepticismen abstract begynder med Tvivlen; tvertimod! den begynder umiddelbart med den umiddelbare Vished; men da den umiddelbare Vished i sig selv har den negative Bestemmelse, at slaae over i Uvished, saa har Tvivlen i denne Uvishedens Tilbliven naturligviis strax en virkelig Begyndelse.

Hvor umiddelbart og distinct Reflexionen i den gamle Skepticisme tog Sigte paa Virkeligheden i dens Mangfoldighed af forskjelligartede Phænomener, sees allerede ved et flygtigt Blik paa dens „Troper“. Hvad der i „Troperne“ anføres om „Dyrenes Forskjellighed“, „Menneskenes Forskjellighed“, „Sandsernes Forskjellighed“, „de forskjellige Omstændigheder, Afstande, Relationer“ o.\ s.\ v., hvori Tingen forefindes, viser tilstrækkelig, hvad det vil sige, at det Problematiske udgaaer fra det Assertoriske, at Uvisheden har sit Udspring fra en umiddelbar Vished.

Opgaven er at tænke det Virkelige, det vil sige: ved Eftertanken at forvisse sig om den umiddelbare Vished; i denne Opgave skjuler sig en Modsigelse. At tænke det Virkelige er, med andre Ord, at forvandle de enkelte Virkelighedsbestemmelser til Tankebestemmelser. Men Virkelighedsbestemmelserne i deres umiddelbare Enkelthed ere tilfældige, Tankebestemmelserne i deres reflecterede Almindelighed nødvendige; Tankebestemmelserne ere logiske, Virkelighedsbestemmelserne antilogiske: her er Modsigelsen. Denne Modsigelse kommer paa det assertoriske Grundlag tilsyne derved, at Gjenstandene i deres tilværende Forskjellighed foranledige de dømmende Subjectiviteter i deres ligeledes tilværende Forskjellighed til indbyrdes afvigende og modstridende Domme: „Vinden er virkelig kold“; „Vinden er virkelig ikke kold“; det er Mod%
% --- p. 214 ---
\opage{214}sigelsen, den gamle Modsigelse, der allerede lod sig see i den protagoriske Sophistik.

Frastødt af Virkeligheden og tilbagebøiet i sig selv, griber nu Reflexionen den assertoriske Dom fra den subjective Side med forhøiet Energie. Ikke nok, at Reflexionen forvandler Dommen til Mening og Meningen til Paastand; det subjective Element forstærkes indtil det yderste derved, at Mening stilles mod Mening, Paastand mod Paastand, og det saaledes, at den ene Mening, den ene Paastand netop erholder samme Gyldighed, altsaa samme Ugyldighed, som den anden. Den umiddelbart bestemte Virkelighed er paa denne Maade blevet til noget aldeles Ubestemt, noget for Reflexionen Ubestemmeligt.

Den skeptiske Reflexion har ikke paa objectiv Maade gjennemtrængt Virkelighedens Væsen, altsaa heller ikke i objectiv Forstand løst Virkeligheden op; tvertimod! hvad det Virkelige i sig selv er, Noget eller Intet, Væsen eller Skin, Substans eller Accidens, derom veed Skeptikeren Intet, slet Intet. Det er altsaa muligt, at det Virkelige i sig selv er Noget; men det er jo ogsaa muligt, at det er Intet; den ene af de to modstridende Antagelser er i Virkeligheden ligesaa mulig som den anden; Virkeligheden med sine umiddelbare Bestemmelser er altsaa kun i sig selv et Ubestemt, forsaavidt den for Subjectiviteten er ubestemmelig: Opløsningen er subjectiv, den skeptiske Mulighed Uvished.

At enhver umiddelbar Bestemthed og derved enhver Virkelighedsbestemmelse i det Øieblik, den berøres af Reflexionen, forvandler sig til Ubestemthed, at der altsaa om virkelige Ting ikke kan gives nogen sand Vished, er nu den Opdagelse, hvortil den skeptiske Reflexion har ført. I hvilken Forstand denne Reflexion kan siges at betegne et Fremskridt, vil være indlysende, naar vi tage Hverdagsbevidstheden til Maalestok. For Hverdagsbevidstheden gjælder det, at vi have umiddelbar Vished om Noget, medens vi ere i Uvished om Andet, saa at
% --- p. 215 ---
\opage{215}den stedfindende Uvished paa visse Betingelser kan hæves til Vished; for den skeptiske Bevidsthed er den umiddelbare Vished een Gang for alle forsvunden, det vil sige: den er forvandlet til Skin, dens Credit er undergravet; den slaaer over i Uvished, den er et Bedrag. Til Afsløring af dette Bedrag anvendte de gamle Skeptikere en Mængde Exempler, sammenhobede fra alle Virkelighedens forskjellige Hjørner og Kanter. Paa disse Exemplificationer med Gjentagelser, Afvexlinger og Ændringer kan der imidlertid, hvad Principet angaaer, ikke lægges synderlig Vægt; det er en Art empirisk Fyldekalk, vidtløftige Opregninger, beregnede paa at erstatte Savnet af fremadskridende Udvikling og udfylde Tomheden.

„Men“, kunde man spørge, „naar den skeptiske Reflexion, frastødt af Virkeligheden, blev stødt tilbage i sig selv, maatte vi da ikke vente, at den, efter at have tabt Visheden om de existerende Tings Natur og Beskaffenhed, med desto større Energie lagde an paa at forstaae sig selv og komme til Vished om sit eget Væsen?“ Hertil maa svares, at den Maade, hvorpaa Reflexionen her opgiver Virkeligheden, netop gjør den det umuligt at blive sig selv gjennemsigtig. Reflexionen er frastødt af den umiddelbare Virkelighed, altsaa af Phænomenet, og har igjen stødt Phænomenet fra sig; just derved er den blevet sig selv et Phænomen, og kan af den Grund ikke opfatte sig som Væsen. Kun et Phænomen kan tilbagestødes af et Phænomen; Væsenet som Væsen kan Phænomenet ikke tilbagestøde. --- Eleaternes Tanke kunde dog paa en Maade fastholde sig selv i det Ene; men den havde heller ikke ladet sig tilbagestøde af det Mangfoldige; den var ikke uvis om det Mangfoldiges Natur, den var tvertimod ganske vis paa, at det Mangfoldigværende var et Ikke-Værende, og just derfor kunde den være vis paa sig selv i det Eneværende. Ogsaa den fichteske Reflexion gjør det paa sin Maade af med den hele udvortes Virkelighed, idet den gjør Tingen til Skranke for Jeget; men denne Opløsning af det virkelig Objective
% --- p. 216 ---
\opage{216}betegner ingen Uvished; tvertimod! Tanken i den subjective Idealisme faaer netop Kraft til at nedsætte Tingen til Skranke derved, at den i sit Andet formaaer at fastholde sig selv: den veed, hvad Skin er, og derfor veed den ogsaa, hvad Væsen er. Anderledes med den skeptiske Reflexion.

Skepticismen opgiver Virkeligheden og erklærer den for Skin; men da Erklæringen kun har subjectiv Betydning, kun er en Betegnelse for Uvisheden, idet det jo ligesaavel er muligt, at det Virkelige i sig selv ikke er Skin, som at det er Skin: saa tilstaaer Skepticismen, at den ikke veed, hvad Skin er. Men en Reflexion, der ikke veed, hvad Skin er, veed heller ikke, hvad Væsen er, og kan umulig gjennemskue sit eget Væsen.

Det er i denne Henseende charakteristisk, at Skeptikerne opfatte det Indvortes og Intellectuelle, Erkjendelse, Forstand o.\ s.\ v., ligesaa phænomenalt, udvortes og umiddelbart som de sandselige Ting. Saaledes f.\ Ex.\ naar de tale om Muligheden af, at Mennesket skulde kunne erkjende sig selv;\footnote{„Skal nemlig Mennesket erkjende sig selv, maa enten det hele Menneske være det erkjendende Subject, eller det hele Menneske Erkjendelsens Gjenstand, eller en Deel af Mennesket det Erkjendende og en anden Deel Erkjendelsens Gjenstand. Men er det hele Menneske det erkjendende Subject, saa er Erkjendelsens Gjenstand ophævet, og er det hele Menneske Erkjendelsens Gjenstand, saa er det erkjendende Subject forsvundet. Er derimod kun en Deel af Mennesket det Erkjendende og en anden Deel Erkjendelsens Gjenstand, saa maa Mennesket, efterdi dets Dele ere Legemet, Sandserne og Forstanden, enten med Legemet erkjende Sandserne og Forstanden, eller med Sandserne og Forstanden erkjende Legemet“. E. Frisenberg Nielsen. Den antike Skepticisme. Kjøbenh.\ 1862. S.~75.}% note *)
{} Forstanden og Erkjendelsen betragtes i slige Undersøgelser paa samme umiddelbare Maade som Blaat og Guult. For en saa umiddelbar Opfattelse er Forstanden kun et Phænomen; om et saadant Phænomen haves kun umiddelbar Vished, den flygtige Vished, der øieblikkelig er Uvished.

% --- p. 217 ---
\opage{217}Uvishedsreflexionen kan iøvrigt drives til en mærkelig Grad af Skarphed og Fiinhed, uden at den reflecterende Subjectivitet derved kommer Væsenserkjendelsen et Haarsbred nærmere. Saaledes f.\ Ex.\ i Reflexionerne over Aarsag og Virkning. „Dersom Noget er Aarsag, maa enten et Legeme være Aarsag til et andet Legeme, eller et Ulegemligt til et Ulegemligt, eller et Legeme til et Ulegemligt, eller et Ulegemligt til et Legeme“. Færdigheden i at opstille og antage forskjellige Tilfælde er iøinefaldende; det er overalt den phænomenale Umiddelbarhed med den tilsvarende umiddelbare Antagelse, der afgiver Udgangspunktet. Antag kun Noget: enhver Antagelse skal strax give Plads for den modsatte Antagelse; sæt kun et muligt Tilfælde, og det modsatte Tilfælde skal vise sig ligesaa muligt; stil hvilketsomhelst bestemt Tilfælde op som virkeligt: det virkelige Tilfælde skal strax forandre sig til et urimeligt, altsaa til et utænkeligt Tilfælde. „Skal et Legeme i og for sig frembringe et andet Legeme, maa denne dets Væren i sig, hvad der er modsigende, frembringe en Væren udenfor sig og saaledes, hvad der er urimeligt, virke udenfor sin egen Natur. Skal et Legeme derimod i Forbindelse med et andet frembringe et tredie, vil der af eet ikke blot blive to eller af to tre, men et uendeligt Antal, hvilket er umuligt“\footnote{E. F. Nielsen. Den antike Skepticisme. S.~117 \&\ 118.}% note *)
{} o.\ s.\ v.

Den dialektiske Behændighed, hvormed de gamle Skeptikere, for Størstedelen ved ligefrem Afbenyttelse af hvad de forefandt hos Forgjængerne, de store Philosopher, forstaae at opvise Grunde for og imod enhver Antagelse, har ledet Hegel til en Overvurdering af det hele Aandsphænomen. „Kun den antike Skepticisme,“ siger han, „er af sand og dyb Natur, medens den nyere snarere er Epikuræisme. Thi medens den antike Skepticisme netop er rettet imod den forstandige Tænken, der lader den bestemte Forskjel gjælde som sidste, som værende, er den nyere Skepticisme derimod rettet imod Tankerne, imod
% --- p. 218 ---
\opage{218}Begrebet og Ideen, altsaa imod den høiere Philosophie. Den nyere Skepticisme, der paa denne Maade lader Tingenes Realitet staae hen ganske ubetvivlet, er da ikke engang en Bondephilosophie; thi selv Bønder vide, at alle jordiske Ting ere forgjængelige, at deres Væren altsaa er ligesaa god som deres Ikke-Væren“.

Af den heri indeholdte Misforstaaelse giver følgende Udtalelse en, som det synes, træffende og ret tilfredsstillende Charakteristik. „For at forklare den hegelske Overvurdering“, hedder det\footnote{E. F. Nielsen. Den antike Skepticisme. S.~190--192.}, % note *)
„maae vi udtrykkelig mærke os den Forskjel, der er imellem en negativ Brug af Kategorierne og en høiere Bevidsthed om Betydningen af det Negative selv. Med Henviisning til syvende Trope, hvor det t.\ Ex.\ hedder: „Viin styrker, naar den nydes med Maade, men svækker og opløser Legemet, naar den nydes i Overmaal,“ fremhæver Hegel, at Skeptikerne have, om end kun indirect, antydet, at Qvantitet og Sammensætning ere saa langt fra at være noget Ligegyldigt med Hensyn paa Qvalitet og Opløsning, at der med Forandringen i Qvantiteten ogsaa følger Forandring i Qvaliteten. Men at Hegel i denne Sammenhæng læser sin egen Logik ud af de skeptiske Reflexioner, er iøinefaldende. Det er vist ligesaalidt faldet Nogen af de gamle Skeptikere ind, at lægge særligt logisk Vægt paa den væsentlige Forbindelse mellem det Qvantitative og det Qvalitative, som det i det Hele er stemmende med deres Aandsretning, at tillægge nogensomhelst Kategorie væsentlig Gyldighed. For endnu at anføre et andet Exempel paa Hegels oversvævende Fortolkning, henvise vi til følgende Udtalelse hos ham: „Skeptikerne søge overalt at ombytte Udtrykket: „det er“, der i sig indeholder noget Fast, Noget, der vedligeholder sig under alle Omstændigheder, med Udtrykket: „det synes“ eller „Alt er bevægeligt.“ Men ophæve de ogsaa derved den sandselige Lighed og Identitet, og følgelig \emph{denne}
% --- p. 219 ---
\opage{219}Almindelighed, saa indtræder en anden; thi Almindeligheden eller det Værende ligger netop deri, at man t.\ Ex.\ veed, at for den Guulsottige viser det sig guult, hvad der for en Anden t.\ Ex.\ viser sig hvidt o.\ s.\ v. Imod en slig Almindelighed, der jo ikke er sand, fordi den er umiddelbar, paavises derfor med fuld Føie indenfor denne selv dens Ikkealmindelighed ved en anden Almindelighed“.%
% PRINTER: one closing mark here for two open quotations ("„maae vi" p. 218 and the nested Hegel quotation "„Skeptikerne søge"); left unbalanced as printed
{} „Men“ --- saa lyder Indvendingen --- „hele Sammenhængen viser dog tydeligt nok, at Skeptikerne ingenlunde stræbe at sætte den ene Almindelighed imod den anden for derved at klare Bevidstheden om en sandere og høiere Almindelighed, men meget mere, i Overeensstemmelse med deres altopløsende Tendens, at stille den ene Antagelse imod den anden, og just derved gjøre det Almindelige ukjendeligt. Idealiteten af den skeptiske Negativitet overskygges tydeligt af det Bevidsthedens Mørke, hvori Skeptikeren bevæger sig. Der er i den opløsende Skarpsindighed en vis ubevidst --- halvbevidst Blindhed, som gjør, at man let tager Skeptikerens Ord bogstaveligt, naar han yttrer, at han ikke engang veed, om han veed Noget“.

„Skepticismen bliver“, ifølge Hegel\footnote{Jvfr. Vorlesungen über die Geschichte der Philosophie. Zweiter Theil. Werke. Vierzehnter Band. S.~475 \&\ 514.}, % note *)
„staaende ved det Negative som Resultat og miskjender derved, at Negationen har et bestemt affirmativt Indhold i sig, hvorimod den speculative Methode bestemmer Negationen som den sig til sig selv refererende Negation, som Negationens Negation, eller nærmere som den uendelige Affirmation. Skepticismen udøver sin Dialektik efter Tilfældighed, da den, alt eftersom Stoffet frembyder sig for den, i dette paaviser, at det i sig er det Negative, medens den speculative Methode derimod udøver sin Dialektik i Kraft af Nødvendighed, idet den godtgjør, at Ideen som speculativ Idee hverken er et ligefrem Positivt eller et ligefrem Negativt, men en saadan Forening af begge, at det Negative og det Positive i denne Forening overalt reflectere hinanden.“

% --- p. 220 ---
\opage{220}Den skeptiske Negation adskiller sig i en Vrimmel af isolerede Negationer, ligesom den skeptiske Tilfældighed i en Vrimmel af særskilte Tilfælde. Til Antagelsen af hvert enkelt Tilfælde svarer det Assertoriske; til Forkastelsen af hvert antaget Tilfælde det Problematiske. Begyndelsen er allevegne, thi her begyndes overalt med et Tilfælde; Enden er allevegne, thi Visheden om ethvert Tilfælde forvandler sig overalt til Uvished. Det er de assertoriske og problematiske Functioners negativt pointerende Vexelspil, der charakteriserer Skepticismen: kun i dette Vexelspil kan Uvisheden uophørlig vorde.

„Vel kunde det synes, som om hans\footnote{E. F. Nielsen. Den antike Skepticisme. S.~3--5.}% note *)
{} (Skeptikerens) Søgen førte idetmindste til det bestemte Resultat, at ingen Sandhed lader sig finde. Men at lade Skeptikeren slaae denne Sætning fast, er atter at misforstaae Skepticismen og at forvexle den med en negativ Dogmatisme. Den thetiske Opfattelse af det Negative er jo netop charakteristisk for den negative Dogmatisme, idet den ikke blot er forvisset om, at Noget er uvist, men tillige, at det er vist, at Noget er uvist. Den negative Dogmatisme veed da, hvad der er vist, og hvad der er uvist, og er saaledes vis paa sin Sætning og dens Modsætning. Skeptikeren derimod seer med sit Dobbeltsyn overalt Fordoblingen, selv i Sandhedens negative Udtryk. Et Udsagn som: Alt er uvist, er derfor, efter ham, selv uvist, og dette Sidste atter uvist, og saaledes i det Uendelige. Paa tilsvarende Maade forholder det sig med Udtrykkene: Alt er falsk, Intet er sandt o.\ s.\ v. \dots\ Men selv om Aporetikeren nok saameget angiver Uvisheden, vil Angivelsen altid faae Skin af en Fastsættelse, en Position, eftersom det jo maa blive ved et Udsagn, da et negativt Udtryk ligesaa godt kan staae fast som et positivt. Han maa dog, som det synes, indrømme og paastaae, at han har en sikker og vis Erkjendelse af Er%
% --- p. 221 ---
\opage{221}kjendelsens Usikkerhed og Uvished. Som den, der nu vil undgaae Beskyldningen for ved sit blotte Udsagn at ophæve Grundenes og Modgrundenes Ligevægt, \dots.\ kalder Skeptikeren sig Ephektiker, \dots.\ thi alene ved at gjøre Epoche, at holde ethvert Bifald, ethvert Udtryk tilbage, d.\ e.\ med andre Ord, at tie, kan han værge sig imod enhver Misforstaaelse og tilfredsstille sin ideale Fordring. Tausheden er altsaa den høieste Opløsning i Uvished. Det er ikke en Taushed, hvori man hviler som i Inderlighed og Selvfordybelse, thi denne rene Taushed er ikke skeptisk; men den Taushed, hvori det Negative fuldbyrdes, idet den viser, at Intet afgjøres, den Taushed, hvori der er Uafgjorthed, Uvished, hvori der er Mulighed.“

Det er Skeptikerens Opgave at bringe Ligevægt mellem de modstridende Antagelser, de indbyrdes modsigende Meninger, en Ligevægt, der betyder, at det Ene er ligesaa rimeligt, og derfor ligesaa urimeligt, som det Andet. Naar denne Ligevægt er tilveiebragt for det enkelte Tilfælde, maa der strax vise sig et nyt Tilfælde, at Forhandlingen kan begynde forfra. Gaves her kun eet eneste Tilfælde, saa maatte jo Opmærksomheden vedblivende fæste sig paa den ved Ligevægten opnaaede Uvished: det maatte da være vist, at Uvished een Gang for alle var opnaaet; denne staaende Uvished er ufordragelig. Ligefrem benegte Uvisheden kan Skeptikeren ikke; thi ved en Benegtelse af Uvisheden, vilde han bekræfte Visheden. Ligefrem bekræfte Uvisheden kan han heller ikke; thi en Bekræftelse af Uvisheden vilde da selv give en Vished. Alle Modsigelser kan Uvisheden udligne; \emph{alene Modsigelsen af Vished og Uvished kan den ikke udligne}: sættes Visheden, saa er Uvisheden udelukket; sættes Uvisheden, saa er Visheden udelukket. Hver Gang Skeptikeren skal acquiescere ved en sidste afgjørende Uvished, gjør han Epoche; han tier, og Tausheden udtrykker, at Uvisheden fastholdes, ikke i Vishedens, men i Uvishedens Form. At tale vilde i slige Tilfælde være en Overilelse; at
% --- p. 222 ---
\opage{222}erklære sig ligefrem for Uvisheden, vilde være at sætte den i Vishedens Form, antage Uvisheden for det Visse, altsaa modsige sig selv.

Alt er uvist; dette er Skepticismens Hovedsætning. Men Visheden af, at Alt er uvist, faaer Ingen ham til at fastsætte. Nei, han fastsætter Intet, han tier; Tausheden betyder, at han aldeles ikke kan indlade sig paa \emph{almindelige} Besvarelser af \emph{almindelige} Spørgsmaal. At indlade sig paa en almindelig Besvarelse af et almindeligt Spørgsmaal vilde for Skepticismen være det Samme som at afgive en definitiv Erklæring; men en saadan eneste definitiv Erklæring kuldkaster hans hele Skepticisme.

Men ligger der da i denne Sætning: Alt er uvist, ikke en definitiv Erklæring? Ingenlunde! at sige i skeptisk Forstand. Istedenfor at indlade sig med det Almene og sige: Uvisheden er almeengjældende, gaaer Skeptikeren løs paa det Enkelte, og trøster sig til, at han i ethvert enkelt Tilfælde nok skal faae Uvisheden ud, nok skal holde den i fortsat Vorden: kun i fortsat Vorden er Uvisheden uendelig.

Den skeptiske Mulighed er uendelig Uvished; \emph{den uendelige, d.\ v.\ s.\ endeløse Uvished, er i uendelig Vorden}. Ender Vorden i Tilværen, faaer Uvisheden, saa at sige, Lov til at staae stille og sætte sig ned i en Tilværelse, saa er det i Reflexionen ude med den: en værende Uvished, en Vished om, at Uvisheden er værende, er en i sig negativt reflecteret Vished. Med Uvishedens Forvandling til negativt reflecteret Vished er Skepticismen forvandlet til negativ Dogmatisme: den skeptiske Mulighed er bleven dogmatisk.

2) \emph{Den dogmatiske Mulighed: Modalitetens objective Side.} Ubestemtheden maa altid have sit Grundlag i en eller anden Bestemthed, Uvisheden i en Vished, de problematiske Functioner i et Ikkeproblematisk. Saalænge det Ikkeproblematiske falder sammen med Virkelighedens umiddelbare Bestemthed og den til samme svarende Vished, er Grundlaget
% --- p. 223 ---
\opage{223}assertorisk. Opgiver den paa sine egne Reflexioner reflecterende Subjectivitet den assertoriske Umiddelbarhed, maa Reflexionen selv sætte og forudsætte sit reflecterede Grundlag.

Det reflecterede Grundlag er nu for det Første betinget af, at det Umiddelbare bortskaffes; dette er Sagens negative Side. Tages den umiddelbare Bestemthed bort, saa haves kun Ubestemthed; tages Visheden bort, saa er her kun Uvished. Hvorledes skal da det Negative nu ved sig selv kunne afgive et Positivt? Den i sig reflecterede Ubestemthed er selv negativ Bestemthed; den i sig bestemte Uvished er negativ bestemt Vished; det nye Grundlag er altsaa en i Ubestemtheden vunden Bestemthed, en i Uvisheden affirmeret Vished: istedenfor Virkelighedens Umiddelbarhed have vi faaet en Tankens Umiddelbarhed; dette er Sagens positive Side.

Det reflecterede Grundlag er paa Grund af Negationen og „Negationens Negation“ uendelig affirmativt, men ikke antilogisk positivt; det er i Virkeligheden selv ikke et Ydre, men et Indre; ikke en Phænomenbestemmelse, men en Væsenhed. Kun i den virkelige Væsenhed kan Virkelighedsreflexionen komme væsentlig til sig selv; den virkelige Væsenhed, den i Reflexion anticiperede Virkelighed, er dogmatisk Mulighed.

Den skeptiske Mulighed var negativ saavel med Hensyn paa Objecterne som paa den reflecterende Subjectivitet; den dogmatiske Mulighed er affirmativ saavel for Objecterne som for Subjectiviteten. Den affirmative Mulighed er objectiv som Tingenes Mulighed, den Mulighed, hvoraf det Virkelige fremgaaer, den Mulighed, der kun optager Betingelserne for at blive virkelig, og kun bliver virkelig, idet den selv paa een Gang naaer og overskrider sin egen Grændse, thi Grændsen for Mulighed er jo Virkelighed. Den affirmative Mulighed er subjectiv som Tingenes Tænkelighed, thi kun i Tænkelighedens, Tankens, Begrebets Form kan en Virkelighed med alle sine Bestemmelser være til som Mulighed. Den affir%
% --- p. 224 ---
\opage{224}mative Mulighed faaer sit særlige dogmatiske Præg derved, at de supponerende Functioner gaae over i de problematiske.

Ligesom den skeptiske Mulighed fremkom derved, at den problematiske Function udsprang af den assertoriske, saaledes fremkommer nu den dogmatiske Mulighed derved, at det Problematiske udspringer af supponerende Functioner. Hvad de assertoriske Functioner ere i Forhold til umiddelbar Virkelighed, ere de supponerende i Forhold til Væsenhedens reflecterede Umiddelbarhed. Bevægelsen er denne: her sættes ved logisk Underskydning en Tankebestemmelse i reflecteret og dog umiddelbar Eenhed med en tilsvarende Værensbestemmelse: den umiddelbare Eenhed af Tænken og Væren er dogmatisk. Er den paagjældende Værensbestemmelse opstillet, saa følger Prøvelsen; Reflexionen stiller sig overfor sin egen Forudsætning som for et logisk Object; det Tænkte forestilles, det Forestillede tænkes; dette er Uroen; i denne Uro ligger Negativiteten: den paa een Gang forestillede og dog tænkte Negativitet vækker det Problematiske.

For den dogmatiske Reflexion er det Mulige i al Almindelighed Eet med det Tænkelige; den supponerede Eenhed af Tænkelighed og Mulighed svarer ligefrem til den supponerede Eenhed af Tænken og Væren: som det Tænkelige er Muligheden subjectiv, som det Mulige er Tænkeligheden objectiv. Deraf kommer det, at Prøven paa det Mulige falder sammen med Prøven paa det Tænkelige, ligesom Prøven paa det Umulige falder sammen med Prøven paa det Utænkelige. Kjendemærket paa det Tænkelige er Frihed for Modsigelse; det samme Kjendemærke gjælder for det Mulige (\textit{possibile est, quod non implicat contradictionem}). Kjendemærket paa det Utænkelige er Modsigelsen; det Samme gjælder for det Umulige (\textit{impossibile est, quod implicat contradictionem}). Men denne supponerede Eenhed af Tænkeligt og Muligt, af Utænkeligt og Umuligt, fremkalder nu Tvivl om Mulighedens objective Væsen.

% --- p. 225 ---
\opage{225}„Alt det, i hvilket ingen Modsigelse ligger, er muligt, enten det saa er til eller ikke. Ja alt det, som er til, har først været muligt, førend det blev til: ellers kunde det aldrig komme til nogen Virkelighed. Det, der har været til, er endnu muligt, skjøndt det ei er længere til, og det bliver muligt uden Ende \dots\ Et Menneske var aldrig bleven til, havde det ikke været muligt, at hans Forældre kunde have født ham til Verden; og, efterat han er død, kan man sige det Samme, han havde aldrig levet, havde hans Liv ikke været muligt. Det Mulige er følgelig evigt, og enhver Mulighed er evig, fordi der lader sig ingen Tid tænke, hverken før eller efter, paa hvilken den enten skulde begynde, eller høre op igjen“. Disse Ord af Eilschow, en selvstændig Tænker\footnote{Fr. Chr. Eilschov. Philosophiske Breve. Kjøbenhavn 1748. S.~147.}% note *)
% NOTE: name spelt "Eilschow" in the text, "Eilschov" in the footnote, both as printed.
{} fra den leibnitz-wolfiske Periode, ere vel skikkede til at charakterisere den dogmatiske Mulighed. Altsaa: Muligheden er evig; den er en Tankens Overeensstemmelse med sig selv i Form af en Værens Overeensstemmelse med sig selv; det i sig Overeensstemmende stemmer i al Evighed overeens med sig selv, enten der er Nogen, der tænker derpaa eller ikke.

Først bliver altsaa det Mulige sat identisk med det Tænkelige; dernæst bliver det Tænkelige, paa Grund af Tankens supponerede Eenhed med Væren, sat i det Muliges Form saa værende, at Dogmatikeren geraader i Tvivl, om han skal lade Muligheden evig beroe paa sig selv, eller gjøre det evig Tænkelige afhængigt af en evig Tænkende. Den evig Tænkende er Gud; det ligger i den Umiddelbarhed, hvormed de supponerende Functioner forudsætte Gud, at Problemet faaer et nok saa theologisk som logisk Anstrøg. Man kunde da fristes til at emancipere Muligheden fra Gud. „Dog skal vi“, siger vor dogmatiske Philosoph, „ikke gjøre enhver Mulighed til en liden Guddom, der har Grunden i sig selv til sig selv, hvilket synes at flyde af det, jeg har sagt. Nogle
% SEAM: quotation „ikke gjøre enhver Mulighed ... continues on p. 226 (no word break at page foot).
% --- p. 226 ---
\opage{226}gamle Skole-Vise ere faldne paa de vrange Tanker, at Mulighederne ikke hængte af Gud eller havde Grund i hans Væsen, fordi man kunde tænke paa dem og bevise dem, uden at forestille sig, at de havde deres Grund i noget evig-nødvendigt Væsen, eller i Gud. En Trekant er mulig, fordi en Figur kan indesluttes af tre Linier: man kan hverken sige, at disse tre Linier eller deres Sammenstødelse er mulig i ligesaa mange Punkter, fordi der er en Gud til. Det synes meget mere, at denne Mulighed var unegtelig, om han end aldrig var til; og en Atheist, som negter Gud, negter dog ikke Trekantens Mulighed. Hvad skal vi da fastsætte?“

Tre foreløbige Muligheder indstille sig nu angaaende de evige Muligheder: enten ere de grundede i sig selv, eller i Guds Villie, eller i den guddommelige Forstand. „Skal vi sige med Cartesius, at de ere grundede i Guds Villie og Velbehag alene, saa at en Ting er mulig, blot fordi Gud vil, at den skal være det, saa maa og tilstaaes, at den samme Ting kan blive umulig, saa snart han finder det for godt. En Trekant er da bleven mulig, da Gud vilde, at tre Linier i ligesaa mange Punkter skulde kunne indslutte et Rum; men hvordan kunde det nu skee, at det Samme kunde blive umuligt eller hvorledes kunde Gud beslutte, at tre Linier ikke længere herefter skulde kunne sammensætte en Figur? Tilstaaer man den sidste Beslutning at være umulig, som man ikke andet end kan tilstaae, saa maa den første, som gjør Trekanten mulig, være evig; bagefter i det Ringeste: det hedder, hvad Gud een Gang har villet at skulle være muligt, maa siden forblive det i alle Evigheder \dots\ Men det er endnu ikke Alt, hvad man kan erindre. Lige saavel som Cartesius maa tilstaae, at det, der en Gang er bleven muligt efter Guds Villie, forbliver det altid, saa kan vi og paastaae, at det, der nu er muligt, har været det stedse; thi ligesaa stor Vold, som man maa gjøre paa Forstanden for at forestille sig, at den Tid kunde komme, da en Trekant ligesaa lidet skulde kunne indeslutte et Rum,
% --- p. 227 ---
\opage{227}som to rette Linier nu kan gjøre det, saa aldeles maa man og voldtage Forstanden for at forestille sig, at der har været en Tid, da Gud kunde have gjort en Tokant mulig og en Trekant umulig. Heraf flyder, at ligesom man maa tilstaae Mulighedernes Evighed bagefter, saa maa man og tilstaae den forved, eller at alle Muligheder ere uden Begyndelse og Ende uforanderlige, nødvendige, evige“.\footnote{Eilschov. Philosophiske Breve. S.~147--49.}% note *)

Hvad her siges om de evige Muligheder passer dogmatisk paa alle de abstract aprioriske Begreber, de saakaldte evige Sandheder (\textit{veritates æternæ}). Men uagtet Muligheden i sit almene Begreb er evig, og i denne sin Evighed ligegyldig mod Forskjellen mellem virkelig Væren og Ikke-Væren, saa kan Muligheden i Forhold til Tingen umulig være ligegyldig mod Tingens mulige Tilværen. % UNSURE: Tilværen or Tilvaren (faint æ, as in "Væren" on the line above)
Gjøres Mulighederne uafhængige af Gud, og ere Tingene kun virkeliggjorte Muligheder, saa synes Tingene jo med det Samme at blive uafhængige af Gud: hvorledes er dette nu muligt? „Tingenes Muligheder ere afhængelige af Gud, men ikke af hans blotte Villie og Velbehag. De ere grundede i hans uendelige Forstand og sammes evige Forestillinger om Alt det, der kunde frembringes ved hans Almagt. Var ikke Gud til, saa var ingen Ting mulig, det er sandt nok; men derfore er en Mulighed ligesaa lidet skabt af Gud, som han kan siges at have skabt sin egen Forstand“.\footnote{Anf. Skr. S.~150.}% note **)
{} --- Den dogmatiske Mulighed er vel et evigt Begreb, men med al sin Evighed dog, mærkelig nok, et uroligt Begreb. Kunne Tingenes Muligheder ikke beroe paa Guds „blotte Villie og Velbehag“, efterdi det evige Begreb er hævet over al Vilkaarlighed; kunne Mulighederne ikke være uafhængige af Gud, efterdi ingen Ting vilde være mulig, dersom Gud ikke var til, saa maae de vel være i den guddommelige Forstand. Men ogsaa ved denne Antagelse støder den
% --- p. 228 ---
\opage{228}dogmatiske Reflexion paa Vanskeligheder. Skulde Tingenes Muligheder afhænge af den guddommelige Forstand, fordi alle mulige Ting afhænge af Gud, da maatte Gud, der har skabt Tingene, jo ogsaa have skabt Mulighederne; men det er aldeles umuligt. Forstanden kan da ikke bestemme Muligheden; Mulighedsbegrebet er objectivt: Forstanden har at respectere Begrebets uforanderlige Objectivitet. „Jeg siger, det er nok, at Tingenes Egenskaber ikke modsige eller omstøde sig selv indbyrdes; derfore ere de mulige, og de behøve ingen anden Grund til deres Mulighed, end just denne Umodsigelse. Eller hvorledes skal jeg begribe Sammenhængen i denne Slutning: den eller den Ting har været, er og bliver evig mulig; derfore forefindes der et Væsen, som evig forestiller sig dens Mulighed. Vore Forfædre har ikke tænkt, at mange af de Ting, som de Nyere har beviist i mechaniske og physiske Videnskaber, har været mulige, førend Kunsten i vore Dage har opdaget dem; men har de derfore ikke været det? De Gamle kjendte, for Exempel, ikkun ganske faa elektriske Legemer, og har ikke forestillet sig, at Vandet lod sig elektrisere \dots.; dog var Naturen da, ligesom nu, opfyldt med elektriske Legemer. Vi vide det; men de vidste det ikke. Deres og vores Forestilling forandrer lige saa Lidet i Naturen, som man kan sige, at et Speil i en Stue gjør Tingene, som findes deri, eller at der er Stol, Bord, Bænke, fordi man kan see alt dette i Glasset \dots. Naar vi sige, at den uendelige Forestillingskraft forestiller sig uendelige Ting, vores derimod ikkun visse Ting \dots. saa sætte vi al den Forskjel, som der er imellem Guds og vores Forstand. Hiin forestiller sig paa een Gang alle mulige Ting, denne ikkun nogle faa, og dem endda efterhaanden; hiin er uindskrænket, denne er indskrænket, o.\ s.\ v. Men i alt dette finde vi endnu ingen Grund i Guds Forstand til Tingenes Mulighed. Saa meget viser den berømte Wolf, at Guds Villie uden foregaaende Nødvendighed ikke kan have opdigtet de Ideer, i hvilke Mulighederne indeholdes; siden de ere ufor%
% --- p. 229 ---
\opage{229}anderlige og derfore evige; men det er langt fra ingen Følge, at de derfore ere grundede i hans Forstand.\footnote{Eilschov. Philosophiske Breve. S.~154--55.}% note *)
“

Af disse Udtalelser vil det kunne sees, hvorledes Forholdet er mellem Dogmatismens underskydende og dens problematiske Functioner. Den dogmatiske Tænkning er hypostaserende. Først gjør den Mulighederne til rene Tankevæsener og giver dem en objectiv Tilværen; ved Siden heraf sætter den Guddommen som den absolute Substans; og nu opkommer det Spørgsmaal, hvorvidt Mulighederne ere i Gud eller udenfor Gud.

Saaledes kan den dogmatiske Reflexion, ifølge Ovenstaaende, opfatte det Mulige som det Tænkelige, stirre paa dette i sig Tænkelige som paa et Værende, uden at indsee Umuligheden af, at det blot Tænkelige skulde kunne have anden Tilværen end den, det erholder ved Tænkningen: Tankens logiske Objectivitet imponerer den dogmatiske Tænker. Det Mulige, jeg tænker mig, siger Dogmatikeren, bliver ikke muligt, fordi jeg tænker det; men jeg tænker det muligt, fordi det --- uafhængigt af min Tænkning --- er muligt. „En Trekant er mulig, endog førend jeg lærer at forestille mig den i Tankerne. Jeg kan siden forestille mig den; men Figurens Mulighed er ældre end min Forestilling, og denne indeholder ikke Grunden til hiin. Skriver jeg den paa en Flade, saa har den mig at takke for det, den er til; men ingenlunde min foregaaende Forestilling for det, den er mulig\footnote{Anf. Skr. S.~155.}% note **)
“. Tænkning og Forestilling brydes med hinanden. Idet Tænkningen objectiverer sig sine Begreber, antage disse Begreber en forestillet Objectivitet; ved den forestillede Objectivitet blive Tænkningens egne Frembringelser uafhængige af Tænkningen: den dogmatiske Tænkning kommer fremmed til sit Eget og bliver derved fremmed for sig selv.

% --- p. 230 ---
\opage{230}Saalænge Forestillingen brydes med Tanken, vil et forandret Udgangspunkt bestandig føre til et forandret Resultat: her er i Dogmatismen som i Skepticismen Mulighed til en Mængde høist forskjellige Udgangspunkter og derfor ogsaa Mulighed til en Mængde indbyrdes modstridende Resultater. Forskjellen er kun den, at medens den skeptiske Reflexion finder sig frastødt af ethvert Resultat, fordi den i ethvert Resultat støder paa det Antilogiske, vil den dogmatiske Reflexion hver Gang finde Selvbekræftelse i sit vundne Resultat, efterdi dette Resultat ikke er Andet end dens egen, ved supponerende Functioner underskudte, logiske Væsenhed. Hvad der uendelig interesserer Skeptikerne, at et Problematisk uophørlig resulterer af enhver mulig Antagelse, det overrasker og piner Dogmatikeren. Den dogmatiske Tænker mener hvert Øieblik, efter endt Beviisførelse, at have fundet noget Fast, Noget, hvorved han kan slaae sig til Ro, og agter ikke paa, at Modsigelsen ulmer i Muligheden, at Muligheden er som en Slange i Græsset.

Gaaer Dogmatikeren f.\ Ex.\ ud fra den Betragtning, at Tingenes Mulighed kun kan være til forud for Tingen derved, at den har Tilværelse i en af Tingens Virkelighed uafhængig Forestilling, saa vil han snart indsee, at det er Forestillingen, der gjør Muligheden uafhængig af Tingen. Men den almindelige Forestilling kan umulig forestille sig selv: enhver Forestilling maa forudsætte en forestillende Subjectivitet. Altsaa, idet den af Tingen uafhængige Mulighed er afhængig af Forestillingen, og Forestillingen af Subjectets Forestillen, saa maa Mulighedens Uafhængighed af Tingen jo netop være begrundet i dens Afhængighed af Subjectiviteten: de evige, af alle mulige Ting uafhængige Muligheder maae være begrundede i en evig af alle mulige Ting uafhængig, tænkende, villende Substans. Nu er det altsaa vist, at de evige Muligheder ere i Gud.

Men i denne Vished skjuler sig igjen adskillig Uvished.
% --- p. 231 ---
\opage{231}Thi naar den guddommelige Villie sættes som dogmatisk positiv Realitet ligeoverfor eller ved Siden af en tilsvarende Realitet, nemlig den guddommelige Forstand: saa er det jo endnu uvist, hvorvidt de evige, i Gud værende Muligheder skulle antages at beroe enten paa Guds Villie eller paa den guddommelige Forstand. Alt kommer her an paa Udgangspunktet: man kan dogmatisk gaae ud fra Forstanden, men man kan jo ogsaa gaae ud fra Villien. Her staae, som vi have seet, Wolf og Cartesius skarpt imod hinanden. Wolf antager, at alle Muligheder ere grundede i den guddommelige Forstand, Cartesius, at de beroe paa Guds Villie. Forfatteren af de „Philosophiske Breve“ slutter sig nærmest til Wolf, men forsøger en egen Begrundelse. „Det hjælper ikke,“ siger han, „at byde en aristotelisk Modstander Spidsen under det sædvanlige philosophiske Harnisk, med hvilket Nogle mener at kunne afværge alle Stød. Det hjælper ikke, vil jeg sige, at komme frem med den almindelige Grundregel: \emph{Alt hvad som er, maa have en tilstrækkelig Grund, hvorfor det heller er, end ikke er}, og derpaa at slutte saaledes: enhver Mulighed maa have en Grund, hvorfor den heller er en Mulighed end en Umulighed: denne kan den ikke indeholde i sig selv; ellers var den Gud, følgelig i noget Andet“ \dots.\ „bemeldte Hovedgrund angaaer alene de virkelige, og ikke de mulige Tings Rige.“ Forfatterens Overeensstemmelse med Wolf er da, ifølge det anbragte Correctiv, begrundet i en Uovereensstemmelse; Resultatet er det samme, kun Veiene ere forskjellige. Men den dogmatiske Adskillelse af Vei og Resultat fører just ind i nye Forviklinger og nye Modsigelser. I det anbragte Correctiv skjuler sig en Modsigelse.

Blendet af det rene Begrebs formelle Objectivitet sætter „Brevenes“ Forfatter den evige Mulighed selvstændig imod al Virkeliggjørelse; stolende paa Modsigelsens Grundsætning mærker han ikke den Modsigelse, at en Mulighed, der mangler Tilværelsesgrund, og derfor, saalænge Grunden mangler, al%
% --- p. 232 ---
\opage{232}drig kan blive virkelig, er umulig. Imellem Muligheden og Virkeligheden indlægger han Grunden, og mærker ikke, at Muligheden selv trænger ligesaa høilig til Grunden, for at blive mulig, som „de virkelige Tings Rige,“ for at være virkeligt. „Alt det, som er, navnlig alt det, som er virkeligt, alt det, som er til, maa have en Grund, hvorfor det er til, og hellere paa den Maade end paa nogen anden; men ikke saalænge det forbliver i Mulighed. --- Saalænge det alene er muligt, har det Grund nok i sig selv til sig selv“; --- \emph{uden Gud ere Tingene mulige}. „Gud forestiller sig nødvendig alle Muligheder, saadanne som de ere, det er som Muligheder, som Ting, i hvis Egenskaber og Beskaffenheder ingen Modsigelse findes: videre Grund behøve de ikke. Skal nogen af dem blive til udenfor Guds Forestilling, saa maa Gud selv bevise, at den er mulig, saa maa han skabe, saa maa han ville, at den skal være til\footnote{Eilschov. Philosophiske Breve. S.~156 flg.}% note *)
“: \emph{uden Gud ere Tingene umulige}. Altsaa: \emph{uden Gud ere Tingene mulige; uden Gud ere Tingene umulige}; dette er Modsigelsen.

Efter at Eilschov paa denne Maade ved at krydse mellem Modsigelser, forfulgt af Modsigelsen, har bragt Bevægelse ind i Wolfs stivnende Syllogisme, gaaer han saa over til at fremsætte sin Anskuelse af Sagen, og kommer til et Resultat, der ret er betegnende for den dogmatiske Reflexion; thi i Resultatet sees igjen Modsigelsen. Han vil, som han udtrykker sig, vise, „at alle Muligheder ere grundede i Guds Forstand, og at de følgelig har Grund i Guds Mulighed alene, over hvilken ingen er.“ Først oplyser han Sagen ved et Exempel af „Tal og Størrelser“, thi herom gjælder det, at „det Mindre altid er indesluttet i det Større indtil den første Eenhed.“ „Eetets Mulighed er ikke grundet uden i sig selv; men det har Grunden i sig til alle andre Tals Mulighed;“ det er derfor saa træffende, „at Pythagoras har kaldet Gud
% --- p. 233 ---
\opage{233}med et Navn, som er fuldt af Eftertryk, \emph{Eenhed}“. Nu følger Forklaringen. Guds Mulighed er ene og alene grundet i sig selv, fordi den er bestemt ved sig selv uden Hensyn til Andet, ligesom Eenheden, den oprindelige Eenhed, er bestemt ved sig selv uden Hensyn til noget endeligt Tal. Skulde nu en Tings Mulighed tænkes uafhængig af Guds Mulighed, da maatte den ligeledes tænkes at være bestemt ved sig selv „uden nogen anden Muligheds Hjælp“: men dette er paa Grund af Tingens Endelighed umuligt. „Naar jeg spørger, om en Ting er mulig, saa bliver jeg urimelig, dersom jeg ikke først sætter den Tings Beskaffenhed, om hvilken jeg spørger.“ Det Endelige er i sine Grændser ubestemt, og maa derfor bestemmes; kun idet det bestemmes, bliver det muligt. Det er vel i Almindelighed muligt, at en Linie kan „ligge jævnt imellem sine Punkter“, men en uendelig ret Linie er umulig. „Naar man da nævner en uendelig ret Linie, saa siger man en Linie, som skal have Punkter, det er Skranker, ellers er der ingen Linie, og som tillige ingen skal have, thi ellers er den ikke uendelig, det er, \emph{man siger sig selv imod}.“ --- Ved disse Betragtninger kommer Forfatteren af de „Philosophiske Breve“ da til det Resultat, at Tingenes Muligheder indeholdes i Guds Mulighed, ligesom de endelige Tal i Eenheden. „Var der ingen uendelig Forestillingskraft mulig, saa var ingen endelig Ting mulig: fordi enhver endelig Mulighed ikke af Fornødenhed er bestemt; \emph{den maa bestemmes først af det Uendelige}. Denne Bestemmelse kan ikke skee uden i Forestillingen, saalænge det Mulige ikke er virkeligt, eller er til. \emph{Følgelig har endelige Tings Muligheder deres Grund i den uendelige Forestillingskraft og ingenlunde i sig selv}\footnote{Eilschov. Philosophiske Breve. S.~160--61.}% note *)
“.

Men en Slutning, som paa denne Maade skal løse Problemet, vil netop paa Grund af Maaden --- ved Func%
% --- p. 234 ---
\opage{234}tionernes Modalitet kommer Alt an paa Maaden --- være en ny Begyndelse til det Problematiske. Dersom nemlig Tingenes Mulighed maa bestemmes ved en uendelig Forstand, fordi det Endelige, det i sig Ubestemte, men dog Bestemmelige, kan bestemmes paa uendelig mange Maader, saa maa der jo, dogmatisk forstaaet, være noget Vilkaarligt i den Maade, hvorpaa den uendelige Forstand udtrykkelig bestemmer hver enkelt Mulighed: ere de enkelte Muligheder afhængige af den guddommelige Forstand, da ere de med det Samme afhængige af den guddommelige Villie. Paa denne Maade, --- og Alt kommer an paa Maaden --- kan Striden mellem Cartesianere og Wolfianere aldrig ende.

\emph{Thesis. Det er Guds Vilkaarlighed, der har bestemt de endelige Muligheder. Beviis.} Det Mulige er i sin endelige Bestemthed modsat det Nødvendige; thi det Mulige bestemmes kun \emph{saaledes}, fordi det ogsaa kunde bestemmes \emph{anderledes}. Hvad der bestemmes \emph{saaledes}, skjøndt det kunde bestemmes \emph{anderledes}, bestemmes med Vilkaarlighed; altsaa \dots

\emph{Antithesis. Guds Vilkaarlighed har ikke bestemt de endelige Muligheder. Beviis.} Nødvendigt er det, hvis Modsatte er umuligt eller modsigende. Nu er det Modsigende, at det Mulige skulde være umuligt; følgelig er enhver Mulighed nødvendig. Det Nødvendige kan ikke bestemmes vilkaarligt; altsaa \dots

Det er charakteristisk, at Kant, under sit Angreb paa Dogmatismen, opstillede „Antinomien“. Antinomien er Dogmatismens Conseqvens; dens Hemmelighed er en Ombytning af væsentlige Udgangspunkter: i Antinomien bliver den dogmatiske Mulighed til Umulighed; Opgaven er at faae denne Umulighed igjen forvandlet til Mulighed.

3) \emph{Mulighedernes Ontologie: dobbeltsidig Modalitet.} Det er blevet betragtet som Tegn paa den hegelske Dialektiks Overlegenhed over „den endelige Forstand“, at den
% --- p. 235 ---
\opage{235}har kunnet løse de kantiske Antinomier. Paa en vis Maade er denne Overlegenhed umiskjendelig. I to Punkter viser Kant sig hildet i dogmatiske Fordomme. Han er fixeret af Umiddelbarheden i de dogmatiske Udgangspunkter. I hans Thesis gjælder et, i hans Antithesis et andet og modsat Synspunkt; mellem de modsatte Synspunkter gives her ingen væsentlig Overgang: man kommer ved et Spring fra det ene til det andet. Han er fixeret af Umiddelbarheden i den dogmatiske Grændse. Den dogmatiske Grændse er en fast, værende Grændse, sikkret ved en Distinction, der støtter sig paa Modsigelsens Grundsætning --- og paa \textit{principium exclusi medii}.

I begge Henseender har Hegels Logik belyst Vildfarelsen og ved denne Belysning indlagt sig udødelig Fortjeneste. Den hegelske Logik kan, hvad Formen angaaer, med fuld Ret kaldes en Mulighedernes Logik. Mulighed er i Grunden Overgang; den, der ikke kan forstaae den hegelske Logik, kan ikke i Begrebet forstaae nogen væsentlig Overgang. Der kan være Meget at indvende imod mange af denne Logiks Overgange i det Enkelte; men Et er den ved bestemte Kategorier bestemte Overgang, et Andet er Overgangsbegrebet: det almindelige Overgangsbegreb er Muligheden. Vilde Nogen indvende, at det ikke er Muligheden, men Vorden, der er den hos Hegel egentlige og første Overgangskategorie, saa maa man hertil svare, at Vorden som Overgang fra Væren til Intet aldeles ikke kan faae nogen Betydning, naar det ikke indsees, hvorledes det er muligt, at Væren kan blive til Intet, og Intet til Væren. Alle, der have gjort Indsigelse mod denne Overgang, have egentlig gjort Indsigelse mod dens Mulighed.

Stridighederne om den logiske Vorden, den immanente Overgang, Modsigelsens Grundsætning, \textit{principium exclusi medii} o.\ dsl.\ kunne alle henføres under det ene afgjørende Synspunkt, Synspunktet for Muligheden. Det er en gammel Vittighed, at en Dialektiker kan gjøre Sort til Hvidt. Forsaavidt her tænkes paa en virkelig Modsætning, er den, der
% --- p. 236 ---
\opage{236}gjør Sort til Hvidt, aabenbart en Sophist. I Virkeligheden er Sort ikke Hvidt, Hvidt ikke Sort; i Virkeligheden gjælder uden Modsigelse \textit{principium exclusi medii}. Men naar Talen er, ikke om det Virkelige, men om det Mulige, mon den da ogsaa skulde gjælde for en Sophist, der paaviste Muligheden af en Overgang fra Sort til Hvidt? Skulde det være sophistisk at paavise, hvorledes de modsatte Begreber med indre Nødvendighed gaae over i hinanden, saa at det ene, væsentlig seet, er det andets Mulighed: da maatte det ogsaa være sophistisk at paavise Sammenhæng, Grund, Begrundelse o.\ s.\ v.

Hvor inderlig Spørgsmaalet om Muligheden hænger sammen med en forandret Opfattelse af Modsigelsens Grundsætning, sees ved at stille Mulighedsbegrebet mod dets Modsætninger. Det Mulige modsættes det Umulige, det Nødvendige og det Virkelige. At det Mulige er det Nødvendige, er jo en Modsigelse; men dersom det Mulige stilles abstract mod det Umulige, saa bliver det, ifølge Wolfs egen Definition, just til det Nødvendige: nødvendigt er det, hvis Modsatte er modsigende; det Muliges Umulighed er modsigende; altsaa er det Mulige nødvendigt. Men det Mulige kan forandres; det Foranderliges Uforanderlighed er modsigende, følgelig kan det Mulige umulig være det Nødvendige. Stilles det Mulige imod det Virkelige, da er det en Modsigelse, at det Mulige skulde være det Virkelige; men dersom det Mulige ikke kan være virkeligt, hvorledes kan det Virkelige da være muligt? Den i Muligheden skjulte Modsigelse have de Gamle belyst fra mangfoldige Sider; men først Hegel har indseet, at Modsigelsens Løsning netop laae i Overgangen, at med andre Ord Mulighedsbegrebet er et Overgangsbegreb.

Modsætningen mellem Muligt og Umuligt er relativ; kun mellem Nødvendigt og Umuligt er Modsætningen absolut; Modsætningen mellem Muligt og Virkeligt er igjen en anden end mellem Muligt og Nødvendigt. Havde Hegel været opmærk%
% --- p. 237 ---
\opage{237}som paa Modsætningernes forskjellige Styrke og forskjelligartede Charakteer, saa vilde han med sine Mulighedsformer have fundet ganske andre Overgange end dem, der nu findes i hans Logik; men at han har løst Modsigelsen op i Muligheden og hævet Mulighedsformen til immanent Overgangsform, maa paaskjønnes. Vi skulle nu prøve de forskjellige Modsætninger.

Modsætningen mellem det Mulige og det Virkelige er, saavel i den nyere som i den ældre Philosophie, paa mangfoldige Maader forqvaklet. Idet Dogmatikerne uden videre identificere det Mulige med det Tænkelige, bliver det Virkelige blot en Modification af det Mulige: „alt det, i hvilket ingen Modsigelse ligger, er muligt, enten det saa er til eller ikke“. I Muligheden er da paa denne Maade Forskjellen mellem Virkeligt og Uvirkeligt aldeles udslettet: i Muligheden er det Uvirkelige virkeligt, thi det er jo virkelig muligt; og det Virkelige uvirkeligt, thi i Muligheden gjælder det kun for et Muligt. Det er den subjective Modalitet, som her uden videre confunderes med den objective. I subjectiv Modalitet er det Mulige ganske vist det Tænkelige, og det Tænkelige indifferent mod Modsætningen af Virkeligt og Uvirkeligt. At tænke det Virkelige er at forvandle det til et Muligt, det er at opfatte det logisk som et Begreb; det er, med andre Ord, at abstrahere fra dets Virkelighed; thi det Virkelige er i Virkeligheden ikke et Begreb, det er en Existens; det er ikke logisk, men antilogisk: Muligheden er logisk, Virkeligheden antilogisk.

Naar det Antilogiske er sat, saa er det Logiske hævet; naar det Virkelige er sat, saa er det Mulige hævet: det Antilogiske har været logisk, thi Existensen har været Begreb, Realiteten har været Idealitet, Magten har været Viden, det Virkelige har været Muligt. Men det Virkelige er som Virkeligt ikke, hvad det har været, før det blev virkeligt. Den, der siger, at det Virkelige som Virkeligt endnu er muligt, han siger jo, at det Virkelige endnu er, hvad det var, før det
% --- p. 238 ---
\opage{238}blev virkeligt; han taler tvetydigt, og bliver let en Sophist. Tvetydigheden er, at den subjective Modalitet udgives for objectiv, Begrebet for dets Existens. Begrebet er ikke ved en Ændring bleven Existens, det Logiske ikke ved en Ændring bleven antilogisk; det er et Magtens Udslag, et reelt Omslag, hvorpaa Forvandlingen beroer; det er et Magtens Udslag, et reelt Omslag, hvorved det Mulige er blevet et Virkeligt: dette er Charakteren af den objective Modalitet. At reducere det Mulige og det Virkelige til Indifferens i det Tænkelige er en Feil, Hegel med al sin Dialektik har tilfælles med Dogmatikerne. Den Eneste, der ret har været opmærksom paa Modsætningen mellem Muligt og Virkeligt, er Kant; men han har ikke kunnet magte den.

Er Modsætningen af Muligt og Virkeligt opfattet i al Skarphed, saa ligger den Indvending nær, „at eftersom der ikke kan gives nogen jævn Overgang fra Mulighed til Virkelighed, saa maa ovenstaaende Anpriisning af det hegelske Mulighedsbegreb som immanent Overgangsbegreb vel kaldes tilbage“. Dette er dog en Misforstaaelse. Vistnok kan det Immanente i Overgangen ikke opfattes saa eensidigt, at Forskjellen mellem det i Modaliteten Subjective og Objective skulde forsvinde i den rene Objectivitet; men Vexelbestemmelsen af Subjectivt og Objectivt, Logisk og Antilogisk er i Overgangen saa uendelig, Omsætningen saa continueerlig, at det Immanente sees i Forvandling: som Overgangsbegreb er Muligheden væsentlig et Forvandlingsbegreb.

Muligheden er et Videns Ord, Virkeligheden et Magtens Ord; Modsætning af Muligt og Virkeligt er en afgjørende Modsætning af Viden og Magt. Men saa er Overgangen fra Mulighed til Virkelighed jo med det Samme bestemt som en Omsætning af Viden i Magt, en uophørlig Forvandling af Begrebsbestemmelser til Existensbestemmelser. Da nu, som vi fra mange Sider have seet, den samme Totalitet af Bestemmelser kan være baade i Videns og i Magtens Form,
% --- p. 239 ---
\opage{239}saa kan Modsætningen af Muligt og Virkeligt ikke ligefrem henføres til en Mangel i Muligheden, at den f.\ Ex.\ har nogle Betingelser, men mangler andre. Ved eensidig at fremhæve dette Skjelnemærke kunde man let komme til Megarikernes Paastand, at den Mulighed, der mangler Betingelser, er paa Grund af de manglende Betingelser umulig. Nei, skal Modsætningen af Mulighed og Virkelighed selv bære Virkelighedens Præg, da maa den i sig optage en anden Modsætning, en Modsætning nemlig mellem Muligt og Muligt, mellem den ene Mulighed og den anden. Hvor en Virkelighed er mulig, er dens Modsatte ogsaa muligt: Muligheden har alternerende Led.

Betydningen af Mulighedens alternerende Led har allerede viist sig ved Anticipationen; den videre Udvikling er betinget af de anticiperende Functioners vexlende Modalitet. Forholdet mellem Muligheden, Anticipationen og Præformationen er paa anskuelig og populær Maade angivet af Leibnitz. I sin Theodicee er Leibnitz, som bekjendt, gaaet ud fra den Forudsætning, at Gud, for at kunne skabe den bedste Verden, har truffet et Valg mellem uendelig mange mulige Verdener. Med Hensyn hertil har saa Kuno Fischer\footnote{Geschichte der neuern Philosophie. Zweiter Band. Mannheim 1855. S.~509.}% note *)
{} meent at kunne paavise „Modsigelsen“ i det leibnitziske „Verdensbegreb“ paa følgende Maade: „Af de talløse Verdener, der ere mulige, vælger Gud den, han skaber. Den virkelige Verden er desaarsag metaphysisk tilfældig og kun moralsk nødvendig. Men hvad ere hine talløse Verdener, der ere mulige i den guddommelige Forstand? Eet lader sig apodiktisk paastaae, at de alle bestaae af Monader, thi selv den guddommelige Forstand kan ikke forestille sig Tingene anderledes end som Monader. Maae nu ikke alle Monader være besjælede Legemer, forestillende Kræfter, analoge Væsener? Og de analoge Væsener maae, forskjellige
% --- p. 240 ---
\opage{240}som de ere, danne en gradviis Orden, en fuldkommen Trinrække, hvori der ikke er ringeste Afbrydelse. Enhver tænkelig Monade er en tænkelig Differens, og denne maa, ifølge Continuitetsloven, være virkelig. Antage vi nu, at de mulige Verdener ikke bestaae af Monader, saa ere de utænkelige, altsaa i logisk, og følgelig ogsaa i guddommelig Forstand umulige. Antage vi, at de ere Monader, der kun existere i den guddommelige Forstand, men ikke i Naturen, saa \emph{mangle} de aabenbart i Naturen, saa er her en \textit{défaut d'ordre}, saa er det virkelige Universum ingen fuldkommen Trinrække, altsaa ikke den fuldkomneste, ikke den bedste Verden. Vor Slutning lyder: enten ere hine talløse Verdener umulige, eller dersom de ere mulige, saa maae de ogsaa være virkelige. I begge Tilfælde er der udenfor den virkelige Verden ingen anden Verden mulig“.

Imod denne Kritik maa her imidlertid, om ikke paa Monadernes, saa dog paa Mulighedernes Vegne, gjøres Indsigelser. Fra Modalitetens objective Side er det jo vistnok fuldkommen rigtigt, at den Mulighed, som enten fordi den i sit Begreb har noget Modsigende, eller fordi den mangler og, ifølge Virkelighedens Vilkaar, for evigt vil komme til at mangle en Betingelse, er umulig. Men stiller Sagen sig nu ogsaa saaledes, naar vi betragte Muligheden fra Modalitetens subjective Side? Ingenlunde. I subjectiv Forstand er det saa langt fra, at det Modsigende, og forsaavidt i objectiv Forstand Umulige, skulde være et Utænkeligt, altsaa et i subjectiv Forstand Umuligt, at Modsigelsen tvertimod som et skarpt betegnet logisk Begreb lader sig tænke, og det saa tydeligt, at det objectivt Umulige bliver et subjectivt Muligt.

„Hvad betyder Modsigelsen? Et tilværende Intet er en Modsigelse; en rund Fiirkant er en Modsigelse; en flyvende Hest er ogsaa en Modsigelse. Ikke destomindre ere disse Modsigelser af forskjellig Art. En flyvende Hest kan man forestille sig, men ikke tænke. Det Utænkelige ligger ikke
% --- p. 241 ---
\opage{241}ligefrem i det Logiske, at Prædicatet „flyvende“ modsiger Begrebet
„Hest“, thi dette Begreb kunde jo udvides; men deri,
at det hele saaledes udvidede Begreb staaer i Strid med
Erfaringens Betingelser og Love. Anderledes med den runde
Fiirkant; her sees Modsigelsen umiddelbart i Forestillingen
selv. At forestille sig Fiirkanten rund er umuligt: at tænke
sig en rund Fiirkant er derimod muligt, forsaavidt det er
muligt at tænke Prædicatets \emph{Strid} med Subjectet. Modsigelsen
er altsaa ikke blot den, at Prædicatets Strid med
Subjectet gjør Dommen meningsløs; men Modsigelsen er i
høiere Forstand denne, at Modsigelsen selv, dette Utænkelige,
dog lader sig tænke“\footnote{Philos. Propæd. S.~43.}% note *)
. Dersom det Urimelige ikke lod sig
% NOTE: the ø-strokes in "gjør", "meningsløs", "høiere" are lost in this faint scan (witness reads o); ø as elsewhere in the book.
antage, dersom det Utænkelige slet ikke lod sig tænke, saa
gaves der ingen Indsigt, ingen Begrundelse, ingen Begriben.
Hvorpaa beroer vel Indsigten i, hvad der er antageligt?
Aabenbart paa en Indsigt i, hvad der er forkasteligt. Hvorpaa
beroer Erkjendelsen af det Rimelige, det i sig Overeensstemmende?
Naturligviis paa en Erkjendelse af det Urimelige, det
i sig Stridige og Uovereensstemmende. Hvorledes er det
muligt at tænke det Tænkelige uden netop ved Modsætningen
til det Utænkelige? Men saa maa jo for Modsætningens, for
det Tænkeliges Skyld det Utænkelige dog paa en Maade
tænkes. Saa smidig er Tanken, at den med Lethed tænker
sit eget Negative, sin egen Negation.

Ifølge Aristoteles skulde man rigtignok antage, at Tanken
om det Falske, det Urimelige, det Ufuldkomne var ligefrem
udelukket fra den fuldkomne Tænkning; men det er en Ferskhed
i den aristoteliske Dogmatisme, en Ferskhed, hvis Conseqvens
er, at Guddommen ikke engang formaaer at tænke det
Mulige. Kan nu Guddommen ikke tænke det Mulige, da
kan den egentlig heller ikke tænke det Virkelige, og hvad kan
% --- p. 242 ---
\opage{242}den saa tænke?\footnote{Grundideernes Logik I. S.~447.}% note *)
{} Det forholder sig ikke saaledes, at den
tænkende Subjectivitet skulde være uvidende om det Falske, det
Utænkelige i samme Grad som den er vidende om det Sande
og Tænkelige; tvertimod, i samme Grad som Indsigten i det
Sande stiger, i samme Grad stiger ogsaa Indsigten i det
Falske; i samme Grad som Mulighederne, de objective Muligheder,
der ere bestemte til at gaae over til Virkelighed, gjennemskues,
i samme Grad gjennemskues ogsaa de blot subjective
Muligheder, de Muligheder, der ere bestemte til at udelukkes
fra Virkeligheden, efterdi de i objectiv Forstand ere lutter
Umuligheder.

Som et af de meest talende Exempler paa Forskjellen
mellem det Umulige og det Mulige, det blot subjectivt og
det objectivt Mulige, kan Mathematikens imaginære Størrelse
anføres. Hører den imaginære Størrelse til det Mulige eller
til det Umulige? Siden det er en Størrelse, maa den naturligviis
høre til det Mulige, den lader sig jo antage, betegne,
man kan jo regne med den: den kan altsaa ikke ligefrem høre
til det Umulige. Men den imaginære Størrelse staaer dog
i en bestemt Modsætning til den reelle; denne Modsætning
vilde jo tabe al Betydning, dersom der ikke i det Reelle var
Noget muligt, som er umuligt i den imaginære. Opgaven
være f. Ex. at bestemme Skjæringen mellem en ret Linie og
en Cirkel. „Kalde vi den rette Linies Afstand fra Centrum
$p$, Radius $r$, Chorden, der afskjæres, $c$, da have vi paa Spørgsmaalet:
hvor stor er Chorden? som Svar: $c = 2\sqrt{r^2 - p^2}$.
Er nu $p > r$, bliver $c$ imaginær, eller, som Newton udtrykker
det, umulig. Det imaginære $c$ tilkjendegiver altsaa
Umuligheden af, at Linien kan af Cirklen afskjære nogen
Chorde. Men hvorfor er nu dette umuligt? Fordi Linien
ligger heelt udenfor Cirklen. Men saa er her jo dog en
Linie, og vi kunne fra en anden Side see noget Reelt for%
% --- p. 243 ---
\opage{243}bundet med det Imaginære. Ved nemlig at forvandle
$2\sqrt{r^2 - p^2}$, for $p > r$, til $2\sqrt{p^2 - r^2}\,\sqrt{-1}$, seer man,
at den imaginære Størrelse har en reel Bestanddeel, og at
denne Deel netop er det Stykke af Tangenten til den givne
Cirkel, som en den givne Linie tangerende Cirkel afskjærer. ---
For nu ret tydelig at see, hvorledes det Umulige og det
Mulige bestemmes ved hinanden, idet det Imaginære vexler
med det Reelle, ville vi endnu betragte det anførte Exempel
lidt nærmere. Sæt, at Opgaven stilles under følgende Form:
Der er givet to concentriske Cirkler, den ene med Radius $p$,
den anden med Radius $r$; man finde Længden af en Chorde
i den yderste, som tangerer den inderste. I Løsningen af
denne Opgave maae vi da skjelne mellem:
\begin{center}
1)\enspace $p < r$, saa at man faaer\\
$c = 2\sqrt{r^2 - p^2}$;\qquad (I)\\
2)\enspace $p > r$, saa at man faaer\\
$c = 2\sqrt{p^2 - r^2}$.\qquad (II)
\end{center}
\noindent Ved Multiplication med den imaginære Størrelse $\sqrt{-1}$,
kan man da let reducere (I) til (II). Heraf følger da, at
(I) er umulig, naar (II) er mulig, og omvendt (II) er umulig,
naar (I) er mulig. Dette Mathematiske er saa tillige i fuldkommen
Overeensstemmelse med det Logiske, thi $p$ kan dog
umulig, --- i slige Tilfælde veed Dialektiken meget vel, hvad
den skal indrømme den formelle Logik, --- paa een Gang
være baade større (hvorpaa (II) giver Svar) og tillige (hvorpaa
(I) giver Svar) mindre end $r$.“\footnote{Om Begrebet Nøiagtighed med særligt Hensyn til Hr. Prof. Steens
„Bidrag“. Efterskrift til „Philosophie og Mathematik“ af R.
Nielsen. Kjøbenhavn 1860. S.~26--27.}% note *)

Med Læren om hvad vi her ville kalde de imaginære og
reelle Muligheder staaer Læren om den guddommelige Forudvidenhed
i nøieste Sammenhæng. Det er under „det logiske
% --- p. 244 ---
\opage{244}Ideal“\footnote{Grundideernes Logik. I. S.~235.}% note *)
{} allerede anført, i hvilken Betydning de gamle
Theologer tillagde Gud en betinget Viden (\textit{scientia media
[sc. de futuro conditionato s. futuribili, hypothetica,
simplicis intelligentiæ], qua Deus perspicit omnia, quæ,
positis quibusdam conditionibus, evenire potuissent}).
Hertil slutte sig saa, naturligviis paa speculativ Maade, Udtalelser
i den speculative Dogmatik. „Den Modsigelse, man
har villet finde mellem Begrebet af et frit Verdensløb og den
guddommelige Alvidenhed“ --- hedder det i Læren om „det frie
Verdensløb og Guds mangfoldige Viisdom“ --- „beroer paa en
eensidig Opfattelse af Alvidenheden som blot \emph{Forudvidenhed},
paa en Miskjendelse af det \emph{Betingede} i den guddommelige
Raadslutning \dots\ Verdensudviklingens Endemaal med
den hele Række af i og for sig nødvendige Udviklingsmomenter
maa tænkes at staae fast i den evige Raadslutning; men den
virkelige Udførelse af den evige Raadslutning, den hele Fylde
af Verdensløbets virkelige Bestemmelser, forsaavidt som disse
ere betingede ved Skabningens Frihed, kunne ogsaa kun være
Gjenstand for en betinget Forudvidenhed, det er, de kunne
kun forud vides som Muligheder (som \textit{Futuribilia}), efterdi
andre Muligheder kunne blive virkelige“.

Dersom den speculative Dogmatik blot meente Noget med,
hvad den saadan siger, dersom man virkelig kunde formaae
den til at lægge videnskabelig Vægt paa sit frie Tankeløb og
sin egen „mangfoldige Viisdom“, saa vilde det nu være
en let Sag at aftvinge den følgende Tilstaaelse: „Dersom
Verdensmulighederne virkelig afhænge af Betingelsernes Vexel,
og de vexlende Betingelser indskrænke Forudvidenheden, saa
maa Conseqvensen være, at enhver Forudviden af det Virkelige
bliver aldeles umulig. Det vilde paa disse Vilkaar ikke
være en Alvidende muligt at vide nøiagtig idag, hvordan
Veiret skal blive imorgen“. Er Sammenhængen mellem Be%
% --- p. 245 ---
\opage{245}tingelserne derimod at opfatte saaledes, at den, der kjender
Verdenstilstanden i en given Tid, deraf kan forudsige Betingelsernes
Vexel for alle kommende Tider, da har „det Betingede“
aldeles ikke sat Forudvidenheden nogen Grændse;
men Muligheden er forvandlet til en futurisk Nødvendighed,
Verdensløbet er ufrit. Hvad meddeler den speculative Dogmatik
nu herom? Tvetydigheder! Ifølge den speculative Dogmatik
skal „det Betingede i Raadslutningen“ til en vis Grad indskrænke
Forudvidenheden; ifølge den orthodoxe Dogmatik skal
Læren om Verdensbetingelserne tvertimod udvide Forestillingen
om Forudvidenheden. Hvad den orthodoxe Dogmatik lærer
\textit{de Futuribilibus}, bliver jo vistnok ikke udviklet til Begreb,
men Antagelsen af en \textit{scientia, qua Deus perspicit omnia,
quæ positis quibusdam conditionibus evenire potuissent},
indeholder dog betydningsfulde Vink til en logisk Adskillelse
mellem imaginære og reelle Muligheder. Den Mulighed,
som \textit{positis quibusdam conditionibus} er imaginær, kan
igjen \textit{positis quibusdam conditionibus} blive reel, og omvendt.

Da Modsætningen imellem de subjective og objective,
de imaginære og reelle Muligheder ikke er absolut, men relativ,
og da denne Relativitet paa Grund af de vexlende Betingelser
kan gjennemløbe utallige Gradationer: saa indsees, at Vexelbestemmelsen
af Muligt og Umuligt, Muligt og Muligt,
Muligt og Virkeligt maa, i Consequens af de anticiperende
og præformerende Functioners dobbeltsidige Modalitet, være
mangfoldig. Er der i den guddommelige Viden assertoriske
Domme i et uendeligt System, da er der ogsaa et uendeligt
System af problematiske Domme: der er i Mulighedernes
Kreds uendelig meget Mere, end der er i Virkeligheden. At
Leibnitz har ladet Guddommen bestemme den virkelige Verden
ved Udelukkelse af utallige mulige, kan fra Mulighedernes
Standpunkt ingenlunde bebreides ham; hvad der paa Logikens
Vegne kan bebreides ham, er Maaden, den Maade, hvorpaa
% --- p. 246 ---
\opage{246}han lader Anticipationen, Valget og Præformationen af „den
bedste Verden“ foregaae.

At Guddommens Valg blandt utallige Muligheder hos
Leibnitz faaer et stærkt anthropomorphistisk Anstrøg, hidrører
jo vistnok for en Deel derfra, at han i sin populære Fremstilling
har maattet lempe sig efter Læserens Fatteevne, og
saavidt muligt undgaae Conflicter med Theologerne, men
ifølge Grundcharakteren af hans Philosophie, dog væsentlig
derfra, at Forestillingen bestandig brydes med Tænkningen.
Medens Skabelsen er for Tanken et Grundforhold, er den
for Forestillingen en overnaturlig Begivenhed. Det Valg,
hvormed Skaberen træffer den bedste Verden, danner hos
Leibnitz et eget Grændseskjel mellem Mulighed og Nødvendighed,
mellem det abstract Subjective og det selvstændigt Objective.

Forud for Valget hersker Subjectiviteten eensidig; Guddommens
Blik seer ligesom et Hav af bølgende Muligheder,
Alt er i Svæven; thi hver Mulighed er, om Virkeligheden
end er nok saa fuldstændig anticiperet, en svævende Bestemthed;
maaskee saaledes, maaskee saaledes! de problematiske
Functioner ere under Anticipationen i uendelig Bevægelse.

Efter Valget er Alt derimod forandret, den bedste Verden
er nu fuldstændig bestemt; alle dens Muligheder ere med eet
Slag forvandlede til lutter eiendommelige, i deres Eiendommelighed
uforanderlige Anlæg, Anlæg, der udvikle sig med indre
Nødvendighed. Fra alle Sider er nu Alt bestemt, forud beregnet
og fastsat saa punktlig, at hele Udviklingen er en vordende, en
sig stadig udfoldende Nødvendighed: her er ingen Plads for
Muligheden.

Hvad der hos Leibnitz skulde været udviklet til Begreb,
er blevet staaende ved en Forestilling. Kun Forestillingen kan
ville fra Evighed af sætte en endelig Grændse imellem Muligt
og Nødvendigt og saa mene, at Sagen dermed een Gang for
alle er afgjort. Begrebet viser, at Forholdet mellem Muligt
% --- p. 247 ---
\opage{247}og Nødvendigt er et uendeligt Vexelforhold; men da Modsætningen
af det Mulige og det Nødvendige er saa absolut, at
det Mulige ligesaa lidt kan gjøres til et Nødvendigt, som det
Nødvendige til et Muligt, vil Begrebsudviklingen fordre en
nærmere Bestemmelse.

\textit{γγ)} \emph{Den apodiktiske Dom: Functioner af
reflecteret Bestemthed.} Den umiddelbare Bestemthed,
Skepticismens Forudsætning, slaaer ved Reflexionen over i
Ubestemthed; den i Ubestemtheden reflecterede Bestemthed forvandler
sig for Anskuelsen til en umiddelbar bestemt Væsenhed
og danner Forudsætning for Dogmatismen. At den reflecterede
Bestemthed gjælder for umiddelbar, gjør den dogmatiske
Reflexion kjendelig: forskjellige Udgangspunkter betinge
forskjellige Resultater; de modstridende Resultater indføre
paany Ubestemtheden, hver Gang den antages for overvunden.
Som Mulighedsphilosophie vil Dogmatismen udvikle et System
af objective Væsensbestemmelser; men da den med alt
sit Blik for Mulighedens Væsen ikke er istand til at gjennemskue
Mulighedens Form, slaaer den sine Væsensbestemmelser
fast som endelige Værensbestemmelser, og maa standse for den
i Antinomien fremstillede Umulighed. Til Mulighedernes Ontologie
har Dogmatismen imidlertid afgivet charakteristiske Bidrag
ved paa den ene Side at søge deres Oprindelse i den
altbestemmende Subjectivitet, paa den anden Side at fordre
det Muliges objective Sammenhæng med det Virkelige og
Nødvendige. Er her end paa denne Maade fremkommet et
Problem, som den dogmatiske Reflexion ikke formaaer at løse,
saa ere dog de Begrebselementer tilveiebragte, ved hvis Bearbeidelse
Problemets Løsning vil opnaaes.

Det er Formen, Mulighedsformen, hvorpaa det ved Elementernes
Bearbeidelse væsentlig kommer an; men naar Eftertrykket
falder paa Form og Formudvikling, ville Mulighed,
Virkelighed og Nødvendighed reflectere hinanden som forskjellige
Former af een Form, særskilte Bestemmelser i een totaliseret
% --- p. 248 ---
\opage{248}Bestemthed. Da nu Formbestemmelsernes totaliserede Eenhed
selv er Form, saa indtræder det for Functionernes Modalitet
betegnende Formspørgsmaal, om det er Muligheden, Virkeligheden
eller Nødvendigheden, der afgiver den omfattende Formeenhed,
eller om Eenheden selv maaskee er en endnu høiere
og concretere Form: vi skulle da undersøge Formeenheden.

I Formudviklingen kan Muligheden, den endnu ubestemte
Bestemthed, bestemmes som Eenhed af Virkelighed og Nødvendighed.
Det Virkelige er, i væsentlig Modsætning til det
Nødvendige, et Foranderligt; det Nødvendige, i Modsætning
til det Virkelige, et Uforanderligt. Ethvert af de saaledes
modsatte Led har paa een Gang sin Modsætning, sit Negative,
i sig og udenfor sig; det har og er sit eget Modsatte.

Som det Foranderliges Form har den umiddelbare Virkelighed
Nødvendigheden, det Uforanderliges i sig reflecterede
Bestemthed, udenfor sig: det Foranderlige er ikke det Uforanderlige,
det Tilfældige ikke det Nødvendige. Hvad er da
Virkeligheden i sin abstracte Umiddelbarhed, i sin rene
Tilfældighed? Dens positive Væren er at bestemme negativt;
det Foranderlige er det Ikke-Uforanderlige, har i Virkeligheden
denne negative Bestemthed til Grund og Grundlag: det
Foranderlige er kun muligt ved det Uforanderlige.

Som det Uforanderliges Form er Nødvendigheden i sig
negativ bestemt; det Uforanderlige kan kun bestaae som det Foranderliges
Negation; men skal Negationen af det Foranderlige
vedblive, maa det Foranderlige selv vedblive. En Negation,
der Intet har at negere, er umulig. Kan det Foranderlige
ikke længere negeres, fordi det er forsvundet, vil
Negationen øieblikkelig angribe det Uforanderlige. Det Uforanderlige
er kun muligt ved det Foranderlige.

Men at det Foranderlige kun er muligt ved det Uforanderlige,
vil jo med andre Ord sige, at det Uforanderlige er
det Foranderliges Mulighed; at det Uforanderlige kun er
muligt ved det Foranderlige, viser da ligeledes, at det For%
% --- p. 249 ---
\opage{249}anderlige er det Uforanderliges Mulighed. Med den Fordobling,
at det Foranderlige selv indeholder den Uforanderlighed,
af hvis Negation det fremgaaer, er Foranderligheden
udtrykkelig bestemt som Virkelighed; med den Fordobling, at
det Uforanderlige selv indeholder den Foranderlighed, paa hvis
Negation dets Selvbekræftelse beroer, er Uforanderligheden udtrykkelig
bestemt som Nødvendighed. At det Virkelige har
Nødvendigheden til Mulighed, vil sige, at Nødvendigheden
formedelst det Virkelige er Mulighed; at det Nødvendige har
Virkeligheden til Mulighed, vil sige, at Virkeligheden formedelst
det Nødvendige er Mulighed. Paa denne Maade er da Mulighedsformen
alle Formers Formeenhed. \emph{Alt er Mulighed,
Alt er muligt}\footnote{Uagtet det kun er Formbestemmelser, rene Reflexionsbestemmelser,
hvorom her handles, maa man dog ikke bilde sig ind, at de ere
uden Betydning for Virkeligheden. Hvad en forskjellig Opfattelse
af Forholdet mellem Mulighed, Virkelighed og Nødvendighed kan
komme til at betyde for Existensen i dens meest afgjørende Kriser,
har S. Kierkegaard mesterlig udviklet i sit Skrift: „Sygdommen
til Døden“. --- „Den Troende eier den evig sikkre Modgift mod
Fortvivlelse: Mulighed; thi for Gud er Alt muligt i ethvert Øieblik.
Dette er Troens Sundhed, der løser Modsigelser. Modsigelsen
er her, at menneskelig talt er Undergangen vis, og at der
saa dog er Mulighed. Sundhed er overhovedet at kunne løse
Modsigelser. Saaledes legemlig eller physisk: Træk er en Modsigelse,
thi i Træk er Kulde og Varme disparat eller udialektisk; men et
sundt Legeme løser denne Modsigelse og mærker ikke Træk. Saaledes
ogsaa med Troen.~\dots\ --- Fatalisten er fortvivlet, har tabt
Gud og saaledes sit Selv. Men Fatalisten har ingen Gud, eller
hvad der er det Samme, hans Gud er Nødvendighed; som nemlig
for Gud Alt er muligt, saa er Gud det, at Alt er muligt.~\dots\
For at bede maa der være en Gud, et Selv --- og Mulighed, eller
et Selv og Mulighed i prægnant Forstand, thi Gud er det, at Alt
er muligt, eller at Alt er muligt, er Gud; og kun den, hvis Væsen
blev saaledes rystet, at han blev Aand ved at forstaae, at Alt er muligt,
% note breaks here: continues at the foot of p. 250
kun han har indladt sig med Gud. Det at Guds Villie er det
Mulige, gjør at jeg kan bede; er den blot det Nødvendige, er
Mennesket ligesaa umælende som Dyret.“ Anf. Skr. S.~35--36.}% note *)
.

% --- p. 250 ---
\opage{250}Er nu virkelig Alt muligt, saa er det ved Formudviklingens
Modalitet da ogsaa muligt, at det Virkelige kan udgjøre
Mulighedens og Nødvendighedens forenende Midte, at
med andre Ord Virkelighedsformen kan blive Hovedform. I
sin reflecterede Modsætning til Virkeligheden er den objective
Mulighed Eet med Virkeligheden; denne Modsigelse giver sig
tilkjende derved, at det Mulige, uagtet det ikke er virkeligt,
dog, forsaavidt det er objectivt, maa have Tilværelse i et eller
andet Virkeligt, altsaa være som et Virkeligt: kun i det Virkelige
er der en virkelig Mulighed. Kun den virkelige Mulighed
er reel, den uvirkelige er imaginær; den imaginære
Mulighed er i objectiv Forstand Umulighed. Maa nu Muligheden,
for at blive virkelig paa en Maade, allerede forud
være virkelig paa en anden Maade, saa følger jo heraf, at
Muligheden er en Modalbestemmelse, altsaa kun en Sidebestemmelse
ved det Virkelige: Virkeligheden er umiddelbart sig
selv og Muligheden.

I umiddelbar Eenhed med sin Mulighed maa det Virkelige
nødvendig reflectere Nødvendigheden. Nødvendigheden
er, som vi have seet, Virkelighedens Mulighed og kun ved
denne Mulighed væsentlig forskjellig fra Virkeligheden; at
Virkeligheden er i umiddelbar Eenhed med sin Mulighed vil
da sige, at den ved Muligheden integrerer sig selv, at den er
i Eenhed --- formedelst Muligheden i reflecteret Eenhed --- med
Nødvendigheden. Ved Reflexionen taber Virkeligheden jo vistnok
sit eiendommelige Præg, forsaavidt den skal være umiddelbar;
men Tabet gjenoprettes derved, at Reflexionen ophæver sig
selv: Reflexionens Selvophævelse i reflecteret Bestemthed er
Nødvendighed. Ved den Umiddelbarhed, hvormed Reflexionen
begynder og afbrydes, ved den umiddelbare Umiddelbarhed,
% --- p. 251 ---
\opage{251}har Virkeligheden i sig optaget Muligheden; ved den Umiddelbarhed,
hvori Reflexionen fuldendes og derfor selv ender,
ved den reflecterede, af Reflexionen ponerede Umiddelbarhed,
har Virkeligheden i sig optaget Nødvendigheden. Paa denne
Maade er det da Virkelighedsformen, der er Formernes Formeenhed:
\emph{Alt er Virkelighed, Alt er virkeligt.}\footnote{„Det er nemlig ikke saa, som Philosopherne forklare, at Nødvendighed
er Eenhed af Mulighed og Virkelighed, nei, Virkelighed er
Eenhed af Mulighed og Nødvendighed \dots\ Allerede i Forhold til
at see sig selv i et Speil, er det nødvendigt at kjende sig selv; thi
gjør man ikke det, saa seer man ikke sig selv, men blot et Menneske.
Men Mulighedens Speil er ikke noget almindeligt Speil,
det maa bruges med den yderste Forsigtighed. Thi om dette Speil
gjælder det i høieste Forstand, det er usandt. At et Selv seer
saaledes og saaledes ud i Muligheden af sig selv er kun halv
Sandhed; thi i Muligheden af sig selv er Selvet endnu langtfra
eller kun halvt sig selv. Saa gjælder det, hvorledes dette Selvs
Nødvendighed bestemmer det nærmere“. S. Kierkegaard. Sygdommen
til Døden. S.~32--33. Men hvor villig man end fra Existensens
og Virkelighedens eget Synspunkt maa indrømme, at Virkeligheden
er Eenhed af Mulighed og Nødvendighed, saa ligger heri ingen
Hindring for, at Nødvendigheden jo ogsaa, ved en Forandring af
Synspunktet, kan opfattes som Mulighedens og Virkelighedens
Eenhed.}% note *)

Det er Umiddelbarheden, den umiddelbare Bestemthed,
hvormed Reflexionen begynder; det er Umiddelbarheden, den
reflecterede Umiddelbarhed, hvori den ender. Saalænge det
Virkelige ved Reflexion kan omsættes i det Mulige, og det
Mulige igjen ved den tilbagevendende Umiddelbarhed blive det
Virkelige, saalænge, med andre Ord, Reflexion og Umiddelbarhed
vedblive at betegne en Omskiftelse, vil Hovedformen
skifteviis antage snart Mulighedens snart Virkelighedens Præg,
og Nødvendigheden være en Biform. Men naar det Umiddelbare
selv er blevet reflecteret, og det Reflecterede umiddelbart,
ender Omskiftelsen, og Nødvendigheden, det Uomskifte%
% --- p. 252 ---
\opage{252}liges Grundform, lader sig tilsyne. Hvorfor kan den endelige
Modsætning af Umiddelbarhed og Reflexion ikke holde sig?
Fordi ethvert af Modsætningens Led er lige omskifteligt.
Umiddelbarheden, der ligesaavel i det Uendelige som i det
Endelige forudsætter Reflexionen, skifter om og bliver reflecteret;
Reflexionen, der i det Endelige som i det Uendelige forudsætter
Umiddelbarheden, skifter om og gaaer op i det Umiddelbare: det
Reflecterede bliver umiddelbart. End ikke Adskillelsen mellem
umiddelbar og reflecteret Umiddelbarhed, en Adskillelse, hvormed
Reflexionen har anstrengt sig til det Yderste, er istand
til at holde sig. Ved Reflexionens Begyndelse bliver den
første, altsaa umiddelbare Umiddelbarhed, Reflex i den begyndende
Reflexion, altsaa \emph{reflecteret}; med Reflexionens
Fuldendelse forsvinder dens sidste Reflex i den tilblevne, altsaa
tilværende, altsaa ikke længere reflecterede, men \emph{umiddelbare}
Umiddelbarhed. Reflexionen kan nu vende og dreie sig, hvordan
den vil: altid faaer den det Samme ud af det Forskjellige
og det Forskjellige ud af det Samme; altid det
Samme og det Samme, det i alle Omskiftelser Uomskiftelige,
det Uforanderlige, der magter alle Forandringer, det Nødvendige.
Som det Muliges og Virkeliges absolute Eenhed
er Nødvendigheden nu Hovedform: \emph{Alt er Nødvendighed,
Alt er nødvendigt.}

At nu ethvert af Formudviklingens Resultater med hele
sin Alsidighed dog --- paa Grund af den skiftende Modalitet ---
selv er eensidig, falder i Øinene. Var Resultatet: Alt er
Virkelighed, Alt er virkeligt, det sidste og høieste, da vilde,
hvad Empirikerne paastaae, de assertoriske Domme udelukkende
afgive Erkjendelsens tilfældige Grundlag; de logiske Functioner
vilde i utallige Ændringer være Functioner af umiddelbar
Bestemthed. Var Resultatet: Alt er Mulighed, Alt er muligt,
Formudviklingens Høieste, da skulde, hvad Skeptikerne antage,
de problematiske Domme udgjøre Erkjendelsens grundløse
Grundlag, og alle logiske Functioner være Functioner af op%
% --- p. 253 ---
\opage{253}hævet Bestemthed. Var endelig Resultatet: Alt er Nødvendighed,
Alt er nødvendigt, det eneste Fornuftige, da maatte,
hvad Dogmatismen forudsætter, de apodiktiske Domme,
--- Axiomer, Definitioner, o.\ s.\ v. --- afgive Erkjendelsens
væsentlige, faste Grundlag; de logiske Functioner maatte da
paa alle mulige Maader ændres til Functioner af reflecteret
Bestemthed.

De vundne Resultater have paa en vis Maade samme
Gyldighed. Den apodiktiske Dom vil ligesaa lidt være istand
til at fortrænge eller beherske enten den assertoriske
eller den problematiske, som nogen af disse nogensinde
vil formaae at gjøre sig til Videns eneherskende. Det
Assertoriske, det Problematiske og det Apodiktiske ere i Modaliteten
tre Former, hvorunder Dommens Gyldighed oprindelig
fremtræder; men i Ideernes Logik --- det ligger i
Logikens Væsen --- er det en Selvfølge, at den apodiktiske
Form maa være den overgribende. Til nærmere Belysning
udhæves da de forskjellige Phaser for det Apodiktiske.

Forsaavidt Nødvendighedsformen fremkommer ved, at Reflexionen
ophæver og standser sig selv, falder det Apodiktiske
sammen med det Logiske: det Apodiktiske er \emph{logisk}. Forsaavidt
Reflexionen kun ved Hjælp af et Umiddelbart kan naae til
Umiddelbarhed, og det Umiddelbare, i hvis Reflex Reflexionen
standser, som Realnegation er Reflexionen modsat, maa det
Apodiktiske falde sammen med det Antilogiske: det Apodiktiske
er \emph{antilogisk}. Skal nu det logisk Apodiktiske ikke geraade i
Strid med det antilogiske, maae begge Former blive Modalbestemmelser
af een absolut overgribende Form, det Rationelle:
det Apodiktiske er baade som logisk og antilogisk \emph{rationelt}.
Altsaa:

1) \emph{Det Logisk-Apodiktiske: væsentlig Nødvendighed.}
I den hegelske Logik afhandles Nødvendigheden paa
tre forskjellige Steder: i „den objective Logik“ under „Virkeligheden“
i „den subjective Logik“ under „Dommen“ (das Ur%
% --- p. 254 ---
\opage{254}theil der Nothwendigkeit) og under „Slutningen“ (der Schluß der
Nothwendigkeit); imod en saadan Fordeling kan der i systematisk
Henseende vistnok Intet indvendes. Charakteristisk er
det derimod, at medens Udviklingen i Logikens objective Deel
er udførlig, skarpsindig og frugtbar, er Behandlingen i den
subjective Deel derimod kort, aphoristisk og gold. Grunden
hertil er let at indsee; det er Miskjendelsen af Subjectiviteten,
som ogsaa i dette Punkt hævner sig.

Naar Nødvendigheden belyses ved Dommen, skulde man
dog vente at see dens Betydning fra den subjective Side, og
dette saameget mere, som dens objective Betydning allerede
forud har gjennemgaaet sin logiske Udvikling. Men nei!
Dommen er jo objectiv, i Subjectivitetens Begreb, saa objectiv,
at den objectivt dømmer sig selv; Nødvendigheden maa
da anden Gang præsentere sig objectivt. Nødvendigheden skulde
nu som Begrebets Nødvendighed erholde en trinviis Udvikling
gjennem den kategoriske, den hypothetiske og den disjunctive
Dom\footnote{„\emph{Dieß ist die Nothwendigkeit des Begriffs}“, siger Hegel
angaaende den disjunctive Dom --- „worin erstens die Dieselbigkeit
beider Extreme einerlei Umfang, Inhalt und Allgemeinheit ist;
zweitens sind sie nach der Form der Begriffsbestimmungen unterschieden,
so daß aber um jener Identität willen diese als bloße
Form ist. Drittens erscheint die identische objektive Allgemeinheit
deswegen als das in sich Reflektirte gegen die unwesentliche Form,
als \emph{Inhalt}, der aber an ihm selbst die Bestimmtheit der Form
hat; das eine Mal als die einfache Bestimmtheit der \emph{Gattung};
das andere Mal eben diese Bestimmtheit als in ihren Unterschied
entwickelt, --- auf welche Weise sie die Besonderheit der \emph{Arten},
und deren Totalität, die Allgemeinheit der Gattung, ist“. ---
At disse Bemærkninger blot anføres i Anledninng af Dommen % PRINTER: printed "Anledninng", sense "Anledning"
og ingenlunde ere Begrebsbestemmelser, der med indre Nødvendighed
udvikles af Dommens eiendommelige Væsen, er iøinefaldende.}% note *)
); men de Bemærkninger, som i Anledning af
Dommens Form anføres, ere saa langt fra at føre Udvik%
% --- p. 255 ---
\opage{255}lingen af Nødvendighedsbegrebet videre, at de tvertimod tjene
til at aflede Opmærksomheden fra Hovedsagen. Kun ved
Modaliteten kan Bestemmelsen af Nødvendigheden vise sig at
staae i væsentlig Afhængighed af Dommen. Afhængigheden
kjendes paa Indsigten: i Indsigten mødes det Subjective med
det Objective.

Det Objective, der svarer til Indsigten, er det Indlysende;
det Indlysende fremstiller sig aldrig i et blot umiddelbart Lys:
Forstandens Lys er altid reflecteret. Og dog er i det reflecterede
Lys Reflexionen ophævet, Indsigten er i sin Reflecteren
umiddelbar, i sin Umiddelbarhed reflecteret. Til Indsigtens
reflecterede Umiddelbarhed svarer det umiddelbart Indlysende;
til det umiddelbart Indlysende det umiddelbart Apodiktiske.
Det umiddelbart Apodiktiske har ved første Øiekast en paafaldende
Lighed med det Assertoriske. „Vandet er flydende“; „den
rette Linie er den korteste Vei imellem to Punkter“; med Hensyn
paa det Umiddelbare i Forvisningen synes begge Domme
at høre under samme Kategorie, og dog er her en Sphæreforskjel:
den første er assertorisk, den sidste er apodiktisk. At
den sidste Dom væsentlig er apodiktisk, sees deraf, at det er
Reflexionen, ikke Sandsningen, der giver os Visheden. Vi
drage i Tankerne utallige Linier, vi sammenligne dem, og
indsee nu, uden at see de utallige Linier umiddelbart,
at den rette maa være den korteste. Forsaavidt Indsigten
beroer paa en Sammenlignen, er Lyset i det Indlysende reflecteret.
Men en egentlig Beviisførelse finder alligevel ikke Sted;
hele Sammenligningen beroer paa et Skjøn: Sætningen er
som Grundsætning umiddelbart indlysende.

At det ved sig selv Indlysende væsentlig beroer paa Reflexion,
at ingen Grundsætning staaer saa fast, at den jo maa
falde for et dialektisk Angreb, er paa Dogmatismens Standpunkt
en logisk Hemmelighed. Og dog kan man gjøre Forsøget
med hvilkensomhelst Grundsætning. „To Størrelser,
som hver for sig ere ligestore med een og samme tredie, ere
% SEAM: p. 255 ends with a whole word ("ere"); the quotation „To Størrelser ... continues on p. 256.
% --- p. 256 ---
\opage{256}indbyrdes ligestore“; denne Sætning vakler øieblikkelig, naar den ret bliver gjennemreflecteret. Er her da virkelig nogen Modsigelse? „For at indsee, hvad det egentlig er, her tales om, maa man udtrykkelig mærke sig, at her ikke tales om Ting, der have Størrelse, men om Størrelser. Tales her, ikke om Størrelserne selv, men om Ting, der have Størrelse, f.\ Ex.\ et Æble, en Trekant, en Time o.\ s.\ v., da er her ganske vist ikke Spor af Modsigelse. At to Trekanter, som hver for sig ere ligestore med een og samme tredie, ere indbyrdes ligestore, forstaaes af sig selv. Saalænge vi have med concrete Størrelser at gjøre, kan Forskjellen og Fleerheden meget vel bestaae med Ligestorheden; men naar der handles om rene Størrelser, stiller Sagen sig anderledes. En Fleerhed af Størrelser forudsætter med Nødvendighed, at den ene maa kunne adskilles fra den anden: uden indbyrdes Forskjel ingen Fleerhed. Men hvorved skulle rene Størrelser vel adskille sig fra hverandre indbyrdes uden ved Storheden? Ere Størrelserne indbyrdes ligestore, saa er jo ethvert Glimt af Forskjel forsvunden, dermed er tillige Fleerheden forsvunden: at flere Størrelser ere indbyrdes lige, vil altsaa med andre Ord sige, at det ikke er flere, men een og samme Størrelse, som, hvormange Gange den end gjentages, dog ikke mangfoldiggjøres, men vedbliver at være den selvsamme“.\footnote{Forelæsninger over „Phil. Propæd.“ Universitetsaaret 1860--61. S.~2--3.}% note *)

Vil Nogen mene, at logiske Undersøgelser af denne Art ere aldeles ufrugtbare, da behøve vi blot at minde om, hvor inderlig Spørgsmaalet om de rene Størrelsers Forskjellighed hænger sammen med Sørgsmaalet
% PRINTER: printed "Sørgsmaalet", sense "Spørgsmaalet"
om \textit{principium indiscernibilium}, et Spørgsmaal, der er afgjørende ved Bedømmelsen af den leibnitziske Philosophie. Vil Nogen mene, at den Dialektik, hvormed dogmatiske Grundsætninger rokkes, absolut maa være sophistisk, da maa herved mærkes, at det ikke er
% --- p. 257 ---
\opage{257}det negative Resultat, men dets Udtydning og Anvendelse, hvorpaa Sophistiken skal kjendes. Naar en Grundsætning er rokket, udtyder Sophisten det saaledes, at det er den apodiktiske Dom, der er kuldkastet, og benytter det negative Resultat til et Tankespil. Den sande Dialektik forstaaer Sagen anderledes. At Grundsætningen kan rokkes, betyder, at det Apodiktiske hviler paa Reflexionen; Sætningens Vaklen er Udtrykket for, at Reflexionen er sat i Bevægelse. Bevægelsen betinges af Abstractionen. Idet den abstracte Sætning opstilles som Grundsætning, maa der altid en eller anden Bibestemmelse til for at holde den i Umiddelbarheden. Saaledes holder Mathematiken de forskjellige ligestore Størrelser ud fra hverandre ved Betegnelser, Biforestillinger om Summer, Produkter (3 + 3 = 3.\,2 = 12\,:\,2), Potenser o.\ s.\ v. Idet nu Dialektikeren lader Abstractionen borttage de paagjældende Biforestillinger og sætter Tanken i Uro, er han sig tilfulde bevidst, hvad han gjør; thi han veed med sig selv, at hele Operationen er meningsløs,
dersom han ikke tilbagegiver Sætningen den fradragne Umiddelbarhed og saaledes igjen bringer Reflexionen til Ro. Paa denne Maade bliver det indlysende, at Sætningen er en apodiktisk Dom og som apodiktisk Dom en sand Grundsætning.

Den sunde Menneskeforstand stoler, og det med Rette, paa det umiddelbart Indlysende; men naar der skal gjøres Adskillelse mellem den assertoriske og den apodiktiske Dom, er Opgaven ikke saa let. At prøve det Apodiktiske ved Hjælp af Modsigelsen faaer kun Betydning, naar man først tilfulde har prøvet Modsigelsen. Hvor vanskeligt det undertiden kan være at paavise og prøve en Modsigelse, sees af mangfoldige Exempler, og minder os, hvad Anstrengelsen med Problemerne hos de Gamle angaaer, i denne Forbindelse særlig om Løgneren (\textit{ψευδόμενος}). „Naar En siger, at han lyver, lyver han saa med dette sit Udsagn, eller siger han Sandhed? Her fordres til Svar enten Ja eller Nei. „Ja, han lyver!“
% --- p. 258 ---
\opage{258}--- Og siger det selv, nævner Tingen med det rette Navn, siger, hvad sandt er, lyver altsaa ikke. „Nei, han lyver ikke, han taler Sandhed!“ --- Man maa altsaa troe, hvad han siger; nu siger han jo netop, at han lyver, saa taler han jo ikke Sandhed“. Enten her svares Ja eller Nei, saa svares her i en apodiktisk Dom. Reflexionsbevægelsen fremkommer derved, at Begreberne Løgn og Sandhed uophørlig ombyttes med hinanden; det er to apodiktiske Domme i indbyrdes Strid, en Antinomie i sophistisk Form.

At her er en Sophisme, er umiddelbart indlysende. Men naar dette umiddelbart Indlysende nu tillige skal give en tilsvarende Indsigt, maa det nøiagtig vises, hvori det Sophistiske bestaaer; kun saaledes kan Sophismen udtrykkelig hæves. Det er to Domme, som her ere slaaede sammen i Eet; en Dom om en Dom er smeltet sammen med den bedømte Dom. Den, der siger: jeg lyver, han dømmer om et Udsagn, han enten har givet eller vil give, altsaa forsaavidt om en Dom. Kalde vi Dommen i Udsagnet x, Dommen om Udsagnet a, saa indsees, at den Talende maa apodiktisk være sanddru i den ene Dom, naar han apodiktisk er Løgner i den anden. Er Dommen x en Sandhed, saa stikker Løgnen i Dommen a; er Dommen a en Sandhed, saa have vi Løgneren i Dommen x. Den, der selv siger, at han lyver, maa absolut enten have sagt Noget, eller i sit Udsagn henvise til Noget, hvori han ikke lyver; han bliver da Løgner ved at være sanddru, og sanddru ved at være Løgner. --- Istedenfor at afvise Sophismen med Foragt, maae vi snarere beundre den Snildhed, hvormed Knuden er knyttet, og opfatte det, ikke som et Tegn paa Aandsforvirring, men som et talende Beviis paa den græske Tænknings logiske Energie, naar vi høre, at Chrysippos skrev sex Bøger om denne Tankegaade, og at Philetas fra Chos paadrog sig en Svindsot ved overdrevne Anstrengelser for at løse den. Det umiddelbart
% --- p. 259 ---
\opage{259}Indlysende maa hæves til Indsigt; det Apodiktiske maa sees at være apodiktisk.

Den assertoriske Dom har den problematiske udenfor sig og hjemfalder just derfor saa let til det Problematiske; den apodiktiske Dom derimod har ved Reflexionen det Problematiske i sig selv, har overvundet det og magter det. At det Apodiktiske er indlysende ved sig selv, viser, at det Problematiske er overvundet; at det i sig selv Indlysende hviler i Reflexion, sees deraf, at det Problematiske, for at være overvundet, maa være forudsat: paa Gjennemsigtigheden af det i det Apodiktiske ophævede Problematiske beroer Indsigten.

Det er da heller ikke blot Dialektiken, der for Begrundelsens Skyld rokker ved Grundsætningerne og giver det Umiddelbare til Priis for Reflexionen; en exact Videnskab som Mathematiken gjør det Samme. Det kan for Mathematiken i flere Tilfælde være tvivlsomt, om en Sætning skal opstilles som Grundsætning, betragtes som Definition eller bevises som Læresætning. Læresætningen og Grundsætningen ere begge apodiktiske Domme; Forskjellen er, at hiin er middelbart, denne umiddelbart indlysende. Men denne Forskjel kan Tænkningen ophæve; thi Umiddelbarheden i det umiddelbart Indlysende er reflecteret, og Reflexionen i det middelbart Indlysende ved Lysningen selv umiddelbar.\footnote{I „Elementær Geometri“ ved C.\ Ramus, Kjøbenhavn 1850, er Sætningen om den rette Linie mellem to Punkter opstillet som \emph{Grundsætning}. „En Grundsætning (\textit{axioma})“ --- hedder det S.~7 --- „fremstiller en Sandhed som umiddelbart given, idet dens Nødvendighed er i sig selv indlysende (f.\ Ex.\ kun een Linie kan være den korteste mellem to Punkter)“. Herved henvises vi til et Foregaaende (S.~5), hvor det allerede er sagt, at den Linie, som angiver den korteste Afstand mellem to Punkter, er den rette Linie. I „Elementær Plangeometrie“ af L.\ Oppermann, Kjøbenhavn 1845, er Definitionen paa den rette Linie derimod angiven for sig (S.~4),
% note breaks here: continues at the foot of p. 260
og Sætningen: „Den korteste Linie mellem to Punkter er ret“, opstilles saa senere hen (S.~43) udtrykkelig som \emph{Læresætning} med udførligt Beviis.}% note *)

% --- p. 260 ---
\opage{260}I den apodiktiske Dom har Reflexionen udlignet Forskjellen mellem subjectiv Vished og objectiv Væsenhed. Den assertoriske Dom er umiddelbart subjectiv, fordi den er umiddelbart objectiv; den apodiktiske er reflecteret objectiv, fordi den er reflecteret subjectiv. I hvilken Betydning det logisk Apodiktiske beroer paa den ved Reflexion fastsatte Bestemthed, sees af Domme, hvori den Vidende besinder sig paa sin Viden. Paa denne Maade kan en hvilkensomselst Dom bestemmes ved en apodiktisk Dom. Tavlen er \emph{virkelig} sort, altsaa: Dommen om en Tavles Farve er \emph{nødvendigviis} assertorisk. Den chemiske Tiltrækning er \emph{maaskee} elektrisk, altsaa: Dommen om Identiteten af elektrisk og chemisk Tiltrækning er --- paa Videnskabens nuværende Trin --- \emph{nødvendigviis} problematisk. Videns objective Væsenhed bestemmes apodiktisk ved Viden selv: der gives i formel Reflexion en apodiktisk Viden af det Assertoriske og det Problematiske ligesaa vel som af det Apodiktiske. Men Et er Objectiviteten i Videns Form, et Andet det Objective i Videns Indhold. Hvad Videns Indhold angaaer, er det kun Ideen, det Absolute selv, der kan give sand apodiktisk Vished.

Naar det synes, som om der i det Endeliges Relationer var Meget, hvorom der kunde haves apodiktisk Vished uden ligefremt Hensyn til Ideen, da maa herved mærkes, at det Apodiktiske i slige Tilfælde er betinget, og det Betingede i alle Reflexioner en Reflex af sit Ubetingede. „Den kritiske Adskillelse af den rene Anskuelse uden Begreb og det rene Begreb uden Anskuelse er i sit Væsen usand; Begreb og Anskuelse ere Momenter for hinanden gjensidig, deres dialektiske Eenhed er deres Sandhed, de ere i Ideen Eet. Den kritiske Adskillelse af Qvantitetens, Qvalitetens, Relationens og Modalitetens
% --- p. 261 ---
\opage{261}Kategorier har i det Endelige sin Betydning, men er, naar det gjælder Væsenet selv, ideeløs og usand. Til Kritikens endelige Adskillelse svarer en ligesaa endelig Forbindelse; men Ideen er ligesaa lidt en Sum af endelig forbundne Kategorier, som en for sig, d.\ e.\ ligeoverfor de øvrige værende Fornuftkategorie. Ideen er som Idee Kategoriernes indre principielle Totalitet, Sandheden i de kategoriske Bestemmelser. Den kritiske Adskillelse af en Ting i sig uden Egenskaber og et Phænomen af Egenskaber uden en Ting i sig er ideeløs og usand; Sandheden er, at Tingen bestemmes ved sine Egenskaber og Egenskaberne ved Tingen: den uendelige Vexelbestemmelse grunder sig paa Ideen. Den kritiske Adskillelse af et Subject for sig uden Object og et Object for sig uden Subject er fra Væsenssiden den høieste Ideeløshed. At negte Identiteten af det Subjective og det Objective er at fornegte Ideen og gjøre al Sandhedserkjendelse umulig.“\footnote{Philos. Propæd. S.~181--82.}% note *)

Hvad den synlige Verden angaaer, er der en væsentlig Forskjel mellem at see Tingene i Lyset og at see Lyset som Lys; i det Intellectuelle er der en endnu større Forskjel mellem at see Sandheden stykkeviis i de sande Sætninger og at see de sande Sætninger som Momenter i Sandheden selv. Sandheden i enhver sand Sætning, i enhver apodiktisk Dom, er et Glimt af Ideens Lys; ligesom ingen Gjenstand kan sees i det rene Mørke, saaledes kan heller ingen Sandhed erkjendes i den rene Ideeløshed: overalt i det Indlysende, være sig det umiddelbart eller middelbart Indlysende, er det en Lysning af Ideen, der sees i Indsigten. Men adspredt i de mange Lysninger, fordeelt paa de mange apodiktiske Domme i deres deels ubetingede Umiddelbarhed deels endelige Betingethed, kunne Straalerne af Ideens Lys endnu ikke samle sig i det Brændpunkt, hvori Ideen belyser sig selv; det relativt Apodiktiske mangler endnu sit Absolute.

% --- p. 262 ---
\opage{262}I denne Sammenhæng kan Forholdet mellem mathematiske og logiske Sandheder, mellem Mathematikens og Logikens apodiktiske Domme nærmere belyses. Enkeltviis og i deres specielle Forbindelser have de mathematiske Sætninger et langt stærkere Præg af apodiktisk Gyldighed end de logiske. De logiske Sætninger befinde sig under den strømmende Bevægelse i bølgende Overgange. Denne Bølgen, denne dialektiske Svingen og Svæven giver dem for den overfladiske Betragter noget Vaklende, noget Slibrigt og Utilforladeligt, der ikke synes at kunne bestaae med Fordringen paa apodiktisk Vished. Anderledes i Mathematiken. De mathematiske Sætninger bære jo vistnok ogsaa Relativitetens Mærke, forsaavidt de i indbyrdes Sammenhæng afhænge, den ene af den anden; men trods denne Afhængighed og gjensidige Betingethed er dog hver Sætning for sig saa exact bestemt, og giver enkeltviis en saa uimodsigelig Vished, at det Apodiktiske derved erholder en punktlig Allestedsnærværelse: det kan bogstavelig paapeges paa ethvert Punkt. Og dog mærker man i Mathematiken Intet til Ideen som Idee; med hvad Ret kan det da siges, at ethvert apodiktisk Glimt er en Fulguration af Ideen?

I Systemet af Mathematikens apodiktiske Domme er Videns Idee nærværende fra Først til Sidst; men at Ideen trods sin apodiktiske Nærværelse ikke lader sig mærke som Idee, har apodiktisk sin Grund i Punktualitetens Eiendommelighed. Den Punktlighed, hvormed hver enkelt Hovedsætning fremhæves og fastholdes for sig, giver Endeligheden, Middelbarheden og Betingetheden en saa afgjørende Overvægt, at den indre Uendelighed maa holde sig skjult. Ja, saa fremmed kan Uendelighedens Væsen være for en Mathematikers Bevidsthed, at han endog opfatter Uendelighedsbegrebet som en Art „\emph{Blandingsbegreb}“ og paastaaer, at der i Mathematiken ikke gives noget Uendelighedsproblem.\footnote{At tænkende Mathematikere, hvormegen Vægt de end iøvrigt lægge paa Operationsmethodens tilfredsstillende Resultater, desuagtet aldrig
% note breaks here: continues at the foot of p. 263
have kunnet faae det i deres Tanker ubegribelige Begreb (Uendelighedsbegrebet) paa en tilfredsstillende Maade adskilt fra den paa dette Begreb hvilende Methode, viser Mathematikens Historie lige fra Leibnitz
% UNSURE: Leibnitz or Leibniz (degraded print; the tz ligature is clear only at p. 256 "leibnitziske")
og Newton (Jvfr.\ Mathem.\ og Dialekt.\ S.~90--91, S.~95 o.\ flg.) indtil den nærværende Tid. „For det Første kunne vi, naar Principet er ubekjendt, jo ikke vide, om ogsaa alle deraf følgende Consequenser lade sig indrømme: vi gaae som i Blinde, og det turde hænde sig, at vi efter møisommelig Anstrengelse tilsidst stødte paa Noget, som netop udviste Principets Natur, d.\ v.\ s.\ paa Noget, vi ikke kunde begribe. For det Andet: er det ikke en skammelig Nedværdigelse af Fornuften, at lade den følge en Sammenhæng af Led, en Række, hvis Begyndelse og Ende den ikke kan gjennemskue?“ Disse Savriens Betænkeligheder forklare tilstrækkelig den Uro og Spænding, hvori det Uendeliges Analyse bestandig har holdt den mathematiske Reflexion (Mathem.\ Dial.\ S.~55--56), hvis Problem --- saaledes som en af Berlinerakademiet i Aaret 1784 udsat Priisopgave charakteristisk nok udtrykker det, „wie aus einer widersprechenden Voraussetzung so viele wahren Sätze haben können hergeleitet werden“ --- ikke vil ophøre at vende tilbage. Forelæsninger over „Phil. Propæd.“ fra Universitetsaaret 1860--61. S.~47--48.}% note *)
{} I en Ideernes Logik kunne de
% --- p. 263 ---
\opage{263}enkelte Sætninger derimod ikke særlig fremtræde med apodiktisk Selvstændighed, da hver relativ Bestemmelse er et Moment, et Glimt af det Absolute. Men just denne de apodiktiske Momenters reflecterede og i Reflexionen gjennemsigtige Form gjør Ideen, det absolut Apodiktiske, synlig.

Modalitetens almindelige Hovedformer: Mulighed, Virkelighed og Nødvendighed, anbefale sig vexelviis til den sunde Sands saaledes, at snart Endeligheden, snart Uendeligheden, snart Virkelighedssiden, snart Ideesiden vender imod Betragteren. Paa Endelighedens Trin gjælder det, at der nødvendigviis maa være Noget, der blot er \emph{muligt}, Noget, der er \emph{virkeligt}, og Noget, der er \emph{nødvendigt}. Seet under Uendelighedens Synspunkt træder Forskjellen mellem Muligt, Virkeligt og Nødvendigt saaledes i Baggrunden, at den ene
% --- p. 264 ---
\opage{264}af Formerne uvilkaarlig faaer Magt med de to andre, og nu hedder det: snart, at Alt er muligt, snart, at Alt er virkeligt, snart, at Alt er nødvendigt. Falder Eftertrykket paa Vorden (Heraklit), saa er Alt muligt; falder Eftertrykket paa Væren og det Værende, saa er Alt virkeligt (Megarikerne); falder Eftertrykket paa Virkelighedens i sig reflecterede Væsen, den uendelige Substans, saa er Alt nødvendigt (Spinoza). I logisk Forstand er Nødvendigheden overgribende, men den overgribende Nødvendighed er væsentlig subjectiv.

2) \emph{Det Antilogisk-Apodiktiske: virkelig Nødvendighed.} Til nærmere Bestemmelse af den væsentlige eller subjective Nødvendighed er det nødvendigt at bestemme den objective eller virkelige; er den subjective Nødvendighed logisk, saa maa den objective være antilogisk. Men i Indsigten maa Nødvendigheden kjendes paa det Apodiktiske; gives der en antilogisk Nødvendighed, saa maa der gives en Vished, som paa een Gang baade er apodiktisk og antilogisk; der maa, med andre Ord, gives en \emph{Indsigt i det Antilogiske}. Al Indsigt er væsentlig logisk; en Indsigt i det Antilogiske er da en haard Modsigelse: et antilogisk Apodiktisk maa ved første Øiekast forekomme Betragteren som en \textit{contradictio in adjecto}.

Den objective, antilogiske Nødvendighed er haard, den haarde Nødvendighed er umiddelbar; men selv i sin haardeste Umiddelbarhed er Nødvendigheden dog væsentlig reflecteret og staaer ved Reflexionen aaben for Indsigten. Det er just den objective Nødvendighed, der med al sin Haardhed gaaer mæglende ind imellem det Logiske og det Antilogiske, og gjør en Indsigt i det Antilogiske mulig. Nødvendighedens logiske Væsenhed reflecterer sig i det Almene; dens haarde Umiddelbarhed sees i det Enkelte. Fæstes Blikket udelukkende paa det Enkelte, saa forsvinder Nødvendigheden; den bliver borte i en Tilfældighed. Det Tilfældige er, fordi det er, har sin Grund i sig selv, eller rettere: det har ingen Grund, fordi det i al Umiddelbarhed er en tilværende Grund og følgelig grundløst;
% --- p. 265 ---
\opage{265}heri ligner Tilfældigheden den haarde, i sin Haardhed blinde Nødvendighed. Det umiddelbart Nødvendige er et Tilfældigt, det Tilfældige er et umiddelbart Nødvendigt: her sees det Antilogiske. Men just, idet den haarde Umiddelbarhed antilogisk støder Reflexionen fra sig, vil denne Tilbagestøden væsentlig være en Tiltrækning. Subjectiviteten, der ikke kan gjennemreflectere det Tilfældige, kan dog reflectere over det, altsaa gjøre det til Gjenstand for en Reflexion.

Som Gjenstand for en Reflexion viser det Tilfældige sig let i sin Forskjel fra og Eenhed med det objectivt Nødvendige. At det Tilfældige i sig selv er grundløst, vil da sige, at det har sin Grund udenfor sig selv; at det er, fordi det er, vil sige, at det er, fordi et Andet er;
% UNSURE: "er;" or "er:" (degraded print; text layer reads a comma)
at det er, fordi et Andet er, vil sige, at det som udvortes Enkelt staaer imod et andet udvortes Enkelt. Anderledes med det Nødvendige. At det Nødvendige er, fordi det er, betyder, at det umulig kan have sin Grund udenfor sig, eftersom det har sine Betingelser og dermed sin tilstrækkelige Grund i sig selv alene. Forsaavidt det Nødvendige fremtræder som det enkelte Tilfælde, reflecterer det sig i den almene Nødvendighed; det har som et Nødvendigt i denne Reflexion det Almene, ikke udenfor sig, men i sig: det de nødvendige Tilfælde iboende Almene er Magten. Alle Udslag af den umiddelbare Magt ere i deres umiddelbare Enkelthed lutter Tilfælde, i deres udvortes Gjensidighed tilfældige; men Magten selv, det i det Enkelte reflecterede og i Reflexionen dog umiddelbart Almene, er Nødvendigheden, den \emph{virkelige} Nødvendighed.

Hvorledes den virkelige, i Magtens Form antilogiske Nødvendighed forholder sig til Indsigten, kan sees af den apodiktiske Dom. Paa det antilogiske Grundlag har den apodiktiske Dom en saa paafaldende Lighed med den assertoriske, at det i de enkelte Tilfælde endog er vanskeligt at skille dem ad. „Dette Legeme er virkelig tungt“; „Legemet er nødvendigviis tungt“; i første Tilfælde fæstes Opmærksomheden paa Legemet
% --- p. 266 ---
\opage{266}i dets umiddelbare Enkelthed, og Dommen maa blive assertorisk; i sidste Tilfælde fæstes Opmærksomheden paa det i Legemet legemlige Væsen, og Dommen maa blive apodiktisk, nemlig antilogisk-apodiktisk. Den assertoriske Dom bestemmes ved Gjenstanden, den antilogisk-apodiktiske ved Magten. Gjenstanden i dens Eiendommelighed er et Magtens Ord; men Gjenstandens objective Mening, den mægtige, almene Tanke, gjør Magtordet til et Videns Ord.

Det store kantiske Spørgsmaal: „hvorledes ere synthetiske Domme \textit{a priori} mulige?“ kan i denne Sammenhæng paralleliseres med det Spørgsmaal: hvorledes ere antilogisk-apodiktiske Domme mulige? At de synthetiske Domme have en fremherskende empirisk, altsaa antilogisk og forsaavidt assertorisk Charakteer, er forudsat i det kantiske Spørgsmaal; dersom det Synthetiske ikke nærmest fremkom \textit{a posteriori}, kunde det jo ikke være et Problem at finde det \textit{a priori}. At de analytiske og forsaavidt apodiktiske Domme ere af apriorisk Charakteer, er ligeledes forudsat. Kant spørger just efter det Synthetiske, ikke efter det Analytiske, fordi de analytiske Domme, hvad Hume allerede har fremhævet, paa Grund af Identitetens formelle Væsen, ikke berige Erkjendelsen.\footnote{Analytische Urtheile (die bejahenden) sind also diejenigen, in welchen die Verknüpfung des Prädicats mit dem Subject durch Identität, diejenigen aber, in denen diese Verknüpfung ohne Identität gedacht wird, sollen synthetische Urtheile heißen \dots\ \emph{Erfahrungsurtheile}, als solche, sind insgesammt synthetisch. Denn es wäre ungereimt, ein analytisches Urtheil auf Erfahrung zu gründen, weil ich aus meinem Begriffe gar nicht hinausgehen darf, um das Urtheil abzufassen, und also kein Zeugniß der Erfahrung dazu nöthig habe. Daß ein Körper ausgedehnt sei, ist ein Satz, der \textit{a priori} feststeht, und kein Erfahrungsurtheil. Denn, ehe ich zur Erfahrung gehe, habe ich alle Bedingungen zu meinem Urtheile schon in dem Begriffe, aus welchem ich das Prädicat nach dem Satze des Widerspruchs nur herausziehen und dadurch zugleich der Nothwendigkeit des Urtheils bewußt
% note breaks here: continues at the foot of p. 267
werden kann, welche mich Erfahrung nicht einmal lehren würde. Dagegen ob ich schon in dem Begriff eines Körpers überhaupt das Prädicat der Schwere gar nicht einschließe, so bezeichnet jener doch einen Gegenstand der Erfahrung durch einen Theil derselben, zu welchem ich also noch andere Theile eben derselben Erfahrung, als zu dem ersteren gehörten, hinzufügen kann. Ich kann den Begriff des Körpers vorher analytisch durch die Merkmale der Ausdehnung, der Undurchdringlichkeit, der Gestalt \&c.,
% printed Fraktur "rc." (= etc.)
die alle in diesem Begriffe gedacht werden, erkennen. Nun erweitere ich aber meine Erkenntniß, und, indem ich auf die Erfahrung zurücksehe, von welcher ich diesen Begriff des Körpers abgezogen hatte, so finde ich mit obigen Merkmalen auch die Schwere jederzeit verknüpft, und füge also diese als Prädicat zu jenem Begriffe synthetisch hinzu. Es ist also die Erfahrung, worauf sich die Möglichkeit der Synthesis des Prädicats der Schwere mit dem Begriffe des Körpers gründet, weil beide Begriffe, ob zwar einer nicht in dem anderen enthalten ist, dennoch als Theile eines Ganzen, nämlich der Erfahrung, die selbst eine synthetische Verbindung der Anschauungen ist, zu einander, wiewohl nur zufälliger Weise, gehören. Kritik der reinen Vernunft. Werke. Zweiter Band. Leipzig 1838. S.~42--44.}% note *)

% --- p. 267 ---
\opage{267}Med den kantiske Modsætning mellem analytiske og synthetiske Domme parallelisere vi endvidere Modsætningen mellem apodiktiske og assertoriske. De assertoriske Domme ere Virkelighedsdomme, altsaa Erfaringsdomme, altsaa synthetiske; de apodiktiske Domme ere Nødvendighedsdomme, altsaa Reflexionsdomme, altsaa analytiske. Men nu har det i Læren om Dommens Form viist sig, at den analytiske Dom nødvendig forudsætter den synthetiske, og omvendt; altsaa maa i enhver Dom det Aprioriske nødvendig forudsætte det Empiriske, og omvendt. Virkelighedsdommen maa, for at være Dom, reflectere Nødvendigheden, og Nødvendighedsdommen, forsaavidt Nødvendigheden er objectiv, Virkeligheden; idet Nødvendighedsdommen reflecterer Virkeligheden, bliver den analytiske Dom empirisk, antilogisk; idet Virkelighedsdommen reflecterer
% --- p. 268 ---
\opage{268}Nødvendigheden, bliver den synthetiske Dom apriorisk, logisk. Forholdet mellem det Empiriske og det Aprioriske i vore Erfaringsdomme bestemmer Kant paa følgende Maade: „alle Erfaringsdomme ere empiriske, men alle empiriske Domme ere derfor ikke Erfaringsdomme. Saadanne Domme som: Værelset er varmt; Sukkeret er sødt, Malurten er bitter, ere empiriske; men som blotte Perceptionsdomme have de kun subjectiv Gyldighed, og ere desaarsag ikke Erfaringsdomme. Perceptionsdommen udtrykker kun, at to Sandsefornemmelser til en given Tid ere forenede i mig, det sandsende Subject.“ Dette er, som vi ovenfor have seet, en Opfattelse af den assertoriske Dom fra den umiddelbart subjective Side. Anderledes forholder det sig med Erfaringsdommen. „Hvad Erfaringen under visse Omstændigheder lærer mig, maa den lære mig og Enhver til enhver Tid“ \dots\ „Derfor udtaler jeg alle slige Domme som objectivt gyldige, som f.\ Ex.\ naar jeg siger: Luften er elastisk, saa er denne Dom nærmest kun en Perceptionsdom; jeg henfører to Sandsefornemmelser til hinanden. Vil jeg, at det skal kaldes en Erfaringsdom, saa fordrer jeg, at Forbindelsen stilles under en Betingelse, der gjør den almeengyldig.“ Men, for at en umiddelbar Sandsningsdom skal kunne blive til en Erfaringsdom, fordres, at Sandsningen subsumeres under et eller andet „Forstandsbegreb“; „Luften, f.\ Ex., hører under Aarsagsbegrebet; dette Begreb bestemmer Dommen om Luften med Hensyn paa Udvidelse som hypothetisk. Derved bliver nu denne Anskuelse forestillet, ikke blot som hørende til min Sandsning af Luften i min Tilstand eller til en Tilstand for Andres Sandsning, men som \emph{nødvendig};“ Dommen bliver almeengyldig, altsaa Erfaringsdom, derved, „at visse Domme, der subsumere Anskuelsen af Luften under Aarsagens og Virkningens Begreb, gaae forud og derved bestemme Sandsefornemmelserne, ikke blot respectivt ved hinanden i mit Subject, men i det Hele med Hensyn
% --- p. 269 ---
\opage{269}paa Dommens Form (her den hypothetiske) og paa denne Maade gjøre den empiriske Dom almeengyldig“.\footnote{Jvfr.\ Kant's kleinere metaphysische Schriften. Werke. Dritter Band. Leipzig 1838. Prolegomena zu jeder künftigen Metaphysik. S.~215 flg.}% note *)

Heraf sees, at den kantiske Erfaringsdom falder mellem den assertorisk-objective og den antilogisk-apodiktiske: i sin Eenhed med Perceptionsdommen udtrykker den det Assertoriske, i sin Eenhed med Forstandskategorien udtrykker den det Apodiktiske. Men i tre Punkter er Misviisningen kjendelig.

For det Første maa Kant indskrænke det Apodiktiske til det psychologisk Subjective, efterdi det Aprioriske hos ham kun har psychologisk-subjectiv Gyldighed. „Vi maae“, siger han, „adskille Naturens empiriske Love, der altid forudsatte
% UNSURE: forudsatte or forudsætte (degraded print; sense "forudsætte")
særegne Sandsninger, fra de rene eller almene Naturlove, som, uden at særegne Sandsninger ligge til Grund, blot indeholde Betingelserne for deres nødvendige Forening i Erfaringen, og i Henseende til de sidste er Natur og \emph{mulig} Erfaring ganske og aldeles eens, og da Lovmæssigheden i denne beroer paa Phænomenernes nødvendige Sammenknytning i en Erfaring, uden hvilken vi aldeles ikke kunne erkjende nogen Gjenstand i Sandseverdenen, altsaa paa Forstandens oprindelige Love, saa lyder det vel i Begyndelsen noget underligt, men er ikke desmindre vist, naar jeg med Hensyn paa de sidste siger: \emph{Forstanden henter ikke sine Love} (\textit{a priori}) \emph{fra Naturen, men foreskriver den disse}.\footnote{Prolegomena z.\ j.\ k.\ Metaphys.\ S.~240.}% note **)
% PRINTER: no closing quotation mark after "disse." for the Kant quotation opened at "Vi maae" / "adskille"; left unbalanced as printed

For det Andet mystificerer Fornuftkritiken Forholdet mellem „Synstingene (\textit{Phænomena})“, der udgjøre „Sandseverdenen“ og „de særskilte Forstandsvæsener (\textit{Noumena})“, der skulle udgjøre „Forstandsverdenen“. Kritiken indskjærper nemlig paa det Eftertrykkeligste, „at vi aldeles ikke vide og ikke kunne vide noget Bestemt om disse rene Forstandsvæsener, efterdi vore rene Forstandsbegreber ligesom ogsaa vore rene
% --- p. 270 ---
\opage{270}Anskuelser kun gaae ud paa Gjenstande for en mulig Erfaring, altsaa paa blotte Sandsevæsener“. Da nu „Forstandsvæsenerne“ antages for ubestemmelige, mangler enhver Mægling mellem de almene Kategorier og de individuelle Gjenstande eller, nøiagtigere, mellem det i Kategorierne Almene og det i Gjenstandene Individuelle.

For det Tredie er den kantiske Kritik ikke istand til at angive, hvorledes det i sig Objective egentlig forholder sig til Sandsning og Tænkning; om det f.\ Ex.\ først gjennem de rene Anskuelser individualiserer og bestemmer det Sandselige og saaledes udfylder Forstandsformerne; eller om det maaskee først bestemmer Forstandsformerne og gjennem disse de til Anskuelserne svarende Sandseindtryk; eller om det paa een Gang, altsaa lige oprindelig, udfylder og bestemmer baade Anskuelsesformer og Kategorier. I første Tilfælde har det ubekjendte Objective størst Affinitet til det Sandselige, i andet Tilfælde til det Intellectuelle, i tredie Tilfælde er det ligelig beslægtet med begge.

Den første Misviisning hæves paa følgende Maade. Reflecteret i den ontologiske Subjectivitet har det Aprioriske ikke blot subjectiv, men ogsaa objectiv Gyldighed, ikke blot Gyldighed i Viden, men ogsaa Gyldighed i Magt. Det er altsaa ikke den menneskelige Forstand, der foreskriver Naturen Love; den erkjender tvertimod i Naturen Lovenes objectivt-aprioriske Gyldighed: til denne Erkjendelse svarer det Apodiktiske i Erfaringsdommene. Enhver Erfaringsdom kan nu være baade assertorisk og apodiktisk, eftersom den tilsvarende Umiddelbarhed er \emph{blot umiddelbar} eller \emph{reflecteret}. Bestemt ved Gjenstanden i al Umiddelbarhed er Dommen assertorisk, er fra den subjective Side, hvad Kant kalder en Perceptionsdom (ein Wahrnehmungsurtheil), og bliver, naar Reflexionen sættes i Bevægelse, let problematisk. Bestemt, ikke ved Gjenstanden i dens Umiddelbarhed, men ved den i Gjenstandsforholdene sig aabenbarende Nødvendighed, er Dommen
% --- p. 271 ---
\opage{271}apodiktisk. Det er ikke Loven som Lov, der i det virkelige
Tilfælde udtrykker Nødvendigheden; thi Loven som Lov forudsætter
en frigjort oversvævende Reflexion i en bestemt Modsætning
til det enkelte Tilfælde, hvori den opfyldes. En saadan
Reflexion finder her ikke Sted; Nødvendigheden reflecterer
sig i Tilfældet selv, og Tilfældet virker paa Indsigten
med slaaende Umiddelbarhed. Naar Bordet hælder, og Glas
og Tallerkener ere i Begreb med at falde, saa gjør Huusmoderen
Anskrig: „Tag Dig iagt! det vælter, det maa absolut
vælte, dersom Du ikke“ o.\ s.\ v. Her have vi et Exempel
paa en antilogisk-apodiktisk Dom. Den ængstede Huusmoder
faaer uvilkaarlig et Indtryk af den betingede Nødvendighed.
Vel seer hun endnu ikke Ødelæggelsen, thi den er jo ikke virkelig
foregaaet; men hun seer den dog for sine Øine, hun
seer den umiddelbart og dog i Reflexion, thi hun indseer
\emph{Nødvendigheden}, „dersom ikke“ \dots\ Her falder i en saadan
Indsigt ikke først et Blik paa de almindelige Naturlove, særlig
paa Faldloven, dernæst et Blik paa de stedfindende Betingelser
og saa tilsidst en Forstaaelse af, at Faldloven i det faretruende
Tilfælde sørgelig vil opfyldes: en saadan Middelbarhed, en
saa endelig Udstykning staaer jo i latterlig Strid med det
Øieblikkelige i Indsigten. Og dog ere, trods det Umiddelbare
og Øieblikkelige, alle ved Udstykningen angivne Momenter
tilstede i Reflexionen; thi Reflexionen er \emph{øieblikkelig}:
det er den øieblikkelige Reflexion, hvori Indsigten seer Nødvendigheden.

Forvexles den øieblikkelige Reflexion med den umiddelbare
Sandsning, saa bliver den antilogisk-apodiktiske Dom, Nødvendighedsdommen
om det enkelte virkelige Tilfælde, forvexlet
med en assertorisk Dom, en umiddelbar Virkelighedsdom, og
Confusionen er i fuld Gang. Gives der ingen apodiktisk
Dom om det enkelte virkelige Tilfælde; er enhver Virkelighedsdom
uden videre assertorisk: saa er Nødvendigheden i sin objectivt
virkelige Nærværelse een Gang for alle gaaet tabt. Vi
% --- p. 272 ---
\opage{272}maae da enten med Hume give Affald paa al virkelig Nødvendighed,
eller med Kant lade den menneskelige Forstand
foreskrive Naturen Love.

Misligheden ved „de særskilte Forstandsvæsener“ lader
sig nu ogsaa hæve. Hvad ere „de særskilte Forstandsvæsener“
vel Andet end Existenskategorier, kategoriske Substantialitetsformer,
objectivt aprioriske Magtbegreber, hvorved Tingen i
existentiel Enkelthed viser sig som en Individuation af sit Almene.
Saadanne Udtryk som Kulstof, Qvælstof, Bor, Silicium
o.\ dsl. betegne paa een Gang Kategorier og existerende
Stoffer. Som Kategorier ere de „særskilte Forstandsvæsener“
objectivt aprioriske Substantialitetsformer (\textit{Noumena}); som
kategorisk bestemte Stoffer ere de existerende Substanser, Synsting
i stofagtig Umiddelbarhed (\textit{Phænomena}). Visheden om,
at Substansen Kul existerer i mangfoldige Gjentagelser af
virkelig tilværende Kulstof, er antilogisk-apodiktisk. Det Antilogiske
bestemmes ved den umiddelbare Kulstof-Existens, og
viser i denne Umiddelbarhed hen paa det Assertoriske; men
det Apodiktiske, der bestemmes ved Kul-Substansen, beroer paa
øieblikkelig Reflexion og sikkrer Indsigten i det Almeengyldige.
Saadanne Domme som: „Kulstoffet forekommer i reen Tilstand
som Diamant og som Graphit, mindre reent som Steenkul,
Anthracit“ o.\ s.\ v., ere fra Chemiens Standpunkt antilogisk-apodiktiske.
Det Samme gjælder om Domme, der
umiddelbart angaae Stoffernes Forbindelser. „Med Chlor og
Fluor danner Boret luftformige Forbindelser, der ryge stærkt
i Luften; med Qvælstof danner det et hvidt, amorpht, ildbestandigt
Legeme, der phosphorescerer stærkt med et grønligt Lys,
naar det opvarmes i Luften“ o.\ s.\ v. Dersom slige Domme
ikke vare \emph{apodiktiske}, vare de uden videnskabeligt Værd,
dersom det Apodiktiske ikke var \emph{antilogisk}, vare de ikke
Virkelighedsdomme, altsaa ikke Erfaringsdomme, ikke betegnende
for en Erfaringsvidenskab.

% --- p. 273 ---
\opage{273}Det ligger i Sagens Natur, at en Virkelighedsdom,
der ifølge den antilogiske Umiddelbarhed er assertorisk, maa
kunne blive problematisk; forsaavidt udgiver Erfaringsvidenskaben
da heller ikke sine Domme om de enkelte „Synsting“
i de enkelte Tilfælde for ufeilbare. Men hvormange Correctiver
Dommene om det Enkelte end undergaae, maa Erfaringsvidenskaben
dog nødvendigviis --- altsaa apodiktisk --- fastholde
Almeengyldigheden af det Grundlag, hvorpaa Bedømmelsen
oprindelig hviler: \emph{de assertoriske Domme om det Enkelte
ere Individuationer af apodiktiske Domme
om det Almene}. Kunne Dommene om Kulstoffets, Borets
og Chlorets Forbindelser i de enkelte Tilfælde med de i det Enkelte
bestemte Betingelser end mangfoldig corrigeres og ændres,
saa er det dog apodiktisk vist, at Kulstof altid er Kulstof,
Bor Bor og Chlor Chlor; det vil sige: enhver af de existerende
Substanser er under alle Forandringer sig selv liig; det Uforanderlige
virkeliggjør sig i sine Forandringer, det Almene realiseres
i det Individuelle, det Nødvendige aabenbarer sig i de
enkelte Tilfælde.

Vi kunne i denne Henseende give Kant Ret, naar han
indskjærper, at de rene Forstandsbegreber ikke have nogen Betydning,
„wenn sie von Gegenständen der Erfahrung abgehen
und auf Dinge an sich selbst (\textit{Noumena}) bezogen werden
wollen“;\footnote{Prolegom. S.~232.}% note *)
{} men ved at udskyde \textit{Noumena} gjør han netop
Forstandsbegrebernes Eenhed med Gjenstandene uforklarlig.
Begreber som: Kulstof, Qvælstof, Bor, Chlor o.\ s.\ v. kunne
umulig gjælde for rene, i sig existerende Væsenheder; det er
saa langt fra, at de særskilte Substanskategorier, Kants „Forstandsvæsener“,
skulde have selvstændig Tilværen som „Dinge an
sich“ (\textit{Noumena}) ligeover for de existerende Stoffer, at den almene
Kategorie i sin Adskillelse fra de særskilte meget mere
selv er en Skygge. Det er just det Særskilte, der mægler
% --- p. 274 ---
\opage{274}Foreningen af det Almene med det Individuelle. Forskjelseenheden
af den almene Substans (Kategorien) med den existerende
Substans (Gjenstanden) formedelst den særskilte Substans,
(„Forstandsvæsenet“), Forskjelseenheden af almindelig Substans
med Kulsubstans i Kulstoffet, af almindelig Substans med
Qvælstofsubstans i Qvælstoffet, af almindelig Substans med
Svovlsubstans i Svovlet, o.\ s.\ v., er en nødvendig, ved sin
virkelige Nødvendighed antilogisk-apodiktisk Eenhed. Fastholder
Reflexionen derimod det subjectivt aprioriske Forstandsbegreb
for sig, medens Sandseanskuelsen for sig er fængslet af det
existerende Stof, og „Forstandsvæsenet“ (\textit{Noumenon}), den
særskilte Substans, udelukkes som et i sig Ubestemmeligt, saa
miskjendes den virkelige Nødvendighed og med den det Antilogisk-Apodiktiske.

Den tredie Mislighed hos Kant er kun en Modification
af de to foregaaende; dens Ophævelse maa derfor betragtes
som en nærmere Belysning af de allerede anbragte Correctiver.
Den nødvendige Eenhed af det Aprioriske og det Empiriske,
det Almene og det Individuelle, det Intellectuelle og
det Sandselige er, kategorisk bestemt, en Magteenhed og kun
som saadan en virkelig Nødvendighedseenhed. Den kategoriske
Dom udtrykker her Ligevægten mellem det Sandselige og det
Intellectuelle; dette (et Sandseligt) er Kulstof (en Substanskategorie).
Medens nu den assertoriske Dom lægger umiddelbart
Eftertryk paa det Sandselige, uden at opgive det Intellectuelle,
lægger den apodiktiske Dom et afgjørende Eftertryk
paa det Intellectuelle, uden dog som antilogisk-apodiktisk at
opgive det Sandselige. Seet under Belysning af begge Domme
er den paagjældende Existens et Magtens Ord, med sin virkelige
Nødvendighed et Element i det kategoriske Magtsprog.
I den virkelige Nødvendighed, Nødvendigheden for den øieblikkelige
Reflexion i det slaaende Tilfælde, er det Intellectuelle
et Sandseligt, det Sandselige et Intellectuelt.

Det er en stor Vildfarelse, en Vildfarelse, som just ikke
% --- p. 275 ---
\opage{275}ligefrem, men dog indirect stammer fra den formelle Logiks
Abstractioner, at det i Virkeligheden Nødvendige skulde lade
sig udfinde ved Induction og Analogie, altsaa ved en sammenlignende
Betragtning af flere enkelte Tilfælde, uagtet hver
enkelt Tilfælde kun gav en assertorisk Vished. Paa denne Maade
skulde da et vist Antal assertoriske Domme omsider kunne
lede til en apodiktisk; hvor tankeløst! Den virkelige Nødvendighed
er i Dommen bestemt ved oprindelig Indsigt i det
Virkelige selv; det Virkelige selv er i antilogisk Umiddelbarhed
det virkelige Tilfælde: den apodiktiske Vished er i Virkeligheden
ligesaa oprindelig som den assertoriske. Den Reflexion,
hvorved den dømmende Subjectivitet umiddelbart indseer Nødvendigheden
i det virkelige Tilfælde, er en \emph{øieblikkelig}; fra den
øieblikkelige, ved det i Virkeligheden umiddelbare Væsensindtryk
bestemte, Reflexion har det Antilogisk-Apodiktiske sit Udspring.\footnote{Hvorledes
Inductionen og Analogien paa det herved givne Grundlag
udfolde deres Betydning, vil blive at paavise i en senere
Sammenhæng.}% note *)

3) \emph{Det Rationelt-Apodiktiske: absolut Nødvendighed.}
Som Nødvendigheden er, saaledes er Indsigten;
som Indsigten i det Nødvendige er, saaledes er Visheden;
som Visheden er, saaledes er det Apodiktiske. Den væsentlige
Nødvendighed, Reflexionsnødvendigheden, er logisk objectiv og
forsaavidt ontologisk subjectiv; den virkelige Nødvendighed, Indtryksnødvendigheden,
er objectiv, factisk og forsaavidt antilogisk:
Eenheden af den væsentlige og virkelige Nødvendighed maa
altsaa være en Eenhed, der i sig selv har det Logiske og
magter det Antilogiske; denne den absolute Nødvendighed fremstillende
Eenhed er den \emph{rationelle}.

Vistnok opfattes „det Rationelle“ ikke sjelden saa eensidigt,
at det derved faaer snart „det Irrationale“, snart „det Suprarationale“,
snart „det Extrarationale“ til Modsætning; men
for Videnskabens Idee ere disse Modsætninger alle forsvindende,
% --- p. 276 ---
\opage{276}thi Videnskabens Idee er absolut rationel. Hegels Ord: „det
Virkelige er det Fornuftige“, er ofte nok blevet misforstaaet,
og Misforstaaelsen atter og atter tilbageviist. At det Virkelige
er fornuftigt, vil jo ingenlunde sige, at ethvert Virkelighedens
Brudstykke, enhver enkelt Tilstand, Begivenhed eller Handling
ligefrem skulde tilfredsstille Fornuften. Tvertimod! Der er
i Endeligheden med dens Udstykninger og Forqvaklinger uendelig
Meget, som umiddelbart staaer i Strid med Fornuftens
evige Fordring; men middelbart hører det Endelige i al dets
Strid og Forvirring naturligviis med til den virkelige Udvikling
af Fornuftens Indhold. Hvad der umiddelbart strider
mod Fornuften, maa middelbart tjene dens Formaal: Fornuften
er ikke blot i Ideen den uendelige Viisdom, men i
Virkeligheden den overgribende Magt, hvormed Ideen underlægger
sig alle Betingelser.

I hvilken Betydning den absolute Nødvendighed er Formen
for Fornuftens Indhold, sees af den Maade, hvorpaa
Alt, alle Forskjelligheder, alle Modsætninger og Modsigelser
indgaae i Nødvendigheden. Saaledes er det Umulige i formel
Forstand et Nødvendigt ligesaavel som det Mulige. Thi
dersom der ingen Grændse var mellem Muligt og Umuligt,
dersom Alt var lige muligt og just derfor lige umuligt, saa
var der ingen Nødvendighed; men dersom Grændsen er nødvendig,
saa er jo det Umulige --- for Grændsens Skyld --- ligesaa
nødvendigt, som det Mulige. Det Uvæsentlige er et Nødvendigt
ligesaa vel som det Væsentlige. Thi dersom der ingen
Forskjel var mellem det Væsentlige og det Uvæsentlige, dersom
Alt var lige væsentligt og netop af den Grund lige uvæsentligt,
saa var der ingen Nødvendighed; men dersom Forskjellen
er nødvendig, saa er jo det Uvæsentlige, om ikke for Andet,
saa for Forskjellens Skyld, nødvendigt. Hvad der er nødvendigt,
maa ogsaa være væsentligt: det Uvæsentlige er altsaa
med Nødvendighed væsentligt. Saaledes kunde man blive ved
i det Uendelige.

% --- p. 277 ---
\opage{277}Det Vexelspil af Former, hvorunder Nødvendigheden
stadfæster og igjen opløser enhver relativ Bestemmelse, har i
Virkeligheden ikke blot formel, men ogsaa reel Gyldighed.
Der maa i det Virkelige selv være noget Uvæsentligt, for at
det Væsentlige kan være. Forsaavidt det Uvæsentlige paa
denne Maade bliver væsentligt, synes det jo vistnok, som om
den virkelige Grændse mellem Væsentligt og Uvæsentligt var
udslettet; men dette er kun en Indbildning. Den Nødvendighed,
der giver det Uvæsentlige en væsentlig Betydning, den
samme Nødvendighed fordrer, at det Uvæsentlige netop faaer
Betydning som et Uvæsentligt. Endnu skarpere og paa en
mere existentiel Maade fremtræder den nødvendige Modsætning
mellem det Nødvendige selv og det Tilfældige. Det er f.\ Ex.
umuligt at drikke \emph{en} Viin, uden at drikke \emph{denne} Viin. Men
hvorfor nu just denne? Det kunde jo ligesaa godt være enten
denne eller denne \dots\ eller denne! Hvilken „\emph{Denne}“ man
vælger, er, forsaavidt det blot kommer an paa at vælge en Viin,
aldeles ligegyldigt: en „Denne“, just en „Denne“ er altsaa tilfældig.
Men hvilken man end vælger, saa er det dog aldeles
uundgaaeligt, at man kommer til at vælge en „Denne“; en
„Denne“, just en „Denne“ er altsaa nødvendig. Den Indsigt,
at det Tilfældige ikke er nødvendigt, er kun fornuftig i Forening
med den Indsigt, at det Tilfældige dog er nødvendigt. At det
Tilfældige er nødvendigt, vil da ikke sige, at Modsætningen
af Nødvendigt og Tilfældigt skal opgives, men tvertimod, at
den skal hævdes. Paa denne Maade har enhver relativ Bestemmelse
sin bestemte Gyldighed d.\ v.\ s. sin Nødvendighed: i
Totaliteten af de relativt nødvendige Momenter bestaaer den
absolute Nødvendighed; den absolute Nødvendighed er Formen
for den uendelige Fornuft.

Vilde Nogen nu spørge, om man da ikke for Fornuftens
Skyld maa fordre noget Ufornuftigt paa samme Maade, som
man for det Væsentliges Skyld maa fordre noget Uvæsentligt
eller for det Nødvendiges Skyld noget Unødvendigt: da ligger
% --- p. 278 ---
\opage{278}Svaret ikke langt borte. Vistnok maa der være en Forskjel,
en væsentlig Forskjel, mellem Fornuftigt og Ufornuftigt, ligesom
mellem Nødvendigt og Unødvendigt; men ligesom den
Forskjel, hvorved det Unødvendige selv bliver nødvendigt, falder
ind under den absolute Nødvendighed, saaledes falder ogsaa den
Nødvendighed, hvormed det Ufornuftige, --- hvormeget det end
strider mod det Fornuftige, --- melder sig, ind under den absolute
Fornuft. Seet i Ideen kan den absolute Fornuft ligesaa lidt
have noget Ufornuftigt imod sig, som den absolute Nødvendighed
kan have noget Unødvendigt udenfor sig. Heller ikke kan det
paa Ideens Standpunkt gaae an at give det Rationelle en
Modsætning i det Antirationelle paa tilsvarende Maade, som
naar det Antilogiske sættes imod det Logiske.

Modsætningen af det Logiske og det Antilogiske betyder,
at det Logiske, hvormeget det end for Reflexionen viser sig som
et uendelig Overgribende, dog i sit Væsen er eensidigt, ligesom
den rene Viden er en Eensidighed ligeoverfor Magten. Anderledes
med det Rationelle; her er Viden i Eenhed med Magten,
thi Magten er i sit Væsen ligesaa rationel som Viden.\footnote{Kun
paa Troens Standpunkt, altsaa i en Sphære, der falder
udenfor Videnskaben, kan der være Tale om det Antirationelle;
det Antirationelle bliver da Kjendemærket paa det Paradoxe.}% note *)
{} Den uendelige Fornuft bestemmer det Logiske ved det Antilogiske,
og omvendt. Visheden om Ideens, d.\ e. den absolute
Fornufts Gyldighed er rationel-apodiktisk. Da nu det Antilogiske
har sin rationelle Gyldighed ligesaavel som det Logiske,
saa har ogsaa Virkeligheden i sin Forskjel fra Muligheden og
Nødvendigheden, Muligheden i sin Forskjel fra Nødvendigheden
og Virkeligheden, Nødvendigheden i sin Forskjel fra Muligheden
og Virkeligheden rationel-apodiktisk Gyldighed i det uendelige
Fornuftrige. De assertoriske Domme kunne i Ideen
ligesaa lidt forvandles til apodiktiske, som de problematiske
kunne det; men alle tre Slags Domme gjælde med relativ
% --- p. 279 ---
\opage{279}Nødvendighed; for den fuldkomne Viden er Gyldigheden af
dem alle rationel-apodiktisk.

For at indsee, hvorledes det er muligt, at Systemer af
assertoriske, problematiske og apodiktiske Domme kunne have
Gyldighed ikke blot udadtil ved Objectiveringen af det Endelige,
men ogsaa indadtil i Subjectivitetens Selvobjectivering,
maae vi fremfor Alt erindre, at den ontologiske Selvobjectivering,
der overalt beroer paa en Objectivering af Andet, ikke
er tom, men uendelig indholdsrig. Indholdet er Fornuften
selv, Ideens absolute Indhold. Reflecteret i Modalitetens
tre Former vil Ideen og med den Subjectiviteten da antage
følgende grundforskjellige Skikkelser. I Forhold til Virkeligligheden,
% PRINTER: printed "Virkelig-|ligheden", sense "Virkeligheden"
altsaa i Systemet af de assertoriske Domme, er
Subjectiviteten anskuende, og Ideen i Virkelighedens Form en
Phænomenets Idee. I Forhold til Muligheden, altsaa i
Systemet af de problematiske Domme, er Subjectiviteten selvbestemmende,
og de forudbestemte Muligheders Idee en Formaalsidee.
I Forhold til Nødvendigheden er Subjectiviteten
tænkende, begribende, og Ideen i Nødvendighedens Form en
almeengjældende Væsensidee. Da Ideen under alle tre Skikkelser
er identisk med sig selv, er det en Selvfølge, at Modalitetens
tre Former maae forudsætte hinanden saaledes, at
den ene ved sin bestemte Fremtræden altid reflecterer de to
andre.

Hvad Phænomenets Idee angaaer, forudsættes i Virkeligheden
en umiddelbar Adskillelse mellem Phænomenet og
Ideen. Her maa i den virkelige Verden kunne gives Phænomener,
der falde tilfældig udenfor Ideen, til Forskjel fra
Phænomener, hvori Ideen umiddelbart lader sig tilsyne. Just
derved, at Ideephænomenet i sin ideelle Gjennemsigtighed træder
i Modsætning til det i sin Tilfældighed ideeløse Virkelighedsphænomen,
danner der sig en assertorisk Dom om det Skjønne,
svarende til Ideephænomenet, altsaa til en Idealform: „dette
er virkelig skjønt“, og en assertorisk Dom om det Uskjønne,
% --- p. 280 ---
\opage{280}svarende til det ideeløse, mod Idealformen stridende Phænomen:
„dette er virkelig uskjønt“. Striden mellem Ideephænomenet
og det ideeløse Virkelighedsphænomen udlignes i Muligheden:
begge Phænomener kunne paa utallige Maader ændres,
idet Ideen, som formedelst Mulighederne i det Uendelige
magter, opfylder og gjennemtrænger de ideeløse Virkelighedsphænomener,
tilfredsstiller den sig selv i Alt og Alt i sig selv
anskuende Subjectivitet. Er Dommen om det Skjønne og
det Uskjønne paa Grund af Phænomenets Umiddelbarhed
assertorisk, saa er Dommen om Skjønheden selv, d.\ v.\ s. den
Dom, hvorved Dommen om Forskjellen, den væsentlige Forskjel,
mellem Skjønt og Uskjønt oprindelig fastsættes, derimod
apodiktisk, rationel-apodiktisk.

Hvad Formaalsideen, de bestemte Muligheders Idee, angaaer,
er det ved sig selv indlysende, at den altformaaende
Subjectivitet formedelst sin Magt over Andetheden, altsaa
over den umiddelbare Virkelighed, maa leve i et uendeligt
System af Muligheder. Men et uendeligt System af Muligheder
er kun tænkeligt i et uendeligt System af problematiske
Domme. „Gud er det, at Alt er muligt“; det vil
her sige: for det absolute Selv staaer intet Endeligt fast;
alle Ting sees i uendelig Vorden, i Overgange, i Muligheder.
Gjennem Mulighedernes uendelige System hævder den guddommelige
Subjectivitet sin uendelige Selvbestemmen ved
særlig Bestemmen af Andet: den forvandler det Virkelige til
Betingelser og sætter den ved Anticipationen bestemte Mulighed
som Formaal. Det fornuftige Formaal kan altsaa realiseres
paa Grund af Mulighederne; men det kan i Virkeligheden
ogsaa deelviis forstyrres, forqvakles og hæmmes paa
Grund af Mulighederne. Dommen om Formaalets, det Godes,
Opnaaelse er i Endeligheden betinget og problematisk; men i
Ideen selv rationel-apodiktisk. Ligesom de assertoriske Domme
paa Virkelighedens Grundlag ere Individuationer af det Apodiktiske,
der svarer til Phænomenalideens Nødvendighed, saa%
% --- p. 281 ---
\opage{281}ledes ere de problematiske Domme paa Mulighedens Grundlag
Specificationer af det Apodiktiske, der i guddommelig Indsigt
svarer til Formaalsideens Nødvendighed.

Hvad endelig Væsensideen, den almene Nødvendighedsidee,
angaaer, vil den tænkende, Alt i Videns og Magtens
Eenhed begribende Subjectivitet objectivere sig Alt, det Mulige
saavel som det Virkelige, det Evige saavel som det Timelige,
det Indre og Ideelle saavel som det Ydre og Reelle i Ordet,
ved Ordet som Ordet. Det eneste for den absolute Fornufts
Selvaabenbarelse absolut Nødvendige er Ordet; i Ordet frembringes
Alt, ved Ordet bestemmes Alt: den almene Nødvendighedsidee
er Ordets Idee. Hvad er vel Dommens Gjenstand?
Et Magtens Ord! Og Dommen selv? Et Videns
Ord! Og Dommens Gyldighed? Et Virkelighedens, Mulighedens
og Nødvendighedens Ord! Eller hvorledes skulde Identiteten
af Dommen og Dommens Gjenstand og dermed Sandheden
kunne indlyse, dersom ikke Gjenstanden, trods sin virkelige
Modsætning til Dommen, var i Ordet selv væsentlig
Eet med Dommen. I Ordet som Magtens Ord gaaer den
anskuende Subjectivitet med alle Magthandlinger og Domme
ind paa den objective, antilogisk herskende Virkelighedsnødvendighed;
i Ordet som Videns Ord gaaer den tænkende Subjectivitet
gjennem alle virkelige Modsætninger ind paa den
logisk-objective, ontologisk-subjective Væsensnødvendighed; i den
reflecterede Eenhed af Magtens og Videns Ord gaaer den
altbegribende Subjectivitet ind paa den absolute Nødvendighed,
et System af rationelt-apodiktiske Domme, hvori Alt har sin
relative Sandhed, fordi Alt har sin relative Nødvendighed.

Dommen i og for sig er ikke Sandheden, thi i Sandheden
maa Viden forudsætte Magten; Gjenstanden i og for
sig er ikke Sandheden, thi i Sandheden maa Magten forudsætte
Viden. Er nu Ordet i Dommen for sig et Videns Ord,
og Ordet i Gjenstanden for sig et Magtens Ord, hvorledes
% --- p. 282 ---
\opage{282}skal da Eenheden af dette i sig fordoblede, ved Fordoblingen sig
selv oprindelig modsatte Ord nærmere bestemmes?

\begin{center}{\bfseries b)\quad Functioner af Slutningen: Modalfunctioner.}\end{center}

Lader man Dommen om en Gjenstand synke ned til
simpel Forestilling og bestemmer Forestillingen, ikke ved et
Ord, men umiddelbart ved Gjenstandens Billede, da skulde
man ved første Blik formode, at Overeensstemmelsen mellem
Forestilling og Gjenstand paa denne Maade vilde være sikkret,
og den saaledes sikkrede Overeensstemmelse den bedste Garantie
for Dommens Sandhed. Men at et Sandhedskriterium paa
disse Vilkaar ikke lader sig opstille, have de gamle Skeptikere,
navnlig ved deres Indsigelser imod Stoikernes \textit{φαντασία
καταληπτική},\footnote{Definitionen paa denne iøvrigt noget tvetydige Kategorie er hos
Sextus (\textit{Adversus Logicos I. §~402.}) angivet i følgende Udtryk:
\textit{ἤν γὰρ καταληπτικὴ φαντασία ἡ ἀπὸ ὑπάρχοντος, καὶ κατ᾿
αὐτὸ τὸ ὑπάρχον ἐναπομεμαγμένη καὶ ἐναπεσφραγισμένη,
% UNSURE: punctuation after ἐναπεσφραγισμένη looks like "." in print; text layer ","
ὁποία οὐκ ἄν γένοιτο ἀπὸ μὴ ὑπαρχοντος}, altsaa: den Forestilling,
der har sin Oprindelse fra den existerende Ting, og i fuldkommen
Overeensstemmelse med denne er saaledes indpræget og indprentet
i Sjælen, at den umulig kan hidrøre fra en ikke existerende
Ting. For Kortheds Skyld bruge vi her Benævnelsen: den gribende
Forestilling.}% note *)
{} tilstrækkelig viist. De Forhandlinger, som
Sextus Empirikus har opbevaret os angaaende dette Punkt, ere
i høi Grad charakteristiske. „Der gives“ hedder det „ingen Objectet
gribende Forestilling saa eiendommelig, at den jo kan hidrøre
fra et andet Object end det, hvortil den just i et givet Tilfælde
henføres; dette have foruden Skeptikerne endogsaa Akademikerne
og da fornemmelig Arkesilaus og Karneades fra
mange Sider stræbt at godtgjøre. Det Bifald (\textit{συγκατάθεσις}),
som Stoikerne kalde Opfatning (\textit{κατάληψις}) og gribende
Forestilling (\textit{φαντασία καταληπτική}) er, ifølge Arkesilaus,
kun et tomt Ord; thi et saadant Bifald, hvis det gives,
% --- p. 283 ---
\opage{283}maa enten gives af den Vise eller af Daaren. Men gives
det af den Vise, er det en Viden (\textit{ἐπιστήμη}), og gives det
af Daaren, er det kun en Mening (\textit{δόξα}). Skal det Bifald, vi
skjænke den gribende Forestilling, være Opfatningen, da har
denne ingensteds hjemme; thi for det Første er det jo ikke
Forestillingen, men Tanken, vi skjænke vort Bifald; for det
Andet findes der ingen Forestilling saa sand, at den jo ogsaa
kunde være falsk“.\footnote{\textit{Sext. Empirici Opera ed. Fabricius. Lipsiæ. Anno MDCCXVIII.
Adversus Logicos. I. §~143--54.}}% note *)
{} At adskille den sande Forestilling fra
den falske er ikke muligt; thi at ville foretage en saadan Adskillelse
ved at sige med Stoikerne: den sande eller gribende
Forestilling hidrører fra en existerende Gjenstand, den falske
eller ikke gribende fra en ikke existerende, er forgjæves, da der
jo gives Indbildninger, som i Klarhed og Livlighed ikke staae
tilbage for de sande, til Virkeligheden svarende Forestillinger.
„I sit Raseri tog Herkules sine egne Børn for Eurystheus's,
og handlede i Conseqvens af denne sin Indbildning, d.\ v.\ s.
han dræbte i Indbildningen sin Fjendes Børn, men i Virkeligheden
sine egne“.\footnote{\textit{Adversus Logicos. I. §~406.}}% note **)

Men selv om man gaaer ud fra den Forudsætning, at
den paagjældende Forestilling har sin Oprindelse fra en virkelig
Gjenstand, saa vil det, antager Skeptikeren, og det fra sit Synspunkt
--- saalænge her blot reflecteres paa endelig Sammenligning
mellem Forestilling og Gjenstand --- med fuldkommen
Ret, dog være umuligt at opvise en Forestilling, der skulde
give et saa tydeligt Afbillede af sin Gjenstand, at man ved
dette Afbillede vilde være istand til at adskille Gjenstanden fra
alle andre. Formerne kunne nemlig være ganske lige, og
Gjenstandene selv endda forskjellige. „Jeg giver“, siger Karneades,
idet han vælger Exemplet af to Æg, der ere hinanden
ganske lige, „Stoikeren at prøve det vexelviis snart med det
% --- p. 284 ---
\opage{284}ene og snart med det andet Æg, hvorvidt den Vise, naar han
seer dem, vil være istand til at sige, hvilket der er hvilket“.\footnote{\textit{Adversus Logicos. I. §~409.}}% note *)

Forsøger man at undvige Vanskeligheden ved saadanne
Udflugter som, at det jo ikke kommer an paa at angive Almeenbestemmelser,
der skulde indeholde Kjendemærker for hvert
enkelt Tilfælde, at Objecter, der ere hinanden saa lige, at de
i Forestillingen selv ikke kunne adskilles, derfor ingenlunde
tabe Charakteren af at være virkelige, o.\ s.\ v.: da gjør Skeptikeren
med stor Lethed et Spring fra det umiddelbart Enkelte
til det abstract Almene; thi hans Færdighed er, som vi
tidligere have seet, at springe fra Extrem til Extrem saaledes,
at han bestandig springer det forbindende Mellemled over.
Var Vanskeligheden i Extremet af det Enkelte den, at Dogmatikeren
med sin „gribende Forestilling“ ikke kunde adskille det
ene Æg fra det andet, saa bliver nu Vanskeligheden i Extremet
af det Almene den, at hans „gribende Forestilling“ kun kan
bestemmes ved det Object, der selv maa bestemmes ved Forestillingen.

Med Hensyn paa det Almene i Vexelforholdet af Forestilling
og Gjenstand lyder Akademikerens Indsigelse paa, at
Stoikerne gjøre sig skyldige i en Diallelus. „Naar vi nemlig
spørge, hvad en gribende Forestilling er, sige de: en saadan,
som hidrører fra en existerende Gjenstand, stemmer overeens
med denne, er aftrykt og afpræget i Sjælen, og af den Beskaffenhed,
at den ikke kan hidrøre fra en ikke existerende
Gjenstand.\footnote{\textit{Adversus Logicos. I. §~426--29.}}% note **)
{} Men da nu enhver Bestemmelse maa gaae ud
fra det Bekjendte, og vi derfor fremdeles spørge: hvad er det
Existerende? vende de tilbage til den første Bestemmelse og
sige: Existerende er det, som fremkalder en gribende Forestilling.
Saaledes viser da Bestemmelsen af den gribende
Forestilling tilbage til Bestemmelsen af det Existerende, og
% --- p. 285 ---
\opage{285}denne igjen til hiin; vi erholde altsaa hverken af Gjenstandsforestillingen
eller af Gjenstanden noget tydeligt Begreb.
Hvorvidt den forestillede Gjenstand har en blot tilsyneladende
eller baade tilsyneladende og virkelig Gyldighed, lader sig derfor
ikke afgjøre ved nogensomhelst Forestilling; skulle Gjenstandene
bedømmes efter den ikke gribende Forestilling, maa
jo denne, hvad Dogmatikerne selv ikke ville indrømme, kunne
bedømme det Sande. At antage den gribende Forestilling for
det Noget, \emph{hvorefter} det Sande skal bedømmes, er latterligt;
thi da der just søges et bestemt Kjendetegn for den gribende
Forestilling, maa der til enhver Forestilling, der opstilles som
gribende, bestandigt kræves en ny Forestilling til Stadfæstelse
af den første, og saaledes i det Uendelige“.\footnote{Jvfr. E.\ F.\ Nielsen. Den antike Skepticisme. S.~91--97.}% note *)

At hæve disse Vanskeligheder er umuligt, saalænge Standpunktet
udelukkende vælges i det Endelige; men for at det
Uendelige kan komme med, maa Forestillingen bestemmes ved
Dommen, og Modsætningen af Dom og Gjenstand opklares i
Ordet. Dette er den Tankegang, vi have fulgt i Læren om
Domsfunctionerne. Det begynder med, at den naturkraftige
Tænkning danner og former Sproget, at den i de udviklede
Sprogformer kan kjende og gjenkjende sig selv. Under denne
Udvikling har den med Sproget og Tankerne arbeidende Subjectivitet
ikke blot objectiveret sig en virkelig Yderverden, men
efterdi Sproget forudsætter og betinger en indre ideel Objectivering,
tillige en indvortes Verden, en Verden af væsentligt
Indhold.

Det væsentlige, ved Sproget articulerede Tanke-Indhold
er i Principet ligesaa oprindelig objectivt som subjectivt;
Identiteten af det subjective og objective Indhold er i Væsenet
selv \emph{Ideen}, men i Formens Umiddelbarhed \emph{Ordet}: de
sprog- og tankeformende Functioner ere paa denne Maade
blevne ideeformende. At de ideeformende Functioner fra deres
% --- p. 286 ---
\opage{286}eiendommelige Udgangspunkt begynde med en ligesaa naturkraftig
Umiddelbarhed som de sprogformende, sees af Mythedannelsen.
Mythologiens væsentlige Indhold udviser Momenter af logisk
Oprindelighed, Formbestemmelser, hvori en evig Fornuft har
fundet sin Afspeiling; vi have seet de ideelle Resultater, hvortil
Ophævelsen af det mythiske Bevidsthedstrin med indre Conseqvens
kan føre. Vel kunde de gnostiske Ideefunctioner ikke
ved sindrig Forening af Phantaseren og Tænken give det Absolute
nogen adæqvat Form; vel maatte de theosophiske Functioner,
der kun lade Tanken gribe det Absolute absolut ved
at underskyde Phantasiens Form og gjøre Formen formløs,
tabe sig i det Mystiske; men paa disse Overgangsstadier viste
det sig dog bestandig klarere og klarere, hvad det vil sige, at
Tanken ved Ordet fatter Ideen. De mythiske Guder, de
gnostiske Æoner, de theosophiske Væsenheder ere forskjellige
% UNSURE: mark over the second e of "Idee" (Ideë?) looks like an ink blemish; set Idee
Idee-Aabenbarelser, der af Tanken bestemmes ved Ordet, leve
og voxe sig i Ordet. Afklaret til Begreb, aabenbarer Ideen
sig da ogsaa tilsidst som Ordets guddommelige Væsen, Ordet
som Ideens evige Skikkelse.

Identiteten af Ord og Idee er tillige Identitet af Dom
og Indhold; Dommen absolut er Ordet, Indholdet absolut
er Ideen. Men Identiteten blev kun opfattet fra Formsiden;
dens almene Væsenhed er endnu uden Virkelighed. Fra Virkelighedens
Standpunkt kunde der da heller ikke begyndes med
en videre Udvikling af Dom og Indhold; der maatte begyndes
med den skarpe Modsætning af Dom og Gjenstand. Idet
Gjenstanden stilles selvstændig mod Dommen, skulde man næsten
frygte for, at den i udvortes Tilfældighed existerende Gjenstand
vilde unddrage sig enhver bestemt og tilforladelig Opfattelse,
saa at de gamle Qvaklerier med \textit{φαντασία καταληπτική}
maatte vende tilbage. Denne Mislighed hæves derved, at den
dømmende Subjectivitet paa overgribende Maade lader Dommen
selv være bestemmende for Forholdet af Dom og Gjenstand.
Den meest umiddelbare Selvstændighed, der kan
% --- p. 287 ---
ind\opage{287}rømmes Gjenstanden, er kjendelig derpaa, at det er Gjenstanden,
der antages at bestemme Dommen; Conseqvensen af
denne Antagelse er, at Alt opløses, baade Tankebestemmelserne
i Dommen og Realitetsbestemmelserne i Gjenstanden. Kan
det nu ikke være Gjenstanden, der bestemmer Dommen, saa
opstaaer det Spørgsmaal, om det da ikke maaskee omvendt er
Dommen, der bestemmer Gjenstanden? Ved denne Antagelse
træde nu vistnok Tankebestemmelserne i Kraft, men saa overgribende
er Idealitetens Magt, at Gjenstanden bukker under,
opløses i Begreb og forvandles til Dommens Indhold.
Gjenstandens umiddelbare Selvstændighed imod Dommen har
altsaa paa dobbelt Maade viist sig som Uselvstændighed: skal
den uden Begreb beroe paa sig selv, da bliver den et løst
Indbegreb af Egenskaber uden Substans, smuldrer hen i Tilfældigheder
og synker sammen; skal den i Kraft af Begrebet
beroe paa Tanken, altsaa paa Dommen, forvandles den til
et Phantom og forsvinder som Dunst.

For at Gjenstanden kan staae ved Magt, maa der være
en ideel Magt i Dommen; den ideelle Magt i Dommen som
logisk Dom forudsætter en reel Magt i den til Dommen
svarende Handlen, thi oprindelig Viden forudsætter Magt, og
Magt Viden. Gjenstanden i sin Umiddelbarhed er ikke Begreb,
men dens Existens har sin Grund i Begrebet; uden
Grund ingen Existens, uden Begreb ingen Gjenstand: som
Grunden er, maa Existensen være, som Begrebet er, maa
Gjenstanden være. Gjenstanden maa altsaa have ideel Tilværen
i sit Begreb, for at faae virkelig Tilværelse i Existensen:
Gjenstanden maa anticiperes; ved Anticipationen griber Dommen
ud over Gjenstanden, ligesom Subjectiviteten griber ud
over det Objective. At Gjenstanden bliver til paa objective
Betingelser, betyder da ingenlunde, at den løsrives fra Begrebet
eller ifølge Tingsaarsagers Samvirken opstaaer ved
Siden af Begrebet; nei, tvertimod! Gjenstandens Virkeliggjørelse
er netop Begrebets Virkeliggjørelse: det er
% --- p. 288 ---
\opage{288}igjen Dommen, der er overgribende; det er i Kraft af den
antilogiske Dom, at Gjenstanden realiseres. I den antilogiske
Dom er Tanken Magt, i den logiske er Magten Tanke; da
nu Tankens Udtryk oprindelig er Ordet, og Ordet articuleres
i Sætningen, saa bliver Gjenstanden til, væsentlig ved en
Tankesætning, virkelig ved en Magtsætning.

Som en i Tankesætning reflecteret Magtsætning er Gjenstanden,
seet i Begrebets Lys, mere end en umiddelbar Ting;
den er i Væsen en objectiv Dom, et Ord, et Magtens Ord:
Dom og Gjenstand forholde sig som Videns og Magtens
Ord. I dette Eenhedsforhold er Modsætningen bevaret; thi
Gjenstand staaer imod Dom som Magt imod Viden. Men
i Modsætningen er tillige Eenheden bevaret; thi Gjenstanden
er Magtsætning, ligesom Dommen er Tankesætning. Den
reelle Tanke, der er i Magtsætningen, svarer i Principet til
den ideelle Magt, der er i Tankesætningen; saaledes bliver i
det Uendelige Gjenstanden bestemt ved Dommen, og Dommen
igjen ved Gjenstanden. Dom og Gjenstand ere Eet i Ordet,
efterdi de ere Eet i Ideen.

Eenheden af Dom og Gjenstand er Modsætningseenhed;
af denne Modsætningseenhed fremlyser Modalitetens Magt.
Kategorierne Virkelighed, Mulighed og Nødvendighed vise ikke
blot hen til tre forskjellige Sider af Tilværelsen, men de vise
tillige, naar vi see nærmere til, at disse tre Sider ere tre
Totaliteter, saa at Synspunkterne vexle med de tre Domme:
Alt (nemlig alt Tilværende) er \emph{virkeligt}, \emph{Alt er muligt},
\emph{Alt er nødvendigt}. De tre Totaliteter danne da tre forskjellige
Sphærer, paa Grund af Modaliteten; men de tre
forskjellige Sphærer ere desuagtet Sidebestemmelser for een og
samme Tilværelse paa Grund af Modaliteten: Dom og
Gjenstand indgaae altsaa i tre forskjellige Forhold paa Grund
af Modaliteten. Vi skulle ved dette Tilbageblik belyse Modalitetsforholdene
endnu noget nærmere.

% --- p. 289 ---
\opage{289}Paa Virkelighedens Trin staaer Gjenstanden, skjøndt
væsentlig selv en objectiv Dom, i umiddelbar Selvstændighed
mod den subjective Dom; den ved Gjenstanden umiddelbart
bestemte Dom er assertorisk. Dersom nu Virkeligheden saaledes
udelukkede Muligheden, at denne kun kom op i Subjectiviteten
og fik et subjectivt uendeligt Spillerum ligeoverfor
hiin, saa vilde den hele Logik hjemfalde til Skepticismen, og
det gamle Spil med \textit{φαντασία καταληπτική} gjentage sig i
nye Vendinger. Men nu har Muligheden formedelst Modaliteten
objectiv Væren ligesaavel som Virkeligheden: „Alt er
muligt“, dette er ikke en abstract subjectiv Dom, der har udelukket
Objectiviteten, men Subjectivitetens oprindelige Dom
om alt det Objective. Det allerede Bestemte lader sig, forsaavidt
det er bestemt, ikke videre bestemme: kun det endnu Ubestemte
kan bestemmes. Staaer det nu fast, hvad ovenfor er viist, at
ingen Existensbestemmelse er tænkelig uden en tilsvarende Begrebsbestemmelse,
ingen Gjenstandsbestemmelse uden en bestemmende
Tanke; staaer det fast, at Grunden til alt Objectivt principielt
maa bestemmes subjectivt og kun kan bestemmes subjectivt,
forsaavidt det Objective endnu er problematisk, da maa
den overgribende Subjectivitet jo vide Oprindelsen til al endelig
Bestemthed som problematisk: „\emph{Alt er muligt}“; den problematiske
Dom er ikke blot menneskelig, men guddommelig.

I den problematiske Dom er Gjenstanden opløst; den er
gaaet under i Ubestemtheden, begravet i Uvisheden; den problematiske
Dom har vel et Indhold, men ingen Gjenstand.\footnote{At der paa Endelighedens Standpunkt kan fældes en problematisk
Dom over en i sig bestemt Gjenstand, følger deraf, at her foruden
Dommen er en Forestilling med, der lægger Gjenstanden selv udenfor
Dommen.}% note *)
Skal Dommen i oprindelig Forstand have baade Gjenstand
og Indhold, en til det bestemte Indhold svarende Gjenstand,
et til den bestemte Gjenstand svarende Indhold; skal, med
% --- p. 290 ---
\opage{290}andre Ord, Modsætningseenheden af Dom og Gjenstand naae
til fuldendt Bestemthed, da maa Modaliteten selv fuldendes
i Nødvendigheden, og den apodiktiske Dom blive den overgribende.
I Kraft af Nødvendigheden er Identiteten af Dommens
Indhold og Gjenstand bestemt paa to væsentlig forskjellige
Maader: Identiteten er sat i \emph{Indholdet}, og hertil
svarer det logisk Apodiktiske; Identiteten er sat i \emph{Gjenstanden},
og hertil svarer det Apodiktiske i antilogisk Form. Begge
Maader afgive med al deres Modsætning to Sider af det
Samme. Fra den antilogiske Side er Nødvendigheden skjult
i Virkeligheden, endnu bestemtere, i Tilfældigheden. Vil
Nogen paastaae, at Nødvendigheden, ved at falde sammen med
Tilfældigheden, enten maa gjøre det Tilfældige nødvendigt, og
altsaa ophæve Tilfældigheden, eller gjøre det Nødvendige tilfældigt,
og saaledes, hvad der er modsigende, afskaffe sig selv:
da kunde han vistnok have Gavn af at betænke, om Dommen:
„det Tilfældige er nødvendigt“, ikke ligesaa fuldt
hævder Betydningen af Tilfældigheden, som Dommen: „det
Nødvendige er tilfældigt“, Betydningen af Nødvendigheden.

Er Nødvendigheden med overalt, i alle Tilfælde og i tilfældige
Phænomener, uden dog --- hvad Dommen: „det Tilfældige
er nødvendigt“, forbyder --- at forvandle det Tilfældige
til Skin: da har Nødvendigheden, ved at give efter for Tilfældet,
respectere Virkeligheden og selv paatage sig alle Tilfældighedens
Skikkelser, netop inddraget alt Tilfældigt under
sig, og derved beviist sin uendelige Overmagt.

Som nu det Virkelige forholder sig til det Nødvendige,
saaledes forholder den assertoriske Dom sig til den apodiktiske.
Vil Nogen berøve den assertoriske Dom sin eiendommelige
Gyldighed og forvandle den til Reflex af den apodiktiske, da
protesterer Nødvendigheden; vil Nogen gjøre den assertoriske
Dom absolut uafhængig af den apodiktiske, da erkjender han
jo netop det Nødvendige i dens eiendommelige Charakteer,
døber den, saa at sige, med Nødvendighedens Daab og
% --- p. 291 ---
\opage{291}lader paa denne Maade det Apodiktiske indestaae for det
Assertoriske.

Hvad dernæst Modsætningen af den problematiske og den
apodiktiske Dom angaaer, er det strax indlysende, at den falder
sammen med Modsætningen af Mulighed og Nødvendighed.
Da nu Muligheden som Mulighed har sin Nødvendighed,
saa er Muligheden selv en latent Nødvendighed, og saa følger
heraf, at der i den problematiske Dom er et apodiktisk Moment,
og at dette apodiktiske Moment, langt fra at tilintetgjøre det
Problematiske, tvertimod hævder det en ontologisk Ret, en i
Ideen guddommelig Gyldighed. Den ved første Øiekast anstødelige
Dom: \emph{Alt er nødvendigt}, kan altsaa erkjendes
for den logisk overgribende, og den til Nødvendigheden svarende
apodiktiske Viden fastholdes som Logikens Ideal, uden
at den Indflydelse, der nødvendig tilkommer det Assertoriske
og det Problematiske, i fjerneste Maade lider derunder.

Er den articulerende Vexelbestemmelse af Assertorisk,
Problematisk og Apodiktisk ved Bestemmelsen af Dom og
Gjenstand tilfulde anerkjendt, saa indsees da ogsaa heraf, hvad
Identiteten af Dommens Indhold med Dommens Gjenstand,
seet under Nødvendighedens Synspunkt, oprindelig har at
betyde. I Kraft af Nødvendigheden er Dommens Gjenstand
identisk med dens Indhold; thi den apodiktiske Dom er i sit
Væsen logisk, og enhver Gjenstandsbestemmelse i den logisk
apodiktiske Dom forvandlet til Tankebestemmelse og derved til
Indholdsbestemmelse. Men i Kraft af Nødvendigheden er da
ogsaa Dommens Indhold identisk med dens Gjenstand; thi
den apodiktiske Dom er i Virkeligheden antilogisk, ja i Phænomenet
assertorisk, og enhver Indholdsbestemmelse en Reflex
af en tilsvarende Gjenstandsbestemmelse.

Det er i Consequens af den overgribende Nødvendighed
med dens Vexelbestemmelser af Logisk og Antilogisk, at den
altformaaende Subjectivitet ved kategorisk Magtsprog forvandler
en Verden af virkelige Gjenstande til en Verden af objective
% --- p. 292 ---
\opage{292}Domme. Hvorledes Begrebet Ting forholder sig til Begrebet
Egenskab, bestemmes ved en Tankesætning, altsaa ved en \emph{subjectiv}
Dom; Forholdet imellem den existerende Ting og de
existentielle Egenskaber er derimod oprindelig bestemt ved en
Magtsætning, en praktisk-antilogisk Sætning, altsaa ved en
\emph{objectiv} Dom.

Undersøgelsen af Forholdet mellem Dom og Gjenstand
har altsaa med indre Conseqvens forvandlet sig til en Undersøgelse
af Forholdet mellem Dom og Dom. Denne Undersøgelse
gaaer i tre Retninger; thi det gjælder nu: deels
at bestemme Forholdet mellem Dom og Dom i subjectiv Forstand,
altsaa at indsee de subjective Dommes Sammenhæng;
deels at bestemme Forholdet mellem Dom og Dom i objectiv
Forstand, altsaa at indsee de objective Dommes Sammenhæng;
og endelig at bestemme Forholdet mellem Dom og Dom i subjectiv-objectiv
Forstand, altsaa at indsee Dommenes oprindelige
Sammenhæng i Ideen.

At bestemme Forholdet mellem Dom og Dom er, hvad
Sammenhængen angaaer, at bestemme, hvilken Dom der er
den oprindelige, og hvilken der er den afledede, det er, med
andre Ord, at aflede Dom af Dom, altsaa at \emph{slutte}: alle
Slutningsfunctioner ere i logisk Forstand Sammenhængsfunctioner.
At Sammenhængskategorien danner Grundlaget,
skjønnes umiddelbart deraf, at den samme Dom kan, uden at
Noget forandres enten i Form eller Indhold, forekomme
baade som Præmis og Conclusion; Forskjellen mellem
Præmis og Conclusion afgjøres ved Sammenhængen. Heraf
sees tillige Forholdet mellem Modalitet og Sammenhæng; er
Dommen i en Sammenhæng Præmis, i en anden Conclusion,
da opfattes den jo i den første Sammenhæng paa en anden
Maade end i den sidste; som Sammenhængsfunctioner ere
Functionerne af Slutningen kategorisk bestemte Modalfunctioner.
Idet Alt her kommer an paa Sammenhængen, kommer Alt
med det Samme an paa Maaden, altsaa paa
% --- p. 293 ---
Sammenhængs\opage{293}maaden; thi Maaden bestemmer Sammenhængen og Sammenhængen
igjen Maaden.

I Modalfunctionernes kategoriske Eiendommelighed ligger
deres Inddelingsgrund. Til den subjective Slutning, Bestemmelsen
af de subjective Dommes Afhængighedsmaade, svare
Functioner af subjectiv Sammenhæng, \emph{Modalfunctioner
af Viden}; til den objective Slutning, Bestemmelsen af de objective
Dommes Afhængighedsmaade, svare Functioner af objectiv
Sammenhæng, \emph{Modalfunctioner af Magt}. Til
Ideeslutningen, Bestemmelsen af de subjectiv-objective Dommes
indbyrdes Afhængighed i Reflexion af det Absolute selv, svare
Functioner af subjectiv-objectiv Sammenhæng, \emph{Modalfunctioner
af Sandhed}.

\medskip
\noindent\textit{α)} \emph{Subjectiv Slutning: Modalfunctioner af Viden.}
\medskip

Til Overskuelse af de subjective Dommes indre Forbindelser
og særskilte Afhængighedsmaader er det fra Begyndelsen
af nødvendigt at gjøre sig Rede for Udgangspunktet, Maalet
og Methoden. Udgangspunktet, eller nøiagtigere Centralpunktet,
det Punkt, hvorom hele Udviklingen dreier sig, er
Subjectiviteten selv, den tænkende, dømmende Bevidsthed.
Maalet, den Grundopgave, til hvis Løsning Modalfunctionerne
i deres forskjellige Endringer samvirke, er en Articulation af
den vidende Subjectivitets Indhold. Da Indholdet, opfattet
som Tankeindhold, bestaaer i en Mangfoldighed af Domme,
kan den Vidende kun stemme overeens med sig selv og besidde
sit Indholds Rigdom i fuld Gjennemsigtighed, forsaavidt de
mangfoldige Domme umiddelbart eller middelbart forbindes
til Eenhed, eller dog forsaavidt de ved den eiendommelige
Sammenhængsmaade vise sig indbyrdes overeensstemmende.
Ved den saaledes bestemte Opgave er nu Fremgangsmaaden,
den logiske Methode, tillige angivet. Idet Sammenhængs- og
Afhængighedsmaaden kjendes af Slutningen, vil den subjective
Slutning naturligviis antage forskjellig Form, eftersom
% --- p. 294 ---
Af\opage{294}hængighedsmaaden bestemmes partielt ved Sammenhængen
mellem to, mellem tre, mellem flere, eller totalt ved Sammenhængen
mellem alle Domme. Bestemmes Afhængighedsmaaden
ved Sammenhængen mellem to Domme alene, er Slutningen
\emph{umiddelbar}; bestemmes den særlig ved Sammenhængen mellem
tre eller flere Domme, er Slutningen \emph{middelbar}; bestemmes
den endelig ved den oprindelige Sammenhæng, der forvandler
de i Forestillingen adskilte, for Anskuelsen umiddelbart opstillede
Domme til Momenter i et alsidigt, ved det Middelbares
og Umiddelbares gjennemsigtige Eenhed totaliseret System,
er Slutningen \emph{immanent}. Overgangen fra den umiddelbare
Slutning gjennem den middelbare til Immanensslutningen
udviser, som vi skulle see, en fremadskridende Bevægelse,
der er betinget af Modalfunctionernes tilsvarende Endringer.

\textit{αα)} \emph{Umiddelbar Slutning: analytiske Modalfunctioner.}
Ligesom man paa eensidig og vilkaarlig Maade
kan indskrænke den logiske Dom til Schemaet: „det Enkelte
er (subsumeres under) det Almindelige“, saaledes kan man da
ogsaa paa en ikke mindre eensidig og vilkaarlig Maade dictere
sig selv et Schema, hvorefter der til enhver logisk Slutning
altid maa kræves en Adskillelse af to Extremer og en Midte.
Naar et saadant Schema skal gjælde for eneste kanonisk, er
det unegtelig en Selvfølge, at der i Logiken ikke mere kan
være Tale om „umiddelbare Slutninger“. Hegel, som med
Hensyn paa dette Schema lader Dommen ved sig selv udvikle
sig til Slutning, saa at de to Extremer svare til Subject
og Prædicat ligesom Midten til Copula,\footnote{„Der Schluß hat sich als die Wiederherstellung des \emph{Begriffes} im
\emph{Urtheile}, und somit als die Einheit und Wahrheit beider ergeben.
Der Begriff als solcher hält seine Momente in der Einheit
aufgehoben; im Urtheil ist diese Einheit ein Innerliches, oder was
dasselbe ist, ein Aeußerliches, und die Momente sind zwar bezogen,
aber sie sind als \emph{selbstständige Extreme} gesetzt. Im \emph{Schlusse}
% note breaks here: continues at the foot of p. 295
sind die Begriffsbestimmungen wie die Extreme des Urtheils, zugleich
ist die bestimmte \emph{Einheit} derselben gesetzt.“ Hegel. Die
subjektive Logik. S. 118.}% note *)
har i den Grad
% --- p. 295 ---
\opage{295}overseet den gamle Lære om „umiddelbare Slutninger“, at
han ikke engang har villet spilde et Ord paa at faae den afviist.

Ikke desto mindre er Læren om de umiddelbare Slutninger,
naar vi see nærmere til, ligesaa uundværlig i Logiken,
som Læren om de middelbare, og det af den gode Grund, at
disse først faae deres Betydning ved hine. Begynde vi med
den middelbare Slutning: „Alle Mennesker ere dødelige; Cajus
er et Menneske; altsaa er Cajus dødelig“, da indsees det let,
at Slutningsfunctionen ikke kan indskrænke sig til at bestemme
Conclusum: „Cajus er dødelig“, efterat de to Præmisser i
Forveien ere opstillede. Ved at nedsætte Functionen til en
saa formel mechanisk Act miskjender man Slutningens logiske
Betydning. Det Logiske i Slutningen udfolder sig først, naar
Foreningen af Følgesætningens Subject og Prædicat udtrykkelig
er sat som Opgave, naar det med andre Ord er bragt til
Bevidsthed, at en i et bestemt Tilfælde given Sætning kun
har Gyldighed, forsaavidt den lader sig fremstille som en til
bestemte Forudsætninger svarende Følgesætning. Under dette
Synspunkt vil den middelbare Slutning fra alle Sider
fordre og forudsætte den umiddelbare. Sætningen: „Cajus
er dødelig“, skal bestemmes som Følgesætning, \emph{altsaa} maa
den have Præmisser; Sætningen: „alle Mennesker ere dødelige“,
er den ene Præmis; \emph{altsaa} maa Sætningen: „Cajus
er et Menneske,“ være den anden; omvendt! Sætningen:
„Cajus er et Menneske“ er den ene Præmis, \emph{altsaa} maa
Sætningen: „alle Mennesker ere dødelige,“ være den anden.
Det er lutter umiddelbare Slutninger; men hvis en eneste af
disse umiddelbare Slutninger faldt bort, hvad blev det saa til
med den middelbare? Man analysere hvilken middelbar
% --- p. 296 ---
Slut\opage{296}ning man vil: det skal, naar man blot ikke glemmer, hvad
og hvormeget der oprindelig hører med i Slutningsfunctionen,
nok vise sig, at det Middelbare i enhver Slutning forudsætter
det Umiddelbare.\footnote{At det ikke skal synes, som om Anpriisningen af den umiddelbare
Slutning var en tom Formalitet, og Slutningen selv uden Indflydelse
paa den videnskabelige Erkjendelses virkelige Fremskridt,
vælge vi strax et Exempel af den exacte Videnskab. Hvis en Mathematiker
begynder Løsningen af et Integrationsproblem saaledes:
det er Udtrykket: $\frac{x^{n}dx}{\sqrt{a+bx^{2}}}$, der skal integreres, \emph{altsaa} er det
Udtrykket: $\frac{x^{n-1}}{b}\cdot\frac{bxdx}{\sqrt{a+bx^{2}}}$, der skal integreres, da vil Enhver,
der besinder sig paa, at en Størrelse ved at multipliceres og divideres
med Ligestort bliver uforandret, vist indrømme, at her begyndes
med en umiddelbar Slutning. Men hvortil nu denne
Slutning? Ja, det seer man af, hvad der for den Sagkyndige
igjen viser sig som en umiddelbar Slutning: den ene Factor,
Udtrykket $\frac{bxdx}{\sqrt{a+bx^{2}}}$, er aabenbart Differentialet af $\sqrt{a+bx^{2}}$;
\emph{altsaa} kan man integrere deelviis o.~s.~v.}% note *)

At Opdageren af den immanente Methode ikke har havt
Øie for den umiddelbare Slutnings gjennemgaaende Betydning,
er jo mærkeligt nok, da han ellers har et saa skarpt
Blik for medierende Overgangsformer; imidlertid er just denne
Omstændighed ingenlunde uden charakteristisk Indflydelse paa
Methoden selv. Dersom Hegel ikke havde anseet saadanne
logiske Bevægelser som fra Væren og Intet gjennem Vorden
til Tilværen, fra Tilværen og Tilværelse gjennem Noget og
Andet til Et (For-sig-væren) o.~s.~v. for saa objective, at de
kunde foretages af Begrebet selv, uden begribende Subjectivitet,
saa vilde han fra Begyndelsen af have indseet, at disse Bevægelser
i Virkeligheden foregik derved, at Dom i uafbrudt Sammenhæng
fulgte paa Dom; men en Udviklingsrække, hvori den
% --- p. 297 ---
\opage{297}ene Dom, continuerlig følger paa den anden, vil tillige vise,
hvorledes den ene Dom ved det frie Tankeblik umiddelbart
følger af den anden; men at den ene Dom umiddelbart
følger af den anden, er jo netop et Udtryk for den umiddelbare
Slutning. Man kan i en hegelsk Logik vælge, hvilke
Exempler paa immanent Tankebevægelse man vil: det skal
ved nærmere Eftersyn nok vise sig, at Overgangene skee ved
umiddelbare Slutninger. „De Ting, som ere lige“, hedder
det i Heibergs „\emph{logiske System}“, „staae derved i Relation
til den fælles Maalestok og følgelig ogsaa til hinanden selv,
eftersom de ellers ikke kunde sammenlignes. Men da al Relation
er Forskjel, Endelighed, Afstand og Bevægelse, saa er
alt dette umiddelbart tilstede i det Oprindeliges absolute Eenhed
og Forskjel,“ o.~s.~v. Det er let at reducere den angivne
Tankebevægelse til en Række umiddelbare Slutninger. Tingene
ere lige, \emph{altsaa} maae de være sammenlignede; Tingene ere
sammenlignede, \emph{altsaa} maae de staae i Relation til den fælles
Maalestok; Tingene staae i Relation til den fælles Maalestok,
\emph{altsaa} staae de i Relation til hinanden selv. Al Relation
er Forskjel, Endelighed, Afstand og Bevægelse, \emph{altsaa} er Endelighed,
Afstand og Bevægelse umiddelbart tilstede i det Oprindeliges
Eenhed og Forskjel, o.~s.~v.

Man vil let overbevise sig om, at dersom der i en saadan
Fremstilling opdages et Punkt, hvor den umiddelbare
Slutning ikke slaaer til, da glipper det ogsaa med Forstaaelsen.
Finder Nogen den Slutning: „Tingene ere lige, altsaa maae
de være sammenlignede“, uholdbar, efterdi der jo, som det
synes, maa kunne gives mange baade lige og ulige Ting, der
endnu ikke ere blevne sammenlignede: da vil han ogsaa ansee
en Begrundelse som den, at de lige Ting maae staae i Relation
til hinanden, „eftersom de ellers ikke kunde sammenlignes,“
for utilstrækkelig.

Vistnok kan den Sammenhæng, hvori de enkelte Sætninger
forekomme, udbrede et Lys over hver især af dem, saa
% --- p. 298 ---
\opage{298}at Overgange, der ere forstaaelige i den flydende Fremstilling,
kunne blive uforstaaelige, naar Sætningerne parviis udrives
af Contexten saaledes, at den ene bliver Præmis, den anden
Conclusion. Af den flydende Fremstilling er det i vort Exempel
klart, at her ved „de Ting, som ere lige“, tænkes paa
Ting, hvis Lighed vi erkjende ved Sammenligning; hvorimod
den umiddelbare Slutning: „Tingene ere lige, altsaa maae de
være sammenlignede,“ har sit Anstødelige deri, at der i Præmissen
tales om objectiv Lighed, \emph{altsaa} --- en umiddelbar
Slutning --- om en Lighed, der ikke skjønnes at kunne beroe
paa Sammenligningens subjective Function. Men den
Omstændighed, at Contexten i det Hele virker med til Forstaaelsen
af det Enkelte, forandrer ikke Sagen; thi Indtrykket
af Contexten danner nu selv en Præmis, der har Forstaaelsen
af det Enkelte til Conclusion, altsaa igjen --- en umiddelbar
Slutning. „Contexten viser, at her tales om en tænkt
% PRINTER: the quotation „Contexten viser ... Lighed, is not closed before the interrupting dash; a fresh „ opens „om en Lighed ...“ (unbalanced as printed)
Lighed, \emph{altsaa} --- en umiddelbar Slutning ---, „om en Lighed,
der nødvendig forudsætter Sammenligning.“

Idet Contexten af et bestemt Sted danner Præmis for
en umiddelbar Slutning, tør Læseren ikke glemme, at den
saaledes dannede Præmis selv maa blive Følgesætning til en
mere omfattende Præmis, og at dette maa skee ved en ny
umiddelbar Slutning. „I en objectiv Logik kan der kun
være Tale om en objectiv Tænken, en objectiv Sammenligning;
\emph{altsaa} er der i denne --- under den objective Logik hørende ---
Context kun Tale om en objectiv Sammenlignen.“ Man
seer, hvorledes det i Forholdet mellem Lighed og Sammenligning
skjulte Problem ved den umiddelbare Slutning uophørlig
flytter sig fra et Standpunkt til det andet. Den umiddelbare
Slutning staaer i nøieste Sammenhæng med, hvad der
i logisk Forstand kaldes Blik, Indblik, Overblik; ja den er
selv Blikkets, Indblikkets og Overblikkets divinatoriske Fremgang
til middelbar Viden.

Spørges der, hvormange væsentlig forskjellige Former
% --- p. 299 ---
\opage{299}den umiddelbare Slutning vil kunne antage, da kommer det
fornemmelig an paa, hvorledes man i det Hele opfatter Forholdet
mellem Form og Indhold. Behandles Formen for sig
i abstract-mechanisk Modsætning til Indholdet, saa fremkommer
Schematismen; udvikles Formen i logisk Eenhed med
Indholdet, vil Formforskjellen beroe paa en Forskjel i Modalfunctionerne
selv; den saaledes udviklede Formforskjel leder os
saa til at indsee den umiddelbare Slutnings Overgang i den
middelbare.

1) \emph{Schematisme: de saakaldte umiddelbare Slutningers
Schema.} At schematisere er den formelle Logiks
stærke Side; Schemaet er dens Ideal, Schemaets Fuldendelse,
dens Blomst og Krone. Overalt, hvor den logiske Formalisme
med dens Inddelinger og Inddelingers Inddelinger ret
er trængt igjennem, er Schemaet foldet ud og forgrenet indtil
det Mikroskopiske. I ophøiet Stræben mod Schematismens
Fuldendelse har den formelle Logik da heller ikke kunnet indskrænke
sig til en Angivelse af „de umiddelbare Slutningers“
Hovedarter; den har ransaget Underarterne af hver Art for
sig, den er for Alvor trængt ind i det Enkelte.\footnote{Et Overblik over de saakaldte umiddelbare Slutningers Hovedarter
kunde jo vel yde en formel Tilfredsstillelse, naar man undgik alle
Vidtløftigheder og oplyste hver Art ved et Exempel. Et Exempel
som: alle Fugle lægge Æg, \emph{altsaa} nogle Fugle lægge Æg, vilde
da charakterisere Underordnings Slutningen (\textit{ratiocinatio subjectionis,
vulgo subalternationis}); et Exempel som: intet Dyr er moralsk;
\emph{altsaa}: ethvert Dyr er et ikke-moralsk Væsen, vilde passe
til Identitetsslutningen (\textit{ratiocinatio æqualitatis, æquipollentiæ,
conclusio ad æpuipollentem}); et Exempel som: alle Dyr have vilkaarlig
% PRINTER: printed "æpuipollentem", sense "æquipollentem"
Bevægelse, \emph{altsaa} er der intet Dyr, der ikke har vilkaarlig
Bevægelse, vilde være betegnende for Modsætningsslutningen (\textit{concl.
ad oppositam}), og et Exempel som: alle Mennesker ere dødelige,
\emph{altsaa} nogle Dødelige ere Mennesker, belyse Omvendings-Slutningen
(\textit{ratioc. per conversionem}) o.~s.~fr.

% note breaks here: continues at the foot of p. 300
Men herved kan den formelle Logik ikke blive staaende. Ifølge
dens scholastiske Princip er Schemaet en Hovedsag, og naar Schemaet
er en Hovedsag, saa gjælder det om Udholdenhed og Conseqvens:
med et Princip bør man ikke blive staaende paa Halvveien.
Det er at betragte som en i sin Art respectabel Conseqvens, at Redacteurerne
af den formelle Logik, efter at de en Tidlang have viist
den Svaghed at give efter for Tidens Fordringer og indskrænke
Schemaet, paany vedkjende sig det gamle Princip og udfolde Schematismen.
Vi skulle nu see, hvorledes her, hvad de umiddelbare
Slutninger angaaer, reflecteres, ikke blot paa Arterne, men (jvfr.
Drbal. Lehrbuch der propædeutischen Logik. Wien 1865. S. 55---71.)
ogsaa paa Underarterne.

Af Conversions-Slutninger f.~Ex. gives der nu to Arter:
Egentlige Conversions-Slutninger og Contrapositions-Slutninger.
Af egentlige Conversions-Slutninger gives der igjen 4 Underarter
1) Conversion af en almindelig bekræftende Dom (Alle Fugle
lægge Æg, altsaa: nogle æglæggende Væsener ere Fugle), 2) Conversion
af en almindelig benegtende Dom (Ingen Lediggjænger
veed at skatte Tiden, altsaa: Ingen, der veed at skatte Tiden, er en
Lediggjænger); 3) Conversion af en particulær bekræftende Dom
(Nogle Hunde ere troe Vægtere, altsaa: nogle troe Vægtere ere
Hunde): 4) Conversion af den particulært benegtende Dom. For
denne Conversion gives der ingen Regler. For at bevise, at der
ingen Regler gives, construerer man tre geometriske Figurer, hver
bestaaende af to Cirkler, som i første Figur skjære hinanden, i
tredie Figur falde udenfor hinanden, medens den større i anden
Figur omgiver den mindre. Betragter man nu første Figur, saa
ligge nogle under Sphæren \textit{S} faldende Gjenstande ikke i Sphæren
\textit{P}, nemlig \textit{x}, \textit{y} o.~s.~v.; ligeledes gives der Gjenstande hørende
under \textit{P}, som ikke ligge i \textit{S}, nemlig \textit{n}, \textit{z}\dots. Derhos gives der
Dele af \textit{S}, nemlig \textit{s}, der ligge i \textit{P}, og Dele af \textit{P}, nemlig \textit{p}, der
ligge i \textit{S}. Altsaa gives her fire Domme: 1. Nogle \textit{S} er \textit{P}; 2.
nogle \textit{S} ere ikke \textit{P}; 3. nogle \textit{P} ere \textit{S}; 4. nogle \textit{P} ere ikke \textit{S}.
Heraf følger nu Conversionen: nogle \textit{P} ere ikke \textit{S}. f.~Ex. Nogle
Pligter ere ikke Tvangspligter (ikke Gjenstand for Tvang), altsaa:
% note breaks here: continues at the foot of p. 301
Noget, der er Gjenstand for Tvang, er ikke Pligt. Nogle Exempler
ere ikke gode: altsaa noget Godt er ikke Exempel\dots. For
at faae noget saadant Poit ud, sætter man uhyre schematiske Vehikler
i Bevægelse, sytten Pompeværker for at tilveiebringe en eneste
Draabe Vand. Ogsaa her gjælder det: \textit{parturiunt montes, nascetur
ridiculus mus}.}% note *)
Det er den
% --- p. 300 ---
\opage{300}rene Iver for den rene Form; Ingen skal sige, at der i saadanne
Slutninger som: hvad \emph{alle} Fugle gjøre, det maae
% SEAM: p. 300 ends with the whole word "maae"; no word break. The footnote anchored on p. 299 runs over pp. 300--301 and is set here in full, so the p. 301 batch must omit its last six footnote lines ("Noget, der er Gjenstand for Tvang ... ridiculus mus.").
% --- p. 301 ---
\opage{301}\emph{nogle} Fugle gjøre; hvad \emph{ingen} Dyr ere, kunne \emph{nogle} Dyr ikke være; nogle Væsener, hvorom det kan siges, der kan siges om alle Mennesker, maae være Mennesker, o.\ s.\ v. --- lægges for eensidigt Eftertryk paa det logiske Indhold; derimod kan en passende Afbenyttelse af universelt bekræftende og universelt benegtende, particulært bekræftende og particulært benegtende Domme til passende Omsætninger, vel siges at gjøre en ypperlig Virkning, naar det gjælder om at faae Schemaet saa omfangsrigt og forgrenet som muligt.
% NOTE: the footnote lines at the foot of p. 301 („Noget, der er Gjenstand for Tvang ... ridiculus mus.“) continue the note anchored on p. 300; they belong to the previous batch.

Skulde der nu virkelig være Mening i en saadan Schematisering, da maatte Schemaet dog vel idetmindste kunne gjøre Tjeneste som et Slags Register eller Tabel. Men hvad logisk Tjeneste skulde det vel være? Som Register gjør Schemaet ingen Nytte. Trangen til et Register vilde forudsætte Vigtigheden af, at Slutningerne, for ret at forstaaes, dannede en vis Række og bleve opstillede i en bestemt Orden; men derpaa ligger her aldeles ingen Vægt. Den Slutning, at hvad \emph{alle} Fugle gjøre, det maae \emph{nogle} Fugle gjøre, er indlysende ved sig selv uden Hensyn til Række eller Register. Det er aldeles vilkaarligt og for den logiske Indsigt fuldkomment ligegyldigt, om man vil stille Equipollens-Slutningen foran Contradictions-Slutningen, denne igjen foran Conversionsslutningen, eller omvendt. Som Tabel gjør Schemaet da heller ingen Nytte. Hvorfor affatter man vel Tabeller? Naturligviis for at sikkre sig Resultater af møisommelige Iagttagelser, indviklede Analyser, vidtløftige Beregninger, o.\ dsl. Saaledes har man jo Stjernetabeller, Logarithmetabeller, chemiske Affinitetstabeller o.\ s.\ v.
% --- p. 302 ---
\opage{302}Men den formelle Logiks umiddelbare Slutninger ere ikke Resultater af vidtløftige Beregninger; man kan virkelig slutte, at hvad \emph{alle} Fugle gjøre, det maae \emph{nogle} Fugle gjøre, uden mange Iagttagelser eller vidtløftige Beregninger.

Fra alle Sider staaer den formelle Logiks Anstrengelse med Schemaet i skrigende Misforhold til det logiske Udbytte: Schemaet er vidtløftigt, det logiske Udbytte $=$ 0.

2) \emph{De analytiske Modalfunctioner.} Den Maade, hvorpaa den formelle Logik behandler de umiddelbare Slutninger, er i betænkelig Strid med dens eget Princip og bringer den tilsidst i Modsigelse med sig selv. Til en umiddelbar Slutning høre jo to særskilte Domme, af hvilke den ene danner Præmissen, den anden Følgesætningen. Men hvordan forholder nu denne Antagelse sig til Logikens øverste Grundsætning, det saakaldte Identitetsprincip? Er her virkelig to forskjellige Domme, saa kan den ene jo ikke, uden at være forskjellig fra sig selv, indeholde den anden. Men kan den ene af to Domme ikke indeholdes i den anden, hvorledes skulde den ene da kunne udledes af den anden? Er det derimod ikke to forskjellige Domme, men kun to identiske Sætninger, vi have for os, da give de to Sætninger jo kun een og samme Dom, udtrykt paa to forskjellige Maader; men hvor der kun er een Dom, kan der umulig være Tale om Slutning. At „Alle“ bestaaer af „Nogle“, at den universelt bekræftende Dom indeholder den particulært bekræftende, men udelukker baade den universelt og den particulært benegtende, er jo Noget, der i den Grad forstaaes af sig selv, at her ikke kan blive Plads for nogen Slutning. Det er saaledes i Conseqvens af den formelle Logiks eget Princip, at dens Lære om umiddelbare Slutninger omsider er blevet forkastet\footnote{Diese angeblichen Schlüsse sind in Wahrheit keine Schlüsse, sondern nur die Hervorhebung (Aussprache) dessen, was zwar nicht ausdrücklich, doch aber unmittelbar im Obersatze mitgesetzt ist, so daß % note breaks here (p. 302/303)
ich mir den Inhalt desselben nur klar zu machen brauche, um es auch in ihm zu finden. Wo der Mittelsatz fehlt, wo es also keiner Vermittelung zwischen der Prämisse und der Conclusion bedarf, da bedarf es offenbar auch keiner Folgerung, d.\ h.\ da findet die Function des Folgerns gar nicht statt, da ist auch überhaupt kein Schluß vorhanden. Mit andern Worten: die Function des Schließens und somit der Schluß als Ausdruck derselben besteht eben nur in der Ableitung der Conclusion aus der Prämisse \emph{mittelst} des Mittelsatzes: wo also der Schlußsatz \emph{nicht} abgeleitet, sondern nur ausgesprochen wird, daß sein Inhalt in der Prämisse mit gesetzt ist, da findet offenbar auch kein Schluß statt.“ Ulrici. Conpendium der Logik. Leipzig 1860. S.~192. Man lægge vel Mærke til, at Forkastelsen ligefrem er grundet paa en Opfattelse af Identitetsprincipet, hvori der i ingen Henseende afviges fra den formelle Logik.}% note *)
% PRINTER: quotation closed with “ after „statt.“ but no opening „ before „Diese“ (unbalanced as printed).
% PRINTER: printed "Conpendium", sense "Compendium".
.

% --- p. 303 ---
\opage{303}Fra den objective Logiks Standpunkt kan der imidlertid ogsaa, ihvorvel Forskjellen her udvikler sig af Identiteten, reises Indsigelse imod Slutninger af een Præmis. Man gjør fra dette Standpunkt gjældende, at her imellem Præmis og Følgesætning kun kan være en Reflexionsforskjel, og at en saadan Forskjel kun bestemmes ved en Reflexion, ikke ved en Slutning. Betragte vi Sagen abstract objectivt, da kan den immanente Overgang fra Væsen til Phænomen, fra Grund til Existens, fra Mulighed til Virkelighed, fra Aarsag til Virkning vistnok ikke fortjene Navn af Slutning; thi enhver af de adskilte Kategorier er, naar den fastholdes for sig, saa eensidig, at den, for at svare til sin Bestemmelse, maa reflectere sig i sit Complement. Den ved Reflexionen betegnede Fremgang er derfor ligesaameget en Tilbagegang; det er en Overgang, som dog ingen Overgang er. Idet Tankebevægelsen gaaer fremad fra Væsen til Phænomen, gaaer den netop fra Phænomen tilbage til Væsen. Her er altsaa ikke en særskilt Dom om Væsenet, hvoraf der skulde kunne udledes en egen Dom om Phænomenet, thi Dommen om Væsenet falder i den objective
% --- p. 304 ---
\opage{304}Reflexion nødvendig sammen med Dommen om Phænomenet; ligesaa lidt kan her tænkes paa en særskilt Dom om Phænomenet, hvoraf en egen Dom om Væsenet lod sig udlede, da Dommen om Phænomenet i objectiv Reflexion nødvendig maa falde sammen med Dommen om Væsenet.

Disse Indvendinger bortfalde, naar det erindres, hvad i en tidligere Sammenhæng er viist, at den objective Reflexion overalt maa beroe paa den subjective. Men det i Reflexionen Subjective beroer igjen --- formedelst Ordet --- paa Dommen. Det bliver da fra dette Synspunkt ikke blot muligt, men nødvendigt at fælde en særskilt Dom om Væsenet for at kunne fælde en særskilt Dom om Phænomenet. Her er for den subjective Reflexion to Spørgsmaal at besvare: hvad er Væsen? hvad er Phænomen? At disse Spørgsmaal maae besvares ved forskjellige Domme, er uimodsigeligt. Da nu den Dom, der fældes om Phænomenet, ligesaa vel er bestemmende for den, der fældes om Væsenet, som omvendt: saa vilde den dømmende Subjectivitet befinde sig i en uendelig Svæven, dersom den ene af disse Domme ikke udtrykkelig blev sat som den oprindelige, den anden som den afledede. Er f.\ Ex.\ Dommen om Phænomenet sat som den oprindelige og forsaavidt umiddelbare, saa maa Dommen om Væsenet rette sig derefter; den maa som Følgesætning ved Analyse udledes af sin Forudsætning. Forudsætningen er Præmis, Følgesætningen Conclusion; Phænomenet er opfattet \emph{saaledes}, \emph{altsaa} maa Væsenet opfattes \emph{saaledes}: her er en umiddelbar Slutning.

Hvad dette vil sige, maa vel falde i Øinene, naar Spørgsmaalet er, hvorvidt der kan sluttes fra Mulighed til Virkelighed og fra Virkelighed til Mulighed. Den formelle Logik er, ifølge det eengang Vedtagne, ligesaa sikker paa, at enhver Slutning fra Virkelighed til Mulighed maa være sand: \textit{ab esse ad posse esse valet consequentia}, som den er forvisset om, at enhver Slutning fra Mulighed til Virkelighed maa være falsk: \textit{a posse esse ad esse non valet conse%
% --- p. 305 ---
\opage{305}quentia}. At en Slutning fra Virkelighed til Mulighed imidlertid ligesaa vel kan være falsk (det Virkelige har været, men er ikke længere, et blot Muligt, efterdi Muligheden er afløst af Virkeligheden), som en Slutning fra Mulighed til Virkelighed kan være sand (det i objectiv Forstand Mulige maa blive virkeligt, thi ellers bliver det Mulige umuligt), er let at indsee. Alt beroer i de givne Tilfælde paa en særskilt Dom om Muligheden og en særskilt Dom om Virkeligheden.

Betydningen af den umiddelbare Slutning bliver endnu mere iøinefaldende, naar vi besinde os paa, at de objective Kategorier kun bestemme hinanden objectivt, idet de vise sig bestemmende for paagjældende Objecter. Først idet Objectet erkjendes at være bestemt ved sin Kategorie, bliver den reflecterende Subjectivitet dømmende; men i de objective Forhold maa den ene kategoriske Bestemmelse nødvendigt drage den anden efter sig, den ene Dom lede til den anden: her er en Virkning (Trætoppene bevæge sig), \emph{altsaa} maa her være en Aarsag (Blæsten); her er et Phænomen (Blegheden og de funklende Øine), \emph{altsaa} maa her være en Grund (stærk Sindsbevægelse); her ere bestemte Midler (Agerdyrkningsredskaber), \emph{altsaa} maa her være et bestemt Øiemed (Agerdyrkning). Ved at opfatte de umiddelbare Slutninger fra det objective, praktiske Synspunkt, vil man overbevise sig om, at de, langt fra at være overflødige, spille en vigtig Rolle ved Subjectivitetens Orientering i den virkelige Omverden. Hvad vilde den blot sandselige Seen og Høren være uden en tilsvarende Mangfoldighed af umiddelbare Slutninger? Eller hvorledes skulde man forklare sig den Kjendsgjerning, at der kan gives Sandsebedrag, optiske, akustiske Illusioner, dersom der ikke gaves umiddelbare Slutninger?

Da nu en umiddelbar Slutning i et givet Tilfælde ligesaa vel kan være falsk som sand, maa Logiken anføre et Kriterium, hvorved det Sande kan adskilles fra det Falske. De er Sammenhængen imellem de to Domme, Præmis og Følge%
% PRINTER: printed "De er Sammenhængen", sense "Det er".
% --- p. 306 ---
\opage{306}sætning, paa hvis Bestemmelse det væsentlig kommer an; man undersøge da, hvorvidt denne Sammenhang lader sig reducere til en hypothetisk Dom eller ikke! Kan denne Reduction foretages, saa er Slutningen paalidelig; hvis ikke, er den upaalidelig. „Dette er en plan Trekant, altsaa er Vinklernes Sum $=$ to Rette“; denne Slutning er sand, thi „\emph{dersom} det er en plan Trekant, \emph{saa maa} Summen af dens Vinkler nødvendig være to Rette.“ Slutter Nogen derimod saaledes: „Familien befandt sig vel imorges (Kl.~8), \emph{altsaa} befinder den sig vel endnu (Kl.~12 Middag)“: da indsees let, at en saadan Slutning ikke er grundet paa Nødvendighed, men kun paa Sandsynlighed og altsaa ikke kan være ganske paalidelig. I Theorien er dette kjendeligt derpaa, at Præmis og Conclusum ikke forholde sig som For- og Eftersætning i en hypothetisk Dom; i Praxis kjendes det derpaa, at Slutningen kun virker beroligende under almindelige Forhold, men ikke, naar der er Fare paa Færde.

De analytiske Modalfunctioner maae altsaa gaae ud paa, ikke blot at angive Forholdet imellem den umiddelbare Slutnings Præmis og Conclusum i Almindelighed, men tillige for hvert eiendommeligt Tilfælde at bestemme, hvorvidt Slutningens Grundlag er Nødvendighed eller Sandsynlighed, og hvis det er en Sandsynlighed, da at angive dens Grader. Skulde det i særlige Tilfælde vise sig, at Grundlaget ikke er værende, men svæver imellem Sandsynlighedens og Nødvendighedens Poler, da er det den dømmende Subjectivitets Opgave at drage en Grændselinie mellem \emph{Slutning} og \emph{Formodning}, en ofte meget fiin Grændselinie, der saavel i Theorie som i Praxis er af afgjørende Betydning for den fra Dom til Dom fremskridende Tænkning.

For at opfatte de analytiske Modalfunctioners Betydning i hele dens Omfang forestille man sig Subjectiviteten i Forhold, ikke blot til en Mangfoldighed af umiddelbare Phænomener, men som dømmende Subjectivitet til en Mangfoldighed
% --- p. 307 ---
\opage{307}af umiddelbare, logiske, antilogiske, kategoriske, assertoriske Domme. Hvordan skulde det nu være muligt at tilveiebringe Eenhed i en saadan Bevidsthed, dersom det ikke ved en analyserende Reflexion var muligt at aflede en Dom af en anden, dersom det, med andre Ord, ikke var muligt at gjøre umiddelbare Slutninger. Kun formedelst den umiddelbare Slutning kommer der væsentlig Sammenhang imellem Dom og Dom, kun formedelst den umiddelbare Slutning er der Continuitet i den fremadskridende Tænkning. En continuerlig fremadskridende Tænkning er en intellectuel Bevægelse; men en saadan Bevægelse kan jo gaae i høist forskjellige Retninger.

At her kan være Tale om en Fleerhed af analytiske Modalfunctioner, beroer da fornemmelig derpaa, at der ikke blot af enhver Dom kan udledes en anden Dom, men at der endog af een og samme Præmis kan uddrages væsentlig forskjellige Slutninger. „Cajus er et Menneske, \emph{altsaa} er Cajus dødelig; Cajus er et Menneske, \emph{altsaa} er Cajus udødelig; Cajus er et Menneske, \emph{altsaa} trængende til Søvn,“ o.\ s.\ v. Enhver af disse Slutninger faaer sin særegne Betydning ved den eiendommelige Retning, hvori Tanken bevæger sig: Physiologen slutter anderledes end Ethikeren, Ethikeren igjen anderledes end Æsthetikeren o.\ s.\ v. „Men“, kunde Nogen spørge, „er her nu ogsaa i logisk Forstand forskjellige Functioner? Det synes dog virkelig, som om Tanken fungerede paa samme Maade, enten her sluttes: \emph{altsaa} er C. dødelig, eller: \emph{altsaa} er C. udødelig; eller: \emph{altsaa} trænger C. til Søvn: forsaavidt det ene Altsaa er forskjelligt fra det andet, angaaer denne Forskjel jo kun Indholdet, ikke Formen.“ Vi gjøre opmærksom paa, at enhver væsentlig Indholdsforskjel altid maa reflectere sig i en tilsvarende Formforskjel. Overalt hvor Indholdsforskjellen unddrager sig Formforskjellen, synker Indholdet ned til at være blot og bart Stof. Er Forskjellen mellem et physiologisk, ethisk og æsthetisk Indhold meer end en Stofforskjel, da er Forskjellen mellem en physiologisk,
% --- p. 308 ---
\opage{308}ethisk og æsthetisk Slutning ogsaa en væsentlig Formforskjel. Forskjellen selv vedkommer ikke Logiken, men det, at her er en Forskjel, en væsentlig Formforskjel, vedkommer Logiken. At den formelle Logik ikke kjender den væsentlige Formforskjel, er forklarligt; den formelle Logik seer paa Resultatet og overseer Tankebevægelsen, den holder paa Schemaet og glemmer Formen. Concentreres hele Opmærksomheden paa det Conclusum, der følger paa den umiddelbare Slutnings \emph{Altsaa}, da er det ene Altsaa ganske vist ligesom det andet, her er i Functionen selv ingen mærkelig Forskjel. Men hvoraf kommer nu det? Det kommer deraf, at man overseer Tankebevægelsen, den Fremgang, hvori Slutningsacten selv betegner en Vending, og kun seer paa Resultatet, om f.\ Ex.\ Slutningssætningen er bekræftende eller benegtende, almindelig benegtende eller o.\ s.\ v., Forskjelsbestemmelser, der eensidig passe for Schemaet, den mechaniske, fra Indholdet løsrevne Form.

Concentreret i det \emph{Altsaa}, der ifølge Schemaet udelukkende støtter sig til den ene Præmis, er Tankebevægelsen reduceret til et forsvindende Minimum, et Element, hvori dens \emph{Hvorfra} og \emph{Hvorhen} er blevet aldeles ukjendeligt. Alle \textit{S} ere \textit{P}, \emph{altsaa}: nogle \textit{S} ere \textit{P}, her er --- under Synspunkt af en Slutning --- en uendelig lille Bevægelse fra, at Alle ere, til, at Nogle ere, en saa lille Bevægelse, at man endog kan stride, om den er fremad- eller tilbageskridende. Det forholder sig med Bevægelsesdifferentialerne i de forskjellige til den umiddelbare Slutning hørende \emph{Altsaa}, som med de forskjellige Differentialer i „det Uendeliges Analyse.“ Elementerne \textit{dx}, \textit{xdx}, \textit{adx}, \textit{cosxdx} o.\ dsl. kunne, naar de betragtes som tilværende, ansees for at være ligestore, efterdi de alle ere uendelig nærved 0; men naar de integreres, saa seer man \emph{Hvorfra} og \emph{Hvorhen}, saa seer man \emph{Bevægelsen}, og saa er Forskjellen gjennemgaaende. Saaledes med det logiske Bevægelsesdifferential, den i \emph{Altsaa} indesluttede Tanke; den faaer sin logiske Betydning ved den hele
% --- p. 309 ---
\opage{309}Tankebevægelse. I Kraft af Tankebevægelsen gjør det en væsentlig Forskjel, om det er et physiologisk, et ethisk eller æsthetisk \emph{Altsaa} vi have for os. Til Grundforskjellen i den umiddelbare Slutnings \emph{Altsaa} svarer Forskjellen i de analytiske Modalfunctioner.

3) \emph{Den umiddelbare Slutnings Overgang i den middelbare.} Det ligger i den umiddelbare Slutnings Begreb, at dens Følgesætning maa udledes af een Præmis; men det ligger ingenlunde i dens Begreb, at den i ethvert Tilfælde absolut skal være umiddelbart indlysende. At denne Fordring ikke kan opstilles som almeengjældende, er en ligefrem Følge af, at det, som for den Ene øieblikkelig er klart, kan for den Anden være dunkelt. Hvad En seer ved første Øiekast, er en Anden maaskee først efter mange Forberedelser og Forklaringer istand til at fatte\footnote{Denne Forskjel er ofte af stor komisk Virkning. I „Hovmesteren i Knibe“ kan Ane strax af den glemte Hat (altsaa ved en umiddelbar Slutning) see, at der har været en Dame i Hovmesterens Værelse; men det koster hende adskillig Anstrengelse, inden hun kan gjøre Phænomenet indlysende for Hans.}% note *)
. Talent for en særegen Videnskab beroer paa et særeget Blik, men det særegne Blik er jo netop kjendeligt paa den Lethed og Sikkerhed, hvormed den i en given Retning Begavede gjør umiddelbare Slutninger. Objectivt eller i og for sig kan Fordringen heller ikke gjælde; thi for det Første gives der ingen Slutning saa absolut objectiv, at den skulde kunne slutte sig selv; for det Andet gives der ingen Slutning saa absolut umiddelbar, at den, naar Reflexionen vækkes, jo strax maa gaae over i det Middelbare\footnote{Paa en egen naiv Maade kommer Middelbarheden i de umiddelbare Slutninger da ogsaa tilsyne derved, at Man for enhver saadan Slutning udtrykkelig kræver Beviis. Man beviser --- tænk engang! --- Gyldigheden af en Slutning som denne: det er sandt, at alle Dyr have vilkaarlig Bevægelse, altsaa er det falsk, at nogle % note breaks here (p. 309/310)
Dyr ikke have vilkaarlig Bevægelse. Beviset lyder (Drbal. Lehrbuch der propädeutischen Logik. S.~59) saaledes: „Naar Dommen: alle Dyr have vilkaarlig Bevægelse er sand, saa maa Dommen nogle Dyr have vilkaarlig Bevægelse, ogsaa være sand. Antages nu, at Dommen: nogle Dyr have ikke vilkaarlig Bevægelse, ligeledes var sand, saa vilde, efterdi Dommene have \emph{samme særlige} Subject, Sætningerne: nogle Dyr have vilkaarlige Bevægelser og netop \emph{disse samme} (Dyr) have ingen vilkaarlig Bevægelse, paa een Gang være sande, hvilket strider imod Modsigelsens Grundsætning.“ Hvor middelbar den umiddelbare Slutning ved en saadan Beviisførelse maa blive, sees endvidere deraf, at Beviset for dens Gyldighed udtrykkelig grunder sig paa Beviset for Gyldigheden af Underordnings-Slutningen. Dette Beviis lyder saaledes: „Naar den universelt bekræftende Dom: Alle \textit{S} ere \textit{P}, er gyldig, saa maa ogsaa den particulært bekræftende Dom: nogle \textit{S} ere \textit{P}, være gyldig, thi det Modsatte vilde, hedder det, være „urimeligt“.% PRINTER: the quotation opened at „Naar den universelt ... has no closing mark of its own; only „urimeligt“ is closed (as printed).
{} At det nu med alle disse Beviisførelser er Fias, maa indrømmes, thi den, der ikke ved umiddelbart Blik kan see Modsigelsen i, at nogle Dyr ikke skulde have vilkaarlig Bevægelse, naar det først er givet, at alle Dyr have vilkaarlig Bevægelse, er en complet Idiot; men en complet Idiot bliver ikke klogere af, at man henviser ham til Modsigelsens Grundsætning.\par Til Vurdering af den formelle Logiks formelle Uklarhed er Beviisførelsen forsaavidt charakteristisk, som det herved bliver aabenbart, at den ikke har gjort sig tilvørligt% PRINTER: printed "tilvørligt", sense "tilbørligt".
{} Rede for den umiddelbare Slutnings Begreb. Skal den umiddelbare Slutning kjendes paa det umiddelbart Indlysende: hvor vil Logiken da hen med sine Beviser? Skal Slutningens Umiddelbarhed kjendes paa, at den kun har een Præmis, hvorfor indskrænkes den da til intetsigende Former og Exempler?}% note **)
. „Linien fra \textit{a} til \textit{b} er ret, altsaa er
% --- p. 310 ---
\opage{310}den den korteste Vei imellem \textit{a} og \textit{b}“. Er denne Slutning nu ved sig selv indlysende eller er den det ikke? Subjectivt seet, gjælder, som i det Foregaaende er viist, begge Opfattelser. Slutningen er umiddelbart indlysende, thi den betegner en Grundsætning; dog nei, Slutningen er ikke umiddelbart indlysende, thi den foregivne Grundsætning er egentlig en Lære%
% --- p. 311 ---
\opage{311}sætning, og en Læresætning kræver Beviis. Objectivt seet, er Slutningen derimod ubetinget umiddelbar; thi hvorpaa beroer vel dens Umiddelbarhed? Ene og alene derpaa, at dens Følgesætning udledes af en eneste Præmis.

Heraf følger saa igjen, at den umiddelbare Slutning kan ikke alene kræve en rig og mangesidig Analyse, men ogsaa have et rigt videnskabeligt Indhold. Slutninger som disse: „Solens og Stjernernes daglige Bevægelse er kun tilsyneladende, \emph{altsaa} maa det være Jordens Bevægelse, der er virkelig“; „en Steen, der udkastes fra et høit Taarn, afviger mod Øst, \emph{altsaa} bevæger Jorden sig fra Vest mod Øst“, kræve til deres rette Forstaaelse unegtelig en Mængde Reflexioner og Mellembestemmelser; men derfor ere de, ifølge det fastsatte Begreb, ligefuldt umiddelbare. At Conclusum beroer paa en eneste Præmis, gjør Slutningen umiddelbar.

Den umiddelbare Slutning er en endelig Form, hvortil et uendeligt Indhold lader sig reducere. Det Resultat, der er Videnskabens sidste og rigeste, maa, saasandt det fremkommer ved en Conclusion, jo ligesaa vel kunne henføres til en eneste Præmis som det, der uddrages af den allerførste og simpleste Grundsætning. Vel maa Præmissen for den sidste Conclusion, naar den analyseres, indeholde en Mangfoldighed af Præmisser; men Fordringen er, at de mange Præmisser med deres Conclusioner reduceres til een Præmis. Det er en stedse rigere Reduction, som, idet den giver Præmissen et stedse rigere Indhold, afgiver Grundlaget for Modalfunctionernes analytiske Virksomhed. Dersom store indholdsrige Vidensresultater ikke kunde concentreres i umiddelbare Slutninger, vilde en sand og omfattende Videnskab være umulig.

Skulde hver enkelt af de i den indholdsrige Præmis liggende Præmisser hver Gang enkeltviis drages frem, gjennemgaaes og forklares, da maatte Videnskaben jo ved Behandlingen af ethvert nyt Punkt idelig begynde forfra, og kunde paa denne Maade aldrig komme til Ende med sine
% --- p. 312 ---
\opage{312}Undersøgelser. Man vilde f.\ Ex.\ neppe drive det ret vidt i Geometrien, dersom man ikke kunde støtte sig til saadanne umiddelbare Slutninger som: „Trekanten er retvinklet, altsaa er Qvadratet paa Hypothenusen liig Summen af begge Catheternes Qvadrater;“ „Figuren er en Cirkel med Radius $=$~\textit{r}, \emph{altsaa} er dens Areal $= \pi r^2$“; „\textit{A} er hele Kuglefladens Areal, altsaa er \textit{A} liigt Productet af Diameteren og Storcirklens Peripherie“, o.\ s.\ v.

Da nu Præmissen i en umiddelbar Slutning kan indeholde en saa stor Mangfoldighed af middelbart vundne Resultater som man vil, saa er Overgangen fra den umiddelbare til den middelbare Slutning med det Samme motiveret: den umiddelbare Slutning er i Grunden middelbar. Skulde der kunne gives en Slutning saa umiddelbar, at den ikke lod sig opløse i en middelbar, da maatte dens Præmis være en absolut, ved sig selv indlysende Grundsætning; men at slige Grundsætninger ikke gives, er viist i det Foregaaende.

Den umiddelbare Slutning er ikke blot den middelbares Forudsætning, men tillige dens Resultat, og kan som Resultat igjen opløses i sin Forudsætning. En Slutning som: „Cajus er et Menneske, altsaa dødelig“, hviler aabenbart paa den stiltiende Forudsætning, at ethvert Menneske er dødeligt; men naar denne Forudsætning udtrykkelig fremhæves, faaer Slutningen jo to Præmisser: ethvert Menneske er dødeligt, Cajus er et Menneske, altsaa --- her er en middelbar Slutning.

\textit{ββ)} \emph{Den middelbare Slutning: synthetiske Modalfunctioner.} Naar den analytiske Modalfunction, hvad Slutningsleddene angaaer, holder paa Identiteten, saa holder den synthetiske derimod paa Forskjellen. Kun ved en egen Samvirken af Tænken og Forestillen kan her tilveiebringes en Ligevægt af Analyse og Synthese, der gjør den middelbare Slutning mulig. Giver man i det ofte anførte Exempel: „Alle Mennesker ere dødelige, Cajus er et Menneske, altsaa er Cajus dødelig,“ den analyserende Tænkning eensidig
% --- p. 313 ---
\opage{313}Overvægt over den af Forestillingen betingede Synthese, saa taber Slutningen sin Betydning. Dette har Hegel gjort, og har desaarsag ikke kunnet anerkjende Slutningen. „Oversætningen“, saa lyder hans Indvending\footnote{Die subjektive Logik. S.~151.}% note *)
, „er kun rigtig, forsaavidt \emph{Følgesætningen er rigtig}; dersom Cajus tilfældigviis ikke var dødelig, saa vilde Oversætningen ikke være rigtig. Den Sætning, der skulde være Følgesætning, maa allerede umiddelbart for sig være rigtig, da Oversætningen ellers ikke kunde indbefatte alle Enkelte; før% UNSURE: før or for (print faint; sense requires før)
{} Oversætningen kan gjælde for rigtig, maa man i Forveien spørge, om ikke hiin Følgesætning selv er en \emph{Instans} imod den“.

Af Indvendingen sees, at Hegel lader den analytiske Modalfunction absorbere den synthetiske og derved ophæve sig selv. Er Cajus et Menneske, saa maa man for at kunne sige, hvad ethvert Menneske er, forud kunne sige, hvad Cajus er. Paa denne Maade reduceres da Identiteten af, hvad Cajus er, og hvad ethvert Menneske er, til en tautologisk Omskrivning: „Cajus er, hvad ethvert Menneske er, naar ethvert Menneske er, hvad Cajus er.“

Denne Formalisme ophører imidlertid, naar den endelige Tænkning bringes i det rette Forhold til Anskuelse og Forestilling. Sættes Individet Cajus imod Begrebet Menneske, saa at Individet som Subject er for Forestillingen, hvad Begrebet som Prædicat er for Tanken: da indsees, at Subjectet, det enkelte Menneske, maa have de almeenmenneskelige Egenskaber paa individuel Maade. Den endelige, af Forestillingen betingede Tænkning opfatter nu dette saaledes, at Individets Egenskaber falde i to Systemer: Systemer af de almeenmenneskelige og Systemer af de individuelle; idet de individuelle Egenskaber da antages at gjælde enten udelukkende for \emph{denne} eller i al Fald kun for \emph{nogle} Enkelte.

% --- p. 314 ---
\opage{314}Først naar Analysen er foretaget paa denne Maade, er her Plads for en tilsvarende Synthese.

Ifølge Adskillelsen af almindelige og særlige Egenskaber ved hvert enkelt Menneske deler Totalerkjendelsen sig da i to partielle Erkjendelser. Paa den ene Side er en for alle Mennesker fælles Egenskab opfattet og bestemt. Denne almeengjældende Egenskab sees under Tænkningens Synspunkt at ligge i Menneskets Væsen; det er forsaavidt Kategorien, Begrebet som Begreb, hvorpaa Almeengyldigheden grunder sig. „Mennesket er --- det ligger i den menneskelige Natur --- dødeligt.“ Men den endelige Tænkning er særlig betinget af Forestillingen; idet den almeengyldige Egenskab fremhæves, fæster Blikket sig ikke paa det rene Begreb, men paa de Individer, hvori Begrebet er virkeliggjort, altsaa ikke paa Mennesket, men paa Menneskene. Saaledes have vi i Oversætningen en Erkjendelse for sig: „\emph{Alle Mennesker ere dødelige}“. Hvorledes komme vi nu til Undersætningen? Aabenbart ikke ved en fortsat Analyse, men ved en Synthese. I den Erkjendelse, at alle Mennesker ere dødelige, er den Erkjendelse, at Cajus er et Menneske, ingenlunde indbefattet. Cajus kunde jo f.\ Ex.\ være enten en Hest eller en Hund eller en Elephant, ja der er i Logiken endogsaa Intet til Hinder for, at Cajus, om man saa vil, kunde være en Stokfisk.

Vil man vende Forholdet om og sige, at det egentlig er den i Undersætningen udtalte Erkjendelse, der indbefatter Erkjendelsen af Oversætningen, da er ogsaa dette en Misforstaaelse. Vistnok er det umuligt at erkjende Cajus for et Menneske uden at see ham deelagtiggjort i visse for Menneskenaturen almeengjældende Egenskaber. Men heraf følger ingenlunde, at den forestillende Tænkning skulde behøve at gjennemreflectere alle Menneskenaturens Egenskaber for at opfatte det enkelte Individ som et Menneske. Tvertimod! her forudsættes udtrykkelig, at Opfattelsen skeer paa den velbekjendte umiddelbare Maade, ved nemlig at anskue den men%
% --- p. 315 ---
\opage{315}neskelige Skikkelse. Undersætningen er i den particulære Erkjendelse altsaa ligesaa uafhængig af Oversætningen som denne af hiin. Idet de to særskilte Erkjendelser sammenbringes, faaer Slutningen sin Betydning, og Modalfunctionerne vise sig at være synthetiske.

De i den umiddelbare Slutnings Præmisser udtalte Erkjendelser ere med al deres respective Uafhængighed dog ingenlunde hinanden uvedkommende. Tvertimod! idet de sammenbringes, altsaa, idet Synthesen foretages, viser det sig, at de ved deres Forening danne en ny Erkjendelse. Følgesætningen, den nye Erkjendelse, kan paa forskjellig Maade være begrundet i Præmisserne; hvorledes Begrundelsen i de særlige Tilfælde nærmere er beskaffen, maa Analysen afgjøre. Det kommer nemlig an paa, hvorledes Slutningselementerne, (\textit{syllogismi elementa}),% UNSURE: closing parenthesis faint or absent in the print; text layer has it
{} de Led, som i Præmisser og Følgesætning udgjøre Subjecter og Prædicater, stille sig til hinanden. Lægges der særlig Vægt paa de af Elementernes vexlende Stilling betingede Formler, have vi \emph{Syllogistiken}; spørges her om de oprindelige Grundbestemmelser, da kommer det væsentlig an paa en \emph{Reduction} af de \emph{syllogistiske Former}; opfattes endelig de særskilte Præmisser som formgivende Momenter af et i Slutningssystemer sig continuerlig udviklende Indhold, er \emph{Syllogistikens Opløsning} logisk motiveret.

\emph{Syllogistik.} Den middelbare Slutning er efter sit Begreb den Tænkningsact, der begrunder en Doms Gyldighed paa en Forening af to eller flere andre Domme, og betegnes særlig med det aristotelisk-scholastiske Navn \emph{Syllogisme}. Syllogistiken eller Læren om Syllogismerne og deres Anvendelse begynder da schematisk med en Adskillelse af den \emph{enkelte} Syllogisme (\textit{syllogismus simplex}), som kun har to Præmisser, og den \emph{sammensatte} (\textit{s.\ compositus}), der danner sig ved Forbindelse af en Række enkelte Syllogismer (Polysyllogisme). Den enkelte, ikke sammensatte Syllogisme
% --- p. 316 ---
\opage{316}adskilles igjen i den \emph{kategoriske}, hvis Præmisser begge ere kategoriske Domme, den \emph{hypothetiske}, hvis ene Præmis idetmindste maa være en hypothetisk Dom, og den \emph{disjunctive}, hvis ene Præmis idetmindste maa være disjunctiv. --- Schematiseringen er atter i fuld Gang\footnote{Det er en rig Høst, den scholastiske Syllogistik tilbyder de Skarpsindige, hvis Tænkningsmechanisme særlig fremhjælpes ved Rubriker og Schemata. Betegne vi, i den kategoriske Syllogisme, Følgesætningens Subject ved \textit{S}, dens Prædicat ved \textit{P} og Middelordet ved \textit{M}, saa giver følgende Schema os de fire Figurer:
\begin{center}\begin{tabular}{c|c|c|c}
Første Figur & Anden Figur & Tredie Figur & Fjerde Figur\\
\textit{M}---\textit{P} & \textit{P}---\textit{M} & \textit{M}---\textit{P} & \textit{P}---\textit{M}\\
\textit{S}---\textit{M} & \textit{S}---\textit{M} & \textit{M}---\textit{S} & \textit{M}---\textit{S}\\ \hline
\textit{S}---\textit{P} & \textit{S}---\textit{P} & \textit{S}---\textit{P} & \textit{S}---\textit{P}\\
\end{tabular}\end{center}
Lægger man nu desuden særlig Vægt paa Dommenes Inddeling i universelt bekræftende (\textit{a}) og particulært bekræftende (\textit{i}), efter de to første Vocaler i Verbet \textit{affirmo}, i universelt benegtende (\textit{e}) og particulært benegtende (\textit{o}), efter Vocalerne i Verbet \textit{nego}, saa faaer man, idet man lader enhver af disse 4 Arter af Domme i enhver af de 4 Figurer vexelviis afgive Over- og Undersætning, 64 mulige Combinationer; naar man saa derhos anvender de tre gyldne Regler: \textit{ex mere negativis nihil sequitur} (intet \textit{M} er \textit{P}, intet \textit{S} er \textit{M} giver ingen Slutning), \textit{ex mere particularibus nihil sequitur} (nogle \textit{M} ere \textit{P}, nogle \textit{S} ere \textit{M} giver ingen Slutning) og \textit{conclusio sequi debet partem debiliorem prræmissarum} % PRINTER: printed "prræmissarum", sense "præmissarum"
(naar een af Præmisserne er negativ, bliver Følgesætningen negativ; er den ene Præmis particulær, bliver Følgesætningen particulær), vil man finde, at kun de 19 Combinationer give sande Slutninger. Disse 19 Combinationer anvendes da i de 4 Figurer saaledes, at hver Figur derved erholder sine gyldige Modi. --- Saaledes er f.\ Ex.\ første Modus i første Figur \textit{Barbara: M (a) P, S (a) M, ergo: S (a) P} (\emph{alt} Immaterielt er uforgjængeligt, \emph{alle} absolut enkelte Væsener ere immaterielle, altsaa: \emph{alle} absolut enkelte Væsener ere uforgjængelige). Tredie Modus i første Figur er \textit{Darii: M (a) P, S (i) M, ergo: S (i) P} (\emph{alle} Syrer farve blaae Plantestoffer \emph{røde}, \emph{nogle} Ilter ere Syrer, altsaa: \emph{nogle} Ilter farve % note breaks here: p. 316 / p. 317
blaae Plantestoffer \emph{røde}). I tredie Figur hedder derimod den fjerde Modus ikke \textit{Darii}, men \textit{Datisi}, en Bemærkning, der, som man seer, er meget vigtig: \textit{M (a) P, M (i) S, ergo: S (i) P} (\emph{alle} Svampe ere Kryptogamer, \emph{nogle} Svampe snylte paa andre Organismer, altsaa: \emph{nogle} paa andre Organismer Snyltende ere Kryptogamer).

Har man først sagt \textit{A} --- det er en anerkjendt Talemaade --- saa maa man ogsaa sige \textit{B}. Vil man virkelig lægge videnskabelig Vægt paa de tre aristoteliske Figurer, da er det ganske i sin Orden, at man ogsaa fordrer Respect for den fjerde, den galeniske. Hvorfor skulde den fjerde Figur ikke, som Figur betragtet, være ligesaa respectabel som de tre andre? Vel sandt, den fjerde Figur har ikke fuldt saamange Modi som den tredie, thi den tredie Figur har sex Modi: \textit{Darapti, Felapton, Disamis, Datisi, Bocardo, Ferison}; men den har dog flere end enten den første eller den anden. Disse have nemlig hver kun fire Modi, den første: \textit{Barbara, Celarent, Darii, Ferio}, den anden: \textit{Cesare, Camestres, Festino, Baroco}, men den fjerde har fem gyldige, for „Videnskaben“ meget vigtige Modi, nemlig: \textit{Bamalip, Calemes, Dimatis, Fesapo, Fresison}. Altsaa: skal den fjerde Figur bort, da bort med dem alle! men er man istand til at holde paa den fjerde Figur, er man ogsaa istand til at holde paa dem alle.

Kun ved at holde paa alle fire Figurer vil man ret kunne fatte den dybe Logik i det sindrige Vers, der skriver sig fra Petrus Hispanus (han blev, uden at vi just tør paastaae, at det var for Versets Skyld, siden Pave under Navn af Johan XXI, og døde 1277); Verset lyder saaledes:
\par\medskip
\hspace*{2em}\textit{Barbara, Celarent primæ, Darii, Ferioque.}\\
\hspace*{2em}\textit{Cesare, Camestres, Festino, Baroco secundæ.}\\
\hspace*{2em}\textit{Tertia grande sonans recitat Darapti, Felapton,}\\
\hspace*{2em}\textit{Disamis, Datisi, Bocardo, Ferison. Quartæ}\\
\hspace*{2em}\textit{Sunt Bamalip, Calemes, Dimatis, Fesapo, Fresison.}
\par\medskip\noindent
Lægger man nu Mærke til --- thi ved et saa skarpsindigt Vers er der Meget at lægge Mærke til og saare Meget at bemærke --- at ikke alene Vocalerne \textit{a, e, i, o}, men ogsaa Begyndelsesbogstaverne \textit{B, C, D, F}, ere betydningsfulde, idet de minde om det Analoge % note breaks here: p. 317 / p. 318
i de under forskjellige Figurer hørende Slutningsarter (\textit{Bamalip} f.\ Ex.\ analog med \textit{Barbara}), saa beaandes man af den logiske „Samvittighed“; og bliver man saa desuden vaer --- og her er Meget at blive vaer --- at endogsaa Consonanterne \textit{s, p, c, m} erholde en dybere Mening derved, at de ifølge Verset:
\par\medskip
\hspace*{2em}„\textit{S vult simpliciter verti, P verte per accid.}\\
\hspace*{2em}\textit{M vult transponi, C per impossibile duci;}“
\par\medskip\noindent
angive, hvilke Forandringer der maa foretages med enhver af Slutningsarterne i de tre sidste Figurer, naar de skulle reduceres til første Figur, saa føler man „\textit{testimonium spiritus}“.}% note *)
.

% --- p. 317 ---
\opage{317}Ihvorvel den udførlige Schematisering af de opstillede Formers Mangfoldighed og indbyrdes Forskjellighed ogsaa,
% --- p. 318 ---
\opage{318}hvad de middelbare Slutninger angaaer, viser sig at staae i Misforhold til det logiske Udbytte, er her dog paa Grund af Synthesen en Fleerhed af Mellembestemmelser og Mellemled, der gjør det forklarligt, at den paa sine Love reflecterende Tænkning kan finde sig tilskyndet til at søge en udtømmende Bestemmelse af Tilfælde og Regler.

Allerede den kategoriske Syllogismes tre Elementer: \emph{Middelordet} (\textit{terminus medius}), det fælles Tredie, hvortil Følgesætningens Subject og Prædicat, hvert for sig, maae henføres, \emph{Underbegrebet} (\textit{terminus minor}), Følgesætningens Subject, og \emph{Overbegrebet} (\textit{terminus major}), Følgesætningens Prædicat, maae frembyde Combinationspunkter og vække Interesse for en Schematisering. Lægger man nu Mærke til, hvorledes \emph{Oversætningen} (\textit{propositio major}), den Præmis, hvori Overbegrebet er forbundet med Middelordet, og \emph{Undersætningen} (\textit{p.\ minor}), den Præmis, hvori Underbegrebet er forenet med Middelordet, dannes i forskjellige Tilfælde: da opdager man let, ot Middelordet kan skifte Plads i Præmisserne og derved give disse en forskjellig Charakteer. % PRINTER: printed "ot", sense "at"
I Slutningen: „alle \emph{Mennesker} ere dødelige, Cajus er et \emph{Menneske}, altsaa er Cajus dodelig“, er Middelordet Subject i Oversætningen, men Prædicat i Undersætningen; i Slutningen: „ingen Rectangel er \emph{ligesidet}, alle Qvadrater ere \emph{ligesidede}, altsaa er intet Qvadrat en Rectangel“, er Mid%
% PRINTER: printed "dodelig" (no stroke on the o at 400 dpi), sense "dødelig"
% --- p. 319 ---
\opage{319}delordet Prædicat i begge Præmisser. Det Spørgsmaal opstaaer da, hvormange Tilfælde er her egentlig, og hvad kunne disse forskjellige Tilfælde vel have at betyde? Har man først udfundet, at der ved Middelordets forskjellige Stilling kan dannes fire syllogistiske Figurer, da ville de under enhver af Figurerne hørende Tilfælde, der snart vise sig skikkede snart igjen uskikkede til en adæqvat Syllogisme, fremkalde Spørgsmaal om de forskjellige Maader, hvorpaa der kan sluttes i hver enkelt Figur. Schematismen viser forsaavidt en Stræben efter at sikkre sig Tilfældene og henføre det Specielle under det Almene; men Bestræbelsen er ufrugtbar, Bevægelsen illusorisk; istedenfor at fordybe sig i sit Væsen driver Tanken et combinatorisk Spil med de paa Overfladen vexlende Former.

2) \emph{Deduction af syllogistiske Former.} Det gjælder om den middelbare Slutning, hvad allerede er sagt om den umiddelbare: Formudviklingen maa gaae i Eet med den logiske Indholdsudvikling. Men Formudviklingen, en Deduction af Slutningens særskilte Former, er betinget af en tilsvarende \emph{Reduction}, en Tilbageføren af det Mangfoldige til Begyndelsen selv, det eenfoldige Første.

Blandt de Logikere, der i den nyere Tid have lagt an paa en fornuftig Reduction af den middelbare Slutnings særskilte Former, maae vi særlig nævne Ulrici\footnote{Compendium der Logik. Leipzig 1860. S.~183--194.}% note *)
. „I det Almenes Forhold til det Enkelte“, siger han, „i Slægtens Forhold til Arten, ligesom Artens til dens Exemplarer, altsaa i Dommens Begreb ligger det umiddelbart, at hvad der gjælder (prædiceres) om det Almene, ogsaa maa gjælde om det under det Almene indbefattede Enkelte, og at hvad der gjælder om Slægten, ogsaa maa gjælde om dens Arter, Underarter og Exemplarer.“ I Consequens af denne Forudsætning reducerer Ulrici Syllogismens Nødvendighed til de første logiske Grundsætninger: „\textit{A} = \textit{A}, om Lige gjælder Lige, hvad der
% --- p. 320 ---
\opage{320}gjælder om Slægten, maa gjælde om dens Arter og disses Exemplarer; dette er den almene Grundsætning for alle \emph{positive} Slutninger; endvidere: \textit{A} er ikke = \textit{non A}, om Ligt gjælder ikke Uligt, og om Uligt ikke Ligt, hvad der ikke gjælder om Slægten, kan ikke gjælde om dens Arter og deres Exemplarer; dette er den almene Grundsætning for alle \emph{negative} Slutninger“.

„Af den almene Grundsætning for alle positive Slutninger følger: 1) at naar en heel Art subsumeres under et almindeligt Prædicats- eller høiere Slægtsbegreb, saa subsumeres ogsaa ethvert af Artens Exemplarer under samme Begreb. Med de efter denne Norm dannede Slutninger er første saakaldte Slutningsfigur implicite given. Af den samme almene Grundsætning følger: 2) at et Subject, der er at subsumere under et høiere Slægtsbegreb, nødvendig maa være at subsumere under et eller andet af de i samme indbefattede Artsbegreber, en Grundsætning, der hører til den anden Slutningsfigur, hvori den især finder negativ Anvendelse. Af Grundsætningen følger endvidere: 3) at naar en Art subsumeres (lader sig subsumere) under to forskjellige Slægter, maa ogsaa enhver af disse tildeels være at subsumere under den anden. De efter denne Norm dannede Slutninger høre til tredie Figur. Endelig er det: 4) en Consequens af den almindelige Grundsætning, at naar et Artsbegreb er at subsumere under et Slægtsbegreb, og dette igjen under et høiere Slægtsbegreb, saa maae ogsaa nogle til sidste Begreb hørende Exemplarer være at subsumere under Artsbegrebet; og af Grundsætningen for alle negative Slutninger, at naar Slægtsbegrebet, under hvilket Artsbegrebet subsumeres, ikke hører under et høiere Slægtsbegreb, intet Exemplar af dette sidste kan subsumeres under Artsbegrebet; samt at, naar Arten ikke subsumeres under det lavere Slægtsbegreb, maae ogsaa nogle Exemplarer af det høiere Slægtsbegreb være udelukkede fra
% --- p. 321 ---
\opage{321}Artsbegrebet. Det er de efter denne Norm dannede Slutninger, der ere betegnede ved fjerde Figur\footnote{Under sit hegelske Synspunkt har Heiberg paa en meget simpel og naturlig Maade deduceret de tre Figurer med Forkastelse af den fjerde, ikke aristoteliske. Efter at have angivet Systemet af de tre Slutninger ved a) \textit{E}---\textit{S}---\textit{A}, b) \textit{S}---\textit{E}---\textit{A} og c) \textit{E}---\textit{A}---\textit{S}, efter den Stilling, det Enkelte, det Særskilte og det Almene hver Gang indtage, vedbliver han: „Da Prædicatet altid er af større Omfang end Subjectet, saa er i den første Slutning ethvert foregaaende Led Subject, og ethvert efterfølgende Prædicat. Men i den anden Slutning er Mellemleddet Subject, og Yderleddene Prædicater, og i den tredie ere Yderleddene Subjecter, og Mellemleddet Prædicat. Eller: i den første Slutning bliver det Led, som først er Prædicat, siden Subject; i den anden er eet Subject fælles for to Prædicater; og i den tredie er eet Prædicat fælles for to Subjecter. Men i Conclusionen af alle tre Slutninger er det første Led Subject, og det tredie Prædicat. [Herved er at bemærke, at ved det fuldendte System af Slutninger blive Subject og Prædicat af lige Omfang; Forholdet mellem Art og Slægt forsvinder; det Almindelige er ikke af større Omfang end det Enkelte. Heri ligger allerede Overgangen til Objectivitet]“. Specul.\ Logik. §~155.}% note *)
“.

Men, saa klar og forstandig denne Deduction i sine Hovedtræk end maa siges at være, saa lider den dog af den formelle Logiks gamle, velbekjendte Grundmangel, at Forskjellen holdes udenfor Identiteten, og Modsigelsen bliver staaende som et logisk Skramsel. Skal Ligningen \textit{A} = \textit{A} udtrykke Identiteten, og den saaledes udtrykte Identitet være Princip, da vil Logiken umulig kunne drive det til en Dom, endsige til en Slutning. Selv om man med Ulrici vilde ansee den Sætning, hvori et Enkelt som Enkelt subsumeres under det Almindelige, for at være den eneste logiske Dom, saa er det dog klart, at man ved en saadan Sætning gaaer langt udenfor den rene Identitet. Kunde et Blik paa Dommen ikke strax lære En, at Identiteten \textit{A} = \textit{A} er en Tautologie, og en Tautologie utilstrækkelig til at begrunde en Dom, saa
% --- p. 322 ---
\opage{322}burde et Forsøg paa at deducere og bestemme Slutningen, navnlig den middelbare, dog kunne det. Det er Vexelbestemmelser af Slægt, Art og Individ, hvorom det Hele dreier sig. Hvorledes er det nu muligt at tænke sig Arten som Mellembegreb mellem Slægt og Individ, uden at see dette Mellemled paa den ene Side identisk med Individet, paa den anden Side med Slægten? Men Slægten er jo forskjellig fra Individet; Mellemleddet er altsaa identisk med to Forskjellige og trods denne Identitetsfordobling identisk med sig selv; her er tre Identiteter i een Identitet: en saadan treenig Identitet har i Sandhed sin Forskjel, ikke udenfor sig, men i sig.

Overseer man den Forskjelseenhed, uden hvilken Mellemleddet umulig vilde kunne forene Yderleddene, da vil man ogsaa blive staaende ved en meget eensidig Opfattelse af Slutningsfunctionen selv. Dette er netop Tilfældet med Ulrici. Istedenfor at erkjende Grundfunctionen, saaledes som den fremtræder i sine særskilte Former, d.\ v.\ s.\ saaledes som den i Virkeligheden fungerer, fastholder han dens abstracte Lighed med sig selv; Function = Function, ligesom \textit{A} = \textit{A}, og faaer paa denne Maade en Function, der ikke kan fungere\footnote{„Es giebt \dots\ vier \dots\ Schlußfiguren, die sich insofern logisch unterscheiden lassen, als sie aus dem allgemeinen Grundsatz alles Schließens in seiner positiven wie negativen Form, der zugleich die Norm für die Gestaltung aller Schlüsse ist, unmittelbar sich ergeben. Sie bilden indessen, wie von selbst einleuchtet, keineswegs vier verschiedene „\emph{Arten}“ von Schlüssen. Denn sie sind keineswegs Modificationen der logischen \emph{Function} des Schließens, sondern nur --- im Grunde unerhebliche --- Specialformen der allgemeinen \emph{Gestalt} der Schlüsse“. Anfr.\ Skr.\ S.~188.}% note *)
.

At Forskjellen ligesaa oprindelig hører med i Grunden, og altsaa i enhver Grundfunction, som Identiteten; at Identiteten selv i Grunden er Forskjel, har „den begrundende Reflexion“ paa sit Sted\footnote{Grundideernes Logik. I. S.~93--99.}% note **)
{} allerede godtgjort; at en Grund%
% --- p. 323 ---
\opage{323}function, naar den fastholdes i abstract Almindelighed, aldeles ikke kan anvendes, altsaa maa forblive i en tom Væren, uden at kunne naae intellectuel Tilværelse, kan endog ligefrem bevises mathematisk.

$F(x) = y$ er enhver mulig Function og netop af den Grund ubestemt: med en saadan Function kan der ikke regnes. Forsaavidt Functionen er \textit{in abstracto} som Grundfunction, er den, ved at være i Grunden, kun i Muligheden. „I Forhold til Reductionen er Muligheden Grund, i Forhold til Deductionen er Grunden Mulighed. Identiteten af Grund og Mulighed sees af den disjunctive Doms forskjellige Former: \textit{A} er hverken \textit{B} eller \textit{C} eller \textit{D} \dots\ $f(x)$ er hverken $ax$, eller $x^a$ eller $a^x$ \dots: de særskilte Bestemmelser ere opløste i det Almene, og derved gaaede til Grunde; \textit{A} er, i Grunden, baade \textit{B} og \textit{C} og \textit{D} \dots\ $f(x)$ er, i Mulighed, baade $ax$, $x^a$ og $a^x$ \dots\ Forsaavidt det i Grundens Indhold Særskilte lader sig tilsyne, reflecteres det i Mulighed; hvor eet Tilfælde er muligt, maae idetmindste to Tilfælde være mulige, thi dersom det mulige Tilfældes Modsatte ikke ogsaa var muligt, var Tilfældet selv nødvendigt: Muligheden har alternerende Led. Skal nu et af Leddene sættes som tilværende, da viser den disjunctive Dom i dens sædvanlige Form: \textit{A} er enten \textit{B} eller \textit{C} eller \textit{D} \dots, $f(x)$ er enten $ax$ eller $x^a$ eller $a^x$ \dots\ Spændingen imellem de alternerende Led\footnote{Mathematik og Dialektik. S.~50--51.}% note *)
“.

Hvad der i saa Henseende gjælder om de mathematiske Functioner, gjælder paa tilsvarende Maade om de logiske. Det er ikke muligt at anvende Slutningens Grundfunction uden at specificere den efter en af de fire Figurer. Er Dommen: „Cajus er dødelig“, fordret som Følgesætning, og Dommen: „alle Mennesker ere dødelige“, given som den ene Præmis, saa maa Dommen: „Cajus er et Menneske“, jo være den anden Præmis; det er da ingenlunde vilkaarligt,
% --- p. 324 ---
\opage{324}men aldeles nødvendigt, at Slutningen svarer til første Figur. Er Dommen: „nogle i Vandet levende Dyr ere Pattedyr“, fordret som Følgesætning, og Dommen: „alle Hvaler leve i Vandet“, given som den ene Præmis, da maa man nødvendigviis have enten: „alle Hvaler ere Pattedyr“; eller: „nogle Hvaler ere Pattedyr,“ til den anden Præmis, og det er da en Selvfølge, at Slutningen falder under tredie Figur. Kjendes nu Slutningsfunctionen, som ovenfor er viist, ikke alene paa den Tænkningsact, hvorved en Følgesætning udledes af givne Præmisser, men tillige paa den synthetiske Tankevirksomhed, hvorved Præmisserne udfindes og opstilles, da ligger det i Sagens Natur, at Slutningsfunctionen modificeres med hver af Slutningens typiske „Specialformer“. Præmissernes Synthese skeer nemlig under andre Betingelser og paa en anden Maade, naar Middelordet skal være Subject i Oversætningen og Prædicat i Undersætningen, end naar det skal være Prædicat baade i Over- og Undersætningen, og i sidste Tilfælde igjen under andre Betingelser og paa en anden Maade, end naar det enten skal være Subject i begge Præmisser eller Prædicat i Oversætningen og Subject i Undersætningen. Overseer man dette, maa Slutningsacten synke ned til en tankelos Mechanisme; men indseer man dette, saa vil man ogsaa indsee, at de synthethiske Modalfunctioner maae forandre sig med den forandrede Synthese. % PRINTER: printed "synthethiske", sense "synthetiske"; printed "tankelos", sense "tankeløs"

3) \emph{Syllogistikens Opløsning.} Men hvilken Betydning have de middelbare Slutninger nu egentlig for Viden og Videnskab? Man har, i nyere Tid især, indvendt mod Syllogistiken, at den endeliggjør Vidensformerne, udstykker Erkjendelsen, og ved Syllogismernes „bekjendte Beqvemhed til at bevise de forskjelligste, ja endog modsigende Ting“, befordrer Scepticismen. Syllogismens Vilkaarlighed ligger deri, „at den Beviisførende kan vælge Mellemordet (\textit{terminus medius}) efter Behag. Da nemlig enhver Ting frembyder flere Synspunkter for Betragtningen, og desto flere, jo con%
% --- p. 325 ---
\opage{325}cretere den er, saa kan den ved enhver enkelt af disse sammensluttes med det meest Forskjellige, ja endog Modsatte. \dots\ F.\ Ex.\ dersom man ved Syllogismer vil bevise, først at Mennesket er et Dyr, dernæst at Mennesket ikke er et Dyr, saa siger man i første Tilfælde: „Mennesket har Sandser, Organer, Circulations-, Digestions-, Nerve-System, o.\ s.\ v., kort, alle Betingelser for det dyriske Liv. Hvad, som har disse Betingelser, er et Dyr. Altsaa er Mennesket et Dyr.“ --- I andet Tilfælde siger man: „Dyret er ikke et Fornuftvæsen. Mennesket er et Fornuftvæsen. Altsaa er Mennesket ikke et Dyr\footnote{J.\ L.\ Heiberg. Speculativ Logik. §~148. Anm.~1.}% note *)
“.

Det Sande i denne Indvending fordunkles imidlertid derved, at Slutning og Beviis uden videre slaaes sammen i Eet. Her er, hvad Slutningen angaaer, kun Spørgsmaal om Erkjendelsens Formside, ikke om dens reelle Indhold. Vel maa den væsentlige Formudvikling ogsaa være en Indholdsudvikling, men det er da et logisk, ikke et reelt Indhold, hvorpaa der tænkes. Det er den articulerende Bestemmelse af Dommenes, altsaa af Erkjendelsernes Sammenhæng, der giver Slutningen, den middelbare saavelsom den umiddelbare, sin væsentlige Betydning At nu% PRINTER: no full stop after "Betydning"
{} en \textit{methodus mathematica}, som binder den fremadskridende Erkjendelse til en Kjæde af fortløbende Syllogismer, ikke kan anprises i Philosophien, er en afgjort Sag; men derfor bør Syllogismens Betydning som videnskabelig Reductionsformel ikke oversees. Det anselmske Beviis for Guds Tilværelse, som, naar det, om vi saa maa sige, sees i forkortet Perspectiv, lader sig reducere til følgende Syllogisme: „det fuldkomneste Væsen (\textit{id quo majus cogitari non potest}) maa have, ikke en blot tænkt, men ogsaa en virkelig Væren; Gud er det fuldkomneste Væsen; altsaa har Gud, ikke en blot tænkt, men en virkelig Væren“, er vistnok, som Beviis betragtet, utilfredsstillende; men som Reductions%
% --- p. 326 ---
\opage{326}formel for et uendeligt rigt Indhold, er Syllogismen ikke desto mindre af gjennemgribende Betydning. Oversætningen indeholder nemlig hele Ontologien, Undersætningen hele Religionsphilosophien, og Slutningen selv den Indsigt, som begge Videnskaber i Forening tilveiebringe. Vistnok kan ikke den enkelte Slutning, og endnu mindre den enkelte Dom, give det rigere Indhold en udtømmende, sikker og alsidig betegnende Form; men derfor er det dog ligefuldt af Vigtighed, at store Totaliteter kunne til punktuel Belysning concentreres i en enkelt Slutning, ja i en enkelt Dom.

Det anføres endvidere som en Mislighed ved Brugen af Syllogismer, at der etableres en uendelig Række, som Videnskaben ikke kan ophæve. „Da de to i Slutningen indeholdte Omdommer, % PRINTER: printed "Omdommer", sense "Omdømmer"
 \textit{E}---\textit{S} og \textit{S}---\textit{A}, fremstilles som isolerede Sætninger, som Præmisser, d.\ e.\ som Noget, der ikke umiddelbart ligger i Slutningen, men gaaer forud for denne, saa fordres, at enhver af disse Sætninger iforveien skal være fremkommen som Slutningssætning i en foregaaende Syllogisme, eller, som man siger, at den skal være beviist. Dette forudsætter to foregaaende Syllogismer, men da disses Præmisser ligeledes maae være beviiste i det Foregaaende, saa forudsætte disse to altsaa fire foregaaende, disse igjen otte foregaaende, og saaledes fremdeles i en uendelig Række, som i sit 21de Led vilde indbefatte omtrent halvanden Million Syllogismer, og derfra hurtigt stige til Billioner og Trillioner \dots\ At troe sig i Besiddelse af Sandheden, idet man befinder sig i et System af Syllogismer, er endnu urimeligere, end at troe sig fri, fordi man har Lov til at gaae omkring i et Buur, thi selv det mindste Fuglebuur er større i Forhold til hele Jorden, end den ved Syllogismer frembragte Erkjendelse i Forhold til den, som de forudsætte\footnote{Heiberg. Speculativ Logik. §~148. Anm.~1.}% note *)
“.

Men ogsaa hvad Indsigelserne i denne Retning angaaer,
% --- p. 327 ---
\opage{327}maa det erindres, at Paaviisning af Misbrug Intet beviser imod den sande Brug. Det gaaer med Syllogistiken i det Hele som med Tankens Endeliggjørelse. Dersom det Endelige ikke var med i Udviklingen, vilde det Uendelige jo slet ikke kunne begribes. Kan den immanente Methode behandle \emph{Reflexionens} endeløse Rækker som Betingelser for det Endeliges Overgang i det Uendelige, hvorfor skulde den da ikke kunne behandle \emph{den middelbare Slutnings} endeløse Rækker paa samme Maade; lade Reflectionsrækkerne sig ophæve, hvorfor skulde de syllogistiske Rækker ikke ogsaa lade sig ophæve?

Der er i hele den hegelske Logik ikke et Afsnit, hvis Indhold jo med Lethed kan reduceres til syllogistisk Form; men kan den immanente Udviklings Indhold reduceres til Syllogisme, da maa Syllogismen jo ogsaa kunne reduceres til immanent Udvikling. Man begynde f.\ Ex.\ med Begyndelsen: „I et System kan der kun gives een Begyndelse, der er den sande; Logiken er et System, altsaa kan der i Logiken kun gives een Begyndelse, der er den sande. Det ligger i Begyndelsens Natur, at den maa være indholdslos; kun „den rene Væren“ er indholdslos; altsaa maa Logiken begynde med den rene Væren. „Den rene Væren“ er absolut indholdslos; det absolut Indholdsløse er Intet; den rene Væren er altsaa Intet.“ --- Paa denne Maade kan man da blive ved gjennem Væren, Væsen og Begreb, og forvandle endog den allersidste Slutning til en Syllogisme% PRINTER: printed "indholdslos" three times in this paragraph (ø without stroke; "Indholdsløse" has ø), sense "indholdsløs"
\footnote{At man underveis kan faae Brug for alle fire Slutningsfigurer, og, naar det endelig kommer an paa at gjøre Kunststykker, endog faae Leilighed til at anvende baade \textit{Darapti} og \textit{Bamalip}, baade \textit{Ferison} og \textit{Fresison}, er en Sag, der forstaaes af sig selv.}% note *)
.

Men hvad lære vi nu heraf? For det Første, at den syllogistiske Form er en for Reductionen charakteristisk Abbreviatur, en Abbreviatur, der ret egner sig til at sammenfatte og
% --- p. 328 ---
\opage{328}opbevare Resultater af en tilbagelagt Udvikling. % PRINTER: below, printed "noiagtig", sense "nøiagtig"
 For det Andet, at den Maade, hvorpaa de vundne Resultater opfattes og fremstilles i de relativt selvstændige Domme, der udgjøre Slutningens Præmisser, noiagtig svarer til den Maade, hvorpaa de umiddelbare Præmisser igjen opløses og afgive deres Indhold til en fra syllogistisk Tvang frigjort Begrebsudvikling, en Begrebsudvikling, der da er identisk med den logiske Sagudvikling. Man overveie blot, hvorledes f.\ Ex.\ de to Præmisser: „den rene Væren er absolut indholdsløs, og det absolut Indholdsløse er Intet,“ blive anlagte, begrundede og gjennemklarede i den immanente Begrebsudvikling, og man vil snart opdage, hvorledes de syllogistiske Rækker lade sig ophæve. For det Tredie, at enhver logisk Indholdsudvikling, hvad enten den iøvrigt foregaaer i den immanente eller syllogistiske Form, nødvendigviis maa gaae igjennem et System af Domme. Men et System af Domme er umuligt uden et tilsvarende System af Slutninger. Hvad er vel Slutningen Andet end det logiske Udtryk for et System af Domme, den umiddelbare Slutning for et System af to Domme, den middelbare for et System af tre Domme, Kjædeslutningen for et System af flere Domme o.\ s.\ v.

„Mange“, siger Heiberg\footnote{Specul.\ Logik. §~148. Anm.~1.}, % note *)
„staae i den urigtige Formening, at naar man i sin Tale anfører Grunde, da ere disse kryptiske Syllogismer. Men snart giver man Grunde, og staaer da i Reflexionen, snart, naar man staaer i Idealiteten, fremfører man Slutninger, men aldrig Syllogismer. Naar man siger: „Cajus er ufuldkommen, fordi han er et Menneske,“ da er det ingen Slutning. Man angiver Grunden til hans Ufuldkommenhed. Hans Menneskenatur er det Væsen, hvis Phænomen er Ufuldkommenhed.“ Dersom disse og lignende Misforstaaelser ikke paa det Allerinderligste hængte sammen med den hegelske Mystification af Dommenes og
% --- p. 329 ---
\opage{329}Slutningernes metaphysiske Objectivitet, skulde vi ikke opholde os videre ved dem; men her er i den opstillede Synsmaade en Skjævhed, som for Principets Skyld absolut maa rettes. Skjævheden er, at Slutningerne ligesom Dommene skulle figurere objectivt paa et Stadium, hvor de endnu ikke engang ere blevne subjective. Siger man: „Cajus er ufuldkommen, fordi han er et Menneske,“ saa seer man vistnok Ufuldkommenheden begrundet i Menneskenaturen, forsaavidt man seer objectivt; men Sætningen: „Cajus er ufuldkommen“, er jo dog en Dom, Sætningen: „Cajus er et Menneske“, ligeledes. For at faae Øie paa Slutningen, maa man see at faae Øie paa Dommen; for at faae Øie paa Dommen, maa man ikke blot see objectivt, men ogsaa subjectivt. Begrundelsen: „Cajus er ufuldkommen, fordi han er et Menneske“, hviler aabenbart paa Slutningen: „Cajus er et Menneske, altsaa ufuldkommen,“ en umiddelbar Slutning, som Reflexionen med uendelig Lethed forvandler til en middelbar: „Ethvert Menneske er ufuldkomment, Cajus er et Menneske, altsaa“ o.\ s.\ v.

Erkjendes det først, at Dom og Slutning maae bestemmes subjectivt, altsaa oprindelig forudsætte en dømmende og sluttende Subjectivitet; erkjendes det, at enhver Tanke, enten den forøvrigt staaer i „Reflexionen“ eller i „Idealiteten“, kun kan bestemmes som Tanke, forsaavidt den i Ordet bestemmes ved en Dom: da vil det ogsaa indsees, at Tankernes Sammenhæng er Eet med Dommenes Sammenhæng, og Dommenes Sammenhæng, selv hvor Slutningsformen ikke udtrykkelig fremhæves, væsentlig Eet med et System af umiddelbare og middelbare Slutninger.

\textit{γγ)} \emph{Den immanente Slutning: dialektiske Modalfunctioner.} Hvorledes Dommen som Dom bliver Form for det absolute Indhold, er allerede oplyst ved „Functioner af Dommen.“ Da nu den enkelte Dom først faaer sin Betydning i Sammenhængen med andre Domme, og reflecterer Sammenhængen desto dybere, jo rigere dens Indhold er, alt%
% --- p. 330 ---
\opage{330}saa i en uendelig Dybde, naar den har Ideen, det Absolute selv, til Indhold: saa følger heraf, at kun et uendeligt System af Domme vil være istand til at rumme det Absolute. Men et uendeligt System af Domme er, naar Mellemleddene betones, \textit{eo ipso} et uendeligt System af Slutninger. Idet Formen for det Absolute udvikler sig til et System af Slutninger, er det en Selvfølge, at al Udstykning og Endelighed maa ophæves. De enkelte Præmisser blive da Momenter for den tilsvarende Slutning, idet hver enkelt Slutning bliver Moment for en relativ Totalitet af Slutninger, og hver relativ Slutningstotalitet igjen Moment for Totaliteternes gjennemgaaende Totalslutning. For at kunne blive Moment, maa ethvert af de særskilte Led, være sig Præmis, enkelt Slutning eller Slutningstotalitet, have særskilt Tilværen; thi Momentet fremkommer kun ved en Forsvinden af den særskilte Tilværen. Ved den saaledes indtrædende Vexelbestemmelse af Leddenes Moment-Væren og umiddelbare Tilværen blive Modalfunctionerne dialektiske.

Forsaavidt Præmisserne til den enkelte Slutning opgive deres særskilte Tilværen og blive Momenter i Slutningsacten, er Syllogistiken hævet: Slutningen er \emph{immanent}. Den immanente Slutning kan altsaa være Moment i en Totalitet af Slutninger og en Slutningstotalitet for sig. Som det Absolutes subjective Form frembyder Immanensslutningen tre Hovedsider, idet den oprindelig svarer til \emph{Subjectiviteten selv}, til Udviklingen af \emph{Subjectivitetens Indhold} og til \emph{den frie Nødvendighed}, hvormed Subjectiviteten besidder sit Indhold. Vi betragte enhver af disse Hovedsider for sig.

1) \emph{Den immanente Slutning som Form for Subjectiviteten selv.} At Subjectivitetens absolut oprindelige Slutning maa svare til en absolut oprindelig Dom, kan udtrykkelig paavises. Skal Schemaet \textit{E}---\textit{S}---\textit{A} gjælde for en oprindelig Subjectivitetsslutning, maa jo ethvert af
% --- p. 331 ---
\opage{331}dens Led have absolut Oprindelighed. E, det Enkelte, kan altsaa ikke betyde et vist Enkelt, dette eller hiint Enkelte, et Enkelt blandt flere andre, men Enkeltheden selv, Enkeltheden som Princip. Men Enkeltheden som Princip er Selvhed: \emph{det absolut Enkelte er det absolute Selv}. A, det Almindelige, kan da fra sin Side heller ikke være en Abstraction af et for Mange Fælles, et abstract Almindeligt, thi et abstract Almindeligt er kun en Relativitet; det absolut Almene er\footnote{Grundideernes Logik. I. S.~181--192.}% note *)
{} uendelig concret: \emph{det i alle Concretioner absolut Almene er det absolute Selv}.

Hvad der gjælder abstract i den yderste Abstraction, at nemlig den uendelige Modsætning er en uendelig Eenhed, og omvendt, det Samme gjælder concret i den fuldendte Concretion: det Enkelte er kun absolut Eet med det Almene ved at være det Almene absolut modsat. Kan nu Eenheden af det Enkelte og det Almene ikke være umiddelbar, efterdi den beroer paa Modsætning, og Modsætningen ikke umiddelbar, efterdi den beroer paa Eenheden: da følger det af sig selv, at her maa gives en Mægling. Først ved Bestemmelsen af denne Mægling kan Forholdet mellem Reflexion, Dom og Slutning nærmere fastsættes.

Ved Reflexionen bliver det absolut Enkelte jo vistnok adskilt fra det absolut Almene; her sættes i een og samme Selvhed en gjennemgaaende Modsætning imellem Selv og Selv. Men denne Modsætning forsvinder igjen for Reflexionen ligesaa hurtig, som den sættes, thi Selvet er dog sig selv, fra hvilken Side det end sees: Selvmodsætningen er Reflex i Eenheden, ligesom Selveenheden fra sin Side er Reflex i Modsætningen. Anderledes i den oprindelige Dom: at Selvet i absolut Forstand baade er Subject og Prædicat, betyder at Modsætningen af Subject og Prædicat er ligesaa uudslettelig som Eenheden. Medens Reflexionen uophørlig udligner Mod%
% --- p. 332 ---
\opage{332}sætningen af Modsætning og Eenhed, sætter Dommen sig med afgjørende Eftertryk mod denne altfor flygtige Udligning, og fordrer udtrykkelig Modsætningens Bestaaen i den totaliserede Eenhed. Da det nu i lige absolut Forstand er Alvor med Leddenes Identitet og Modsætning, saa danne de modsatte Led, Subjectet (det absolut Enkelte) og Prædicatet (det absolut Almene) Extremer, hvis Mellemled, S, \emph{det absolut Særskilte}, skyder op af Identitetens Grund. Den oprindelige Dom udvikler sig med uafviselig Conseqvens til immanent Slutning.

Ifølge den immanente Slutning er Subjectiviteten paa een Gang den oprindelige Doms absolute Forudsætning og dens Resultat. Det dømmende Selv kan umulig afledes af Dommen; Dommen forudsætter en Domsact, Domsacten en dømmende Selvhed: paa den dømmende Selvhed beroer Selvets absolute Prioritet. Men som det første logiske Udtryk for den oprindelige Selvhedsact gaaer Dommen ikke ud paa denne eller hiin endelige Gjenstand, har ikke dette eller hiint relative Indhold; nei, Dommens, den oprindelige Doms oprindelige Indhold er Subjectiviteten selv, Domsacten en første absolut nødvendig Selvsætningsact, en primitiv Handling, hvorved den i Begrebet Dømmende oprindelig skifter og deler imellem Selvhed og Andethed\footnote{Das Urtheil (ἀπόφανσις, als Bestandtheil des Schlusses πρότασις = \textit{propositio}) ist gleichsam der richterliche Bescheid, welchen das Denken über das Verhältniß zweier Begriffe ertheilt; daher ist das Wort „\emph{urtheilen}“ nicht, wie Einige behaupten, von „\emph{Ur- und Theilen}“ abzuleiten, sondern es kommt von dem mittelhochdeutschen Worte „\emph{urteilen}“; niederdeutsch: \emph{ordelen} (wovon „Ordalie“) \emph{aussagen}, zu Recht erkennen, rechtlich entscheiden.“ Drbal. Lehrbuch der propädeutischen Logik. S.~29. Kom det nu i en Logik virkelig an paa at tilveiebringe Overeensstemmelse mellem lingvistiske, og dialektiske Ordforklaringer, saa vilde det vist ikke være vanskeligt at vise, at „\emph{ordelen}“ i Grunden er „urtheilen“ og „urtheilen“ „ordelen“.}% note *)
% PRINTER: the quotation from Drbal closes with “ after "entscheiden." but has no opening „ before "Das Urtheil"; left unbalanced as printed.
.

% --- p. 333 ---
\opage{333}Den Dom, hvori den evige Subjectivitet udtaler sit Væsen, er ikke, hvad Fichte antager, en tautologisk Sætning: Jeg = Jeg; en saadan tautologisk Sætning er en yderste Abstraction af psychologisk Tænkning, men ingen \textit{in concreto} oprindelig Dom. I den concret-oprindelige Dom er Selvet, det absolut Enkelte, som Subject stillet imod Selvet, det absolut Almene, som Prædicat. Den gjennemgaaende Modsætning, Subjectivitetens Selvmodsætning, er sat og bestemt derved, at Subjectet i Dommen selv er det dømmende Subject og Dommens Prædicat det oprindelige Begreb, hvori Subjectet griber sig selv som Object.

At den Grundact, hvorved Subjectiviteten gjør sig selv til Subject og Prædicat, saaledes at Prædicatet bliver det objectiverede Subject, ikke er enten en umiddelbar Selvspeiling, eller en abstract Reflecteren, men en logisk Dom, faaer afgjørende Betydning ved Bestemmelsen af den immanente Slutning. For at udtale sit Væsen i Dommen, og ved denne Udtalelse adskille sig i Subject og Object, maa Subjectiviteten nødvendigviis oversvæve sin egen Selvadskillelse og forsaavidt sætte sig, ikke blot som det i Dommens Subject \emph{Enkelte} og det i Prædicatet \emph{Almene}, men tillige som det i den dømmende Selvhed med Subject og Prædicat Identiske og dog fra begge Forskjellige, altsaa som det \emph{Særskilte}. Det er ved denne tredobbelte Selvadskillelse (E--A, E--S, S--A), at Subjectiviteten forvandler sit Væsens oprindelige Dom til en oprindelig immanent Slutning.

Men Subjectivitetens tredobbelte Selvadskillelse, som ved en fortsat Oversvæven af den sig oversvævende Selvhed kan fordobles og gjentages i det Uendelige, vilde dog nedsynke til en uendelig tom Formalisme, dersom her ikke i Subjectiviteten selv gaves et Middel, hvorved Extremer og Midte under deres indbyrdes Overgang kunde holde sig adskilte. Men et saadant Middel er \emph{Ordet}. Dersom Ordet ikke var, vilde den subjective Tænken og dermed Subjectiviteten evig fortære
% --- p. 334 ---
\opage{334}sig selv; dersom Ordet ikke var, vilde Anskuelsens umiddelbare Andet staae som et fremmed, uforklarligt Object, paa hvilket den stumme Subjectivitet maatte stirre sig blind; dersom Ordet ikke var, vilde Selvet, det Enkelte, saaledes gjennemskinnes af det Særskilte og det Almene, at Reflexionens rene Lys, den ubestemte, i sig ubestemmelige Klarhed tilsidst forsvandt i forskjelsløs Eenhed og tabte sig i det rene Mørke.

Ved at mægle mellem Tænken og Anskuen og saaledes gjøre den principielle Selvbestemmelse mulig, sætter Ordet den immanente Slutning i Kraft. I Ordet har Tænkningen sine elementære Andethedsbestemmelser, sit uendelige Forraad af Billeder og Forestillinger. Der er i den virkelige Verden, i Realiteternes uendelige Rige ikke en Gjenstand, intet Hele, ingen Deel, ikke en Ting, ikke en Egenskab, uden at Ordet jo ved Billede og Forestilling formaaer at optage og gjengive det deri Eiendommelige. Beriget med Reflexer af alle mulige Realbestemmelser har Ordet nu den Idealitet, at det i Reflexionen opgiver sin Umiddelbarhed og offrer sig for Tanken saa inderlig og absolut, at den Tænkende gjennem Ordet tænker sig selv i Alt og i Alt sig selv, og begriber sin Subjectivitet som \emph{det absolut Almene}.

Det absolut Almene er jo vistnok en Tanke; men det er ikke en fattig Tanke, thi Tanken lever af Ordet, og Ordet har en uendelig Rigdom. Det absolut Almene er jo vistnok, for sig betragtet, et Ord; men det er ikke et tomt Ord, thi Ordet, der nærer Tanken, beriges af Tanken, og Tanken er uendelig rig.

I den ved Ordet berigede Tænkning begriber Subjectiviteten sig selv under alle tre Skikkelser, Enkelthedens, Særskilthedens og Almindelighedens, men er dog, idet det Almene gjennem det Særskilte slutter sig sammen med det Enkelte, endnu kun sat i Tænkningens, i det Almenes eensidige Form: Subjectiviteten er tænkende. Modsætningen af Subject og Object er vel tænkt som en gjennemgaaende, absolut Mod%
% --- p. 335 ---
\opage{335}sætning, men det er dog endnu kun en tænkt Modsætning; Eenheden af Subject og Object er jo vistnok tænkt som en absolut, ved den særskilte --- paa een Gang oversvævende og dog Extremerne iboende --- Selvhed bestemt Eenhed, men det er med alt dette dog kun en tænkt Eenhed. Den hele immanente Slutning er ved den Alt gjennemtænkende Tænknings Eensidighed selv blevet Extrem.

Mod Selvheden i sit Extrem af det Almene staaer nu ikke, som før, Extremet af det Enkelte, men et Selvhedsextrem af det Særskilte. Denne Extremforandring er afgjørende bestemt ved Ordet. I Tænkningen er Ordet negativt, dets Umiddelbarhed forbrænder i Reflexionen; kun Mangfoldigheden af Lysninger, Afskygninger, Skjær og Reflexer udgjør den sig i Alt og Alt i sig tænkende Subjectivitets Rigdom. At nu Reflexionen ikke kan fortære denne Rigdom og styrte Subjectiviteten i en saadan Fattigdom, at den seer ud som en Abstractionens Skygge, et psychologisk Gjenfærd af Jeg = Jeg, er forebygget ved Ordet; thi Ordet, der er negativt i dets Eenhed med Tanken, er positivt i dets Eenhed med Anskuelsen. Vel maa Ordet ophæve al sandselig Umiddelbarhed og forsaavidt bevirke en Forvandling af alle Sandsebilleder; men de billedlige Elementer, der fortæres af Tænkningen, forklares i Anskuelsen og bestemmes ved Ordet. Forsaavidt Ordet er negativt, gaaer Billedet gjennem Forestillingen over i Tanken; men forsaavidt Ordet formedelst Anskuelsen er positivt, gaaer Tanken op i Forestillingen, og Forestillingen i et forklaret Billede, et Tankesyn, en intellectuel Anskuelse. I uendelig Mangfoldighed af Syner, Billeder og Forestillinger er den anskuende Subjectivitet Alt, hvad den saavel i det Almene som i det Enkelte er og kan være, i Ordet, ved Ordet, som Ordet. Ved sit Positive i Anskuelsen og sit Negative i Tænkningen er Ordet det Særskilte. Det er ikke et enkelt Extrem af Subjectiviteten, men den hele Subjectivitet i Extremer og Midte, der forklares af Ordet. Den tænkende, i Ordet
% --- p. 336 ---
\opage{336}sig objectiverende Subjectivitet er altsaa bestemt som \emph{det Særskilte}: formedelst det Særskilte som Object er Selvet som Subject det Almene ved at være det Enkelte, og det Enkelte ved at være det Almene.

Og dog er her endnu en Ubestemthed, en Svæven, et Anstrøg af Formalisme, der truer med at udhule Slutningen og forvandle Subjectivitetslæren til logisk Phraseologie. „Hvorledes kan Subjectiviteten, som det i Ordet Særskilte, paa een Gang baade holde sig i Modsætning til det Enkelte paa den ene og det Almene paa den anden Side, trods denne Modsætning tillige være identisk saavel med det Enkelte som med det Almene, ja ved denne Identitet bevirke, at det Enkelte selv er det Almene, det Almene det Enkelte?“ Et saadant Spørgsmaal er fra den formelle Logiks Synspunkt ikke blot naturligt, men i nærværende Sammenhæng fuldkommen berettiget. Ved Indsigt i Subjectiviteten og den immanente Slutning kan Sagen imidlertid opklares.

Al Subjectivitetsbestemmelse er væsentlig Selvbestemmelse, men en Selvbestemmelse, der forudsætter Andet og betinges af Andet. At den fuldstændige Subjectivitetsbestemmelse antager en tredobbelt Form, har netop sin Grund i, at Bestemmelsen af Andet skal gaae op i Bestemmelsen af Selvet. Dette er nu paa Relativitetens og Endelighedens Trin en Umulighed. Bestemmelsen af Andet staaer i psychologisk Forstand som en Skranke, en Hindring for Selvbestemmelsen. Ligeoverfor Andet, den mægtige Skranke, er Selvbestemmelsen afmægtig, abstract og formel. Afmagten sees da ogsaa baade paa Tanken og Ordet, thi den rene Selvtænken er en tom Reflecteren, og Ordet et Navn, der Intet formaaer over Gjenstanden. Henført til den i sin Indskrænkning formelle og abstracte Subjectivitet, maa den immanente Slutning med dens Extremer og Midte naturligviis ogsaa blive formel. Forskjellen mellem det sig tænkende og det tænkte Selv, mellem den sig objectiverende og den objectiverede Subjectivitet, er her en
% --- p. 337 ---
\opage{337}Ordforskjel, altsaa kun en Navneforskjel: Bestemmelserne ere nominelle. Men naar Talen er, ikke om den psychologiske, men om den ontologiske, ikke om den relative, men om den absolute Subjectivitet, maa Nominalismen forsvinde.

Det er det almægtige Ords absolute Idealitet, hvorpaa Eftertrykket falder. Allerede i Læren om de ideeformende Functioner blev den tænkende Subjectivitet opfattet som Ordets Ophav, den anskuende Subjectivitet som Ordet selv, og den sig ved Tanken og Ordet bestemmende Subjectivitet som udviklet Eenhed af begge. Men saalænge Dommen ikke var gaaet over i Slutning, vare Vidensordet med sin Idealitet og Magtordet med Realiteten endnu saa endelig adskilte, at Magten eensidig lød paa Andethed, naar Viden lød paa Selvbestemmen og Selvhed. Under Functionernes fremadskridende Udvikling blev det i Modsætningen Eensidige efterhaanden hævet; ved Dommens Overgang til Slutning blev Videns og Magtens Idealitet udtrykkelig sat i Ordet. Idet nu Ordet saaledes faaer subjectiv, ideel Betydning baade fra Magt og Viden, kommer dets mægtige, Alt formaaende, i Alt overgribende Idealitet tilsyne.

Hvad Ordet bestemmer om Subjectivitetens tredobbelte Selvadskillelse i den immanente Slutnings tredobbelte Form, bestemmer det ikke blot paa Videns, men ogsaa paa Magtens, ikke blot paa Selvhedens, men ogsaa paa Andethedens Vegne: Subjectivitetens Bestemmelse ved Ordet er dens egen absolute Selvbestemmelse, thi Ordets Ophav er i sin Modsætning til Ordet identisk med Ordet.

Lade vi nu \emph{Middelordet} angive Forskjellen mellem Slutningssystemets tre Slutninger, saa bliver Schemaet E--A--S betegnende for Subjectiviteten som \emph{det absolut Almene}, E--S--A for Subjectiviteten som \emph{det absolut Særskilte}, og S--E--A for Subjectiviteten som \emph{det absolut Enkelte}. At ethvert af Totalslutningens Led selv er en Slutning, og enhver af de særskilte Slutninger igjen Mo%
% --- p. 338 ---
\opage{338}ment i Totalslutningen, er betinget af Ordet; vi skulle nu see, hvorledes Ordet virker som Magtens Idealitet.

Den tænkende Subjectivitet bruger, som sagt, Ordet negativt. I det umiddelbare Videns og Magtens Ord er Altilværelsens ideelle Rigdom indesluttet; men Umiddelbarheden offres paa Tanken. Hvad det vil sige, at Ordet i den uendelige Tænken maa bruges negativt, kan endnu oplyses fra en anden Side. Det gjør jo en væsentlig Forskjel, om vi ved saadanne Udtryk som: det Enkelte, det Særskilte, det Almene, reflectere paa Ordet eller paa Tanken. Reflectere vi eensidig paa Ordet, saa er her i det Enkelte, det Særskilte og det Almene tre forskjellige Ord. Disse tre forskjellige Ord ere i al Overfladiskhed tre forskjellige Benævnelser; hver Benævnelse kræver da, hvis den skal betyde Noget, en særegen Forestilling, en Anskuelse, et Billede af det, der benævnes: vi føres ved Ordets Umiddelbarhed bort fra Tanken over i Billedernes, Anskuelsernes og de endelige Forestillingers Kredsløb. Anderledes naar her udtrykkelig reflecteres paa Tanken, den ved Ordets Negativitet almægtige Tanke. De tre Ord ere nu tre ideelle Magtbestemmelser, tre substantielle Begreber; det gaaer i denne Forstand med Ordet som med Lyset: vi see Gjenstandene i Lyset og glemme Lysstraalerne; vi see Tankerne i Ordet og glemme Benævnelserne.

Subjectiviteten, det tænkende Ene (E), faaer ved Ordets negative Virken ingen Billeder, ingen Syner, ingen endelige Forestillinger, men en uophørlig Forsvinden af utallige Billeder, Syner og Forestillinger, utallige særskilte Tanker. I særskilte Tanker sætter Subjectiviteten en væsentlig Forskjel imellem den Tænkende, Tænkningen og det Tænkte; den tænker ved denne Forskjel sig selv særskilt. Den ved Ordet tænkende (E) og den i Ordet tænkte Subjectivitet (S), danne nu substantielle Extremer; men i Ordets negative Functioner, i den uendelige Tænken, have disse Extremer deres væsentlige Midte. Forsaavidt den uendelige Tænken er Extremernes Midte, altsaa den virkende Form, hvori Tænkeren selv er det Tænkte, er Subjectiviteten bestemt ved at bestemme sig som den uendelige, Alt gjennemtænkende,
% --- p. 339 ---
\opage{339}Alt i Tanken formaaende Eenhed, altsaa bestemt som \emph{det absolut Almene} (A). Schemaet E--A--S er fyldestgjort.

Ved den uendelige Selvtænken kan Subjectiviteten jo vistnok siges at grunde, fordybet i sin absolute Eenhed; Eenheden er et uendeligt Dyb, men Dybet er Klarhed: magiske Syner, gnostiske Æoner, Drømmebilleder af βυθὸς og σιγὴ, προπάτωρ og ἔννοια, forsvinde for Tankens Energie i den immanente Slutning. Tankens Energie er Ordets negative Virken. I sin negative Virken er Ordet skjult; men det skjulte Ord maa blive aabenbart, dets negative Virken gaae over i den positive.

Den tænkende, i sig fordybede Subjectivitet, der holder Ordet skjult, er \emph{Ordets Ophav}. Ordet kommer ikke Tanken til Hjælp udenfra; det for Tanken bestemmende Ord er oprindelig bestemt ved Tanken selv: Ordet er med Nødvendighed i Tanken, det fødes af Tanken. Født af Tanken faaer Ordet først positiv Betydning som Ord: den tænkende Subjectivitet bliver i Reflexionen anskuende, og Slutningens Schema forandres til E--S--A.

Med Ordets reflecterede Umiddelbarhed træder \emph{det Særskilte} i Kraft. Vistnok maa Tænkningen, som ovenfor viist, forudsætte baade det Særskilte og det Enkelte; men i Tænkningen som Tænkning kunne de nævnte Extremer kun virke negativt gjennem Reflexer; først naar det Umiddelbare dukker op af Reflexionens Strøm, og Tanken farves af Anskuelsens Lys, bryder Forskjellen frem, og det Særskilte kommer til sin Ret. Med Præg af umiddelbar Anskuelighed er det Særskilte tillige bestemt som det Objective, det udvortes Mangfoldige; den anskuende Subjectivitet (E) fæster nu Blikket paa det Særskilte (S) og seer en Mangfoldighed af Syner. Men det er ikke Sandsebilleder; det er Syner i Ord, en Mangfoldighed af særskilte Domme, her staaer for Subjectivitetens Blik. I enhver af disse Domme er Ordet særskilt; Umiddelbarheden er i Subjectet, Reflexionen i Prædicatet: Subject
% --- p. 340 ---
\opage{340}og Prædicat bestemmes særskilt ved særskilte Ord, Umiddelbarhedens og Reflexionens Eenhed maa da ogsaa bestemmes ved et særskilt Ord. Det særskilte Ord er relativt; Relativiteten ligger i Mangfoldigheden. Ved sin Relativitet har Selvhedens Ord optaget Andetheden; ved Andetheden hæfter Objectiviteten; det i objective Bestemmelser særskilte Ord staaer altsaa som Subjectivitetens Andet mod Subjectiviteten selv: ved det særskilte Ord bliver den anskuende Subjectivitet selv et Særskilt (S).

Stilles Subjectiviteten skarpt imod Ordet, da viser det sig klart, hvorledes det Særskilte fremkommer; det er jo ikke blot Subjectiviteten, der bliver særskilt formedelst Ordet, det er, ifølge Ordets og Tankens oprindelige Eenhed, ligesaameget Ordet, der bliver særskilt formedelst Subjectiviteten. Er Subjectiviteten i Ordet \emph{subjectivt} dobbeltsidig ved at være sig selv og sit Andet, da er Ordet i Subjectiviteten \emph{objectivt} dobbeltsidig ved at være sig selv og sit Andet; er Subjectiviteten med sine relative, ved Relativiteten særskilte, Tænkningsacter en absolut Eenhed, da er Ordet med sine relative, ved Relativiteten særskilte, Elementer ligeledes en absolut Eenhed. Med Blikket fæstet, ikke paa Ordet i sig, men paa \emph{Ordets Andet}, paa det i Mangfoldigheden og Forskjelligheden \emph{særskilte} Ord, seer den anskuende Subjectivitet kun \emph{sit Andet}, ikke \emph{sig selv}: Særskiltheden er paa begge Sider relativ. Ved at fortabe sig i det Relative vilde den anskuende Subjectivitet da ikke blot blive særskilt mod sit Andet i Ordet, men tillige særskilt mod sig selv i Tanken; den anskuende Subjectivitet vilde, med andre Ord, ikke være tænkende, den tænkende ikke anskuende; ved en endelig Adskillelse af Tanken og Anskuen vilde den uendelige Subjectivitet synke ned i Endelighed og Afmagt; men det er en Modsigelse, ja en uendelig Selvmodsigelse.

For at hæve Selvmodsigelsen maa den anskuende Subjectivitet altsaa hæve sig over det Relative, forklare Tænk%
% --- p. 341 ---
\opage{341}ningen ved Anskuelsen, Anskuelsen ved Tænkningen, og gribe sig selv absolut ved at begribe Ordet. Som Subjectiviteten er, saaledes er Ordet; som Ordet er, saaledes er Subjectiviteten. I det absolute Ord anskuer den tænkende Subjectivitet sit eget Billede. Der er i dette Billede en objectiv Umiddelbarhed, thi den tænkende Subjectivitet er \emph{anskuende}; men Subjectivitetens Billede er Subjectiviteten selv, altsaa et ubilledligt Billede, et almægtigt Tankebillede, den i Ordet og som Ordet objectiverede Subjectivitet. I den absolute Selvobjectivering gaaer altsaa Tænkningen sammen med Anskuelsen, det Almindelige (A) sammen med det Særskilte (S); hermed har Slutningen naaet sin Fuldendelse i Særskilthedens Form.

„Men“ --- kunde man spørge --- „er der nu ogsaa nogen sand logisk Nødvendighed for, at Slutningen ogsaa skal gjentage sig i Enkelthedens Form? Det synes dog virkelig, som om her allerede var sagt Alt, hvad der fra Subjectivitetens almindelige Synspunkt paa nogen Maade kan siges om den immanente Slutning. Er det fra Begyndelsen af givet, at Subjectiviteten selv er baade det absolut Enkelte, det absolut Særskilte og det absolut Almene, da maa det jo være en ligesaa let som ufrugtbar Kunst at omsætte de tre Led paa tre forskjellige Maader, retfærdiggjøre disse paa Middelordet beregnede Omsætninger ved passende Omskrivninger, og lade hegelske Combinationer af E, S og A træde istedenfor \textit{Barbara, Celarent} og \textit{Baroco}“.

Det er imidlertid ingen tom Schematisme, men en tydelig Fremstilling af Ordets og Tankens absolute Energie, hvorpaa her i nærværende Sammenhæng lægges an. Dersom det nu blot kom an paa Betegnelsen, det sproglige Udtryk, den kategoriske Formside, da er Ordet Energie hos Aristoteles og efter ham hos Neoplatonikerne, Scholastikerne og Leibnitz fra mange Sider blevet opfattet, bestemt og forklaret. Men naar det ikke er Ordet Energie, men Energien i Ordet, paa hvis Erkjendelse her væsentlig sigtes, ville Definitioner, Om%
% --- p. 342 ---
\opage{342}skrivninger og formelle Forklaringer ikke strække til. Ordets primitive Energie, Energien i det guddommelige Ord, er bestemt ved Modsætningseenhed af Viden og Magt: i Ordet og ved Ordet er Viden forskjellig fra Magten, i Ordet og ved Ordet er Viden identisk med Magten. Heraf følger, at, naar Eenheden er sat, saa bryder Modsætningen frem med uimodstaaelig Kraft: i denne Uimodstaaelighed er Energien; omvendt, naar Modsætningen er sat, hæver Eenheden sig med overgribende, uimodstaaelig Overlegenhed: i denne Overlegenhed sees Energien. Da nu Ordet, i absolut Forstand, er identisk med Subjectiviteten, saa maa Ordets Energie tillige være Subjectivitetens Energie; men Subjectivitetens absolute Energie er dens absolute Selvbestemmelse.

Ved absolut Selvbestemmelse griber da den med Ordet identiske Subjectivitet sig selv i Alt, Alt i sig selv, og bliver derved absolut bestemt som \emph{det absolut Enkelte}. I Bestemmelsen af Subjectiviteten som det absolut Enkelte virker Ordet paa een Gang negativt og positivt. Det virker negativt; den absolute Selvbestemmelse er negativ Energie, thi Selvet kan kun være det absolut Enkelte ved at være det absolut Almene. Men da det i Tænkningens Form Almene har sin mægtige Modsætning i Anskuelsen af det Særskilte, bliver Selvbestemmelsen kjendelig paa den Energie, hvormed Tænkningens negative Ord magter og hæver Umiddelbarheden. Dog, Ordet kan kun virke negativt for Tænkningen derved, at det med lige Energie virker positivt for Anskuelsen. Den absolute Selvbestemmelse er positiv Energie, thi Selvet kan kun være det absolut Enkelte ved at sætte sig mod alt Andet og alt Andet imod sig, altsaa kun ved at blive det absolut Særskilte. Den negative Energie, hvormed Subjectiviteten udelukker Andet, for at sætte sig selv, er en Kraftmaaler for den positive Energie, hvormed Subjectiviteten sætter Andet, idet den sætter sig selv, og omvendt. I energisk Eenhed af negativ og positiv Virken er den absolute Selvbestemmen en absolut
% --- p. 343 ---
\opage{343}Villen: kun den villende Subjectivitet kan i egentlig Forstand siges at være \emph{det absolut Enkelte}. Med det absolut Enkelte til Mellemled er den immanente Slutning udviklet til at være den ontologiske Subjectivitets adæqvate Form.\footnote{„Men“, vil den opmærksomme Læser maaskee spørge, „er dette nu ikke et maskeret Forsøg paa at forklare Treenigheden, og saaledes gjøre Indgreb i den speculative Dogmatiks al philosophisk Concurence udelukkende Kirkeprivilegium?“ Vi svare, at de rationelle Former, hvormed Dogmatiken, for ikke ligefrem at indskrænke sig til en Sammenstilling af Skriftsteder, pleier at smykke sin Treenighedslære, alle ere laante, stundom endog sammenrapsede, fra forskjellige philosophiske Systemer. En alsidig Udvikling af den immanente Slutning som Subjectivitetens oprindelige Form falder derimod ganske og aldeles under Logiken; her laanes Intet udenfra, allermindst fra Dogmatiken. Hvad der i den objective Treenighedslære er ægte theologisk og forsaavidt udelukkende tilhører Dogmatiken, skal Philosophien, og da især Logiken, vel vogte sig for at begjære. Hertil høre saadanne „dogmatiske Oplysninger“ som f.\ Ex.: „Kan der siges, at den ene Gud ligesom seer ind i Verden med tre Ansigter (τρία πρόσωπα), da maa der ogsaa siges, at disse Ansigter ikke blot vende ud mod Verden, men at de ogsaa vende ind mod sig selv, at de see sig selv i gjensidig Speiling. Ellers vilde de ikke være Aabenbaringen af Guds sande Indre, men kun skuffende Masker“. Det er en vigtig, en meget vigtig „dogmatisk Oplysning“, thi deraf sees, at hvert Ansigt maa, for ikke at være „en skuffende Maske“, paa een Gang vende baade indad og udad, følgelig være et Dobbeltansigt. For ikke at have „tre skuffende Masker“, ved blot at have tre udadvendte Ansigter, maa Dogmatikens Treenige altsaa have, i det Allermindste, sex Ansigter. Endvidere: „den hele Trinitet staaer i eet nærværende Nu, tre evige Flammer i det ene Lys“. Endvidere: „den sig selv aabenbare Gud, for hvem Alt er aabenbart, \emph{har} ikke et Øie, men han \emph{er} Øie“. Altsaa: Dogmatikens Treenige \emph{har} sex Ansigter, men ikke et Øie. Thi han \emph{er} Øie, et treenigt Øie: tre Flammer i eet Lys, sex Ansigter i eet Øie! At faae et saadant Øie med saadanne Flammer, i et saadant Lys, med saadanne Ansigter ind i Videnskaben dertil udfordres overnaturlige Kræfter, ja dertil fordres Tro og \textit{testimonium spiritus}.}% note *)
% PRINTER: printed "Concurence", sense "Concurrence".
{} I
% --- p. 344 ---
\opage{344}denne Form har Subjectiviteten Hvile, thi Extremer og Midte ere simultane; men denne Hvile er en evig Bevægelse, thi Extremer og Midte sættes og omsættes i det Uendelige.

2) \emph{Den immanente Slutning som Form for Subjectivitetens Indhold.} I den absolute Subjectivitet er Indholdet selv absolut; Forskjellen imellem Subjectivitet og Indhold er altsaa en Forskjel i det Absolute selv, en Forskjel imellem Absolut og Absolut. Men Forskjellen i det Absolute selv er Eet med Relativiteten. Da nu det Absolute absolut maa have det Relative som immanent Bestemmelse, kan Forskjellen mellem Subjectiviteten og Indholdet alene beroe paa det dobbelte Forhold, hvori det Relative staaer til det Absolute: \emph{enten gaaer det Relative op i det Absolute, eller omvendt det Absolute i det Relative}. Gaaer det Relative op i det Absolute, saa giver den immanente Slutning os, som ovenfor er viist, Subjectiviteten selv; gaaer det Absolute derimod op i det Relative, saa giver Slutningen, hvad i det Følgende skal vises, den immanente Form for Subjectivitetens Indhold.

Forskjellen i det Relatives Forhold til det Absolute er kjendelig paa ethvert af Slutningens Led. At det Enkelte, forsaavidt det Relative gaaer op i det Absolute, maa blive Subjectiviteten selv, er viist; men naar det nu er det Absolute, der gaaer op i det Relative, hvad bliver der saa af det Enkelte? En Mylder af Enkelte! At det Almene, forsaavidt det Relative gaaer op i det Absolute, maa blive Subjectiviteten selv, er viist; men naar det Absolute nu gaaer op i det Relative, hvad bliver der saa af det Almene? En Fleerhed af Almene! Og hvad bliver det paa samme Vilkaar til med det Særskilte? En Mangfoldighed af Særskilte!

I Subjectiviteten selv, den absolute Subjectivitet, er der ingensomhelst endelig Forskjel imellem Slægt, Art og Individ. Først med Relativiteten kommer Adskillelsen. Det er Relativiteten, der i Slægten bevirker Adskillelsen af de mange
% --- p. 345 ---
\opage{345}Arter, og i hver Art igjen Adskillelsen af de mange Individer. Udsprungen af det Absolute er Relativiteten principiel, ja selv et Princip: Relativitetsprincipet er Særskilthedens Princip; Andethed i Tanke og Ord er Principets Energie.

Dersom Særskilthedens Princip ikke ved Andetheden gjennemvirkede Subjectivitetens uendelige Indhold, vilde det subjectivt Objective, det forskjelsløse, alle Tanker, alle Begreber og Former udjævnende Almene hvile i uendelig Eensformighed og fortætte sig til evigt Mørke. Men naar Særskilthedens Princip bringer det Almene i Zittren, naar Tanken lyser i Ordet, og Ordet med Tankens Lys i Relativitetens Farveskjær adspreder Mørket og aabner Erkjendelsen Indgang: hvilke Syner fremstille sig da for Anskuelsen, hvilke Skikkelser vil Tankens Blik da møde i Ideernes Rige? Det skulde dog vel aldrig være „den speculative Dogmatiks“ „legende Muligheder“, „plastiske Forbilleder“, „magiske Virkeligheder“, en „himmelsk Hærskare“ af de „Syner“, der ikke kunne nøies med en „Aabenbaring \textit{ad intra}“, men ere i paatrængende Spænding for at slippe ud og „loslades til Aabenbaring \textit{ad extra}“? Ingenlunde! For at see slige Syner maa Speculationen i Selskab med Troen først indtage et Qvantum kirkelig \textit{spiritus}, thi i ædru Tilstand seer Tanken noget ganske Andet. „Magiske Syner“, speculative Hallucinationer egne sig ikke for en guddommelig Subjectivitet, og maadelig nok for en menneskelig. Med en Hærskare af slige Syner kan Individualiteten maaskee blive en gjærende Mystiker, men ikke en Digter, endnu mindre en Tænker.

De Syner, der oprinde for Digteren, de Forbilleder, der foresvæve Tænkeren, have som Tilstrømninger fra Ideernes Kilde altid, om Strømmen end gaaer nok saa stærk, det Særkjende, at de komme med Ordet, ordne sig i Tanker og forklares i Domme; selv i den mythiske Bevidsthed maae Synerne trods al deres Magie underordne sig Ordets Magt. Men Syner, der underordne sig Ordets Magt, have ikke den
% --- p. 346 ---
\opage{346}„creaturlige“ Nærgaaenhed, den apokryphiske Paatrængenhed, at de umiddelbart skulde „begjære at løslades til en Aabenbaring \textit{ad extra}“; nei, tvertimod! de klare sig just til Aabenbaring \textit{ad intra}, thi Videns Ord i Ideernes Dyb er Aabenbaring \textit{ad intra}.

I Ideernes Rige, i Guddommens Dyb er der intet Gjærende og intet umiddelbart Værende: Alt er bestemt ved Tanken og Ordet, Alt bestaaer i uendelig Reflecteren, Dømmen og Slutten. I den uendelige Dommen og Slutten have de tre, for Indholdets absolute Totalitet gjældende Grundbestemmelser: det Enkelte, det Særskilte og det Almene, eller: Individet, Arten og Slægten, samme oprindelige Gyldighed. Det er den tænkende Subjectivitet nødvendigt at begribe sit Indhold i Form af det Almene, efterdi den selv er det absolut Almene; det er den anskuende Subjectivitet nødvendigt at opfatte sit Indhold i Form af det Særskilte, efterdi den i sin reflecterende Anskuen selv er det Særskilte; det er endelig den villende, ved Selvbestemmen mod Andet bestemte Subjectivitet nødvendigt, at den maa forestille sig sit Indhold articuleret til lutter bestemte Enkelte, efterdi den selv er det absolut Enkelte. Subjectivitetens Indhold er altsaa ikke udstykket, saa at en Deel for sig skulde falde under det Almene, en anden under det Særskilte, en tredie under det Enkelte; ved en saadan Udstykning vilde det Absolute, istedenfor at gaae op i det Relative, gaae tabt i lutter endelige Relativiteter. Men naar det hele udeelte, i sig udeelbare Indhold skal fordeles paa de tre Former, saa kan den Modsigelse, at det Udeelte dog underligger en Grunddeling, kun løses derved, at to af Formerne skifteviis danne Extremer, der have den tredie til Midte, altsaa kun ved et System af immanente Slutninger.

Under Systemets Udvikling gaaer nu Bevægelsen: deels fra det Enkelte gjennem det Særskilte til det Almene, hvorved Slutningsrækkerne blive kategoriske, og Eftertrykket falder paa \emph{Subsumtionen}; deels fra det Almene gjennem det Sær%
% --- p. 347 ---
\opage{347}skilte til det Enkelte, hvorved Slutningsrækkerne blive disjunctive, og Eftertrykket falder paa \emph{Disjunctionen}; endelig indsees, hvad det vil sige, at Bevægelsen paa een Gang gaaer i begge Retninger, saa at Slutningsrækkerne totaliseres, og det afgjørende Eftertryk falder paa \emph{Individuationen}.

\emph{Kategoriske Immanensslutninger: Subsumtionsproblemet.} I den formelle Logik ulmer en mærkelig Strid imellem Reflexionen og Dommen, en Strid, der netop angaaer det \emph{at subsumeres} under Noget og \emph{at være} Noget. Ifølge Dommen er Individet sin Art, Arten sin Slægt. Men ifølge Reflexionen kan Individet umulig være sin Art, Arten umulig være sin Slægt; thi Arten, der har større Omfang end Individet, er hvad det enkelte Individ ikke er; Slægten, der har større Omfang end Arten, er ved sit større Omfang, hvad Arten ikke er, og Arten, der har større Indhold end Slægten, ved sit større Indhold, hvad Slægten ikke er. Det Enkelte kan for den endelige Reflexion umulig være det Særskilte, saalidt som det Særskilte kan være det Almene.

Ihvorvel Reflexionen i den formelle Logik ellers er ivrig nok for at sikkre Dommen sin almindelige Form, at Copula \emph{er} ikke skal forvexles enten med Verbet \emph{er} eller med Realiteten \emph{er}, saa forbeholder den sig dog ikke desto mindre selv Ret til at afgjøre, hvad Mening der tør ligge i Dommens Copula. Den til sin typiske Form reducerede Dom siger ligefrem: det Enkelte er det Almene; dette er en Hest, Individet er sin Art; Hesten er en Tykhud, Arten er sin Slægt; men Reflexionen forsikkrer, at det umulig kan være meent saaledes. Det Enkelte er ikke det Almene, men Meningen er, at det subsumeres under det Almene; Individet er ikke Art, men Meningen er, at det subsumeres under Arten o.\ s.\ v.

Charakteristisk er det, at Dommen, naar den bogstavelig svarer til den af Logiken selv bestemte Form, ikke kan forstaaes bogstavelig. Hver Gang Dommen antager sit logiske Præg,
% --- p. 348 ---
\opage{348}bliver Omskrivningen: „\emph{subsumeres under}“ fortrængt af Copula „\emph{er}“; hver Gang Reflexionen skrider ind med sin Forklaring, bliver „\emph{er}“ omskrevet til „\emph{subsumeres under}“. Af denne Spænding fremgaaer da \emph{Subsumtionsproblemet}.

I den formelle Logik kan Subsumtionsproblemet ikke komme frem for lutter subjectiv Reflecteren, i den „speculative“ Logik kan det ikke komme frem for lutter objectiv Væren. Fra Endelighedens og de eensidige Begrebers Standpunkt er det ganske vist en Modsigelse, at Subjectet skulde være sit Prædicat. Her er jo ikke blot en Forskjel, men en \emph{endelig} Forskjel mellem det Enkelte, det Særskilte og det Almene, mellem Ting og Begreb, mellem Individ, Art og Slægt. For nu ret at udhæve den endelige Forskjel uden at opgive Eenheden, lægges der Eftertryk paa Subsumtionen. Anderledes fra Ideens Synspunkt. I Ideen er der vel ogsaa Forskjel, en gjennemgaaende Forskjel mellem Individet og Arten, mellem Arten og Slægten; men i den gjennemgaaende Forskjel er Endeligheden hævet. Individet er paa \emph{individuel} Maade, hvad Arten er paa \emph{særskilt}, og Slægten paa \emph{almindelig} Maade. Der er ikke andre eller flere Begrebsbestemmelser i Individet end der er i Arten: Identiteten af Individ og Art er absolut, Individet \emph{er} sin Art. Men Begrebsbestemmelserne, der i Individet alle ere individuelle, Enkeltmomenter i det Enkelte, ere i Arten alle særskilte, Reflexionsmomenter i det eiendommelige Tankebillede: Forskjellen mellem Individ og Art er absolut, Individet \emph{subsumeres} under sin Art.

Hvad dernæst Forholdet mellem Arten og Slægten angaaer, kunde det vel ved første Øiekast synes, som om Arten indeholdt flere og følgelig andre Bestemmelser end Slægten; men denne Synsmaade har kun Gyldighed for Abstractionen. Seet i Ideen har Arten samme Totalitet af Bestemmelser som Slægten; Forskjellen beroer alene paa Bestemmelses%
% --- p. 349 ---
\opage{349}maaden.\footnote{Jvfr. Grundideernes Logik. I. S.~218--20.}% note *)
{} I Modsætning til Slægten som det Almene er Arten det Særskilte, og derfor have Artsbestemmelserne \emph{Særskilthedens}, Slægtsbestemmelserne \emph{Almindelighedens} Form. Lade vi f.\ Ex.\ Kategorierne: Sviin, Tapir, Hest og Elephant, forestille Arterne, Kategorien: Tykhud, Slægten, da er det let at vise, at Arten umulig kan indeholde flere Bestemmelser end Slægten, aldenstund Slægten ligesaavel som Arten er, ikke et tomt Schema, men et relativt Absolutum. Hesten er ikke en Tykhud og saa noget Mere, men Tykhuden selv i en særegen Skikkelse. Hesten og Elephanten forholde sig da heller ikke til hinanden saaledes, at de skulde have en fuldkommen Lighed i det at være Tykhud og ved Siden heraf Ulighed i de Formbestemmelser, der faldt udenfor det at være Tykhud. Nei, Ligheden er gjennemgaaende, efterdi Uligheden er gjennemgaaende. Vel er der Forskjel imellem Elephantens og Hestens Fodform, men denne Forskjel falder netop i Slægten, just fordi den falder i Arten: Elephantfoden er Tykhudens Fod ligesaa vel som Hestefoden. Vel er Slægten forskjellig fra sine Arter derved, at Arterne ere forskjellige fra hinanden indbyrdes; men Slægten er ved denne Forskjellighed tillige identisk med sine Arter. Reflecteres her altsaa paa Identiteten, saa \emph{er} Arten sin Slægt; reflecteres her paa Forskjellen, saa maa Arten \emph{subsumeres under} Slægten.

Det dobbeltsidige Forhold mellem Individ, Art og Slægt finder sin logiske Udvikling og Stadfæstelse i den kategoriske Immanensslutning, som vi da ogsaa kunne kalde \emph{Subsumtionsslutningen}. Naar Dommen har sat abstract Eenhed og derved abstract Forskjel mellem det Enkelte og det Almene, saa indfører Slutningen det Særskilte som Mellemled, og forliger ved dette Mellemled Subsumtionen med Væren. Det gjør i denne Henseende en væsentlig Forskjel, om en Dom opfattes umiddelbart som Dom eller som Conclusum. Sige
% --- p. 350 ---
\opage{350}vi: „Hesten (Individet) er en Tykhud“, „Elephanten (Individet) er en Tykhud“, saa fristes vi til at opstille de to umiddelbare Domme som Præmisser, og uddrage den Slutning, at Hesten er hvad Elephanten er, at Hesten er Elephant, og Elephanten Hest. Anderledes derimod, naar enhver af de anførte Domme opfattes som Slutningssætning, som Conclusum. Hesten (Individet) er da ikke uden videre en Tykhud, men subsumeres under Slægten Tykhud; Individet Hest er ikke uden videre Arten Hest, men subsumeres under Arten. Dersom Individet umiddelbart subsumeredes under Slægten, vilde Subsumtionen eensidig fortrænge Væren; men idet Arten træder imellem saaledes, at Individet subsumeres under Arten, og Arten under Slægten, dukker Væren frem af Subsumtionen, thi Individet \emph{er}, som vi have seet, væsentlig sin Art, Arten væsentlig sin Slægt.

Den umiddelbare Dom: „Hesten (Individet) er en Tykhud“, er da, naar vi see nærmere til, en fra sine Præmisser løsreven Følgesætning. I reflecteret Eenhed med Præmisserne vil Dommen som Følgesætning lyde saaledes: dette Individ er som Hest en Tykhud; altsaa, hvad de øvrige Domme angaaer, dette Individ er som Sviin en Tykhud; dette Individ er som Elephant en Tykhud, o.\ s.\ v. Men paa denne Maade er Artsforskjellen lagt ind i Slægten. Tykhuden er ikke uden videre Tykhud; thi Tykhuden Hest er forskjellig fra Tykhuden Elephant, fra Tykhuden Sviin o.\ s.\ v. Slægten er altsaa kun identisk med sig selv derved, at den sætter og afsætter Arten, ligesom Arten kun er identisk med sig selv derved, at den sætter og afsætter Individet; det Almene er kun identisk med sig selv derved, at det sætter og afsætter det Særskilte, ligesom det Særskilte kun er identisk med sig selv ved at sætte og afsætte det Enkelte. Men at Slægten afsætter Arten, og Arten igjen Individet; at det Almene afsætter det Særskilte, og det Særskilte igjen det Enkelte, er logisk begrundet i Disjunctionen.

% --- p. 351 ---
\opage{351}\emph{Disjunctive Immanensslutninger: Disjunctionsproblemet.} Den formelle Logiks Schema: \textit{A} er enten \textit{B} eller \textit{C}, nu er \textit{A B}, altsaa ikke \textit{C}, eller: nu er \textit{A} ikke \textit{C}, altsaa er \textit{A B}, giver et for Erkjendelsen tarveligt Udbytte, saalænge det fastholdes i tom Abstraction. Som Form for Ideen, det absolute Indhold, faaer den disjunctive Slutning derimod en uendelig Betydning. Er Subjectet (\textit{A}) det Almeen-Absolute, saa maa Prædicatets Led (\textit{B} og \textit{C}) være Særskilt-Absolute. Særskiltheden angives udtrykkelig ved Disjunctionen: enten \textit{B} eller \textit{C}, og falder forsaavidt udenfor \textit{A}, thi \textit{A} er umiddelbart som \textit{A} ikke forskjelligt fra sig selv, ikke et i sig Særskilt: \textit{A} er kun, hvad \textit{A} er; i \textit{A} er det Særskilte ikke. Og dog falder, ifølge den disjunctive Dom, Disjunctionen netop indenfor \textit{A}. Det er jo netop \textit{A} som \textit{A}, der er enten \textit{B} eller \textit{C}. Her er Dommen igjen i Modsigelse med Reflexionen. Idet Dommen udtrykkelig siger, at \textit{A} er enten \textit{B} eller \textit{C}, forudsætter den at \textit{A} i Grunden er baade \textit{B} og \textit{C}, nemlig i det ene Tilfælde \textit{B}, i det andet \textit{C}. Ængstet af Modsigelsen vil Reflexionen nu gjerne forvandle den disjunctive Dom til en blot Mulighedsdom, en Dom, der determinerer Muligheden ved at begrændse og angive dens Alternativer. Udenfor Ideen, i de tomme Formers ideeløse Mangfoldighed, kan det nu rigtignok saa omtrent være det Samme, enten man tager den disjunctive Dom for en Omskrivning af den problematiske eller opstiller den for sig, thi det er vist, at Disjunctionen selv forstaaer man ikke. Seet i Ideen, er Disjunctionen absolut, og den disjunctive Doms Forskjel fra den problematiske netop kjendelig paa Slutningen. Særkjendet for den disjunctive Dom er just dette, at den efter sin oprindelige Betydning danner Præmis for ligesaa mange positiv-negative Slutninger, som Prædicatet indeholder alternerende Led, et uimodsigeligt Beviis for, at Disjunctionen er kategorisk. Opfattes Dommen blot som Udtryk for Muligheden, saa kan her, hvormange Led Prædicatet end indeholder,
% --- p. 352 ---
\opage{352}kun haves Præmis for een eneste positiv Slutning. \textit{A} er enten \textit{B} eller \textit{C} eller \textit{D}, det vil sige: \textit{A} er maaskee \textit{B}, maaskee \textit{C}, maaskee \textit{D}. Hvorfra komme de tre Alternativer? Alene fra Uvisheden! Alternativerne ere kun subjective, ikke objective. I objectiv Forstand er her kun et eneste Led; naar dette er sat, saa ere de øvrige udelukkede, deres Intethed er bleven aabenbar; ere omvendt de øvrige udelukkede, saa maa det resterende Led sættes; hermed er Sagen afgjort. Det Positive i Slutningen er da med det Samme skilt fra det Negative saaledes, at der til Formernes Formerelse fremkomme to særskilte Slutninger, en positiv ved den saakaldte \textit{modus tollendo ponens}, og en negativ ved \textit{modus ponendo tollens}. Anderledes i det Absolute; her er Disjunctionen, hvad Dommen siger, objectiv, d.\ v.\ s.\ logisk objectiv. Det i sig udeelte, i endelig Forstand udeelbare Absolute er ved den principielle Disjunction adskilt i særegne, hinanden udelukkende Absolute: enten \textit{B} eller \textit{C} eller \textit{D}; ethvert af de særskilte Led har samme objective Gyldighed. Positionen af ethvert Led er reflecteret i Negationen af de øvrige, og omvendt; derfor maa hvert Led fattes ved en egen Position og fordre en egen Slutning. Da Positionen af hvert enkelt Led reflecterer sig i Negationen af de øvrige, er hver Slutning for sig paa een Gang baade negativ og positiv: \textit{modus ponendo tollens} reflecterer sig i \textit{modus tollendo ponens}, og omvendt. I Positionen \textit{A} er \textit{B} reflecterer sig Negationen, altsaa er \textit{A} hverken \textit{C} eller \textit{D}; omvendt, i Negationen: \textit{A} er hverken \textit{C} eller \textit{D}, reflecterer sig Positionen: altsaa er \textit{A B}.

Og hvad er saa disse Slutningers væsentlige Resultat? lutter kategoriske Domme: \textit{A} er \textit{B}, \textit{A} er \textit{C}, \textit{A} er \textit{D}; disse paa disjunctive Slutninger hvilende Domme ere nu Udtryk for, at det Almene er et Særskilt.

Med Hensyn paa Subsumtionen søgte vi en Begrundelse af den Dom: det Særskilte er et Almeent, med Hensyn paa Disjunctionen have vi fundet lige det Omvendte, at nemlig det
% --- p. 353 ---
\opage{353}Almene er et Særskilt; Subsumtion og Disjunction falde endnu ud fra hinanden, et Tegn paa, at Problemerne endnu ikke ere fuldstændig løste. Problemernes Løsning vanskeliggjøres derved, at Slutninger og almindelige Reflexionsbevægelser kun altfor let forvexles. Sige vi: „Alt er almindeligt, det Almindelige er Alt“, da gjør det en væsentlig Forskjel, om det er en Reflexionssætning eller et System af Slutninger, hvorpaa her skal tænkes. Som Reflexionssætning er det en Omskrivning af Grundens første Sætning: „Alt er i Grunden Eet“, vi have for os, og den til Sætningen svarende Bevægelse er Reductionen. Naar Tilværelsens ydre Mangfoldighed føres tilbage til Grundens Eenhed, er Grundeenheden det concret Almene, men som logisk Function vil Reductionen herved tillige betegne, at det Mangfoldige opløses i Grundeenheden, at det Tilværende gaaer tilgrunde. Anderledes derimod, naar Reflexionsbevægelsen er udviklet til Slutning. Istedenfor Reductionen gjælder da Subsumtionen. Subsumtionen opløser Intet; den bevarer Alt: det Enkelte gaaer ikke tilgrunde i det Særskilte, for at det Særskilte igjen kan gaae tilgrunde i det Almene; nei, det Enkelte bevares i det Særskilte derved, at det Særskilte bevares i det Almene. Slutningen er væsentlig Sammenslutning, og Sammenslutningen et kategorisk Udtryk for, at Extremerne med deres Midte fuldstændig bevares. Forsaavidt Slutningen: „Alt er almindeligt“, udtrykkelig opfattes som Slutning, er Slutningen selv et System af kategoriske Immanensslutninger, ja et System af Systemer i det Uendelige. Bevægelsen fra det Enkeltes Mangfoldighed gjennem det Særskilte til det Almene er en Bevægelse fra Peripherie til Centrum, en Slutningsbevægelse, der i den absolute Viden kan gjennemløbe utallige Mellemstadier. Blot en Sætning som: „dette (Individet Hest) er et levende Væsen“, indeholder allerede, naar den skal reflecteres som Følgesætning i sin tilsvarende Slutning, en Mængde væsentlige Mellemled mellem Subject og Prædicat: „dette er
% --- p. 354 ---
\opage{354}en Hest, Hesten er en Tykhud, Tykhuden er et Hovpattedyr, Hovpattedyret er et Pattedyr, Pattedyret er et Hvirveldyr, Hvirveldyret er et Dyr, Dyret er et besjælet Væsen; det besjælede Væsen er et levende Væsen; altsaa: dette Individ er et levende Væsen“. Siger jeg: „dette (Skriverbordet) er et Tilværende,“ hvilken Afstand synes her da ikke at være imellem det Individuelle i Subjectet og det Almene i Prædicatet, naar vi betænke, at \emph{dette} (Huset) og \emph{dette} (Stjernen) og \emph{dette} (Keiseren af China), kort sagt, alle Individer, Ting og Egenskaber, maae subsumeres under det Tilværende. Sættes nu Tilværen som det Almene, det Tilværende derimod som Betegnelse for ethvert bestemt Individ, et hvilketsomhelst Enkelt, da indsees let, at Veien fra det Tilværende (Subjectet) til Tilværen (Prædicatet) er paa een Gang uendelig lang og dog uendelig kort; en lige og dog uendelig forgrenet Linie.

Den tænkende Subjectivitet, der har sit hele Indhold i det Almindelige, seer alt det Enkelte i eet Blik, fordi den i eet Blik seer det uendelige System af Slutninger, hvori hvert Enkelt gjennem sine Mellemled gaaer sammen med det Almene.

Men nu vilde det jo være umuligt, at et uendeligt System af Slutninger skulde kunne tænkes uden et tilsvarende System af Mellemled. Hvad er nu Mellemleddet i kategorisk Bestemthed? Altid et Særskilt! Enten vi vælge Schemaet: \textit{E}---\textit{S}---\textit{A} eller \textit{S}---\textit{E}---\textit{A} eller \textit{E}---\textit{A}---\textit{S}, Midten er i Characteer af Mellemled altid det Særskilte, thi den reflecterer jo i sin Forskjel fra Extremerne Extremernes relative Eenhed. Men derved skeer det da ogsaa, at Extremerne i deres relative Eenhed med Midten komme til at reflectere deres indbyrdes Forskjel, at Extremerne altsaa, med andre Ord, formedelst det Særskilte i Midten selv blive særskilte. Men naar Extremerne blive særskilte ved at være Extremer, efterdi Midten er særskilt ved at være Midte, saa blive da paa denne Maade alle Indholdsformerne særskilte: Alt er et Særskilt, det Særskilte er
% --- p. 355 ---
\opage{355}Alt. Forvexle vi nu Slutningen med en Reflexionsbevægelse, da er Reflexionsbevægelsen en immanent, men tvetydig Omsætning af Reduction i Deduction, og omvendt. Ved Reductionen opløses Alt, det Enkelte med det Særskilte, i det Almindelige; men ved at optage alle Bestemmelser af det Enkelte og Særskilte i sig, maa det Almindelige selv gaae i Eet med Totaliteten af Særskilte. Men ligesom det Almene ikke har anden Bestaaen end dets Reflexion i de ved Opløsning af de Enkelte fremkomne Særskilte, saaledes have disse Særskilte da heller ingen anden Bestaaen end den ved Opløsningen bestemte Reflexion i det Almene. Grundreflexionen er en gjensidig Opløsning. Det Almindeliges Eenhed opløses og gaaer tilgrunde i Mangfoldigheden af det Særskilte, idet denne Mangfoldighed opløses og gaaer til Grunde i Eenheden. Hertil svarer saa Grundens anden Hovedsætning: „Alt er i Grunden forskjelligt.“ Ved Forskjelligheden falder Eftertrykket paa Deductionen, og ved Deductionen fremkommer det Særskilte; men ved Eenheden falder Eftertrykket igjen paa Reductionen, og det Særskilte drukner i Grunden.

Kan nu den rene Reflexion ikke drive det videre end til en Udvikling af den Sætning, at Alt i Grunden er et Særskilt, fordi Alt i Grunden er et Almindeligt, og omvendt, da er det jo med det Samme indlysende, at den fuldstændige Bestemmelse af det Særskilte maa beroe paa Slutningen. I Slutningen, d.\ v.\ s.\ i et uendeligt System af Slutninger, seer den anskuende-reflecterende Subjectivitet overalt det Særskilte: det Almindelige begrunder sig selv, idet det begrunder det Særskilte, og formedelst det Særskilte tillige det Enkelte. Den Bevægelse, der saaledes gaaer fra det Almindelige gjennem det Særskilte til det Enkelte, en Slutningsbevægelse fra Centrum til Peripherie, er bestemt ved Disjunctionen, altsaa ved et System af Systemer af disjunctive Slutninger. Et Exempel: „Det Tilværende er i sin Almindelighed baade et Levende og et Ikkelevende (\textit{A} er baade \textit{B} og \textit{C}), men i sin Særskilt%
% --- p. 356 ---
\opage{356}hed enten et Levende eller et Ikkelevende (\textit{A} er enten \textit{B} eller \textit{C}). Det Ikkelevende er \emph{sat}, altsaa er det Levende \emph{udelukket}; det Levende er \emph{sat}, altsaa er det Ikkelevende \emph{udelukket}. Det Levende og det Ikkelevende ere begge ved gjensidig Udelukkelse og tilsvarende Vexelbestemmelse satte som Særskilte. Formedelst den gjennemførte Disjunction er det Almene sluttet og afsluttet i sig saaledes, at det Særskilte med det Samme enkeltviis er sluttet og afsluttet. % PRINTER: the quotation opened at „Det Tilværende (p. 355) is never closed

Da nu Disjunctionen i ethvert af de særskilte Systemer maa fortsættes lige ned til Extremet af det Enkelte (det levende Væsen er enten Plante eller Dyr, Dyret enten Bløddyr eller Leddyr eller Hvirveldyr, Hvirveldyret enten o.\ s.\ v.\ ned til Individerne), saa maae vi erkjende, at i Sætninger som: „det Tilværende er (i nærværende Tilfælde) enten en Hest eller en Oxe“, maa Veien fra Subjectet (det Tilværende) til Prædicatet (Hest) være for Dommens Skyld uendelig kort, for Slutningens Skyld uendelig lang.

\emph{Slutningssystemernes Totalisering: Individuationsproblemet.} Hvorfra har det Individuelle vel sin første og egentlige Oprindelse? Realismen (\textit{univers. ante rem}) svarer: fra det Almene, og anstrenger sig med at pidske Begreberne op til en saadan Concentration af det Særskilte i det Enkelte, at det for Reflexionen kan tage sig ud, som om det Individuelle i sin Umiddelbarhed virkelig var en intellectuel Fortætning af det Almene. I Modsætning til Realismen gaaer Nominalismen (\textit{univers. post rem}) ud fra den Forudsætning, at det i sig Individuelle, langt fra at have nogen egentlig Oprindelse, meget mere selv er det Oprindelige, og giver det saa ved Abstractionens Hjælp Udseende af, at det er det Almene, der er et Afledet. Realisterne misbruge \emph{Disjunctionen}, Nominalisterne \emph{Subsumtionen}: hine fortabe sig i en Individualisering, der aldrig naaer det Individuelle selv; disse fastholde et Individuelt, der maa udelukke enhver Individualisering. Fra begge Sider miskjendes da ogsaa det
% --- p. 357 ---
\opage{357}Særskilte; thi det oprindelig Særskilte forudsætter, at det Almene ligesaa principielt bestemmes ved det Individuelle som dette ved hiint, at det Almene og det Individuelle altsaa have samme Oprindelighed.

Da alle tre Grundformer ere lige oprindelige, og Indholdet i enhver af dem identisk med sin Form, maa det Almene, det Særskilte og det Individuelle have samme absolut ideelle Realitet. Her kan da ligesaa lidt være Tale om at universalisere det Individuelle som om at individualisere det Almene. Derimod følger det af Relationsprincipet, at de tre ideelle Realiteter absolut maae bestaae formedelst hinanden indbyrdes: det relativt Almene bestaaer formedelst det Særskilte i oprindelig Eenhed med det Enkelte; det Særskilte bestaaer formedelst det Enkelte i oprindelig Eenhed med det Almene, og det Enkelte formedelst det Almene i oprindelig Eenhed med det Særskilte. Den fuldt udviklede Relation er altsaa Slutning. Men en Bestaaen i Slutning er en Bestaaen i ideel Energie, en Bestaaen i uendelig Subsumtion og Disjunction: den articulerende Vexelbestemmelse af Subsumtion og Disjunction er Individuation.

Det Individuelle bliver ikke til ved endelig Individualisering; det \emph{bliver ikke til}, thi det \emph{er}. Og dog kan det Individuelle ikke bestaae uden Individuation, thi Individuationen er den fuldendte Slutten, og det Individuelle bestaaer i evig Slutning. Her er i Bestemmelsen af det relativt Absolute ikke en individualiserende Virksomhed ved Siden af en universaliserende, thi Individuationen, der giver det Individuelle Bestaaen, giver med det Samme det Universelle en fra det Individuelle forskjellig og just derved individuel Væren; kun en Raisonnementsphilosophie, der pjasker i Ord, kan bilde sig ind, at den virkelig taler om Noget, naar den i al Almindelighed taler om det „Altuniversaliserende“, „Altindividualiserende“ o.\ s.\ v.

Er nu enhver af Grundformerne formedelst Individua%
% --- p. 358 ---
\opage{358}tionen en i sig udviklet Totalitet, der omfatter hele Subjectivitetens Indhold, saa maae den totaliserede Slutnings Extremer og Midte nødvendigviis udgjøre \emph{tre immanente Regioner}, og hver Region for sig i sin Form omfatte det hele uendelige System af kategoriske og disjunctive Slutninger. Til Regionen af det Almene svarer da Subjectiviteten som tænkende, til Regionen af det Særskilte Subjectiviteten som reflecterende-anskuende, til Regionen af det Enkelte Subjectiviteten som villende. Saalidt som den absolute Subjectivitet kan siges at opgive sin Eenhed, fordi dens Modsætning af Tænken, Reflexionsanskuen og Villen adskiller sig i tre Phaser, ligesaalidt kunne de tre Immanensregioner bevirke en Spaltning af det oprindelige Fornuftindhold, som om Indholdet var lagdeelt, og Idealverdenen udstykket i uligeartede Provindser. Saa gjennemgaaende Modsætningen er imellem Tænken, reflecterende Anskuen og definitivt bestemmende Villen, ligesaa gjennemgaaende er ogsaa Modsætningen af den for Indholdet gjældende Totalslutning, eftersom den foregaaer i det Almenes, i det Særskiltes eller i det Enkeltes Region.

Fra Overblikkets Synspunkt vil Regionen af det Almene nærmest kunne henføres til, hvad Neoplatonismen paa den ene Side har opfattet som det rene Lys (τὸ φῶς), paa den anden Side som den guddommelige Fornuft (ὁ νοῦς). Men uagtet de Skarpsindigste blandt Neoplatonikerne, Plotin f.\ Ex., bestemte det Absolute selv som den uendelige Energie\footnote{Εἰ οὖν τελειότερον ἡ ἐνέργεια τῆς οὐσίας, τελειότατον δὲ τὸ πρῶτον, πρῶτον ἂν ἐνέργεια εἴη· ἐνερνήσας οὖν ἤδη ἐστὶ τοῦτο, καὶ % PRINTER: printed "ἐνερνήσας", sense "ἐνεργήσας"
οὐκ ἔστιν ὡς πρὶν γενέσθαι ἦν· ὅτε γὰρ οὐκ ἦν πρὶν γενέσθαι, ἀλλ᾽ ἤδη πᾶς ἦν. Ἐνέργεια δὴ οὐ δουλεύουσα οὐσίᾳ καθαρῶς ἐστιν ἐλευθέρα· καὶ οὕτως αὐτὸς παρ᾽ αὐτοῦ αὐτός. Πλωτῖνος. Jvfr. Kirchner. Die Philosophie des Plotin. Halle 1854. S.~35--46.}% note *)
, saa vare de dog ikke istand til at gjennemskue Forholdet imellem „det
% --- p. 359 ---
\opage{359}Ene“ og „Fornuften“, efterdi de endnu ikke vare istand til at indsee Fornuftvirksomhedens Identitet med den immanente Slutning\footnote{„Aus der reinen Einheit geht zunächst der Geist des Universums hervor, dem alle einzelnen Geister entspringen \dots\ Das erste Sein (οὐσία) und die erste Thätigkeit (ἐνέργεια) lebt er in ruhender That (ἐνέργεια ἑστῶσα), in sich ruhend, aber um jenes sich bewegend (ἐνεργῶν). Man darf die göttliche Vernunft nicht mit der menschlichen verwechseln, die einem äußeren Sein gegenübersteht. Es ist von keiner Vernunft die Rede, die durch Berechnungen und Schlüsse zur Wahrheit gelangt, sondern von der, die Alles hat und Alles ist und ewig bei sich selbst bleibt“. Kirchner. Anfr. Skr. S.~46.}% note *)
.

Henføre vi „den første Væren“ (οὐσία), den „i sig hvilende“ Grundenergie (ἐνέργεια ἑστῶσα) under Slutningsregionen af det Almene, saa finde vi let Udtrykket for den absolute, over alle særskilte Bestemmelser ophøiede Eenhed i den tænkende Selvhed, den i al Tænken med sig identiske Subjectivitet. Det sit Indhold tænkende Selv tænker ved Ordets negative Functioner det Enkelte og Individuelle ligesaa skarpt og tydeligt som enten det Særskilte eller det Almene. Det tænkende Selv tænker ved sit negativt bestemmende Ord Totaliteten af alle Slutningssystemer gjennem alle Individuationer i alle Retninger, uden at Mangfoldigheden i mindste Maade enten smitter eller svækker den oprindelige Eenhed. Det er just den til Slutningsenergien svarende Gjennemsigtighed af Led og Slutninger, som i denne Region gjør Fornuftriget til et evigt Lysrige.

Henføre vi dernæst „den Udfoldning af den høieste Eenhed“, hvori Neoplatonikerne sætte „Ideernes“ Oprindelse, under Regionen af det Særskilte, da blive vi istand til at ombytte de forblommede, figurlige og halvmystiske Antydninger med rene logiske Bestemmelser. For det anskuende-reflecterende Blik fremstiller enhver Bestemmelse sig med Nødvendighed
% --- p. 360 ---
\opage{360}som et Særskilt. At tale om det Særskiltes Udtræden af det Almene, Farvestraalernes Udskillelse af det rene Lys, det Eensformiges Overgang i det Mangeformige (τὸ ποικίλον) er at tale om Tænkningens Overgang til reflecterende Anskuen, en Overgang, der skeer ved Ordet, idet Ordet ombytter de negative Functioner med positive. Den Anskuelsens Umiddelbarhed, som herved fremkommer, indfører just i Reflexionen det Skjær af Endelighed, hvori Neoplatonikerne saae en charakteristisk Forskjel mellem et intellectuelt Stof (ὕλη νοητή) og en Mangfoldighed af vexlende Former\footnote{Blandt flere af de hos Plotin forekommende Steder, som Kirchner anfører, er især følgende (Anfr. Skr. S.~48) charakteristisk: Εἰ οὖν πολλὰ τά εἴδη, κοινὸν μέν τι ἐν αὐτοῖς ἀνάγκη εἶναι, καὶ δὴ καὶ ἴδιον, ᾧ διαφέρει ἄλλο ἄλλου. % PRINTER: printed "τά" (acute before a word), sense "τὰ"; UNSURE: comma or full stop after εἴδη
Τοῦτο δὲ τὸ ἴδιον καὶ ἡ διαφορὰ ἡ χωρίζουσα ἡ οἰκεία ἐστὶ μορφὴ· εἰ δὲ μορφὴ, ἔστι καὶ τὸ μορφούμενον, περὶ ὅ ἡ διαφορὰ. Ἔστιν ἄρα καὶ ὕλη ἡ τὴν μορφὴν δεχομένη καὶ ἀεί τὸ ὑποκείμενον.}% note *)
. % PRINTER: printed "ἀεί" (acute before a word) in the note, sense "ἀεὶ"

Under Særskilthedens Region høre nu Slægter og Arter (γένη og εἴδη). Dette er dog ikke at forstaae, som om det var en ubestemt Mangfoldighed af „Syner“ og „plastiske Forbilleder“, der uvilkaarlig dukkede op af Reflexionernes evigt bølgende Hav; nei Individuationens bestemte Forskjel mellem Slægt og Art er i alle Slægter og i alle Arter en udtrykkelig Slutningsbestemmelse. Ved Slutningsbestemmelsen er Individet overalt med i Arten, efterdi Arten er med i Slægten; det særegne Individ (det individuelt Særskilte) slutter sig formedelst Arten (det i særlig Forstand Særskilte) sammen med Slægten (det særskilt Almene).

At Neoplatonikerne med Plato antage Ideerne (εἴδη) for „Urbilleder“, og at Plotin med Aristoteles anseer ethvert εἶδος for en individuel Væsenhed (οὐσία), kan i en vis Forstand bifaldes; men Grundfeilen er, at de individuelle Væsenheder opfattes paa en blot umiddelbar og derfor overfladisk
% --- p. 361 ---
\opage{361}Maade. Opfattelsen er lige overfladisk, enten vi med Plato lade disse Væsener hvile i evig Ubetingethed, eller med Aristoteles fortabe sig i en Formen af Stoffet. For at magte sit Stof maa det formende Væsen tillige oversvæve Stoffet. Det sit Stof oversvævende Enkeltvæsen kan ikke bestemmes individuelt uden at bestemmes særskilt; kan ikke bestemmes særskilt som Art uden at bestemmes særskilt som Slægt; kan ikke paa een Gang bestemmes som Individ, Art og Slægt uden ved en immanent Slutning.

Henføre vi endelig hele Subjectivitetens Indhold under Enkelthedsregionen, Regionen af det Individuelle, da maa det ret blive os indlysende, hvad Individuationen oprindelig har at betyde. Den neoplatoniske Reflexion kommer, hvad Principet angaaer, ikke ud over εἶδος; men εἶδος tilhører væsentlig Særskiltheden: Neoplatonismen naaer ikke den punktuelle Bestemthed, men forbliver i Speculationen. Det er et resolut Skridt af Herbart, at han, for at komme ud af Speculationen uden at hilde sig i Empiriens metaphysiske Modsigelser, lægger „Realerne“ til Grund, og ved denne Punktualitet vil sikkre sig det absolut Individuelle. Men som intellectuelle Atomer ere Realerne selv behæftede med altfor skrigende Modsigelser, til at de skulle kunne gjøre Immanensslutningen overflødig. Leibniz med sine Monader er, hvad Individuationen angaaer, aabenbart den, der kommer det Sande nærmest. Monaden er virkelig et Individ, der, istedenfor at udelukke Arten og Slægten, tvertimod har trukket Slægten og Arten ind i sig: Enkelthedsregionen kan forsaavidt gjerne opfattes som en Region af Monader. Den villende Subjectivitet er ved den definitive Modsætning af det bestemte Selv og det bestemte Andet \emph{exact forestillende}; Forestillingerne, de absolute Positioners ideelle Gjenstande, tage sig unegtelig ud som Monader.

Men som rene Forestillingsproducter uden Arter og Slægter kunne Monaderne hverken tænkes eller anskues. Skal
% --- p. 362 ---
\opage{362}Monaden tænkes, maa den absolut bestemmes som det Almene. Monaden er da efter sit Begreb een i alle; det Monadiske er i alle Monader sig selv liigt. Disjunctionen: A er enten B eller C eller D, maa forvandles til: A er enten A\textsubscript{1} eller A\textsubscript{2} eller A\textsubscript{3}, thi A\textsubscript{1} er som Monade væsentlig identisk med A\textsubscript{2} og med A\textsubscript{3}. Disjunctionen \emph{er tænkt}; men da Forskjellen bliver ved det Negative, kan Særskiltheden ikke træde særlig frem. At den disjunctive Dom bliver Præmis for en Slutning, at A ved Udelukkelse af A\textsubscript{2} og A\textsubscript{3} udtrykkelig bliver sat som A\textsubscript{1}; ved Udelukkelse af A\textsubscript{1} og A\textsubscript{3} som A\textsubscript{2} og ved Udelukkelse af A\textsubscript{1} og A\textsubscript{2} som A\textsubscript{3}; at altsaa baade A\textsubscript{1}, A\textsubscript{2} og A\textsubscript{3} kunne subsumeres under A, viser imidlertid, at Subsumtionen ligeledes \emph{er tænkt}. Den tænkte Subsumtion er en ligesaa negativ Bestemmen som den tænkte Disjunction; det Individuelle er en Negation af det Almene, fordi det Almene er en Negation af det Individuelle. Men naar alle monadiske Bestemmelser ere negative, efterdi de ere tænkte og som tænkte væsentlig almene, saa er den tænkte Monade selv i alle Individuationer negativ: \emph{Monaden er det Almene; det Almene er Monade}.

Gjøres Monaden til Gjenstand for reflecterende Anskuen, saa træder det Særskilte særlig i Kraft. Disjunctionen: A er enten B eller C eller D, forvandler sig til Udtryk for en Articulation af Særskiltheden: S (Monaden i sin Særskilthed) er enten S\textsubscript{1} eller S\textsubscript{2} eller S\textsubscript{3}. Det i Særskiltheden Særskilte eller den udtrykkelige Forskjel mellem S\textsubscript{1}, S\textsubscript{2} og S\textsubscript{3} giver nu tillige Subsumtionen sin særegne Betydning: S, det almindeligt Særskilte, hvorunder S\textsubscript{1}, S\textsubscript{2} og S\textsubscript{3}, de særligt Særskilte, subsumeres, er bestemt som Slægt: \emph{Monaden er Slægt, Slægten er Monade}; de særligt Særskilte ere med det Samme bestemte som Arter: \emph{Monaden er Art, Arten er Monade}.

Bliver endelig Monaden ved oprindelig Forestillen udtrykkelig bestemt monadisk, da er det en Selvfølge, at Mo%
% --- p. 363 ---
\opage{363}naden som Individ kun har Bestaaen ved Monaden som Art og Monaden som Slægt. Den absolute Forestillen forudsætter nemlig den reflecterende Anskuen med samme Nødvendighed, som den forudsætter den almene Tænkning. Dersom det er nødvendigt, at Monaden forestilles, da er det ogsaa nødvendigt, at den reflecteres; altsaa: dersom det er nødvendigt, at den opfattes som Individ, da er det ogsaa nødvendigt, at den opfattes som Art og Slægt. Individ, Art og Slægt ere monadisk forskjellige og monadisk identiske i een og samme Slutning. Men ligesom den monadiske Forskjellighed holder sig i Extremet af det Individuelle formedelst Forestillingsacterne, den villende Subjectivitets absolute Positioner, saaledes holder den monadiske Identitet sig i Extremet af det Almene formedelst Tænkningen, Systemet af Subjectivitetens negative Functioner.

Monaden som det Individuelle og Monaden som det Almene udgjøre altsaa to Regioner; men de to Regioner ere Extremer, der have Monaden i Særskilthedens Form til Midte. Sandheden, den gjennemførte Individuation, er ligesaalidt i Extremerne for sig som i den isolerede Midte, ligesaalidt i det Individuelle for sig som i det Særskilte eller i det Almene; Sandheden i Subjectivitetens Indhold er identisk med Sandheden i Subjectiviteten selv, thi Sandheden er en Slutning, en i Slutningssystemer totaliseret Slutning.

3) \emph{Den immanente Slutning som Form for den frie Nødvendighed.} I Individuationen har Vexelbestemmelsen af Selvhed og Andethed fundet et betegnende Udtryk. Al individuel Bestemmen er oprindelig Selvbestemmen; men al Selvbestemmen er igjen en Bestemmen af Andet ved Andet. Reflecteres her nu paa Selvheden, saa er Individualbestemmelsen et Udtryk for \emph{Friheden}; reflecteres her paa Andetheden, er Individualbestemmelsen et Udtryk for \emph{Nødvendigheden}. Ligesom Friheden væsentlig svarer til Subjectiviteten,
% --- p. 364 ---
\opage{364}fordi den afspeiler Selvheden, saaledes svarer Nødvendigheden til Objectiviteten, fordi den svarer til Andetheden.

Forsaavidt her sees paa det i Tanken, Dømmen og Slutten Subjective, lægges her Vægt paa Selvbestemmelsen, og derved paa Friheden. Forsaavidt her sees paa Tankens, Dommens og Slutningens Indhold, altsaa paa det, hvorved Tanke, Dom og Slutning væsentlig bestemmes, sees her paa det Objective, altsaa paa Nødvendigheden. Saalænge den Nødvendighed, hvorved den tænkende Subjectivitet oprindelig bestemmes, og den Frihed, hvormed Subjectiviteten i al Tænken bestemmer sig selv, ikke ere erkjendte i deres principielle Udspring, vil Forholdet imellem Guddommen og Ideerne ikke kunne opklares.

Naar Plato i æsthetisk Naivetet snart taler om, at Gud har skabt Ideerne, snart igjen lader den skabende Gud see hen til Ideerne som til evige Forbilleder, saa er ved det første Udtryk Friheden, ved det sidste Nødvendigheden antydet;\footnote{Plato, eller rettere den platoniske Sokrates, udtaler sig undertiden saaledes, at Gud maa antages at have skabt Ideerne. Ζωγράφος δή, κλινοποιός, θεός, τρεῖς οὗτοι ἐπιστάται τρισὶν εἴδεσι κλινῶν. Ὁ μὲν δὴ θεός, εἴτε οὐκ ἐβούλετο, εἴτε τις ἀνάγκη ἐπῆν μὴ πλέον ἤ μίαν ἐν τῇ φύσει ἀπερνάσασθαι αὐτὸν κλίνην, οὕτως ἐποίησε μίαν μόνον αὐτὴν ἐκείνην ὃ ἔστι κλίνη· \textit{Platonis Dialogi Lipsiæ MDCCCLX.} Πολιτειασ ί, S.~290. Derimod i Τιμαιος (anfr. Udgave S.~333): εἰ μὲν δὴ καλός ἐστιν ὅδε ὁ κόσμος ὅ τε δημιουργὸς ἀγαθός, δῆλον ὡς πρὸς τὸ ἀϊδιον ἔβλεπεν·}% note *)
% UNSURE: τρεῖς or τρεὶς (accent on the scan is unclear); ἤ or ἢ
% PRINTER: printed "ἀπερνάσασθαι", sense "ἀπεργάσασθαι"
% PRINTER: printed "Πολιτειασ" (medial sigma, no accent), sense "Πολιτείας"
% PRINTER: printed "Τιμαιος" (no accent), sense "Τίμαιος"
{} men naar den høiere Eenhed af Nødvendighed og Frihed forklares derhen, at Friheden overtaler Nødvendigheden, da svigter det Logiske. Plato har ved sine naive Yttringer antydet et Problem, paa hvis Løsning Philosophien først kan begynde at arbeide, naar den har opfattet det i hele dets Dybde. I Middelalderen fik Problemet en særlig Betydning i Scholastikernes Strid om Forholdet imellem det Godes Idee og den
% --- p. 365 ---
\opage{365}guddommelige Villie. Ifølge den thomistiske Sætning: \textit{Deus vult, quia bonum est}, er det Ideen, der bestemmer Villien, det Objective, der bevæger Subjectiviteten, Nødvendigheden, der er Frihedens Grund; ifølge den schotistiske Sætning: \textit{Bonum est, quia Deus vult}, er det omvendt Villien, der bestemmer Ideen, Subjectiviteten, der behersker det Objective, Friheden, der opløser Nødvendigheden. Hvorledes Problemet endnu i den allernyeste Tid bevæger sig igjennem Ethik og Religionsphilosophie, sees blandt Andet af den Maade, hvorpaa Feuerbach har stillet Idee og Gud imod hinanden. Istedenfor at gjøre f.~Ex. Retfærdighedens Idee til Prædicat for Gud som Subject, burde man, efter Feuerbachs Anviisning, gjøre det Guddommelige til Prædicat og Retfærdighedsideen til Subject; istedenfor den religiøse Dom: „Gud er retfærdig“, burde man sætte den ethiske Dom: „Retfærdigheden er guddommelig“. Ikke blot i Ethik og Religionsphilosophie, ogsaa i Æsthetiken maa Problemet komme frem, naar Spørgsmaalet om Skjønhedens Realitet i sidste Instans skal besvares. Men at Problemet rører sig saa stærkt i Ideelærens concretere Sphærer, indeholder just en Opfordring til at undersøge dets Udspring i Logiken; thi dersom det fra Grunden af skal løses \textit{in concreto}, maa det først være løst \textit{in abstracto}.

At den principielle Modsætning mellem Frihed og Nødvendighed har sit logiske Udspring i Modsætningen mellem Tænkeren og Tanken, mellem det Subjective og Objective i Tænkningens eget Væsen, kan ikke betvivles. Al Tænken forudsætter med Nødvendighed en tænkende Subjectivitet; men den tænkende Subjectivitet maa da ogsaa med Nødvendighed tænke Noget. Bliver nu den Tænkende paa denne Maade ligesaa afhængig af Tanken, som Tanken af den Tænkende, saa synes jo vistnok den gjensidige Afhængighed og dermed Nødvendigheden at være det Oprindelige; men den saaledes formede Nødvendighed er Frihedens Præformation. Subjectivitetens logiske Væsen er Tanken; den Tænkende bestemmer
% --- p. 366 ---
\opage{366}i alle Tanker og i al Tænken væsentlig sig selv: enhver Tænkningsact er oprindelig en Frihedsact. Vel har Tanken Objectivitet og i Objectiviteten Nødvendighed; men det i Tanken Objective er saa langt fra at tynge paa Tanken, at det tvertimod er et Vehikel for Tænkningens Energie og derved for Friheden.

Uagtet Friheden er ligesaa væsentlig forskjellig fra Nødvendigheden, som det Subjective er forskjelligt fra det Objective, saa har dog hverken Friheden nogen endelig Grændse i Nødvendigheden, eller Nødvendigheden i Friheden. Alt skeer oprindelig med Frihed, fordi Alt skeer med Nødvendighed, og omvendt. Hver Tanke, hvori Subjectiviteteten % PRINTER: printed "Subjectiviteteten", sense "Subjectiviteten"
sætter sin Frihed, har som Conclusion sin nødvendige Præmis i en allerede forudsat Tanke, og kommer derved til at støtte sig paa Nødvendigheden. Den forudsatte Tanke, der udøver Nødvendighedens Magt over den efterfølgende, er selv et ved en forudgaaende Nødvendighed betinget Product af Friheden. Umiddelbart seet er Præmissen fri, Conclusionen nødvendig; men da enhver Præmis i det uendelige Kredsløb selv er Conclusion, og enhver Conclusion Præmis, saa er Friheden i hver Slutning Nødvendighedens subjective, Nødvendigheden Frihedens objective Moment: Nødvendigheden er fri, Friheden nødvendig.

At den immanente Slutning er adæqvat Form for den frie Nødvendighed (\textit{libera necessitas}), skal nu kjendes derpaa, at den formaaer at udtrykke Frihedens og Nødvendighedens absolute Modsætning i absolut Identitet. Forsoningen af den absolute Modsætning med sin tilsvarende Identitet kan, ifølge vort Foregaaende, ligesaa lidt søges i en umiddelbar Væren som i en endelig Vorden; Alt kommer an paa den Energie, hvormed Bevægelse og Hvile i det Uendelige forudsætte og ophæve hinanden.

I denne Ideebevægelse kunne de forskjellige under „Relationen“ hørende Slutningssystemer, hypothetiske, kategoriske og disjunctive, indgaae. Fra Synspunktet af det Absolute,
% --- p. 367 ---
\opage{367}hvori Subjectiviteten selv gjennemtrænger sit Indhold, vil den hypothetiske Slutning tjene til at udvise Forholdet mellem Nødvendighed og Frihed. „Dersom A er B, saa maa C være D“, ved denne Dom har den dømmende Subjectivitet anerkjendt Nødvendigheden og dog reserveret sig Friheden. Reservationen ligger i Forsætningen; hvorvidt Forsætningen skal antages, beroer jo udelukkende paa, hvorvidt Subjectiviteten selv vil sætte den. Slutningens \textit{p. minor}: „Nu er A B, altsaa“ \dots\ betegner en Domsact, der begynder med Frihed og ender med Nødvendighed. Bestemmelsen af det Individuelle, det Ikke-Nødvendige, er i Ideeriget en ubetinget Frihedsact; at denne Frihedsact igjen afgiver et Element for Nødvendigheden, sees deraf, at den Præmis, der udtrykker Friheden, kun faaer Betydning ved det i Conclusionen, der udtrykker Nødvendigheden.

Det er Forholdet imellem Subjectiviteten og dens Indhold, som ved Ideebevægelsen nærmere skal bestemmes. Forsaavidt Bevægelsen gaaer fra Selv til Selv i Subjectiviteten, fra Selvet, det Almene, fra Selvet, det Særskilte, fra Selvet, det Enkelte, og gjennem Midtens Fordobling tilbage, er Bevægelsen central: i Centralbevægelsen er Nødvendigheden Moment for Friheden. Forsaavidt Bevægelsen derimod foregaaer i Indholdet, er Friheden Moment for Nødvendigheden. Bevægelsen i Indholdet er, ligesom Centralbevægelsen, tredobbelt: en Bevægelse fra Centrum (Regionen for det Almene) gjennem Midte og Peripherie tilbage til Centrum; en Bevægelse fra Midten (Regionen for det Særskilte) gjennem Centrum og Peripherie tilbage til Midten; endelig en Bevægelse fra Peripherien (Regionen for det Enkelte) gjennem Centrum og Midte tilbage til Peripherien. Under disse Bevægelser vexle Disjunction og Subsumtion med hinanden i det Uendelige. Bevægelsen fra Centrum gjennem Midten til Peripherien skeer ved Disjunctionen --- i Systemer af disjunctive Slutninger; Bevægelsen fra Peripherien gjennem Midten til Centrum skeer ved Subsumtionen --- i Systemer af kategoriske Slutninger. Bevægelsen fra Midte
% --- p. 368 ---
\opage{368}gjennem Centrum til Peripherie og fra Peripherie gjennem Centrum til Midte medfører en Krydsning af Individuationer i de allerforskjelligste Retninger.

Ved at gjøre Immanensslutningen til Form for den frie Nødvendighed eller, hvad der i denne Forbindelse er det Samme, for Nødvendighedens Eenhed med Friheden, have vi, overfladisk betragtet, ikke gjort Andet end udvikle en Form for Former, en Formernes Form; thi Friheden er jo Form ligesaavel som Nødvendigheden. Det kan da atter fremhæves, at den logiske Formudviking % PRINTER: printed "Formudviking", sense "Formudvikling"
altid maa give en logisk Indholdsudvikling; hver gjennemgaaende Formbestemmelse er i det Absolute selv en Indholdsbestemmelse.

At Nødvendighed og Frihed ere to indbyrdes modsatte Former for det Absolute, vil med andre Ord sige, at det Absolute som det Nødvendige og det Frie er sig selv oprindelig modsat; at her søges en Eenhedsform for Nødvendighed og Frihed, faaer da den Betydning, at her søges en Eenhedsform for eet og samme, sig selv modsatte Indhold. Hvad dette vil sige, bliver os ret indlysende, naar vi nu, istedenfor at behandle Indholdet fra Formsiden, vende Forholdet om og behandle Formen fra Indholdssiden, istedenfor at lade Indholdet reflectere sig i Formen, lade Formen reflectere sig i Indholdet.

Naar Formen, den absolute Form, reflecteres i sit absolute Indhold, have vi Ideen. Opkastes det Spørgsmaal: er Ideen væsentlig Form eller Indhold? endnu tydeligere: er Ideen det Indhold, der har det Absolute til Form, eller er maaskee det Absolute det Indhold, der har Ideen til Form? da mærkes det strax, at en endelig Adskillelse af Form og Indhold her er umulig. Skal Ideen som Idee absolut være Indhold, saa forvandler den sig til Form; skal den absolut være Form, saa forvandler den sig til Indhold. Gribe vi Ideen som Indhold i Formen, ville de tre Grundbevægelser
% --- p. 369 ---
\opage{369}i den foregaaende Formudvikling give os Ideen i Skikkelse af tre Grundideer: \emph{Sandhed, Skjønhed, Godhed}.

Det er da for det Første indlysende, at Formbevægelsen i sin Slutning fra Centrum til Centrum, fra det Almene gjennem det Særskilte og Enkelte til det Almene, A--(S, E)--A, maa give os Ideens Indhold i Almindelighedens og Nødvendighedens Form, altsaa \emph{Sandhedens Idee}.

Det er for det Andet indlysende, at Formbevægelsen i sin Slutning fra Midte til Midte, fra det Særskilte gjennem det Enkelte og Almene til det Særskilte, S--(E, A)--S, maa give os det Enkelte forklaret i det Almene og det Almene individualiseret i det Enkelte; en i Særskilthed saaledes reflecteret Ligevægt af det Frie og det Nødvendige, at det Frie er nødvendigt og det Nødvendige frit, svarer da kategorisk til \emph{Skjønhedens Idee}.

Det er for det Tredie indlysende, at Formbevægelsen i sin Slutning fra Peripherie til Peripherie, fra det Enkelte gjennem det Almene og Særskilte til det Enkelte, E--(A, S)--E, maa give os det absolute Indhold i Enkelthedens, i Selvbestemmelsens og Frihedens overgribende Form, altsaa Villiesideen, Friheden selv, det \emph{Godes Idee}.

Hvor adæqvat den til første Grundbevægelse svarende Slutning netop svarer til Sandhedens Idee, maa falde i Øinene. At Sandheden er een, at den har sit Væsen i det Almene, det Nødvendige, det Evige, er almindelig anerkjendt. Men det er da ogsaa almindelig anerkjendt, at Alt har sin Sandhed, at der er mange og mange Slags Sandheder, at den evige Sandhed ikke er tom, men uendelig indholdsrig. Det er tillige anerkjendt, at Sandheden, uagtet den hviler i evig Væren og behersker Alt med Nødvendighedens uimodstaaelige Magt, desuagtet indgaaer i det Timelige, indlader sig med det Vordende; thi naar Alt i al Evighed har sin Sandhed, maa Tiden, Vorden og det Vordende jo ogsaa have sin Sandhed. Det er endelig anerkjendt, at Sandheden,
% --- p. 370 ---
\opage{370}skjøndt i Væsen ubegrændset, dog har en virkelig Grændse i Usandheden, Illusionen, Løgnen o.~s.~v.

Disse forskjelligartede Bestemmelser lade sig alle forlige og bringe til Overeensstemmelse i Slutningen. Idet Slutningen viser, hvorledes det Almene selv ved at udvikle det Særskilte og Enkelte som Midte skiller sig i Extremer, viser den med det Samme, hvad det vil sige, at der af den ene absolute Sandhed kan fremgaae mangfoldige relative Sandheder, og at disse relative Sandheder, fastholdte hver for sig umiddelbart og endeligt, kunne blive til Usandheder, Illusioner, Bedrag o.~desl.

Det er i det Særskilte og Enkelte, at Sandheden har sin uendelige Rigdom af Indholdsbestemmelser. Saaledes udvikler enhver Fagvidenskab Sandheden fra sin Side, i sit særegne System af Erkjendelser. Men fra Endelighedens Standpunkt er det umuligt at komme ud over Adskillelsen af Extremer og Midte. Vistnok skeer Bevægelsen fra det Almene gjennem det Enkelte og Særskilte til det Almene; men Midten bliver staaende i udvortes Selvstændighed; Extremerne lade sig ikke sammenslutte. Det Almene, der udgjør det første Extrem, og hvorfra der gaaes ud, opfattes i svævende Anskuelser, tomme Abstractionsbegreber, hvormed der skrides til Undersøgelse af Kjendsgjerninger og Phænomener, altsaa af det Enkelte. Anskuelsen af det Enkelte skal da opløses i sine særskilte Bestemmelser, for at disse særskilte Bestemmelser igjen kunne optages i det Almene; Bevægelsen skeer fra Peripherie til Centrum: Methoden er \emph{analytisk}. Men det Enkelte og Særskilte lader sig ikke opløse; det Almene, der er Bestræbelsens Maal, bliver staaende som et sidste Extrem; Maalet, det concret Almene, naaes kun i en Tilnærmelse. Det opnaaede Almene kan dernæst fastholdes i sin Almindelighed som et første Grundlag, det sidste Extrem kan gjøres til det første. Det Særskilte og Enkelte er da i dette første Grundlag anticiperet saaledes, at det Almene ved indre Selvadskillelse kan afgive, omslutte og sammenholde sit
% --- p. 371 ---
\opage{371}Indholds forskjelligartede Dele; Bevægelsen gaaer fra Centrum til Peripherie: Methoden er \emph{synthetisk}\footnote{I væsentlig Overeensstemmelse med det hegelske Forbillede har Heiberg (Speculativ Logik. Prosaiske Skrister. % PRINTER: printed "Skrister", sense "Skrifter"
Første Bind. S.~354--62) skarpt og træffende angivet den analytiske og synthetiske Idee-Erkjendelse, nemlig under den Forudsætning, at den Erkjendende er, ikke en guddommelig, men menneskelig Subjectivitet.\par
„Erkjendelsen, som Livets fornyede Umiddelbarhed, er selv i sin Begyndelse umiddelbar, nemlig som \emph{Anskuelse}, og staaer saaledes under Schemaet af den første (umiddelbare) Slutning, E--S--A, hvis Betydning er, at Subjectet, formedelst Erkjendelsen, slutter sig sammen med den objective Sandhed.~\dots\ Men Anskuelsens Umiddelbarbed % PRINTER: printed "Umiddelbarbed", sense "Umiddelbarhed"
maa, som enhver anden, ophæves i sin Modsætning, som er dens dialektiske Udvikling. Herved fremkommer den \emph{analytiske Erkjendelse}, som staaer under Schemaet af den anden Slutning, S--E--A, hvis Betydning er, at Erkjendelsen selv, formedelst Subjectet, slutter sig sammen med den objective Sandhed eller det Almindelige; d.~e. at det erkjendende Subject gjør sin Erkjendelse til Gjenstand, eller stiller sig selv imellem sin Erkjendelse og dens Gjenstand, og i denne Stilling seer Gjenstanden sammensluttet med Erkjendelsen. Dette skeer ved \emph{at analysere} Erkjendelsen, d.~e. formedelst Adskillelse og Abstraction at bringe den i Formen af det Abstract-Almindelige. Den saaledes opnaaede Form er \emph{Definitionen}, hvorved den analytiske Erkjendelse bliver staaende, som ved sit Maal.“\par
„Den i Anskuelsen indeholdte umiddelbare Vished om Subjectets Identitet med Objectet i Erkjendelsen, er ophævet i Analysen, hvor Subjectet, som den forbindende Midte mellem sin Erkjendelse og dens Gjenstand, derved tillige adskiller disse, saa at deres Forening blot finder Sted i Subjectet. Ophævelsen af den umiddelbare Vished har kun den Bestemmelse, selv at ophæves, og vende tilbage til sit Udspring. Dette skeer i den synthetiske Erkjendelse, som staaer under Schemaet af den tredie Slutning, E--A--S, hvis Betydning er, at Objectet selv er den forenende Midte, som slutter Subjectet sammen med dets Erkjendelse. Ligesom derfor den analytiske Erkjendelse er Opløsningen af sig selv i Formen af det Abstract-Almindelige, saa gaaer den synthetiske
% note breaks here: continues at the foot of p. 372
Erkjendelse den modsatte Vei, og bringer det Abstracte i Formen af det Concrete; d.~e. den construerer Objectet, som Analysen betragtede som givet. Og ligesom den analytiske Erkjendelse ender med Definitionen, saa begynder den synthetiske Erkjendelse med denne.“\par
Manglen ved denne Fremstilling er imidlertid den, at her ikke vises, hvorledes det bliver den erkjendende Subjectivitet muligt at fuldende enten Analysen eller Synthesen og ved denne Fuldendelse tilegne sig Ideens Indhold.}% note *)
.

% --- p. 372 ---
\opage{372}Modsætningen imellem analytisk og synthetisk Methode forbliver i Reflexionen og kan ikke hæves til immanent Slutning; den tilhører Endeligheden og naaer ikke det Absolute. At den analytiske Methode kun har Ideen i Tilnærmelse er kjendeligt paa det Misforhold, hvori det Individuelle, ifølge denne Methode, kommer til at staae til det Særskilte og det Almene. Det Individuelle, hvormed her begyndes, er en umiddelbart concret Existens, en Ting med mange Egenskaber. Hvad gjør nu Analysen? Den ombytter Tingen med dens Begreb, og bestemmer Begrebet ved en Definition. Tingen, det umiddelbart Existerende, falder imidlertid udenfor Begrebet, som det Antilogiske falder udenfor det Logiske; om ogsaa de ved Analysen vundne Almeenbestemmelser nøiagtig svare til Enkeltbestemmelserne i det Existerende, er her dog en gabende Kloft mellem Existens og Begreb: Existensen er sandselig og concret; Begrebet er intellectuelt og abstract.

At Tingen, det Individuelle, ved Abstractionen er skilt fra sit Begreb, sees udtrykkelig af Definitionen. Det er nemlig ikke det Individuelle, men det Særskilte, ikke Tingen, men Begrebet, der bliver defineret. I Definitionens \textit{genus proximum} og \textit{differentia specifica} reflecteres jo vistnok Forholdet imellem det Almene og Særskilte, forsaavidt Begrebet bestemmes som Art, og Arten som individualiseret Slægt; men da Slægten er et Abstractum, og Arten, istedenfor at optræde særskilt som Mellemled mellem Slægten og Indi%
% --- p. 373 ---
\opage{373}videt, synker ned til et logisk Mærke, der vedføies Slægtsbegrebet: kan Definitionen ikke optræde som Slutning, men antager Præg af en i egentlig Forstand \emph{identisk Dom}\footnote{Vistnok har Heiberg i sin speculative Logik (Pros. Skr. 1ste Bind S.~356) bestræbt sig for at bestemme Definitionen som en Slutning; men Bestræbelsen ender i megen Tvetydighed, da han endog tre Gange skifter Synspunkt. Først hedder det, at Definitionen, ligesaavel som Omdømmet, ikke maa „betragtes blot subjectivt, som en menneskelig Tankeoperation, men objectivt“. Ikke desto mindre skal det dog være „det erkjendende Subject“, der, selv naar Definitionen er objectiv, udgjør Slutningens Mellemled. „Ligesom enhver Ting er et Omdømme, en Slutning og et System af Slutninger, saaledes er ogsaa enhver Ting en Definition i det erkjendende Subject, nemlig en Subsumtion af det Særskilte under det Almindelige, formedelst det Enkelte; d.~e. Erkjendelsens Subsumtion under det Abstract-Almindelige, formedelst det erkjendende Subject.“ Denne Synsmaade bliver saa tilsidst forandret derhen, at „det erkjendende Subject“ bortfalder som Mellemled og det Enkelte repræsenteres ved „Subsumtionens Ophævelse“; hele Forklaringen ender, trods den omhyggelige Exemplification, i noget aldeles Uforklarligt. „Fremdeles“, saa lyder nemlig Forklaringen, „er det formedelst sin Enkelthed, at det Særskilte subsumeres under det Almindelige, saa at Definitionen i sig selv er en Slutning, der staaer under Schemaet af anden Figur: S--E--A. Enhver Definition har derfor tre Led: 1) Det, som skal defineres, S; 2) Subsumtionen under det Almindelige, A; og 3) denne Subsumtions Ophævelse, idet den bestemmes som Enkelthed, E \dots\ F.~Ex. den Definition: „Mennesket er et Fornuftvæsen“, indeholder de tre Led: 1) Mennesket: 2) et Væsen (Subsumtion under det Almindelige); 3) men et Væsen, som er fornuftigt (Subsumtionens Ophævelse, idet det Almindelige bestemmes som Enkelt).“ Uklarheden kan neppe være større. Thi hvordan forholder nu det Enkelte i Subsumtionens Ophævelse sig til det Enkelte som det erkjendende Subject? Er her ved Subsumtionens Ophævelse to Enkelte eller kun eet? Hvorledes kan „det Almindelige bestemmes som det Enkelte“ i Prædicatet, og Prædicatet blive identisk med Subjectet, naar Subjectet selv ikke er det Enkelte, men det Særskilte?}% note *)
% PRINTER: in the note, "1) Mennesket: 2)" printed with a colon, sense ";"
.

% --- p. 374 ---
\opage{374}At den synthetiske Methode fra sin Side heller ikke formaaer at beseire Endeligheden og tilegne sig Ideen, er ligeledes indlysende. Naar den analytiske Methode, langt fra at magte Subsumtionen og naae fra det Individuelle gjennem det Særskilte til det i Slutningssystemet concret Almene, hilder sig i Abstractionen og ender med Definitionen: saa begynder den synthetiske Methode med Definitionen, ikke med det i Slutningssystemet concret Almene, bevæger sig, ikke igjennem principielle Disjunctioner, men igjennem formelle Inddelinger, og naaer saaledes ikke til Individerne, de concrete Existenser, men kun til Læresætninger og formelle Beviser. Definitionen, der skulde bestemme det \emph{Almene}, falder udenfor Inddelingen, der skulde give det \emph{Særskilte}; Inddelingen falder igjen udenfor Læresætningen, og en immanent Totalslutning viser sig at være umulig.

Idet den endelige Viden, hvad enten den iøvrigt foretrækker den analytiske eller den synthetiske Vei, forfeiler den nødvendige, gjennemgaaende Vexelbestemmelse af Individ, Art og Slægt, forfeiler den Sandheden og hilder sig i en Usandhed. Bestemmes Sandheden som nødvendig Eenhed af Tænken og Væren, Begreb og Realitet, da maa denne Eenhed selv udtrykkelig bestemmes som Slutningseenhed. Den aristoteliske Fornegtelse af Slægternes og Arternes ideelle Selvstændighed er ligesaa afvigende fra Sandheden som den platoniske Hypostaseren. Vistnok gives der Arter og Slægter, der ikke have evig Gyldighed (Arten: Hvirvelvind, af Slægten: Vind; Arten: Sandbunke, af Slægten: Bunke); men hvad beviser det? At de til slige Arter og Slægter svarende Individer (denne Hvirvelvind, denne Sandbunke) heller ikke have evig oprindelig Gyldighed! Seet i Ideens Lys have Individ, Art og Slægt samme Nødvendighed, om end i forskjellig Form, samme Selvstændighed, om end i forskjellige Regioner. Enhver Idee er med Nødvendighed Individ, Art og Slægt i absolut Modsætning, og just derfor Individ, Art og Slægt i absolut
% --- p. 375 ---
\opage{375}Eenhed: Ideens Selveenhed er en immanent Slutningseenhed.

At Individet har sin Nødvendighed i Art og Slægt, Arten sin Nødvendighed i Individ og Slægt, Slægten sin Nødvendighed i Individ og Art, er et Beviis paa Frihedens Eenhed med den rationelle Nødvendighed, og denne Eenheds oprindelige Eenhed med Sandheden. Friheden viser sig nemlig deri, at ethvert af Slutningens Led bevarer sin Selvstændighed; men da Extremernes Selvstændighed er Eet med den selvstændige Midte, og omvendt: saa er den gjensidige Selvstændighed identisk med den gjensidige Afhængighed, og Friheden sat som Moment i Nødvendigheden. Paa Frihedens fuldstændige Opgaaen i Nødvendigheden beroer Sandheden.

Men hvorpaa beroer saa Usandheden? Usandheden kan antage meget forskjellige Former. Den kan nominalistisk vise sig deri, at der tillægges Individet og det Individuelle en ubetinget Selvstændighed paa Artens og Slægtens Bekostning; den kan realistisk vise sig deri, at Slægter og Arter erholde Selvstændighed paa Individernes Bekostning; den kan vise sig i en Overvurdering af Tænkningen paa Anskuelsens, af den anskuende Reflecteren paa Forestillingens Bekostning o.~s.~v. Men hvad der især blender og forvirrer den endelige Subjectivitet, er Biforestillinger af sandselig Umiddelbarhed, af endelig existentiel Bestemthed.

Ved at indblande Biforestillinger af det sandseligt Concrete og det intellectuelt Abstracte bliver den endelige Viden bestandig hildet i den Illusion, at det Individuelle er timeligt, og kun det Almene evigt, at Individerne maae forgaae, og kun Arterne, Slægterne eller endnu almenere Væsensformer bestaae. Herved fremkommer en utaalelig Strid imellem det Virkelige og det Sande. Individet er det Virkelige; men Individet er i Virkeligheden forgjængeligt, og det Virkelige forsaavidt et Usandt. Det Almene (Arten, Slægten o.~s.~v.) er i Væsen det Uforgjængelige og forsaavidt det Sande.
% --- p. 376 ---
\opage{376}Men det Almene existerer ikke; det har sin Væren i Tanken, ikke i Virkeligheden: det Sande er altsaa ikke virkeligt, og det Virkelige ikke sandt; saa haard er Modsigelsen.

I den immanente Slutning ere derimod det Individuelle og Særskilte fra begge Sider omsluttede af det Almene. Med det Enkelte og Særskilte til Mellemled danner Sandhedsideen af sit Almene to Extremer, hvis Sammenslutning netop er den sande concrete Eenhed; paa denne Maade faaer Sandheden Magt med Virkeligheden: alle Virkelighedens Momenter faae deres Sandhed.

Ligesom det Almene vilde være eensidigt og usandt, dersom det i den immanente Slutning ikke udtrykkelig havde det Særskilte til Modsætning, saaledes vilde Sandhedsideen, Ideen for den væsentlige Objectivitet, selv blive eensidig og usand, dersom den ikke havde Skjønhedsideen, Ideen for den phænomenale Objectivitet, til absolut Modsætning. At den absolute Modsætning maa grunde sig paa absolut Eenhed, vil i nærværende Sammenhæng endvidere fremlyse deraf, at Skjønhedsideen kun er sand ved at være Sandhedsideen modsat.

Fra to Sider bliver nemlig Skjønhedsideen formørket, idet den enten ligefrem smelter sammen med Sandheden og taber sig i Nødvendigheden, eller ved umiddelbar Frihed adskilles fra Sandheden og priisgives Tilfældigheden. Neoplatonikerne, navnlig Plotin, identificere Skjønhedsprincipet med Skjønheden selv, den over alle skjønne Former ophøiede Skjønhed, en Skjønhed, der dog ikke selv kan være skjøn, efterdi den er forskjellig fra det Skjønne. „Det, som giver Skjønhed, er selv uden Skikkelse; thi det, vi kalde Skikkelse, er et i sig Formløst, der har modtaget Form, og saaledes har Andeel i Skjønheden, men ikke selv er den“.\footnote{Ἀρχὴ δὲ τὸ ἀνείδεον, τὸ μὴ μορφῆς δεόμενον, ἀλλ᾽ ἀφ᾽ οὗ πᾶσα μορφὴ νοερά. Καὶ τὸ κάλλος αὐτοῦ ἄλλον τρόπον καὶ κάλλος ὑπὲρ κάλλος· οὐδὲν γὰρ ὂν τί κάλλος; ἐράςμιον δὲ % note breaks here: p. 376 / p. 377
% PRINTER: printed "ἐράςμιον" (final-form ς inside the word), sense "ἐράσμιον"
ὂν τὸ γεννῶν ἂν εἴη τὸ κάλλος. Δύναμις οὖν παντὸς καλοῦ ἄνθος ἐστὶ κάλλους καλλοποιόν. Plotin. Jvfr. Kirchner. Anfr.\ Skr.\ S.~36--39.}% note *)
 Det er, med
% --- p. 377 ---
\opage{377}andre Ord, den tænkte Skjønhed, Skjønheden i Form af Sandhed, der her er sat som det Høieste. Men naar det nu er Ideen som Sandhed, ikke særlig som Skjønhed, der giver Skjønhed, saa at det at „have Andeel i Skjønheden“ maa blive det Samme som at have Andeel i Sandheden: hvorledes forholder Skjønhedsprincipet sig da til den skjønne Skikkelse? Idet Skikkelsen dannes, maa der foruden Formen være Noget, der modtager Formen, nemlig Stoffet. Dersom Stoffet ikke blev formet, kunde Skikkelsen ikke fremkomme; men Stoffet som Stof er jo i et tilfældigt Forhold til den \emph{bestemte} Form, selvfølgelig ogsaa til den \emph{skjønne} Form. Der maa da, ifølge den neoplatoniske Synsmaade, være noget Tilfældigt i den skjønne Skikkelse; men idet der træder et Tilfældigt ind imellem den skjønne Skikkelse og Principet, er Skjønheden selv sunken ned i det Tilfældige og dermed i det Forgjængelige.

„Macht' ich doch, sagte der Gott, nur das Vergängliche schön“. Disse Ord, hvormed Digteren lader Zeus trøste „den forgjængelige Skjønhed“, indeholde en væsentlig Krænkelse af Skjønhedens Idee. Kun den jordiske Skjønhed, den Skjønhed, der i de vexlende Phænomener maa indlade sig med det Stofagtige og Tilfældige, kan være forgjængelig; Skjønhedsideen selv kan umulig være forgjængelig. Det Andethedselement, som Skjønheden oprindelig kræver, kan ikke falde udenfor Ideen; det maa med Nødvendighed være et Ideen selv iboende Moment. Og dog sætter den reflecterende Anskuen en væsentlig Forskjel imellem Selvbestemmelsen og Bestemmelsen ved Andet. Den oprindelige Eenhed af Ideen og Skikkelsen, der udtrykker Skjønheden, er netop en reflecteret
% --- p. 378 ---
\opage{378}Eenhed af Selvbestemmelse og Bestemmelse ved Andet, en reflecteret Eenhed af Frihed og Nødvendighed; en reflecteret Eenhed, der er saa gjennemgaaende, at Reflexionen selv opløser sig i forklaret Umiddelbarhed. Skikkelsen maa bestemmes, thi Skjønhedens Idee maa have sin Skikkelse; er Bestemmelsen abstract Selvbestemmelse, saa er Skikkelsen vilkaarlig; er den \textit{in abstracto} en Bestemmelse ved Andet, saa er Skikkelsen tilfældig. Imod Vilkaarligheden og Tilfældigheden staaer Nødvendigheden. I Adskillelsen af det abstract Vilkaarlige, det abstract Tilfældige og det abstract Nødvendige er Reflexionen endeliggjort; i den endeliggjorte Reflexion sees Uskjønheden, men i Uskjønheden sees igjen det Særskilte.

Det Uskjønne og det Skjønne have altsaa det tilfælles, at de begge beroe paa det Særskilte. Der gives en Særskilthed, som er uskjøn, efterdi den falder udenfor Ideen, og en Særskilthed, der er skjøn, efterdi den udtrykker Ideens og Skikkelsens oprindelige Eenhed og saaledes er identisk med Ideen. At det Særskilte udtrykker Ideen, er betinget deraf, at det Virkelige og det Tilfældige opløser sig i det Nødvendige; ved denne Opløsning fremkommer Friheden: Ideeskikkelsen er bestemt med Frihed, efterdi den er bestemt med Nødvendighed.

Naar Talen er om skjønne Gjenstande (en skjøn Hest, en skjøn Blomst, en skjøn Qvinde), da indsees det let, at her ikke er Tale om forskjellige Phænomener, som hvert for sig have sin Deel i een og samme Skjønhed, men om forskjellige Arter af Skjønhed. Og dog maa Skjønheden i alt det Skjønne jo være sig selv liig: hvorledes forholder da Skjønheden sig til de forskjellige Arter af Skjønhed? Ethvert philosophisk System, der bliver i Reflexionen uden at naae Slutningen, maa vistnok erklære dette Problem for uopløseligt. For at en Gjenstand skal kunne kaldes skjøn, udfordres for det Første, at den maa være en individuel Afspeiling af sin Art, thi Arten er dens Forbillede; for det Andet, at Arten selv maa have utve%
% --- p. 379 ---
\opage{379}tydig Bestemthed i Modsætning til andre Arter, og altsaa være en typisk Individuation af Slægten. Uskjønheden kan med Hensyn hertil fremkomme paa to Maader: deels derved, at Individet staaer i æsthetisk Misforhold til Arten (en Vanskabning), deels derved, at Arten selv er en æsthetisk Tvetydighed (en Skruptudse).

At nu de mange Arter af Skjønhed antages for særskilte Aabenbarelser af Skjønheden selv, den ene og samme Skjønhed, opfatter Reflexionen i sin Eensidighed saaledes, at Skjønheden absorberes af Sandheden, Friheden af Nødvendigheden. At de forskjellige Arter af Skjønhed træde selvstændigt frem som særlige Typer, guddommelige Forbilleder, udtyder den paa Slutningen usikkre Reflexion saaledes, at Anskuelsen har sat sig op imod Tanken, at den med Sandheden identiske Skjønhed vil adskille sig fra Sandheden. Skilsmissen kan iøvrigt foregaae paa to Maader, enten derved, at de guddommelige Forbilleder antages for evige, i sig selvstændige Væsener, og Frigjørelsen af det Særskilte maa da beroe paa en evig Tilfældighed; eller derved, at Guddommen, den absolute Subjectivitet, vilkaarlig bestemmer de særskilte Typer: „Gud skaber Ideerne“, saa at Frigjørelsen er vilkaarlig.

Disse Abstractioner og Modsigelser forsvinde, naar Ideen udtrykkelig opfattes som Slutning. I Slutningen, ikke i Tænkningens, men i Anskuelsens immanente Slutning er Blikket særlig fæstet paa det Særskilte, og det Særskilte udtrykkelig bestemt som Ideeskikkelse, som Type (παράδειγμα).

Ideeskikkelsen er med Nødvendighed; thi det Almene og det Individuelle maae --- det er Ideens uafviselige Fordring --- reflectere sig i Skikkelsen saaledes, at det Almene med det Individuelle til sit Moment bliver Slægt (γένος), det Individuelle med det Almene til Moment bliver Art (εἶδος), og Arten den Midte, hvori det Almene er sammensluttet med det Individuelle. I denne Henseende gjælder det, at Subjectiviteten ikke vilkaarlig frembringer, men oprindelig anskuer de
% --- p. 380 ---
\opage{380}evige, i sig bestemte Typer (δῆλον ὡς πρὸς τὸ ἀΐδιον ἔβλεπεν).

Men Ideeskikkelsen er da ogsaa med Frihed. Ideen kan kun antage Anskuelsens Form derved, at Subjectiviteten anskuer; men den anskuende Subjectivitet kan kun bestemmes til Anskuelse af Andet ved Andet, idet den bestemmer sig selv. I utallige Anskuelser, altsaa i utallige Skikkelser, kan det Almene slutte sig sammen med det Individuelle. Hvorfor nu just i \emph{denne}, ikke i en anden Skikkelse? Her er ikke først en i sig objectiv, af Anskuen uafhængig Skikkelse, hvorefter den anskuende Subjectivitet „sig haver at rette“; nei Formens Gyldighed er kun objectiv, forsaavidt den er subjectiv; den Nødvendighed, hvormed Ideeskikkelsen anskues, bestemmes oprindelig ved Anskuelsen selv, altsaa ved Subjectiviteten, altsaa ved Friheden (οὕτως ἐποίησε μίαν μόνον αὐτὴν ἐκείνην).

Den immanente Slutning er den frie og dog nødvendige Form, hvori den det Særskilte anskuende Subjectivitet frembringer og bevarer Skjønhedens evige Typer.

Hvad endelig den tredie Grundbevægelse, Bevægelsen fra det Enkelte gjennem det Særskilte og Almene tilbage til det Enkelte, altsaa fra Peripherie gjennem Midte og Centrum til Peripherie, angaaer, skal det nu vises, at den hertil svarende Slutning er en Villiesslutning, der oprindelig bestemmer det Godes Idee.

Det villende Selv maa absolut ville Noget; idet Selvheden er bestemt som Villie, og Villien som det subjectivt Enkelte, er Andetheden med det samme bestemt som Villiens Indhold, Indholdet igjen som lutter ideelle Gjenstande, en Mangfoldighed af objectivt Enkelte. Idet her begyndes fra Peripherien, begyndes her med den haarde Modsætning af Enkelt mod Enkelt, af det villende Selv mod Villiens Gjenstand. Det Haarde i Modsætningen bliver at opfatte fra to Sider. Forsaavidt Villiens forestillede Gjenstand staaer Villien imod, falder den udenfor Villien og bliver staaende som uover%
% --- p. 381 ---
\opage{381}vindelig Skranke. Ligeoverfor Skranken er Villien tom, og den tomme Villie afmægtig. Forsaavidt Villiens Gjenstand derimod er sat og bestemt ved Villien selv, falder den ind under Villien og kan umulig holde sig som Skranke. Den skrankeløse, i sig uindskrænkede Villie bestemmer nu sit Indhold enkeltviis i hvert enkelt Tilfælde, som den selv vil; Villiens Indhold, det Gode, er da fra dette Synspunkt absolut bestemt ved en absolut Vilkaarlighed: \textit{bonum est, quia Deus vult}.

Men den absolute Vilkaarlighed er, naar vi see nærmere til, en utaalelig Modsigelse. Vilkaarlighedens Magt skulde jo netop vise sig deri, at den Villie, der magter alle Skranker, udtrykkelig sætter og bestemmer det objectivt Enkelte, og gjør med dette Alt hvad den vil. Det er Villien, der i sin absolute Vilkaarlighed giver sit Indhold objectiv Tilværen i det objectivt Enkelte; det er Villien, der i sin absolute Vilkaarlighed kan tilbagekalde dette Enkelte og gjøre det til Intet. Hvad følger nu heraf? Heraf følger med uafviselig Consequens, at det Enkelte ingen objectiv Gyldighed har; at Villien ved i alle Enkeltobjecter at fortære sin objective Modsætning fortærer sit Indhold og tærer paa sig selv. Den absolute Vilkaarlighed, der skulde være et Udtryk for Villiens Almagt, bliver, væsentlig seet, et Udtryk for dens Afmagt.

For at Modsigelsen skal kunne løses, maa det villende Selv fatte sig som tænkende. Idet Villen gaaer i Eet med Tænken, bliver den tænkende Villie en objectiv Almeenvillie, en Villie, som i Eet og Alt lader Indholdets egen væsentlige Udvikling gjælde. Men Indholdets væsentlige Udvikling er en Udvikling i Tænkningens Form; i Tænkningens fuldendte Form er det objective Indhold ved den immanente Slutning absolut bestemt som Sandheden. Det er altsaa Sandheden, der antager Villiens Form, idet Almeenvillien antager Tænkningens Form. I Tænkningens Form er den almene Villie naturligviis absolut bestemt ved sit Indhold, thi Indholdet er
% --- p. 382 ---
\opage{382}jo udtrykkelig sat som Sandhedens Indhold. Fra dette Synspunkt er det altsaa ingenlunde Villien, der vilkaarlig bestemmer det Gode, men omvendt det Gode, der væsentlig bestemmer den gode Villie: \textit{Deus vult, quia bonum est}.

At denne Modsigelsens Løsning dog heller ikke er tilfredsstillende, er lige iøinefaldende, enten vi nærmest betragte det Godes Indhold eller dets Form. Hvad Indholdet angaaer, er det Gode ligefrem ombyttet med det Sande, idet Villien, hvad Formen angaaer, ligefrem er blevet ombyttet med Tænkningen: den almene Villie er en reen Tankevillie. Som Tankevillie er Villien contemplativ og theoretisk; den theoretiske Villies contemplative Væsen er kjendeligt derpaa, at den, istedenfor at bestemme sit Indhold, Punkt for Punkt lader sig bestemme af Indholdet. Men en Villie, hvis almene Bestemmen bestaaer i at lade sig bestemme, er en passiv, villieløs Villie, altsaa en Villie, som ingen Villie er. Ved at tilhylle Modsigelsen kan man dog endnu gjøre et tilsyneladende Fremskridt, idet man lader Villien, efter at den formedelst sin Eenhed med Tænkningen er blevet Almeenvillie, tillige formedelst sin Eenhed med Anskuelsen blive en særskilt Villie. Det er da paa denne Maade ikke blot Sandhedens, men ogsaa Skjønhedens Idee, der umiddelbart opfattes som en Villiesidee og gjøres til Eet med det Gode.

Den ulogiske Amalgameren af det Sande, det Skjønne og det Gode er i Almindelighed charakteristisk betegnet ved en æsthetisk-overfladisk Opfattelse af den guddommelige Viisdom. Ikke nok, at Grækerne i al Naivetet identificerede det Skjønne med det Gode (καλοκαγάθία); % PRINTER: printed "καλοκαγάθία" (acute on both α and ι), sense "καλοκἀγαθία"
„den speculative Dogmatik“ vedbliver endnu i sin logiske Impotens at gjøre det Samme. „Hvad Speculationen kalder Ideen, den verdensdannende Tanke, betegnes i den hellige Skrift som Viisdommen, der evindelig legede for den Høiestes Ansigt, og i hvilken han havde sin Lyst. Og den beskrives ikke blot som indre Speilbillede i Gud, men ogsaa som den virkende, den altformende
% --- p. 383 ---
\opage{383}Tanke. Thi Viisdommen (Ideen, den guddommelige Sophia, den himmelske Jomfru, som Theosopherne have kaldt den) er alle Tings \emph{Værkmesterinde}“ o.\ s.\ v. „Den formende Tanke“, der peger mod Sandhedens Idee, „den himmelske Jomfru“, der peger mod Skjønhedens Idee, „Viisdommen“, der peger mod „et uendeligt Villiesformaal“, altsaa mod det Godes Idee, disse evige Visere paa Logikens Uhr rager „den christelige Speculation“ i den allerchristeligste Mening sammen i een Æske.

Ved at fremhæve Villiens Eenhed med Tænken og Anskuen er dog allerede Saameget vundet, at det objectivt Enkelte, Villiens ideelle Gjenstand, nu kan danne det Mellemled, hvori den villende Subjectivitet\footnote{Den logisk-ontologiske Modsætning af de tre Grundideer: det Sande, det Skjønne og det Gode, er af afgjørende Betydning saavel for Livet som for Videnskaben. Dersom Sandhedens Idee ikke havde sin Selvstændighed ligeoverfor det Gode, vilde Videnskaben aldrig være istand til at hævde sig et af den praktiske Livsanskuelse uafhængigt Princip; dersom det Godes Idee ikke havde ubetinget Selvgyldighed i Modsætning til Sandhedsideen, vilde Ethik og Religion umulig kunne hævde sig et særligt, af Videnskaben uafhængigt Princip. Middelalderen holdt paa Troesprincipet og kuede Videnskaben; den gjorde Sandhedsideen ligefrem afhængig af det Godes Idee: det var barbarisk. Nutiden gaaer til den modsatte Yderlighed; den holder paa Videnskaben og vil ved Videns Midler opløse Troen; den gjør uden videre det Godes Idee afhængig af Sandhedsideen: den sidste Misforstaaelse er ligesaa iøinefaldende som den første.}% note *)
 lader det Almene slutte sig sammen med det Særskilte. Idet Mellemleddet, det objectivt Enkelte, af Villien selv bestemmes som dens udtrykkelige Gjenstand, ere de to Extremer, Tænkningens Almene og Anskuelsens Særskilte, med det Samme bestemte som den tænkende og anskuende Villies objective Indhold. Ved Extremerne bestemmes da Nødvendigheden, ved Midten bestemmes Friheden; ved Slutningen selv bestemmes Eenheden, den for det Godes
% --- p. 384 ---
\opage{384}Idee afgjørende Eenhed af Nødvendighed med Frihed i Frihedens Form.

\medskip
\noindent\textit{β)} \emph{Objectiv Slutning: Modalfunctioner af Magt.}
\medskip

Det blev i „Subjectivitetens Logik“ fra forskjellige Sider viist, at den hegelske „Slutning“ var for objectiv til at være subjectiv, og igjen for subjectiv til at være objectiv. Grundfeilen var: 1) at Forskjellem imellem det Enkelte, det Særskilte og det Almene blev sat og hævet paa en altfor flygtig og just derfor overfladisk Maade;\footnote{Grundideernes Logik. I. S.~530--32.}% note *)
 % PRINTER: printed "Forskjellem", sense "Forskjellen"
2) at Extremernes ved Mellemleddet medierede Eenhed opfattedes som en Værens-Eenhed, istedenfor at opfattes som en ved Tænken og Virken bestemt Væsenseenhed, d.\ v.\ s.\ som en Bevægelseseenhed;\footnote{Anfr.\ Skr.\ S.~216.}% note **)
 3) at der, efterdi ethvert Hensyn til Bevægelsesenergien manglede, ikke kunde paavises nogen sand Forskjel imellem Slutningen i Idealsystemet, den ontologisk subjective Slutning, og Slutningen i Realsystemet, den ontologisk objective Slutning,\footnote{Anfr.\ Skr.\ S.~224.}% note ***)
 saa at Slutningens Fuldendelse i Teleologie kun fik en schematisk, ikke en for Indholdet væsentlig Betydning.\footnote{Anfr.\ Skr.\ S.~224--28.}% note †)
 Resultatet blev, at Leddenes Charakteer og indbyrdes Forhold maatte i hver Sphære og for hvert System angives paa en til Functionernes Energie nøiagtig svarende Maade.\footnote{Anfr.\ Skr.\ S.~433--34.}% note ††)

I nærværende Sammenhæng, hvor Opgaven er at belyse den objective Slutning fra de logiske Functioners Synspunkt, vil Indvendingen mod den hegelske Slutnings Objectivitet væsentlig lyde paa, at Objectiveringens Charakteer er overseet og Functionen forfeilet. Den hegelske Sætning: \emph{Alt er en Slutning}, lader sig i umiddelbar objectiv Forstand
% --- p. 385 ---
\opage{385}ligesaa lidt forsvare, som den Sætning, at \emph{Alt er en Dom}. Vel er den existerende Gjenstand en Dom, forsaavidt den i Væsen er et Magtens Ord; men saaledes viser den sig først, naar dens umiddelbare Objectivitet igjen bliver objectiveret: de logiske Functioner af Magtens Ord ere alle Functioner af fordoblet Objectivering. For den umiddelbare Objectivering ere Gjenstandene i deres virkelige Forhold og Forbindelser Kjendsgjerninger, ikke Begreber og Domme, endnu mindre Slutninger. Ved den fordoblede Objectivering blive Gjenstandene til objective Domme, antilogiske Sætninger i det kategoriske Magtsprog: til det kategoriske Magtsprog, de objective Dommes System, svare de interpreterende Functioner.

At nu de objective Slutninger ligesaa lidt som de objective Domme kunne opfattes og bestemmes uden Objectiveringsfordobling og de til en saadan Fordobling svarende Functioner, er en Selvfølge, og lader sig uden Vanskelighed godtgjøre. Naar den objective Slutning, saaledes som hos Hegel, udelukkende skal opfattes fra Værens Standpunkt og bestemmes ved indicative Functioner, synker den ned til et tomt Schema, et Schema, hvis objective Betydning ovenikjøbet kommer til at beroe paa subjectiv Reflexion; hvorimod Systemet af de tre Slutninger, ogsaa med den Form, hvortil det indskrænkes i den hegelske Logik, faaer væsentlig Betydning, naar det bæres, ikke af indicative, men af explicative og interpreterende Functioner. Vi skulle ved Exempler oplyse dette.

Om hvad han kalder Produktet af den formale mechaniske Proces udtaler Hegel sig paa følgende Maade: „\emph{Unmittelbar} ist das Objekt \emph{vorausgesetzt} als Einzelnes, ferner als Besonderes gegen andere, drittens aber als Gleichgültiges gegen seine Besonderheit, als Allgemeines. Das \emph{Produkt} ist jene \emph{vorausgesetzte} Totalität des Begriffes nun als eine \emph{gesetzte}. Er ist der \emph{Schlußsatz}, worin das mitgetheilte Allgemeine durch die Besonderheit des Objekts mit der Ein%
% --- p. 386 ---
\opage{386}zelnheit zusammengeschlossen ist; aber zugleich ist in der Ruhe die \emph{Vermittelung} als eine solche gesetzt, die sich \emph{aufgehoben} hat, oder doß das Produkt gegen dieß sein Bestimmtwerden gleichgültig und die erhaltene Bestimmtheit eine äußerliche an ihm ist“.\footnote{Wissenschaft der Logik. Zweiter Theil. S.~189.}% note *)
 % PRINTER: printed "doß", sense "daß"

Det er nu let at vise, hvorledes denne Udtalelse forandrer sin Charakteer, eftersom det er indicative, explicative eller interpreterende Functioner, hvorpaa den antages at hvile. Gaaer man ud fra Hegels egen Forudsætning, at „Produktet af den mechaniske Proces“ skal fastholdes i umiddelbar Objectivitet, da er et saadant Object i antilogisk Virkelighed en udvortes Gjenstand, ikke et Magtens Ord, ikke en Dom, endnu mindre en Slutning. Som existerende Gjenstand er Objectet aldeles ligegyldigt mod den Proces, hvis Produkt det sættes at være. Først fremstilles Objectet, som om dets Tilværelse beroede paa en i Slutning udtrykt Mediation; dernæst fremstilles det som et mod denne Mediation ligegyldigt Produkt, altsaa som et Object, hvis Tilværelse er ligegyldig mod den foregivne Slutning. Det mechaniske Object er altsaa en Slutning i den Forstand, at det ikke er nogen Slutning; det Skrigende i denne Modsigelse er kun dæmpet ved den tvetydige Vending: „aber zugleich ist in der Ruhe die \emph{Vermittelung} als eine solche gesetzt, die sich \emph{aufgehoben} hat.“ Skal Tvetydigheden hæves, saa viser sig igjen den gamle Confusion af objectiv Sagbevægelse med subjectiv Tankebevægelse. At et mechanisk Object fremkommer som Produkt af en mechanisk Proces, et Aggregat f.\ Ex.\ af en Aggregationsproces, er logisk i sin Orden; ligesom det da ogsaa er i sin Orden, at en saadan Proces kan være standset, naar Produktet er færdigt, og „Mediationen“ i denne Hvile forsaavidt være sat som den, der „har ophævet sig selv“. Men dette vedkommer aldeles ikke den Proces, hvorom her tales. Den mechaniske
% --- p. 387 ---
\opage{387}Proces, som her skal gjælde for en objectiv Slutning, er en saa aldeles subjectiv Tankeproces, at den angaaer det Intellectuelle i Objectiveringen, uden i nogen Maade at angaae det virkelige Object.

Sættes derimod Gjenstandens Begreb imod Gjenstanden, saaledes at der udtrykkelig reflecteres paa Begrebet, da blive Functionerne explicative, og da vil ogsaa den hegelske Udtydning vise sig i et bedre Lys. Hvad der siges ligefrem om det mechaniske Object, at det er „\emph{forudsat} som et Enkelt“, at det er „som et Særskilt mod andre“, at det „som et mod denne Særskilthed Ligegyldigt er et Almindeligt“, er virkelig tre Hovedsider af det mechaniske Objects Begreb, og den begribende Subjectivitet, der anlægger Begrebets Maalestok paa Objectet selv, er fra dette Synspunkt da ogsaa ganske rigtig opfordret til at lade Begrebets tre Hovedsider adskille sig med en til „Mechanismen“ svarende Umiddelbarhed og derpaa lade de saaledes adskilte Hovedsider slutte sig sammen som to Extremer med en Midte, og forsaavidt give en til Objectets Begreb svarende Slutning, altsaa en objectiv Slutning. Men Slutningsbegrebet i denne Slutning kan desuagtet ikke fuldendes, med mindre der udtrykkelig reflecteres paa den Tilfredsstillelse, det i Extremernes Sammenslutning erkjendte Begreb yder den begribende Subjectivitet, d.\ v.\ s.\ med mindre Subjectiviteten veed sin Begriben af det givne Object identisk med den Objectet frembringende Begriben, og saaledes gjennem Objectet corresponderer med den frembringende Subjectivitet; men paa disse Betingelser blive Functionerne interpreterende.

Den oprindelige Sammenhæng mellem den objective Slutning og de interpreterende Functioner bliver endnu mere indlysende, naar vi, istedenfor en mechanisk formel Proces, vælge en chemisk til Behandling. Dersom der til en objectiv Slutning kun udfordres en Sammenslutning af to Extremer med en Midte, vil enhver chemisk Proces kunne tjene som Exempel. Er Ilten det negative, Brinten det positive Extrem,
% --- p. 388 ---
\opage{388}saa er deres Affinitet den væsentlige Midte, og Vandet det Produkt, der betegner den fuldendte Slutning. Paa denne Maade bliver den chemiske Proces udtydet hos Hegel: „Er (der Proceß) beginnt mit der Voraussetzung, daß die gespannten Objekte, so sehr sie es gegen sich selbst, es zunächst eben damit gegen einander sind; --- ein Verhältniß, welches ihre \emph{Verwandtschaft} heißt \ldots\ Die Mitte, wodurch nun diese Extreme zusammengeschlossen werden, ist \emph{erstlich} die \emph{ansichseyende} Natur beider, der ganze beide in sich haltende Begriff. Aber \emph{zweitens}, da sie in der Existenz gegen einander stehen, so ist ihre absolute Einheit auch ein \emph{unterschieden} von ihnen \emph{existirendes}, noch formales Element; das Element der Mittheilung (Hegel tænker herved særlig paa Vandet) \ldots\ Das Verhältniß der Objekte ist als bloße Mittheilung in diesem Elemente einer Seits ein ruhiges Zusammengehen, aber anderer Seits ebenso sehr ein \emph{negatives Verhalten}, indem der konkrete Begriff, welcher ihre Natur ist, in der Mittheilung in Realität gesetzt, hiermit die \emph{realen Unterschiede} der Objekte zu seiner Einheit reducirt werden“.\footnote{Wissenschaft der Logik. Zweiter Theil. S.~202--203.}% note *)

Men hvormegen Priis man end maa sætte paa den dialektiske Begrebsudvikling, som gaaer igjennem hele dette Partie af den hegelske Logik, saa er det for det Første indlysende, at den beroer paa explicative Functioner; dens Afvigelse fra Sagudviklingen kan sees ved Henblik til den Fremstilling, der gives i Chemien. Dernæst er det øiensynligt, at Resultatet af den chemiske Proces umulig kan fremstilles som Slutning (Conclusum), saalænge Begrebsudviklingen ikke gaaer ud over det blot Explicative. For at det chemiske Produkt skal forvandles til Slutning, maa det udtrykkelig forudsættes, at Resultatet er anticiperet i sine Forudsætninger, at det kommer frem som et oprindelig Tilsigtet, at Processen med dens Extremer og Midte først ved Resultatets Opnaaelse
% --- p. 389 ---
\opage{389}viser sig i sin egentlige Betydning; men ved denne Opfattelse bliver Slutningen teleologisk, og Functionerne interpreterende.

Begrebets Bevægelse i Hegels Logik fra „Mechanisme“ gjennem „Chemisme“ til „Teleologie“ udviser en interessant Blanding af Sandt og Falsk, forsaavidt den beroer paa en Sammenblanding af Sagbevægelse, objectiv Begrebsbevægelse og Bevægelse i subjectiv Begriben, altsaa, hvad Tankevirksomheden angaaer, paa en Sammenblanding af indicative, explicative og interpreterende Functioner. Fremstillingen af „Chemismen“ slutter med følgende Ord: „Der Begriff, welcher hiermit alle Momente seines objektiven Daseyns als äußerliche aufgehoben und in seine einfache Einheit gesetzt hat, ist dadurch von der objektiven Aeußerlichkeit vollständig befreit, auf welche er sich nur als eine unwesentliche Realität bezieht; dieser objektive freie Begriff ist der Zweck“.\footnote{Anfr.\ Skr.\ S.~208.}% note *)
 Kan der gives nogen tydeligere Udtalelse af, at Sagbevægelsen i Chemismen har været, ikke en Realbevægelse, men en Begrebsbevægelse, og den hele Forstaaelse bygget paa explicative Functioner? Fremstillingen af Teleologien begynder saaledes: Wo Zweckmäßigkeit wahrgenommen wird, wird ein Verstand als Urheber derselben angenommen, für den Zweck also die eigene, freie Existenz des Begriffes gefordert“.\footnote{Anfr.\ Skr.\ S.~209.}% note **) (the note itself is printed with "*)")
 % PRINTER: no opening quotation mark before "Wo Zweckmäßigkeit"; closing „…“ unbalanced as printed
% PRINTER: second footnote on p. 389 printed with mark "*)", sense "**)" (anchor in text is "**)")
Kan der gives nogen tydeligere Udtalelse af, at Forstand paa Øiemedet med Tingene beroer paa en Forstaaelse af den Forstand, der frembringer Tingene, at altsaa de for Teleologien afgjørende Functioner ere interpreterende? Her er ikke først et Par Systemer af objective, mechaniske og chemiske Slutninger, der paa immanent objectiv Maade danne Overgang til og fuldende sig i Teleologie. Nei, Sagen forholder sig omvendt! Her begyndes paa de lavere Trin med en stiltiende Ombytning af Gjenstandens Objectivering med en tilsvarende Objectivering af Gjenstandens
% --- p. 390 ---
\opage{390}Begreb, altsaa med en Ombytning af indicative Functioner med explicative. Ved explicative Functioner, idet Opmærksomheden stadig er fæstet paa Begrebsobjectet, opdager saa Tænkeren i de mechaniske og chemiske Processer, vel ikke egentlige Slutninger, thi Slutninger i egentlig Forstand forudsætte en dømmende og sluttende Subjectivitet, men dog Noget, der ligner Slutninger, Sammenslutninger nemlig af Extremer med en Midte; ved at forfølge disse Slutningsanalogier med stigende Eftertryk paa Begrebsobjecterne og de tilsvarende Explicationer, naaer Tænkeren omsider et Vendepunkt, da der gaaer et Lys op for ham, og det bliver klart, hvad Begrebet af den objective Slutning egentlig er: at det er Øiemedet, Formaalet, Hensigten, der giver Extremernes objective Sammenslutning Præg af Slutning. Med Blikket fæstet paa Hensigten og „det Hensigtsmæssige“ gjør Tænkeren tilsidst den Opdagelse, at det, hans Forstand i alle lavere Instanser har været beskæftiget med, ikke er det virkelige Object, ikke engang Objectets objective Begreb for sig, men Begrebets oprindelige Mening, den guddommelige, i Begreb og Gjenstand sig udtalende Forstand: enhver objectiv Slutning er uden Forskjel teleologisk.

Med denne Forudsætning skulle vi da i det Følgende udvikle den objective Slutnings \emph{Formalbegreb}, \emph{dens Forhold til Magtens Ord og til det kategoriske Magtsprog}.

\noindent\textit{αα)} \emph{Den objective Slutnings Formalbegreb.}\\
Forskjellen imellem den subjective og den objective Slutning er oprindelig bestemt ved Forskjellen mellem Functioner af Viden og Functioner af Magt. Den objective Slutning bestemmes ved objective Domme, de objective Domme ved det kategoriske Magtsprog. Functionerne af den subjective Slutning ere alle indicative; Forskjellen imellem indicative, explicative og interpreterende Functioner, en Forskjel, der bestemmes ved Dom og Gjenstand, kan først faae sin fulde Betydning i den objective Slutning. Her svare da de indicative Functioner
% --- p. 391 ---
\opage{391}til den første umiddelbare Objectivering, den Objectivering, hvorved Gjenstanden opfattes som Gjenstand; til den fordoblede Objectivering, den Objectivering, hvorved Gjenstandens Begreb bliver Gjenstand, nemlig intellectuel Gjenstand, svare de explicative Functioner: ved de explicative Functioner fremkommer et System af objective Domme. Endelig, idet Objectiveringen gaaer bestemmende tilbage paa den objectiverende Subjectivitet, der ikke blot søger Begrebet i Kjendsgjerningen, men i Kjendsgjerningens Begreb den oprindelige Mening, viser det sig, at Functionerne af den objective Slutning væsentlig ere interpreterende.

Overseer man Functionernes eiendommelige Charakteer som interpreterende, da skal det ikke være muligt at opvise et eneste Exempel paa en sand objectiv Slutning. Man vælge Solsystemet, Forholdet mellem Sjæl og Legeme, Staten med dens Borgere, Regjering og Love, kort, hvilket Exempel, man vil: saalænge det valgte Exempel behandles umiddelbart objectivt, d.\ e.\ saalænge der umiddelbart sees paa de existerende Led og disses stedfindende Relationer, vil det være umuligt at komme ud over Kategorier som Hele og Dele, Aarsag, Virkning og Vexelvirkning. Her er i Factorernes Gjensidighed ingen objective Domme, og hvor der ingen objective Domme er, der er der heller ingen objective Slutninger. Det Høieste, man driver det til, er at paavise en totaliseret Modsætningseenhed, en Vexelvirkning af tre Totaliteter, hvori den ene Totalitet afgiver Eenheden, naar de to andre danne Extremerne.\footnote{Trendelenburg gjør med Hensyn paa Hegels Exempler, adskillige træffende Bemærkninger; hvad navnlig Staten som et System af tre Slutninger angaaer, siger han: „Nach dieser Ansicht wächst das Ganze dadurch kräftig zusammen, daß das Besondere, Einzelne und Allgemeine wechselsweise und gegenseitig Grund und Folge wird. Daß sich die Thätigkeiten des Ganzen und der Theile innig % note breaks here: p. 391/392
durchdringen, das bildet allerdings die Selbsterhaltung des organischen Ganzen. Will man die zusammenwirkenden Glieder das Allgemeine, Besondere und Einzelne nennen: so hat auch das im Sprachgebrauche einen Grund. Aber man verwirrt die Sache, wenn man das Analogon eines Schlusses bildet“ (Logische Untersuchungen. Zweiter Band. S.~278--79). Men idet han saa skal til at gjøre Rede for, hvori det Utilfredsstillende da egentlig bestaaer, geraader han aabenbart paa Afveie. Han gjør nemlig Syllogismen til en saa typisk Form, at ingen Forening af Extremer og Midte fortjener Navn af Slutning, med mindre den i alle Maader indrettes efter Reglerne for en Syllogisme. At nu Staten ikke kan blive til en Syllogisme, er øiensynligt, „denn die Bedürfniße sind nicht als Art der Allgemeinheit des Staates, noch die einzelnen Bürger als Individuen oder Art eines Geschlechts den Bedürfnissen subsumirt“. Af samme Grund kan Maalet med dets Midler da heller ikke opfattes som en Slutning. „Im Zwecke \dots\ soll das Mittel den \textit{terminus medius} bilden, durch den sich die subjective Vorstellung mit der Objectivität zusammenschließt. Die drei Termini des Schlusses wären in einem einfachen Beispiele: deutlich sehen wollen das eine Extrem, das optische Glas der Mittelbegriff, das wirkliche deutliche Bild das andere Extrem. Mag man hier vergleichungsweise sagen, daß sich das Subject mit der objectiven Welt, der es seinen Zweck abgewinnt oder einbildet, zusammenschließt: dies Bündniß ist noch kein logischer Schluß. Wie will man in den genannten drei Terminis das deutlich sehen dem optischen Glas als den Umfang dem Inhalt unterordnen? oder gar das wirkliche deutliche Bild den deutlich sehen wollen so subsumiren, wie sonst der Unterbegriff in den logischen Umfang des Oberbegriffs fällt? Man kann doch die wirkliche Ausführung nicht als eine Art der vorgestellten betrachten. Wenn der reale Schluß, wie er von Hegel in die Objectivität eingeführt ist, wirklich dem logischen entspräche: so müßte er sich in die vollständige Form eines Syllogismus fassen lassen (Anfr.\ Skr.\ S.~276).}% note *) % PRINTER: the quotation opened at „Im Zwecke is never closed; % PRINTER: printed "Anfr. Skr.", sense "Anf. Skr." (anførte Skrift)

% --- p. 392 ---
\opage{392}Fra de interpreterende Functioners Standpunkt vil det derimod vise sig, at et hvilketsomhelst Exempel, det være ved første Øiekast nok saa ufordeelagtigt, naar det blot kan afgive
% --- p. 393 ---
\opage{393}to Extremer og en Midte, lader sig opfatte som en objectiv Slutning. Vi vælge et af de Exempler, Trendelenburg har benyttet imod Hegel. Svovlsyre og Kalk være den \textit{terminus medius}, altsaa det Særskilte, hvorved det Almene, Begrebet Gips, slutter sig sammen med det Enkelte, Existensen Gips. Her skal, vel at mærke, ikke føres Beviis for, at Gips nu ogsaa i endelig Forstand er en objectiv Slutning; her skal blot vises, hvad det vil sige, at betragte Gipsdannelsen under Synspunkt af Slutning. Fra de interpreterende Functioners teleologiske Synspunkt bliver da Existensen Gips udtrykkelig opfattet som det Produkt, paa hvis Dannelse den forudgaaende Dannelse af Svovlsyre og Kalk er anlagt; \emph{Meningen} med Svovlsyren og Kalken er at søge i Gipsen. I Begrebet er Gipsen anticiperet; med denne Anticipation er Begyndelsen skeet. Begyndelsen er Slutningens første Præmis; thi Gipsens Begreb kunde umulig bestemmes, med mindre Bestemmelsen af Svovlsyre og Kalk var med i Anticipationen. Svovlsyre og Kalk i særskilt Tilværelse afgive nu en Midte, ikke i abstract objectiv Forstand, men i Kraft af det Forhold, hvori de særskilt Tilværende ere satte paa den ene Side til den bestemte Begyndelse, paa den anden Side til Enden, Resultatets Opnaaelse, Virkeliggjørelsen af det anticiperede Begreb. Begrebet forekommer her tre Gange og paa tre forskjellige Maader: 1) som anticiperet Begreb, det Almene, Begrebet i subjectiv Mulighed, 2) som Begreb i de objective Elementer, det Særskilte, Begrebet i objectiv Mulighed, og 3) som realiseret Begreb, det Enkelte. Det er, ifølge vort Foregaaende, uimodsigeligt, at ethvert Begreb, der skal realiseres, maa gjennemgaae disse tre Phaser.

Intet Begreb kan realiseres umiddelbart. Skulde et Begreb realiseres umiddelbart, da maatte det vel skee paa een af to Maader: enten maatte dets subjective Væren i Anticipationen ved et absolut Magtsprog, altsaa uden Betingelser og Mellemled, forvandles til objectiv-virkelig Tilværen, hvad
% --- p. 394 ---
\opage{394}der da vilde forudsætte et Mirakel; eller Begrebet maatte, uden al Anticipation, begynde med at være saaledes nedsænket i de objective Elementer, at det realiseredes i blind Objectivitet, og fremkom alene som Resultat, som umiddelbart Produkt. Men Produktet er som Produkt kun Gjenstand, ikke Begreb. Som Resultat af Svovlsyrens Forening med Kalken er Gipsen et chemisk Produkt, et Salt; men dette Salt er i sin Umiddelbarhed intet Begreb. For at Begrebet kan komme tilsyne i Gipsen, maa Gipsen som Foreningens Produkt reflectere sig i Foreningen selv og de forenede Bestanddele; i denne Reflexion skinner Subjectiviteten og med Subjectiviteten det Bestanddelene og deres Forening bestemmende Begreb igjennem: kun under Forudsætning af, at her er et i subjectiv Form bestemmende Begreb, og at de Elementer, der underlægges Begrebsbestemmelsen have objectiv Tilværelse, kan der være Tale om et Begrebs Realisering.

Men lægges nu afgjørende Eftertryk paa, at det er det anticiperede, ved Anticipationen bestemte, sine objective Elementer ved et System af executive Functioner bestemmende Begreb, der realiseres, saa er da ogsaa Realiseringen med det Samme bestemt som Slutning. Det Objective i Slutningen er, at Processen hører op, at det tager en Ende, at der kommer et Resultat; det Subjective, der bestemmer Realiseringsprocessen som en Begrebsbevægelse og giver den Præg af Slutningsact, er, at det subjective Begreb møder sig selv i det objective Produkt; gjennem Midten gaaer Enden tilbage til sin Begyndelse; her er ingen Afbrydelse, men et fuldendt Kredsløb, et realiseret Formaal, en Slutning. Hvor beslægtet den objective Slutning er med den virkelige Afslutning og denne igjen med det opnaaede Øiemed, kan endogsaa sees af Talemaader som: „naar Enden er god, saa er Alting godt“, \textit{respice finem}, o.\ dsl. I disse Talemaader har den sunde Sands med kraftig Eensidighed fremhævet snart den objective, snart den subjective Side, og just derved viist, at
% --- p. 395 ---
\opage{395}Sagen maa sees fra begge Sider. Hvad det første Udsagn angaaer, kan jo vistnok det Objective fremhæves saa umiddelbart, at den gode Ende, der gjør Alting godt, bliver en eventyrlig Hændelse, et lykkeligt Udfald, der pludselig bringer de haardt Prøvede til at glemme alle Gjenvordigheder. Men en saadan Ende, et saadant Udfald er kun et Spil af Tilfældet, et Skjæbnens Indfald, men ingen Slutning. At her ingen Slutning er, men kun en Afbrydelse, en Standsning, en vilkaarlig Ende, er indlysende saavel fra objectiv som subjectiv Side. Objectivt seet, er Manglen af Slutning kjendelig derpaa, at Ulykkernes Afbrydelse igjen kan afbrydes, og den lykkelige Ende vise sig som Begyndelse til nye Ulykker. Subjectivt seet, er Manglen af Slutning kjendelig derpaa, at her intet Begreb, ingen Plan er, der realiseres. Naar Enden i Sandhed skal være god og gjøre Alting godt, da maa den i Udvikling og Forklaring svare til sin Midte og Begyndelse; den maa da, med andre Ord, være en Slutning. Thi hvad er vel Begyndelsen? Det anticiperede Begreb! Og Midten? Begrebets objective Elementer! Og Enden? Det i Elementerne realiserede Begreb! Under Begrebets Virkeliggjørelse er Afslutningen en Slutningsact, en ideel Magthandling, en logisk-antilogisk Act af den sluttende Subjectivitet.

Da Slutningen, den sande objective Slutning, skal samle og forene, hvad der er spredt i dens Midte og anlagt i dens Begyndelse, saa indsees heraf, at det saavel i Digt som i Virkelighed maa høre til de store Opgaver, at naae en Slutning. I hvilket Afhængighedsforhold den sande Slutning maa staae til den sande Begyndelse, skjønnes af de bekjendte Talemaader: „Vel begyndt er halvt fuldendt“; „en god Begyndelse faaer aldrig nogen slet Ende“, o.\ dsl. Det er just det Mærkværdige ved disse Talemaader, at naar de tages fra Endelighedens Side og sees overfladisk, ere de kun halvsande, ja vel endog usande. I den endelige Tilværelse falde Begyndelse, Midte og Ende jo netop, ifølge den for Endeligheden
% --- p. 396 ---
\opage{396}charakteristiske Udstykning, saaledes ud fra hinanden, at Begyndelsen kan, som det hedder, tegne godt, og endda kan der under Sagens Fremgang indtræde høist sørgelige Tilfælde, ligesom det om Morgenen kan tegne til at blive skjønt Veir, og dog kan, naar Middagen kommer, Himlen være overtrukken, og Dagen ende med Storm og Uveir. Naar Forholdet mellem Enden og Begyndelsen derimod sees fra Væsenets Side, som i en objectiv Slutning, er Talemaaden fuldkommen sand; Begyndelsen er i den objective Slutning ikke Andet end den anticiperede Ende. I det ved sin Forudsætning bestemte Resultat, i den ved sin Begyndelse præformerede Ende kan Intet fremkomme, uden hvad Forudsætningen og Begyndelsen selv indeholder. Den gode Begyndelse er Anlæget til den gode Ende; udvikles Anlæget, naaer Begyndelsen gjennem Midten til Slutningen, maa Enden absolut blive god: det er i denne Forstand bogstavelig sandt, at en god Begyndelse aldrig kan faae nogen slet Ende.

I Talemaaden: \textit{quidquid agis, prudenter agas, et respice finem}, betyder Ordet \textit{finis}, ligesom ogsaa det græske \textit{τέλος}, paa een Gang baade Ende og Formaal. Men at Enden er Formaal, og det realiserede Formaal Ende, betyder jo netop, at Enden er Slutning, og Slutningen teleologisk. Der gives altsaa ikke, hvad Hegel antager, abstract objective, mechaniske og chemiske, Slutninger, som, idet de selv falde udenfor det Teleologiske, skulde indlede Teleologien; thi enhver Slutning, hvor objectiv den end i sine Kjendsgjerninger mon være, er dog kun Slutning formedelst det i Slutningsacten selv Subjective og Teleologiske.

Sammenstille vi nu, efter disse Oplysninger, den objective Slutning med den subjective, da er Forskjellen og Modsætningen mellem dem saa fremtrædende, at den Maade, hvorpaa Hegel har villet hævde deres Identitet, maa ansees for uholdbar. Ikke desmindre er her i Modsætningen en Eenhed; det er en Modsætningseenhed af Viden og Magt, hvortil vi
% --- p. 397 ---
\opage{397}igjen ere henviste. Ligesom Functionerne af den subjective Slutning have deres Grundlag i en med den ideelle Magt identisk Viden, saaledes have Functionerne af den objective Slutning deres Grundlag i en med Viden identisk Magt. Den subjective Slutningsact er en theoretisk Tankehandling med Præmisser af logiske Domme; den objective Slutningsact er, fra den frembringende Subjectivitets Synspunkt, en praktisk Tankehandling med antilogiske Domme til Præmisser. Fra den betragtende Subjectivitets Standpunkt bliver den i Frembringelsen udtalte Slutning bestemt ved interpreterende Functioner. Ved den subjective Slutning gjælder det om en middelbar Viden af det i Præmissernes Følgesætning udtrykte Begreb: den subjective Slutning er udelukkende theoretisk. Den objective Slutning er derimod baade af praktisk og theoretisk Betydning. For den frembringende Subjectivitet er Slutningen praktisk, for den betragtende derimod theoretisk: absolut forstaaet maa den frembringende Subjectivitet være betragtende, den betragtende frembringende. Det er den middelbare Virkeliggjørelse af det anticiperede Begreb, hvorpaa det praktisk, og en Indsigt i denne Virkeliggjørelse, hvorpaa det theoretisk kommer an.

I de anførte Træk er Modsætningen fremhævet; det vil nu være et Spørgsmaal, hvorvidt Identiteten paa Grundlag af denne Modsætning lader sig paavise. Opgaven er at undersøge, hvorvidt den subjective Slutning lader sig bringe ind under Bestemmelser, der oprindelig gjælde for den objective, og omvendt.

Synspunktet for den objective Slutning er, som vi have seet, afgjørende teleologisk; dette Synspunkt kan i formel Forstand ogsaa gjælde for den subjective Slutning. Stilles Forbindelsen af Subject og Prædicat i Følgesætningen som en Opgave, saa er jo vistnok Løsningen af denne Opgave Slutningens Formaal, og hele Synsmaaden forsaavidt teleologisk. Spørges her, for at blive ved det forslidte Exempel, om Sub%
% --- p. 398 ---
\opage{398}jectet Cajus lader sig forene med Prædicatet dødelig, saa ere Præmisserne Midler, Følgesætningen Maal. Men dette er, som sagt, kun formelt.

Synspunktet for den subjective Slutning er afgjørende logisk, ved Reduction til sit endelige Schema syllogistisk. Det Logiske og Syllogistiske kjendes, navnlig i den kategoriske Syllogisme, derpaa, at Subjectet i enhver af Præmisserne subsumeres under sit Prædicat, saa at Følgesætningen bliver et nødvendigt Udtryk for Eenheden af den dobbelte Subsumtion. At Subjectet Cajus subsumeres under Prædicatet dødelig, er et Udtryk for, at Cajus er subsumeret under det (Mennesket), som er subsumeret under dødelig. Spørgsmaalet er nu, om et saadant System af Subsumtioner ogsaa passer paa den objective Slutning. Ved Anvendelse af explicative og interpreterende Functioner gaaer det vel an. I Gipsens Dannelse slutter den bestemte Saltart (Gipsen) sig sammen med Slægten (Iltesalt), \textit{S}---\textit{E}---\textit{A}; men, da Præmisserne ere antilogiske, maa Subsumtionen skee i omvendt Orden. Ifølge Realiseringens første Præmis, kan Slægten kun virkeliggjøres i sine Arter; den virkeliggjorte Art bliver altsaa at subsumere under en Virkeliggjørelse af Slægten. Ifølge Realiseringens anden Præmis, kan Arten kun virkeliggjøres i sine Individer; det virkelige Individ (Gipsexemplaret) bliver altsaa at subsumere under en Virkeliggjørelse af Arten. Syllogistisk: En Realisering af Gipsbegrebet er en særlig Realisering af det almene Saltbegreb; i den virkelige Gips er Gipsbegrebet realiset; % PRINTER: printed "realiset", sense "realiseret"
altsaa: i den virkelige Gips er det almene Saltbegreb paa eiendommelig Maade realiseret. Praktisk: Først Slægtsbegrebet, derpaa Artsbegrebet, og saa Begrebets Virkeliggjørelse!

Saa tvungen og haartrukken denne Opfattelse end er, saa viser den dog Saameget, at man ikke, hvad den formelle Logik er saa tilbøielig til, kan afgrændse Subsumeringen paa mechanisk Maade. Opfatter jeg f.\ Ex.\ Sætningen: Gips er en chemisk Forbindelse af Svovlsyre og Kalk, som en Definition paa Gipsen, da har Prædicatet jo vistnok ikke større Omfang end
% --- p. 399 ---
\opage{399}Subjectet, og her kan forsaavidt ikke være Tale om Subsumering; men Sætningen kan jo ogsaa sees fra et andet Synspunkt; jeg kan nemlig opfatte Gipsen i en Enkeltforestilling for sig, og betragte Svovlsyrens Forbindelse med Kalken som en af mange mulige Forbindelser, altsaa som det Særskilte. Endelig er der da heller Intet til Hinder for, at den i Begrebet særskilte Forbindelse stilles imod sin Virkeliggjørelse saaledes, at Virkeliggjørelsen sees som det Almindelige, hvorunder ikke blot Forbindelsen af Kalk og Svovlsyre, men alle chemiske Forbindelser, saa sandt de skulle gjælde for mere end Muligheder, maae subsumeres. Men naar dette er indseet, skal det for Resten indrømmes, at ligesom det teleologiske Synspunkt er uadæqvat for den subjective, syllogistiske Slutning, der mangler de objective Mellembestemmelser, saaledes er den syllogistiske Subsumering tvungen og uadæqvat for den objective, i objective Elementer særlig begrundede Slutning.

Opfattes den objective Slutning, ikke fra den betragtende, men fra den frembringende Subjectivitets Synspunkt, da blive dens Præmisser antilogiske Domme. Gipsen, som anticiperet Begreb \emph{det Almene}, bestemmes da ved Svovlets, Iltens og Kalkmetallets særegne Forbindelse, \emph{det Særskilte}; denne Forbindelse bestemmes igjen ved sin Virkeliggjørelse til \emph{det Enkelte}, den existerende Gips. Dog heller ikke i denne Henseende mangler det paa Analogie med den subjective Slutning. Den subjective Slutning betinges nemlig, som vi have seet, ikke blot af Subsumtionen, men ogsaa af Disjunctionen. \textit{A}, det Almene, er enten \textit{B} eller \textit{C}, det Særskilte --- thi Særskiltheden ligger i den endnu uafgjorte Vexelbestemmelse; ved negativ Bestemmelse, d.\ e.\ ved at udelukke det ene Led: \textit{A} er ikke \textit{D}, bliver Særskiltheden indskrænket og saaledes individualiseret til det Enkelte: altsaa \textit{A}, det Almene, er \textit{C}, det i sin Existeren Enkelte.

Dersom den antilogiske Bevægelse ikke kunde eftergjøres logisk, kunde der ikke gives nogen objectiv Slutning; den objective Slutningsact maa nødvendig reflectere sig subjectivt i
% --- p. 400 ---
\opage{400}den frembringende Subjectivitet, men den subjective, Frembringelsen oversvævende Reflexion er logisk. Dersom den logiske Bevægelse ikke kunde eftergjøres antilogisk, vilde Overgangen fra de theoretiske til de praktiske Functioner, fra Functioner af Viden til Functioner af Magt, være umulig. I Anticipationen er Alt logisk, og dog ere alle logiske Bestemmelser saa antilogisk vendte, at det i Udførelsen virkelig Antilogiske er, hvad Bevægelser og Overgange angaaer, præformeret i det Logiske. Her er i Værkmesterens Viden ikke en blot Række logiske Domme ved Siden af eller parallelt løbende med en Række antilogiske; her er tillige et System af Domme, nemlig logisk-antilogiske Domme, der afgiver Bindeled for Extremerne. Til Systemet af de logisk-antilogiske Domme svarer nu et System af \emph{logisk-antilogiske Slutninger}, et Mellemsystem, hvori de subjective og objective Slutninger uophørlig omsættes i hinanden.

Hvad denne Omsætning har at betyde, kan belyses ved en nærmere Betragtning af Forholdet mellem Extremer og Mellemled. Er det nødvendigt, at Extremer og Mellemled maae i enhver, saavel objectiv som subjectiv, Slutning udgjøre tre Totaliter? % PRINTER: printed "Totaliter", sense "Totaliteter"
Det synes ved første Øiekast, som maatte Spørgsmaalet besvares med Nei. Mangfoldige Præmisser til subjective Slutninger ere jo Domme med eensidige Prædicater, rene Egenskabsprædicater, som f.\ Ex.\ alle Mennesker ere dødelige. Ligeledes ere Leddene i mange af de for objective Slutninger gjældende Præmisser saa elementære og eensidige, at de aldeles ikke kunne tage sig ud som Totaliteter. Sættes Gipsens Begreb imod den existerende Gips, kan man vel sige, at Extremerne ere Totaliteter; men hvorledes kan det i to Elementer faldende Mellemled siges at udgjøre en Totalitet? Vanskeligheden bliver ikke mindre, om vi, efter Hegels Exempel, sætte Svovlsyren som det ene Extrem, Kalken som det andet, Affiniteten som Midte, og Gipsen som den virkelige Slutning.

% --- p. 401 ---
\opage{401}Ogsaa hvad det Simultane og Successive angaaer, reise sig Betænkeligheder. Idet den subjective Slutning opfattes som Syllogisme, dannes Præmisserne successivt, Conclusionen med Følgesætning følger bagefter: Successionen er umiskjendelig. Forvandles den subjective Slutning derimod til Immanensslutning, saa forsvinder Successionen; Bevægelsen bliver Hvile, og Slutningsleddene simultane. Noget Lignende sees paa den objective Slutning. Vælge vi Gipsdannelsen til Exempel, saa ere Leddene successive; vælge vi Ligevægten, ere de simultane.

Ved nøiere Undersøgelse af disse tvivlsomme Operationsphænomener, vil Omsætningen af det Logiske i det Antilogiske og af dette i hiint efterhaanden blive paaviist. Hvad nu Forholdet mellem Leddenes Simultaneitet og Succession angaaer, for at begynde med dette, vil det snart vise sig, at Modsætningen af subjectiv og objectiv Slutning er begrundet i Identitet. Successionen i det Subjective, det Ideelle, forudsætter sit Simultane; thi det, der efterhaanden indlyser, indlyser jo som det, der allerede var; Sandhedens Vorden, d.\ v.\ s.\ dens successive Kommen til Bevidsthed, reflecterer sig i dens evige Væren. Hvad Subjectiviteten efterhaanden faaer at vide, faaer den \textit{a priori} at vide som Noget, der forud \emph{har været} og derfor er i evig Viden. Successionen i det Objective, det Reelle, er vistnok saaledes endeliggjort, at den endog bliver timelig; men just fordi det Successive er udstykket i Timelighedens Schema, er Simultaneiteten udvortes realiseret i Rummets Schema. Enten vi altsaa betragte Slutningen fra det subjective eller det objective, fra det logiske eller det antilogiske Synspunkt, see vi i begge Tilfælde en Vexelbestemmelse af det Simultane og det Successive.

At Vexelbestemmelsen i det Logiske er Vexelbestemmelsen i det Antilogiske modsat, forstaaer sig; men det ligger i Slutningens Væsen, at Identiteten mægles ved Modsætningen. Totalbevægelsen i det Logiske gaaer fra det Ideelle gjennem
% --- p. 402 ---
\opage{402}det Reelle tilbage til det Ideelle; her er med andre Ord to ideelle Extremer med et reelt Mellemled: \emph{Anticipationen er i denne Belysning meer end en Dom, den er en Slutning}.

Anticipationsslutningens første Extrem bestaaer i den Totalitet af logiske Bestemmelser, der oprindelig bestemme det anticiperede Begreb. Men da Anticipationen selv bestemmes ved objective Hensyn, saa vil Totaliteteten % PRINTER: printed "Totaliteteten", sense "Totaliteten"
af de for disse Hensyn objective Elementer afgive Slutningens reale Midte. Den reale Midte sees ved første Blik, ikke som Midte, men som Extrem. Fra ethvert af de objective Elementer udgaaer en antilogisk Reflex, der fra den objective Side er ligesaa bestemmende for den anticiperende Dom, som de logiske Momenter ere det fra den subjective Side. Desuagtet er det for det objective Syn ikke Realiteten i dens Umiddelbarhed som Realitet, hvorved her acquieseres. Realitetshensynet er kun Vehikel for Begrebsudviklingen: de antilogiske Bestemmelser ere kun Reflexer, Skjær og Farvetoner i det Logiske. Bevægelsen gaaer altsaa gjennem det Reelle tilbage til det Ideelle; det sidste Extrem, det ved de antilogiske Momenter articulerede Begreb, er nu formedelst den reale Midte forskjelligt fra og dog Eet med sin ideelle Begyndelse: Anticipationsslutningen er fuldendt.

At den fuldstændig udviklede Anticipation absolut maa ende som Slutning, kunde allerede forudsees i Læren om Dommen. Ved den for Anticipationen afgjørende Vexelbestemmelse af den disjunctive og den hypothetiske Dom, blev der lagt særdeles Vægt, ikke blot paa Valget af objectivt bestemte Elementer, men, hvad her fra dette Synspunkt er afgjørende, paa den Indflydelse, som Valget af de objective Elementer maatte udøve, ikke paa Begrebets Virkeliggjørelse, thi det forstaaer sig af sig selv, men paa Anticipationsbegrebets egen anticiperede Bestemthed. I Anticipationsslutningens reale Midte gaaer der altsaa en Strøm af Bevægelser fra
% --- p. 403 ---
\opage{403}det Ideelle mod det Reelle, fra det Reelle mod det Ideelle med overgribende Idealitet, altsaa en Bevægelse fra det Logiske mod det Antilogiske, og omvendt, med Overgriben af det Logiske. Anticipationsslutningen er altsaa, naar den sees i sin Midte, et fuldstændigt Udtryk for den Subjectivitetsbevægelse, hvori det Logiske uophørlig vender sig for at eftergjøre det Antilogiske, og derved slutte sig sammen med sig selv.

Fæste vi Blikket paa Virkeliggjørelsen, er det da ikke iøinefaldende, at \emph{Begrebets Realisering er meer end en Dom, at den fra sin objective Side er en fuldstændig Slutning}? At det anticiperede, i Anticipationen fuldt udviklede Begreb imidlertid danner, ikke, som det ved første Øiekast kunde synes, Slutningens første Extrem, men dens Midte, maa i nærværende Sammenhæng udtrykkelig fremhæves. Saalænge vi endnu kun betragtede Begrebets Virkeliggjørelse, i Eenhed med dets Anticipation, som een eneste objectiv Slutning, maatte Anticipationsbegrebet, d.\ v.\ s.\ det anticiperede Begreb, naturligviis blive Slutningens første Extrem. Nu derimod, efter at Anticipationen er bestemt som en Slutning for sig, maa Virkeliggjørelsen som særegen Slutning nødvendig udgaae fra Virkelighed og vende tilbage til Virkelighed; Realisationsslutningen maa, med andre Ord, have to reelle Extremer med en ideel Midte.

Den ideelle Midte er netop det anticiperede Begreb. Thi hvad er vel i objectiv Forstand det Første og Vigtigste for den, der har et Begreb, en Plan, han vil realisere? Dog vel aldrig at sidde og stirre paa Planen og repetere, hvilke og hvormange Fordringer det nu egentlig var, der laae i Begrebet; ved idelig Gjentagelse af slige Betragtninger kommer den frembringende Subjectivitet ingen Vegne med det Objective. Skal den Frembringende ikke blive hængende i anticiperende Reflexioner som Fluen i Vævet, da maa han, idet han skrider til Udførelsen, øieblikkelig og umiddelbart tage fat paa de objective Elementer, vælge sit Materiale; her er altsaa det
% --- p. 404 ---
\opage{404}første Extrem. I sin udvortes objective Tilværelse er Materialet ikke blot aldeles fremmed for, men, som det synes, endog aldeles ueensartet med det Begreb, som deri skal realiseres; nu begynder Bevægelsen fra Materialet mod Begrebet, fra det første Extrem mod Midten.

Lader os betragte denne Bevægelse lidt nærmere. Extremet er reelt, Midten er ideel; skal det være muligt at naae fra Extrem til Midte, maa det altsaa være muligt at omsætte Extremets Realitet i ideelle Reflexer, d.\ v.\ s.\ at opløse dets antilogiske Bestemthed i logiske Bestemmelser. Idet her begyndes med den umiddelbare Realitet, begyndes her just med et System af logiske Domme. Materialet bestaaer f.\ Ex.\ af Jern, Steen og Træ; hvor er det nu muligt for den Frembringende at gjøre Noget ud af et saadant Materiale, dersom han ikke fra sit Synspunkt kjender dets Beskaffenhed; paa Beskaffenheden beroer jo netop Brugbarheden. Men en udtømmende Erkjendelse af Materialets Beskaffenhed forudsætter et System af logiske Domme. Enhver Bestanddeel, ethvert af de særskilte Stoffer afgiver her Subjectet for en Mængde Prædicater, hvortil den Frembringende under Udførelsen af sin Plan nødvendigviis maa tage Hensyn.

Det System af Prædicater, hvori Materialets Begreb udvikles, mødes nu med det System af Prædicater, hvorved det anticiperede Begreb er bestemt; paa dette Møde beroer Gyldigheden af den ideelle Midte. Gjennem Prædicaternes logiske Væsenhed gaaer Formaalsbegrebet ind i sine objective Elementer, bliver bestemmende for deres Behandling; Bevægelsen gaaer altsaa fra det Ideelle mod det Reelle, fra det Logiske mod det Antilogiske. Den Elementernes Realitet, der i Form af Materiale afgav Slutningens første Extrem, svarer nu til den Realitet, det virkeliggjorte Begreb, der danner det sidste Extrem. Formedelst den ideelle Midte er det sidste Extrem identisk saavel med Begrebet selv som med det første Extrem, er, med andre Ord, det første Extrems Identitet med
% --- p. 405 ---
\opage{405}det virkeliggjorte Begreb. Bevægelsesfordoblingen i den ideelle Midte er igjen et Exempel paa, at det Antilogiske i enhver objectiv Slutning eftergjør det Logiske, for at det Logiske kan eftergjøre det Antilogiske.

Det begyndte med, at Begrebets Realisering viste sig som een Slutning; den ene Slutning er nu ved immanent Fordobling bleven til to Slutninger; hvad er herved vundet? Den objective Slutning har udviklet sit subjective Moment! Den Modsigelse, at den objective Slutning som Function af Slutningsacten nødvendig maatte være subjectiv, er opløst i Fordoblingen: Anticipationsslutningen har viist, hvorledes det Objective som Moment opgaaer i det Subjective; Virkeliggjørelsesslutningen har viist, hvorledes det Subjective bliver Moment for det Objective. Men det Subjective kan --- i Slutningen --- kun blive Moment for det Objective paa det udtrykkelige Vilkaar, at det Objective er blevet Moment for det Subjective. De to Slutninger, Anticipationen og Virkeliggjørelsen, kunne da ikke stilles som indbyrdes uafhængige overfor hinanden; Anticipationen har kun Betydning som Forudsætning for Virkeliggjørelsen, og denne kun Betydning som Følgesætning af hiin: de \emph{to Slutninger danne altsaa Extremer for en tredie}, der hævder begges Gyldighed ved at optage dem begge som Momenter.

Den tredie, de to særskilte Slutninger omsluttende Slutning maa igjen have sin Midte og kjendes paa Midten. Er nu Anticipationsslutningen med sin reelle Midte det subjective Extrem, Virkeliggjørelsesslutningen med sin ideelle Midte det objective, da ligger det i Sagens Natur, at den til disse Slutningers Slutning svarende Midte maa være de to særskilte Midters concrete Eenhed. Men enhver af de to særskilte Midter er som Midte selv den hele særskilte Slutning, den reelle Midte hele Anticipationsslutningen, den ideelle hele Virkeliggjørelsesslutningen; de to Midter danne ved deres \emph{Forskjellighed} den tredie Slutnings Extremer. Men heraf
% --- p. 406 ---
\opage{406}følger igjen, at de to Midter maae ved deres \emph{Identitet} tillige danne den tredie Slutnings Midte.

Og hvad skal saa Alt dette sige? Det skal sige, at den objective Slutning er et System af tre Slutninger, at Bevægelsen i dette System er uendelig Hvile. Bevægelsen gaaer fra Extrem til Extrem; endelig seet, gaaer Hvilen over i Bevægelse. Men ethvert Extrem er selv den hele Slutning: uendelig seet, gaaer Bevægelsen over i Hvile. Det Subjective gaaer fra Extrem til Extrem gjennem den reelle Midte, det Objective fra Extrem til Extrem gjennem den ideelle Midte: det Subjective og det Objective gaae i modsat Retning. Men den ideelle Midte er reel, den reelle er ideel: det Subjective og det Objective gaae med modsat Retning dog i samme Retning. \emph{Umiddelbart} er dette en Modsigelse; men i den objective Slutning er Alt \emph{uendelig middelbart}, og Modsigelsen hævet.

Idet det Logiske paa uendelig Maade eftergjør det Antilogiske, er Tankehandlingen fra alle Sider gaaet sammen med Magthandlingen; Viden er efter sin Realitet blevet identisk med Magten. Idet det Antilogiske paa uendelig Maade eftergjør det Logiske, er Magthandlingen i alle Punkter gaaet sammen med Tankehandlingen; Magten er efter sin Idealitet blevet identisk med Viden. Den teleologiske Identitet af Tanke- og Magthandling er gjennemklaret.

Anvendes nu den saaledes gjennemførte Formudvikling paa Bestemmelsen af de endelige Extremer, hvorledes disse end stilles til den respective Midte, da er Totaliteten i alle tre Led let at opdage. Hvad Exemplet med Gipsen angaaer, kan her naturligsviis ikke være Tvivl om, at Realbegrebet Gips er en Totalitet, og den virkelige Gips ligeledes. Tvivlen falder alene paa Midten. Men see nu Følgende: „Begrebet Gips er et --- ved Forbindelse af Svovlsyre og Kalk --- eiendommelig bestemt Realbegreb“; her have vi Anticipationsslutningen med de to ideelle Extremer og den
% --- p. 407 ---
\opage{407}reelle Midte. „Elementerne: Svovl, Ilt og Kalkmetal forenes --- ifølge Gipsbegrebets Fordring --- til virkelig Gips“; her have vi Virkeliggjørelsesslutningen med de to reelle Extremer og den ideelle Midte. Lade vi nu, i immanent Reflexion af Slutningens Væsen, den ideelle og den reelle Midte med deres indbyrdes Modsætning danne Extremer, og ved deres Identitet igjen Midte for en ny Slutning, saa er Totaliseringen af Slutningens Led alsidig gjennemført. Det samme Resultat udkommer, naar vi opfatte Gipsslutningen saaledes, at Svovlsyren bliver det ene Extrem, Kalken det andet, og deres Affinitet den væsentlige Midte.

Fra et chemisk Standpunkt vilde det, reentud sagt, være % PRINTER: printed "reentud" without a space, sense "reent ud"
Galimathias, om Nogen paastod, at enten Svovlsyren for sig eller Kalken for sig i Grunden var Gips; men fra et logisk Synspunkt bliver her, hvad Gipsslutningen angaaer, noget Andet at tage Hensyn til. Ikke Svovlsyren for sig som existerende Syre, men Svovlsyren som et Gipsslutningens Extrem er det, som her paa eensidig Maade repræsenterer Begrebets Totalitet. Kun i Forhold til den hele Slutning kan det enkelte Led blive Extrem. Idet Svovlsyren er sat som Extrem for Gipsen, er dens modsvarende Extrem, Kalken, udtrykkelig forudsat. Det for ethvert Extrem gjældende System af Egenskaber faaer alene Betydning derved, at begge Egenskabssystemer reflecteres i eet System, altsaa derved, at to eensidige Totaliteter reflecteres i een alsidig Totalitet; altsaa derved, at denne ene alsidige Totalitet særlig reflecteres i to eensidige, ikke blot \emph{ved} hinanden, men ogsaa \emph{for} hinanden bestemte Totaliteter. Denne Synsmaade tjener da tillige til Bekræftelse paa, at den objective Slutning væsentlig er teleologisk. Som Extrem i Gipsslutningen har Svovlsyren ikke blot visse Egenskaber, \emph{saa at} den i Forbindelse med Kalken \emph{kan}, men tillige udtrykkelig, \emph{for at} den under visse Betingelser \emph{skal} kunne danne Gips. Hvor en saadan Teleologie hverken gjælder eller kan gjælde, der har det (som f.~Ex.
% --- p. 408 ---
\opage{408}i Chemien) heller ingen Betydning at tale om objective Slutninger; hvor Teleologien derimod gjælder, gjælder ogsaa den objective Slutning.

\textit{ββ)} \emph{Magtens Ord i den objective Slutning.} Hvad Magtens Ord betyder, saavel efter dets Idealitet som efter dets Realitet, i hvilket Forhold de anticiperende, executive og explicative Functioner maae tænkes at staae til Magtordets Udtalelse i Dom og Gjenstand; paa hvilken Maade de ideale Magtfunctioner som imperative, indicative og interpreterende indgaae i Bestemmelsen af Dommens og Gjenstandens Idee; hvorledes den anticiperende Dom udvikler sig til anticiperende Slutning: disse og lignende Hovedpunkter ere i vort Foregaaende allerede afhandlede og belyste fra saa mange Sider, at en videre Udførelse i samme Retning kun vilde tage sig ud som ufrugtbare og trættende Gjentagelser.

Udviklingen af den objective Slutning er allerede ført til det Vendepunkt, at de logiske Functioner nu mødes med Functionernes Logik. I Functionernes Logik fik den objective Slutning særlig Betydning derved, at „Totaliteternes System“\footnote{Jvfr. Grundideernes Logik. I. S.~434 flgd.}% note *)
{} særlig maatte bestemmes ved systematiske Forandringer i de tilsvarende Functioners Energie. Paa den væsentlige Forandring i Functionernes Energie blev da den objective Forskjel imellem den mechaniske, den mechanisk-dynamiske og den organiske Slutning saaledes kjendelig, at det Objective i Overgangen fra den ene af disse Slutninger til den anden med det Samme lod sig tilsyne. Her, under Synspunktet af logiske Functioner, maa den objective Slutning i dens Forhold til Magtordet opfattes og bestemmes saaledes, at det kan indsees, hvilken Forskjel der er imellem at henføre de objectiverede, i Objectiveringen realiserede Begreber til den psychologiske og til den ontologiske Subjectivitet.

% --- p. 409 ---
\opage{409}Saalænge de objectiverede, i Gjenstandene virkeliggjorte Begreber, hvad Begribningen angaaer, udelukkende henføres til den psychologiske Subjectivitet, vil den endelige Adskillelse mellem Begreber og Gjenstande, mellem aprioriske og empiriske Begreber, mellem Tænken og Væren aldrig kunne hæves. Paa dette Standpunkt kan den Erkjendende vel ved subjective, paa Kjendsgjerninger støttede Slutninger arbeide sig frem til objective Resultater; men i de objective Resultater kan den endelige Viden aldrig naae de objective Slutninger. For at klare de objective Slutningers Realitet og derved ophæve den endelige Udstykning af aprioriske og empiriske Erkjendelser, maa Logiken henføre de objectiverede, i Objectiveringen realiserede Begreber, ikke til den psychologiske, men til den ontologiske, i Viden og Magt overgribende Subjectivitet. En Belysning af de objective Slutningers Realitet tjener da med det Samme til at opklare det Logiske i Forskjellen mellem Philosophiens og Fagvidenskabernes Realerkjendelse. Hvad dette betyder, skulle vi i det Følgende stræbe at paavise; vi vælge til Exempler \emph{Tyngden}, \emph{Affiniteten} og \emph{den organiske Proces}.

1) \emph{I Magtens Ord er Tyngden en objectiv Slutning.} Saavel i Logiken som i Erfaringsvidenskaben maa der udtrykkelig gjøres Forskjel imellem Tyngdens Begreb og den virkelige Tyngde; men i Erfaringsvidenskaben bestemmes Forskjellen fra et ganske andet Synspunkt end i Logiken. Hvad lærer Physiken om Tyngden? „Ethvert Legeme, som er understøttet af et andet“, siger Physiken,\footnote{Naturlærens mechaniske Deel af H.~C. Ørsted. Tredie Udgave ved Holten. Kjøbenhavn 1859. S.~154.}% note *)
{} „udøver et Tryk paa samme, men naar det ikke er understøttet, falder det lodret ned mod Jordens Overflade. Denne Legemernes Egenskab kaldes Tyngde“. Tyngden er altsaa ifølge denne Forklaring en \emph{Egenskab}. Strax efter
% --- p. 410 ---
\opage{410}hedder det\footnote{Anfr. Skr. S.~155.}% note *)
{} „Tyngden er ikke en Jorden eiendommelig Kraft, men Følgen af en almindelig Tiltrækning, der finder Sted mellem alle Legemer“. Fra dette Synspunkt bliver Tyngden altsaa bestemt ved to nye Kategorier, Kategorien \emph{Kraft} og Kategorien \emph{Følge}, en Følge, hvis virksomme Grund er en \emph{almindelig Tiltrækning}. Naar her saa til Slutning tilføies, „at den Kraft, vi paa Jorden kalde Tyngde, er det almindelige Kraftbaand, hvormed det store Verdensalt sammenholdes“,\footnote{Anfr. Skr. S.~174.}% note **)
{} saa glider ved dette „almindelige Kraftbaand“ Kategorien Tyngde umærkelig over i Kategorien Magt; Tyngden er altsaa en \emph{Magt}.

At nu Tyngden virkelig lader sig opfatte fra alle tre Sider, som \emph{Egenskab}, som \emph{Kraft} og som \emph{Magt}, uden at den ene af disse Opfattelser paa Fagvidenskabens Standpunkt geraader i Strid med den anden, skal villig indrømmes. Men Et er Realbegrebets Dannelse i Fagvidenskaben, et Andet er dets Udvikling i Logiken. Naar Logiken skal til at undersøge, hvorledes det Empiriske i det naturvidenskabelige Tyngdebegreb forholder sig til det Aprioriske, indstille sig Vanskeligheder, som paa den endelige Videns Trin umulig lade sig hæve.

For at forklare det Almene, det Nødvendige og forsaavidt Aprioriske i Tyngden, subsumerer man de paa Jorden forekommende Tyngdephænomener under „den almindelige Tiltrækning“; for at forklare den almindelige Tiltrækning, beraaber man sig paa den alle materielle Dele iboende Tiltrækningskraft; for at forklare den almindelige Tiltrækningskraft, henviser man til den for alle Legemer fælles Egenskab; men for at forklare denne Egenskab, hvad gjør man saa? Dommen om Legemernes fælles Egenskab, den Dom, at
% --- p. 411 ---
\opage{411}alle Legemer ere tunge, hører, efter hvad vi paa et tidligere Stadium\footnote{See ovf. S.~264--75.}% note *)
{} have oplyst, til de antilogisk-apodiktiske Domme.

Hvorledes det Empiriske og forsaavidt Tilfældige i de antilogisk-apodiktiske Domme strider med det Nødvendige og forsaavidt Aprioriske, er da kjendeligt nok paa den mechaniske Physiks Forklaring af Tyngden. Skulde en for alle Legemer fælles Egenskab virkelig lade sig opfatte som Tyngdephænomenernes Grundaarsag, da maatte denne Egenskab \textit{a priori} lade sig udlede af Materiens Begreb. Men til Materiens Begreb som Materie hører jo, ifølge Erfaringslæren, kun det, at opfylde Rummet. „Det Rumopfyldende i Legemerne kalde vi \emph{Materien}, det begrændsede opfyldte Rum et \emph{Legeme}“. Men af det blotte Begreb om Rumopfyldning vil en almindelig Tiltrækningskraft aldrig med logisk Nødvendighed lade sig aflede. Skulde det være muligt at aflede den almindelige Tiltrækning af Materiens Begreb, da maatte det jo ogsaa være muligt at aflede Tiltrækningsloven af Materiens Begreb. Men hvor misligt det paa den endelige Videns Standpunkt er at ville bestemme Tyngdeloven \textit{a priori}, sees da blandt flere Exempler tydeligt nok af den bekjendte Strid imellem Buffon og Clairaut. Buffon har Ret i at gaae ud fra den Antagelse, at Tyngdeloven maa svare til et med objectiv Nødvendighed gjældende Fornuftbegreb, altsaa til et Begreb, hvis Gyldighed er \textit{a priori}. Men det Uholdbare i hans saavel logiske og metaphysiske som mathematiske Argumenter viser, at han har Uret i at ville construere Naturens aprioriske Fornuftbegreber \textit{a priori}.\footnote{Om Stridens logiske Momenter jvfr. Forelæsn. over „Phil. Propæd.“ fra Universitetsaaret 1861--62. S.~116--118.}% note **)
% NOTE: a speck before the full stop after "118" (text layer "118..") is not a printed point.

Et Realbegreb som Tyngden kan ikke sammenstykkes af ueensartede Bestanddele, saa at en objectiv Forudsætning og en subjectiv Tanke, et empirisk Element og en apriorisk
% --- p. 412 ---
\opage{412}Kategorie, et tilfældigt Indhold og en nødvendig Form skulde ved sindrige Combinationer kunne fyldestgjøre Videns og Videnskabens høieste Fordring. Paa alle Punkter fordrer Tanken det Aprioriske, i alle Punkter standser Anskuelsen ved det Empiriske. De ueensartede Bestanddele, hvoraf et Erfaringsbegreb ved første Diekast synes at være sammensat, maae i Begribningen selv forvandles til eensartede Momenter: i Erfaringsbegrebet som Begreb ere alle Bestemmelser aprioriske; i Begrebet som Gjenstand, Erfaringsbegrebets Existens, ere alle Bestemmelser empiriske\footnote{Hvis man overseer dette, vil man neppe kunne unddrage sig Conseqvensen af følgende Discussion: „A. Du vil da vel indrømme mig, at det Aposterioriske i det Hele forholder sig til det Aprioriske som Sandsning til Tænkning? B. Naturligviis. A. Jo mere Sandseligt der er i et Begreb, desto mere indeholder det af et Aposteriorisk. B. Forstaaer sig. A. Men det Sandselige er jo, som vi have seet, et Utænkeligt, Noget, der ikke er til for Tanken, ikke lader sig tænke? B. Maa indrømmes. A. Jo Mere af det Aposterioriske et Begreb indeholder, desto mere Utænkeligt maa det altsaa indeholde; det vil sige: et Begreb er i samme Grad utænkeligt, som det er aposteriorisk; det vil sige: aposterioriske Begreber ere utænkelige“. Forelæsn. over phil. Propæd. fra Universitetsaaret 1861--62. S.~107--108.}% note *)
.

Eenheden af den virkelige Tyngde og Tyngdens Begreb er en Slutningseenhed, en Eenhed af Extremer, der selv ere Slutninger. Men det er, vel at mærke, den ontologiske, ikke den psychologiske, den guddommelige, ikke den menneskelige Subjectivitet, der magter en saadan Slutning. Det subjective Extrem er Magtordets Idealitet; i dette Extrem ere alle Tyngdebestemmelser anticiperede Vidensbestemmelser: Realbegrebet med Led og Forhold bæres \textit{a priori} af imperative Functioner. Det objective Extrem er Magtordets Realitet; i dette Extrem ere alle Tyngdebestemmelser realiserede Existensbestemmelser: Begrebet er empirisk, og det Empiriske en Objectivering ved indicative Functioner.

% --- p. 413 ---
\opage{413}Fra den endelige Videns Standpunkt kan som sagt Extremernes Sammenslutning ikke erkjendes; men just i den Maade, hvorpaa de stilles imod hinanden, viser den menneskelige Fornuft sin naturlige Trang til at hævde dem begge.

Retter Tanken sit Blik mod det subjective Extrem, da lyder Forklaringen saaledes: „Enten vi betragte Tyngden i sin simpleste Yttring som Tiltrækning mellem to i Rummet adskilte Legemer, eller fra et høiere Synspunkt som det Kraftbaand, der sammenholder Verdenssystemerne, saa er det os lige utænkeligt, at Tyngden skulde beroe paa en objectiv, Materien tilfældig iboende Egenskab. Men dersom Tyngden i sin objective Tilværen ikke skal beroe paa en blot factisk, i sin Facticitet tilfældig Egenskab, saa er det jo uimodsigeligt, at en almægtig Fornuft, en guddommelig Viisdom, maa have indplantet Materien denne Egenskab. Da det nu endvidere er aldeles utænkeligt, at en Grundegenskab skulde kunne indplantes i en allerede forud tilværende Materie, saa følger tillige heraf med Nødvendighed, at den samme guddommelige Fornuft, der har bestemt Materiens Grundegenskab, maa have frembragt Materien selv. Materien er altsaa skabt tilligemed Tyngden. Den almindelige Tiltrækning er det phænomenale Udtryk for Delenes Sammentænkning i den guddommelige Viden. Det „almindelige Kraftbaand“ er et almægtigt Tankebaand. Vi kunne ikke begribe det, men vor Fornuft nødes til at antage det“.

Fæster Tanken derimod sit Blik paa det objective Extrem, da gaaer Forklaringen i modsat Retning: „Hvorledes vi end opfatte Tyngdebegrebet i og for sig, saameget er vist, at Tyngden er en Kjendsgjerning. Som Kjendsgjerning er Tyngdens Virkelighed, dens Objectivitet, uomtvistelig. Men naar Tyngden virkelig er objectiv, kan de materielle Deles gjensidige Tiltrækning umulig forklares af, at en oversvævende Fornuft skulde sammentænke Delene, bringe dem sammen ved subjective Tænkningsacter. Det ligger i Materiens Natur, det ligger i Delene selv, at de gjensidig tiltrække hinanden.
% --- p. 414 ---
\opage{414}Som Kjendsgjerning har Tyngden ene og alene sin Grund i Materiens virkelige Tilværelse. I den virkelige Tilværelse er Grunden Eet med den objective Nødvendighed. Mod en objectiv Nødvendighed kan der ikke raisonneres; i den objective Nødvendighed ligger just Grunden til, at Tyngden som en for alle Legemer fælles Egenskab ikke lader sig begrunde: Tyngdebegrebet er og bliver aposteriorisk“.

Denne, som det for en kritisk Philosophie kunde synes, uopløselige Antinomie finder paa en simpel og naturlig Maade sin Løsning i den objective Slutning, en Løsning, der er betinget af, at det psychologiske Synspunkt ombyttes med det ontologiske. For den ontologiske Subjectivitet har ethvert af Extremerne samme oprindelige Gyldighed. Det Aprioriske gaaer ikke i Tiden forud for det Empiriske; de imperative Functioner komme ikke til at virke for sig ved Siden af de indicative; Extremernes Midte tilveiebringes ikke ved endelige Combinationer af subjective og objective Elementer; nei, Midten er selv Extremernes articulerede Identitet: alle Slutningens Led ere simultane. Sammenslutningen er en uendelig Vexelbestemmelse af Extremernes Eiendommelighed, altsaa af det Aprioriske med det Empiriske, af det Subjective med det Objective, af Frihed med Nødvendighed, af Magtordets Idealitet med dets Realitet, en Vexelbestemmelse, der netop grunder sig paa den Fordobling, at begge Slutningens Extremer selv ere Slutninger.

2) \emph{I Magtens Ord er Affiniteten en objectiv Slutning.} I den almindelige Tyngdeslutning er al endelig Succession udelukket: Extremer og Midte fremstaae som med eet Slag. Just fordi Tyngden som Egenskab er en almindelig, Materien som Materie iboende Egenskab, kan her, hvad Tyngdebegrebets Realisering angaaer, ikke være Tanke om endelig Succession. Men hvorledes forholder det sig nu med Affinitetsbegrebet? Staaer det fast, at den ontologisk mægtige Subjectivitet ikke ved nogensomhelst af sine
% --- p. 415 ---
\opage{415}imperative Functioner kan paatvinge den objectivt tilværende Materie en eneste Egenskab, men maa indicativt opfatte og bestemme hver enkelt Egenskab som objectivt Udtryk for Materiens eget objective Væsen: da følger det af sig selv, at de til Affiniteten svarende Egenskaber ligesaalidt kunne være Materien vilkaarlig indpræntede som Tyngden.

Ikke desmindre frembyder Affinitetsbegrebets Realisation eiendommelige, for den objective Slutning charakteristiske Hovedsider. Tyngden er i sit Væsen eensformig; Affiniteten derimod mangeformig. Det er saa langt fra, at den virkelige Forskjel imellem forskjelligartede Legemers Vægt og Vægtfylde skulde betegne en væsentlig Forskjel i deres Tyngde, at den, ifølge Physiken, tvertimod tjener til Beviis for, at alle Legemer ere lige tunge. Anderledes med Affiniteten. Ere Brinten og Iridiet end lige tunge, uagtet deres Vægtfylde er saa forskjellig (Brintens 0,06, Iridiets omtrent 23,0), saa kan man dog ikke sige, at deres Affinitet er den samme, uagtet de indgaae i meget forskjellige Forbindelser. Af den ene og samme almindelige Tyngde kunne de særskilte Vægtforhold, paa givne Betingelser, udledes; men af den almindelige Affinitet kunne de særskilte Affiniteter, hvormegen Vægt man end iøvrigt maatte tillægge Betingelserne, paa ingen Maade afledes.

I Læren om „Functionernes Logik“, særlig om „Functioner af den blandede Grund“,\footnote{Grundideernes Logik. I. S.~316--329.}% note *)
{} er det allerede viist, hvor umuligt det er, under Functionernes concret-physiske Vexelforhold, at aflede de særskilte Affiniteter af en almindelig Affinitets Grundvæsen. Vel kan man opstille en Theorie, som lærer, at Affiniteten i hvert Stof altid ved samme Temperatur maa have samme Intensitet; men hvilken Intensitet Affiniteten i det givne Stof, ved den givne Temperatur, nu virkelig har, er en empirisk Bestemmelse, der ligesaa
% --- p. 416 ---
\opage{416}lidt lader sig aflede af det almindelige Begreb om Affinitet som af det almindelige Begreb om Stof.

Det Empiriske i Affinitetsbegrebet peger imidlertid paa en ganske særegen Maade hen paa det Aprioriske. Ligesom nemlig den ene Tings Egenskab i Almindelighed først kommer tilsyne derved, at den, under Tingenes Vexelvirkning, reflecterer sig i en anden Tings Egenskab, saaledes gjælder det i særegen Forstand, at et Stofs eiendommelige Affinitet kun kan sees, forsaavidt den under Forbindelsen gjør den tilsvarende Affinitet i et andet Stof kjendelig. De forskjellige Forhold, i hvilke Kulstof forener sig med Brint, vise os med det Samme de forskjellige Forhold, i hvilke Brinten forener sig med Kulstoffet: Kulstoffets Affinitet bliver aabenbar ved Brintens, idet Brintens bliver aabenbar ved Kulstoffets. De herunder hørende Kjendsgjerninger udvise en mærkværdig Vexelbestemmelse af Tilfældigt og Nødvendigt, af Mulighed og Anlæg.

Opfattede efter deres Phænomen ere de forskjelligartede chemiske Stoffer, Metaller og Metalloider, jo ganske uafhængige af hinanden. Eller hvorledes skulde det vel kunne paavises, at Tilværelsen af det ene Stof, saasandt det virkelig er et Grundstof, skulde med Nødvendighed forudsætte Tilværelsen af det andet? Ilten kunde jo dog meget vel tænkes at være til uden Brinten, Brinten uden Qvælstoffet o.~s.~v. Forsaavidt nu den for sig existerende Ilt er uafhængig af den for sig existerende Brint, falde Iltens og Brintens Egenskaber jo ganske udenfor hinanden: Iltens Egenskaber vedkomme aldeles ikke Brinten, Brintens ikke Ilten; de to Egenskabssystemer ere altsaa tilfældige for hinanden. Betragte vi dernæst Affiniteten, idet vi søge at trænge ind i Stoffernes „Natur og Væsen“, hvad opdage vi saa? Ja, saa opdage vi, at Stoffernes chemiske Egenskaber, langt fra at være hinanden ligegyldige og tilfældige, tvertimod ere bestemte ved hinanden og for hinanden, og saaledes staae i et indbyrdes nødvendigt Vexelforhold. Hvad Ilten i sin Forbindelse med Brinten
% --- p. 417 ---
\opage{417}betyder, beroer paa Brinten; hvad Brinten betyder, bestemmes igjen ved Ilten: dersom Ilten og Brinten ikke oprindelig vare bestemte for hinanden, vilde Vandet aldrig være blevet til.

Hvad Vexelbestemmelsen af Tilfældigt og Nødvendigt (her tænkes paa den virkelige Nødvendighed) er med Hensyn paa Phænomenet, er igjen Vexelbestemmelsen af \emph{Mulighed} og \emph{Anlæg} med Hensyn paa Væsenet. Henført til sin Mulighed er det særegne Produkt et tilfældigt eller, nøiagtigere, tilfældig-nødvendigt Resultat; henført til sit Anlæg er Produktet et realiseret Formaal. For den objectivt fornuftige Betragtning ville de chemiske Produkter i det Hele tage sig ud, snart som tilfældig-nødvendige Resultater, snart som realiserede Formaal. At Betragteren hverken kan opgive den ene eller den anden af disse Synsmaader, og dog heller ikke paavise deres indbyrdes Grændse, viser igjen, hvorledes den psychologiske Subjectivitet svinger imellem den objective Slutnings Extremer, uden at magte Slutningen selv.

Lade vi den saakaldte sunde Menneskeforstand (eller „den sunde Fornuft“) vælge Udgangspunktet for sine Reflexioner først i det subjective, dernæst i det objective Extrem, da ville de modstridende, under Forhandlingen om det Aprioriske og det Empiriske i Tyngdebegrebet angivne Synsmaader udvikle sig i videregaaende, til Affinitetsbegrebet særlig svarende Conseqvenser. Fæstes altsaa Blikket eensidig paa det subjective Extrem, vil Betragtningen gaae i følgende Retning: „Vistnok kan det for vor indskrænkede Viden i mangfoldige Tilfælde tage sig ud, som om de chemiske Produkter kun vare tilfældig-nødvendige Resultater af stedfindende, ved Stoffernes Beskaffenhed og Omverdenens Betingelser bestemte, Affinitetskræfter; men see vi nærmere til, da kunne vi umulig unddrage os Indtrykket af, at hvad der i Kjendsgjerningen seer ud som Produkt af en udenfra betinget Mulighed, ved nærmere Eftertanke viser sig at være en virkelig Udtalelse af guddommelig Plan og Anlæg. Muligheden af Iltens Forbindelse med Brinten er kun
% --- p. 418 ---
\opage{418}tænkelig under den Forudsætning, at den samme guddommelige Viisdom, der \textit{a priori} har bestemt Iltens, ogsaa har bestemt Brintens chemiske Beskaffenhed. Men ere de nævnte Stoffers chemiske Egenskaber bestemte ved en over begge ophøiet Fornuft, er, med andre Ord, Iltens Eiendommelighed fastsat, ikke blot saaledes, at den virkelig \emph{kan}, hvad der kun antyder \emph{Mulighed}, men ogsaa, hvad der ligefrem vidner om \emph{Anlæg}, saaledes, at den paa givne Betingelser udtrykkelig \emph{skal} forene sig med Brinten: hvo vil da kunne negte, at den samme guddommelige Forstand, der har bestemt Ilten til Forbindelse med Brinten, ved denne Forbindelse tillige forud har bestemt, og altsaa tilsigtet, Vandets Dannelse. Hvor charakteristisk ethvert Anlæg til chemisk Forbindelse maa være, udtaler sig jo aabenbart i den mangfoldigt varierende Modsætning af Syre og Base: disse Forbindelser ere Syrer, efterdi disse ere Baser. Ved den forud fastsatte Adskillelse af Syrer og Baser er ikke blot Muligheden af en stor Mængde Saltdannelser given; men der er ved den i Anlæget forud bestemte Adskillelse af Syrer og Baser udtrykkelig lagt an paa Dannelsen af bestemte Salte.“ Det er Leibnitzes Lære om Muligheder, Anlæg og forud bestemte Harmonier, som ved Reflexioner over den objective Slutnings subjective Extrem atter og atter vil finde Stadfæstelse.

Fæstes Blikket ligesaa eensidig paa det objective Extrem, gaaer Betragtningen naturligviis i modsat Retning: „Saa utænkeligt som det er, at Tyngden ved et guddommeligt Valg, en almægtig Vilkaarlighed, kan være Materien paatvungen, ligesaa utænkeligt er det, at Sligt skulde kunne skee med nogen anden Egenskab. Ikke ifølge guddommelig Overveielse, Beregning og Forudbestemmelse kan der indpodes snart disse, snart hine mere eller mindre hensigtssvarende Egenskaber i Materien; thi Materien med sine virkelige og mulige Egenskaber er med eet Slag, hvad den er og som den er: her gjælder hverken Hensigt eller Valg eller Vilkaarlighed. For%
% --- p. 419 ---
\opage{419}skjellen imellem Mulighed og Anlæg er en Reflexionsforskjel i den betragtende Subjectivitet, og har Intet med det Objective at gjøre. Vistnok spiller Modsætningen af Syrer og Baser en vigtig Rolle i Naturens „Huusholdning“, og hører forsaavidt væsentlig med i Tingenes Orden. Men hvad har vel den objective Tingenes Orden med saadanne Subjectivitetskategorier som Anticipation, Præformation, Anlæg og mechanisk-vilkaarlig Forudbestemmelse at gjøre? Der er i Naturen en Tingenes Orden, ikke fordi der er et særligt ordnende Princip, men fordi en Tingenes Uorden er umulig. Naturen er et uendeligt Hele, der bestaaer i sine Dele. Som Delene ere, maa ogsaa det Hele være, og omvendt: som det Hele er, maae Delene være. Havde vi et andet Naturhele, saa vilde vi ogsaa have andre Dele: Naturens Orden beroer just paa, at Delene ikke kunne Andet end svare til det Hele, og det Hele til Delene.“\footnote{Jvfr. Phil. Propæd. S.~189.}% note *)
{} Det er Empirismen med sine Kjendsgjerninger og sin factiske Nødvendighed, der atter og atter vil søge sit Støttepunkt i den objective Slutnings objective Extrem. Men hverken den abstracte Idealisme med sit Aprioriske eller den abstracte Realisme med sit Empiriske vil være istand til at magte Slutningen selv.

Til Indsigt i Forholdet mellem Magtordet og den objective Slutning er derimod Alt i det Foregaaende forberedt. Med Magtens Ord i dets Idealitet ere alle Affinitetens specifiske Bestemmelser \textit{a priori} satte og forudsatte; med Magtens Ord i dets Realitet ere de samme Bestemmelser empirisk realiserede. Misforstaaelsen af den forud bestemte Harmonie i det subjective Extrem er hævet i Læren om de anticiperende og præformerende Functioner. Misforstaaelsen af den empiriske Bestemthed i det objective Extrem er hævet i Læren om de objective Muligheder og deres Afhængighed af subjectivt bestemmende Functioner. Slutningens Midte er
% --- p. 420 ---
\opage{420}altsaa ligesaa oprindelig som dens Extremer, og Extremerne ligesaa oprindelige som Midten: det Aprioriske og det Empiriske gjennemtrænge hinanden i alle Punkter.

3) \emph{I Magtens Ord er den organiske Proces en objectiv Slutning.} For den raisonnerende Betragter falde den objective Slutnings to Extremer ud fra hinanden som to adskilte Hemisphærer; er den ene over, saa er den anden under Erkjendelsens Horizont.

Fordybet i det subjective Extrem, fordrer den endelige Viden jo vistnok det Objective; men da Objectiviteten ikke faaer Lov at slaae ud i sit Extrem, vil dens væsentlige Gyldighed, dens relative Selvstændighed, aldrig kunne hævdes. Istedenfor at have sine Tingsaarsager, sine Grunde og Muligheder i sig selv, faaer alt Objectivt tvertimod sin sidste Aarsag og dermed sin oprindelige Grund og Mulighed udenfor sig selv, nemlig i Subjectiviteten.

Imponeret af det objective Extrem, endog i den Grad, at den ikke engang kan mærke det som Extrem, er den endelige Viden naturligviis saa optagen af Kjendsgjerninger, Sandsninger og empiriske Forudsætninger, at den ikke har et Blik tilovers for det oprindeligt Subjective. Er den objective Verden med sine Kjendsgjerninger ikke selv et Extrem, men noget i sig Absolut: hvorledes skulde den da kunne taale, endsige fordre, Subjectiviteten som et modsvarende Extrem? Vel maa der en Subjectivitet til, for at betragte det Objective og danne sig et Begreb om det Virkelige; men i Virkeligheden er det Objective selv ligesaa ligegyldigt mod Betragteren som mod Betragtningen og det ved Betragtningen vundne Begreb.

Fra intet af de extreme Synspunkter er det muligt at løse Objectiveringens Problem. I det subjective Extrem anskues den objectiverende Subjectivitet; men Objectiveringen mangler sit Objective. I det objective Extrem anskues Objectiviteten, det i Virkeligheden Objectiverede, men den objectiverende Subjectivitet er kun en Skygge. Først i den objective
% --- p. 421 ---
\opage{421}Slutnings alsidige Udvikling kan Objectiveringsproblemet finde sin Løsning.

\emph{Tyngden} og \emph{Affiniteten} have, som ovenfor viist, det Charakteristiske ved sig, at de, forsaavidt de falde under den endelige Videns Synskreds, hver paa sin Maade kategorisk holde Slutningens Extremer saaledes adskilte, at den empiriske Forsker udelukkende maa rette sit Blik mod det Objective, og lade den speculative Tænker famle efter det Subjective. Med den organiske Proces er Phænomenet blevet et andet. Kjendsgjerningen giver nu selv en ligelig Speiling af begge Extremer: den organiske Proces er i sit Phænomen en objectiv Slutning. Som objectiv Slutning bliver den organiske Proces dog aabenbart mistydet saavel af den empiriske Physiologie som af den immanente Speculation.

Som Fagvidenskab holder Physiologien sig til det objective Extrem, naturligviis ikke som Extrem, men som slet og ret Objectivitet; den sammenligner Organismen og dens enkelte Dele med en Maskine: „Hjertet ligner et med passende Ventiler forsynet Pumpeværk, Arterierne elastiske Rørledninger“ o.~s.~v. Den organiske Proces forklares af mechaniske og chemiske Functioner: Betragtningen er aldeles objectiv. Bestemmes nu Organernes Virksomhed udelukkende ved deres „anatomiske, physiske, chemiske og deraf afledede eller medvirkende Egenskaber, % PRINTER: opening „ before "anatomiske" has no closing mark
da er det jo en Selvfølge, at „den specielle Anatomie“, Læren om Vævenes Bygning, „Physiken“ og „Chemien“ maae afgive Grundlaget for Physiologien. Men nu kan, hvad allerede er viist i „Functionernes Logik“, en organisk Function paa ingen Maade afledes eller ligefrem forklares af blot mechaniske og chemiske Virksomheder. Physiologiens Forklaringer løbe da ogsaa tilsidst ud paa, at Videnskaben i Grunden Intet kan forklare. „Vi vide“, siger Physiologen, „ikke hvad Livet er“. Men dersom det virkelig forholder sig saa, at Physiologien ikke kan lede til Indsigt i, hvad Livet er, saa kan den heller ikke lede til Indsigt i, hvad en Orga%
% --- p. 422 ---
\opage{422}nisme er, og da heller ikke til Indsigt i, hvad en organisk Proces er.\footnote{„Livet“, siger Physiologen, „maa forholde sig som Organisationen. En enkelt Celles Liv er \emph{Indbegrebet af denne Celles Virksomhed} og er simpelt i samme Forhold som Organisationen er simpel. Livet hos et Dyr med en høit udviklet Organisation er Indbegrebet af alle de Kræfter, der røre sig i den, og er compliceret i samme Forhold som Organismen er compliceret“. Det Eensidige viser sig deri, at Physiologen, der paa denne Maade gjør Livet til „\emph{en Følge af Organismen}“, glemmer, at den paaberaabte Idealitet er begrundet i en væsentlig Forskjel, at Livet kun er en Følge af Organismen, forsaavidt Organismen selv er en \emph{Følge af Livet}“. % PRINTER: closing “ after "Livet" has no opening mark
Forelæsn. ov. „Phil. Propæd.“ fra Universitetsaar 1860--61. S.~223.}% note *)

Den organiske Proces bliver heller ikke klaret i den immanente Speculation. Vel har Speculationen Ret, naar den imod Erfaringslæren gjør gjældende, at Begrebet som Begreb netop realiseres i Organismen; men den har Uret i at lade, som om den virkelig kunde bevise, at det sig i Organismen realiserende Begreb, var det samme som det paa Mechanismens og Chemismens Trin med Betingelserne kæmpende, gjennem lavere Forudsætninger sig oparbeidende Begreb. Ved denne, i det Foregaaende fra flere Sider paaviste Mistydning, kommer „Begrebets“ speculative Philosophie i en ligesaa skjæv Stilling til den organiske Proces som Physiologien.

Physiologiens Misviisning er, hvad Begrebet angaaer, den, at den organiske Proces blot skulde være Resultat af mechaniske og chemiske Functioners Samvirken. Speculationens Misviisning er, hvad Kjendsgjerningen angaaer, den, at den organiske Proces skulde med uafviselig Conseqvens lade sig i Logiken udvikle og bestemme som en naturnødvendig Yttring af det med sit Indhold identiske Begreb.\footnote{Daß das lebendige Individuum sich in sich selbst gestaltet, damit spannt es sich gegen sein ursprüngliches Voraussetzen, und stellt sich % note continues on p. 423
als an und für sich seyendes Subjekt der vorausgesetzten objektiven Welt gegenüber. Das Subjekt ist der Selbstzweck, der Begriff, welcher an der ihm unterworfenen Objektivität sein Mittel und subjektive Realität hat; hierdurch ist es als die an und für sich seyende Idee und als das wesentliche Selbstständige konstituirt, gegen welches die vorausgesetzte äußerliche Welt nur den Werth eines Negativen und Unselbständigen hat. Sein Trieb ist das Bedürfniß, dieß Andersseyn aufzuheben, und sich die Wahrheit jener Gewißheit zu geben“. % PRINTER: closing “ has no opening mark (none before "Daß")
Hegel. Die subjective Logik. S.~255--56.}% note **) -- PRINTER: anchor in text printed "*)", note printed "**)"

% --- p. 423 ---
\opage{423}Vistnok bliver den organiske Proces i Hegels Logik opfattet og bestemt som objectiv Slutning. Men Opfattelsen er saa eensidig, at hvad her sættes som System af objective Slutninger, Systemet nemlig af de i Organismen organiske Processer, nærmere beseet, mangler Fuldstændighed og falder eensidigt ind under den objective Slutnings objective Extrem. Hvorledes det hermed forholder sig, kunne vi let oplyse lidt nærmere.

Allerede i Functionernes Logik\footnote{Grundideernes Logik I. S.~434--446.}% note *)
{} er Organismen som „Objectiveringens tredie Grundtotalitet“ bestemt ved en objectiv Slutning: S\,--\,A\textsubscript{o}\,--\,O\textsubscript{o} er Slutningens Schema. Betragte vi dette Schema noget nøiere, da er det klart, at den organiske Proces foregaaer i O\textsubscript{o} og kun henføres til A\textsubscript{o}, forsaavidt O\textsubscript{o} selv fremgaaer af A\textsubscript{o}. Da nu O\textsubscript{o} i Systemet af alle sine Slutninger udtrykkelig betegner den objective Slutnings objective Extrem, saa følger heraf, at den organiske Proces vel er en objectiv Slutning, men en Partialslutning, der opgaaer som Moment i det objective Extrems Totalslutning.

Denne Belysning har nu en paa Tyngde og Affinitetsslutning tilbagevirkende Kraft. Ligesom nemlig Tyngdeslutningen i sine særskilte Phænomener gaaer ind under Systemet af de abstract-mechaniske Functioner, altsaa ind under „Objectiveringens første Grundtotalitet“, saaledes gaaer Affinitets%
% --- p. 424 ---
\opage{424}slutningen, specificeret i sine chemiske Fremtrædelser, ind under Systemet af de mechanisk-dynamiske Functioner, altsaa ind under „Objectiveringens anden Grundtotalitet“.

Og hvad lære vi saa heraf? Vi lære heraf, at den objective Slutning ligesaa vel er at udvikle i Systemer af Total- og Partialslutninger som den subjectiv-immanente, at Functionerne af Magt forsaavidt ere ligesaa dialektiske som Functionerne af Viden. Magtfunctionernes Dialektik kan da udvikles ved en nærmere Bestemmelse af

\textit{γγ)} \emph{Den objective Slutning og det kategoriske Magtsprog.} --- Den objective Slutning er en teleologisk Act, en virkelig Udtalelse af den begribende, sit Begreb i Gjenstanden fremstillende Subjectivitet, en Udtalelse, som, naar den ved interpreterende Functioner er blevet bestemt, med Rette kan siges at give Gjenstanden sin egentlige Betydning, og forsaavidt godtgjøre, at der virkelig er Mening i Tingen, efterdi der oprindelig er meent Noget med Tingen.

Den Vildfarelse ligger imidlertid nær, at det subjective Hensyn nu med Eet skulde faae en ulogisk, vilkaarlig Overvægt over det objective. Man kunde jo raisonnere saaledes: „Er enhver Ting et realiseret Begreb, ethvert realiseret Begreb bestemt ved en objectiv Slutning, og enhver objectiv Slutning en Udtalelse af den guddommelige Fornufts Øiemed, dens Mening om Tingen og med Tingen: hvilken Vrimmel af Øiemed, subjective Antydninger og guddommelige Meninger maa der da ikke findes i Tingenes Verden! Skal det nu være den menneskelige Videns høieste Maal at efterspore, opdage og bestemme disse Meninger, saa er der vist ingen Tvivl om, at vi snart maae komme tilbage til en Teleologie, en i Sandhed theologisk Teleologie, der er ligesaa dybsindig og i sin Art ligesaa meningsløs, som den, der satte en Athanasius Kircher istand til --- i „Lithotheologiens“, „Petinotheologiens“, „Græshoppetheologiens“ Tidsalder --- at opdage 6561 Beviser for Guds Tilværelse.“

% --- p. 425 ---
\opage{425}At dette Raisonnement beroer paa en Misforstaaelse, er let at eftervise. For det Første er der jo paa et tidligere Stadium gjort en Adskillelse, en væsentlig, kategorisk bestemt Adskillelse imellem den guddommelige Mening med Tingene og Menneskenes Meninger om Tingene. For det Andet gjælder det paa et høiere Trin om det Teleologiske, hvad der paa et lavere gjælder om det Væsentlige og Nødvendige. At Tingenes Væsen gjennemtrænger Tingene, at Alt, hvad der skeer i Tingenes Verden, skeer med en vis Nødvendighed, kan jo hverken udelukke Bestemmelsen af noget Uvæsentligt eller ophæve Betydningen af noget Tilfældigt. Paa tilsvarende Maade forholder det sig med det Teleologiske. Fordi Alt i Tingenes Verden underlægges Teleologien, efterdi Alt i Tingenes Verden underlægges den objective Slutning, deraf følger ikke, at hver enkelt Ting i al Umiddelbarhed skulde som særlig opnaaet Øiemed have særlig teleologisk Betydning.

Hvad altsaa Forskjellen mellem Nødvendighed og Tilfældighed er for Reflexionen, det er Modsætningen mellem Formaal og Midler for den objective Slutning. I Systemer af objective Slutninger har den altformaaende Subjectivitet udtalt sin Videns Indhold i Magtens Ord: et rigt Indhold, et mangfoldigt Ord! Men til Forstaaelse af dette Magtens Ord er det fremfor Alt nødvendigt at gjøre Forskjel mellem \emph{Hovedtanker}, \emph{Mellemtanker} og \emph{Bitanker}.

I enhver af „Objectiveringens Grundtotaliteter“ udtaler sig en guddommelig \emph{Hovedtanke}; en saadan Hovedtanke er, med Hensyn paa Formen, objectivt udtrykt i Totalslutningens objective Extrem. Hver Hovedtanke er, hvad det Enkelte og Særskilte angaaer, udviklet i et Indbegreb af charakteristiske \emph{Mellemtanker}: den bestemte Totalslutning bestaaer i et System af eiendommelige, for Totaliteten nødvendige Particularslutninger. Endelig, idet Particularslutningerne indgaae i Endeligheden, optage det Tilfældige og ændres i det Udvortes, kunne disse
% --- p. 426 ---
\opage{426}paa een Gang væsentlige og dog uvæsentlige Ændringer vel siges at udtrykke Systemets \emph{Bitanker}.

„Men“, vil man maaskee spørge, „er denne Forklaring vel meer end en sindrig Anvendelse af interpreterende Functioner; vil Læren om det kategoriske Magtsprog virkelig kunne udøve nogen befrugtende Indflydelse paa Videnskaben?“ Vi svare: Ja! Læren om det kategoriske Magtsprog er af afgjørende Betydning ved Bestemmelsen, saavel af Philosophiens Maal, som af dens Forhold til Fagvidenskaberne. Vi skulle strax prøve dens Betydning paa Objectiveringens tre Grundtotaliteter, den mechaniske, den mechanisk-dynamiske og den organiske.

1) \emph{Den mechaniske Totalitet i kategorisk Mening}. I en Totalitet som vort Planetsystem er et i sig fuldstændigt Begreb blevet realiseret. Heri ligger, at dette Begreb maa tage sig ud som et System af Kjendsgjerninger, og at Kjendsgjerningernes System omvendt maa lade sig forklare i et Begreb. Fra Physikens Standpunkt maae Kjendsgjerningerne sees som det Oprindelige, Begrebet som det Afledede. Functionerne, som her henføres til den psychologiske Subjectivitet, ere indicative og explicative, men ikke interpreterende: Physikens Opfattelse af Magtbegrebet er og maa være uteleologisk. Fra Naturphilosophiens Synspunkt maa Begrebet derimod som Begreb være det i guddommelig Forstand Oprindelige, Kjendsgjerningernes System det Afledede; Functionerne, som her henføres til den ontologisk mægtige Subjectivitet, ere ikke blot indicative, men ogsaa imperative, ikke blot explicative, men ogsaa interpreterende: Naturphilosophiens Opfattelse af Magtbegrebet er og maa være teleologisk.

Striden imellem den mathematiske Physiks mechaniske og den immanente Speculations teleologiske Opfattelse af Planetsystemet er paa en ret uforbeholden og slaaende Maade kommet til Udtalelse i Hegels Naturphilosophie.

Hegels speculative Indsigelser imod Newtons \textit{principia
% --- p. 427 ---
\opage{427}philosophiæ naturalis mathematica}\footnote{Vorlesungen über die Naturphilosophie. Werke 7ter Band. 1. Abtheil. Berlin 1842. S.~98. Jvfr. Forelæsn. ov. „Phil. Propæd. Universitetsaar 1861--62.“ S.~57--66.}% note *)
{} lyde paa: \textit{a}) at Newtons mathematiske Fremstilling af de keplerske Love er en Forvanskning af Keplers paa Erfaring byggede Theorie, \textit{b}) at Newtons mathematiske Beviser i Grunden Intet bevise, og endelig \textit{c}) at Mathematikernes af Newton indførte Lære om Kræfternes Parallelogram, og hvad dermed hænger sammen, er uden Anvendelse paa Himmellegemernes Bevægelse og indbyrdes Forhold. Og hvad vil saa Hegel sætte istedenfor de mechaniske Kræfters mathematiske Analyse? Slutninger, et System af objective, teleologiske Slutninger!

De colossale Misgreb, hvori Hegel ved denne Leilighed har gjort sig skyldig, lade sig alene forklare deraf, at Forholdet imellem den objective Slutning og det kategoriske Magtsprog har været ham aldeles uklart.

Hegel har et skarpt Øie for den tresidige Form af den objective Totalslutning, men naar han har faaet Totalslutningens trefoldige Schema opstillet og i Anledning af hver Figur fremført nogle almindelige Bemærkninger, saa er han færdig. Paa Physikens Iagttagelser og Beregninger kan han aldeles ikke indlade sig; de mathematiske Begrebsbestemmelser ere i den Grad uforenelige med hans absolute Begreb, at han ligefrem bekæmper dem og det endog, hvad der ingenlunde har været til Philosophiens Fordeel, uden at have forstaaet dem.

Den mathematiske Physik er fra sin Side conseqvent; den betragter hele Planetsystemet, ikke som et System af objective Slutninger, men som et af givne Forudsætninger og Betingelser fremgaaet Resultat; den er da i god Overeensstemmelse med sig selv, naar den opfatter Systemets enkelte
% --- p. 428 ---
\opage{428}Dele fra samme Synspunkt. Hegel derimod er i høieste Grad inconseqvent. Han seer Systemet i dets Totalitet under Synspunktet af objective Slutninger; men, mærkværdigt nok, han seer ikke, at hvad der gjælder om det Totale ogsaa maa gjælde om det Particulære; han seer ikke, at Totalslutningen kun kan bestaae i et System af Particularslutninger.

Udtale Planeternes, Biplaneternes og Kometernes Bevægelse om Solen som Centrallegeme en i den objective Slutnings objective Extrem realiseret Hovedtanke, da er der heller Intet iveien for, at Jordens Bevægelse om Solen jo kan tydes som en til Hovedtanken svarende Mellemtanke, og Maanens Bevægelse om Jorden som en til Mellemtanken svarende Bitanke. At Kategorierne Hovedtanke, Mellemtanke og Bitanke imidlertid maae være uendelig relative, sees umiskjendelig deraf, at enhver af Mellemtankerne igjen kan sættes som en Hovedtanke med tilsvarende Mellemtanker og Bitanker.

Denne Synsmaade afgiver en for Realbegrebets philosophiske Udvikling meget vigtig Conseqvens. Conseqvensen er, at Naturphilosophiens Analyser af Particularslutningerne ligesaavel kunne gaae i det Enkelte som den mechaniske Physiks Analyser af de hypothetisk-nødvendige Resultater. Ja, hvad Mere er, ved denne og kun ved denne Synsmaade bliver Naturphilosophien istand til at tilegne sig de mathematisk-physiske Analysers for Idee-Erkjendelsen væsentlige og uundværlige Udbytte.

Den dialektiske Vexelbestemmelse af Hovedtanker, Mellemtanker og Bitanker bliver nemlig ikke at fastsætte efter et speculativt Skjøn; et saadant Skjøn er for ubestemt og subjectivt til at træffe Afgjørelser i Videnskaben. Det Dialektiske maa her nødvendigviis have det Exacte til Grundlag. Den naturphilosophiske Opgave vil altsaa bestaae i at vise, hvorledes den mathematiske Physiks hypothetisk-nødvendige Resultater i det Enkelte lade sig forvandle til betegnende Particularslut%
% --- p. 429 ---
\opage{429}ninger, idet disse Particularslutninger igjen vise sig at være Momenter i en tilsvarende Totalslutning.\footnote{Hvad hermed menes, kan foreløbig oplyses ved et Par Exempler. „Dersom et materielt Punkt paavirkes af en accelererende, mod eet og samme faste Midtpunkt rettet Kraft, og meddeles en Begyndelseshastighed, da ville de af Radius vector i lige Tider beskrevne Arealer være ligestore.“ Gjøres denne Dom til Oversætning i en formel subjectiv Slutning med kategorisk Undersætning, erholde vi et \emph{hypothetisk-nødvendigt Resultat}. Naar dette Resultat derimod i sin Sammenhæng med den i Totaliteten udtalte Hovedtanke viser sig som en realiseret Mellemtanke, forvandles det til en objectiv Particularslutning, antager et teleologisk Præg, og afgiver saaledes en til det kategoriske Magtsprog hørende Sætning. Den samme Forvandling lader sig naturligviis foretage med alle lignende Resultater. „Dersom Tiltrækningen skeer efter Newtons Lov, maa Banen blive et Keglesnit;“ „dersom Begyndelseshastigheden er lille i Forhold til Attractionskraften, maa Keglesnittet blive en Ellipse“. „Dersom forskjellige materielle Punkter bevæge sig i Ellipser om samme Centrum, vil, idet $\frac{\mathrm{T}^2}{a^3}$ er constant, Omløbstidernes Qvadrater forholde sig som Cuberne af de halve Storaxer.“ Idet disse hypothetisk-nødvendige Resultater forvandles til objective Particularslutninger, sees ethvert af dem at være et særligt, under Totalanlæget A\textsubscript{z} % UNSURE: subscript z or x (small italic, reads as z)
hørende Anlæg, et antydet Øiemed, en egen Sætning med sin ved Magtsproget bestemte Mening.}% note *)

Physikens logiske Functioner ere væsentlig explicative, Naturphilosophiens interpreterende: det kategoriske Magtsprog, der har forvandlet Begreberne til Kjendsgjerninger og netop derved anviist Physiken sin Opgave, har med det Samme lagt en guddommelig Mening i Kjendsgjerningerne, og derved sat en Opgave for Philosophien.

2) \emph{Objectiveringens anden Grundtotalitet er bestemt ved det kategoriske Magtsprog}. „Idet Geologien forudsætter Geognosien, maa den nødvendigviis tage sin Begyndelse i det Endelige, i Virkelighedens Stykværk. At
% --- p. 430 ---
\opage{430}samle og undersøge Mineralier, at charakterisere Bjergarter, at paavise Forholdet mellem Slenter, gjennembrudte Masser, neptunske, plutoniske og vulkanske Dannelser o.~s.~v. hører i den Grad under Stykværkets Kategorie, at man ved foreløbig Betragtning neppe skulde antage, at saa uensartede Bestanddele kunde sammenarbeides til et videnskabeligt Hele, endsige da deres Bestanddele føre til en Erkjendelse af Ideen. Ikke desto mindre er Geognosien blevet en Videnskab; jo mere denne Videnskab magter Kjendsgjerningerne, jo grundigere den bearbeider Endelighedselementerne, desto mere viser den ud over sig selv, thi dens Indhold er af Ideen \dots\ Seet fra Geognosiens Standpunkt er Jorden kun et sammenstykket Hele; fra Geologiens Standpunkt er Jordlegemet derimod at opfatte som en elementær Organisme, et væsentligt Hele, hvis Dele efterhaanden have udviklet sig af et indre Anlæg. I Anlægets Oprindelighed have vi da igjen det Principielle, i Delenes Udsondring og Heelhedens planmæssige Articulation det Typisk-Teleologiske, d.~e. vi have, dersom Erkjendelsen af den uendelige Formidling blot var fuldkommen, den sig i Jorddannelsen realiserende Idee“.\footnote{Phil. Propæd. S.~186.}% note *)

Vexelbestemmelsen af Geognosie og Geologie viser, naar Kategorierne tages i ovenangivne Mening, hvorledes Naturvidenskabens alsidige Bearbeidelse af en elementær Organisme med indre Nødvendighed maa føre fra den abstract objective, ikke teleologiske, Synsmaade over i Teleologien. Det er nemlig ikke Geologen mulig at bedømme Forholdet mellem Tilstandene i de forskjellige Jordperioder og charakterisere disse som en Række af lavere og høiere Udviklingstrin, medmindre han forudsætter et Slags Endemaal, en Modenhedstilstand, hvortil den hele Udvikling stræber hen som til en Afslutning. Men Grundforholdet imellem det physisk Objective og det Teleologiske vil Geologen aldrig være istand til at
% --- p. 431 ---
\opage{431}gjennemskue: det, der blender ham, er atter det kategoriske Magtsprog.

I Kraft af det kategoriske Magtsprog er Geologiens Begreb realiseret i virkelige Kjendsgjerninger, i Kjendsgjerningernes System sees overalt Tingsaarsager; men den med Hensigtsaarsagen identiske Grundaarsag bliver borte. Hvor huul og phantastisk Naturphilosophien maa blive, naar den for at bøde paa Empiriens Eensidighed springer af fra Tingsaarsagerne (\textit{causæ efficientes}) og udelukkende stirrer paa Grundaarsagen (\textit{causa finalis}), kan sees af Hegels Forklaring. „Wie die Luft sich als allgemeines Element ihre Spannung aus der Erde nimmt, so das Meer seine Neutralität. Die Erde dünstet gegen die Luft aus, als Meer; gegen das Meer aber ist die Erde der Krystall, der das überflüßige Wasser aus sich abscheidet, in Quellen, die sich zu Flüssen sammeln. Aber dieß ist, als süßes Wasser, nur die abstracte Neutralität, das Meer dagegen die physische, in die der Krystall der Erde übergeht. Der Ursprung der uversiegbaren % PRINTER: printed "uversiegbaren", sense "unversiegbaren"
Quellen darf also nicht auf mechanische und ganz oberflächliche Weise als ein Durchsikern % PRINTER: printed "Durchsikern", sense "Durchsickern"
dargestellt werden, so wenig als nach der andern Seite das Entstehen der Vulkane und heißen Quellen; sondern wie die Quellen die Lungen und Absonderungsgefäße für die Ausdünstung der Erde sind, so sind die Vulkane ihre \emph{Leber}, indem sie dieß Sich an ihnen selbst Erhitzen darstellen.“ Det, der blender den speculative Naturphilosoph, er atter det kategoriske Magtsprog. Han seer i Kjendsgjerningernes Totalitet ikke blot et Begreb, men i Begrebet en Mening. Meningen søger han ganske rigtig i Øiemedet; saaledes faae da i hans Øine Jordens Kilder og Vulkaner deres Tilværelse fra deres Betydning og deres Betydning fra den Tjeneste, de yde Jordlegemet: de forholde sig til Jordlegemet som Organer til Organisme. Men ved Bestemmelsen af det Enkelte glemmer den speculative Philosoph, at det Ord, hvori Begrebets Mening objectivt er udtalt,
% --- p. 432 ---
\opage{432}er et Magtens Ord, og at Magtens Ord forkyndes i virkelige Kjendsgjerninger, Vexelvirkninger af virkende Aarsager.

Det er „Naturens Huusholdning“, paa hvis Udtydning saavel Naturforskeren som den speculative Naturphilosoph, hver fra sin Side, lægge an. Men fra begge Sider vil Udtydningen mislykkes, saalænge Forholdet mellem hypothetisk-nødvendige Resultater og objective Particularslutninger ikke er klaret.\footnote{Hvorledes Geognosiens hypothetisk-nødvendige Resultater lade sig i Erkjendelsen omdanne til Particularslutninger, er ovenfor viist ved Exempler af Chemien.}% note *)
{} Først i Vexelvirkningen af Particular- og Totalslutninger bliver Magtsproget forstaaeligt, og Udtrykket „Naturens Huusholdning“ at opfatte, ikke som en figurlig Talemaade, men som et naturvidenskabeligt Begreb.

3) \emph{Objectiveringens tredie Grundtotalitet er bestemt ved det kategoriske Magtsprog}. At den levende Organisme maa opfattes saavel under Synspunktet af Hensigtsaarsager som af Tingsaarsager, er almindelig anerkjendt. Leibnitzes Ord: \textit{finalem causam non tantum prodesse ad virtutem et pietatem, in Ethica et Theologia naturali, sed etiam in ipsa Physica ad inveniendum et detegendum abditas veritates}, skulle vistnok, om ikke de egentlige Physikere, saa dog idetmindste Botanikere og Zoologer villig godkjende. I Udtydning af et Hvirveldyrs Beenbygning, af dets Tænder, Kløer o.~s.~v. maa Zoologen nødvendig støtte sig til teleologiske Forudsætninger. Ved denne Udtydning synes altsaa Fagvidenskaben og Philosophien at nærme sig hinanden, hver fra sit Udgangspunkt. At Tilnærmelsen ikke kan føre til Sammensmeltning, er imidlertid afgjort derved, at Fagvidenskaben, som igjennem det Empiriske opfatter det Aprioriske, aldrig kan igjennem Tingsaarsagerne naae Grundaarsagen, medens Philosophien, der ifølge sin Methode maa begynde med den ved Kjendsgjerninger bestemte Grundaarsag, aldrig
% --- p. 433 ---
\opage{433}ved nogensomhelst Begrebsudvikling vil være istand til at gjøre Forholdet imellem Grundaarsagen og de virkelige Tingsaarsager \textit{in concreto} gjennemsigtigt.

Hvad er da i denne Sammenhæng Philosophiens Betydning? At den med klar Erkjendelse af den menneskelige Videns Grændser tillige kjender og erkjender en guddommelig Viden, en Viden, der, i sin Eenhed med Magten, udtaler sit Indhold i Magtens Ord, i et kategorisk Magtsprog. At Philosophien nu virkelig kjender og erkjender en guddommelig Viden, kan naturligviis ikke sees deraf, at den umiddelbart forudsætter og ved supponerende Functioner beraaber sig paa Eenheden af Guds Alvidenhed og Almagt; men det sees deraf, at Philosophien, efter mange famlende Forsøg, omsider er blevet sig sin Methode bevidst, idet den er blevet sig bevidst, hvorledes den, uden at miskjende den menneskelige Videns Grændse, vil være istand til at udvikle og tydeliggjøre Begrebet af en guddommelig Viden\footnote{Jvfr. Grundideernes Logik I. S.~230 flgd.}% note *)
.

Organismen, den levende Organisme, er af alle Existensformer den, hvori Realiteten af det kategoriske Magtsprog allertydeligst er udtalt. Kun forsaavidt er den menneskelige Viden istand til at forstaae dette Sprog, som Philosophien under fremadskridende Begrebsudvikling er istand til stedse dybere, grundigere, alsidigere at gjøre Rede for Viden og Magt i deres gjennemgaaende Modsætning og dog oprindelige Eenhed. Med den almindelige Eenhed til Grundlag er Modsætningen af Viden og Magt udviklet i Læren om \emph{Dom og Gjenstand}; med den udviklede Modsætning til Grundlag er Modsætningseenheden i Magtens Form udviklet og ved det kategoriske Magtsprog articuleret i \emph{den objective Slutning}. Men hermed er da den objective Slutning tillige udviklet til en Sammenslutning med den subjective.

% --- p. 434 ---
\medskip
\noindent\opage{434}\hspace*{1em}\textit{γ)} \emph{Ideeslutning: Modalfunctioner af Sandhed.}
\medskip

At Udtrykket: Functioner af Sandhed, i nærværende Sammenhæng ikke særlig kan bestemmes ved den contradictoriske Modsætning til Functioner af Usandhed, men ved den ontologiske Modsætning til Functioner af Viden og Functioner af Magt, forstaaes vel af sig selv. Som Form for den absolute Subjectivitet med dens Indhold og frie Nødvendighed er den subjective Immanensslutning en absolut Form for det Absolute selv, en adæqvat Ideeform, og forsaavidt en Ideeslutning. Videns Indhold er fra dette Synspunkt identisk med Magtordets Idealitet, og Ideeslutningens subjective Functioner Modalfunctioner af Viden. Som Form for Magtordets Realitet i virkelige Kjendsgjerninger er den objective Slutning en adæqvat Form for Udviklingen af Existensens Absolute, og fra sin Side ligeledes en Ideeslutning. Slutningens Indhold er Videns Gjenstand; Videns Gjenstand objectiveres i et kategorisk Magtsprog; Slutningens Functioner ere med Hensyn hertil bestemte som Modalfunctioner af Magt.

Under Udviklingen saavel af den subjective som af den objective Slutning viste sig imidlertid den mærkelige Fordobling, at enhver af Extremerne selv blev en Slutning. Denne Fordobling gjentager sig, og den ved Fordoblingen betingede Integreren af Eensidigheder i enhver Retning er nu et Særkjende for Modalfunctioner af Sandhed. Saaledes bliver i denne Sammenhæng Subjectivitetens alsidig udviklede Immanensslutning udtrykkelig sat som Ideeslutningens subjective Extrem, idet Objectivitetsslutningen med Realbegreber, imperative, executive og interpreterende Functioner o.~s.~v. fra sin Side afgiver Ideeslutningens objective Extrem. Da begge Slutningsextremer ere Ideeslutninger, maa den concrete Idee nødvendigviis selv afgive den Midte, hvori de sammensluttes. Først ved denne Sammenslutning er Begrebet \emph{Ideeslutning} alsidig opfattet og bestemt.

% --- p. 435 ---
\opage{435}Men Begrebet er endnu kun bestemt i Almindelighed; det i sin Almindelighed fastsatte Begreb fordrer en bestemt, gjennem det Særskilte og Enkelte gaaende Udvikling. Under denne Udvikling vil da Fordoblingen bestaae deri: at den hele subjectiv-objective Ideeslutning sættes i sit subjective Extrem og characteriseres som \emph{Idealitetsslutning}; at dernæst den hele, objectiv-subjective Ideeslutning sættes i sit objective Extrem og characteriseres som \emph{Realitetsslutning}; at endelig Idealitetsslutningens subjective og Realitetsslutningens objective Extrem sammensluttes i en Alt forklarende Midte, der saaledes danner en \emph{Afslutning}.

\textit{αα)} \emph{Idealitetsslutning: Ideernes Præexistens.} Opkastes i Logiken det Spørgsmaal, hvilken af de to Kategorier: \emph{Væren} eller \emph{Existeren}, der i strengere Forstand bliver at sætte som den for Ideen oprindelige, kan Besvarelsen vel fordre en nærmere Overveielse. Forsaavidt Ideen som Idee er det Almene, synes Væren unegtelig at maatte foretrækkes for Existeren; thi Væren udtrykker det Almindelige. Man maatte da sige: Ideen er, men den \emph{existerer} ikke; thi det Almindelige er, men det existerer ikke. Men nu er Ideen som Idee jo ikke blot det absolut Almene; den er, efter hvad vi i det Foregaaende fra saamange Sider have oplyst, ligesaa vel det absolut Særskilte og det absolut Enkelte. Forsaavidt Ideen da er bestemt som det absolut Enkelte, er der al Grund til at lægge Eftertryk paa Ideens Existens: Ideen er, fordi den \emph{existerer}.

Og dog er, umiddelbart seet, Kategorien Existens ligesaa uadæqvat for den absolute Idee som Kategorien Væren. Væren er tvetydig, fordi den ligesaa vel lyder paa det Abstracte som paa det Concrete, ligesaavel paa det Uselvstændige som paa det Selvstændige, ligesaavel paa det Uvæsentlige som paa det Væsentlige: Væren er ligegyldig mod Forskjellen mellem Abstract og Concret, Væsentligt og Uvæsentligt, Selvstændigt og Uselvstændigt. Denne Ligegyldighed staaer i af%
% SEAM: word continues on p. 436 as "gjørende"
% --- p. 436 ---
\opage{436}gjørende Strid med Ideens Fordring; thi Ideen er og vil være, ikke det Abstracte, men det absolut Concrete, ikke det Uvæsentlige, men det absolut Væsentlige, ikke det Uselvstændige, men det absolut Selvstændige: Væren er uadæqvat for Ideen. Hvad Væren mangler, besidder Existensen. Ved Kategorien Existens sigtes der udtrykkelig paa en Bestaaen for sig, en særlig, selvstændig Tilværen: det Existerende er ikke blot et Tilværende, men et concret, for sig bestaaende Enkeltvæsen. Og dog har Existenskategorien, umiddelbart forstaaet, Sidebestemmelser, der ikke lade sig forene med Ideen som Idee. En Ting existerer, just fordi det er en Ting: Tingen og det Existerende ere identiske Udtryk. At gjøre Ideen til et Existerende er, med andre Ord, at gjøre den til en Ting; men Ideen er ligesaalidt Ting som Tingen Idee.

Saalænge Væren og Existeren falde paa endelig Maade ud fra hinanden, vil ingen af disse Kategorier tilfredsstille Ideens Krav paa alsidig Bestemthed; for at blive adæqvate, maae de blive identiske; for at blive identiske, maae de afgive Sidebestemmelser for en høiere Kategorie. Den høiere, Væren som Existeren og Existeren som Væren indesluttende, Kategorie kan betegnes ved en \emph{absolut reflecteret Bestaaen}.

I den absolut reflecterede Bestaaen er Væren med; thi det Bestaaende er et Værende. I den absolut reflecterede Bestaaen er Existensen ligeledes med; thi det Bestaaende er et Selvstændigt, og det Selvstændige et i og for sig Existerende. Den forhen uadæqvate Væren er nu ved at være optagen i en absolut Bestaaen, blevet adæqvat, thi den er kommen til at lyde paa afgjørende Selvstændighed; den forhen uadæqvate Existens er ligeledes, efter at have opgivet sin Umiddelbarhed og være forvandlet til absolut reflecteret Bestaaen, blevet adæqvat, thi Ideen har sin Bestaaen, sin absolute Selvstændighed i absolut Reflecteren.

Men hvorledes bestemme vi nu den absolute Reflecteren? Ifølge vort Foregaaende som en uendelig Energie! Ideen \emph{er},
% --- p. 437 ---
\opage{437}det vil sige: den \emph{bestaaer} i uendelig Energie; Ideen \emph{existerer}, det vil sige: dens uendelige Energie er Eet med \emph{dens absolute Selvstændighed}. Den uendelige Energie er, subjectivt seet, en uendelig \emph{Begriben}, \emph{Dømmen}, \emph{Slutten}. I den uendelige Begriben, Dømmen, Slutten er Ideen paa een Gang baade den begribende, dømmende, sluttende Subjectivitet og tillige Begrebets, Dommens, Slutningens Indhold: Energien er, at Ideens Selvadskillelse beroer paa dens absolute Selveenhed, og denne igjen paa hiin.

Subjectiviteten har i sin trefoldige Væren, som tænkende, anskuende, villende, en trefoldig Existens: den absolute Subjectivitet existerer i tre Existenser. At de tre Existenser i Sagen selv ere mere end tre Modificationer af een og samme Existeren, kan sees paa Subjectivitetens trefoldige Indhold. Det trefoldige Indhold er udtrykkelig bestemt ved de tre absolute Ideer: det Sande, det Skjønne, det Gode. Den tænkende Subjectivitet existerer med det Almenes Energie i den evige Sandhed, den anskuende Subjectivitet med det Særskiltes Energie i den evige Skjønhed, den villende Subjectivitet med Individualitetens Energie i det evige Gode. Her er i Ideen tre Centra, tre Regioner: den ideale Verden (\textit{κόσμος νοητὸς}) er oprindelig adskilt i tre Riger. Men den trefoldige Existens er i Subjectivitetens og Ideernes Slutningseenhed desuagtet een i sig reflecteret Bestaaen, een energisk Væren, een eneste uendelig concret Existeren.

Den ideale Verden med sit uendelige System af Ideer og Ideeslutninger er i sin reflecterede Bestaaen en relativ-absolut Existens; den er ved Objectiveringens Energie bestemt som Objectivitetens Idee, Ideernes Idee i objectiv Form, det vil sige som \emph{Verdensidee}. Som Guddommens Ord er Verdensideen Eet med og dog forskjellig fra Guddommen selv. Den er Eet med Guddommen, thi i al Objectivering af Andet objectiverer Subjectiviteten oprindelig sig selv. Den er forskjellig fra Guddommen, ja Guddommen modsat som
% line-end hyphen in "relativ-absolut" (p. 437) kept as a compound hyphen
% --- p. 438 ---
\opage{438}Andetheden i sin absolut gjennemgaaende Relativitet kan være den oprindelige, sig selv i Alt og Alt i sig selv bestemmende, Selvhed modsat.

Som Objectivitetens Idee har Verdensideen sin selvløse Selvhed i Relativitetens Princip: Objectivitetsideen er i sit objective Væsen Andethedens Idee. Heri ligger, at den ene oprindelige Verdensidee maa have sin reflecterede Bestaaen i Systemer af relative Ideer. Ifølge Magtordets Idealitet udgjør de relative Ideers systematiske Eenhed \emph{Idealsystemet}; ifølge Magtordets Realitet udgjør den \emph{Realsystemet}. Idealsystemet afgiver nu Ideeslutningens subjective, Realsystemet dens objective Extrem. Da Ideeslutningens subjective Extrem imidlertid selv er en Slutning, en \emph{Idealitetsslutning}, saa følger heraf, at hele Realsystemet med sine Magtbestemmelser og antilogiske Momenter allerede maa, paa ideel-concret Maade, være anticiperet i den ideale Ideeslutning.\footnote{Om Ideal- og Realsystemernes oprindelige Modsætning og teleologiske Eenhed jvfr. Grundideernes Logik I. S. 216--229.}% note *)

Til alsidig Bestemmelse af Idealitetsslutningen, den ideale Ideeslutning, er det altsaa nødvendigt, at Totaliteterne i det objective Extrem maae være gjennemsigtige i det subjective: det objective Extrems reale Totaliteter ere anticiperede og ved Anticipationen medbestemmende for Ideeslutningens subjectivt absolute Indhold. Tænkes Realsystemets medbestemmende Reflexer borte, er Idealsystemet kun en tom Abstraction. Idealitetsslutningen er paa een Gang Extrem og Slutning. Som Slutning har den ved Integrationer fra Realiteten hele Verdensideen til sit aprioriske Indhold; som Extrem har den sit Indholds uendelige Rigdom i absolut reflecteret Bestaaen, i en med evig Virken identisk Væren. Det anticiperede Realsystem er ved Idealitetsslutningen bestemt paa een Gang baade som \emph{Mulighed} og som \emph{ydre Virkelig%
% --- p. 439 ---
\opage{439}hed}, som Mulighed formedelst Anticipationen, som ydre Virkelighed i Charakteer af objectivt Extrem.

Ved det objective Extrems Sammenslutning med det subjective er Verdensideen, Ideernes Verden, uendelig concret, en virkelig Verden \textit{ad intra}. Men da Anticipationen af det objective Extrem kun har Gyldighed formedelst Extremets egen objective Realitet, kan Ideernes Verden kun have Bestaaen i energisk reflecteret Existens under den ufravigelige Forudsætning, at den reale Tingenes Verden har udvortes umiddelbar og materiel Existens. I den indadgaaende, reflecterede Existens har al Vorden i alle Virksomheder, Dannelser og Bevægelser Præg af \emph{uforanderlig Væren}: Idealexistensen udtrykker det Evige. I den udadgaaende, umiddelbare og materielle Existens antager al Væren og Tilværen derimod Præg af Endelighed, Udstykning, udvortes Forandring, \emph{uophørlig Vorden}: Realexistensen udtrykker det Timelige.

Saa nødvendigt som det nu er, at en Slutnings Extremer --- langt fra at være hinanden ligegyldige --- maae som Extremer have deres Bestaaen i indbyrdes Modsætning og gjensidig Spænding, ligesaa nødvendigt er det da ogsaa, at Ideal- og Realexistensen maae hævde deres Gyldighed ved Modsætning og Spænding.

Men hvorledes bliver nu Idealexistensen at bestemme ved sin afgjørende Modsætning til Realexistensen? Som Evighedsexistens, altsaa som Præexistens! I Modsætning til Idealexistensen bliver saa den reale Existens i Timelighed og Vorden at opfatte som en afledet Existens, en Existens, i hvilken ethvert existerende Led har sin endelige Begyndelse. I Spændingen mellem Idealitetsslutningens to Extremer ligger „Tidens Fødsel“.\footnote{Om Theologernes Qvaklerier med at bestemme Tidens og Verdens Begyndelse jvfr. D. Strauß. Den christelige Troeslære. Oversat af H. Brøchner. Første Bind 1842. S. 557--71.}% note *)

% --- p. 440 ---
\opage{440}Enhver Realexistens begynder i Tiden; men Tiden selv begynder ikke. Forudsætter Realexistensen et Anlæg, saa forudsætter Anlæget igjen en præexisterende Idee. \emph{Præexistensen er deels absolut, deels relativ.}

I Modsætning til Realexistensens mangfoldige Verdener og Verdensbygninger har den evige Verdensidee \emph{en absolut Præexistens}. Men de i Virkeligheden forskjellige Verdensbygninger vilde dog umulig, saasandt hver enkelt skal udvikle sig af sit Princip efter sin Type, kunne fremgaae af den absolut præexisterende Verdensidee, dersom ikke særskilte, til særskilte Verdener svarende Ideer ved en principiel Disjunction udviklede sig af Grundideen selv: de saaledes udviklede Ideer have i Forhold til Realudviklingens særskilte Verdener ikke en absolut, men \emph{en relativ Præexistens}.

Hvorledes det oprindelig maa forholde sig med Præexistensens Relativitet og den principielle Disjunction, sees tydelig nok af Forholdet mellem \emph{Idee og Anlæg}. At en Verdensbygning udvikler sig af et Anlæg, at en Jordklode udvikler sig af et Anlæg, at en Organisme udvikler sig af et Anlæg, betyder, ifølge vort Foregaaende, at Verdensbygningen, Jordkloden og Organismen, hver paa sin Maade, virkeliggjøre en Idee. Hvad Anlæget er i objectiv Væren, det er Ideen i subjectiv Reflecteren; Ideen er altsaa Eet med og dog forskjellig fra Anlæget. Forskjelseenheden bestemmes derved, at Anlæget er den Realbegyndelsen iboende Idee, Ideen det ved Præformationen fuldendte, i Anticipationen præexisterende Anlæg. Identiteten af Anlæg og præexisterende Idee er, naar den nærmere bestemmes, et kategorisk Udtryk for Idealitetsslutningens oprindelige Midte. Det kommer i nærværende Sammenhæng særlig an paa en Tydeliggjørelse af, hvad der væsentlig ligger i denne Midte.

For det Første udtrykker da Midten, at Extremerne maae gjælde ved deres indbyrdes Modsætning, thi dersom Modsætningen ikke var, saa var jo Midten overflødig. For det
% --- p. 441 ---
\opage{441}Andet udtrykker Midten, at de indbyrdes modsatte Extremer maae være identiske; thi dersom Identiteten ikke var, saa var Midten uden Betydning. Vi ville nu begynde med at reflectere paa Modsætningen for sig og see, hvad deraf følger. Det subjective Extrem vil da ved sin Modsætning til det objective paa eensidig Maade fremstille Friheden, medens det objective fra sin Side paa en ligesaa eensidig Maade vil fremstille Nødvendigheden.

Ifølge Frihedsextremet maatte da Realsystemets principielle Totaliteter med deres Anlæg umiddelbart henføres til relativt præexisterende Ideer, og hver enkelt af disse Ideer ligefrem bestemmes ved imperative Functioner af den guddommelige Subjectivitet. Saaledes maatte, hvad f.\ Ex. Organismen og dens Typer angaaer, ikke alene det organiske Livs Fremkomst paa Jordkloden, ikke alene den typiske Grundforskjel imellem Plante- og Dyreorganisme, men tillige enhver af de typiske Former indenfor Plante- og Dyreriget enkeltviis og umiddelbart henføres til en guddommelig Villiesact. Dette er Extremet i det Subjective. Om man iøvrigt helst vil forestille sig Sagen saaledes, at den guddommelige Subjectivitet i sit Ideerige forefinder de præexisterende Typer og saa med disse for Øie danner Realexistensens forskjellige Anlæg, eller foretrække at mene, at den guddommelige Villie først frembringer Paradigmerne og derefter de tilsvarende Anlæg, er Videnskaben ganske ligegyldigt; thi Tanken svigter: Begrebet er gaaet ud, det objective Extrem er krænket.

Ifølge Nødvendighedsextremet maatte ikke blot de præexisterende Ideer bortfalde som orkesløse, deuteroskopiske Phantasier, men selv de i Realsystemet characteristiske Anlæg maatte nødvendigviis opgive deres principielle Betydning og forvandle sig til bare Muligheder; altsaa en Verden af Muligheder uden Anlæg! Det er utroligt, hvilke Anstrængelser Erfaringsvidenskaben, ja selv Philosophien, har gjort, for at komme tilrette med en saadan Verden af Muligheder uden Anlæg.
% --- p. 442 ---
\opage{442}Hvor Anlæget mangler, maa der naturligviis særegne Omstændigheder og gunstige Betingelser til, for at dække Manglen. Saaledes skal da det organiske Liv paa Jorden ingenlunde have sin Oprindelse fra et guddommeligt Anlæg, men alene fra en under særegne Omstændigheder og gunstige, særdeles gunstige, Betingelser indtrædende Mulighed. Da nu ingen Mulighed kan virkeliggjøres uden tilstrækkelige Betingelser, medens enhver tilstrækkelig betinget Mulighed altid virkeliggjøres med Nødvendighed, saa hersker Nødvendigheden. Dette er Extremet i det Objective. Hvorledes man iøvrigt vil udmale sig de særegne Omstændigheder og gunstige Betingelser, om man vil tilkjende den guddommelige Viden en contemplativ, til indicative Functioner alene indskrænket Rolle, eller ligefrem forkaste enhver Forestilling om guddommelig Viden, er ligegyldigt, ganske ligegyldigt; thi Begrebet er gaaet ud, det subjective Extrem er krænket.

Ved deres indbyrdes Modsætning ere Extremerne altsaa satte som eensidig Frihed mod eensidig Nødvendighed. Herved er Midten, den totaliserede Identitet, udtrykkelig bestemt. At Friheden i det abstract Subjective er eensidig, vil jo med andre Ord sige, at den maa have sin reale Mulighed i den objective Nødvendighed; at Nødvendigheden i det abstract Objective er eensidig, betyder, at den maa have sin væsentlige Grund i den subjective Frihed. Friheden i Idealsystemet bliver altsaa at integrere ved Nødvendigheden i Realsystemet, og omvendt: her er Midten. Realiteten tynger med Nødvendighedens Tryk paa Friheden; men ved selvbestemmende Reaction faaer Friheden Bugt med Nødvendigheden. Det er Nødvendighedens Tryk og Subjectivitetens Reaction, der give Slutningens Midte den eiendommelige Spændkraft. Paa denne Spændkraft beroer Real- og Idealsystemets teleologiske Eenhed; paa denne Eenhed beroer Vexelforholdet mellem reel og ideel Disjunction; paa dette Vexelforhold beroer Ideens guddommelige Productivitet.

% --- p. 443 ---
\opage{443}I en Realverden med Muligheder uden Anlæg er Naturen uden Productivitet. Saalidt som det „høstløse“ Hav kan siges at have Productivitet, fordi det paa Betingelse af Blæst og Storm kan frembringe høie Bølger, ligesaalidt vil en Realverden med Muligheder uden Anlæg kunne siges at have Productivitet, fordi den paa Betingelse af Stoffer med bestemte Smeltepunkter og en vexlende Temperatur kan frembringe faste, flydende, luftformige Masser. Først med det principielle Anlæg kommer Productiviteten. Vandets Overgang fra draabeflydende Tilstand til Iis kan ikke i kategorisk Forstand kaldes productiv Frembringelse; i det organiske Livs Fremkomst, i Tilblivelsen af enhver ny organisk Type lægger Naturen derimod en væsentlig Productivitet for Dagen. Hver productiv Form er bestemt ved sit eiendommelige Anlæg; hvert eiendommeligt Anlæg er igjen bestemt ved en reel Disjunction: enten Plante eller Dyr, enten Bløddyr eller Leddyr o.\ s.\ v. De productive Overgange beroe paa lutter Totalitetsadskillelser, lutter principielle Disjunctioner. I slige Disjunctioner udtaler Naturen sig productivt.

Og dog taler den productive Natur jo, egentlig talt, ikke ud af sig selv, saalidt som noget Sprog enten taler ud af sig selv eller selv taler. Nei, det er Subjectiviteten, der taler i og ved Sproget, idet den taler ud af sig selv: den guddommelige Subjectivitet taler i og ved Naturen, idet den taler af sig selv. Realsystemets Productivitet er paa alle Dannelsernes Trin en Udtalelse af Magtens Ord, et kategorisk Magtsprog i disjunctive Domme og disjunctive Slutninger: den principielle Disjunction er reel.

Men den principielle Disjunction kunde umulig være reel, dersom den ikke oprindelig var ideel; den kunde umulig foregaae i Realsystemet, dersom den ikke oprindelig foregik i Idealsystemet; den kunde umulig indgaae i Nødvendigheden, dersom den ikke oprindelig udsprang af Friheden. Var den guddommelige Selvhed ikke uendelig productiv, vilde den tynges af Andetheden og qvæles af Objectiviteten.

% --- p. 444 ---
\opage{444}Det ligger i Selvhedens Væsen, at den kun kan bestemmes af Andet, idet den bestemmer sig selv; det ligger i Magtens Ord, at det umulig kunde være reelt, dersom det ikke oprindelig var ideelt; det ligger i Disjunctionen, at den ikke kan have Realitet i et Anlæg, hvis den ikke har Idealitet i en præexisterende Idee. Paa Forholdet imellem det bestemte Anlæg og den bestemmende Idee beroer Forholdet imellem den ideelle eller indadgaaende og den reelle eller udadgaaende Productivitet. Den præexisterende Idee, Magtidealitetens særskilte Ord, er ikke orkesløs: den bærer den objective Totalitets Anlæg, idet den selv bæres af Subjectiviteten; den bevirker Anlægets reelle Udvikling i det Objective og bestemmer Ordet i Frembringelsens kategoriske Magtsprog, idet den af Subjectiviteten faaer sine til Realudviklingens antilogiske Magtbestemmelser svarende Indholdsmomenter udviklede i reflecteret Bestaaen.

Den udadgaaende Productivitet svarer til Nødvendigheden, den indadgaaende til Friheden; men i Extremernes Midte --- og Midten er her det Afgjørende --- er den ideelle Productivitet nødvendig formedelst den reelle, og denne fri formedelst hiin. Men al Productivitet udtrykker, ifølge sit Begreb, at Subjectiviteten paa een Gang baade anerkjender og magter sit Objective. I Productiviteten er Friheden nødvendig, og Nødvendigheden fri; men da den reelle Productivitet afhænger af den ideelle som det bestemte Anlæg af den bestemmende Idee, saa har Friheden i den frie Nødvendighed Nødvendigheden selv under sig.

At Jordudviklingen, geologisk seet, viser hen til et Anlæg, betyder da, logisk forstaaet, at den foregaaer i Conseqvens af en præexisterende Idee; her kan altsaa med fuld Ret tales om \emph{Jordlegemet} selv som elementær Organisme, \emph{Jordlegemets Anlæg} og \emph{Jordlegemets} Idee. Udviklingens Begyndelse er præformeret derved, at Anlæget i Totaliteten af de objective Elementer alsidig er bestemt ved den i Subjectivitetens productive Energie reflecterede Idee.

% --- p. 445 ---
\opage{445}Dersom Anlæget ikke blev baaret af Ideen, vilde det forsvinde som en reen Modsigelse; thi hvilken forunderlig Fordring er der dog ikke lagt i et Anlæg! I Anlæget skal, væsentlig seet, det Hele være færdigt fra Begyndelsen af, og dog fordrer Anlæget just, at der maa foregaae en virkelig Udvikling: i Jordens oprindelige Anlæg ligger paa een Gang alt Jordisk, og dog skal dette jordiske Alt først successivt komme frem, og forsaavidt \emph{efterhaanden} blive til under Udviklingen. Denne Modsigelse kan kun løses derved, at det samme Indhold er paa een Gang baade \emph{udadtil}, altsaa i Magtordets Realitet som Anlæg, og \emph{indadtil}, altsaa i Magtordets Idealitet som præexisterende Idee.

„Men“, kunde Nogen spørge, „opkommer nu ikke den samme Modsigelse i Ideen som i Anlæget? Hvad vinder man vel ved at flytte den modsigende Fordring fra Anlæget ind i Ideen?“ Vi maae i denne Anledning indskærpe, at Præexistensen kategorisk er bestemt som Bestaaen i energisk Reflecteren, som reflecteret Bestaaen: i reflecteret Bestaaen har Ideen alle Indholdets Momenter paa een Gang; den energiske Reflexion er i denne Bestaaen Anticipation, og i Anticipationen maa det virkeligt Tilkommende jo være væsentligt nærværende. Hvad Anticipationen er i Ideen, det er nu \emph{Driften} i Anlæget. Den udadgaaende Productivitet, den virkelige Jordudvikling, viser os kun Driften, ikke Anticipationen; og dog er al Dannelsesdrift i det Objective kun et Magtens Ord, der stadfæster Anticiperingens Energie i det Subjective.

Af Anticiperingens Energie sees da endvidere, at Totalitetsudviklingen i den udadgaaende Productivitet Punkt for Punkt maa svare til Ideeudviklingen selv i den indadgaaende. Den præexisterende Jordidee har jo vistnok fra Begyndelsen af alle Udviklingsmomenter i reflecteret Eenhed, og Realudviklingen kan forsaavidt ikke berige Ideen med Nogetsomhelst; men da Ideemomenterne selv ere anticiperede, saa maa Realudviklingen i den ydre Totalitet, i Jordlegemet selv, nød%
% --- p. 446 ---
\opage{446}vendigviis være Betingelse for en tilsvarende Idealudvikling i Jordlegemets Idee: de executive Functioner reflectere sig i de anticiperende. At Jorden paa et vist Udviklingstrin naaer, som det hedder, en vis Modenhed, vil da, logisk talt, sige, at den præexisterende Jordidee har naaet en vis bestemt Energie.

Hvorfra komme saa paa Jordklodens relative Modenhedstrin de nye Anlæg, Anlæg nemlig til Planter og Dyr? Naturligviis fra en forudgaaende Idee-Udvikling. I Forhold til Paradigmerne for det organiske Liv er Jordideen præexisterende; i Forhold til de organiske Individer ere Livsideerne præexisterende. Den frie Nødvendighed, hvormed de typisk forskjellige Livsideer udtræde af Jordideen, er i det Foregaaende allerede bestemt ved Læren om de immanente Subsumtions- og Disjunctionsslutninger og de til samme svarende Individuationer. Vi skulle endnu betragte Vexelforholdet mellem Individuationerne og Grundideens uendelige Prægnans.

Det ligger i Sagens Natur, at Verdensideen, Andethedens, Objectivitetens og Relativitetens absolute Idee, maa være uendelig prægnant. Andethedsformen er endelig; Indholdet uendeligt: i enhver af de endelig bestemte Former giver Indholdet altsaa et uendeligt Overskud. Dette Misforhold vedbliver paa alle Dannelsens Mellemtrin, indtil Andethedsformen igjen kan optages i Selvhedsformen. I den første Slutning, den for Verdensbygningen fundamentale Slutning, der ifølge Magtordets Idealitet, har Grundideen selv til sit ideelle, og ifølge Magtordets Realitet, denne bestemte virkelige Verden til sit materielle Extrem, ligger den første Mellemidee i Særskilthedens Form, Ideen af \emph{denne Art} Verdensbygning. Denne af Grundideen afledede, men for den virkelige Verdens Anlæg bestemmende Mellemidee, en Anticipationsidee, som mod sin bestemte Realverden selv har relativ Præexistens, antager nu Særskilthedens bestemte Form, altsaa en Form med bestemt Indskrænkning, bestemt Grændse; den skal da med denne
% --- p. 447 ---
\opage{447}Form være en Individuation af Grundideen, skal rumme dens hele Indhold, altsaa et uendeligt Indhold i en endelig Form: her maa blive et uendeligt Overskud, den særskilte Idee maa være uendelig prægnant.

Med Fundamentalslutningen er Productivitetsformen eiendommelig begrændset. Hvor ofte den samme Art Slutning end gjentages, hvormange endelige Verdensbygninger der end føies til de allerede givne: Productiviteten i den mechaniske Form vil bestandig lide af den Eensidighed, at Ideen, saavel efter dens ideelle som efter dens reelle Udvikling, forbliver i Andetheden og formedelst Andetheden i endelig Modsætning til Subjectiviteten.

At den sit Indhold mægtige Subjectivitet netop derved, at Andetheden drives til Extrem, magter Andetheden, beroer i Selvhedens Dyb paa Productivitetens Uudtømmelighed. Men den uudtømmelige Productivitet beroer igjen derpaa, at en høiere Form bestandig afløser en lavere, indtil Andethedsformen omsider gaaer op i Selvhedens adæqvate Form. I Realsystemet bliver enhver ny Productivitetsform bestemt ved et nyt Anlæg, altsaa i Idealsystemet ved en ny Idee. Den nye, for det nye Anlæg afgjørende Idee skabes ikke af Intet, beroer heller ikke paa et guddommeligt Indfald, men fremgaaer ved ideel Disjunction af en allerede given Grundlagsidee.

Grundlagsideen kan udtrykkelig bestemmes som den nærmest foregaaende Slutnings \textit{terminus medius}. Betragte vi Slutningen: \textit{A} (Verdensideen i sin Almindelighed) --- \textit{S} (den særskilte Verdensidee) --- \textit{E} (dette virkelige, bestemte Solsystem), da er den særskilte Verdensidee, der afgiver Slutningens Midte, netop den Grundlagsidee, af hvis Dyb Ideen for den nye Productivitet, den til Systemet svarende Planetidee fremgaaer. Den nye med Planetideen begyndende Productivitet danner nu Extremet for den nye Slutning: \textit{A}\textsubscript{2} (Planetideen i Almindelighed) --- \textit{S}\textsubscript{2} (den særskilte Planetidee,
% --- p. 448 ---
\opage{448}Jordens Idee f.\ Ex.) --- \textit{E}\textsubscript{2} (denne virkelige bestemte Jord). Ved denne Slutning gjør Productiviteten allerede et væsentligt Fremskridt, thi Andetheden begynder nu at indgaae i Selvheden: som en sig realiserende Idee er den virkelige Jord en elementær Selvhed.

At Productivitetsformen imidlertid ikke kan være afsluttet med Slutningen \textit{A}\textsubscript{2}---\textit{S}\textsubscript{2}---\textit{E}\textsubscript{2}, at \textit{S}\textsubscript{2} maa, saa at sige, bugne af Overskud, er en ligefrem Følge af den i Andetheden liggende Selvhedsdrift. For Selvhedsdriften i den nye Form maa \textit{S}\textsubscript{2} da igjen blive Grundlagsidee; med andre Ord: det nye Anlæg i Jordlegemet, Anlæget til det organiske Liv, maa forudsætte en præexisterende, af Jordideen sig udviklende Livsidee. Med Slutningen: \textit{A}\textsubscript{3} (Livsideen i jordisk Almindelighed) --- \textit{S}\textsubscript{3} (Livsideen i jordisk Særskilthed, Plante- og Dyreliv) --- \textit{E}\textsubscript{3} (den i Individet sig realiserende Livsidee), og de i dette System indeholdte Systemer vinder Selvheden omsider Seier over Andetheden, og Productiviteten naaer sin fuldendte Form.

Hvad Andetheden er i den ydre Uendelighed, det er Selvheden i den indre; heraf sees, at Productivitetsformen først kan fuldendes med den til et selvbevidst Selv udviklede Selvhed; kun Selvbevidstheden med sin indre Uendelighed kan være den uendelige Productivitets adæqvate Form: Stjernehimlen er for lille til at rumme, hvad der kan boe i en Menneskesjæl.

Hvorledes Nødvendigheden i den uendelige Productivitet forener sig med Friheden, i hvilket Forhold Videns indicative Functioner, der lade Alt beroe paa objective Magtconseqvenser, oprindelig staae til Subjectivitetens imperative Functioner, der lade Alt afhænge af Selvbestemmelsen, fremlyser nu af vort Foregaaende.

Idet den nye Idee udtræder af sit ideelle Grundlag, udtræder den ifølge en principiel Disjunction. Ved en principiel Disjunction er Livsideen fremgaaet af den præexisterende
% --- p. 449 ---
\opage{449}Jordidee, denne igjen af den præexisterende Planetidee, o.\ s.\ v. Den principielle Disjunction er i sin Energie væsentlig betinget af en tilsvarende Subsumtion. Udvortes seet, forudsætter Subsumtionen det allerede existerende Individ (dette Existesterende er en Hest, dette Existerende er en Elephant); Subsumtionen gaaer fra Existens til Kategorie, fra Individ til Idee.
% PRINTER: printed "Existe-sterende" (line-end division), sense "Existerende".

Anderledes, naar her sees paa Idealiteten selv, paa det Indvortes. Den Anticipationens Energie, hvorpaa Disjunctionen, den oprindelige Udsondring af den nye Idee, hver Gang beroer, kan jo vistnok ikke sætte de to Led af Slutningen uden at forudsætte det tredie, kan umulig sætte Grundlagsideen som det Almene (Slægten), og den nye Idee som det Særskilte (Arten), uden at forudsætte Individet som Individ; men Forudsætningen er her en Anticipation. Ved Anticipationsenergien er det altsaa ligesaavel Arten og Slægten, der blive bestemte ved Individet, som omvendt Individet, der bliver bestemt ved Art og Slægt. Fordrer nu, hvad vi ovenfor have viist, Individuationen, at den immanente Slutnings tre Led maae i Principet selv have samme Oprindelighed: saa følger med Nødvendighed heraf, at det Individ, hvori Arten skal realiseres, maa ved anticiperende Subsumtion være en for Artsideen medbestemmende Individualidee: kun i Reflexion af sine Individualideer er Artsideen Paradigme.

Ifølge Subsumtionen bliver den præexisterende Idee bestemt ved objective Hensyn, det vil sige: den bliver bestemt som Anlæg; her er, f.\ Ex. paa dette Jordklodens Udviklingstrin Mulighed til organisk Liv, altsaa skal denne objective Mulighed nu i og ved Livsideen forvandles til Anlæg. Men Livet paa Jorden er umuligt uden en levende Skikkelse; den levende Skikkelse maa være en jordisk-individuel Skikkelse, ifølge de objective Betingelser, enten denne eller denne Skikkelse. Har den ideelle Disjunction først Muligheden for flere Anlæg, saa sætter Anticipationens Energie denne, just denne
% --- p. 450 ---
\opage{450}Skikkelse i individuel Bestemthed; det forudsatte Individ kan nu subsumeres under sin Art, og Artsideen præexistere som Paradigme.

Forsaavidt det altsaa er den ved Realsystemet betingede Subsumtion, der motiverer den ideelle Disjunction, viser Subsumtionen jo vistnok hen til Nødvendigheden, den ideelle Disjunction til Friheden; men da Subsumtionen ligesaa oprindelig er motiveret ved Disjunctionen som denne ved hiin, gaaer Friheden sammen med Nødvendigheden: Friheden sat i Nødvendighedens Form er Anlæg; Nødvendigheden sat i Frihedens Form er præexisterende Idee. Eenhed af Anlæg og præexisterende Idee er en Slutningseenhed, den Ideeslutningens Eenhed, der fuldkommer Idealiteten.

\textit{ββ)} \emph{Realitetsslutning.} \emph{Ideernes Realexistens.} I Idealitetsslutningen, Ideeslutningens subjective Extrem, er Andetheden et i Selvheden gjennemsigtigt Moment; Realiteten er her en Reflex i Idealiteten; Magtens Ord forklares i Videns Form; Productiviteten \textit{ad extra} er kun til som Forudsætning for Productiviteten \textit{ad intra}. Dersom vi nu forsøgte at fastholde det subjective Extrem alene for sig, altsaa ikke som Extrem, men som en i al Evighed selvstændig Slutning, saa vilde det unegtelig være en Abstraction, og denne Abstraction just give os det Forhold i Gud, hvilket Hegel\footnote{Philosophie der Religion. Zweiter Theil. Berlin 1840. S. 248.}% note *)
{} har betegnet som „en Kjærlighedens Leg med sig selv, hvori det ikke kommer til alvorlig Andethed, til Adskillelse og Splid“. Skal det subjective Extrem have Gyldighed, maa det oprindelig have sit objective Extrem, netop det Extrem, hvori det kommer til „alvorlig Andethed, til Adskillelse og Splid“.

Dersom den speculative Dogmatik kjendte Ideeslutningen, vilde den kunne spare sig al objectiv Uleilighed med at aflede Naturens Forkrænkelighed af Djævelens Princip og Synde%
% SEAM: word continues on p. 451 as "faldet" (printed "Synde-" at the foot of p. 450); written whole here, without the following full stop; p. 451 should begin with ". Forstaaer sig, ..." after dropping "faldet"
% --- p. 451 ---
\opage{451}faldet. Forstaaer sig, Afledningen er jo rigtignok derefter. „Det er“, hedder det i den objective Behandling af „Døden og Alskabningens Forkrænkelighed“, „vanskeligt at negte, at der i Naturen ere Phænomener, der selv fremkalde Problemet om det Ondes Indtrængen i Naturen. Vi ville ikke blot paaberaabe os, at hvad der rører sig i den menneskelige Sjæls Mørke, det hvisler gjennem Skovens Løv i Slangens Bugten, det hyler gjennem Ørkenen i vilde Dyrs Blodtørst. Men vi ville isærdeleshed pege hen paa, at der i Naturen aabenbarer sig en gaadefuld Modsigelse mod Naturens egen indre Teleologie, en Modsigelse mod Naturens egen indre Hensigtsmæssighed ... Denne Forstyrrelse af Naturens egne Formaal formaae vi ikke at erkjende som grundet i en sand Nødvendighed; Naturphænomener, der i sig selv bære et forstyrrende Princip, formaae vi ikke at anerkjende som normale“.

Det tager sig jo vistnok paa Afstand høitideligt ud, naar den theologiske Forsker gaaer i Rette med den ikke theologiske Logiker, og henstiller til ham at overveie, om han „fortrøster sig til at godtgjøre, at Forløsningen af den faldne Menneskenatur ikke ogsaa omfatter Tankens Forløsning, at dette ikke er Christendommens Hensigt, fuldstændig at restituere den sande Humanitet, men at lade væsentlige Elementer --- og Tanken hører dog vel til de væsentlige Elementer --- forblive udenfor sin helliggjørende og fuldkommende Virksomhed“. Ja, i et saadant Sprog lyder Talen om Tankens Forløsning endogsaa opbyggeligt; men er det dog ikke besynderligt? Hver Gang den „forløste Tanke“ virkelig skal til at løse et Problem, hvor gaaer det da yderlig smaat med den videnskabelige Forløsning! Naar en Jacob Bøhme, uagtet han anerkjendte, „at et Nei udfordredes, for at Ja'et kunde aabenbare sig selv, og at Modsætningen var nødvendig for at Livet kunde fornemmes“, alligevel ikke troede at kunne forklare sig „den nuværende Naturs forunderlige Regimente, Materiens Haandgribelighed og
% --- p. 452 ---
\opage{452}Dunkelhed, alle Elementers og Kreaturers Kamp og Strid mod hverandre“, uden den Forudsætning, „at Lucifers Hovmod havde stukket den oprindelige Skabnings rene ætheriske Væsen i Brand, hvis Slagger danne den nuværende Verden“:\footnote{See D. F. Strauß. Den christelige Troeslære. Oversat af H. Bröchner. Kbhvn. 1843. S.~13.}% note *)
{} da er det Videnskabens Opgave i en saadan Mystik at adskille, hvad der tilhører Phantasien og hvad der tilhører Tanken. Men paa en saadan Adskillelse skal den, efter Dogmatikens Anviisning, forløste Tanke vist neppe indlade sig. Men hvor er saa Tegnet paa Tankens Forløsning? Thi at Dogmatikeren ikke formaaer at erkjende den naturlige Forstyrrelse af Naturens egne Formaal som „grundet i en sand Nødvendighed“, kan dog umulig opstilles som et Beviis for Tankens Forløsning.

Betragte vi den virkelige, for vore Sandser nærværende Natur som Ideeslutningens objective Extrem, et Extrem, der selv er en Slutning, ja, som viist er i vort Foregaaende, et System af Slutninger; lægge vi, hvad denne Realitetsslutning med dens Systemer angaaer, behørig Vægt saavel paa \emph{den naturlige Andethed} som paa \emph{den naturlige Selvhed}, turde det vel lykkes at erkjende den naturlige Forstyrrelse af Naturens egne Formaal som begrundet i sand Nødvendighed, og tyde den æsthetisk-ethiske Parallel imellem det, „der rører sig i den menneskelige Sjæls Mørke“, og det, der ængster os i Slangens Hvislen og i Rovdyrets Hylen.

Hvad da for det Første Andetheden angaaer, da kan dens elementære Charakteer kun komme i Strid med Selvheden og betegne Forkrænkeligheden paa det Punkt, hvor Livet begynder, altsaa først i Planteorganismen. Betragtes Jordkloden for sig, ville selv de største Omvæltninger ikke betegne nogen Forstyrrelse: Orkaner og Jordskjælv ere for den elementære Selvhed ligesaa adæqvate som Solskin og Blikstille. I Forhold til Organismen derimod begynder Spændingen mellem
% --- p. 453 ---
\opage{453}Princip og Betingelser, mellem Idee og Virkelighed. „I Plantens naturlige Skikkelse er Forskjelseenheden af det Hele og dets Dele, det Indre og Ydre, Totalindividet og dets Partialindivider ... alsidig bestemt ved den articulerende Modsætning af Selvhed og Andethed ... Spørge vi, hvilken af de to Magter der er den stærkeste, Selvheden eller Andetheden, vil en Dialektik af Magt og Afmagt vise sig paa begge Sider. Væsentlig forstaaet er Selvheden, der jo principielt bærer Andetheden i sig, det Oprindelige, det Høiere, og forsaavidt det Mægtigere; Andetheden, der har sit Integrerende udenfor sig, er derimod det Afledede, det Lavere, og forsaavidt det Afmægtige, det Svagere. Men falder nu, hvad Virkeligheden jo ligeledes fordrer, Eftertrykket paa Phænomenet, det Factiske, den i Kjendsgjerningen slaaende Umiddelbarhed, saa vender Forholdet sig om, saa maa Andetheden, der jo netop har sin Styrke i det udvortes Umiddelbare, nødvendigviis vise sig som det Oprindelige, det Stærkere, og det ikke blot \emph{uagtet} den er, men \emph{just fordi} den er det efter sit Væsen Afledede, det Svagere; ligeledes maa Selvheden, hvis Gyldighed beroer paa Reflexionsenergie, nødvendigviis fremstille sig som det Afledede, det Abstraherede, og det ikke blot \emph{uagtet} den er, men just \emph{fordi} den er det i Grunden selv Oprindelige, det i objectiv Reflexion Evige og Uforkrænkelige“.\footnote{Forelæsninger over „Phil. Propæd.“ fra Universitetsaaret 1860--61. S.~317--18.}% note *)

Forkrænkeligheden, der paa mange Maader lader sig tilsyne i Planteverdenen, potenseres i Dyreverdenen til Lidelse, Smerte og Død. Dette skeer ikke ved nogen Indblanding af Lucifer, men af „en sand Nødvendighed“. Det besjælede Individ er i sin umiddelbare Selvfornemmelse sit eget Formaal; for dette Formaal maa den virkelige Omverden afgive Midler og Betingelser, altsaa bestemmes som det Underordnede. Men Omverdenen, der selv underligger den elementære
% --- p. 454 ---
\opage{454}Andethed, staaer Selvet imod, og saa begynder Conflicten. „Hver Dyreart er ligesom hver Planteart bestemt til at leve paa sin Maade. Men det Besjælede lever paa høiere Vilkaar end det Ubesjælede; det Ubesjælede lever kun, forsaavidt det i al Umiddelbarhed er et Levende; men det Besjælede er et Levende, der udtrykkelig lever sit eget Liv. At leve for sig sit eget Liv er at \emph{prøve} Livet; at prøve Livet er i umiddelbar Forstand: at \emph{nyde} og at \emph{lide} ... Maalet er ikke Livsnydelse for sig, saalidt som det er Lidelse; thi det Samme, der gjør, at Individet kan nyde, gjør ogsaa, at det kan lide; den virkelige Nydelse forudsætter Muligheden af en virkelig Lidelse, ligesom omvendt den virkelige Lidelse forudsætter Muligheden af en tilsvarende Nydelse“.\footnote{Anfr. Skr. S.~392.}% note *)
% PRINTER: the footnote *) "Anfr. Skr. S. 392." has no anchor mark in the body of p. 454; placed at the end of the only quotation on the page.

Allerede den hermed charakteriserede Kamp mellem Selvheden og den elementære Andethed vil være tilstrækkelig til at forklare, hvorledes det gaaer til, at der i Naturen kan aabenbare sig „en gaadefuld Modsigelse mod Naturens egen indre Hensigtsmæssighed.“ „Men“, vedbliver Dogmatikeren med „den forløste Tanke“: „Naturphænomener, der i sig selv bære et forstyrrende Princip, formaae vi ikke at erkjende som normale ... En Theodicee, som her vilde berolige os med den almindelige Forestilling om Endeligheden og Livets nødvendige Udvikling gjennem Modsætninger, vilde kun omgaae Problemet. Thi baade den sande og den falske Udvikling bevæger sig igjennem Modsætninger, og Spørgsmaalet er netop, om alle de Modsætninger, som Naturlivet viser os, ere normale“. Vi svare „den forløste Tanke“, at Ideeslutningens objective Extrem naturligviis er og bliver normalt som Extrem; men at det er en Modsigelse, ja en grov Modsigelse, at ville forlange af et Extrem, hvad der først kan naaes i en Midte. Lad „den forløste Tanke“ engang for Alvor gjøre et videnskabeligt Forsøg med at kaste Realextremet bort, og den skal
% --- p. 455 ---
\opage{455}nok see, hvor den, trods alle Declamationer om det Normale og det Abnorme, den sande og den falske Modsætning, tilsidst kommer hen.

Den i Dyreverdenen voxende Strid mellem Selvhed og Andethed naaer i Menneskelivet sit Maximum. Idet Bevidstheden, som i Dyret kun er umiddelbar Sandsebevidsthed, i Mennesket fordobler sig til Selvbevidsthed, bliver ved denne Fordobling ikke blot Glæden og Nydelsen, men ogsaa Smerten og Lidelsen uendelig fordoblet: Døden, der kun rammer Dyret, bliver til for Mennesket. Med Menneskets opvaagnende Selvbevidsthed vaagner Aanden, Villiens Idee med Frihedens Fordring i Frihedens Form. Den naturlige Selvhed er Egoisme; den gjennem hele Naturen gaaende Strid mellem Selvhed og Andethed forvandler sig da i Mennesket til en Frihedens Kamp imellem Egoismen og Aanden. Egoismen er for Aanden en Brøde; Anelsen om Brøde er Angst. Seet fra dette Synspunkt maa det vel for Enhver, hvis Tanke ikke er saaledes „forløst“, at den er blevet ganske løs, uden Lucifers Hjælp være indlysende, hvorledes det kan gaae til, at Noget, der hvisler i Løvet, og Noget, der hyler i Ørkenen, kan ængstende minde om Noget, „der rører sig i den menneskelige Sjæls Mørke“.

Det er det Tiltrækkende ved den oprindelige Mystik, at den overalt seer Naturen i Aandens og Aanden i Naturens Lys; men naar Forstandsreflexionen begynder, maa Tvetydigheden bort. Hvad der som i hellig Dæmring lyser med Primitivitetens Metalglands, kan hverken erstattes eller gjengives ved en overfladisk „Tankeforløsnings“ æsthetiske Politur. At en Jacob Böhme ved et af Aandens stærke Lyn kan mene at see Alskabningens Forkrænkelighed forklaret i det røde Skjær af Lucifers Flammer, kunne vi forstaae; vi kunne uvilkaarlig gribes af det Dybsindige. Men hvor mat og tvetydig bliver ikke Forklaringen, naar den „forløste Tanke“ ved objective Grunde paa objectiv Maade for Alvor paatager sig at føre
% --- p. 456 ---
\opage{456}Beviis for, at Verdensegoismen er det kosmiske Princip, og det kosmiske Princip den vordende Djævel.

Aabenbaring og Videnskab have med al deres Modsætning dog det tilfælles, at de begge forkaste det Tvetydige: en tvetydig Aabenbaring er en ligesaa stor \textit{contradictio in adjecto} som en tvetydig Videnskab. For at hæve Tvetydigheden, maa man fra begge Sider gjøre Skilsmisse mellem Natur og Aand. Men naar Skilsmissen er skeet, hvor dyb er da ikke den Kløft, der adskiller Troen fra Viden! Thi det staaer fast, naar man blot ikke vil lade Videnskabens „Forløsning“ gøgle med Troen, at der har været Smerte og Død til i Millioner Aar, før Mennesket fremtraadte paa Jorden. Og det staaer fast, naar man blot ikke vil lade Troens „Forløsning“ gøgle med Viden, at Døden er Syndens Sold, og Naturens Forkrænkelighed en Aandens Brøde.

At Realitetsslutningen er et Extrem, har bedaaret den troende Speculation, som i Conseqvenser af Andethedens umiddelbare Magt har meent at spore Djævelens Gjerninger; at Realitetsextremet selv er en i Slutningssystemer bestaaende Slutning, har bedaaret Immanensphilosophien, som ved at lade det Ideelle gaae op i det Reelle har taget et Slutningsextrem for den fuldendte Slutning.

\textit{γγ)} \emph{Afslutning: Idee-Existensens Fuldendelse.} Med Idealitets- og Realitetsslutningens fuldstændig udviklede Midte og begge Extremers ved Midten bestemte Sammenslutning naaer Ideeslutningen selv omsider sin Afslutning. Den herved antydede Succession er imidlertid begrundet i en tilsvarende Simultaneitet; Ideebevægelsen er Hvile, Hvilen selv en Bevægelse. Thi det forholder sig ikke saaledes, at den fuldkomne Midte nu skulde afskaffe de to ufuldkomne Extremer, den Idealitetens og Realitetens totaliserede Eensidighed, hvorpaa Dualismen oprindelig hviler; nei tvertimod! som sand og fuldkommen Midte maa den just lade Extremerne i al deres Eensidighed staae ved Magt. En anden
% --- p. 457 ---
\opage{457}Sag er det, at Midten, der selv er en Slutning, et System af Slutninger, ja et System af Slutningssystemer, igjen kan udfolde Extremer, som i deres Fuldkommenhed selv ere Mellemled, idet de fuldkomne Mellemled selv ere Extremer.

Det er \emph{Andetheden} (Grændsen, Realnegationen, Differensen), som herved fremstiller sig i et nyt Lys. I Idealsystemet er Andetheden ideel, og Negationen immanent; Differenserne dannes af de Vexelbestemmelser, som Ordet i den ontologiske Subjectivitets uendelige Tænken, Anskuen og Villen oprindelig medfører. I Realsystemet er Andetheden reel, ved Magtens umiddelbare Udslag endog sandselig og materiel. Differenserne tabe sig i overfladiske Forskjelligheder, timelige Udstykninger, udvortes og tilfældige Relationer.

Dersom det nu, hvad der i „Begrebets Philosophie“ dog saa ofte er paastaaet, virkelig var umuligt, at den reale Andethed kunde hæves til et høiere Trin end det, der charakteriserer denne for os nærværende Virkelighed, saa kunde Idealerne rigtignok, som det hedder, aldrig naaes, saa maatte Idee og Existens befinde sig i evig Splid. Men en Forvandling af Andetheden og dermed en Ophævelse af Forkrænkeligheden er, logisk seet, saa langt fra at være en Umulighed, at den tvertimod er en uafviselig Nødvendighed. Den Splid imellem Ideal og Virkelighed, der er saa betegnende for det objective Realextrem, er ingenlunde et Tegn paa Ideens relative Afmagt; den er meget mere et Beviis paa den idee-opfyldte Subjectivitets absolut overgribende Magt.

En mat Idee har naturligviis matte Extremer med svag Midte og mat Sammenslutning; derfor er den immanente Speculation altid saa ængstelig for hvad der kunde ligne „en Dualisme“. Den mægtige, i Væsen guddommelige Idee har derimod mægtige Extremer; dens Udslag i kraftig Dualisme er kun Forudsætning for den uendelige Overlegenhed, hvormed den fuldfører Slutningen; Idee-Existensens Fuldendelse i concret Afslutning maa nødvendig have en altforklarende
% --- p. 458 ---
\opage{458}Midte, en sidste \textit{terminus medius}, der magter alle Extremer. Da nu den altforklarende Midte ligesaavel maa være et Ideerige som ethvert af Extremerne, saa har Ideeslutningens tre \textit{termini} altsaa med indre Conseqvens forvandlet sig til tre særskilte Verdener: en Idealverden, en Realverden og en ideal Realverden.

„Men“ --- spørger den ædruelige Kritik --- „hvor komme vi nu hen? Skal Logiken virkelig give sig af med at belyse den ideale Realverden, hiint overnaturlige Rige, hvor Idealerne kun ere virkelige, fordi Virkeligheden selv er ideal, saa maa Slutningssystemernes Afslutning jo ende med en formelig Eschatologie. Men saa gjøres der jo dog Indgreb i Theologiens velerhvervede Rettigheder; saa kommer Logikeren jo dog tilsidst til at anstille saadanne Betragtninger i Oververdenen, som kun en speculativ Dogmatiker, en mystisk dybsindig Genius, der i et hiinsidigt „Klosterliv“ havde viet sig ene og alene til „eschatologiske Betragtninger“, vilde være istand til at fuldføre i Underverdenen!“

Vi kunne i denne Henseende strax berolige Kritiken ved at henvise til Problemet: „To Verdener“.\footnote{Forelæsninger over „Phil. Propæd.“ fra Universitetsaar 1861--62. S.~123--137.}% note *)
{} Discussionen af dette Problem ender just med et aldeles negativt Resultat. „Fastholde vi“, siges der, „Ligheden mellem Her og Hisset, saa overraskes vi af Uligheden; fastholde vi Uligheden, saa paatrænger sig igjen Ligheden. Holde vi os til Virkeligheden, saa kommer den i Strid med Idealet; skal Idealet selv være virkeligt, saa faae vi to Slags Idealer; Idealer, som vi kjende, men som aldrig blive virkelige; Idealer, som vi ikke kjende, og om hvis Virkelighed vi ikke kunne danne os nogen Forestilling, endsige noget klart Begreb. Men naar jeg ikke er istand til at indsee et Ideals Mulighed, hvorledes skulde
% --- p. 459 ---
\opage{459}jeg da kunne antage dets Virkelighed? \textit{A non posse esse ad non esse valet consequentia}“.

Hvad det dog er for et i Videnskaben misligt Forehavende, at skulle ved Analogier fra denne nærværende Verden udmale sig Tilstanden i det tilkommende Liv, kan sees af mange Exempler; navnlig ere Theologernes Forsøg med den objective Opstandelseslære lærerige, ogsaa i logisk Henseende. Den i Dogmatiken objectivt „forløste“ Tanke har nemlig i Forening med den ligesaa objectivt „forløste“ Phantasie taget mange og forskjellige og charakteristiske Tilløb til en videnskabelig Forklaring af den i Opstandelsen forklarede Natur. Forsøgene, der ere udførte snart med scholastisk Dristighed, snart med en moderne Bevidstheds afglattende Forsigtighed, snart uæsthetisk og smagløst, snart igjen med Takt og Smag, bevæge sig alle paa Grundlag af een og samme Modsigelse, Modsigelsen af \emph{en naturløs Natur}.

At den philosophiske Kritik, hver Gang den har givet sig af med at belyse det Fade, Flade og Selvmodsigende i Dogmatikernes Forsøg, har kunnet, ikke alene præstere \textit{præstanda}, men feire Triumpher, er forstaaeligt. Men det bør da heller ikke forties, at Kritiken, blendet af adskillige Godtkjøbstriumpher, ikke sjeldent har kritiseret sig saa langt ind i endelige Forstandsmodsigelser, at den \textit{eo ipso} har kritiseret sig ud af Ideen. Ved at bevise Umuligheden af at anskueliggjøre Tilstanden i den tilkommende Verden er den negative Kritik, ifølge en let forklarlig Selvskuffelse, bragt til at antage, at den med det Samme beviste en tilkommende Verdens Umulighed. Men dette er en Overilelse.

Sættes den tilkommende Verden identisk med den ideale Realverden og denne identisk med Ideeslutningens Midte: da kan der ikke længere reises Tvivl om dens Gyldighed; dens Nødvendighed er given. Forsaavidt det da gjælder, ikke om at gjøre dens Virkelighed anskuelig i Form af Tilstand --- thi dette er, som sagt, umuligt --- men om at be%
% --- p. 460 ---
\opage{460}stemme dens Begreb som nødvendig Form for Ideeslutningens Fuldendelse, kan Reflexionen i dette Tilfælde netop støtte sig til den formelle Logiks bekjendte Sætning: \textit{a necesse esse ad posse esse valet consequentia}.

Der kan ligesaa lidt være Tale om noget Spring eller udvortes Overgang fra denne til hiin Verden, som der kan være Tale om Spring eller udvortes Overgang fra en Slutnings Extrem til dens Midte. Vel er der en Sphæreforskjel mellem Her og Hisset, men saaledes er der jo ogsaa en gjennemgaaende Modsætning og qvalitativ Forskjel mellem Midte og Extrem. Just fordi Modsætningen af Extrem og Midte i Ideeslutningen naaer sit Høieste, er Identiteten desto stærkere. Det kan paa alle Hovedpunkter i Extremet sees, at Eensidigheden, om end paa eensidig Maade, integreres fra Midten: gjennem Midten gaaer en continueret Bevægelse fra Extrem til Extrem.

Betragte vi nu de tre Regioner, Ideal- og Realsystemet som Extremer og det ideale Realsystem som Midte, da er det jo fra Grunden af Ideen, den ene og samme Idee, der med alle særskilte Ideer gaaer igjennem alle tre Regioner. I Idealsystemet har Ideen som Idee sin Præexistens, i Realsystemet sin umiddelbare Realexistens; men enhver af disse Existensformer er eensidig og i sin Eensidighed utilfredsstillende. Præexistensen mangler Realitet, Realexistensen mangler Idealitet; dersom disse Mangler ikke kunde afhjælpes, maatte Ideen forblive i evig Strid med sig selv: Strid og Modsigelse vilde da i meningsløs Mening være Tilværelsens Princip (πόλεμος πάντων πατήρ).

Men i ethvert af Extremerne er Midten selv præformeret. Den ideelle Disjunction, hvorved den særskilte Idee erholder sin Præexistens, gaaer netop, som ovenfor er viist, i Eet med den reelle Disjunction, hvorved Realexistensen erholder sit Anlæg. Det Reelle virker oprindelig paa det Ideelle, for at det Ideelle kan virke paa det Reelle; der er i Ideen Intet
% --- p. 461 ---
\opage{461}saa ideelt, at der i dets Energie jo maa være et Moment, en Reflex af reel Umiddelbarhed, og der er i Ideen Intet saa reelt, at det jo har en Idealitet i sin Mulighed. Ved den Fordobling, der saaledes er i ethvert af Extremerne, kan Midten faae Magt; thi hvad er nu Midten? Ideens alsidig bestemte Existens, den Existensens Fuldendelse, der alene kan give Extremerne deres Betydning. Midten er Maal, Maalet er Afslutning, Afslutningen er alle Slutningers Sandhed; fra dette Synspunkt kan det da sees, at Ideeslutningens Functioner væsentlig ere Modalfunctioner af Sandhed.

\emph{Midten er Sandheden.} Vistnok have Extremerne ogsaa deres Sandhed, men kun i Forening med Midten. Tænkes Midten borte, saa er Andetheden paa begge Sider uadæqvat, og Tankeslutningen E--A--S (som Bevægelse A--(E, S)--A) ikke istand til at fyldestgjøre Ideen. I Idealextremet er Andetheden abstract ideel, Concretionen i E og S mangler Realitet; men at mangle Realitet er at mangle Sandhed. I Realsystemet er Andetheden eensidig elementær, Concretionen i E og S mangler Idealitet; men at mangle Idealitet er at mangle Sandhed. I Midten derimod er Andetheden adæqvat, og Sandhedsideen fyldestgjort.

\emph{Midten er Skjønheden}; thi Skjønheden er det Ideelles og Reelles anskuelige Harmonie. I Idealsystemet er, hvad Anskuelsesslutningen E--S--A (som Bevægelse S--(E, A)--S) angaaer, E med al sin Gjennemsigtighed saa mat, at S netop derved kommer til at mangle den slaaende Umiddelbarhed, som det Skjønne aldrig tør savne. I Realextremet har E Umiddelbarhed nok; men da Andetheden er elementær og ved sin elementære Umiddelbarhed tilfældig, er det i Brud med A og uigjennemsigtigt for S. Af det anskuelige Misforhold mellem E, S og A fremgaaer Uskjønheden. I Midten derimod er Andetheden adæqvat; E har Gjennemsigtighed nok til at anskueliggjøre A og Umiddelbarhed nok
% --- p. 462 ---
\opage{462}til individualiserende at optage Skikkelsen S: Skjønhedsideen er fyldestgjort.

\emph{Midten er det Gode}. I Idealextremet savner Villien sin nødvendige Skranke, den virkelige Modsætning, ved hvis Ophævelse den beviser sin Energie. Villiesslutningen S--E--A (som Bevægelse E--(S, A)--E) er tom, thi E er en Skygge, Andetheden er kun en Reflex af Selvheden. I Realsystemet er Andetheden stærk nok til at gjøre Villien Modstand; men nu er Modstanden saa elementær, at Villien og Skranken befinde sig i uophørlig Brydning. I Eenhed med Midten ere derimod begge Extremerne gode; den alsidige Energie, hvormed Villien virker i og ved Ligevægten af den ideelle og reelle Andethed, der charakteriserer Midten, er netop betinget af den Eensidighed, hvormed den virker i Extremerne.

Hvad der gjælder om Idealideerne, det Sande, det Skjønne, det Gode, i al Almindelighed, det gjælder naturligviis i særegen Forstand om Realideerne, Grundlagsideerne og de fra samme opadstigende Selvhedsideer. For Individualexistensen af disse Ideer faaer den elementære Andethed en dobbelt Betydning, først, i Almindelighed, ved det Realextrem, hvortil den hører; dernæst, i Særdeleshed, ved det Substrat den afgiver for Individualideens egen reale Udvikling. Saaledes har f. Ex. Verdensbygningens Idee ikke blot en i det subjective Extrem reflecteret Bestaaen, en under Idealsystemet hørende Præexistens, men i Realextremet, det objective Extrem, har den jo med det Samme en materiel og umiddelbar Bestaaen som denne bestemte Verdensbygning (dette Solsystem). Hvad betyder nu Midten? Den til alsidig Existens udviklede og forklarede Verdensbygning! Her gjælder det om at være opmærksom paa Realextremets dobbeltsidige Rolle.

Opfattes det reale Extrem individuelt i individuel Modsætning til det ideale, altsaa denne existerende Verdensbygning i individuel Modsætning til sin præexisterende Idee, da er Midtens Betydning umiskjendelig. Dersom Midten ikke var,
% --- p. 463 ---
\opage{463}vilde Ideen i sin Præexistens mangle Realitet, og den umiddelbare Realexistens mangle Idealitet; imellem slige Extremer kunde der da ikke tænkes noget Vexelforhold: begge Extremer vilde være uden Sandhed. At Midten er Extremernes Sandhed, vil altsaa sige, at Verdensbygningen kun kan have en umiddelbar og materiel Existens under Forudsætning af, at den med det Samme har en i sin Afslutning reflecteret og ved Reflexionen idealiseret Realexistens. Der kan altsaa, reentud sagt, kun existere en sandselig Verdensbygning under Forudsætning af, at en oversandselig med det Samme kommer til Existens.

Men som præexisterende Idee har Verdensbygningen en Begyndelse i Evigheden, som Realexistens har den en Begyndelse i Tiden; Idealextremet har Evighedens, Realextremet Timelighedens Præg: hvorledes forholder det Evige sig i Afslutningen, Slutningens Midte, til det Timelige? Dersom Begyndelsen i Evigheden ikke oprindelig reflecterede en Tidsbestemmelse, vilde her ikke kunne være Tale om relativ Præexistens. Den relative Præexistens angiver jo netop en Ideens Vorden af sit ideelle Grundlag, en ved Disjunctionens \textit{punctum saliens} betegnet Tilbliven, en i Principet stigende Reflexionsudvikling, altsaa en Succession indenfor det Evige. Omvendt, vilde den i Tiden begyndende Realexistens umulig kunne begynde som virkeligt Anlæg, dersom der i den timelige Begyndelse ikke var noget Evigt. I Midten selv, den ideale Realexistens, vil det Timelige altsaa paa uendelig Maade omsætte sig i det Evige, og dette i hiint.

Men paa dette Punkt indtræder en Tvetydighed. At Midten bestaaer saalænge Extremerne bestaae, er en Selvfølge; men dersom nu et af Extremerne forgaaer, hvad bliver der saa af Midten? At en oversandselig Verden existerer, saalænge denne materielle Verdensbygning existerer, maa, saa%
% --- p. 464 ---
\opage{464}sandt her da ved det Oversandselige ikke tænkes paa en ny udvortes og umiddelbar, men paa en væsentlig, indvortes reflecteret Existens, jo vistnok indrømmes; men Alt, hvad der opstaaer i Tiden, maa jo forgaae med Tiden; naar da denne sandselige Verdensbygning engang forgaaer: hvad bliver der saa af den oversandselige? Til Besvarelse af slige Spørgsmaal er det nødvendigt at stille sig det for Ideens Prægnans og den uendelige Productivitet afgjørende Forhold mellem Selvhed og Andethed klart for Øie.

I alle Regioner, under alle Skikkelser har \emph{Andetheden den dobbeltsidige Betydning: at være Selvheden paa een Gang baade coordineret og subordineret}. At Andetheden er Selvheden \emph{coordineret}, at der ikke kunde være nogen Subjectivitet, dersom det Objective manglede; at det vidende Selv hverken i Tænken, Anskuen eller Villen kunde have nogen Energie, dersom det Magtens Udslag, der giver Andetheden sin Realitet, tænktes borte, dette er Altsammen oplyst i vort Foregaaende. Men at Andetheden med alt dette ikke desto mindre er og maa være Selvheden \emph{subordineret}; at Subjectiviteten, som igjennem sin reale Modsætning altid forholder sig til sig selv, just udtrykker sin Anerkjendelse af det Objective derved, at den i uendelig Productivitet paatrykker det sit Segl, og saaledes besegler dets Objectivitet, maa da ogsaa antages at være tilstrækkelig oplyst.

Den i Andetheden skjulte Modsigelse, at den paa een Gang er Selvheden baade coordineret og subordineret, bliver kun løst ved den dobbelte Rolle, som Andethedselementerne øiensynlig spille i Realextremet.

At et Realextrem i al Almindelighed er nødvendigt for at danne Modsætning til det ideale Extrem, vil sige, at Idealverdenen maatte forsvinde, dersom denne virkelige, materielle Verdensorden i det Hele kunde forsvinde; i denne
% --- p. 465 ---
\opage{465}Henseende staaer det umiddelbart Reelle lige med det Ideelle: Andetheden er Selvheden coordineret. Men hvad den individuelle, til den bestemte Idee svarende Realexistens angaaer, er den elementære Andethed, trods al sin tilsyneladende Selvstændighed, dog Selvheden og Subjectiviteten væsentlig underordnet.

Enhver Idee, selv den i Andetheden allerdybest nedsænkede Grundlagsidee, er dog efter sit oprindelige Væsen en Selvhed. Der kunde altsaa ingen Verdensidee komme til virkelig Existens i den virkelige Verdensbygning, dersom Ideen under sin reale Selvudvikling ikke var istand til at magte Andethedens Elementer.

I Realextremet kunne Andethedens Elementer aldrig --- det strider udtrykkelig mod Realextremets Charakteer --- forvandles til reale Ideemomenter. Men uden at udvikle et iboende System af reale Ideemomenter kunde Ideen umulig udvikle sig som Idee. Da nu de reale Ideemomenter, ihvorvel de fra Begyndelsen af maa articuleres ved Andethedens materielle Elementer, ingenlunde ere Eet med disse: saa kunne Andethedselementerne igjen hjemfalde til det almene Realsystem, og den existerende Idee endda bevare Udviklingens ideel-reelle Resultat og med dette Resultat gaae op i den afsluttende Midte. I denne Midte er den reelle Andethed, selv om dens Realitet iøvrigt er nok saa stor, selv om den i det nye System paany skulde kunne løslades til elementære Extremer --- ja, om det saa skulde tordne midt i „Evighedens Sommer“ --- dog paa alle Maader saa igaaet med det Ideelle, Selvheden i alle Henseender saa absolut underlagt, at Idee-Existensen med fuld Ret maa siges at være fyldestgjort: den oversandselige Verden er uforkrænkelig.

\begin{center}---\end{center}
% END OF VOLUME: p. 465 closes with a short printed rule. PDF 484 carries only a
% handwritten library shelfmark, PDF 485 is blank, PDF 486 is the marbled endpaper;
% no printed end mark, colophon or advertisement follows p. 465.

\end{document}
